% 本文件由 LaTeX 排版 Agent 根据有效 translation.json 自动生成。
% author=Clement of Alexandria; work_id=0210; title=杂篇

\chapter{杂篇}

% structure: p0001 source=h1 target=omitted reason=page-title-already-rendered-by-parent

卷第一\newline{}卷第二\newline{}卷第三\newline{}卷第四\newline{}卷第五\newline{}卷第六\newline{}卷第七\newline{}卷第八\par

\SourceRecord{Translated by William Wilson.；Ante-Nicene Fathers；Vol. 2.；Edited by Alexander Roberts, James Donaldson, and A. Cleveland Coxe.；Buffalo, NY: Christian Literature Publishing Co.,；1885.；<http://www.newadvent.org/fathers/0210.htm>.}

% structure: p0001 source=h1 target=section reason=page-title

\section{\WorkTitle{杂集（第一卷）}}

% structure: p0002 source=h2 target=subsection reason=relative-heading-rank; base=h2

\subsection{第一章 序言——作者的目的——文字作品之益处。}

\EditorialNote{开头缺失}……使你们能随手阅读，并加以保存。到底是否完全不应当留下文字作品？若应留下，该由谁来写？若是前者，何需文字作品？若是后者，写作之事应交给热心之人，还是相反？若有人指责热心之人的写作，却赞同那些非热心之人从事写作，未免荒谬。\Person{特奥庞普斯}、\Person{蒂迈欧}编造寓言与诽谤，无神论的领袖\Person{伊壁鸠鲁}、\Person{希波那克斯}、\Person{阿尔基洛科斯}，都可按他们那可耻的方式写作；然而宣扬真理的人却被阻止留下有益后世的著作？我想，为后世留下好孩子是一件好事——这是我们身体的孩子；但言语是灵魂的孩子。因此我们称教导我们的人为父亲。智慧是具传播性且爱人的。所以\Person{撒罗满}说：\BibleQuote{我儿，你若接受我诫命的言语，将它藏在你心里，你的耳朵必听智慧。}\BibleRef{箴 2:1} 他指出，所播种的言语被藏在学习者的灵魂里，如同在土里，这是属灵的栽种。因此他又加上：\Quote{你当倾心于悟性，并为你儿子的劝戒而应用它。} 因为我认为，灵魂与灵魂结合，精神与精神结合，在言语的播种中，将使所播种的生长发芽。凡受教者，在顺服方面，是教导者的儿子。\Quote{我儿，}他说，\Quote{不要忘记我的律法。}\BibleRef{箴 3:1}\par

倘若知识不属于众人（如俗语所说：对牛弹琴），文字作品却是为多数人而写的。\Quote{例如，猪喜爱污秽胜于清水。} 因此主说：\Quote{因此，}\BibleQuote{我对他们讲话，用比喻，是因为他们看却看不见，听却听不见，也不明白。}\BibleRef{玛 13:13} 并非主使人无知（如此想是亵渎的），而是祂预言性地揭露了存在于他们当中的无知，并暗示他们不会明白所说的话。如今救主显示祂自己，从祂的丰富中，按接受者的能力分施恩赐给祂的仆人，好使他们藉着活动而增加，而后回来与他们算账；那时，祂称赞那些赚了钱的人，那些在小事上忠信的人，命令他们管理许多事，并叫他们进入主的喜乐。但对于那将托付给他生息的银子埋藏起来，如数归还而无增加的人，祂说：\Quote{你这邪恶懒惰的仆人，你应当把我的银子交给银行家，到我来时，我连本带利收回。} 因此那无用的仆人\Quote{将被丢到外面的黑暗里。} \Quote{所以你当坚强，}保禄说，\BibleQuote{在基督耶稣内的恩宠中。你从我许多见证人中所听见的，要传给忠信的人，他们也能教导别人。}\BibleRef{弟后 2:1} 又说：\Quote{你当竭力在天主面前蒙悦纳，作无愧的工人，正确地分解真理的话。}\par

那么，若两者都宣讲圣言——一个藉着写作，一个藉着言语——岂不都该被认可，因为他们以爱德使信德活跃？是由于自己的过错，人不去选择最好的；天主是无可责备的。就当前问题而言，有些人的职责是将话语放贷生息，另一些人的职责是考验它，或选择或不选择。判断在他们自己心中。但有一种知识是传报者的特征，另一种可说是信使的特征，无论以何种方式运作——藉着手或舌头——它都是有用的。\BibleQuote{因为那为圣神而播种的，将由圣神收获永生。我们行善不要疲倦。}\BibleRef{迦 6:8} 那藉着天主眷顾而遇到它的人，领受极高的益处——信德的开端、接受正确生活方式的准备、对真理的推动、探究的运动、知识的痕迹；总之，它给予救恩的方法。那些在真理的话语中受正确培育、并获得永生准备的人，展翅飞向天堂。因此，宗徒极佳地说：\Quote{我们在各样事上表明自己是天主的仆人；似乎贫穷，却使许多人富足；似乎一无所有，却拥有万物。我们的口向你们是张开的。} \Quote{我嘱咐你，}他写信给弟茂德说，\BibleQuote{在天主、基督耶稣及蒙选的天使面前，你要遵守这些事，不可偏袒，办事不可有私心。}\BibleRef{弟前 5:21}\par

因此，两者都必须省察自己：一人，是否有资格讲论并留下文字记录；另一人，是否处于正确的状态去聆听和阅读。同样，有些人在施行圣体圣事时，按惯例吩咐众人各自领受自己的部分。自己的良心最适于准确选择或回避。而它的坚实基础是正直的生活，加上适当的教导。但效法那些已经受过考验、生活正直的人，对于理解和实践诫命是最为卓越的。\BibleQuote{所以无论谁，若不配地吃主的饼，喝主的爵，就是干犯主的体血。但人应省察自己，然后才吃这饼，喝这爵。}\BibleRef{格前 11:27}因此，凡承担促进邻人益处的人，都应当考虑自己是否出于好胜而轻率地从事教导；他传布圣言是否出于虚荣；他所得的唯一报酬是否就是听者的得救，他的言语是否不为讨人欢心。若是如此，那以文字宣讲的人就免除了唯利是图的指责。\Quote{因为我们从未用过谄媚的话，这是你们知道的，}宗徒说，\Quote{也没有藉贪婪掩饰；天主是见证。我们也没有向人求荣耀，不从你们，也不从别人，我们作为基督的宗徒，本可加重你们的负担，但我们在你们中间却是温和的，如同乳母抚育自己的儿女。}\par

同样，那些参与圣言的人，应当谨防自己像建造城市那样出于好奇来审视它们；他们不应为了获得世俗之物而从事此事，因为他们已确知那些献身于基督的人必会分给生活必需品。但这样的人应被视为假善人而辞退。若有人不愿看似正义，而愿成为正义的，那么认识至善之事便属于他。因此，如果\BibleQuote{庄稼多，工人少}，我们就应当\BibleQuote{祈祷}，使工人在数量上尽可能多。\BibleRef{玛 9:37}\par

但耕种有两种——一种是不成文的，另一种是成文的。无论主的工人以何种方式播下好麦种，使其生长并收割麦穗，他必显为真正神圣的农夫。\BibleQuote{你们不要为那必坏的食物劳力，}主说，\BibleQuote{要为那存留到永生的食物劳力。}\BibleRef{若 6:27}养料既来自饼，也来自言语。的确，\BibleQuote{缔造和平的人是有福的，}\BibleRef{玛 5:9}他们教导那些在生命和错误中争战的人，引导他们归回圣言中的和平，并藉分饼滋养那些\Quote{饥渴慕义的人}。因为每个灵魂都有其适当的养料：有些藉知识和科学成长，另一些则以希腊哲学为食，而希腊哲学整个就像坚果，不可食用。\BibleQuote{那栽种的，那浇灌的，都是为那使之生长者服务的，在服务中合而为一。但是每个人将按自己的劳苦领取自己的赏报。我们原是主的农夫，主的农田。你们是主的建筑物，}\BibleRef{格前 3:8}照宗徒所说。因此，听者不得运用比较的检验。那为探究而赐下的圣言，也不可托付给那些在各种言语的巧计和夸大论证的权势中长大的人，他们的思想已被占据，未曾事先清空。但任何选择在信德上用餐的人，都能稳妥地接受圣言，因为他们已凭理性获得了作为判断能力的信德。因而，丰沛的说服力临于他。这正是先知预言之言的含义：\BibleQuote{如果你们不信，你们将不能明了。}\BibleRef{依 7:9}\BibleQuote{所以，我们一有机会，就应向众人行善，尤其应向有信德的人。}\BibleRef{迦 6:10}愿他们中的每一个人，按照有福的\Person{达味}，歌唱感谢。\Quote{求祢以牛膝草洒我，我就洁净；求祢洗涤我，我就比雪更白。求祢使我听到喜乐和欢欣，使祢所压伤的骨头雀跃。求祢转面不看我的罪过，涂抹我的一切过犯。天主，求祢在我内再造一颗纯洁的心，在我的五内注入新的正直之神。求祢不要从祢面前抛弃我，不要从祢身上收回祢的圣神。求祢再赐我祢救恩的喜乐，以慷慨之神扶持我。}\par

凡向在面前的人发言者，必以时间考验他们，以判断甄别他们，并从众人中辨别出能听者；他观察言辞、举止、习惯、生活、动作、姿态、眼神、声音；以及道路、岩石、践踏之路、肥沃之地、树木茂盛之区、丰饶美丽且耕种的处所，能使种子倍增。但那借书卷发言者，在天主前奉献自己，以写作这样呼喊：不为利益，不为虚荣，不为偏私所胜，不为恐惧所奴役，也不为快乐所鼓动；只为收获读者的救恩——他现在并未分享，而是期待那必将由应许按时赐给工人合适酬报的那一位所给予的报偿。但名列人中者不应渴望报偿，因为自夸善行者以荣耀为酬报；凡为报偿而行义务者，岂非被世间的习俗所束缚——或如行善者急于领受赏报，或如作恶者逃避惩罚？我们必须尽可能地效法主。凡遵行天主旨意者将会如此：白白领受，白白施予，并以公民身份本身作为相称的赏报。经上说：\Quote{娼妓的酬金不可进入圣所}，因此也禁止将狗的价银带上祭台。凡灵魂之眼被不良教养和教训蒙蔽的人，让他前进到真光、真理那里去，这真理借写作昭示那未写之事。依撒意亚说：\BibleQuote{你们凡口渴的，都到水边来} \BibleRef{依 55:1}。撒罗满劝告说：\BibleQuote{你要喝自己池中的水} \BibleRef{箴 5:15}。同样，在\Quote{\WorkTitle{法篇}}中，那从希伯来人学习的哲学家柏拉图命令农夫，在未先向下挖掘自己地中所谓处女土而发现其干燥之前，不得从他人处引水或取水。因为供给需要是对的，但支持懒惰并不好。因为毕达哥拉斯曾说过：\Quote{虽然分担重担合乎理性，但拿开它却不是本分。}\par

圣经点燃灵魂的活火花，引导眼睛于默观；或许像农夫接枝时插入一些东西，但按照神圣宗徒的意见，激发灵魂内已有的东西。\Quote{因为我们中间确实有许多软弱的、有病的，许多人在睡觉。但我们若分辨自己，就不致受审判。}如今我这部写作并非为炫耀而精心构思；我的笔记是被储存起来以备年老，作为对抗遗忘的药方，实在是那些我有幸聆听的活泼有力之言论，以及那些有福且真正卓越之人的形象和轮廓。\par

他们中一位在希腊（伊奥尼亚人），另一位在大希腊：第一位来自科埃勒-叙利亚，第二位来自埃及，其他人则在东方。一位生于亚述地，另一位是巴勒斯坦的希伯来人。\par

当我找到最后一位（他在能力上居首），在埃及将他追踪隐藏时，我得了安息。他，真正的西西里蜜蜂，采集了先知和宗徒草地上的花朵，在听众的灵魂中生成了不朽的知识元素。\par

他们保存了来自圣宗徒伯多禄、雅各伯、若望和保禄的蒙福教义的传统，儿子从父亲领受（但像父亲的人很少），也按天主的意愿来到我们这里，播下那些祖传和宗徒的种子。我很清楚他们会欢欣——我不是说因这敬意而喜悦，而仅仅因真理的保存，正如他们所交付的那样。我想，这样的纲要对于渴望不让蒙福传统逃脱的灵魂将是愉快的。\BibleQuote{爱慕智慧的儿子使父亲喜乐} \BibleRef{箴 29:3}。井水汲出后变得更纯净；无人享用的水会变为腐败。使用使钢更亮，废弃则生锈。总之，锻炼使灵魂和身体都获得健康。\BibleQuote{没有人点灯放在斗底下，而是放在灯台上，使那些被视为配得上宴席的人得见光明} \BibleRef{玛 5:15}。因为智慧若不能使能听者变聪明，有什么用处呢？因为救主仍然拯救，\Quote{并且常工作，如祂看见父所作的一样}。因为藉着教导，人学得更多；演讲时，人常与听众一同聆听。因为说话者和听者的老师是同一位——祂浇灌理智和言语。因此，主在守安息日时并不阻止行善，反而允许我们将那些神圣奥迹和那圣光传给能领受的人。祂当然没有把不属于多数人的事启示给多数人，而是启示给少数祂知道属于他们、有能力领受并按照它们塑造自己的人。但隐秘的事托付给言语而非文字，正如天主的情况一样。\par

若有人说经上记载：\Quote{没有隐藏的事不被显露，也没有秘密不被披露}，请也听我们说：对于暗中听的人，那隐秘的事要被彰显。这正是这神谕所预言的。对于能暗中观察所交付之事的人，那遮蔽的事要被显明为真理；那向多数人隐藏的事，要向少数人分明显现。因为为何并非所有人都认识真理？为何公义不被爱，如果公义属于所有人？但奥迹被奥秘地传递，使所说的话在说话者的口中，不如说不在他的声音里，而在他的理智中。\BibleQuote{祂赐给教会：有些做宗徒，有些做先知，有些做传福音者，有些做牧者和教师，为成全圣徒，为服务的工作，为建立基督的身体} \BibleRef{弗 4:11}。\par

我深知，我所写的这些笔记，与我有幸听到的那充满恩宠之灵相比，是软弱的。但它将是一个形象，使那被酒神杖击中的人想起原型。因为 \Quote{说，} 经上说， \Quote{对智慧者，他将更智慧；凡有的，还要加给他。} 我们自称不是充分地解释隐秘之事——远非如此——而只是回忆它们，无论我们是否遗忘，或为了不遗忘。许多事，我深知，因时间长久而遗漏，未写成文字。因此，为了帮助我记忆的软弱，并为我预备一个系统编排的章节作为有益的备忘，我不得不采用这种形式。有些事我们并不记得；因为那些有福的人身上的力量是巨大的。有些事长久未加注意，现已遗忘；另有些事被抹去，在心灵本身中褪色，因为这样的工作对无经验者不易；这些我在注释中重新唤起。有些事我故意省略，行使智慧的选择，害怕写下我防备去说的：并非吝惜——因为那是错的——而是为读者担心，恐怕他们因错误理解而跌倒；正如谚语所说，我们会被发现 \Quote{把剑递给孩童。} 因为所写的不可能不被泄露，即使我没有出版。但它们总是被反复思考，只使用写作这一唯一声音，对查询超出所写内容的人不作回答；因为它们必然需要某人——或原作者，或另一个跟随他足迹的人——的协助。我的论文将暗示一些事；对一些事将停留；一些事仅提及。它将试图不知不觉地说，秘密地展示，并沉默地证明。将引用著名派别教导的教义；并根据对知识最深奥的默观，一切应当预备的反对它们，这默观，当我们前进到那从世界创造开始就出现的著名而可敬的传统规范时，将呈现在我们眼前；摆在我们面前的是根据自然默观必须预先处理的事，并清除阻碍这一安排的事。这样我们就能使耳朵准备好接受真知识的传统；土壤被农夫预先清除荆棘和杂草，为栽种葡萄。因为有竞赛，以及竞赛的前奏；有些奥秘在另一些奥秘之前。\par

我们的书不会惧怕利用哲学和其他预备学识中最好的部分。按照宗徒，\BibleQuote{不但为希伯来人和那些在法律之下的人，} \BibleQuote{也成为犹太人，也成为希腊人，为赢得希腊人，好赢得所有的人。} \BibleRef{格前 9:20} 在 \WorkTitle{致哥罗森人书} 中也写道：\BibleQuote{劝诫每一个人，用一切智慧教训每一个人，为把每一个人在基督内成全起来。} \BibleRef{哥 1:28} 精细的思辨也适合我注释中所呈现的纲要。在这方面，学识的资财就像拌入运动员食物中的调味品，他并非耽于奢侈，而是怀着对卓越的崇高渴望。\par

音乐使我们和谐地缓解过度的严肃紧张。正如那些想要对民众讲话的人，常通过传令官来传达，使所说的话更容易被听见；在这事上也是如此。因为我们有那在共同传统之前曾对多人讲过的言语。因此，我们必须陈述那些个别向他们呼喊的意见和言辞，使听者更容易回转。\par

而且，简言之：在许多小珍珠中，有那一颗；在一大网鱼中，有那条美鱼；借着时间和辛劳，真理必将闪耀，如有好的帮手在手。因为多数恩惠是由天主通过人来提供的。我们所有使用眼睛的人都看到眼前呈现的事物。但有些人观察对象是为了一个理由，另一些人是为了另一个。例如，厨师和牧羊人查看羊的方式不同：一个检查它是否肥壮；另一个观察它是否品种优良。若有人需要食物，就挤羊奶；若需要衣物，就剪羊毛。就让我这样结出希腊学问的果实吧。\par

因为我不以为有任何作品能如此幸运，以致无人对它非议。但只有那些无人有理由非议的，才应被视为合乎理性。那应被赞许的行为和选择，并非毫无瑕疵，而是无人有理性地挑剔它。因为一个人若并非有意地完成某事，并不等于他是在环境的强迫下作的。但他会凭天主赐予的智慧，并适应环境而完成它。因为拥有德行的人不需要通往德行的道路，正如强壮的人不需要恢复健康一样。因为如同农夫预先灌溉田地，我们也是如此用希腊学问的流动之水浇灌其中尘土的成分，使它能够接受撒在其中的精神种子，并易于滋养它。\WorkTitle{杂录集}将包含掺杂在哲学教义中的真理，或者更确切地说，覆盖并隐藏着的真理，如同坚果壳中可食的部分。因为在我看来，真理的种子只应保存给信仰的农夫，而非他人。我没有忘记某些人在无知中因每一声响而惊慌，他们胡说八道，说我们应当专注于最必要且包含信仰的事物，而应当越过那些多余且无益、无谓消耗并拖延我们、对伟大终局毫无贡献的事物。另一些人认为哲学是由一个邪恶的发明者出于邪恶的影响引入生活，为的是毁灭人类。但我将在整部\WorkTitle{杂录集}中表明，邪恶具有邪恶的本性，绝不可能成为任何美善之事的产出者；同时指出，哲学在某种意义上是一项天主眷顾的工作。\par

% structure: p0020 source=h2 target=subsection reason=relative-heading-rank; base=h2

\subsection{第二章 反对本书中哲学著作摘录的数量之预判与答复}

关于这些注释，其中按情况需要包含希腊见解，我对那些爱挑剔的人说这样的话。首先，即使哲学无用，对无用性的证明若有益处，它仍然是有用的。其次，那些仅凭道听途说了解希腊人见解、并未对各领域作细微探究以认识它们的人，不能谴责希腊人。因为基于经验的驳斥是完全可信的。对被谴责之物的认识，被证明是最完全的论证。因此，许多事物虽对最终结果无直接贡献，却装备了艺术家。此外，学识会推荐那陈述最本质教义的人，在听众中产生说服力，在受教者中引发钦佩，并引导他们走向真理。这样的说服是令人信服的，使爱好学习的人接受真理；因此哲学并非因成为错误行为与卑劣行径的始作俑者而毁灭生活——尽管有人诬蔑它，尽管它是真理的清晰图像、是赐予希腊人的神圣恩赐——哲学也并未使我们脱离信仰，仿佛我们被某种欺骗术所迷惑，而是通过使用更迂回的路径，获得一种共同操练以证明信仰。更进一步，教义的并列对照，通过比较，保全了真理，由此带来知识。\par

哲学之所以产生，并非为了自身，而是为了我们从知识中获得益处，即通过理解心智所领会之事物，我们获得了对真实知觉的坚定确信。因为我不提\WorkTitle{杂录集}这一包含多样学识的文集，巧妙地希望隐藏知识的种子。因此，就如好猎者经过搜寻、追踪、嗅探、用犬追捕而捕获猎物，真理在经劳苦寻求得到后，就显为甘美。你或许会问：为何你认为在你的笔记中采用这样的编排是合适的？因为将真哲学的奥秘泄露给那些不爱惜、无正义地反对一切，并毫无体统地喊叫各种名字和话语，欺骗自己、迷惑那些依附他们的人，是非常危险的。\Quote{因为犹太人要求的是神迹，}正如宗徒所说，\Quote{希腊人寻求的是智慧。}\BibleRef{格前 1:22}\par

% structure: p0023 source=h2 target=subsection reason=relative-heading-rank; base=h2

\subsection{第三章 反对诡辩家}

这一类别的人很多：其中一些沉溺于享乐、愿意怀疑，嘲笑那值得一切尊敬的真理，拿它的野蛮性取乐。另一些人则自高自大，试图找出对我们话语的诽谤性反对意见，提出苛责的问题，搜寻琐屑的言辞，操弄可怜的伎俩，好争辩，纠缠于缠绕不清之点，正如那阿布德拉人所说：——\par

\Quote{因为凡人的舌头是滑的，上面有许多言语；\newline{}而且在这处那处有各种话语的广阔天地。}\par

以及——\par

\Quote{不论你说出什么样的话，也必听见同样的话。}\par

那些可怜的诡辩家夸耀于他们的这套技艺，用他们自己的术语喋喋不休；终生劳苦于名称的区分、句子成分和连词的性质，显得比斑鸠更饶舌；依我之见，他们不像男子汉那样，而是搔痒、呵痒那些希望被呵痒的人的耳朵。\par

\Quote{一条愚蠢言语的河——不是一滴；}\par

正如旧鞋，当其余部分都磨损破碎，只剩下鞋舌时。雅典的梭伦极出色地展开论述，写道：——\par

\Quote{要看那油嘴滑舌之人的舌头和言语，\newline{}但你却看不到任何已完成的工作；\newline{}但你们每个人都循着狐狸的脚步行走，\newline{}在你们所有人心里是空虚的意念。}\par

我认为，救主的这句话所表明的就是：\BibleQuote{狐狸有洞，但人子却没有枕头的地方。} \BibleRef{玛 8:20} 因为唯有那完全与其他人——圣经称之为野兽——分离的信徒，宇宙的头颅、仁慈而温柔的圣言才得以安息，\Quote{他以自己的智慧擒获智者；因为主知道智者的思念，原是虚妄的；} 圣经称那些擅长言辞和技艺的人为\OriginalTerm{智者}{\GreekText{σοφούς}}、\OriginalTerm{诡辩者}{\GreekText{σοφιστάς}}。因此，希腊人也把智慧的称号赋予那些精通某种技艺的人，称他们为智者和诡辩者（\OriginalTerm{智者}{\GreekText{σοφοί}}、\OriginalTerm{诡辩者}{\GreekText{σοφισταί}}）。据此，\Person{克拉提努斯}在\WorkTitle{阿尔基洛科斯}中列举诗人之后说：—\par

\Quote{你考察了这样一群诡辩者。}\par

同样，喜剧诗人\Person{伊奥丰}在\WorkTitle{吹笛的萨提尔}中说：—\par

\Quote{因为进来了\newline{}一队诡辩者，全副武装。}\par

对于这些以及类似专注于空谈的人，神圣的圣经最好地说：\BibleQuote{我要摧毁智者的智慧，废除明达者的明达。} \BibleRef{依 29:14}\par

% structure: p0037 source=h2 target=subsection reason=relative-heading-rank; base=h2

\subsection{人为艺术与神圣知识皆出自天主}

\Person{荷马}称工匠为智者；关于\WorkTitle{马尔吉特斯}（如果是他的作品），他这样写道：—\par

\Quote{那时，众神使他既非掘地者，也非耕田者，\newline{}在其他方面也不智慧；但他错过了每一门技艺。}\par

\Person{赫西俄德}还说音乐家\Person{利诺斯}是\Quote{精通一切智慧}；并且毫不犹豫地称水手为智者，因为他写道：—\par

\Quote{在航海方面没有智慧。}\par

\Person{达尼尔}先知说：\BibleQuote{君王询问的奥秘，不是智者、贤士、占卜者、加色丁人所能告诉君王的；唯有在天上的天主才能启示。} \BibleRef{达 2:27}\par

此处他称巴比伦人为智者。圣经将每一门世俗科学或艺术统称为智慧（除此以外，还有人类理性发明的其他艺术和科学），并且艺术性和技巧性的发明来自天主，如果我们引用以下陈述，这一点将变得清楚：\BibleQuote{上主对\Person{梅瑟}说：看，我召叫了犹大支派中\Person{乌黎}的儿子、\Person{奥尔}孙的\Person{贝匝肋耳}；我以天主的神充满了他，使他有智慧、聪明、知识，能设计并执行各种工作，能雕刻、设计，用金、银、铜、蓝色、紫色、朱红色线，以及雕刻宝石的工作和雕刻木头的工作，} 甚至 \BibleQuote{一切工作。} \BibleRef{出 31:2} 然后他补充了一般性的理由：\BibleQuote{我也赐给一切有智慧的心智慧；} \BibleRef{出 31:6} 即赐给每一个能够通过辛劳和练习获得智慧的人。又明确地以上主的名写道：\BibleQuote{你应吩咐所有有智慧的人，就是我用知觉的灵所充满的。} \BibleRef{出 28:3}\par

有智慧的人具有某种与生俱来的特质；那些显明自己有能力的人，从至高智慧那里获得双倍的知觉之灵。因为从事普通技艺的人，在感官方面极具天赋：在听觉上，是通常被称为音乐家的人；在触觉上，是捏塑泥土的人；在声音上，是歌手；在嗅觉上，是调香师；在视觉上，是雕刻印章图案的人。那些从事教学的人，则训练敏感性，据此，诗人容易受到韵律的影响；诡辩者领会表达；辩证法家领会三段论；哲学家则能够进行以自身为对象的沉思。因为敏感性发现并创造；因为它恳切地劝勉人应用。实践将增加以知识为目的的应用。因此，宗徒有理由称天主的智慧为\BibleQuote{丰富多彩的}，并且它\BibleQuote{在诸多领域和诸多方式中}显示了它的能力 \BibleRef{弗 3:10}——通过艺术、知识、信仰、预言——为了我们的益处。\BibleQuote{因为一切智慧皆来自上主，并且永远与祂同在，} 正如\WorkTitle{耶稣的智慧}所说。 \BibleRef{德 1:1}\par

因为如果你大声呼唤智慧和知识，寻求它如同寻觅银子，急切地追踪它，你就会明白虔诚，找到属神的知识。 \BibleRef{箴 2:3} 先知这样说，与依据哲学的知识相对立，哲学教导我们以高尚和尊贵的方式探究，以促进我们虔敬的进步。因此，他将专注于虔敬的知识与之对立，当论及知识时，他这样说：\Quote{因为天主亲口赐予智慧，并赐予知识和聪明，为义人积蓄援助。} 因为对于那些被哲学成义的人，那导向虔敬的知识已被储备为帮助。\par

% structure: p0046 source=h2 target=subsection reason=relative-heading-rank; base=h2

\subsection{哲学：神学之婢}

因此，在主来临之前，哲学对外邦人是为正义所必需的。而现在，它有助于虔敬，成为那些通过论证获得信仰之人的某种预备训练。\BibleQuote{因为你的脚，} 经上说，\BibleQuote{必不绊跌，如果你把好的事物，无论是属于希腊人的还是属于我们的，都归之于天主眷顾。} \BibleRef{箴 3:23} 因为天主是一切美善之因；但有些是首要的，如旧约和新约；有些是间接的，如哲学。或许，哲学也是直接和首要地赐给希腊人的，直到主呼召希腊人。因为这曾是启蒙导师，将\Quote{希腊人的心意}，正如法律带领希伯来人一样，领到\BibleQuote{基督} \BibleRef{迦 3:24}。因此，哲学是一种预备，为那在基督内得以成全的人铺路。\par

\Quote{现在，}\Person{撒罗满}说，\Quote{守护智慧，它必高举你，并以喜乐的华冠遮蔽你。}\BibleRef{箴 4:8}因为当你藉着哲学的外衣，并以正当的付出加强了智慧，你就能保存它，不受诡辩者的攻击。因此，真理之路是唯一的。但众流从各方流入其中，如同流入一条永不止息的河流。因此，受默感而这样说：\Quote{我儿，要听受我的言语，使你的道路成为多条生命之路。因为我教导你智慧之道，使泉源不致断绝，}这些泉源从大地本身涌出。祂不仅列举了为任何义人的几条救恩之路，还添上了许多其他义人的多条路，这样说：\BibleQuote{义人的道路好像黎明的光，越来越亮。}\BibleRef{箴 4:18}诫命和预备训练的方式应被视为生命之路和工具。\par

\BibleQuote{耶路撒冷，耶路撒冷，我多次愿意聚集你的子女，好像母鸡聚集小鸡！}\BibleRef{玛 23:37}至于耶路撒冷，按解释，是\Quote{和平的异象}。因此祂预言性地指出，那些和平默思神圣事物的人，以多种方式被训练归向他们的召叫。那么是什么？祂\Quote{愿意}，却不能。多少次，在何处？两次：藉着先知，和藉着降临。因此，\Quote{多少次}这一表述显示出智慧是多样的；它藉所有方式和数量及质量，无论如何都拯救一些人，既在时间中，也在永恒中。\Quote{因为主的神充满了大地。}若有人强解这是指希腊文化，当它说：\Quote{不要留心邪恶的女人，因为妓女的嘴唇滴下蜂蜜，}让他听下面的话：\Quote{她暂时润滑你的喉咙。}但哲学并不谄媚。那么，祂指谁行了邪淫？祂明说：\Quote{因为愚昧的脚引导使用她的人，死后进入阴间。但她的脚步并不稳固。}因此要使你的路远离愚昧的快乐。\Quote{不要站在她家的门口，免得你将生命交给别人。}祂又作证：\Quote{那时你会在老年后悔，当你肉体的皮肉耗尽的时候。}因为这就是愚昧快乐的结局。事实正是如此。当祂说：\Quote{不要与陌生的女人多交往，}祂告诫我们要使用世俗文化，但不要逗留和花费时间在其中。因为每一世代所获得的有利和合时的预备，都是主圣言的前期训练。\Quote{因为已经有一些人，被婢女的魅力所诱，藐视了他们的配偶哲学，并且老去，其中一些沉溺于音乐，一些沉溺于几何，另一些沉溺于文法，多数沉溺于修辞学。}\Quote{但正如通识学科服务于哲学，哲学是它们的主母；同样，哲学本身也有助于获得智慧。因为哲学是智慧的学习，而智慧是对神圣事物和人类事物及其原因的知识。}因此，智慧是哲学的女王，正如哲学是预备文化的女王。因为若哲学\Quote{宣称控制舌头、肚腹及肚腹以下的部分，它本身就值得选择。但若为着天主的尊荣和知识而培育，则显得更值得尊敬和卓越。}圣经将在下文为所说的事提供证据。撒辣曾一时不生育，是亚巴郎的妻子。撒辣无子，将她的婢女，名叫哈加尔的埃及人，给了亚巴郎以便得子。因此，与具有信德的人（亚巴郎被算作有信德和正义的）同在的智慧，在那时代仍是不生育、无子女的，没有为亚巴郎生出任何与德行相关的事。按适当的方式，她认为他既处在进步的时期，应先与世俗文化交往（埃及象征地表示世界），然后按照天主的眷顾接近她，并生依撒格。\par

\Person{斐罗}解释哈加尔意为\Quote{寄居}。因为与此相关，经上说：\Quote{不要与陌生的女人多交往。}他解释撒辣意为\Quote{我的王权}。因此，那接受了预先训练的人，可以自由地接近最崇高的智慧，以色列民族就从这智慧成长起来。这些事表明，那智慧是可以藉着训导获得的，亚巴郎从默观天上事物转而达到符合天主的信德和正义，就获得了这智慧。依撒格被显示意为\Quote{自学}，因此他被发现是基督的预像。他娶了一妻黎贝加，她们翻译为\Quote{忍耐}。雅各伯据记载与多位结合，他的名字解释为\Quote{操练者}。操练是藉着许多不同的教义进行的。因此，那真正\Quote{赋有看见的能力}的人被称为以色列，他多有经验，且适合操练。\par

三位圣祖也可能显示了另一件事，即：知识的确切印记是由本性、教育和操练构成的。\par

对于前面所说的，你也可以有另一个形象：塔玛尔坐在路旁，扮作妓女的模样；好学的犹大——他的名字解释为\Quote{有能力的}——什么都不放过、什么都探究，看见了她，就转过去到她那里，同时保持他对天主的忠诚。因此，当撒辣因哈加尔受宠而嫉妒时，亚巴郎——好像只拣选世俗哲学中有益的东西——说：\BibleQuote{看，你的婢女在你手中；你随意待她吧！}\BibleRef{创 16:6}显然意思是：\Quote{我拥抱世俗文化如同青年和婢女；但你的知识我尊敬且珍视如真妻。}撒辣就虐待她，这等于纠正和劝诫她。所以这话说得很好：\BibleQuote{我儿，不要轻视上主的惩戒，也不要因祂的责斥而灰心；因为上主所爱的，祂必惩戒，祂鞭打所收纳的每一个儿子。}\BibleRef{箴 3:11}上述圣经在别处查考时，将显示其他奥秘。因此我们在这里仅断言，哲学的特征是探究真理和事物的本性（这就是主亲自说的真理：\BibleQuote{我是真理}\BibleRef{若 14:6}）；而且，为在基督内安息所作的预备训练，借着真正的哲学——这哲学由入门者获得，或者说从真理本身接受——锻炼心智，激发理智，并产生探索的敏锐。\par

% structure: p0053 source=h2 target=subsection reason=relative-heading-rank; base=h2

\subsection{第六章 文化的益处}

先前训练所获得的敏捷，大大有助于感知那些必要的事物；但那些只能被心智感知的事物，则是心智特有的锻炼。就其性质而言，根据我们考虑它们的数量、大小和可述谓的性质，有三种类型。因为由论证构成的话语，在追随者心中植入清晰的信仰，以致他无法想象所证明的事物会不同；因此它不容许我们屈服于那些用欺诈攻击我们的人。在这样的学习中，灵魂从可感觉的事物中被净化，并被激发，以便能够清晰地看见真理。因为营养和温和的训练造就高贵的本性；而高贵的本性在接受了这样的训练后，在其他方面变得比以前更好，尤其在生产力方面，正如其他受造物一样。因此经上记着说：\Quote{懒惰的人哪，你去察看蚂蚁，就可得智慧。蚂蚁在收割时预备丰富的、多样的食物，以应对冬天的严酷。}或者去看看蜜蜂，学习她多么勤劳；因为她以整个草地为食，却产出一个蜂巢。如果你\Quote{进入你的内室祈祷}，如主所教导的，\Quote{以心神敬拜}\BibleRef{玛 6:6}，你的管理就不再只专注于房子，也专注于灵魂：什么该给予灵魂，以及如何给予、给予多少；什么该搁置、珍藏于灵魂；何时该产生、给谁。因为人成为高尚和善良，不是由于本性的，而是由于学习，正如人成为医生和舵手一样。例如，我们都看到葡萄树和马。但农夫会知道葡萄树在结果实方面是好是坏；骑手会轻易区分无精神和敏捷的动物。但有些人天生比别人更有德性的倾向，某些天生如此倾向者的追求就显示出这一点。然而，德性的完美并非专属那些本性更好的人，这一点得到证明：因为那些天性对德性不利的人，在获得适当的训练后，大多能达到卓越；另一方面，那些天资适合的人，因忽视而变坏。\par

再者，天主按本性造我们为合群和正义的；因此，正义不能说是仅仅来自植入。而是，受造所赋予的善，当被视为由诫命所激发；灵魂受训练，愿意选择最高贵的事物。\par

但是，正如我们说一个人可以不经学习而成为信徒，我们也断言，不经学习就不可能理解信仰中所宣告的事物。然而，采纳说得好的，而不采纳相反的，不仅是由信仰引起的，而且是由信仰加上知识引起的。但如果无知是缺乏训练和教导，那么教导就产生对神圣和人类事物的知识。正如在世物的匮乏中有可能正当地生活，在丰裕中也是如此。我们宣称，通过先前的训练，人将更轻松、更迅速地达到德性。但这并非没有训练就无法达到；而是只有当他们学习了、感官受了锻炼之后才能达到\BibleRef{希 5:14}。因为撒罗满说：\Quote{仇恨引起争端，但教导守护生命的路途}——这样我们就不会被那些为了伤害听者而操练卑劣技艺的人欺骗或迷惑。经上说：\Quote{教导却无过地游荡。}\BibleRef{箴 10:19}我们必须熟习推理的艺术，以反驳诡辩家的欺骗性意见。因此，阿那克萨尔科斯在他的著作\WorkTitle{王权}中写得好而恰当：\Quote{学问大大有益于拥有它的人，也大大有害；它帮助配得上的人，却伤害那些轻易说出每句话、在全体民众面前说话的人。必须知道时间的尺度，因为这是智慧的终结。那些在门前唱歌的人，即使唱得巧妙，也不被视为明智，反得愚昧之名。}赫西俄德说：\par

\Quote{缪斯们，使人能言善辩、神圣、声音悦耳。}\par

因为对于言辞流利的人，他称为\Quote{能言善辩}；对于聪明的人，称为\Quote{声音悦耳}；而\Quote{神圣}则指有技巧的人、哲学家、通晓真理的人。\par

% structure: p0059 source=h2 target=subsection reason=relative-heading-rank; base=h2

\subsection{第七章 折衷哲学为神圣美德铺路}

希腊的预备文化，连同哲学本身，被证明是从天主降下给人，并非有明确方向，而是如同雨露降在好地、粪堆和房屋上。同样，草和麦子发芽；无花果和任何其他肆意生长的树木在坟墓上生长。生长之物，显现为真理的预像。因为它们享受同样的雨水之恩。但它们不像那些在肥沃土壤中生长的植物那样有同样的恩宠，因为它们是枯萎或被拔除。在此，我们得到主所解释的撒种比喻的帮助。因为人在其中的土壤的农夫只有一位：祂从起初，从世界的根基，撒下了滋养的种子；祂在每个时代降下了主，即圣言。但接受这些礼物的时间和地点造成了存在的差异。此外，农夫不仅撒麦子（其中有许多品种），也撒其他种子——大麦、豆子、豌豆、野豌豆以及蔬菜和花卉种子。同样的农耕也包含种植以及苗圃、花园、果园中必要的操作，以及各种树木的规划与培育。\par

同样，不仅是羊的照管，还有牲畜的照管、养马、养狗、养蜂，一切技艺，且概括地说，牧群照管和动物饲养，彼此或多或少不同，但对生活都有用。哲学——我并非指斯多葛学派、或柏拉图学派、或伊壁鸠鲁学派、或亚里士多德学派，而是这些派别各自所正确言说的，即讲授正义并带着一种以虔诚贯穿的知识——这种折中的整体我称之为哲学。但人类推理的那些结论，若被人加以切割和歪曲，我绝不会称之为神圣的。\par

现在我们也必须看看这一点：若那些不知如何行善的人生活得好，因为他们偶然碰上了善行。有些人也通过理解正确瞄向了真理之言。\Quote{但亚巴郎不是因行为成义，而是因信德。}\EditorialNote{罗马书第4章}因此，他们在生命终结之后没有任何益处，即使他们现在行善，若没有信德。所以圣经也被翻译成希腊文，以便他们永远不能以无知为借口，因为他们也能听到我们手中所有的，只要他们愿意。一个人以一种方式谈论真理，而真理以另一种方式解释自身。猜测真理是一回事，真理本身是另一回事。相似是一回事，事物本身是另一回事。前者来自学习和实践，后者来自能力和信德。因为虔敬的教导是恩赐，但信德是恩宠。\Quote{因为遵行天主的旨意，我们就知道天主的旨意。}\BibleRef{若 7:17}\Quote{那么，敞开}圣经说，\Quote{正义的门，我要进去，向上主称谢。}\Quote{这是上主的门，正义者要从它进去。}虽然有许多门敞开，但在基督内的正义之门，所有有福者由此进入，并在知识的圣德中指引他们的脚步。现在，克莱孟在\WorkTitle{致格林多人书}中，解释那些在教会内受悦纳之人的差异时，明确地说：\Quote{一人可以是信徒；一人可以是强有力于讲论知识者；一人可以是智慧于分辨言辞者；一人可以是可畏于行为者。}\par

% structure: p0063 source=h2 target=subsection reason=relative-heading-rank; base=h2

\subsection{第八章 诡辩之术无用}

但希腊人培养的诡辩术，是一种幻想的技巧，通过言语使错误的意见显得像真的一样。因为它产生修辞学以达说服，产生辩论术以达争论。因此，这些技巧若不与哲学结合，对每个人都将有害。因为柏拉图公开称诡辩术为一种邪恶的艺术。而亚里士多德追随他，证明它是不诚实的艺术，以貌似有理的方式抽离了智慧的全部事务，并声称拥有其未曾研究的智慧。简言之，修辞学的开端是或然之物，过程是尝试性证明，目的是说服；辩论术的开端是意见之物，过程是争辩，目的是胜利。同样，诡辩术的开端是表象，过程是双重的：一是修辞学的，连续而详尽；二是逻辑的，是审问式的。其目的是钦佩。另一方面，学校中流行的辩证法，是哲学家在意见之事上的操练，为了辩论的能力。但真理根本不在这些之中。因此，尊贵的宗徒有理由贬低这些无用的、专注于言辞的技巧，说：\Quote{若有人不听从健康的言辞，反而因某种教导而自大，一无所知，却\OriginalTerm{患病}{\GreekText{νοσῶν}}于问题和言语之争，由此产生争竞、嫉妒、毁谤、恶意的猜疑、心术败坏、丧失真理之人的乖谬争论。}\par

你看他如何动怒反对他们，称他们那套逻辑技巧——那些喜爱这饶舌有害技巧的人，不论希腊人或外邦人，都以之为荣——为一种\OriginalTerm{疾病}{\GreekText{νοσος}}。因此，悲剧诗人欧里庇得斯在\WorkTitle{腓尼基妇女}中说得非常美——\par

\Quote{但错误的言论，\newline{}本身便是病态，需要精妙的药物。}\par

因为救恩的圣言被称为\Quote{健康的}，祂就是真理；而健康的（有益健康的）永远不死。但远离健康而神圣的便是亵渎，是致命的疾病。这些是披着羊皮的贪婪豺狼，拐卖人口者，花言巧语引诱灵魂者，暗中却是强盗；以欺诈和暴力试图捕捉我们这些没有经验且言辞能力较弱的人。\par

\Quote{人常因辞令不足，表达真理时，不及雄辩者有力。\newline{}但如今在滔滔不绝的口中，\newline{}他们掩饰最重大的真理，\newline{}使其看起来不似应然。}\par

悲剧作家说。这些争论者就是如此，无论他们是追随各学派，还是操练可怜的辩证法技艺。他们就是那些\Quote{张了经却织不出什么}的人，圣经说；从事一项徒劳的任务，宗徒称之为\Quote{人的诡计，用来欺骗}。\BibleRef{弗 4:14}\Quote{因为有许多不受约束的、虚浮的谈论者和欺骗者}，他说。\BibleRef{铎 1:10}因此，并非对所有人都说了\BibleQuote{你们是地上的盐}。\BibleRef{玛 5:13}因为甚至有些听圣言的人，也像海里的鱼，从出生就在咸水中长大，却还需要盐来调味才能吃。因此我完全赞同这悲剧作家的话：—\par

\Quote{啊，儿子，虚假的话可以说得很好，\newline{}真理可能被言辞之美打败。\newline{}但这并非最正确，而是自然与正义；\newline{}练习口才的人固然聪明，\newline{}但我认为行动总是优于言语。}\par

我们不可追求取悦大众。因为我们所行的不是讨他们喜欢的事，而是我们知道远离他们心意的事。宗徒说：\BibleQuote{我们不要贪图虚浮的荣耀，彼此惹气，彼此嫉妒。}\BibleRef{迦 5:26}\par

于是爱好真理的柏拉图，仿佛受到默启，说：\Quote{既然我如此，除了那经过反省后在我看来最好的话，我什么都不服从。}\par

因此他指责那些没有智识就相信意见的人，无端抛弃了正确健全的理性，而相信那个参与谎言的人。因为欺骗自己离开真理是坏的；但讲说真理，并把实在的现实作为我们的见解，是好的。\par

人们不知不觉地被剥夺了善。然而，他们被剥夺，或是由于受欺骗、被诱惑，或是由于被迫而不信。不信的人已经使自己成为自愿的俘虏；改变信念的人是被欺骗了，因为他忘记了时间不知不觉地夺走一些事物，理智夺走另一些。在意见被接受之后，痛苦和焦虑，以及另一方面的好争和愤怒，迫使人改变。最重要的是，那些被快乐迷惑或被恐惧吓倒的人，是被诱骗的。所有这些都是自愿的改变，但知识永远不会通过其中任何一种获得。\par

% structure: p0075 source=h2 target=subsection reason=relative-heading-rank; base=h2

\subsection{第九章 理解圣经所需的人类知识}

有些人自以为天赋异禀，不愿接触哲学或逻辑；不仅如此，他们也不愿学习自然科学。他们只要求单纯的信德，仿佛希望不费心照顾葡萄树，就直接从开头收取果实。主被比喻为葡萄树，要借着辛劳和农艺，按照圣言，才能收取果实。\par

我们必须修剪、挖掘、捆绑，并做其他操作。修枝刀、镐和其他农具，我想，对于栽培葡萄树是必要的，好让它结出可吃的果实。如同在农业中，在医学中也是如此：他学习以达到目的，练习了各种课程，以便能够栽培和医治。同样在这里，我称他真正有学问，他把一切带到真理上；因此，从几何、音乐、语法、哲学本身，摘取有用的，他守护信德免受攻击。正如已经说过的，运动员若不为比赛做好准备，就会受到鄙夷。例如，我们赞扬见过众多城市的经验舵手，以及经验丰富的医生；有些人也这样描述经验主义者。他把一切都引向正直的生活，从希腊人和外邦人那里获取例子，这人就是经验丰富的真理探索者，实际上是个足智多谋的人，像试金石（即吕底亚石），据说具有区分真金和假金的能力。而我们无所不知的\OriginalTerm{真知者}{gnostic}可以区分诡辩与哲学、装饰艺术与体操、烹饪与医学、修辞学与辩证法，以及那些按照蛮族哲学的其它派别与真理本身。渴望参与天主大能的人，通过哲学思考来处理知性主题，这是多么必要！区分模糊的表述，以及那些在圣约中同义使用的表述，是多么有用！因为主在受试探时，巧妙地用歧义表述来应对魔鬼。在这一点上，我还不明白，像某些人认为的，哲学和辩证法的发明者怎么能被言语形式的歧义所迷惑？如果先知和宗徒不知道那些展示哲学操练的艺术，那么先知和教导之神的意念秘密发出，因为并非所有人都有理解的耳朵，这要求有技巧的教学模式以便清晰的阐述。因为先知和圣神的门徒无误地知道他们的意念。因为他们借着信德知道，以一种别人不容易明白的方式，正如圣神所说。但对那些没有学会这样领受的人来说，是不可能的。\Quote{写，}它说，\Quote{要在谋划和知识中写下双重的诫命，使你能回答那些打发人问你的人的真理之言。}那么，回答的知识是什么？或者提问的知识是什么？是辩证法。那么，怎样？说话岂非我们的事？行为岂非来自圣言？因为如果我们不为圣言而行动，便是违背理性。但理性的工作借天主完成。\Quote{并且，}它说，\Quote{没有一样不是借着祂造的}——天主的圣言。\BibleRef{若 1:3}\par

主岂非藉圣言造成万物么？连野兽也在迫切的畏惧驱动下工作。那些被称为正统的人，岂非也在行善，却不知自己所作的是什么？\par

% structure: p0079 source=h2 target=subsection reason=relative-heading-rank; base=h2

\subsection{第十章 善行重于善言}

因此救主拿起饼来，先说话并祝福。然后擘开饼，递给我们，使我们能按照理性吃，并因认识圣经而顺服地行走。正如那些言语邪恶的人并不比行为邪恶的人更好（因为诽谤是刀剑的仆役，恶言伤人；由此生出生活中的灾祸，这些正是恶言的效果）；同样，那些善于言谈的人也接近那些成就善行的人。因此，言谈滋润灵魂，引向高尚；那能双手并用的人是有福的。所以，能善行的人不应被能善言的人贬低；能善言的人也不应被能善行的人轻视。但各人当行其本性所适合的事。一人展示实际所行，另一人则以言语预备善行的道路，引导听者实践善。因为既有拯救之言，也有拯救之行。义德原非没有言语。如果我们废除了善行，领受善行也就被废除了；同样，如果我们不带着命令和解释命令的人同行，服从与信德也就被废除了。但现在，我们藉言语和行为彼此互相受益；但我们必须完全摒弃争辩与诡辩之术，因为这些诡辩家的言论不仅迷惑和欺骗众人，有时还以暴力赢得一场卡德摩斯的胜利。因为那圣咏的话最是真实：\Quote{义人要活到末后，因为他见不到朽坏，当他看见智者死去时。} 他所称的智者是谁？请听\WorkTitle{耶稣的智慧}所说：\BibleQuote{智慧不是关于邪恶的知识。}\BibleRef{德 19:22} 这正是他所谓言语和讨论的艺术所发明的东西。\BibleQuote{你因此要在恶人中寻求智慧，却找不到。}\BibleRef{箴 14:6} 如果你再问这是怎样的智慧，你会被告知：\BibleQuote{义人的口要滴出智慧。}\BibleRef{箴 10:31} 同样，诡辩之术也被称为智慧。\par

但我认为，我的目的——并非没有理由——是依照圣言生活，理解所启示的真理；但绝不追求雄辩，只满足于表明我的意思。至于用什么措辞来表达我所要呈现的内容，我并不介意。因为我深知，获得救恩并帮助那些渴望得救的人，是最好的事，而不是像装饰品一样编造琐碎的句子。\Quote{如果你警惕对术语的过分关注，}\Person{柏拉图}\WorkTitle{政治家}中的毕达哥拉斯派说，\Quote{你将在老年时更富智慧。}在\WorkTitle{泰阿泰德篇}中你又会找到：\Quote{对名称、措辞的漫不经心，以及缺乏精细审视，大都不算粗俗和鄙陋，反倒是其反面，有时还是必要的。}圣经以最简洁的方式表达了这一点，它说：\Quote{不要过多地专注于言辞。}因为表达就像身体上的衣服。实质是血肉和筋骨。所以我们不应更关心衣服胜过身体的安全。因为不仅生活方式要简朴，说话风格也要避免多余和矫饰，这是那选择了真实生命的人所当培养的，如果我们应当摒弃奢侈，因其诡诈和放荡，正如古\Place{拉刻代孟}人弃绝香膏和紫色，正确地将它们称为诡诈的衣服和诡诈的香膏；因为那种调味品多于营养的烹饪方式不对，那种只求愉悦而非裨益听者的说话风格也不对。\Person{毕达哥拉斯}劝我们考虑缪斯比塞壬更令人愉悦，教导我们培养脱离愉悦的智慧，而揭露另一种吸引灵魂的方式是欺骗性的。因为驶过塞壬时，一人有足够的力量，而回答斯芬克斯则需要另一个人，或者，如果你愿意，甚至一个人也不够。我们决不可出于虚荣心而放宽经文匣。对于\OriginalTerm{诺斯替者}{gnostic}而言，只要找到一个听者就足够了。因此你可以听到\Place{波奥提亚}的\Person{品达}写道：\Quote{不要向所有人泄露古老的言语。沉默之道有时最为稳妥。最有力的话语是战斗的刺激。}同样，蒙福的宗徒非常恰当而迫切地劝勉我们\Quote{不要作无益的言辞之争，以致倾覆听者，但要远避世俗的空谈，因为它们使人更不虔敬，他们的话如同毒疮蔓延。}\par

% structure: p0082 source=h2 target=subsection reason=relative-heading-rank; base=h2

\subsection{第十一章 宗徒要我们避免的哲学是什么？}

这样看来，\BibleQuote{这世界的智慧在天主面前是愚拙的，}以及那些\BibleQuote{智慧人的，主知道他们的思念是虚妄的。}\BibleRef{格前 3:19} 因此，没有人应在人的思想出众上夸耀。因为在耶肋米亚中写得很好：\BibleQuote{智慧人不要夸耀自己的智慧，勇士不要夸耀自己的勇力，富人不要夸耀自己的财富；夸耀的，只应夸耀自己明了和认识我，我是主，在地上施行慈爱、公义和正直，因为我喜爱这些事——主说。}\BibleRef{耶 9:23} \BibleQuote{我们不要信赖自己，而要信赖那使死人复活的天主，}宗徒说，\BibleQuote{他由这样大的死亡中拯救了我们，为使我们的信德不建立在人的智慧上，而建立在天主的德能上。} \BibleQuote{属神的人能审断一切，但他却不受任何人的审断。} 我也听见他的话：\BibleQuote{我说这些话，免得有人用花言巧语迷惑你们，或有人进来掳掠你们。}\BibleRef{哥 2:4} 又说：\BibleQuote{你们要谨慎，免得有人用哲学和虚妄的欺骗，依照人的传授，依据世俗的原理，而不是依据基督，把你们掳去。}\BibleRef{哥 2:8} 这不是斥责一切哲学，而是斥责伊壁鸠鲁派——\Person{保禄}在\WorkTitle{宗徒大事录}中提及的\BibleRef{宗 17:18}——这派否认天主的眷顾，尊崇享乐为神，以及任何其他崇敬元素，却不将动力因置于元素之上，也不认识造物主。\par

他也提到的斯多葛派，他们说天主是物体，渗透在最卑劣的物质中，这说得不好。他把逻辑的戏法称为\Quote{人的传授。}因此他也加上：\Quote{要避免幼稚的问题，因为这样的争辩是幼稚的。}\Quote{德性不爱少年，}哲学家\Person{柏拉图}说。而且，据\Person{莱昂提尼的高尔吉亚}说，我们的搏斗需要两种德性：勇敢和智慧——勇敢以面对危险，智慧以理解奥秘。因为圣言如同奥林匹亚公告：呼唤愿意的人，并加冕那些能够坚定不移直至真理的人。事实上，圣言不希望那信了的人懒惰。因为他说：\BibleQuote{你们寻找，就必找到。}\BibleRef{玛 7:7} 但寻求以找到为结束，驱除虚妄的琐事，并认可那确认我们信德的默观。\BibleQuote{我说这话，免得有人用花言巧语迷惑你们，}\BibleRef{哥 2:4}宗徒说，他显然学会了辨别他所说的话，并受教以应对异议。\BibleQuote{你们既然接受了基督耶稣为主，就当在他内行走，在他内生根并被建造，在信德中得以坚定。}\BibleRef{哥 2:6} 现在，说服就是在信德中得以坚定的方法。\BibleQuote{你们要谨慎，免得有人用哲学和虚妄的欺骗，剥夺你们对基督的信德，}这哲学否认天主的眷顾，\Quote{是按照人的传授；}因为那合乎神圣传授的哲学确立并证实天主的眷顾；如果眷顾被否认，救主的救恩计划就成为一个神话，而我们则\Quote{按照世俗的元素，而不是按照基督}受到的影响。\BibleRef{哥 2:8} 因为那合乎基督的教导神化造物主，在具体事件中追溯天主的眷顾，知道元素的本性是可以变化和产生的，并教导我们应当力求上升到那与天主相似的德能，并优先选择那占据首要地位并超乎一切训练的救恩计划。\par

元素被崇拜——\Person{第欧根尼}崇拜空气，\Person{泰勒斯}崇拜水，\Person{希帕索斯}崇拜火；还有那些认为原子是万物本原的人，妄称哲学家之名，实则是沉溺于享乐的可怜虫。\Quote{因此我祈求，}宗徒说，\Quote{使你们的爱德越发增多，充满知识和各种识见，好使你们能辨别卓绝之事。}\BibleRef{斐 1:9} \Quote{自从我们做孩童时，}同一位宗徒说，\Quote{我们就受世界的基本要素奴役。那孩子虽然是继承人，却与奴仆毫无区别，直到父亲指定的时间。} 如此看来，哲学家就是孩童，除非他们被基督造就成人。\BibleQuote{因为如果婢女的儿子不能与自由妇人的儿子一同继承产业，}\BibleRef{创 21:10} 至少他是亚巴郎的后裔，虽非应许之子，却能以无偿恩赐获得应有之物。\BibleQuote{但干粮属于成年人，就是那些由于习惯，感官得以锻炼，能辨别善恶的人。}\BibleRef{希 5:14} \BibleQuote{因为凡是吃奶的，都不熟练正义之言，因为他是婴孩，}\BibleRef{希 5:13} 尚未熟悉那使他相信并行动的圣言，也不能在心中给出理由。\Quote{凡事考验，}宗徒说，\Quote{持守良善的，}\BibleRef{得前 5:21} 这是对属神的人说的，他们依真理判断所讲论的一切，不论其看似真实还是确实持守真理。\Quote{不受纪律纠正的人犯错，鞭打和责备赐予智慧的纪律，}显然是指出于爱的责备。\BibleQuote{正直的心寻求知识。}\BibleRef{箴 15:14} \Quote{因为寻求主的，必寻得知识与正义；正确寻求的，已寻得平安。} \Quote{我要知道的，}经上说，\Quote{不是那些自高自大之人的言语，而是能力。} 这是斥责那些外表聪明、自以为聪明、实际不聪明的人，他写道：\Quote{因为天主的国不在于言语，}\BibleRef{格前 4:19} 不在于虚假之物，只在于出于意见的可能性；但他说\Quote{在于能力，}因为只有真理是强有力的。又说：\Quote{如果有人自以为知道什么，他还不知道他该知道的。} 因为真理绝不只是意见。但\Quote{知识的臆想令人膨胀，}充满骄傲，\Quote{但爱德却建立人，}它不依赖臆想，而依赖真理。因此经上说：\Quote{如果有人爱，他就被认识。}\par

% structure: p0086 source=h2 target=subsection reason=relative-heading-rank; base=h2

\subsection{第十二章 信仰的奥秘不可向所有人揭示}

既然这传统并非仅为领悟圣言之宏伟的人独自公布，因此有必要将天主圣子所教导的智慧隐藏在奥秘之中。如今，依撒意亚先知以火净化他的舌头，以便能讲述神视。我们也必须净化，不仅是舌头，还有耳朵，如果我们试图参与真理。\par

这些是我写作时遇到的障碍。如今我仍然害怕，正如经上说：\BibleQuote{不要将珍珠投在猪面前，免得它们用脚践踏，并且转过来撕碎你们。}\BibleRef{玛 7:6} 因为向猪猡般未经训练的听众展示关于真光的真正纯洁清晰的言语是很困难的。因为几乎没有什么比这些对他们的听闻更可笑，另一方面，也没有什么比这些对高尚本性的人更令人赞叹或更鼓舞人心。\BibleQuote{但属血气的人不接受天主圣神的事，因为在他看来是愚妄。}\BibleRef{格前 2:14} 但明智的人不会用口说出他们在商议中所推理的。\Quote{你们在暗中听到的，}主说，\Quote{要在屋顶上宣讲出来。}\BibleRef{玛 10:27} 命令他们接受真知的秘密传统，并在高处公开讲解；并且我们如何从耳中听到，就照样传授给那些需要的人；但不是命令我们不加区别地告诉所有人那以比喻对他们所说的话。笔记中只有一种描绘，真理被稀疏而广泛地播撒，以躲过那些像寒鸦般捡拾种子的人的注意；但一旦找到好农夫，它们每一颗都会发芽并结出谷粒。\par

% structure: p0089 source=h2 target=subsection reason=relative-heading-rank; base=h2

\subsection{第十三章 所有哲学派别都包含真理的萌芽}

因此，既然真理是唯一的（因为虚假有万千歧途）；正如酒神女信徒撕碎\Person{彭透斯}的肢体，野蛮人和希腊人的哲学派别对真理所做的也是如此，各自夸耀其分得的一份为全部真理。但依我之见，众人都被光之晨曦照亮。因此，让所有追求真理的人——无论是希腊人还是野蛮人，无论是拥有不少真理还是只拥有一部分的人——都献出他们所拥有的真理之言。\par

永恒，例如，在一瞬间呈现未来的和现在的，以及过去的时间。但真理，远比无限的绵延更强有力，能够收集它自己的种子，即使它们已落在异土。因为我们将发现，那些尚未变得完全无知、也未被从自然秩序中切割出去（像寓言中的妇人肢解那人一样，切割基督）的教派所持的许多教义，尽管表面互不相同，但在其起源上与整个真理相符。因为它们在一件事上相符，或是作为部分，或是作为种类，或是作为属。例如，虽然最高的音符与最低的音符不同，但两者构成一个和谐。在数字中，偶数与奇数不同，但两者都适用于算术；图形也是如此，圆形、三角形、正方形，以及互不相同的任何图形。同样，在整个宇宙中，所有部分虽然彼此不同，却都保持着与整体的关系。因此，蛮族哲学和希腊哲学从永恒真理上撕下了一块碎片——不是从狄俄尼索斯的秘教，而是从永生活言的圣言神学。而那将分离的碎片再次带来并使之成为一的，将毫无危险，确定无疑，默观那圆满的圣言，即真理。因此，在\WorkTitle{训道篇}中记载：\BibleQuote{我比所有在我以前在耶路撒冷的人更增添了智慧；我的心见了许多事物；此外，我认识了智慧和知识、比喻和了解。这也是精神的选择，因为智慧的丰盈就是知识的丰盈。}与各种智慧相熟悉的人，将是一位卓越的\OriginalTerm{真知者}{gnostic}。又记载：\BibleQuote{智慧知识的丰盈将赐生命给拥有它的人。}再者，这句话更清楚地确认了所说的话：\BibleQuote{一切都在明悟者眼前}——一切事物，包括希腊的和蛮族的；但二者之一并非全部。\BibleQuote{它们对于愿意接受了解的人是正确的。选择教训，而非银子，选择知识，而非精炼的金子，}并更愿意以理智换取纯金；\BibleQuote{因为智慧比宝石更好，任何珍贵之物都不配与它相比。}\par

% structure: p0092 source=h2 target=subsection reason=relative-heading-rank; base=h2

\subsection{第十四章 希腊哲学家传承}

希腊人说，在\Person{俄耳甫斯}、\Person{利诺斯}以及他们当中出现的远古诗人之后，所谓的七贤是最早因其智慧而受敬仰的。其中四位来自亚细亚：\Person{米利都的泰勒斯}、\Person{普里耶涅的彼亚斯}、\Person{米提利尼的皮塔库斯}、\Person{林多斯的克莱俄布卢}；两位来自欧罗巴：\Person{雅典的梭伦}、\Person{斯巴达的喀隆}；第七位，有些人说是\Person{科林斯的佩里安德}，另一些人说是\Person{锡西厄的阿那卡西尔斯}，还有一些人说是\Person{克里特的厄庇米尼得斯}——\Person{保禄}认识他为一位希腊先知，他在\WorkTitle{弟铎书}中提及他，这样说：\BibleQuote{他们中间有一位自己的先知说：\Emph{克里特人常说谎，是恶兽，是贪馋的肚腹。}这见证是真的。}\BibleRef{铎 1:12}你看，他甚至将某些真理归给希腊的先知，并不羞于在为了某些人的建树、为了另一些人的羞愧而论述时，使用希腊诗歌。同样，在\WorkTitle{格林多前书}中（这并非唯一一次），他论述死人复活时，使用了一行悲剧的抑扬格诗句，他说：\BibleQuote{如果死人没有复活，那我有什么益处呢？让我们吃喝吧，因为明天我们就要死了。不要受骗；恶友败坏善德。}\BibleRef{格前 15:32}其他人将\Person{阿尔哥斯的阿库西劳斯}列为七贤之一；另一些人将\Person{叙罗斯的斐瑞居德}列在其中。\Person{柏拉图}则用\Person{克尼的米松}替代了\Person{佩里安德}，他认为\Person{佩里安德}因其僭主统治而不配得到智慧。希腊中的智者鼎盛于摩西之后，稍后将证明。但他们中间哲学的风格，如同希伯来式和隐晦式的，现在要考虑。他们采用简洁性，适合劝勉，并且最有用处。甚至\Person{柏拉图}也说，古时这种风格在全体希腊人中普遍流行，特别是在斯巴达人和克里特人中间，他们享有最好的法律。\par

这句话，\Quote{认识你自己，}有些人认为是\Person{喀隆}的。但\Person{卡梅利翁}在他的书\WorkTitle{论神明}中将其归于\Person{泰勒斯}；\Person{亚里士多德}归于\Place{皮提亚}。这可能是一条追求知识的命令。因为不认识整体的本质，就不能认识部分；人必须研究宇宙的生成，以便我们能够学习人的本性。同样，对于\Person{斯巴达的喀隆}，他们归给他\Quote{凡事勿过度。}\Person{斯特拉托}在他的书\WorkTitle{论发明}中，将这格言归于\Person{泰格亚的斯特拉托德穆斯}。\Person{狄迪莫斯}将其归于\Person{梭伦}；如同归于\Person{克莱俄布卢}的那句话：\Quote{中庸为佳。}还有一句话，\Quote{立下誓约，祸患将至，}\Person{克莱奥梅尼}在他的书\WorkTitle{论赫西俄德}中说，\Person{荷马}在先前的诗句中已经说了：—\par

\Quote{悲惨的誓约，只为悲惨的人而立。}\par

\Person{亚里士多德学派}认为这是\Person{奇伦}的箴言；但\Person{狄迪莫斯}说这劝戒是\Person{泰勒斯}的。接着，那则格言\Quote{凡人都恶}或\Quote{大多数人恶}（因为同一句箴言有两种说法），\Place{拜占庭}的\Person{索塔德斯}说这是\Person{比阿斯}的。而那格言\Quote{操练征服一切}，他们认为是\Person{佩里安德尔}的；同样，劝戒\Quote{识时务}，被认为是\Person{皮塔科斯}的说法。\Person{梭伦}为\Place{雅典}人立法，\Person{皮塔科斯}为\Place{米蒂利尼}人立法。后来，\Person{毕达哥拉斯}，\Person{费雷西德斯}的弟子，第一个称自己为哲学家。因此，在上述三人之后，出现了三个哲学学派，以其所在地命名：\Place{意大利}学派源于\Person{毕达哥拉斯}，\Place{伊奥尼亚}学派源于\Person{泰勒斯}，\Place{埃利亚}学派源于\Person{色诺芬尼}。\Person{毕达哥拉斯}是\Place{萨摩斯}人，\Person{姆涅萨尔科斯}的儿子，据\Person{希波博图斯}说；据\Person{阿里斯托克塞诺斯}在其所著关于\Person{毕达哥拉斯}、\Person{阿里斯塔库斯}和\Person{特奥蓬波斯}的传记中，他是\Place{托斯卡纳}人；据\Person{尼安特斯}说，他是\Place{叙利亚}人或\Place{提尔}人。所以，\Person{毕达哥拉斯}据大多数人说是蛮族出身。\Person{泰勒斯}，据\Person{利安德}和\Person{希罗多德}记述，是\Place{腓尼基}人；有些人认为，是\Place{米利都}人。他似乎是唯一见过\Place{埃及}先知的人。但无人被描述为他的老师，也没有人被认为是\Place{叙罗斯}的\Person{费雷西德斯}的老师，\Person{费雷西德斯}收\Person{毕达哥拉斯}为徒。但\Place{意大利}哲学，即\Person{毕达哥拉斯}学派，在\Place{意大利}的\Place{梅塔蓬图姆}渐渐老去。\Place{米利都}的\Person{阿那克西曼德}，\Person{普拉克西亚德斯}的儿子，继承了\Person{泰勒斯}；之后是\Place{米利都}的\Person{阿那克西美尼}，\Person{欧吕斯特拉托斯}的儿子；再后来是\Place{克拉佐美奈}的\Person{阿那克萨戈拉}，\Person{赫格西布洛斯}的儿子。他将他的学派从\Place{伊奥尼亚}迁至\Place{雅典}。他由\Person{阿凯劳斯}继承，\Person{苏格拉底}是\Person{阿凯劳斯}的弟子。\par

\Quote{避开了这些人，那石匠；\newline{}谈论法律者；希腊人的蛊惑者。}\par

据\Person{蒂蒙}在其\WorkTitle{讽刺诗集}中说，因为他放弃了物理学而转向伦理学。\Person{安提斯泰尼}在成为\Person{苏格拉底}的弟子之后，引入了犬儒哲学；\Person{柏拉图}则退入学园。\Person{亚里士多德}在师从\Person{柏拉图}研究哲学之后，退入\Place{吕刻昂}，创立了逍遥学派。他由\Person{泰奥弗拉斯托斯}继承，\Person{泰奥弗拉斯托斯}由\Person{斯特拉托}继承，\Person{斯特拉托}由\Person{吕孔}继承，然后\Person{克里托拉俄斯}，然后\Person{狄奥多罗斯}。\Person{斯彪西波}是\Person{柏拉图}的继承人；他的继承人是\Person{色诺克拉底}；后者的继承人是\Person{波莱莫}。\Person{波莱莫}的弟子是\Person{克拉特斯}和\Person{克兰托尔}，由\Person{柏拉图}创立的旧学园至此终止。\Person{阿塞西劳斯}是\Person{克兰托尔}的同僚；从他开始，直到\Person{赫格西劳斯}，中学园兴盛。然后\Person{卡尔内亚德}继承了\Person{赫格西劳斯}，后继者接续而来。\Person{克拉特斯}的弟子是\Place{基提翁}的\Person{芝诺}，斯多葛学派的创始人。他由\Person{克莱安西斯}继承；后者由\Person{克吕西波}继承，之后还有其他人。\Place{科洛封}的\Person{色诺芬尼}是埃利亚学派的创始人，据\Person{蒂迈欧}说，他生活在\Place{西西里}的君主\Person{希伦}和诗人\Person{埃庇卡摩斯}的时代；\Person{阿波罗多罗斯}说，他生于第四十奥林匹克赛会，一直活到大流士和居鲁士的时代。因此，\Person{巴门尼德}是\Person{色诺芬尼}的弟子，\Person{芝诺}是他的弟子；然后来了\Person{留基伯}，然后是\Person{德谟克利特}。\Person{德谟克利特}的弟子是\Place{阿布德拉}的\Person{普罗泰戈拉}和\Place{希俄斯}的\Person{梅特罗多罗斯}，后者的弟子是\Place{士麦那}的\Person{第欧根尼}；他的再传弟子是\Person{阿那克萨尔科斯}，他的又是\Person{皮浪}，他的又是\Person{瑙西法涅斯}。有人说\Person{伊壁鸠鲁}是他的学生。\par

以上是希腊哲学家的传承概要。现在需要依次确定他们哲学创始人的年代，以便通过比较显示\Place{希伯来}哲学更早许多世代。\par

据说\Person{色诺芬尼}是\Place{埃利亚}哲学的创始人。\Person{欧德摩斯}在\WorkTitle{星象史}中说，\Person{泰勒斯}预言了日食，发生于\Place{米底}人和\Place{吕底亚}人交战时，当时\Person{基亚克萨雷斯}（\Person{阿斯提阿格斯}的父亲）统治\Place{米底}，\Person{阿吕阿特斯}（\Person{克洛伊索斯}的儿子）统治\Place{吕底亚}。\Person{希罗多德}在其第一卷中同意他的说法。日期大约在第五十奥林匹克赛会。\Person{毕达哥拉斯}被确定为生活在僭主\Person{波吕克拉特斯}时代，大约第六十二奥林匹克赛会。\Person{姆涅西菲洛斯}被描述为\Person{梭伦}的追随者，与\Person{泰米斯托克利}是同时代人。因此\Person{梭伦}兴盛于大约第四十六奥林匹克赛会。\Person{赫拉克利特}，\Person{包索}的儿子，说服僭主\Person{墨兰科马斯}放弃王权。他藐视\Person{大流士}王的邀请，拒绝前往\Place{波斯}。\par

% structure: p0101 source=h2 target=subsection reason=relative-heading-rank; base=h2

\subsection{希腊哲学大部分源自蛮族}

这些是希腊最古老智者和哲学家的年代。他们大多数是蛮族出身，并且在蛮族中受训，这有什么必要说呢？\Person{毕达哥拉斯}被证明是\Place{托斯卡纳}人或\Place{提尔}人。\Person{安提斯泰尼}是\Place{弗里吉亚}人。\Person{俄耳甫斯}是\Place{奥德里西亚}人或\Place{色雷斯}人。大多数人还表明\Person{荷马}是\Place{埃及}人。\Person{泰勒斯}生为\Place{腓尼基}人，据说他与\Place{埃及}的先知交往；同样\Person{毕达哥拉斯}也与同一批人交往，并被他们施行割礼，以便进入至圣所，从埃及人那里学习神秘哲学。他与\Place{迦勒底}人和术士的首领交谈；他暗示了如今所谓的教会，在他所维持的公共大厅中。\par

柏拉图并不否认他从外邦人那里获得了哲学中最卓越的部分；他承认自己曾到过埃及。因此，在 \WorkTitle{斐多篇} 中，他写道哲学家可以从各方面得到帮助，他说：\Quote{克贝啊，希腊确实广大，到处都有善人，外邦人的种族也很多。} 所以柏拉图认为有些外邦人也是哲学家。但相反，伊壁鸠鲁认为只有希腊人才能从事哲学。而在 \WorkTitle{会饮篇} 中，柏拉图称赞外邦人显著地实践哲学，真实地说：\Quote{还有许多其他例子，无论是在希腊人中还是在外邦人中，为这样的儿子所建的庙宇已经很多了。} 显然，外邦人极大地尊崇他们的立法者和导师，称他们为神。因为，根据柏拉图，\Quote{他们认为善灵魂离开超天区域后，会来到这个塔尔塔罗斯，并披上身体，分担由出生带来的所有苦难，这是出于他们对人类种族的关怀；} 这些灵魂制定法律，传播哲学，\Quote{从神那里从未有过，也将不会有比这更大的恩赐给人类。}\par

在我看来，正是由于认识到通过智者所带来的巨大益处，智者本身受到尊崇，哲学也被所有婆罗门、奥德里西人和格泰人公开地实践。这些人被埃及人、迦勒底人、称为幸福的阿拉伯人、以及居住在巴勒斯坦的人、波斯人的很大一部分、以及除此之外的无数其他种族严格地奉为神。众所周知，柏拉图不断颂扬外邦人，记得他和毕达哥拉斯都是从外邦人那里学到了他们最重要、最高贵的教义。因此，他在 \WorkTitle{斐德若篇} 中称外邦人的种族为\Quote{外邦哲学家的种族}，并认出埃及王，向我们展示他比透特更智慧，他知道透特就是赫尔墨斯。但在 \WorkTitle{卡尔米德篇} 中，很明显他知道某些色雷斯人据说能使灵魂不朽。据报道，毕达哥拉斯是埃及大先知松切斯的弟子；柏拉图是赫利奥波利斯的塞克努菲斯的弟子；欧多克索斯是克尼多斯的科努菲斯的弟子，后者也是埃及人。在他的 \WorkTitle{论灵魂} 一书中，柏拉图再次明确承认预言，当他引入一位先知宣告拉刻西斯的言语，向那些命运已定的灵魂发出预言。在 \WorkTitle{蒂迈欧篇} 中，他引入了非常智慧的梭伦，从外邦人那里学习。那宣告的要旨如下：\Quote{梭伦啊，梭伦，你们希腊人总是孩子。没有希腊人是老人。因为你们没有因年岁而灰白的学问。}\par

德谟克利特借鉴了巴比伦的伦理论述，据说他在自己的作品中结合了阿基卡鲁斯石柱的翻译。你可以发现他写作时标明出处：\Quote{德谟克利特如是说。} 关于他自己，当他夸耀自己的博学时，他说：\Quote{我游历了比我同代任何人更广的土地，考察了最偏远的地区。我见过最多的天空和陆地，听过极多的有学问的人。在创作上，没有人超越我；在论证上，甚至那些埃及人中被称为阿耳珀诺达普特的人也没有，我与他们所有人一起流亡了八十年。} 因为他去了巴比伦、波斯和埃及，向玛吉和祭司学习。\par

毕达哥拉斯表明玛吉琐罗亚斯德是波斯人。那些追随普罗迪库斯异端的人自称拥有此人的秘传书籍。亚历山大在其 \WorkTitle{论毕达哥拉斯象征} 一书中叙述，毕达哥拉斯是亚述人纳扎拉图斯（有人认为是厄则克耳，但并非如此，后文将说明）的学生，并断言除此之外，毕达哥拉斯还听过加拉太人和婆罗门的教诲。逍遥派克莱尔库斯说他认识一位与亚里士多德交往的犹太人。赫拉克利特说，西彼拉说话不是出于人为，而是出自天主的帮助。因此，据说在德尔斐，神谕旁有一块石头，据说第一位西彼拉曾坐在上面，她来自赫利孔山，由缪斯抚养长大。但有人说她来自米利亚，是西顿的拉弥亚的女儿。塞拉皮翁在其史诗中说，西彼拉即使在死后也没有停止预言。他写道：她死后进入空气的东西以声音和预兆的方式发出神谕；她的身体变为泥土，草自然地从中长出，任何恰好在那个地方吃草的野兽，它们的五脏六腑都向人精确地显示未来的知识。他还认为月亮中看到的脸就是她的灵魂。关于西彼拉，就说这么多。\par

罗马王努马是毕达哥拉斯学派，并借助梅瑟的训诫，禁止制作人形或活物形状的天主像。因此，在头一百七十年里，虽然建造了圣殿，但他们没有铸造或雕刻任何像。因为努马秘密地告诉他们，至善的存在只能通过心智来领悟。因此，哲学，一件极其有益的事物，在古代外邦人中兴旺发达，将光芒普照万民。后来它传到了希腊。在其行列中，首先有埃及的先知；亚述人中的迦勒底人；高卢人中的德鲁伊特；巴克特里亚人中的萨马纳人；凯尔特人的哲学家；波斯人的玛吉，他们预言了救主的诞生，并在星光的引导下进入犹太地。印度的祼形智者也在其中，以及其他外邦哲学家。他们有两类，一些称为萨尔马那人，另一些称为婆罗门。萨尔马那中称为希洛比的人，既不居住在城市，也没有屋顶遮蔽，穿着树皮，以坚果为食，用手捧水喝。像今天被称为\Quote{刻苦者}的人一样，他们不知道婚姻和生育子女。\par

有些印度人也遵奉\Person{佛陀}的训诫；因他超凡的圣德，他们将他提升至神明的尊荣。\par

\Person{阿那卡尔西斯}是斯库提亚人，据记载，他胜过许多希腊哲学家。\Person{赫拉尼科斯}记载说，极北族人居住在\Place{里佩山}之外，他们倡导正义，不吃肉，只吃坚果。他们把年满六十岁的人带到城门外，将其处死。在日耳曼人中也有被称为圣女的妇女，她们通过察看河流的漩涡和回水，并倾听溪流的声音，预言和预测未来之事。这些妇女不允许男子在新月出现之前与\Person{凯撒}作战。\par

在这些民族中，犹太民族是最古老的；而毕达哥拉斯派的\Person{斐罗}详尽地证明，他们那成文的哲学优先于希腊哲学；此外，逍遥学派的\Person{阿里斯托布鲁}以及另外几位——为免浪费时间，不逐一列举名字——也证明了这一点。与\Person{塞琉古·尼卡诺尔}同时代的作者\Person{麦加斯梯尼}在其著作\WorkTitle{印度事务}第三卷中非常清楚地写道：\Quote{古人关于自然所说的一切，那些在希腊之外从事哲学的人也都说过：其中一些是印度婆罗门所说的，另一些是叙利亚所谓的犹太人所说的。}还有一些人更加虚构地说，所谓的伊达克泰利人中的某些人是第一批智者；据说他们发明了所谓的\Quote{\OriginalTerm{厄弗所字母}{Ephesian letters}}和音乐中的数字。因此，音乐中的达克提利（dactyls）得名。而伊达克泰利人是弗里吉亚人和蛮族人。\Person{希罗多德}记载，\Person{赫拉克勒斯}在成为哲人和物理学研究者之后，从弗里吉亚蛮族\Person{阿特拉斯}那里接受了宇宙的柱石；这个寓言的意思是，他通过教导获得了关于天体的知识。\Person{贝里图的赫米普斯}称\Person{半人马卡戎}为智者；关于他，那位创作\WorkTitle{泰坦之战}的人说：\Quote{他首先引导凡人种族走向正义，教导他们誓言的庄严、赎罪祭献和奥林帕斯的形象。}\Person{阿喀琉斯}（他曾参加特洛伊战争）曾受教于他。而半人马卡戎的女儿\Person{希波}，她与\Person{埃俄洛斯}同住，教给了他她父亲的学问，即物理学知识。\Person{欧里庇得斯}也这样为\Person{希波}作证：——\par

\Quote{他首先通过神谕，\newline{}并以星辰升起，预言神圣之事。}\par

这位\Person{埃俄洛斯}在\Place{特洛伊}被攻陷之后接待了\Person{奥德修斯}作为客人。要将这些时代与\Person{摩西}的时代，以及与他所传播的哲学的高古相比较，来标记这些时代。\par

% structure: p0113 source=h2 target=subsection reason=relative-heading-rank; base=h2

\subsection{第十六章 其他技艺的发明者大多是蛮族人}

蛮族不仅是哲学的发明者，也几乎是所有技艺的发明者。\Place{埃及}人最先在世人中引入占星术。\Place{加色丁}人也是如此。\Place{埃及}人首先示点燃灯盏之法，又将一年分为十二个月，禁止在圣殿内与妇人交合，并规定凡与妇人同寝者，不沐浴不得进入圣殿。再者，他们是几何学的发明者。有些人说，\Place{卡里亚}人发明了星象预兆。\Place{弗里吉亚}人最先关注鸟类的飞行。而\Place{意大利}的邻族\Place{图斯卡}人精于肠卜之术。\Place{伊扫里亚}人和\Place{阿拉伯}人发明了占卜，\Place{特尔梅索斯}人则发明了梦占。\Place{厄特鲁里亚}人发明了号角，\Place{弗里吉亚}人发明了笛子。因为\Person{奥林普斯}和\Person{马尔叙亚斯}都是\Place{弗里吉亚}人。据\Person{欧福洛斯}所说，在\Place{希腊}人中发明文字的是\Person{加德摩斯}，他是\Place{腓尼基}人；因此\Person{赫罗多托斯}也写道，这些文字被称为\Place{腓尼基}文字。又有人说，\Place{腓尼基}人和\Place{叙利亚}人最先发明了文字；而\Person{伊娥}来到\Place{埃及}之前，\Place{埃及}土著\Person{阿彼斯}发明了医术。但后来据说\Person{阿斯克勒庇俄斯}改进了这项技艺。\Place{利比亚}人\Person{阿特拉斯}最先建造船只并航海。\Person{克尔米斯}和\Person{达玛内乌斯}，\Place{伊达山}的达克提利人，最先在\Place{塞浦路斯}发现了铁。另一位\Place{伊达山}人发现了黄铜的淬火；据\Person{赫西俄德}说，是一个\Place{斯基泰}人。\Place{色雷斯}人最先发明了所谓的\OriginalTerm{弯刀}{\GreekText{ἅρπη}}——是一种曲剑——并最先在马背上使用盾牌。同样，\Place{伊利里亚}人发明了\OriginalTerm{盾牌}{\GreekText{πέλτη}}。此外，据说\Place{图斯卡}人发明了制陶术；\Person{伊塔努斯}（他是\Place{萨莫奈}人）最先制作了\OriginalTerm{长盾}{\GreekText{θυρέος}}。\Place{腓尼基}人\Person{加德摩斯}发明了石雕，并发现了\Place{潘盖翁}山上的金矿。此外，另一个民族\Place{卡帕多奇亚}人最先发明了称为纳布拉的乐器，\Place{亚述}人同样发明了双弦琴。\Place{迦太基}人最先建造了三层桨战船；由土著\Person{波福罗斯}建造。\Place{科尔基斯}人\Person{埃厄忒斯}的女儿\Person{美狄亚}最先发明了染发。此外，\Place{诺罗佩斯}人（他们是\Place{佩奥尼亚}种族，现称\Place{诺里基}人）加工铜，并最先精炼铁。\Place{贝布里西}人的国王\Person{阿米库斯}最先发明了拳击手套。在音乐方面，\Place{密细亚}人\Person{奥林普斯}使用了吕底亚和声；称为穴居人的民族发明了桑布卡琴。据说弯笛是\Place{弗里吉亚}人\Person{萨提洛斯}发明；同样，自然音阶和声也是\Place{弗里吉亚}人\Person{希阿格尼斯}发明；音符由\Place{弗里吉亚}人\Person{奥林普斯}发明；弗里吉亚和声、半弗里吉亚和半吕底亚和声则由与上述人物同属一地的\Person{马尔叙亚斯}发明。多利亚和声由\Place{色雷斯}人\Person{塔米里斯}发明。我们听说，\Place{波斯}人最先制造马车、床和脚凳；\Place{西顿}人最先建造了三层桨战船。靠近\Place{意大利}的\Place{西西里}人最先发明了福明克斯琴，它不比里拉琴逊色多少。他们还发明了响板。在\Place{亚述}女王\Person{塞米拉米斯}时代，据说亚麻衣物被发明。\Person{赫拉尼库斯}说，\Place{波斯}王后\Person{阿托萨}最先写了一封信。这些事由\Place{米蒂利尼}的\Person{斯卡蒙}、\Place{厄弗所}的\Person{泰奥弗拉斯托斯}、\Place{曼提尼亚}的\Person{库狄普斯}，以及\Person{安提法奈斯}、\Person{阿里斯托德摩斯}和\Person{亚里士多德}记载；此外还有\Person{菲洛斯特法诺斯}，以及逍遥学派的\Person{斯特拉托}在其著作\WorkTitle{论发明}中提及。我从他们那里引用了一些细节，以证实蛮族那富于创造且实际有用的才智，希腊人曾受益于他们的研究。若有人反对蛮族的语言，\Person{阿那卡西斯}说：\Quote{所有希腊人对我讲的都是斯基泰语。}正是这位令希腊人钦佩的人曾说：\Quote{我的衣服是斗篷；我的晚餐是奶和酪。}你可见蛮族的哲学注重行动而非言语。宗徒如此说：\BibleQuote{同样，你们若不用舌头说出容易明白的话，怎能知道所说的是什么呢？你们将是对空气说话。世上或许有种种声音，但没有一种是没有意义的。因此，我若不晓得那声音的意义，我就成了对说话者说话的蛮族，而说话者也成了给我说话的蛮族。}又说：\BibleQuote{说语言的，当求能解释。}\par

不仅如此，论述的传讲与写作传到希腊，已是较晚的事情。\Place{克罗顿}的\Person{佩里图斯}之子\Person{阿尔克迈翁}，首先撰写了论自然的论文。又有记载说，\Place{克拉佐美尼}的\Person{黑格西布洛斯}之子\Person{阿那克萨戈拉}，第一个出版了成书的著作。第一个将音乐配合诗歌创作的是\Place{安提萨}的\Person{特尔潘德}，他还将斯巴达人的法律谱上乐曲。\Place{赫耳米奥涅}的\Person{拉索斯}发明了狄提兰博斯；\Place{希梅拉}的\Person{斯特西科罗斯}发明了赞美歌；斯巴达人\Person{阿尔克曼}发明了合唱曲；\Place{忒奥斯}的\Person{阿那克里翁}发明了情歌；\Place{忒拜}人\Person{品达}发明了伴舞歌。\Place{米利都}的\Person{提摩泰}是第一个用里拉琴演奏称为\OriginalTerm{乐歌}{\GreekText{νόμοι}}的音乐作品并伴以舞蹈的人。此外，抑扬格由\Place{帕洛斯}的\Person{阿喀罗科斯}发明，而跛足抑扬格由\Place{以弗所}的\Person{希波那克斯}发明。悲剧起源于雅典人\Person{忒斯庇斯}，喜剧起源于\Place{伊卡里亚}的\Person{苏萨里翁}。它们的年代由文法家们流传下来。但要准确说明它们，则过于繁琐；稍后，我们将表明，庆祝狄奥尼索斯节庆的\Person{狄奥尼索斯}，比\Person{摩西}更晚。他们说，\Place{拉姆诺斯}的\Person{索菲罗斯}之子\Person{安提丰}，首先发明了学园论述和修辞格，并且是第一个为报酬而辩护、并撰写供当庭发表的法庭演说辞的人，如\Person{狄奥多罗斯}所说。\Place{库迈}的\Person{阿波罗多罗斯}首先采用了\Quote{批评家}这一名称，并被称为文法家。有些人说，\Place{昔兰尼}的\Person{埃拉托斯特尼}是第一个获此称呼的人，因为他出版了两本书，名为\WorkTitle{Grammatica}。第一个被称为文法家（正如我们现在使用该术语）的人，是\Place{米提利尼}的\Person{狄奥尼索芬尼斯}之子\Person{普拉克西法涅斯}。\Place{洛克里斯}的\Person{泽琉科斯}据说是第一个（以书面形式）制定法律的人。另有人说，是\Person{林叩斯}时代\Person{宙斯}之子\Person{墨诺斯}。他在\Person{达那俄斯}之后，从\Person{伊那科斯}和\Person{摩西}算起第十一代；我们稍后将予以说明。\Person{来库古}在特洛伊陷落多年后生活，在第一次奥林匹亚竞技会之前一百五十年为斯巴达人立法。我们之前已经谈到过\Person{梭伦}的年代。\Person{德拉古}（他也是一位立法者）被发现生活在大约第三〇九届奥林匹亚竞技会期间。此外，\Person{安提洛科斯}，他撰写了从\Person{毕达哥拉斯}时代到\Person{伊壁鸠鲁}之死的学者列传，\Person{伊壁鸠鲁}死于盖墨利翁月第十天，共计算了三百一十二年。此外，有人说，\Person{伊卡里俄斯}之妻\Person{法诺忒亚}发明了英雄史诗六音步；另有人说，是提坦女神之一的\Person{忒弥斯}。然而，\Person{狄迪莫斯}在他的著作\WorkTitle{论毕达哥拉斯哲学}中叙述说，\Place{克罗顿}的\Person{忒阿诺}是第一位从事哲学并创作诗歌的女性。因此，按照某些人的说法，希腊哲学偶然地、模糊地、部分地把握了真理；而按照另一些人的说法，哲学是由魔鬼推动的。有几人认为，某些来自天上的力量启发了整个哲学。但是，如果希腊哲学并不包含真理的全部，并且缺乏实践主之诫命的力量，那么它仍为真正君王的教导预备了道路；它以某种方式训练、塑造品格，并使相信天主眷顾的人适合接受真理。\par

% structure: p0116 source=h2 target=subsection reason=relative-heading-rank; base=h2

\subsection{第十七章 论救主的话：\Quote{凡在我以前来的都是贼和强盗。}}

但是，他们说，经上记载：\Quote{凡在主来临以前的都是贼和强盗。}那么，凡在圣言内的（因为这些是在圣言降生成人以前的人）都可被一般地理解为这样。但是先知们，既然受主派遣和默感，就不是贼，而是仆人。因此圣经说：\BibleQuote{智慧派遣了她的使女，大声呼唤，邀请到酒杯旁。}\BibleRef{箴 9:3}\par

但是，据说，哲学不是主派遣的，而是偷来的，或由贼给的。那么，是某个力量或天使，学到了真理的一点，但没有停留在真理中，启发并教导了这些事情，不是主不知道，主事先知道每个本质的构成和未来的结果，而是没有禁止。\par

因为当时发生在人间的盗窃，有一些益处；但并非行窃者着眼于利益，而是天主眷顾将这一大胆行为的结果导向益处。我知道许多人不断用这样的指控攻击我们，说没有阻止事情发生，就是事情发生的原因。因为他们说，不采取措施防盗或不阻止盗窃的人，就是盗窃的原因；正如没有在起初扑灭火势的人是火灾的原因；没有收帆的船主是海难的原因。当然，法律会惩罚这些事件的肇因。对于有能力阻止的人，他要对发生的事负责。我们对他们说，原因在于行事、工作、行动；而不阻止在此方面是不作为的。此外，原因与活动相关联；例如造船工与船的产生的关系，建筑工与房屋的建造的关系。但不阻止与被发生的事情是分离的。因此，结果将会实现；因为那个原本能阻止的东西既不行动也不阻止。因为不阻止的东西究竟行使了什么活动？如果他们说，创伤的原因不是投枪，而是没有阻止投枪穿透的盾牌；如果他们没有责怪贼，而是责怪没有阻止盗窃的人，那么他们的论断就会沦为荒谬。那么让他们说，烧毁希腊船只的不是\Person{赫克托耳}，而是\Person{阿喀琉斯}；因为他有能力阻止\Person{赫克托耳}却没有阻止；反而由于愤怒（而愤怒与否取决于他自己）没有制止大火，并成为联合原因。如今，魔鬼拥有自由意志，既能悔改也能偷窃；他是盗窃的创造者，而不是没有阻止他的主。但这份礼物并非有害，以致需要干预阻止。\par

但如果要严格准确地处理它们，让他们知道，那并不阻止我们断言在盗窃中所发生的事情，根本不是原因；但阻止的原因被卷入被指控为原因之中。因为用盾牌保护的人是他所保护的人不受伤害的原因；正如他那样阻止他受伤。因为\Person{苏格拉底}的神灵是一个原因，不是通过不阻止，而是通过劝勉，即使（严格来说）他没有劝勉。而且当灵魂没有倾向和反对的权力，而邪恶是不自愿的时候，赞美或谴责，奖赏或惩罚都不正确。因此，阻止者是原因；而不阻止者公正地判断灵魂的选择。所以天主绝不是邪恶的创造者。但因为自由选择和倾向产生罪过，有时错误的判断占上风，由于这是无知和愚昧，我们不愿退避，惩罚便正确地施加。因为发烧是不自愿的；但当一个人因自己的过错，由于过度而发烧时，我们责备他。因此，既然邪恶是不自愿的——因为没有人喜欢邪恶本身；而是被其中的快乐所诱导，并想象它是善的，认为它是可欲的——既然如此，使自己摆脱无知、邪恶和纵欲的选择，尤其是拒绝对那些迷惑人的幻象表示同意，有赖于我们自己。魔鬼被称为\Quote{贼和强盗}；把假先知与真先知混杂在一起，如同稗子与麦子。\Quote{凡在主以前来的，都是贼和强盗；}并非绝对所有的人，而是所有的假先知，以及所有没有被祂恰当派遣的人。因为假先知不诚实地拥有先知的名称，他们是先知，但属于说谎者的先知。因为主说：\BibleQuote{你们是出于你们的父亲魔鬼，并且你们愿意满足你们父亲的欲望。他从起初就是杀人者，不站在真理上，因为在他内没有真理。他撒谎时，是出于自己说的，因为他是撒谎者，也是撒谎者的父亲。}\BibleRef{若 8:44}\par

但在谎言中，假先知也说了一些真实的事。实际上，他们\Quote{在神魂超拔中}说预言，作为叛教者的仆人。而牧者，即悔改的天使，对\Person{赫尔马斯}谈到假先知说：\Quote{因为他讲了一些真理。因为魔鬼用他自己的灵充满他，也许他能把任何人从正道上拉下来。}因此，一切事从天主安排都是为了善，\BibleQuote{为使教会知道天主各种的智慧，按照祂在基督内所定的永远计划。}\BibleRef{弗 3:10}没有东西能抗拒天主：没有东西能反对祂：因为祂是主，是全能的。此外，那些叛逆者的计谋和活动是片面的，出于恶劣的倾向，如同身体疾病出于恶劣的体质，但被普遍的眷顾引导到有益的结局，即使原因是产生疾病的。因此，天主眷顾的最大成就，就是不使从自愿的背叛中产生的邪恶保持无用、毫无益处、并且完全有害。因为天主的智慧、卓越和能力的工程，不仅是行善（因为这样，可以说是天主的本性，如同火加热、光照亮），而尤其是确保那些通过任何人的邪恶所酿成的事情能产生有益和有用的结果，并利用那些看似邪恶的事物，正如由试探而来的见证。\par

因此，在哲学中，虽像\Person{普罗米修斯}偷来的火，却有一线微弱的火花，能够被煽成火焰，一道智慧的痕迹和来自天主的推动。好吧，就算\Quote{那些贼和强盗}是希腊的哲学家，他们在主来临之前从希伯来先知那里得到了一些真理的片段，并不完全认识，却把这些当作自己的教导，隐瞒了一些观点，用他们的聪明才智诡辩地处理了其他观点，并发现了其他事物，因为他们也许曾拥有\BibleQuote{灵明之性}\BibleRef{出 28:3}。\Person{亚里斯多德}也同意圣经，并宣称诡辩术偷窃了智慧，正如我们之前所指出的。宗徒说：\BibleQuote{我们讲说这些事，不是用人的智慧所教导的言语，而是用圣神所教导的言语，}\BibleRef{格前 2:13}因为关于先知们，经上说：\BibleQuote{我们都从祂的圆满中领受了，}\BibleRef{若 1:16}即从基督的。所以先知们不是贼。\Quote{我的教导不是我的，}主说，\Quote{而是那派遣我来的父的。}至于那些偷窃的人，祂说：\Quote{但那凭自己说话的，寻求自己的荣耀。}希腊人就是如此，\BibleQuote{爱自己，自夸自大，}\BibleRef{弟后 3:2}。当圣经谈到这些人有智慧时，并不是指责真正有智慧的人，而是指责那些外表有智慧的人。\par

% structure: p0123 source=h2 target=subsection reason=relative-heading-rank; base=h2

\subsection{第十八章 他阐明了宗徒的话：\BibleQuote{我要摧毁智慧人的智慧。}}

关于这样的人，经上说：\Quote{我要消灭智者的智慧，废除明达者的明达。}宗徒相应地补充道：\Quote{智者在哪里？经师在哪里？今世的辩士在哪里？}他将经师、今世的辩士和外邦人的哲学家相对照。\Quote{天主岂不是使今世的智慧变得愚妄了吗？}\BibleRef{格前 1:19}这等于说，天主显明它（今世的智慧）是愚妄的，并不像他们所想的那样真实。你若问他们表面智慧的原因，他必说：\Quote{因为他们心地昏昧；}因为\Quote{在天主的智慧里}，即如先知所宣告的，\Quote{世界不认识}，在\Quote{那藉先知发言的}智慧中，\Quote{祂}，即天主——\Quote{天主喜悦藉着宣讲的愚妄}——这在希腊人看来是愚妄——\Quote{拯救那些相信的人。因为犹太人要求神迹}为要相信；\Quote{希腊人寻求智慧，}显然是指那些被称为\Quote{不可辩驳}的推理，和那些被称为三段论的推理。\Quote{但我们宣讲被钉在十字架上的耶稣基督；对犹太人来说是绊脚石，}因为他们虽然知道预言，却不相信这事件的发生：\Quote{对希腊人是愚妄。}因为在那些自以为有智慧的人看来，天主子藉着人说话，天主有了一位儿子，尤其是那位儿子竟然受苦，这都是神话。因此他们先入为主的观念使他们倾向于不信。因为救主的来临并没有使人变得愚妄、心硬和不信，反而使他们变得明达、易于说服和相信。但那不愿相信的人，因他们与那些听从者的自愿依附分离，就被证明是没有明悟、不信和愚笨的。\Quote{但对那些蒙召的人，无论是犹太人或希腊人，基督是天主的德能和天主的智慧。}难道我们不应（更好地）理解所翻译的那句话\Quote{天主岂不是使今世的智慧变得愚妄了吗？}为否定意义：\Quote{天主并没有使今世的智慧变得愚妄吗？}——这样他们心硬的原因就不会显得来自天主\Quote{使今世的智慧变得愚妄}了。因为在各方面，他们本是明智的，却不相信这宣讲，因而招致更大的责罚。因为对真理的偏爱和选择是自愿的。但那宣告\Quote{我要消灭智者的智慧}则宣告祂发出了光明，祂同时提出了被鄙视和轻蔑的蛮族哲学作为对照；正如灯在太阳照耀下，据说就熄灭了，因为这时它不再发出同样的光芒。既然所有人都蒙了召，那些愿意听从的人就被称为\Quote{蒙召的人}。因为天主并不偏不倚。两个民族中相信的人，都成了\Quote{特选的子民。}\BibleRef{铎 2:14}在宗徒大事录中，你将看到这原话：\Quote{于是，凡接受他的话的人，都领了洗；}\BibleRef{宗 2:41}但那不愿听从的人则自行疏远。先知对这些说：\Quote{如果你们愿意，听从我，你们将享用地上的美物；}\BibleRef{依 1:19}证明选择或拒绝在于我们自己。宗徒将符合主的教义称为\Quote{天主的智慧}，以表明真正的哲学是由圣子传达的。此外，那有智慧表象的人，被宗徒嘱咐了某些劝勉：\Quote{你们应穿上新我，这新我是按照天主在正义和真正的圣洁中更新的。所以你们要弃绝谎言，各人与近人说实话。不要给魔鬼留余地。那偷窃的，不要再偷，反倒要劳苦，亲手做正经事}（工作就是在寻求真理中劳苦，因为它伴随着理性的善行），\Quote{为能使你有所给与那有需要的人，}无论是世俗的财富，还是属神的智慧。因为他希望道理被教导，金钱经过准确检验后存入银行，以积累利息。因此他接着说：\Quote{不要让腐败的言语从你们口中出来，}——即\Quote{腐败的言语}，出自骄傲——\Quote{只要按需要说造就人的好话，为能给听众带来恩宠。}善良的天主的话必须也是善良的。那拯救人的，怎么可能不是善良的呢？\par

% structure: p0125 source=h2 target=subsection reason=relative-heading-rank; base=h2

\subsection{第十九章 哲学家们也获得了部分真理}

既然\Person{保禄}被证明确实提出了一些真确的观点，我们便可由此一瞥这些见证。\Person{保禄}在\WorkTitle{宗徒大事录}中，记载他曾向亚略巴古人说：\BibleQuote{我看你们在各方面都更敬畏神明。因为我经过各处，细看你们所敬拜的，也见到一座祭坛，上面写着\InnerQuote{向未识之神}。现在，我就将你们所敬拜而不认识的这位，传告给你们。创造宇宙及其中万物的天主，既是天地的主，就不住人手所建的殿宇，也不受人手的侍奉，好像需要什么似的，而是他将生命、呼吸和一切赏给了众人。他由一个人造了全人类，使他们住在全地面上，给他们立定了年限，和他们所居处的疆界；好使他们寻求天主，或者可以摸索而找到他；其实他离我们每人并不远。因为我们生活、行动、存在，都在他内，正如你们自己的诗人也说过：\InnerQuote{我们也是他的子孙。}} \BibleRef{宗 17:22}由此显而易见，宗徒引用\Person{阿拉托斯}的\WorkTitle{《现象》}中的诗例，赞许了希腊人说得好的部分；并暗示，希腊人通过那位未识之神，以一种迂回的方式敬拜了造物主天主；但有必要通过正面的知识，藉着圣子来领会和认识祂。\BibleQuote{所以，我打发你到外邦人那里去，}经上说，\BibleQuote{开启他们的眼，叫他们从黑暗中转向光明，从撒殚的权下皈依天主；好使他们因信我而获得罪赦，并在因信我而成圣的人中得份。} \BibleRef{宗 26:17}如此，瞎子的眼开了。藉着圣子认识圣父，就是领悟那\Quote{希腊式迂回表达}；而从撒殚的权下转向，就是脱离罪，因为罪带来束缚。我们确实不全盘接受所有哲学，只接受\Person{苏格拉底}在\Person{柏拉图}著作中所说的那种。\Quote{有人说，在神秘仪式中，有许多手持酒神杖的人，但很少是酒神狂女；}意思是，\BibleQuote{许多人被召，但很少被选。}他于是清楚加上：\Quote{在我看来，这些人无非是那些正确地进行哲学思考的人；为了成为他们中的一员，我一生已尽我所能，无不努力。我们是否努力正确并有所成就，当我们到了那里之后，如果天主愿意，不久便会知道。}他岂不是似乎从希伯来文圣经中宣告了义人在死后因信德而有的希望吗？在\WorkTitle{《德摩多库斯》}（如果那确实是\Person{柏拉图}的作品）中：\Quote{不要以为我称那在技艺中消磨一生或学习许多东西为哲学思考，而是有所不同；至少我会认为那是一种耻辱。}我想，他知道\Person{赫拉克利特}所说的\Quote{博学并不教育心智。}在\WorkTitle{《理想国》}第五卷中，他说：\Quote{\InnerQuote{那么，我们要称所有这些以及研究这类事物的人，还有那些致力于较低技艺的人，为哲学家吗？}\InnerQuote{绝不，}我说，\InnerQuote{但像哲学家。}\InnerQuote{那么你称谁，}他说，\InnerQuote{为真正的呢？}\InnerQuote{那些，}我说，\InnerQuote{喜爱默观真理的人。因为哲学不在于几何学及其公设和假设；也不在于音乐，那是臆测的；更不在于天文学，充满物理的、流动的和可能的原因。但善本身和真理本身的知识是必需的——善是一回事，通往善的路是另一回事。}}因此，他不认为训练课程足以达到善，而是协助激发和训练灵魂朝向理性对象。那么，如果他们说希腊人偶然说出了某些真正哲学的言语，那是神圣管理下的偶然（因为没有人会为了与我们争论而神化偶然）；或者说是好运，好运并非不可预见。另一方面，如果有人要说希腊人拥有对这些事物的自然概念，我们知道自然界的唯一造物主；正如我们也称义德是自然的；或者他们拥有共同的理智，让我们思考谁是它的父亲，以及在理智的运作中什么是义德。因为如果有人称其为\Quote{预言}，并归因于\Quote{联合表达}，他是在指定预言的类型。此外，其他人会认为某些真理是由哲学家在外表上说出的。\par

神圣的宗徒论到我们说：\Quote{因为我们是借着镜子观看。}\BibleRef{格前 12:12} 我们在其中借着反射认识自己，同时也尽我们所能，从我们内那神圣的根源，默观那动力因。因为经上说：\Quote{你看见了你的兄弟，你就看见了你的天主；}我想，现今那拯救人的天主已向我们宣告。但在脱去肉身后，\Quote{面对面}——那时，当心灵变得纯洁，我们才能确定而全面地看见。借着反思和直观，那些在希腊人中精通哲学的人看见了天主。因为由于我们的软弱，我们真实的见解就像水中所见的影像，又如我们透过清澈透明的物体看见事物。因此，\Person{所罗门}说得极好：\Quote{播种正义的，工作信德。}\BibleRef{箴 11:21} \Quote{有些人散播自己的，反而增多了。}\BibleRef{箴 11:24} 又说：\Quote{你要留意平原上的青草，你必割草并收集成熟的干草，好使你有羊群作衣服。}\BibleRef{箴 27:25} 你看，对外衣和保养必须多么留心。\Quote{你必明智地认识你羊群的灵魂。}\BibleRef{箴 27:23} \Quote{因为当外邦人没有法律，却顺着本性遵行法律上的事，他们虽然没有法律，对自己就是法律；未受割礼的遵守法律的规条，}按照宗徒，这是法律之前和来临之前的事。\BibleRef{罗 2:14} 圣言彷彿对沉溺于哲学的人和被称为异端者作比较，极清楚地说道：\Quote{近处的朋友胜过远方的兄弟。}\BibleRef{箴 27:10} \Quote{信赖谎言的，靠风为生，追逐飞鸟。}\BibleRef{箴 9:12} 我不认为哲学直接宣告圣言，尽管在许多情况下哲学试图并说服性地教导我们可能的论证；但它攻击那些教派。因此又加上：\Quote{因为他离弃了自己葡萄园的道路，在他自己耕种的路径上徘徊。}这些就是离弃了原始教会的教派。如今，陷入异端的人穿越干旱的旷野，弃绝唯一真天主，缺乏天主，寻求无水之水，到达无人居住和干渴之地，亲手收集荒芜。而那些缺乏谨慎的人，即陷入异端的人，智慧说：\Quote{我吩咐你们，说：偷来的水是甜的，偷来的饼是美味的。}\BibleRef{箴 9:17} 圣经显然将饼和水这些术语应用于不是别的，正是那些在祭献中按照教会规条使用饼和水的异端。因为有些人只用水举行圣体圣事。\Quote{但你离去吧，不要停留在她的地方；}此处\Emph{地方}指会堂，而非教会。祂用含义模糊的名词\OriginalTerm{地方}{place}来称呼它。然后祂接着说道：\Quote{因为如此，你必经过别人的水；}视异端的圣洗为非正当的、不真实的水。\Quote{你必经过别人的河，}那河流奔腾而下入海；凡偏离真理的稳定性的人，便被投入其中，堕入外邦人混乱的生命波涛中。\par

% structure: p0128 source=h2 target=subsection reason=relative-heading-rank; base=h2

\subsection{哲学有助于理解神圣真理}

如同许多人拖曳船只，不能称为许多原因，而是一个由众多组成的原因——因为每一人单独并非船被牵引的原因，而是与其余人一起——同样，哲学作为对真理的探求，有助于对真理的理解；不是作为理解的原因，而是与其他事物一起的原因和合作者；或许也是共同原因。又如各种美德是个人幸福的原因；又如太阳、火、洗澡和衣服是使人暖和的原因：同样，虽然真理是唯一的，许多事物有助于对它的探究。但它的发现在于圣子。如果我们仔细思量，德行在潜能上是一个。但事实上，当显现在某些事物中时，它被称为明智，在其他事物中称为节制，在其他事物中称为勇敢或义德。同样的类比，虽然真理是唯一的，在几何学中有几何学的真理；在音乐中有音乐的真理；在正确的哲学中，将有希腊真理。但那唯一的、无可辩驳的真实真理，是由天主圣子教导我们的。同样，我们说，同一块银币，交给船主时称为船费；交给税吏时称为税；交给地主时称为租金；交给教师时称为学费；交给卖主时称为定金。而每一个（无论是美德还是真理），虽名称相同，却只是自身特定效果的原因；它们的混合产生了幸福的生活。因为我们不是单凭名称而成为幸福的，当我们说善生就是幸福，且灵魂以美德装饰的人是幸福的时候。但如果哲学通过对真理的种种探讨，寻求那贴近真理的知识（即我们所拥有的知识），间接有助于发现真理，那么它帮助那按圣言致力于把握真理的人领悟知识。但是希腊真理与我们持有的真理不同（尽管名称相同），无论是在知识的广度、论证的确实性、神圣德能等方面。因为我们是由天主所教导的，由天主圣子以真正的\Quote{神圣的经文}训诲我们。因此我们所提及的那些人，并非以我们的方式影响灵魂，而是藉不同的教训。并且，如果为了那些喜欢挑剔的人，我们必须加以区别，说哲学是真实领会的协同和合作原因（它是真理的探求），那么我们将承认它是为\OriginalTerm{有真知者}{\GreekText{τοῦ} \GreekText{γνωστικοῦ}}的预备训练；不把仅仅是共同原因的东西当作原因；也不把仅仅是合作的东西当作维持原因；也不把哲学置于\OriginalTerm{必要条件}{sine quâ non}的地位。因为几乎所有我们这些人，没有经过技艺、科学和希腊哲学的训练，有些甚至完全没有学问，却因一种神圣而野蛮的哲学的影响和德能，透过信德接受了关于天主的圣言，由自动运行的智慧所训练。但是那与其他事物共同作用，本身不能单独作用的东西，我们称之为合作的、共同原因的，并说它之所以成为原因，只因为它是共同原因，并且只由于它与其他事物共同作用而获得原因的名称，但它本身不能产生应有的效果。\par

虽然哲学曾一度使希腊人成义，但并不引领他们达到那完全的义德（经证实它与这义德合作），正如第一和第二级台阶帮助你上楼，文法家帮助哲学家一样。并非仿佛由于缺少它，完美的圣言就会不完整，或真理就会消亡；因为视觉、听觉和声音也有助于真理，但心智才是认识真理的适当官能。但在那些合作的事物中，有些贡献较大的力量；有些较小。因此清晰有助于传递真理，逻辑有助于防止我们陷入攻击我们的异端。但是按照救主而来的教导，本身是圆满的、没有缺陷的，是\Quote{天主的德能和智慧}\BibleRef{格前 1:24}；希腊哲学并不因其介入而使真理更有力；而是使诡辩对真理的攻击失去效力，挫败针对真理的阴险阴谋，才被称为葡萄园正当的\Quote{篱笆和墙垣}。而按照信德的真理，对生命是必需的，如同面包；而预备性的训练如同调味汁和甜食。底比斯的品达说：\Quote{筵席结束时，甜点令人愉快。}圣经也明确说：\Quote{无知者因明达而成为智者，智者将获得知识。}\BibleRef{箴 21:11}\Quote{那凭自己说话的，}主说，\Quote{寻求自己的光荣；但那寻求派遣他来者的光荣的，是真实的，在他内没有不义。}\BibleRef{若 7:18}因此，反之，那占取属于野蛮人的东西，并夸耀为已有的人，是作恶，增加自己的光荣，歪曲真理。圣经称这样的人为\Quote{贼}。因此经上说：\Quote{儿子，不要作说谎者；因为谎言导致偷窃。}然而贼确实拥有他不正当地据为己有的东西，无论是金子、银子、言论还是教义。因此他们偷窃的、部分真实的观念，是通过猜测和必要的逻辑推论知道的：因此当他们成为门徒时，将以理智的领悟认识它们。\par

% structure: p0131 source=h2 target=subsection reason=relative-heading-rank; base=h2

\subsection{第21章 犹太人的制度和法律远比希腊哲学古老}

关于哲学家们从希伯来人那里剽窃教义的问题，我们稍后再加讨论。但首先，按照应有的顺序，我们现在必须谈谈\Person{梅瑟}的时代，由此将无可辩驳地证明，希伯来人的哲学是一切智慧中最古老的。\Person{塔提安}在其著作\WorkTitle{致希腊人}中，以及\Person{卡西安}在其\WorkTitle{释经学}第一卷中，都已精确地讨论过这一问题。然而，我们的注释也要求我们扼要重述已说过的话。语法学家\Person{阿皮翁}，绰号\Quote{常胜者}，在其\WorkTitle{埃及史}第四卷中，尽管身为埃及人，对希伯来人怀有如此敌意，以至于写了一部反对犹太人的作品，但在提及埃及国王\Person{阿摩西斯}及其事迹时，他引用了\Person{门德斯的托勒密}作为见证。他的说法大致如下：\Person{阿摩西斯}生活在\Place{阿尔戈斯}的\Person{伊纳库斯}时代，推翻了\Place{阿提里亚}，\Person{门德斯的托勒密}在其\WorkTitle{编年史}中如此记载。这位\Person{托勒密}是一位司祭；他在整整三卷书中叙述了埃及国王的事迹，他说犹太人在\Person{梅瑟}率领下离开埃及，发生在\Person{阿摩西斯}为埃及国王之时。由此可见，\Person{梅瑟}活跃于\Person{伊纳库斯}时代。在希腊诸城邦中，最古老的是\Place{阿尔戈斯}城邦，即起源于\Person{伊纳库斯}的那一个，正如\Place{哈利卡纳苏斯}的\Person{狄奥尼修斯}在其\WorkTitle{年代记}中所教导的。而由\Person{刻克洛普斯}建立的\Place{阿提刻}，比它晚四十代——\Person{刻克洛普斯}是一位具有双重血统的原住民，\Person{塔提安}明确这样说；由\Person{佩拉斯戈斯}建立的\Place{阿卡迪亚}，则晚九代；据说他也是一位原住民。而比这再晚五十二代的，是由\Person{丢卡利翁}建立的\Place{弗提奥提斯}。从\Person{伊纳库斯}时代到特洛伊战争，共计算了二十代或更多；让我们说，四百多年。如果\Person{克特西亚斯}说亚述帝国比希腊早许多年，那么\Person{梅瑟}离开埃及将看来发生在亚述帝国第四十二年，\Person{贝洛库斯}统治第三十二年，正值埃及的\Person{阿摩西斯}和\Place{阿尔戈斯}的\Person{伊纳库斯}时代。在希腊，在\Person{伊纳库斯}的继位者\Person{福洛纽斯}时代，发生了\Person{奥吉格斯}洪水；\Place{西库翁}最先由\Person{埃吉阿琉斯}统治，然后是\Person{欧罗普斯}，再是\Person{忒尔刻斯}；\Place{克里特}则由\Person{克瑞斯}统治。因为\Person{阿库西劳斯}说\Person{福洛纽斯}是第一个人。因此，\WorkTitle{福洛尼德}的作者也说他是\Quote{凡人之父}。随后，\Person{柏拉图}在\WorkTitle{蒂迈欧}中追随\Person{阿库西劳斯}写道：\Quote{他想引导他们讨论古代事物，便说他敢于谈论这座城最古老的往事，关于被称为第一个人的\Person{福洛纽斯}，以及\Person{尼俄柏}，还有洪水之后发生的事。}在\Person{福耳巴斯}时代，有\Person{阿克泰俄斯}，\Place{阿刻泰亚}即\Place{阿提刻}之名源于他；在\Person{特里俄帕斯}时代，有\Person{普罗米修斯}、\Person{阿特拉斯}、\Person{埃庇米修斯}、双重血统的\Person{刻克洛普斯}和\Person{伊诺}。在\Person{克罗托普斯}时代，发生了\Person{法厄同}的焚烧和\Person{丢卡利翁}的洪水；在\Person{斯忒涅洛斯}时代，是\Person{安非克堤翁}的统治和\Person{达那俄斯}到达\Place{伯罗奔尼撒}；在\Person{达耳达诺斯}时代，建造了\Place{达耳达尼亚}，\Person{荷马}他——\par

\Quote{集云者\Person{宙斯}最初所生，}—\par

以及从\Place{克里特}迁徙到\Place{腓尼基}。在\Person{林叩斯}时代，发生了\Person{珀耳塞福涅}被劫、\Place{厄琉西斯}圣域的奉献、\Person{特里普托勒摩斯}的农耕、\Person{卡德摩斯}到达\Place{忒拜}，以及\Person{米诺斯}的统治。在\Person{普洛托斯}时代，发生了\Person{欧摩尔波斯}与雅典人的战争；在\Person{阿克里西俄斯}时代，发生了\Person{珀罗普斯}从\Place{弗律基亚}迁来、\Person{伊翁}到达\Place{雅典}；出现了第二个\Person{刻克洛普斯}，\Person{珀耳修斯}和\Person{狄俄尼索斯}的事迹发生，\Person{俄耳甫斯}和\Person{穆赛俄斯}在世。在\Person{阿伽门农}统治第十八年，\Place{特洛伊}被攻陷，正值\Place{雅典}的\Person{忒修斯}之子\Person{得摩丰}统治第一年，塔尔格利翁月第十二日，据\Place{阿尔戈斯}的\Person{狄奥尼修斯}所说；但\Person{埃吉阿斯}和\Person{得耳库洛斯}在第三卷中说，是帕涅摩斯月最后一段的第八日；\Person{赫拉尼科斯}说是塔尔格利翁月第十二日；而一些\WorkTitle{阿提刻}的作者说，在\Person{墨涅斯透斯}最后一年的最后一段的第八日，正值满月。\par

\Quote{那是午夜，}\par

\WorkTitle{小伊利亚特}的作者说，\par

\Quote{月光皎洁。}\par

其他人说，它发生在同一日期（斯基洛福里翁月）。但\Person{赫拉克勒斯}的对手\Person{忒修斯}比特洛伊战争早一代。因此，\Person{赫拉克勒斯}之子\Person{特勒波勒摩斯}被\Person{荷马}提到曾在\Place{特洛伊}服役。\par

因此，\Person{梅瑟}被证明比\Person{狄俄尼索斯}的神化早六百零四年，如果\Person{狄俄尼索斯}在\Person{珀耳修斯}统治第三十二年神化，正如\Person{阿波罗多洛斯}在其\WorkTitle{编年史}中所说。从\Person{巴克科斯}到\Person{赫拉克勒斯}以及乘阿尔戈船随\Person{伊阿宋}航行的首领们，共包含六十三年。\Person{阿斯克勒庇俄斯}和\Person{狄奥斯库里}与他们一同航行，\Place{罗德岛}的\Person{阿波罗尼奥斯}在其\WorkTitle{阿尔戈英雄纪}中为此作证。从\Person{赫拉克勒斯}在\Place{阿尔戈斯}的统治到\Person{赫拉克勒斯}和\Person{阿斯克勒庇俄斯}的神化，共包含三十八年，按编年史家\Person{阿波罗多洛斯}；从这到\Person{卡斯托尔}和\Person{波鲁克斯}的神化，五十三年。此时\Place{特洛伊}被攻陷。如果我们相信诗人\Person{赫西俄德}，让我们听他：——\par

\Quote{然后\Person{迈亚}，\Person{阿特拉斯}的女儿，为\Person{朱庇特}生了著名的\Person{赫耳墨斯}，\newline{}不死者的信使，当她登上神圣的床榻。\newline{}\Person{卡德摩斯}的女儿\Person{塞墨勒}也生了一个优秀的儿子，\newline{}\Person{狄俄尼索斯}，令人喜悦的，当她与他结合在爱中。}\par

塞墨勒的父亲卡德摩斯，在林叩斯时代来到特拜，他是希腊字母的发明者。特里俄帕斯是伊西斯的同时代人，为伊那科斯之后的第七代。据说伊西斯与伊俄相同，得名于她漫游整个大地的\OriginalTerm{行走}{\GreekText{ἰέναι}}。伊斯忒洛斯在其著作《埃及人的迁徙》中称她为普罗米修斯的女儿。普罗米修斯生活在特里俄帕斯时代，即梅瑟之后的第七代。因此，按照希腊人的年代学，梅瑟似乎在人类出生之前就已兴盛。研究埃及神祇的莱昂说，伊西斯被希腊人称为刻瑞斯，她生活在林叩斯时代，即梅瑟之后的第十一代。据阿里斯提珀斯在\WorkTitle{Arcadica}第一卷所述，阿尔戈斯国王阿庇斯建造了孟斐斯。阿尔戈斯的阿里斯忒阿斯说他被称为塞拉庇斯，并且埃及人所崇拜的就是他。安菲波利的宁佛多洛斯在\WorkTitle{亚洲的制度}第三卷中说，阿庇斯公牛死后被放入\OriginalTerm{棺材}{\GreekText{σορός}}，安放在该处所崇拜的\OriginalTerm{神明}{\GreekText{δαίμονος}}的神殿中，因而被称为索罗阿庇斯，后来按照本地习俗称为塞拉庇斯。阿庇斯是伊那科斯之后的第三代。此外，勒托生活在提堤俄斯时代。\Quote{因为他拖住了拉托娜，宙斯光彩的配偶。}现在提堤俄斯与坦塔罗斯同时代。因此，波奥提亚的品达正确地写道：\Quote{后来阿波罗诞生了；}这并不奇怪，因为他被发现与赫拉克勒斯一同服侍阿德墨托斯\Quote{长达一年。}音乐的发明者泽托斯和安菲翁生活在卡德摩斯时代左右。如果有人声称斐摩诺厄是第一个向阿克里西俄斯吟唱预言诗的人，那么他应当知道，斐摩诺厄之后二十七年，俄耳甫斯、穆赛俄斯以及赫拉克勒斯的老师利诺斯才出现。荷马和赫西俄德比特洛伊战争晚得多；在他们之后，希腊的立法者更晚，如来库古、梭伦、七贤、叙罗斯的斐瑞居德斯以及伟大的毕达哥拉斯，他们后来大约在奥林匹亚时代，正如我们所展示的。我们也证明了梅瑟不仅比希腊人中所谓诗人和智者更古老，甚至比他们的大多数神祇更古老。不仅他，而且西比尔也比俄耳甫斯更古老。因为据说，关于她的名称和神谕有好几种说法：她是弗里吉亚人，被称为阿尔忒弥斯；当她到达得尔福时，她吟唱道——\par

\Quote{哦，得尔福的人们，远射的阿波罗的侍从，\newline{}我前来宣告持神盾的宙斯的心意，\newline{}尽管我对我自己的兄弟阿波罗感到愤怒。}\par

还有另一位，埃律特赖亚的，名叫赫洛菲勒。这些由本都的赫拉克利德斯在其著作\WorkTitle{论神谕}中提及。我略过埃及的西比尔和意大利的西比尔，后者曾居住在罗马的卡尔门塔来，她的儿子是厄万德耳，他在罗马建造了称为卢佩卡尔的潘神神庙。\par

至此，我们有必要考察梅瑟之后希伯来人其他先知的年代。梅瑟的生活结束后，若苏厄接替领导民众；他作战六十五年后，又在福地安息了二十五年。正如《若苏厄书》所述，上述之人是梅瑟的继承人，共二十七年。然后希伯来人犯罪，被交给美索不达米亚的库沙察尔八年，如《民长纪》所提及。但后来他们恳求上主，便得到首领哥特尼耳，他是加肋布之弟，属犹大支派；他杀了美索不达米亚王后统治民众四十年。他们又犯罪，被交到摩阿布王厄革隆手中十八年。但他们悔改后，厄胡得——一位双手均能使用的厄弗辣因支派人——作他们的领袖八十年。正是他杀了厄革隆。厄胡得死后，他们再次犯罪，被交到客纳罕王雅宾手中二十年。之后，拉丕多特的妻子德波辣（属厄弗辣因支派）说预言，同时希厄稣的儿子奥齐雅是大司祭。在她的激励下，纳斐塔里支派贝内尔的儿子巴辣克率军与雅宾的元帅息色辣交战并战胜了他。此后，德波辣治理并审判民众四十年。她死后，民众再次犯罪，被交到米德杨人手中七年。这些事后，约阿士的儿子基德红（属默纳协支派）率他的三百人作战，杀了十二万人，统治四十年；之后阿彼默肋客的儿子（统治）三年。继位的是贝丹的儿子、卡兰的孙子波肋阿斯（属厄弗辣因支派），他统治二十三年。此后，民众再次犯罪，被交到阿孟子民手中十八年；悔改后，由基肋阿得人依弗大（属默纳协支派）统领，他统治六年。之后，犹大支派白冷人阿巴堂统治七年。然后则步隆人厄布龙统治八年。再后，厄弗辣因人厄革隆统治八年。有人将阿巴堂的七年加上厄布龙的八年。此后，民众再次犯罪，落入外邦人培肋舍特人的权下四十年。但他们归向上主后，由丹支派的三松带领，他在战斗中战胜了外邦人。他统治二十年。之后，没有首领时，司祭厄里审判民众四十年。继位的是先知撒慕尔；与撒慕尔同时，撒乌耳作王，统治二十七年。他傅油立了达味。撒慕尔在撒乌耳前两年去世，当时阿彼默肋客是大司祭。他傅油立撒乌耳为王，后者是继民长之后第一位统治以色列的王。从民长到撒乌耳的全部时间，共四百六十三年零七个月。\par

然后在《\BibleBook{列王纪}{上}》中，\Person{撒乌耳}作王二十年，这是他更新之后执政的时期。\Person{撒乌耳}死后，犹大支派\Person{叶瑟}的儿子\Person{达味}，继而在\Place{赫贝龙}作王四十年，正如《\BibleBook{列王纪}{下}》所记载的。\Person{阿彼雅塔}，\Person{阿希默肋客}的儿子，\Person{厄里}家族的司祭，担任大司祭。在他那时，\Person{加得}和\Person{纳堂}说预言。从\Person{农}的儿子\Person{若苏厄}那时起，直到\Person{达味}承受王国，按某些人的说法，中间隔了四百五十年。但依照所陈列的编年，到\Person{达味}去世为止，共计五百二十三年零七个月。\par

此后，\Person{达味}的儿子\Person{撒罗满}作王四十年。在他之下，\Person{纳堂}继续说预言，他也曾就建造圣殿劝勉\Person{撒罗满}。\Place{史罗}的\Person{阿希雅}也说预言。\Person{达味}和\Person{撒罗满}两位君王都是先知。大司祭\Person{匝多克}是第一位在\Person{撒罗满}所建的圣殿中供职的，他是自第一位大司祭\Person{亚郎}以来的第八位。从\Person{梅瑟}到\Person{撒罗满}的时代，按某些人的说法，是五百九十五年，而按另一些人的说法，是五百七十六年。\par

如果你将\Person{若苏厄}到\Person{达味}的四百年五十年，加上\Person{梅瑟}治理的四十年，以及\Person{梅瑟}在希伯来人出离埃及之前另外的八十年，你总共会得到六百一十年。但若是在到\Person{达味}去世的五百二十三年零七个月上，加上\Person{梅瑟}的一百二十年和\Person{撒罗满}的四十年，我们的编年会更为准确。因为你总共会得到，直到\Person{撒罗满}去世时，六百八十三年零七个月。\par

\Person{希兰}将他的女儿给了\Person{撒罗满}，大约是在\Person{墨涅拉俄斯}到达\Place{腓尼基}的时候，即\Place{特洛伊}被攻陷之后，据\Person{佩尔伽蒙的墨南德}和\WorkTitle{腓尼基}的作者\Person{莱图斯}所说。\Person{撒罗满}之后，他的儿子\Person{勒哈贝罕}作王十七年；\Person{匝多克}的儿子\Person{阿彼默肋客}是大司祭。在他的统治期间，王国分裂，\Place{厄弗辣因}支派的\Person{雅洛贝罕}，\Person{撒罗满}的臣仆，在\Place{撒玛黎雅}作王；\Place{史罗}人\Person{阿希雅}继续说预言；还有\Person{阿玛默}的儿子\Person{撒玛亚}，以及从犹大来到\Person{雅洛贝罕}那里、预言反对祭坛的那位。此后，他的儿子\Person{阿彼雅姆}在位二十三年；同样，他的儿子\Person{阿撒曼}。后者在年老时脚上有病；在他的统治期间，\Person{哈纳尼}的儿子\Person{耶胡}说预言。\par

此后，他的儿子\Person{约沙法特}作王二十五年。在他的统治期间，\Place{特士贝}人\Person{厄里亚}说预言，还有\Person{耶布拉}的儿子\Person{米加雅}，以及\Person{哈纳尼}的儿子\Person{敖巴狄雅}。在\Person{米加雅}的时代，还有假先知，\Person{科纳罕}的儿子\Person{漆德克雅}。随后，\Person{约沙法特}的儿子\Person{约兰}作王八年；其间，\Person{厄里亚}说预言；\Person{厄里亚}之后，是\Person{沙法特}的儿子\Person{厄里叟}。在他的统治期间，\Place{撒玛黎雅}的人民吃鸽粪和自己的儿女。\Person{约沙法特}的时期从《\BibleBook{列王纪}{上}》的结束延续到《\BibleBook{列王纪}{下}》。在\Person{约兰}统治期间，\Person{厄里亚}被提升天，\Person{沙法特}的儿子\Person{厄里叟}开始说预言，他预言了六年，那时他四十岁。\par

然后\Person{敖哈齐雅}作王一年。在他那时，\Person{厄里叟}继续说预言，还有\Person{阿达多纳乌斯}与他一起。此后，\Person{敖哈齐雅}的母亲\Person{哥托里雅}作王八年，她杀死了她兄弟的儿女。因为她属于\Person{阿哈布}家族。但\Person{敖哈齐雅}的姐姐\Person{约沙贝雅}偷走了\Person{敖哈齐雅}的儿子\Person{约阿士}，后来使他登上王位。在\Person{哥托里雅}那个时期，\Person{厄里叟}仍在说预言。之后，如我先前所说，由大司祭\Person{约达尔}的妻子\Person{约沙贝雅}救出的\Person{约阿士}作王，他一共活了四十年。\par

因此，从\Person{撒罗满}到\Person{厄里叟}先知去世，按某些人的说法，共有一百零五年；按另一些人的说法，是一百零二年；而根据我们面前的编年，从\Person{撒罗满}的统治算起，是一百八十一年。\par

现在，从\Place{特洛伊}战争到\Person{荷马}的出生，据\Person{菲洛科罗斯}说，过了一百八十年；他生活在爱奥尼亚移民之后。但\Person{阿里斯塔科斯}在\WorkTitle{阿耳希洛科斯回忆录}中说，他生活在爱奥尼亚移民期间，那是在\Place{特洛伊}被围困一百二十年之后。而\Person{阿波罗多罗斯}声称是在爱奥尼亚移民一百二十年之后，当\Person{多律塞乌斯}的儿子\Person{阿革西拉乌斯}是\Place{拉栖代梦}人的王时：所以他使立法者\Person{吕库古}，当时还是年轻人，靠近他。\Person{欧梯美涅斯}在\WorkTitle{编年史}中说，他与\Person{赫西俄德}同时代，在\Person{阿卡斯托斯}时代，生于\Place{希俄斯}，大约在\Place{特洛伊}被攻陷后四百年。\Person{阿尔喀玛科斯}在其\WorkTitle{优卑亚史}第三卷中持有此观点。因此，他和\Person{赫西俄德}都晚于\Person{厄里叟}先知。如果你追随文法学家\Person{克拉特斯}，说\Person{荷马}生于\Place{赫拉克利德}远征时期，大约在\Place{特洛伊}被攻陷后八十年，那么会发现他又晚于\Person{撒罗满}，因为如前面所说，在\Person{撒罗满}时代发生了\Person{墨涅拉俄斯}到达\Place{腓尼基}。\Person{埃拉托斯特尼}说\Person{荷马}的时代在\Place{特洛伊}被攻陷后二百年。此外，\Person{特奥蓬波斯}在\WorkTitle{腓力史}第四十三卷中记载，\Person{荷马}生于\Place{特洛伊}战争后五百年。\Person{欧福里翁}在他的\WorkTitle{阿留阿代斯}一书中主张，他生在\Person{古格斯}时代，\Person{古格斯}在第十八届奥林匹亚开始为王，他也说\Person{古格斯}是第一个被称为\OriginalTerm{僭主}{\GreekText{τύραννος}}的人。再者，\Place{拉科尼亚}的\Person{索西比乌斯}在其\WorkTitle{日期记录}中将\Person{荷马}下推到\Person{波吕得克图斯}的儿子\Person{卡里鲁斯}统治的第八年。\Person{卡里鲁斯}在位六十四年，此后\Person{尼坎德尔}的儿子统治三十九年。据说在他统治的第三十四年，第一届奥林匹亚竞技会创立；这样，\Person{荷马}比奥林匹克运动会的引入早九十年。\par

\Person{约阿士}之后，他的儿子\Person{阿玛齐雅}继位，统治了三十九年。同样，他的儿子\Person{敖齐雅}继位，统治了五十二年，死于癞病。在他那时，\Person{亚毛斯}说预言，还有他的儿子\Person{依撒意亚}，以及\Person{贝黎}的儿子\Person{欧瑟亚}，\Person{阿玛泰}的儿子\Person{约纳}，他来自\Place{加特赫费尔}，向\Place{尼尼微}人讲道，并经过鲸鱼之腹。\par

然后，\Person{乌齐雅}的儿子\Person{约纳堂}登基，在位十六年。在他时期，\Person{依撒意亚}仍在说预言，还有\Person{欧瑟亚}、摩辣士提人\Person{米该亚}，以及\Person{贝突耳}的儿子\Person{岳厄尔}。\par

继承他的是他的儿子\Person{阿哈兹}，在位十六年。在他时期，第十五年，\Place{以色列}被掳到\Place{巴比伦}。亚述王\Person{撒玛纳色}将\Place{撒玛黎雅}的人民掳到\Place{玛待}地区和\Place{巴比伦}。\par

阿哈兹又被\Person{曷舍亚}继位，他统治了八年。随后是\Person{希则克雅}，共二十九年。因他的圣德，当他临近结局时，天主藉\Person{依撒意亚}允许他再活十五年，并赐下太阳后退的征兆。直到他时期，\Person{依撒意亚}、\Person{欧瑟亚}和\Person{米该亚}继续预言。\par

据说这些人生于\Person{莱喀古斯}时代之后，\Person{莱喀古斯}是\Place{拉刻代蒙人}的立法者。因为\Person{狄奥奇达斯}在\WorkTitle{墨伽里克斯}的第四卷中，将\Person{莱喀古斯}时代定为在\Place{特洛伊}陷落后大约二百九十年。\par

希则克雅之后，他的儿子\Person{默纳舍}在位五十五年。然后他的儿子\Person{阿孟}在位两年。他之后，他的儿子\Person{约史雅}在位三十一年，以谨守法律著称。他\Quote{将人的尸首放在偶像的尸首上}，正如\WorkTitle{肋未纪}所载。\BibleRef{肋 26:30} 他在位第十八年举行了逾越节，自\Person{撒慕尔}的日子以来中间时期未曾守过。\BibleRef{列下 23:22} 然后司祭\Person{希耳克雅}，先知\Person{耶肋米亚}的父亲，遇见了藏在圣殿中的法律书，读了以后就去世了。\BibleRef{列下 22:8} 在他时期，\Person{胡耳达}说预言，还有\Person{索福尼雅}和\Person{耶肋米亚}。在\Person{耶肋米亚}时期，有\Person{阿组尔}的儿子\Person{阿纳尼雅}，是假先知。他不听从\Person{耶肋米亚}先知，在\Place{幼发拉底河}被埃及王\Person{法郎乃苛}杀死，因他遭遇了后者，后者正进军\Place{亚述人}。\par

\Person{约史雅}由他的儿子\Person{耶苛尼雅}继位，也称为\Person{约阿哈次}，他统治了三个月零十天。埃及王\Person{法郎乃苛}将他捆绑，带到\Place{埃及}，而改立他的兄弟\Person{约雅金}为王代替他，\Person{约雅金}作他的藩属十一年。他之后，同名者\Person{约雅金}统治三个月。然后\Person{漆德克雅}统治十一年；直到他时期，\Person{耶肋米亚}继续说预言。与他一同说预言的还有\Person{步齐}的儿子\Person{厄则克耳}，以及\Person{撒玛伍斯}的儿子\Person{乌黎雅}，还有\Person{哈巴谷}。\Place{希伯来}君王到此结束。\par

那么，从\Person{梅瑟}诞生直到这次被掳，共九百七十二年；但按严格年代精度，则为一千零八十五年六个月零十天。从\Person{达味}登基到被\Place{加色丁人}掳去，共四百五十二年六个月；但按我们核对日期所得的精度，则为四百八十二年六个月零十天。\par

在\Person{漆德克雅}在位第十二年，即\Place{波斯人}统治前四十年，\Person{拿步高}对\Place{腓尼基人}和\Place{犹太人}开战，如\Person{贝洛索斯}在其\WorkTitle{加色丁史}中所断言。\Person{若阿巴斯}在写\Place{亚述人}事时，承认他从\Person{贝洛索斯}得到历史，并为其准确性作证。因此，\Person{拿步高}弄瞎了\Person{漆德克雅}的眼睛，将他带到\Place{巴比伦}，并将全体人民掳走（被掳持续七十年），只有少数逃往\Place{埃及}的人除外。\par

\Person{耶肋米亚}和\Person{哈巴谷}在\Person{漆德克雅}时期仍继续说预言。在他统治第五年，\Person{厄则克耳}在\Place{巴比伦}说预言；他之后，\Person{纳鸿}，然后是\Person{达尼尔}。他之后，\Person{哈盖}和\Person{匝加利亚}又在\Person{达理阿一世}时期预言了两年；然后是小先知中的天使。\Person{哈盖}和\Person{匝加利亚}之后，\Person{乃赫米雅}——\Person{阿塔薛西斯}的司酒长，\Place{以色列}人\Person{阿赫里}的儿子——建造了\Place{耶路撒冷}城，并修复了圣殿。在充军期间，生活了\Person{艾斯德尔}和\Person{摩尔德开}，他们的书至今存留，还有\WorkTitle{玛加伯书}。在这个充军期间，\Person{米沙耳}、\Person{阿纳尼雅}和\Person{阿匝黎雅}，因拒绝朝拜偶像而被投入火窑，因天使显现而得救。那时，因蛇的缘故，\Person{达尼尔}被投入狮子坑；但藉天主眷顾由\Person{哈巴谷}保护，他在第七天被释放。这段时期也发生了\Person{约纳}的征兆；\Person{多俾亚}在天使\Person{辣法耳}的帮助下，娶了\Person{撒辣}为妻，那魔鬼已杀死了她七个未婚夫；\Person{多俾亚}结婚后，他的父亲\Person{托彼特}重获视力。那时，\Person{则鲁巴贝耳}凭智慧战胜对手，获得\Person{达理阿}允许重建\Place{耶路撒冷}，与\Person{厄斯德拉}一同回到故乡；藉着他，人民的救赎、默感神谕的修订和恢复得以完成；并举行了救赎的逾越节，废除了与异族的婚姻。\par

\Person{居鲁士}曾下诏谕令\Place{希伯来人}重建家园。他的应许在\Person{达理阿}时期实现，举行了献殿节，以及帐棚节。\par

总之，从\Person{梅瑟}诞生直到人民复国，包括被掳时期，共计一千一百五十五年六个月零十天；而从\Person{达味}登基算起，据有些人说是四百五十二年；更准确说是五百七十二年六个月零十天。\par

从\Person{耶肋米亚}先知时代发生的\Place{巴比伦}囚掳起，\Person{达尼尔}先知所说的预言便应验了：\BibleQuote{为你的人民和你的圣城已确定了七十个星期，为终止过犯，为结束罪恶，为赎尽不义，为带来永远的正义，为应验神视和预言，为给至圣所傅油。所以你应知道，也应明白：从发出命令重建\Place{耶路撒冷}，直到受傅的君王，是七个星期；再六十二个星期，城中的广场和沟渠要重新建造，且是在困难时期。六十二个星期之后，受傅者将被铲除，不再有判决；他要与来临的领袖一起毁灭城市和圣所。他们必在洪水中被毁灭，直到战争结束，荒凉之事已定。他要在许多人的一周内坚定盟约；在一周的中间，祭献和素祭将被废除；在圣所中要有荒凉的可憎之物，直到终结的时刻，终结将归于荒凉。在一周的中间，他要使献香的祭献停止，并因毁灭的翅膀，直到终结，如同祭献的毁灭。}可见圣殿是在七个星期内建成的，因为\BibleBook{厄斯德拉}{}书上这样记载。这样，基督在七个星期完成后，在\Place{耶路撒冷}为王，成了\Place{犹太}人的君王。在六十二个星期内，整个\Place{犹太地}平静无战事。而基督我们的主，\BibleQuote{至圣者}，既已来临，应验了神视和预言，便因圣父的圣神在祂的肉身内受傅。在先知所说的\BibleQuote{六十二个星期}中，以及在\BibleQuote{一个星期}内，祂是主。在一半的星期中，\Person{尼禄}掌权，并在圣城\Place{耶路撒冷}放置了可憎之物；在一半的星期结束时，他被除去，接着是\Person{奥托}、\Person{加尔巴}和\Person{维特里乌斯}。\Person{维司帕西安}登上了最高权力，毁灭了\Place{耶路撒冷}，使圣所荒凉。事实如此，正如先知所说，有识者自能明白。\par

第十一年结束时，次年年初，在\Person{约雅金}执政期间，\Person{拿步高}王在\Place{亚述}七年、\Place{埃及}\Person{瓦弗勒斯}二年、\Place{雅典}执政官\Person{菲利普}元年、第四十八届\Place{奥林匹亚}第一年，将人民掳往\Place{巴比伦}。囚掳持续七十年，结束于\Person{达理阿·希斯塔斯普}二年，他成了\Place{波斯}、\Place{亚述}和\Place{埃及}的国王；我在前面说过，在他统治期间，\Person{哈盖}、\Person{匝加利亚}和十二先知的天使发了预言。大司祭是\Person{约匝达克}的儿子\Person{耶叔亚}。\Person{达理阿}在位第二年（\Person{赫罗多德}说他消灭了\Person{玛占}的势力），\Person{沙耳提耳}的儿子\Person{则鲁巴贝耳}被派去重建并装饰\Place{耶路撒冷}的圣殿。\par

\Place{波斯}时期的总年数如下：\Person{居鲁士}在位三十年；\Person{冈比西斯}十九年；\Person{达理阿}四十六年；\Person{薛西斯}二十六年；\Person{阿尔塔薛西斯}四十一年；\Person{达理阿}八年；\Person{阿尔塔薛西斯}四十二年；\Person{奥库斯}或\Person{阿尔塞斯}三年。\Place{波斯}帝国总年数为二百三十五年。\par

\Place{马其顿}的\Person{亚历山大}在此期间除掉了这位\Person{达理阿}，开始统治。同样，\Place{马其顿}诸王时期计算如下：\Person{亚历山大}十八年；\Person{拉果}的儿子\Person{托勒密}四十年；\Person{托勒密·斐拉德尔弗斯}二十七年；之后是\Person{厄厄尔革特斯}二十五年；之后是\Person{斐罗帕托尔}十七年；之后是\Person{厄丕法乃斯}二十四年；继位者\Person{斐罗密托尔}在位三十五年；之后是\Person{菲斯贡}二十九年；之后是\Person{拉图鲁斯}三十六年；之后是绰号\Person{狄奥尼索斯}的二十九年；最后是\Person{克娄巴特拉}在位二十二年。此后是\Place{卡帕多细亚}人统治十八天。\par

因此，\Place{马其顿}诸王时期总共为三百一十二年零十八天。\par

因此，那些在达理阿·希斯塔斯普时代，约在他统治第二年预言的人——\Person{哈盖}、\Person{匝加利亚}和十二位先知的天使，他们在第四十八届奥林匹亚的第一年预言——被证实比\Person{毕达哥拉斯}更古老，据说\Person{毕达哥拉斯}生活在第六十二届奥林匹亚，也比\Person{泰勒斯}更古老，\Person{泰勒斯}是希腊智者中最年长的，大约生活在第五十届奥林匹亚。那些与\Person{泰勒斯}同列的智者们当时是同时代的，正如\Person{安德龙}在\WorkTitle{三脚架}中所说的。因为\Person{赫拉克利特}晚于\Person{毕达哥拉斯}，在他的书中提到\Person{毕达哥拉斯}。因此，无可争议的是，第一届奥林匹亚（被证实比特洛伊战争晚四百零七年）被发现早于上述先知以及所谓七贤的时代。由此可见，\Person{撒罗满}生活在\Person{墨涅拉俄斯}时代（\Person{墨涅拉俄斯}活跃于特洛伊战争期间），比希腊的智者们早了许多年。至于\Person{梅瑟}比他早多少年，我们已在前面所说的话中表明了。绰号波利希斯托尔的\Person{亚历山大}，在他关于犹太人的著作中，抄录了\Person{撒罗满}给埃及国王\Person{瓦弗瑞斯}和提洛的腓尼基王的信件，以及他们给\Person{撒罗满}的回信；其中显示，\Person{瓦弗瑞斯}派了八万埃及人给他建造圣殿，另一位王也派了同样多的人，还有一位\Place{提洛}的工匠，一位犹太母亲所生的儿子，属于丹支派，如那里所记，名叫\Person{许佩龙}。此外，\Place{雅典}的\Person{奥诺玛克里托斯}，据说被认为是\Person{俄耳甫斯}诗作的作者，被确定生活在庇西特拉图家族统治时期，大约第五十届奥林匹亚。而\Person{俄耳甫斯}曾与\Person{赫拉克勒斯}航海，是\Person{穆赛俄斯}的弟子。\Person{安菲翁}比特洛伊战争早两代。\Person{德摩多科斯}和\Person{费弥俄斯}晚于特洛伊陷落；因为他们以弹奏里拉琴闻名，前者在费埃克斯人中，后者在求婚者中。而归于\Person{穆赛俄斯}的\WorkTitle{神谕}据说是\Person{奥诺玛克里托斯}的作品，\Person{俄耳甫斯}的\WorkTitle{调酒缸}是\Place{赫拉克利亚}的\Person{佐皮鲁斯}的作品，\WorkTitle{下冥府}是\Place{萨摩斯}的\Person{普罗迪科斯}的作品。\Place{希俄斯}的\Person{伊翁}在\WorkTitle{三篇}中叙述说，\Person{毕达哥拉斯}将他自己的某些作品归于\Person{俄耳甫斯}。\Person{埃庇革涅斯}在他的\WorkTitle{俄耳甫斯诗作}一书中说，\WorkTitle{下冥府}和\WorkTitle{神圣讲论}是毕达哥拉斯派\Person{刻克洛普斯}的作品；而\WorkTitle{佩普洛斯}和\WorkTitle{自然论}是\Person{布戎提诺斯}的作品。有些人也将\Person{忒耳潘德}视为古老者。因此，\Person{赫兰尼可斯}叙述说他生活在\Person{弥达斯}时代；但\Person{法尼阿斯}将\Place{勒斯沃斯}的\Person{勒斯刻斯}置于\Person{忒耳潘德}之前，认为\Person{忒耳潘德}比\Person{阿耳喀罗科斯}年轻，并叙述\Person{勒斯刻斯}与\Person{阿克提诺斯}竞赛并获胜。\Place{吕底亚}的\Person{克珊托斯}说他大约生活在第十八届奥林匹亚；同样，\Person{狄奥尼修斯}说\Place{塔索斯}岛建于第十五届奥林匹亚：所以清楚的是，\Person{阿耳喀罗科斯}在第二十届奥林匹亚之后已经为人所知。他因此叙述了\Place{马格内特斯}最近被毁灭。\Person{西摩尼得斯}被安排在\Person{阿耳喀罗科斯}时代。\Person{卡利诺斯}不算太老；因为\Person{阿耳喀罗科斯}提到\Place{马格内特斯}已被毁灭，而后者提到它处于繁荣。\Place{科林斯}的\Person{欧墨洛斯}更古老，据说他遇到了建立\Place{叙拉古}的\Person{阿耳喀亚斯}。\par

我们之所以提到这些，是因为史诗循环诗人被置于最远古之列。此外，在希腊人中，据说已经出现了许多占卜者，例如巴基德斯，一位是\Place{玻俄提亚}人，另一位是\Place{阿卡迪亚}人，他们向许多人发布了众多预言。凭借\Place{雅典}的\Person{安菲勒托斯}的建议，他指出了进攻的时机，\Person{庇西特拉图}也巩固了他的统治。因为我们可以略过\Place{克里特}的\Person{科墨忒斯}、\Place{塞浦路斯}的\Person{喀倪剌斯}、\Place{色萨利}的\Person{阿德墨托斯}、\Place{库瑞涅}的\Person{阿里斯泰阿斯}、\Place{雅典}的\Person{安菲阿拉俄斯}、\Place{科基拉}的\Person{提摩克塞俄斯}、\Place{福基斯}的\Person{德马涅托斯}、\Place{特斯皮亚}的\Person{埃庇革涅斯}、\Place{卡律斯托斯}的\Person{尼基阿斯}、\Place{色萨利}的\Person{阿里斯托}、\Place{迦太基}的\Person{狄奥尼修斯}、\Place{科林斯}的\Person{克勒奥丰}、\Person{喀戎}的女儿\Person{希波}、\Person{玻奥}、\Person{曼托}，以及众多西比尔——萨默斯人、科洛丰人、库迈人、埃律特莱人、皮提亚人、塔拉克桑德里亚人、马其顿人、色萨利人和塞斯普罗提亚人。还有\Person{卡尔卡斯}和\Person{莫普索斯}，他们生活在特洛伊战争期间。然而\Person{莫普索斯}更年长，曾随阿尔戈英雄航海。据说\Place{库瑞涅}的\Person{巴图斯}编写了所谓的\Person{莫普索斯}\WorkTitle{占卜术}。\Person{多罗透斯}在第一卷\WorkTitle{总论}中叙述说，\Person{莫普索斯}是\Person{阿尔库翁}和\Person{科罗涅}的弟子。伟大的\Person{毕达哥拉斯}始终专注于预测，还有\Place{许珀耳玻瑞亚}的\Person{阿巴里斯}、\Place{普罗孔涅索斯}的\Person{阿里斯泰阿斯}、来到\Place{斯巴达}的\Place{克里特}人\Person{厄庇墨尼得斯}、\Place{米底}的\Person{琐罗亚斯德}、\Place{阿格里真托}的\Person{恩培多克勒}、\Place{拉刻代蒙}的\Person{福尔弥翁}；还有\Place{塔索斯}的\Person{波吕阿拉托斯}和\Place{叙拉古}的\Person{恩培多提摩斯}；此外特别还有\Place{雅典}的\Person{苏格拉底}。\Quote{因为，}他在\WorkTitle{泰格斯}中说，\Quote{我自孩童时代就由天命获得了一个超自然的提示。这是一种声音，当它来临，阻止我将要做的事，却从不劝勉。}福基斯人的僭主\Person{厄克塞克斯托斯}戴着两枚魔戒，通过它们互相撞击发出的声音来确定行动的恰当时机。然而，他仍然被背叛谋杀而死，尽管事先被声音警告，正如\Person{亚里士多德}在\WorkTitle{福基斯政制}中所说。\par

在这些曾以凡人身份生活在埃及人中、却被人的意见奉为神祇的，有底比斯的\Person{赫尔墨斯}和孟菲斯的\Person{阿斯克勒庇俄斯}；又如\Person{欧里庇得斯}所说，底比斯的\Person{提瑞西阿斯}和\Person{曼托}。还有伊利昂的\Person{海伦诺斯}、\Person{拉奥孔}、\Person{俄诺涅}和\Person{克瑞努斯}。因为据说\Person{克瑞努斯}是\Person{赫拉克勒斯}后裔之一，是一位著名的先知。另一位是埃利斯的\Person{伊阿摩斯}，\Person{伊阿摩斯}族人便出自他；还有阿尔戈斯和墨伽拉的\Person{波吕伊多斯}，悲剧中曾提到他。何必列举独眼巨人的先知\Person{忒勒摩斯}，他预言了\Person{波吕斐摩斯}关于\Person{奥德修斯}漂泊之事；或是雅典的\Person{奥诺玛克里托斯}；或是与七将一同征战底比斯的\Person{安菲阿拉俄斯}，据传比特洛伊陷落早一代；或是刻法勒尼亚的\Person{特奥克吕墨诺斯}；或是卡里亚的\Person{特尔弥索斯}；或是西西里的\Person{伽勒乌斯}？\par

除了这些，还有：与阿尔戈英雄同行的\Person{伊德蒙}；德尔斐的\Person{斐摩诺厄}；潘菲利亚的\Person{阿波罗}与\Person{曼托}之子\Person{莫普索斯}；西里西亚的\Person{安菲阿拉俄斯}之子\Person{安菲洛科斯}；阿卡尔纳尼亚人的\Person{阿尔克迈翁}；提洛岛的\Person{阿尼阿斯}；与\Person{亚历山大}同行的泰尔墨索斯的\Person{阿利斯坦德}。\Person{菲洛科罗斯}也在其著作\WorkTitle{论占卜}第一卷中记载，\Person{俄耳甫斯}是一位先知。\Person{特奥彭波斯}、\Person{埃福罗斯}和\Person{提迈乌斯}则写到一位名叫\Person{奥尔塔戈拉斯}的先知；正如萨摩斯的\Person{皮托克勒斯}在\WorkTitle{意大利志}第四卷中写到\Person{盖乌斯·尤利乌斯·尼波斯}一样。\par

但这些\Quote{窃贼和强盗}（正如圣经所说）大多是靠观察和或然率预言，如同医生和占卜者凭自然征兆判断；另有一些是被魔鬼激动，或受水、烟气、特殊空气干扰。但在希伯来人中，先知们是受天主的德能和默感所驱。在法律之前，\Person{亚当}就女人和给受造物命名的事作了预言；\Person{诺厄}宣讲了悔改；\Person{亚巴辣罕}、\Person{依撒格}和\Person{雅各伯}就未来和当前的事给出了许多清晰的言论。与法律同时代的有\Person{梅瑟}和\Person{亚郎}；此后预言的有\Person{若苏厄}（纳温之子）、\Person{撒慕尔}、\Person{加得}、\Person{纳堂}、\Person{阿希雅}、\Person{舍玛雅}、\Person{耶胡}、\Person{厄里亚}、\Person{米该亚}、\Person{阿布迪乌}、\Person{厄里叟}、\Person{阿巴多奈}、\Person{亚毛斯}、\Person{依撒意亚}、\Person{欧瑟亚}、\Person{约纳}、\Person{约厄耳}、\Person{耶肋米亚}、\Person{索福尼亚}（布齐之子）、\Person{厄则克耳}、\Person{乌黎雅}、\Person{哈巴谷}、\Person{纳鸿}、\Person{达尼尔}、\Person{米撒尔}（他写了三段论）、\Person{哈盖}、\Person{匝加利亚}，以及十二小先知中的\Person{天使}。总共是三十五位先知。女性中（因为她们也发了预言）有\Person{撒辣}、\Person{黎贝加}、\Person{米黎盎}、\Person{德波辣}和\Person{胡耳达}。\par

在同一时期内，\Person{若翰}预言直到救恩的洗礼；基督诞生后，\Person{亚纳}和\Person{西默盎}也发了预言。因为\Person{若翰}的父亲\Person{匝加利亚}，据福音记载，在他儿子之前就已预言。现在让我们从\Person{梅瑟}开始编制希腊人的年表。\par

从\Person{梅瑟}诞生到犹太人出离埃及，八十年；到\Person{梅瑟}去世，另有四十年。出埃及发生在\Person{伊那科斯}时代，早于\Person{索提斯}的漂流，\Person{梅瑟}离开埃及在三百四十五年前。从\Person{梅瑟}统治和\Person{伊那科斯}时期到\Person{丢卡利翁}的洪水（我指第二次淹没）和\Person{法厄同}的烈火（这些事件发生在\Person{克罗托普斯}时代），共列举了四十代（每三代算一个世纪）。从洪水到伊达山的烈火和铁器的发现以及伊达山的达克提利人，按照\Person{特拉叙洛斯}，是七十三年；从伊达山烈火到\Person{伽倪墨得斯}被掳，六十五年。从此时到\Person{珀耳修斯}远征（\Person{格劳科斯}为纪念\Person{墨利刻耳忒斯}设立了地峡运动会），十五年；从\Person{珀耳修斯}远征到特洛伊建城，三十四年。从此时到阿尔戈号航行，六十四年。从此时到\Person{忒修斯}和\Person{弥诺陶洛斯}，三十二年；然后到七将攻底比斯，十年。到\Person{赫拉克勒斯}为纪念\Person{珀罗普斯}设立的奥林匹克竞赛，三年；到阿玛宗人远征雅典和\Person{忒修斯}掳走\Person{海伦}，九年。从此时到\Person{赫拉克勒斯}被神化，十一年；然后到\Person{亚历山大}掳走\Person{海伦}，四年。从特洛伊陷落到\Person{埃涅阿斯}下船并建立剌维尼乌姆，十年；到\Person{阿斯卡尼俄斯}统治，八年；到赫拉克勒斯后裔的迁移，六十一年；到\Person{伊菲托斯}的奥林匹克，三百三十八年。\Person{厄拉托斯特涅斯}如此记载日期：\Quote{从特洛伊陷落到赫拉克勒斯后裔迁移，八十年。从此时到伊奥尼亚殖民地建立，六十年；之后到\Person{吕库尔戈斯}的监护期，一百五十九年；到第一届奥林匹亚的第一年，一百零八年。从这届奥林匹亚到\Person{薛西斯}入侵，二百九十七年；从那时到伯罗奔尼撒战争开始，四十八年；到战争结束和雅典人战败，二十七年；到琉克特拉战役，三十四年；此后到\Person{腓力}去世，三十五年；之后到\Person{亚历山大}去世，十二年。}\par

再者，从第一届奥林匹亚到罗马建城，有些人认为是二十四年；此后到驱逐国王、建立执政官时，约二百四十三年。从巴比伦陷落到\Person{亚历山大}去世，一百八十六年。从此时到\Person{奥古斯都}获胜（\Person{安东尼}在亚历山大港自杀），二百九十四年，当时\Person{奥古斯都}第四次担任执政官。从此时到\Person{多米提安}在罗马设立的赛会，一百一十四年；从第一届赛会到\Person{康茂德}去世，一百一十一年。\par

有些人从\Person{刻克洛普斯}到\Person{马其顿的亚历山大}计为一千八百二十八年；从\Person{德摩丰}计为一千二百五十年；从特洛伊陷落到\Person{赫拉克勒斯后代}的远征计为一百二十年或一百八十年。从此时到\Person{欧埃涅托斯}在\Place{雅典}担任执政官——据说亚历山大正是在此期间进军亚洲，据\Person{法尼亚斯}——为七百五十年；据\Person{埃福罗斯}为七百三十五年；据\Person{提迈奥斯}和\Person{克利塔克斯}为八百二十年；据\Person{埃拉托色尼}为七百七十四年。同样，\Person{杜里斯}从特洛伊陷落到亚历山大进军亚洲计为一千年；从那时到\Person{赫格西亚斯}担任执政官——亚历山大在此期间去世——为十一年。从这一日期到\Person{日耳曼尼库斯·克劳狄·恺撒}统治时期为三百六十五年。从那时起到\Person{科莫多}去世的累计年数显而易见。\par

在希腊时期之后，根据异族所计算的年代，应分配非常长的间隔。\par

从\Person{亚当}到\Emph{洪水}共计二千一百四十八年零四天。从\Person{闪}到\Person{亚伯拉罕}为一千二百五十年。从\Person{以撒}到分地为六百一十六年。然后从民长到\Person{撒慕尔}为四百六十三年零七个月。民长之后，列王时期为五百七十二年零六个月零十天。\par

在这些时期之后，波斯帝国为二百三十五年。然后马其顿帝国，直到\Person{安多尼}去世，为三百一十二年零十八天。此后，罗马帝国直到\Person{科莫多}去世持续了二百二十二年。\par

然后，从七十年充军和人民回归故土，到\Person{维斯帕先}时期的充军，共计四百一十年。最后，从\Person{维斯帕先}到\Person{科莫多}去世，确定为一百二十一年零六个月零二十四天。\par

\Person{德默特琉}在其著作\WorkTitle{论犹太列王}中说，\Person{犹大}、\Person{本雅明}和\Person{肋未}支派没有被\Person{散乃黑黎布}掳去；但从这次充军到\Person{拿步高}从\Place{耶路撒冷}进行的最后一次充军，有一百二十八年零六个月；从十个支派从\Place{撒玛黎雅}被掳到\Person{托勒密四世}时期，为五百七十三年零九个月；而从\Place{耶路撒冷}充军发生之时起，为三百三十八年零三个月。\par

\Person{斐罗}本人对列王的记载与\Person{德默特琉}不同。\par

此外，\Person{欧颇勒摩}在一部类似著作中说，从\Person{亚当}到\Person{托勒密·德默特琉}第五年——他在\Place{埃及}统治了十二年——的所有年份相加为五千一百四十九年；从\Person{梅瑟}带领犹太人出\Place{埃及}到上述日期，共计二千五百八十年。从这时到\Place{罗马}\Person{卡尤斯·多米提安}和\Person{卡西安}担任执政官，计算为一百二十年。\par

\Person{欧福洛斯}和许多其他历史学家说有七十五个民族和语言，因为他们听到了\Person{梅瑟}的陈述：\BibleQuote{凡从\Person{雅各伯}所生、下到\Place{埃及}的灵魂，共七十五人。}根据真实的计算，似乎有七十二种通用方言，正如我们的圣经所记载。其余通俗语言是由两种、三种或更多种方言混合而成的。方言是一种言语方式，展现出某一特定地区的特色，或展现出某一民族特有或共同的言语方式。希腊人说，他们中间有五种方言：阿提卡方言、伊奥尼亚方言、多利亚方言、伊奥利亚方言，第五种是通用方言；而异族的语言数不胜数，不称为方言，而称为语言。\par

\Person{柏拉图}也将一种方言归诸神，主要是从梦和神谕，尤其是从附魔者中形成这一推测，这些人不说自己的语言或方言，而是说占据他们的魔鬼的语言。他还认为无理性的生物也有方言，同类者能理解。因此，当一头大象掉进泥里并吼叫时，任何附近的大象看到发生了什么，不久后返回，带来一群大象，救出掉进去的那头。据说在\Place{利比亚}，一只蝎子如果没有成功蜇到人，就会离开并带回更多；它们一个接一个像链条一样挂着，以此方式试图成功实现其狡猾的企图。\par

无理性的生物并不使用模糊的暗示，或通过采取特定姿态来暗示其意义，而是（我认为）使用它们自己的方言。还有一些人说，如果一条鱼被钓上后因断线而逃脱，当天在同一地点就不会再钓到同种鱼。但最初和通用的异族方言具有天然的术语，因为人们也承认用异族语言念诵的祈祷更有效力。\Person{柏拉图}在\WorkTitle{克拉底鲁}中，想要解释\OriginalTerm{火}{\GreekText{πῦρ}}时，说它是一个异族语词。他因此证实，弗里吉亚人使用这个词，只有轻微变化。\par

在我看来，在这些细节之后，无需阻碍陈述罗马皇帝的在位时期，以证明救主的降生。\Person{奥古斯都}四十三年；\Person{提比略}二十二年；\Person{卡尤斯}四年；\Person{克劳狄}十四年；\Person{尼禄}十四年；\Person{伽尔巴}一年；\Person{维斯帕先}十年；\Person{提图斯}三年；\Person{多米提安}十五年；\Person{内尔瓦}一年；\Person{图拉真}十九年；\Person{哈德良}二十一年；\Person{安敦尼}二十一年；同样再次，\Person{安敦尼}和\Person{科莫多}三十二年。从\Person{奥古斯都}到\Person{科莫多}，总共二百二十二年；从\Person{亚当}到\Person{科莫多}去世，为五千七百八十四年零两个月零十二天。\par

有些人以下列方式记载罗马皇帝的在位时期：—\par

盖乌斯·尤利乌斯·恺撒（\Person{盖乌斯·尤利乌斯·恺撒}）三年四个月零五天；此后，奥古斯都（\Person{奥古斯都}）在位四十六年四个月零一天；然后提比略（\Person{提比略}）二十六年六个月零十九天；继位的是盖乌斯·恺撒（\Person{盖乌斯·恺撒}），在位三年十个月零八天；克劳狄（\Person{克劳狄}）继位十三年八个月零二十八天；尼禄（\Person{尼禄}）在位十三年八个月零二十八天；伽尔巴（\Person{伽尔巴}）七个月零六天；奥托（\Person{奥托}）五个月零一天；维特里乌斯（\Person{维特里乌斯}）七个月零一天；韦斯帕芗（\Person{韦斯帕芗}）十一年十一个月零二十二天；提图斯（\Person{提图斯}）两年零两个月；图密善（\Person{图密善}）十五年八个月零五天；涅尔瓦（\Person{涅尔瓦}）一年四个月零十天；图拉真（\Person{图拉真}）十九年七个月零十天；哈德良（\Person{哈德良}）二十年十个月零二十八天；安东尼（\Person{安东尼}）二十二年三个月零七天；马可·奥勒留·安东尼（\Person{马可·奥勒留·安东尼}）十九年零十一天；康茂德（\Person{康茂德}）十二年九个月零十四天。\par

因此，从尤利乌斯·恺撒到康茂德去世，共计二百三十六年零六个月。而从奠基罗马的罗慕路斯到康茂德去世，总计九百五十三年零六个月。我们的主诞生于奥古斯都统治下首次人口普查令颁布的第二十八年。为证明此事真实，路加福音中记载如下：\BibleQuote{在提比略·恺撒执政第十五年，主的话传到匝加利亚的儿子若翰那里。}同书中又记载：\BibleQuote{耶稣受圣洗时，年约三十岁，}等等。此外，他只需宣讲一年的必要性也有记载：\BibleQuote{他派遣我向贫穷人传报喜讯，宣布主恩赐的禧年。}先知如此说，福音也是如此。因此，在提比略的第十五年和奥古斯都的第十五年中，共满了三十年，直到他受难之时。从他受难到耶路撒冷毁灭，共四十二年零三个月；从耶路撒冷毁灭到康茂德去世，共一百二十八年十个月零三天。因此，从基督诞生到康茂德去世，总计一百九十四年一个月零十三天。另有一些人不仅确定了我们主诞生的年份，也确定了日子；他们说那发生在奥古斯都在位第二十八年，帕克洪月二十五日。而巴西里德的追随者将他的圣洗日作为节日庆祝，整夜诵读经文。\par

他们说那是提比略·恺撒在位第十五年，图比月十五日；另有人说是在同月的十一日。论及他的圣洗日，有些人极其精确地指出那发生在提比略在位第十六年，法墨诺斯月二十五日；另有人说法尔穆提月二十五日；还有人说救主在法尔穆提月十九日受难。此外，另有人说他诞生于法尔穆提月二十四日或二十五日。\par

我们仍需将以下内容加入我们的年表——即达尼尔所指从耶路撒冷荒凉算起的日子，以及韦斯帕芗统治的七年七个月。因为奥托、伽尔巴和维特里乌斯的十七个月零十八天又加了两年；结果共三年六个月，正是达尼尔先知所说的\BibleQuote{半周，}因为他预言，从尼禄的可憎之物立于圣城直到城被毁灭，共有二千三百天。因为下文所附的宣告如此显示：\BibleQuote{这异象要持续多久？即除掉常献的祭，设立那招致荒凉的可憎之物，并将圣所与军旅践踏在脚下？他对我说：直到晚上和早晨，共二千三百天，然后圣所将被夺去。}\BibleRef{达 8:13}\par

因此，这二千三百天成了六年零四个月，其中一半由尼禄掌权，那是半周；另一半由韦斯帕芗与奥托、伽尔巴和维特里乌斯统治。为此，达尼尔说：\BibleQuote{凡到达一千三百三十五天的，是有福的。}\BibleRef{达 12:12}因为战争持续到这些日子为止，之后便止息了。这一数字从后续章节得到证明，如下所说：\BibleQuote{从除掉常献的祭、设立那招致荒凉的可憎之物的时候起，必有一千二百九十天。凡等待而到达一千三百三十五天的，是有福的。}\BibleRef{达 12:11}\par

犹太史学家弗拉维奥·约瑟夫斯（\Person{弗拉维奥·约瑟夫斯}）在计算年代时说，从梅瑟到达味共五百八十五年；从达味到韦斯帕芗第二年共一千一百七十九年；从那以后到安东尼第十年共七十七年。因此，从梅瑟到安东尼第十年，总共二千一百三十三年。\par

另一些人从伊那科斯（\Person{伊那科斯}）和梅瑟算起直到康茂德去世，有的说共三千一百四十二年，有的说二千八百三十一年。\par

而在玛窦福音中，从亚巴郎开始的族谱一直延续到主的母亲玛利亚。\BibleQuote{因为从亚巴郎到达味共十四代；从达味到流徙巴比伦共十四代；从流徙巴比伦到基督又是十四代，}——三个神秘间隔，完成于六周。\BibleRef{玛 1:17}\par

% structure: p0199 source=h2 target=subsection reason=relative-heading-rank; base=h2

\subsection{第二十二章 论旧约的希腊文译本}

以上便是许多人以不同方式陈述、并由我们记录的有关日期的细节。\par

据说，法律书和先知书的圣经都是在拉戈斯之子\Person{托勒密}（另有人说是\Person{托勒密·费拉德尔弗斯}）统治时期，由希伯来方言译成希腊语言的；\Person{德米特里·法勒雷乌斯}以最大的热忱投入此项工作，并在翻译素材上力求精确。当时，马其顿人仍然占据着亚细亚，国王渴望用所有著作来装饰他在亚历山大里亚的图书馆，便请求耶路撒冷人民将他们拥有的预言书译成希腊方言。他们作为马其顿的臣民，从他们当中那些品格最高尚的人中选出了七十位精通圣经、熟练希腊方言的长老，将神圣的经卷送给了他。每位长老分别翻译了各卷先知书，所有译文互相比较后，它们在意义和措辞上都一致。因为这是天主的计划，为希腊人的耳朵而实现。这并非有违于赐予预言的天主的默感，祂也产生了译文，使之犹如希腊的预言。由于圣经在\Person{拿步高}被掳时已遭毁灭，在波斯王\Person{阿尔塔薛西斯}时代，司祭兼肋未人\Person{厄斯德拉}在预言的默感下，恢复了全部古代圣经。\Person{阿里斯托布鲁}在写给\Person{斐洛米托}的第一卷书中写道：\Quote{\Person{柏拉图}遵循了赐予我们的法律，并显然研究了其中所说的一切。}在\Person{德米特里}之前，早在\Person{亚历山大}和波斯人的统治之前，另有人翻译了我们同胞希伯来人离开埃及的事迹、所有发生之事的传闻、他们占领土地的经过以及整套法典，所以很清楚，上述哲学家从这一来源获益甚多，因为他非常博学，\Person{毕达哥拉斯}也是如此，他从我们的书中将许多内容移植到了自己的教义体系中。毕达哥拉斯派的哲学家\Person{努美纽斯}明确写道：\Quote{\Person{柏拉图}若不是讲阿提卡希腊语的\Person{梅瑟}，又是什么呢？}这位\Person{梅瑟}是一位神学家和先知，并且如同一些人所说，是神圣法律的解释者。他的家族、事迹及一生，均由可信赖的圣经本身所记述，不过我们也需尽己所能加以陈述。\par

% structure: p0202 source=h2 target=subsection reason=relative-heading-rank; base=h2

\subsection{第二十三章 \Person{梅瑟}的年龄、诞生与生平}

梅瑟原属加色丁家族，生于埃及，其祖先因长期饥荒从巴比伦迁入埃及。他出生在第七代，接受了王室教育，其历史事迹如下。希伯来人在埃及繁衍成为一大民族，该国国王因惧怕他们人数众多而造反，遂下令希伯来人所生女婴皆可抚养（因女子不适宜作战），但男婴则要消灭，因他猜忌强健的年轻人。但那孩子俊美，父母出于亲情胜过君王的残暴，暗中抚育了三个月。最后，因担心自己与孩子同遭毁灭，他们用当地生长的纸莎草编了一个篮子，将孩子放入其中，置于沼泽河畔。孩子的姐姐远远站着，要看有什么事发生。在这危急时刻，国王的女儿——她已久未怀孕，渴望得子——那天来到河边沐浴盥洗；听见孩子啼哭，便命人抱来；心生怜悯，便寻找乳母。此时孩子的姐姐跑上前说，若她愿意，可以找来一位刚生产不久的希伯来妇人作乳母。公主同意并命她去做，她便带来孩子的母亲，约定工价，假充他人作乳母。于是王后依词源学为孩子取名\OriginalTerm{梅瑟}{Moses}，因他从\Quote{水}中被拉上来——埃及人称水为\OriginalTerm{慕}{mou}——他原被弃在水中等死。他们称梅瑟为\Quote{从水中（被取）而呼吸的人}。显然，父母在行割损礼时已给孩子取名，名为\OriginalTerm{约阿基姆}{Joachim}。他升天后还有第三个名字，奥秘派称其为\OriginalTerm{默尔希}{Melchi}。到了适当年龄，他向埃及精通这些技艺的人学习算术、几何、诗歌、和声，此外还有医学和音乐；还学习了以符号传达的哲学，即象形铭文中所指示的。其余常规课程，希腊人在埃及如王室子弟般教授他，正如斐洛在其\WorkTitle{梅瑟传}中所说。他还学习了埃及的文学，并从加色丁人和埃及人那里获得天文知识；因此在\BibleBook{宗徒}{大事录}中，他被记载\Quote{学习了埃及人的一切智慧}\BibleRef{宗 7:22}。尤波勒摩斯在其\WorkTitle{论犹大地的诸王}一书中说：\Quote{梅瑟是第一位智者，也是第一位将文法传授给犹太人的，腓尼基人从犹太人那里接受，希腊人又从腓尼基人那里接受。}他投身于他们的哲学，智慧得以增长，热切依恋亲族和祖先的传统教育，直至他击杀了那无理攻击希伯来人的埃及人。奥秘派说他仅以言语就杀了那埃及人；诚然，\BibleBook{宗徒}{大事录}记载伯多禄以言语击杀了那些私吞田价并说谎的人\BibleRef{宗 5:1}。阿尔塔帕诺斯在其\WorkTitle{论犹太人}一书中也记载：\Quote{梅瑟因人民要求离开埃及而被埃及王凯内弗瑞斯囚禁，夜间监狱因天主的介入而开启，梅瑟出来，来到王宫，站在睡着的国王面前，唤醒了他；国王被此事所震惊，命梅瑟告诉他那派他来的天主的名字；梅瑟俯身在他耳边说了；国王听后倒地无声，但被梅瑟扶起，恢复了知觉。}至于梅瑟的教育，我们会在厄则克耳——犹太悲剧作家——的戏剧\WorkTitle{出谷纪}中找到一致的描述。他以梅瑟的口吻写道：\par

\Quote{因见吾族繁盛，\newline{}法老设诡谋陷害，\newline{}残忍迫我等于砖窑，\newline{}或于苦役建筑中折磨；\newline{}修筑城塔，耗竭民力。\newline{}复向希伯来族颁令：\newline{}凡男婴皆投尼罗深流。\newline{}慈母怀胎，藏我三月\newline{}（后亲自告我），然后\newline{}以华服装饰，置于\newline{}尼罗河畔芦苇深处，\newline{}我姊米黎盎远望。\newline{}国王之女偕侍女\newline{}前来清流沐浴，\newline{}看见即取我抱起，\newline{}认出我是希伯来人。\newline{}我姊米黎盎跑向公主，\newline{}说道：\InnerQuote{可否容我去寻\newline{}一乳母，为你抚养此子\newline{}于希伯来妇人中？}公主\newline{}应允。少女速往母处\newline{}告知，母亲即刻前来。\newline{}我生母抱我于怀。\newline{}国王之女说：\InnerQuote{为我乳养此子，\newline{}我必付你工价。}\newline{}她赐我名梅瑟，\newline{}因她从河岸水中\newline{}拉出拯救了我。\newline{}童年岁月飞逝，\newline{}我母带我入王宫\newline{}公主居所，\newline{}陈明我身世及天主鸿恩。\newline{}少时我便受王室养育，\newline{}习练诸般王子技艺，\newline{}宛若公主亲子。\newline{}待时日循环，\newline{}我离开王宫。}\par

继而，在叙述了希伯来人与埃及人的搏斗以及将埃及人埋在沙中之后，关于另一场争执，他说：\par

\Quote{为何击打比你软弱者？\newline{}他答：谁立你作我们的法官\newline{}或首领？难道你要杀我，\newline{}如昨日杀那人？\newline{}我惊恐说：这事怎会知道？}\par

然后他逃离埃及，牧放羊群，如此预先为牧养统治而受训练。因为对于注定要统治和平羊群——人类——的人来说，牧羊人的生活是为王权作的准备，正如狩猎对于天性好战者一样。然后天主带领他领导希伯来人。埃及人虽屡受警告，仍不聪明；希伯来人却旁观他人所受的灾祸，在安全中学习天主的能力。当埃及人因愚顽不信而忽视那能力的果效时，正如经上说：\Quote{孩子们知道了}所行的事；后来希伯来人出发，从埃及人那里带走了大量战利品，并非出于贪婪（如诽谤者所说），因为天主并没有说服他们贪恋别人的东西。首先，他们取了长期以来为埃及人服务的报酬；其次，他们以某种方式报复了埃及人，通过夺取战利品来惩罚他们的贪婪，正如埃及人通过奴役他们来对待希伯来人一样。那么，无论是如战争中所行的那样，他们行使征服者的权利，夺走敌人的财产（正如胜利者对待失败者一样——而且有正当的敌对理由：希伯来人是因饥荒作为恳求者来到埃及人那里；埃及人却将他们的客人沦为奴隶，强迫他们像俘虏一样服侍，不给他们任何报酬）；还是在和平时期，他们违背那些长期不给报酬反而掠夺他们的人的意愿，以战利品作为工资，[都是一样的。]\par

% structure: p0208 source=h2 target=subsection reason=relative-heading-rank; base=h2

\subsection{第二十四章 梅瑟如何履行军事领袖的职责}

因此，我们的梅瑟是一位先知、立法者、精通战术和战略者、政治家、哲学家。他如何是先知，我们将在论及预言时再述。战术属于军事指挥，统帅军队的能力是王权属性之一。立法也是王权职能之一，如同司法权一样。\par

王权有一种是神圣的——即按照天主及其圣子，由之供给地上美物以及外在和圆满福乐的那一种。\Quote{因为，}经上说，\Quote{你们要求大事，小事也会加给你们。}第二种王权低于纯粹理性和神圣的管理，它只是在治理任务中运用灵魂的高度勇气；以此方式，\Person{赫拉克勒斯}治理了阿尔戈斯人，\Person{亚历山大}治理了马其顿人。第三种只追求一个目的——征服和颠覆；但将征服用于善恶目的，不属于此类统治。波斯人进攻希腊时就是这种意图。因为好斗仅是激情的产物，只为统治而获取权力；而爱好善是运用高贵精神于高尚目的之灵魂的特征。第四种，最坏的一种，是按照激情冲动而施行的统治，如\Person{萨尔达纳帕鲁斯}那样，以极度满足激情为目的。但王权的工具——既是凭德性获胜，也是凭武力获胜的工具——是管理能力（或策略）。它随本性和质料而变化。在武器和战斗动物的情况下，管理的能力是灵魂和心智，通过有生命和无生命的手段；在我们以德性征服的灵魂激情情况下，管理的能力是理性，通过盖上节制和克己的印记，连同圣洁和真知，使整体结果归于对天主的虔敬。因为智慧是那些如此实践德性之人的规范；神性之事由智慧管理，人事由政治管理——万有皆由王权管理。因此，按照法律治理、拥有管理自愿臣民技能的人就是君王。这样的主，接受所有相信祂并通过祂而来的人。因为圣父将一切交托并臣服于我们的君王基督，\BibleQuote{以致耶稣之名，天上、地上和地下的一切，无不屈膝，并无不承认耶稣基督是主，以光荣天主圣父。}\BibleRef{斐 2:10}\par

将军领导权包含三个概念：谨慎、进取以及二者的结合。每个概念又由三件事组成：它们通过言语、行动或二者同时发挥作用。所有这些都可以通过说服、强迫、或以惩罚应受罚者造成伤害的方式来实现；而这可以通过做正确的事、说假话、说真话、或同时采用其中任何手段来实现。\par

现在，希腊人曾享有从梅瑟领受这一切以及有关如何使用其中每一种知识的益处。为举例说明，我将引述一两个领导权的实例。梅瑟在领百姓出来时，猜疑埃及人会追赶，就离开那条简短直接的路，转向旷野，且多半在夜间行军。因为希伯来人正是在大旷野中，经过漫长的时间，受忍耐的智慧训练（他们所处的环境逼使他们如此），被陶冶成相信独一天主。因此，梅瑟的谋略显示出：必须在危险来临之前辨明什么会有助益，然后加以应对。结果正如他所猜疑的，因为埃及人带着马匹和战车追赶，但很快就被海水冲来淹没，连同马匹和战车一同覆灭，连一个余留者也没有留下。之后，伴随他们的火柱（因为它走在前头作向导）在夜间引领希伯来人穿越无人踏足之地，以辛劳和困苦训练并磨炼他们，使他们刚毅和忍耐，以致在他们经历了看似可怕的困难之后，那从无路旷野中他正领他们前去的福地之益处，便显明了。此外，他击败并杀灭了那地敌对的居民，从荒芜崎岖的行军路线（他的将军才能如此卓越）向他们突袭。因为夺取那些敌对部落之地，是一项技巧和策略的工作。\par

察觉到这一点，雅典将军米太亚德——曾在马拉松战役中击败波斯人——以如下方式效法了它。他穿过一片无路的旷野，在夜间率领雅典人前进，避开了派来监视他的蛮族人。因为从雅典人叛逃的希庇阿斯由于熟悉地形，带领蛮族人进入阿提卡，夺取并占据了有利据点。因此任务就是躲避希庇阿斯。所以米太亚德正确地穿过旷野，在夜间攻击由达特斯指挥的波斯人，带领他的士兵取得胜利。\par

还有，当特拉绪布洛斯从费勒带领流亡者回来，并希望不被察觉时，一根柱子成了他穿过无路之地的向导。给特拉绪布洛斯在夜间，天空无月且暴风雨，有一团火显现引路，它安全地引导他们，把他们留在慕尼基亚附近，那里现时是光明赐予者（福斯福洛斯）的祭坛。\par

因此，借着这样一个实例，让我们对希腊人说的话变得可信，即：全能天主有可能使那在行军时作向导的火柱在夜间走在希伯来人前面。也有人说，在某神谕中——\par

\Quote{给特拜人的一根柱子是令人喜悦的巴克科斯，}\par

出自希伯来人的历史。欧里庇得斯也在\WorkTitle{安提俄佩}中写道——\par

\Quote{在内室中，那牧人，\newline{}戴着常春藤的花冠，是欧伊俄斯神的柱子。}\par

那柱子表明天主是不可描绘的。光柱除了指明天主不可表现之外，也显示神性的稳固与恒久存续，以及祂不变而不可言说的光。因此，在制造像的形态发明之前，古人竖起柱子，并尊崇它们为神性的雕像。所以，撰写了\WorkTitle{福罗尼斯}的人写道——\par

\Quote{卡利托埃，奥林匹克女后之持钥者；\newline{}阿尔戈斯的赫拉，她首先用头带和流苏\newline{}将那女后的高大柱子四围装饰。}\par

此外，\WorkTitle{欧罗皮娅}的作者叙述德尔斐的阿波罗雕像是这样一根柱子，用以下词句：——\par

\Quote{使我们将初果和什一奉献\newline{}挂在神圣的柱子和高大的圆柱上，献给神。}\par

阿波罗，奥秘地解释为\Quote{多数的缺失，}意指独一天主。那么，那如火一样的柱子，以及旷野中的火，是神圣之光的象征，这光借由木（十字架）从大地经过又返回天上；也借由它，理智视觉的恩赐被赐予我们。\par

% structure: p0224 source=h2 target=subsection reason=relative-heading-rank; base=h2

\subsection{第二十五章 柏拉图在制定法律方面是梅瑟的模仿者}

哲学家柏拉图，借助梅瑟的书籍辅助立法，谴责了米诺斯和吕库古的政体，因为它们只以勇敢为目标；而赞扬表达某种统一事物并依照一条法则指导的政体，认为它更为合宜。因为他说，我们理应以力量、尊严和智慧来从事哲学——对同一事物坚持同一见解，毫不动摇，以合乎天界的尊严。因此，他解释了律法中的内容，命令我们注视独一天主并主持正义。关于政治，他说有两种——法律部门和政治部门，严格意义上的政治。\par

他以至高无上的方式，把创造者称为\OriginalTerm{政治家}{\GreekText{ὁ} \GreekText{πολιτικός}}，在他的同名著作\WorkTitle{政治家}（\OriginalTerm{}{\GreekText{ὁ} \GreekText{πολιτικός}}）中；他把那些实践行动与正义生活，兼具沉思之人称为\OriginalTerm{政治家们}{\GreekText{πολιτικοί}}。政治学中被称为\Quote{法律}的那个部分，他划分为行政宽宏与私人良好秩序（他称之为条理），以及和谐与节制——当统治者适合其臣民、臣民顺从统治者时便可看到这些；\Person{摩西}体系孜孜以求达成这种结果。此外，\Person{柏拉图}藉所受的助益断言，法律的部门立足于生成，政治的部门立足于友谊与同意；因此，他随之把\WorkTitle{厄庇诺米斯}中那位通晓一切生成进程（此进程藉行星完成）的哲学家，和另一位天文学家、星体运动及其交感与关联的研究者\Person{提迈欧}，与\Quote{政制}（或\Quote{共和国}）结合在一起。那么，依我之见，政治家与按法律生活之人的目的都是沉思。因此，公共事务必须正确管理，但从事哲学最好。因为智慧之人将全力专注于知识，以善行引导生活，轻视相反之事，追求有助真理的事务。法律并非由法律所规定者（因为所见并非景象），也非每种意见（当然不是恶的意见）。法律乃是善的意见，善即真，真即发现\Quote{真实存有，}并达到它。\Quote{我是自有者，}\Person{摩西}说，\Quote{派遣了我。}依照这善的意见，有些人将法律称为正确理性，它命令当做之事，禁止不当做之事。\par

% structure: p0227 source=h2 target=subsection reason=relative-heading-rank; base=h2

\subsection{第二十六章  \Person{摩西}被恰当地称为神圣立法者，虽次于\Person{基督}，却远胜于希腊的伟大立法者\Person{米诺斯}和\Person{来古格士}。}

因此，法律恰当地说是由\Person{摩西}赐予的，是是非的准则；我们可精确地称它为\OriginalTerm{神圣法令}{\GreekText{θεσμός}}，因为它是由天主藉摩西赐予的。它因此导向天主。\Person{保禄}说：\Quote{法律是因过犯而设立的，直到那蒙应许的后裔来到。}然后，他似乎为解释其意而补充道：\Quote{但在信德来临之前，我们被看守在法律之下，被封闭，}显然是因恐惧，由于罪恶，\Quote{直到那将要显示的信德。因此，法律成了引领我们归向基督的启蒙师，使我们因信德而成义。}真正的立法者是那为灵魂的每一部门分配适合它及其运作之物的人。现在，综言之，\Person{摩西}是一活的法律，由仁慈的圣言治理。因此，他提供了良好的政体，即人在社会生活中正确的纪律。他也处理了司法行政，即那门涉及为正义而纠正违犯者的知识。与之协调的是惩处能力，即对于惩处中应遵守的适度尺度的知识。惩处，就其本身而言，是灵魂的纠正。一言以蔽之，\Person{摩西}的整个体系适合训练那些能够成为良善高尚之人，并搜捕像他们一样的人；这就是指挥的艺术。那能够正确对待被圣言捕获之人的智慧，立法智慧。因为极其君王的性质，拥有并运用这种智慧，\par

因此，惟有智慧的人才是哲学家们所宣称的君王、立法者、将军、正义者、圣洁者、天主所爱者。如果我们从圣经本身发现\Person{摩西}具有这些品性，我们就可以以最确定的信念宣称摩西是真正智慧的。正如我们所说，牧人的技艺在于照顾羊群；因为\BibleQuote{善牧为羊舍掉自己的性命。}\BibleRef{若 10:11}所以我们也要说，立法，由于它主持并照顾人类的羊群，就建立人的德行，尽可能点燃人性中所有的善。\par

而且，如果比喻性地称为属于主的羊群，只不过是一群人的话，那么祂自己就是那唯一羊群的善牧和立法者，\Quote{那听从祂的羊}，那照顾他们的，\Quote{寻找，}并藉法律和圣言找到\Quote{那失去的；}因为事实上，法律是属神的，导向福乐。因为凡由圣神而生的，就是属神的。那真正是立法者的人，不仅宣告善与高贵，而且理解它们。这拥有知识之人的法律是救恩的诫命；或者更确切地说，法律是知识的诫命。因为圣言就是\BibleQuote{天主的德能和天主的智慧。}\BibleRef{格前 1:24} 再者，法律的阐释者就是赐予法律的那一位；神圣诫命的首位阐释者，祂揭示了圣父的怀，即独生子。\par

那么，那些服从法律的人，因为对祂有一些认识，就不能不信或无知于真理。但那不信并表现出厌恶从事法律之工的人，无论别人如何，确实承认他们对真理的无知。\par

那么，希腊人的不信是什么？岂不是他们不愿相信那宣告法律是由天主藉\Person{摩西}赐予的真理，而他们在自己的著作中却尊敬摩西？他们叙述\Person{米诺斯}在九年中通过经常出入宙斯的洞穴从宙斯领受了法律；\Person{柏拉图}、\Person{亚里士多德}和\Person{埃福罗斯}写道，\Person{来古格士}通过不断前往\Place{德尔斐}的\Person{阿波罗}处而接受了立法训练。\Person{赫拉克利亚的卡梅莱昂}，在其著作\WorkTitle{论醉酒}中，以及\Person{亚里士多德}在\WorkTitle{洛克里人的政制}中，提到洛克里的\Person{扎琉库斯}从\Person{雅典娜}那里领受了法律。\par

但那些尽其所能抬高希腊立法权威的人，将其归于神圣根源，模仿梅瑟预言，却因不承认真理及他们当中所述的原型而愚昧无知。\par

% structure: p0234 source=h2 target=subsection reason=relative-heading-rank; base=h2

\subsection{第二十七章　法律即使在纠正与惩罚时，也以人的益处为目标}

那么，谁也不应诋毁法律，仿佛因刑罚之故，它便不美不善。因为驱逐身体疾病的人，难道不被称为恩人吗？那么，那试图将灵魂从不义中解救出来的人，岂不更应被视为朋友，如同灵魂比身体更宝贵一样？此外，为了身体的健康，我们甘受切开、烧灼和药饮；施行这些的人被称为救主和医治者，即便他切除部分，并非出于对病人的怨恨或恶意，而是按医术原则，以免健康部分随之毁灭；无人控诉医术为邪恶。那么，为了灵魂的缘故，我们难道不也同样应当接受流放、惩罚或束缚吗？只要我们能从不义达到正义。\par

因为法律，出于对顺从者的关怀，培育虔诚，规定当行之事，约束各人远离罪恶，甚至对较小的罪也施加刑罚。\par

但当它看到有人处于似乎不可救药、趋向邪恶至极的状态时，则出于对其余人的关怀，以免他们被其毁灭（正如从全身切除一部分），便判处此人死刑，作为最有益健康的方式。\Quote{我们受主审判，}宗徒说，\Quote{受惩戒，以免与世人一同被定罪。}\BibleRef{格前 11:32}因为先知曾预言：\Quote{主惩戒了我，却没有将我交于死亡。}\Quote{为了教导你们认识祂的正义，}经上说，\Quote{祂惩戒、考验你们，使你们在旷野之地饥饿干渴；为使你们心中明白祂的一切法令和判断，正如我今天吩咐你们的；并使你们心中知道，正如人惩戒自己的儿子，上主你们的天主也必惩戒你们。}\par

为证明榜样能纠正人，他直截了当地说：\Quote{精明人见恶人受罚，必严厉自责，因为敬畏上主是智慧之源。}\BibleRef{箴 22:3}\par

但最高、最完美的善，是能引导任何人从行恶转向美德与善行，这正是法律的职分。因此，当人陷入某种无可救药的恶——例如被错谬或贪婪所占据——若将他处死，反而对他有益。因为法律是慈善的，能够使一些人从不义变为义，只要他们肯听从；并能使另一些人脱离现世的恶；对于那些选择度节制与公正生活的人，它引领他们进入不朽。认识法律是好品性的标志。又说：\Quote{恶人不明白法律；但寻求上主的人，必通达一切美善。}\BibleRef{箴 28:5}\par

当然，那管理万物的眷顾，必须是至高且良善的。因为正是两者的能力分施救恩——一种以惩罚来纠正，作为至高者；另一种以行善显示恩惠，作为恩主。你本可以不做悖逆之子，而能从黑暗转入生命，侧耳听受智慧，起初成为天主的法律之奴，然后成为忠信的仆人，敬畏主天主。若有人更进一步，便列入子嗣之中。\par

但当\Quote{爱德遮盖许多罪过，}\BibleRef{伯前 4:8}藉着蒙福希望的圆满实现，我们便可欢迎他，因他已在爱中致富，并被收纳为蒙拣选的嗣子，称为天主所爱的人，同时高唱祈祷：\Quote{愿上主作我的天主。}\par

法律慈善的功效，宗徒在论及犹太人的段落中显示了，如此写道：\Quote{看，你称为犹太人，依靠法律，以天主自夸，知道祂的旨意，能辨别是非，因受法律教导，自信是盲人的向导、黑暗中人的光明、愚昧人的训导者、小孩子的教师，在法律上具有知识和真理的规范。}\BibleRef{罗 2:17}因为人们承认法律有这般力量，尽管那些行为不合法的人，假装好像生活在法律中。\Quote{找到智慧的人是有福的，获得明哲的人是有死的；因为从其口，}显然指智慧的口，\Quote{生出正义，它舌上载有法律和仁慈。}因为法律与福音都是同一个主的能力，祂是\Quote{天主的德能和智慧；}法律所产生的恐惧是仁慈的，且为救恩。\Quote{不可让慈善、信实和真理远离你，要将它们系在你的颈上。}\BibleRef{箴 3:3}同样，保禄也曾预言责备百姓不明白法律：\Quote{毁灭与痛苦布满他们的道路，和平的道路他们不认识。}\BibleRef{依 59:7}\Quote{他们眼中没有敬畏天主之情。}\Quote{他们自称为智者，反而成了愚拙。}\BibleRef{罗 1:22}\Quote{我们知道法律本是好的，只要人用得合法。}\BibleRef{弟前 1:8}\Quote{他们愿意做法律的教师，却不明白，}宗徒说，\Quote{自己说的什么，又证实的是什么。}\BibleRef{弟前 1:7}\Quote{命令的目的就是爱德，出于纯洁的心、良好的良心和无伪的信德。}\BibleRef{弟前 1:5}\par

% structure: p0243 source=h2 target=subsection reason=relative-heading-rank; base=h2

\subsection{第二十八章　梅瑟法律的四重划分}

梅瑟哲学相应地分为四部分：历史部分和特别称为立法部分，这两部分属伦理论述；第三部分是关于祭献的，属自然科学；第四部分，超越一切的，是神学领域，\Quote{直观，}柏拉图称之为真正的大奥秘。亚里士多德称此类为形而上学。按柏拉图的说法，辩证法正如他在\WorkTitle{政治家}中所言，是一门致力于发现事物解释的科学。智者应获得它，并非为了说或行我们在人中所见的任何事（如从事诡辩的辩证论者所做的那样），而是为了能够尽可能地说和行天主所喜悦的事。但真正的辩证法是混合了真理的哲学，通过考察事物、检验力量和权能，逐步上升至对最卓越本体的认识，并试图超越至宇宙的天主，它不宣称关于必死之事的知识，而是关于神圣和天上事物的科学；据此，在言语和行为上，甚至在人间事务中，也遵循适当的实践。因此，圣经渴望使我们成为这样的辩证论者，正确劝勉我们：\Quote{要做精明的兑换银钱者，}拒绝一些事物，但持守那善的。因为这真正的辩证法是一门分析思想对象、抽象地并单独地展示存在物的个别基质、或区分事物为类属的能力的科学，它下降至最特殊的属性，并呈现每个个别对象，使之单纯地如其所是地被沉思。\par

因此，唯独它引导至真正的智慧，这是神圣的能力，处理作为实体的实体的知识，它把握完美，并脱离一切情欲；这并非没有救主，祂借著圣言，撤去因邪恶教育而来的无知之黑暗，这黑暗曾遮蔽灵魂之眼，并赐下最好的恩赐。\par

\Quote{好使我们清楚认识天主或人。}\par

是祂真正地展示了我们如何认识自己。是祂向祂愿意的人启示宇宙之父，并启示到人性所能领悟的程度。\BibleQuote{除了父，没有人认识子；除了子和子所愿意启示的人，也没有人认识父。}\BibleRef{玛 11:27}因此，宗徒正确地说是藉著启示知道了这奥秘：\BibleQuote{正如我先前略略写过的，你们照着能明白的，就可以晓得我对基督奥秘的知识。}\BibleRef{弗 3:3}\Quote{照着你们能明白的，}他说，因为他知道有些人只接受了奶，尚未接受干粮，甚至也不是单纯的奶。法律的意义可以三重方式理解——或作为象征，或作为正确行为的规诫，或作为预言。但我深知，分辨并宣告这些事属于成熟之人。因为全部圣经的意义并非如俗语所说的一个单一的米科诺斯，而是那些追捕神圣教导联系的人，必须以逻辑能力的最大完美来接近它。\par

% structure: p0248 source=h2 target=subsection reason=relative-heading-rank; base=h2

\subsection{第二十九章 希腊人不过是孩童，与希伯来人相比}

因此，柏拉图中的埃及司祭最美丽地说：\Quote{啊，梭伦，梭伦，你们希腊人永远是孩童，在你们灵魂中没有一个从古时传统而来的古老意见。而且希腊人中没有一个是老人；}我想，他说的老人是指那些知道更遥远古代之事的人，即我们的经典；年轻人是指那些处理更近代的、被希腊人研究为题材的事——昨天和近日之事，仿佛它们是古老和古代的。因此他补充说：\Quote{也没有皓首穷经的研究；}因为我们以某种野蛮的方式，使用粗糙的隐喻。因此，那些心智正确的人，以毫无诡计的方式接近解释。关于希腊人，他说他们的意见与神话相差无几。因为幼稚的寓言或儿童中流传的故事都不适合聆听。他称神话本身为\Quote{孩童，}仿佛它们是希腊人中那些自以为是、见识短浅的智者之产物；他所谓的\Quote{皓首穷经的研究}是指蛮族所拥有的、源自最高古代的真理。与此相对，他用了\Quote{孩童寓言}一词，指责现代人尝试的神话性质，如同孩童，没有什么成年的东西，并断言他们的寓言和言论都是幼稚的。\par

因此，神圣地，那藉著启示向赫尔马说话的能力说：\Quote{异象和启示是为那些心怀二意的人，他们心里疑惑这些事是或不是。}\par

同样，来自学识资源的论证也加强、证实并确立论证性推理，因为人的心智像年轻人一样处于摇摆状态。\Quote{善的诫命，}按圣经所说，\BibleQuote{是灯，法律是光，因为训诲纠正生活的道路。}\BibleRef{箴 6:23}\Quote{法律是万物的君王，凡人和神明的君王，}品达说。然而，我理解这些话是指制定法律的那一位。而我认为赫西俄德的以下言论，尽管是诗人随意说出、并无理解，却是关于万有之天主的：\par

\Quote{因为克罗诺斯之子为人类制定了这法律：\newline{}鱼类、走兽和有翼的鸟类可以互相吞食，\newline{}因为她们没有正义的法则；\newline{}但正义（远为最好）祂赐给了人类。}\par

那么，无论所指的是与生俱来的本性律法，还是后来赐下的律法，它当然都出自天主；本性律法与训诲律法实为同一。同样，柏拉图在\WorkTitle{政治家}中说立法者是一位；在\WorkTitle{法律篇}中说懂得音乐者是一位；以此教导：圣言是一位，天主是一位。梅瑟明称主为盟约：\Quote{看，我与你的盟约}，此前曾告知他不可寻求写在文字中的盟约。因为这是万有的缔造者天主所立的盟约。天主被称为\OriginalTerm{天主}{\GreekText{Θεός}}，源出于\OriginalTerm{放置}{\GreekText{θέσις}}，即秩序与安排。在\WorkTitle{伯多禄的宣讲}中，你会发现主被称为律法和圣言。至此，让我们这第一本根据真哲学的诺斯替笔记杂集告一段落。\par

\SourceRecord{Translated by William Wilson.；Ante-Nicene Fathers；Vol. 2.；Edited by Alexander Roberts, James Donaldson, and A. Cleveland Coxe.；Buffalo, NY: Christian Literature Publishing Co.,；1885.；<http://www.newadvent.org/fathers/02101.htm>.}

% structure: p0001 source=h1 target=section reason=page-title

\section{\WorkTitle{杂集（卷二）}}

% structure: p0002 source=h2 target=subsection reason=relative-heading-rank; base=h2

\subsection{第一章 引言}

由于圣经称希腊人为外邦哲学的剽窃者，接下来需要考虑如何简要证明这一点。因为我们不仅要表明他们模仿并抄袭了我们经卷中所记载的奇迹；此外，我们还要证明，他们剽窃并篡改了他们所持守的主要教义（我们的著作更为古老，如我们已表明的），这些教义涉及信德、知识、学问、望德、爱德，以及悔改、节制和敬畏天主——这真是一大群真理的德行。\par

凡处理当前问题所必需的解释，都当包含在内，尤其是外邦哲学中隐晦的部分，即象征与谜语的领域；那些将古代教义置于系统的哲学研究之下的人，对此极为珍视，甚至认为对于认识真理而言绝对必要。此外，在我看来，接下来适宜为之辩护的，是那些招致希腊人攻击我们的教义；他们利用了少数圣经经文，但愿犹太人也愿倾听，能安静地从他已相信的转向他尚未相信的那一位。然后，那些真诚的哲学家将适当地被引入一场友善的揭露，既涉及他们的生活，也涉及他们发现新教义的过程；这并非为了我们报复诽谤者（对于那些学会了祝福诅咒我们的人，即使他们无端向我们发出亵渎之言，我们也远非如此），而是为了他们的悔改；如果这些精通智慧的人在某一程度上感到羞愧，因外邦人的论证而醒悟；从而能虽晚却清楚地看到，他们飘洋过海所追求的学识究竟是什么。他们所偷窃的将被指出，从而我们可以挫败他们的傲慢；而对于那些他们自夸通过探究而发现的东西，我们将提供驳斥。因此，我们还必须论述所谓的学科课程——它们的用处有多大；以及占星术、数学、魔术和巫术。因为所有希腊人都夸耀这些是最高的科学。\Quote{大胆责备的人乃是缔造和平者。}我们已屡次说过，我们既不实践也不学习用纯正的希腊语表达自己；因为那适合那些引诱群众离开真理的人。但真正的哲学证明将有益于听众的理智，而非他们的口舌。而且，在我看来，关心真理的人不应以技巧和刻意来组织他的语言，而只应尽力表达他的意思。因为那些对言辞吹毛求疵、为之耗费时间的人，反而错失了实质。园丁能从荆棘中无伤地摘下玫瑰，工匠能从牡蛎肉中找出珍珠，这乃是适合他们的技艺。据说，禽鸟在食物不丰、不得不以爪刨地艰难觅食时，肉质最为鲜美。如果有人以类比的方式，想要从众多希腊的貌似真实中找到真相，如同面具下的真实面孔，他必将付出艰辛才能觅得。因为在异象中显现给\Person{赫尔马}的那位能力者说：\Quote{凡启示给你的，必要启示给你。}\par

% structure: p0005 source=h2 target=subsection reason=relative-heading-rank; base=h2

\subsection{第二章 只有藉着信德才能认识天主}

\BibleQuote{不要因你的智慧而得意洋洋，}\BibleBook{}{箴言}说。\BibleQuote{在一切道路上都要承认她，好使她指引你的道路，使你的脚不致失足。}这些话语意在表明，我们的行为应当合乎理性，并进一步显示，我们应当从一切文化中选取并拥有有益的东西。智慧的道路多种多样，都直指真理之路。信德就是那条路。\BibleQuote{你的脚绝不致失足}是指有些似乎反对天主眷顾统一治理的人。因此，接着就说：\BibleQuote{不要自以为有智慧，}这是针对那些悖逆天主治理的不敬思想。\BibleQuote{但要敬畏天主，}唯有祂是全能的。由此自然得出，我们不可反对天主。下文尤其清楚地教导：\BibleQuote{敬畏天主就是远离邪恶；}因为经上说：\BibleQuote{并要远离一切邪恶。}这就是智慧的管教（\BibleQuote{因为主所爱的，祂必管教}），带来痛苦以产生悟性，并恢复和平与永生。因此，我们所追随的蛮族哲学，实际上是完美且真实的。\BibleBook{}{智慧篇}中也这样说：\BibleQuote{因为祂赐给我了万有的正确知识，使我知道世界的结构，}等等，直到\BibleQuote{以及根部的效力。}在这一切中，他包含了自然科学，即处理感官世界一切现象的知识。接着，他又在下文中提及理智的对象：\BibleQuote{我知道隐藏的和显露的事，因为创造万有的智慧教导了我。}简而言之，这就是我们哲学所宣称的目标；对这些知识的学习，若伴随着正当的行为，便通过创造万有的智慧，引向万有的统治者——一个难以把握和理解的存在，始终从追求者那里退避和隐去。然而，那在远处的祂——哦，不可言喻的奇迹！——却已非常临近。\BibleQuote{我是临近的天主，}主说。就本质而言，祂是遥远的；\BibleQuote{因为受生者如何能接近非受生者呢？}但祂却因那怀抱万物的能力而非常临近。\BibleQuote{人能在隐秘处行事，而我竟看不见他吗？}\BibleRef{耶 23:23}因为天主的能力永远临在，与我们接触，行使监察、施恩、教导之职。因此，梅瑟确信天主不能为人智所知，便说：\BibleQuote{求祢将祢的荣耀显给我看；}\BibleRef{出 33:18}并奋力进入天主发声的幽暗之中——即进入关于存有的无法接近且不可见的理念。因为天主不在黑暗中，也不在空间里，而是超越空间、时间以及事物的属性。因此，祂也绝不在任何部分之中，不论是作为包容者还是被包容者，不论是受限制还是被分割。\BibleQuote{你们要为我建造怎样的殿宇呢？}主说。\BibleRef{依 66:1}祂甚至没有为自己建造一座，因为祂不能受限制。即使天被称为祂的宝座，祂也并非因此受限制，而是以受造物为乐。\par

因此，显然，真理已向我们隐藏；如果这一点已通过一个例子表明，我们稍后将通过更多例子来确立。按照\Person{撒罗满}的话，那些既愿意学习又能够\BibleQuote{认识智慧和教诲，领悟智慧的话语，接受难解之言，并辨识真正的义德}的人是多么完全值得称赞啊，此外还有一种[义德]并非依照真理，而是由希腊法律和其余哲学家所教导的。经上说：\BibleQuote{并要指引判断}——不是指法庭上的判断，而是指我们必须保持我们内在的判断官能健全而无错误——\BibleQuote{使我能将机敏赐给愚钝者，将见识和领悟赐给少年人。}\BibleQuote{因为智慧人}既被说服遵守诫命，\BibleQuote{听了这些事，将因\OriginalTerm{知识}{\GreekText{γνῶσις}}而变得更智慧}；并且\BibleQuote{聪明人将获得治理权，并会理解比喻和隐晦的话，智慧人的言论和谜语。}\BibleRef{箴 1:2}因为那些受天主感动的人以及被他们争取过来的人所提出的并非虚假的言辞，也非大多数诡辩家用来纠缠青年、虚度光阴的陷阱。相反，那些拥有圣神的人\BibleQuote{探究天主的深奥事理}\BibleRef{格前 2:10}——即把握预言中的奥秘。禁止\BibleQuote{将圣物给狗}，只要它们仍是野兽。因为那些嫉妒、烦乱、行为仍然不信、无耻地吠叫以妨碍探究的人，绝不应浸入神圣而清澈的活水之流。\BibleQuote{不可让你的泉源之水溢出，要让你的水流在自己的街市上。}\BibleRef{箴 5:16}因为能理解所遇到之事的人并不多；即使学习后，他们也不知道，尽管他们自以为知道，正如可敬的\Person{赫拉克利特}所说的。他甚至不也似乎在责备那些不信的人吗？先知说：\BibleQuote{现在我的义人将因信德而生活。}\BibleRef{哈 2:4}另一位先知也说：\BibleQuote{除非你们相信，你们将不能理解。}\BibleRef{依 7:9}因为当不信在内心挣扎时，灵魂如何能接受如此高深主题的超然默观呢？但\OriginalTerm{信德}{\GreekText{πίστις}}——希腊人轻视它，认为它无用且野蛮——是一种自愿的预想，虔诚的赞同——按照神圣的宗徒，是\BibleQuote{所希望之事的实质，是未见之事的确证}。正因如此，尤其\BibleQuote{先人们获得了美誉。但没有信德，不可能中悦天主。}其他人将信德定义为对未见之物的联合赞同，正如对未知之事的确证是一种显然的赞同。如果它是选择，渴望某物，那么在此情况下渴望是理智的。既然选择是行动的开始，信德被发现是行动的开始，成为理性选择的基础，对于任何通过信德向自己展示先前证明的人而言。自愿追随有益之事，是领悟的首要原则。坚定不移的选择，则给予知识相当大的推动力。信德的操练直接成为知识，建立在稳固的基础上。因此，知识被哲学家们定义为一种不能被理性推翻的习惯。除了虔诚——唯独圣言是它的导师——之外，还有任何其他这样的真实状况吗？我想没有。\Person{泰奥弗拉斯托斯}说感觉是信德的根源。因为从感觉，基本原理解释扩展到我们内在的理性和理解。那么，以确定的判断相信神圣圣经的人，在赐予圣经的天主的声音中，接受了一种不可辩驳的证明。因此，信德不是由证明建立的。\BibleQuote{那些没有看见而相信的人，是有福的。}塞壬的歌声展现出超越人类的力量，迷惑了那些靠近的人，几乎违背他们的意愿，使他们接受所讲的话。\par

% structure: p0008 source=h2 target=subsection reason=relative-heading-rank; base=h2

\subsection{第三章　信德非本性的产物}

现在，\Person{巴西里德}的追随者视\OriginalTerm{信德}{\GreekText{πίστις}}为本性的，他们也将其归于\OriginalTerm{选择}{\GreekText{προαίρεσις}}，[认为]它通过理智理解而无需证明就能发现观念；而\Person{瓦伦廷}的追随者将信德归给我们这些简单人，但主张\OriginalTerm{知识}{\GreekText{γνῶσις}}在他们自己（那些通过本性得救的人）内萌生，借着优越之种的优势，说它与信德的差距如同属神者与属血肉者之差。此外，巴西里德的追随者说，信德和选择在每个层次上都是恰当的；并且由于超世俗的选择，世俗的信德伴随整个本性，而信德的恩赐符合各人的希望。那么，如果信德是一种本性优势，它就不再是自由选择的直接结果。\par

那不信的人，既不是他[自己不信]的作者，也不会得到应有的报应；而那相信的人也不是他[相信]的原因。如果我们正确反思，信与不信的全部特性和差异将不受任何赞扬或责备，因为其中附有来自全能者先行的本性必然性。如果我们像无生命之物一样被本性力量的提线牵引，那么意愿、不意愿以及作为两者先导的冲动，都只是多余。就我而言，我完全无法想象这样一种动物，其欲望由外部原因推动，却受必然性支配。那么，曾是不信者的人的\OriginalTerm{悔改}{\GreekText{μετάνοια}}——由此而来罪的赦免——还有什么位置呢？因此，\OriginalTerm{圣洗圣事}{\GreekText{βάπτισμα}}不再合理，荣福的\OriginalTerm{印记}{\GreekText{σφραγίς}}也不合理，圣子和圣父也不合理。但天主，我认为，结果成了本性力量的分配者，而这种分配不以自愿的\OriginalTerm{信德}{\GreekText{πίστις}}作为救恩的基础。\par

% structure: p0011 source=h2 target=subsection reason=relative-heading-rank; base=h2

\subsection{第四章　信德是一切知识的基础}

然而，我们因圣经得知，自主的选择与拒绝已由主赐予人，我们安息于信德无误的准则，显明甘愿的精神，因为我们选择了生命，并藉祂的声音相信天主。凡相信圣言的人，知道这事是真的；因为圣言就是真理。凡不信那说话者，就是不信天主。\par

\BibleQuote{因着信德，我们知道诸世界是藉天主的话造成的，以致看得见的，不是从显明之物造出来的。}宗徒说。\BibleQuote{因着信德，亚伯尔向天主献了比加音更丰富的祭献，藉此他得了证明，他是义人，天主为他的供品作证；因这信德，他虽死了，仍发言。}等等，直到\BibleQuote{暂且享受罪中之乐。}因此，信德使这些在律法之前的人成义，使他们成为天主应许的继承者。那么，我为何还要回顾并从我们手中的历史中引述更多信德的见证呢？\BibleQuote{若述说基德红、巴辣克、三松、依弗大、达味和撒慕尔及众先知，时间就不够了，}等等。\BibleRef{希 11:32}现在，既然真理存在于四件事中——感觉、理解、知识、意见——理智的领悟在自然顺序中是最先的；但就我们而言，并相对于我们自己，感觉是最先的，而知识本质由感觉和理解形成；证据对理解和感觉是共同的。感觉是知识的阶梯；而信德沿着感觉对象的路径前进，撇下意见，奔向不欺骗的事物，并安息于真理。\par

若有人说知识是基于推理的论证，愿他听：第一原理是不能论证的；因为它们既非藉艺术也非藉智慧被认识。后者涉及可变化的对象，而前者仅是实践性的，非理论性的。因此，人们认为宇宙的第一原因只能藉信德理解。因为一切知识都是可教授的；可教授的基于先已知的东西。但宇宙的第一原因先前未被希腊人认识；\Person{泰勒斯}因而认为水是第一原因；后继的自然哲学家亦然，直到\Person{阿那克萨戈拉}首先将心智置于物质事物之上。但他也未保存适合动力因的尊严，因为描述了一些愚蠢的漩涡，连同心智的惰性甚至愚昧。因此圣言也说：\BibleQuote{不要称地上的人为师傅。}\BibleRef{玛 23:9}因为知识是由论证产生的精神状态；但信德是一种恩宠，从不可论证的引导至普遍而单纯者，那既非物质、亦非物质、亦不在物质之下的。但那些不信的人，如所预料，将一切从上天下地、不可见区域拖到地上，\Quote{绝对用手抓住岩石和橡树，}按\Person{柏拉图}所说。因为，执着于这一切，他们断言唯有可触摸可把握的才存在，将物体与本质等同：他们自相矛盾，却非常虔诚地捍卫某些理智的、无形体的形式从天上某处、从不可见世界降临，强烈主张有真正的本质。\BibleQuote{看，我创造新事，}圣言说，\BibleQuote{眼所未见，耳所未闻，人心所未想到的。}\BibleRef{依 64:4}用新的眼、新的耳、新的心，凡可见可听的，都要藉主的门徒的信德与理解来领悟，他们属神地说话、听闻、行动。因为有真币，也有假的；假的同样欺骗非专业人士，却不骗兑换银钱者；他们因学会分辨和区分伪造与真的而知道。所以兑换银钱者只对非专业人士说钱币是假的。但原因，只有银行学徒和受过此训练的人学习。\par

\Person{亚里士多德}说，跟随知识的判断其实就是信德。因此，信德是优于知识的东西，并是知识的标准。猜测，只是一种微弱的假设，模仿信德；如同谄媚者模仿朋友，狼模仿狗。又如工匠看到通过学某些东西成为匠人，舵手通过受教于技艺而能操舵；他不认为仅仅渴望成为卓越和良善就足够，而是必须藉服从的操练来学习。但服从我们称之为导师的圣言，就是相信祂，丝毫不反对祂。因为我们怎能与天主敌对呢？因此，知识以信德为特征；信德藉一种神圣的相互对应，也变得以知识为特征。\par

同样，伊壁鸠鲁——他极其偏好快乐胜过真理——认为信德是心灵的一种预知；他将预知定义为对明显之物的把握，以及对事物之清晰理解；并断言，没有预知，无人能探究、怀疑、判断，甚至辩论。一个人若对他所追求之物没有预知，如何能认识他所研究的事物呢？再者，已经学习的人已将其预知转化为理解。若学习者并非没有包含所表达之内容的预知而学习，那人就有耳可听真理。那对能听之人的耳朵说话的人有福了；同样，那顺服之子也确实有福。现在，听就是理解。所以，如果信德无非是对所谈论之事物的一种心灵预知，而顺服、理解和说服也是如此称呼；那么没有人能在无信德中学到什么，因为没有人能在无预知中学到什么。因此，有更充分的证明表明先知所说的话完全真实：\Quote{除非你们相信，你们也不会理解。} 以弗所的赫拉克利特在释义这句神谕时说：\Quote{如果一个人不盼望，他就不找到那不被盼望的东西，因为它无法探究且难以接近。} 哲学家柏拉图在\WorkTitle{法律篇}中也说：\Quote{那愿意有福和快乐的人，必须从一开始就正直地领受真理，以便尽可能长久地真实生活；因为他是一个有信德的人。但不信者的人是那种自愿的虚假令其愉悦的人；而那种非自愿的虚假令其愉悦的人是无知的；这两者都不值得向往。因为缺乏友善的人，是无信和无知的。} 他在\WorkTitle{欧绪德谟篇}中岂不是隐晦地说这就是\Quote{王者的智慧}吗？在\WorkTitle{政治家篇}中他明确地说：\Quote{因此，真正君王的智慧是君王般的；无论是一位王子还是一位平民，拥有它的人，因这一行为，必定正当地被称为王族。} 现在，那些相信基督的人既是也是被称为\OriginalTerm{善人}{Chrestoi}，如同那些受真正君王照顾的人是君王般的。因为正如智者因智慧而智慧，守法者因法律而守法；同样，属于君王基督的人是君王，而属于基督的人是基督徒。接着，他继续清楚地说：\Quote{正确的将变为合法，法律按其本性是正当理性，它并非存在于文字或其他地方。} 而埃利亚的陌生者宣称那君王般的和政治家般的人是\Quote{一部活的法律}。这样的人就是那成全法律、\Quote{实行天父旨意}\BibleRef{玛 21:31}的人，记载在高柱上，为所有有视力之人树立为神圣美德的榜样。希腊人熟悉拉刻代蒙的埃福尔之杖，上面刻有法律于木头上。但我的法律，如前所述，既是君王般的又是活的；它就是正当理性。\Quote{法律，是万物的君王——必死和不死的，}正如波奥提亚的品达所唱。因为斯珀西波斯在第一卷反对克勒欧丰的书中，似乎像柏拉图那样写道：\Quote{因为如果王权是一件好事，而智者才是唯一的君王和统治者，那么法律，即正当理性，是好的；}事实正是如此。斯多葛派教导与此一致，将王权、司祭职、预言、立法、财富、真正美丽、高贵出身、自由只归于智者。但他们也承认，要找到这样的人极其困难。\par

% structure: p0017 source=h2 target=subsection reason=relative-heading-rank; base=h2

\subsection{第五章 他通过数个例子证明希腊人汲取自圣经作者}

因此，所有上述教义似乎都是从伟大的梅瑟传给希腊人的。万物都属于智者，这在以下话语中得到教导：\Quote{因为天主向我显示仁慈，我拥有万有。} 而智者是天主所爱的，天主在祂说\Quote{亚巴郎的天主、依撒格的天主、雅各伯的天主}\BibleRef{出 3:16}时暗示了这一点。因为第一位被发现明确地被称为\Quote{朋友}\BibleRef{雅 2:23}；第二位被显示接受了新名字，意为\Quote{看见天主的人}；而依撒格，天主亲自选择他作为一尊奉献的祭品，成为我们救恩计划的预像。\par

今在希腊人中，歌颂统治九年之久的米诺斯王，他是宙斯的密友；他们曾听见天主如何与梅瑟交谈，\BibleQuote{如同人与朋友交谈一样。} \BibleRef{出 33:11} 梅瑟乃智者、君王、立法者。然而我们的救主超越一切人性。祂如此可爱，以至于惟独祂被我们爱慕，我们的心专注于真正的美，因为\BibleQuote{祂是那真光。} \BibleRef{若 1:9} 祂被显明为君王，被天真的孩童、不信且无知的犹太人如此欢呼，并被先知所宣告。祂如此富有，以至轻视整个大地，以及其上其下的黄金，连同一切荣耀，当这些被仇敌献给祂时。何需多说祂是惟一的\OriginalTerm{大司祭}{High Priest}，惟独祂拥有敬拜天主的真知？祂是默基瑟德，\BibleQuote{和平之王} \BibleRef{希 7:2}，最适合作人类之首。祂也是一位立法者，因为祂藉先知的口赐下律法，明确命令并教导何事当做，何事不当做。谁的世系比祂更高贵？祂惟一的父就是天主。来吧，让我们引出柏拉图赞同这些教义。他在\WorkTitle{费德鲁斯}中称智者为富有，说道：\Quote{亲爱的潘，以及这里无论什么别的神，请赐我内在成为美；凡外在之物，愿其与内在合宜。我将视智者为富有。} 而那位雅典陌生人责备那些认为拥有许多财产的人便是富有者，如此说：\Quote{至富之人同时也成为善是不可能的——我指的是那些群众视为富有的人。他们称那些在少数人中拥有最值钱财物的人为富有；任何恶人都可能拥有这些。} \Quote{全部财富都属于信徒，}撒罗满说，\Quote{而不信者则一文不名。} 因此，圣经更应被相信，它说：\BibleQuote{骆驼穿过针眼，比一个富人} \BibleRef{玛 19:24} 过哲学的生活更容易。但另一方面，它祝福\BibleQuote{贫穷的人} \BibleRef{玛 5:3}；正如柏拉图理解的那样，他说：\Quote{贫穷不在于减少资源，而在于增加不知足；因为并非微薄的资金构成贫穷，而是不知足；善人摆脱了不知足，也将是富有的。} 在\WorkTitle{阿尔基比亚德}中，他称恶为奴役之物，德性为自由人的属性。\BibleQuote{从我这里取下重轭，背起轻省的轭，} \BibleRef{玛 11:28} 圣经说；同样，诗人们也称（恶）为奴役之轭。而\Quote{你们把自己卖给了你们的罪}这一表达，与上文一致：\BibleQuote{所以，凡犯罪的，就是奴隶；奴隶不会永远住在家里。但如果圣子使你们自由，你们就真是自由了，真理必使你们自由。} \BibleRef{若 8:32}\par

再者，柏拉图断言智者是美的，正如有人说某些人是正义的，即使他们也许其貌不扬。这样论及杰出的正直品格，若有人据此称他们为美，也不会被认为言过其实。而\Quote{祂的容貌比世人子都卑微}，这是预言的预告。\par

此外，柏拉图在\WorkTitle{政治家}中称智者为王。此论已在上文引述。\par

这些要点既已论证，让我们回到关于信德的话语上。柏拉图以最充分的论证证明，处处都需要信德，同时赞扬和平：\Quote{因为在叛乱中，若没有全备的德行，没有人会是可靠而健全的。有许多雇佣兵好斗，愿意在战争中死去；但除了极少数例外，他们大多是亡命之徒和恶棍，傲慢而无知。} 如果这些观察正确，\Quote{每一位稍有见识的立法者，在制定法律时，都会着眼于最大的德行。这德行就是忠诚，}我们无论在和平还是战争时期，以及一生其余时刻都需要它，因为它似乎包含了其他德行。\Quote{最好的事既非战争也非叛乱，因为这些的需要是可憎的。但彼此和平与善意才是最好的。} 从这些言论来看，显然最伟大的祈祷是获得和平，根据柏拉图。而信德是诸德最伟大的母亲。因此，撒罗满的话说得对：\BibleQuote{智慧是在信者口中。} \BibleRef{德 15:10} 因为色诺克拉底在其\WorkTitle{论理智}一书中也说\Quote{智慧是对第一原因和理智本质的知识。} 他认为理智是双重的：实践的与理论的，后者是人的智慧。因此智慧是理智，但并非所有理智都是智慧。并且已经证明，对宇宙第一原因的认识属于信德，而非推论。因为如果撒摩斯的毕达哥拉斯学派的追随者，拒绝讨论问题的证明，而认为单纯的\OriginalTerm{本人之言}{ipse dixit}是相信的基础，并且这一说法就足以确认他们所听到的；而奉献于真理默观的人，竟敢于不相信那位可信的导师、唯一救主天主，反而要求祂为自己所说的提供检验，那就奇怪了。但祂说：\Quote{有耳可听的，就应当听。} 这是谁呢？让厄庇卡摩斯来说：\par

\Quote{心在看，心在听；其余的都是聋子瞎子。}\par

赫拉克利特将一些人斥为不信者，说：\Quote{不知道如何听或说；}无疑有撒罗满的助益，他说：\BibleQuote{你若喜爱听，就必领悟；你若侧耳倾听，就必成为智慧。} \BibleRef{德 6:33}\par

% structure: p0025 source=h2 target=subsection reason=relative-heading-rank; base=h2

\subsection{第六章 信德的卓越与益处}

\BibleQuote{主，谁相信了我们的报告？} \BibleRef{依 53:1} 依撒意亚说。因为\BibleQuote{信德来自听觉，而听觉来自天主的圣言，}宗徒说。\BibleQuote{那么，他们怎能呼求他们所不信的那一位？他们怎能相信从未听过的那一位？没有宣讲者，他们怎能听到？除非被派遣，他们怎能宣讲？正如经上所载：那传播喜讯者的脚步是多么美丽！}你看他是如何将因听觉而来的信德和宗徒们的宣讲上溯至主的圣言和天主子。我们尚未明白主的圣言就是论证。\par

那么，正如打球不仅依赖于一个人熟练地掷球，还需要另一个人灵巧地接球，以使游戏按照规则进行；同样，教导也是可靠的，当听者的信德——可以说是一种自然的技艺——促进了学习过程。同样，大地也通过其生产力合作，适合于播种种子。因为没有学习者接受能力的运用，即使是最好的教导也无用，甚至预言也是如此，当听者缺乏受教的态度时。因为干枯的枝条，准备好接收火的力量，极易点燃；那著名的石头通过亲和力吸引钢铁，如同琥珀泪滴吸引木屑，而磁石使谷壳运动。被吸引的物质服从它们，受到一种微妙精神的推动，不是作为原因，而是作为协因。\par

既然存在着两种恶——一种以诡诈和隐秘为特征，另一种以暴力和驱使为主——圣言呼喊，召集众人；深知那些不会服从的人；然而，既然服从与否在于我们自己的力量，只要我们不能以无知为借口。祂发出公正的召唤，并按照各人的能力要求。因为有些人既能且愿，通过实践和净化达到了这一点；而另一些人，如果尚不能，却已有此意愿。意愿是灵魂的行为，但行动不能没有身体。行动也不仅凭其结果来评估；也根据各人自由选择的因素来判断——他是否轻易选择，是否悔改自己的罪，是否反省自己的失败并\OriginalTerm{悔改}{\GreekText{μετέγνω}}，即（\OriginalTerm{后来知道}{\GreekText{μετὰ} \GreekText{ταῦτα} \GreekText{ἔγνω}}）\Quote{后来知道}。因为悔改是迟来的知识，原初的无罪是知识。因此，悔改是信德的效果。因为除非人相信他所沉溺之事是罪，他不会放弃它；如果他不相信惩罚即将临于犯法者，而且救恩是那按诫命生活者的分，他不会改过。\par

望德也基于信德。因此，巴希理德的追随者将信德定义为灵魂对任何那些不以临在而影响感官之事物的赞同。而望德是对拥有善的期待。因此，期待必然奠基于信德。忠信的人是不折不扣地守护所托付给他的人；而我们被托付了关于天主的话语和神圣言语，以及命令连同施行训令。这就是忠信的仆人，受主称赞。当它说\BibleQuote{天主是忠信的}，表明祂在宣告任何事情时是值得相信的。如今，祂的圣言宣告；而\Quote{天主}本身是\Quote{忠信}。那么，如果相信就是假设，哲学家们怎能认为出自他们自身的东西是确实的呢？因为对先前论证的自愿赞同不是假设，而是对确实之物的赞同。谁比天主更有能力？不信是对抗祂的人软弱的否定假设：正如怀疑是一种难以容纳信德的状态。信德是对预理解的自愿假设和预期。期待是对未来的意见，而关于其他事物的期待是关于不确定性的意见。确信是对某物的强烈判断。因此我们相信祂，我们对祂有确信，以达至天主的光荣和救恩。我们信赖祂，祂是唯一的天主，我们所认识的，祂为我们高贵地应许了那些事物，并为此目的仁慈地创造和赐予了我们，这些事物绝不会落空。\par

善意为他人着想而愿他得福。因为祂一无所缺；从主涌出的恩惠与慈善终结于我们，乃是神圣的善意，并善意导致慈善。如果亚巴郎因相信而算为正义；如果我们也是亚巴郎的后裔，那么我们也必须因听而信。因为我们是以色列人，不是因神迹而被说服，而是因听。因此经上说：\BibleQuote{不生育的妇女，欢呼吧！未曾受产的，高喊吧！因为被遗弃的比有丈夫的儿女更多。} 见：\BibleRef{依 54:1} \Quote{你曾为人民的屏障而生活，你的儿女在他们祖先的帐幕中得蒙祝福。} 如果同样的居所藉预言应许给我们和祖先，那么两约的天主便是同一位。因此更加清楚地补充说：\Quote{你继承了以色列的盟约，} 这是对从外邦中蒙召的人说的，她们曾是不生育的，从前没有这位丈夫——即圣言——，也曾是孤寂的，没有新郎。\Quote{义人将因信德而生活，} 这是按照盟约和诫命；虽然它们名称和时间是两个，但按照[神圣]经济赐下——在能力上原是一个——那旧的和新的，都由同一个天主通过圣子分施。正如宗徒在\WorkTitle{罗马书}中也说：\Quote{因为天主的正义在它内由信德而显明，为引人进入信德，} 教导那从预言到福音由同一位主完成的唯一救恩。\Quote{这一嘱咐，}他说，\Quote{我交托给你，儿子弟茂德，照从前指着你的预言，使你凭这些预言打那美好的仗；持守信德和良好的良心；有些人丢弃良心，便在信德上翻了船，} 见：\BibleRef{弟前 1:18} 因为他们因不信玷污了来自天主的良心。因此，信德不再有理由被轻率地贬低为简单粗俗、人人都有。因为如果它只是人类的习惯，如同希腊人所认为的，它早已消失了。但如果它生长，且无处不有，那么我断言，信德，无论是建立在爱上还是恐惧上（如诋毁者所说），都是神圣的；它既不会被世俗的友谊撕裂，也不会因恐惧的存在而瓦解。因为爱因与信德友善结盟，使人为信徒；而信德作为爱的根基，转而引入行善；同样，恐惧，法律的训导者，被那些相信它的人当作恐惧。因为如果它的存在在其运作中显现，它既在将要行动和威胁时被相信，也在不行动和显现时被相信；而它被相信存在，并不自身产生信德，而是由信德检验并证明可靠。因此，从不信到信德——以及对望德和恐惧的信赖——这样的转变是神圣的。事实上，我们确信信德是迈向救恩的初次运动；之后恐惧、望德和悔改，与节制和忍耐一同前进，引领我们达到爱和知识。因此，宗徒巴尔纳伯正确地说：\Quote{我从领受的部分中尽力一点一点地寄给你们，使你们与信德一同也要有圆满的知识。恐惧和忍耐是你们信德的帮助者；我们的盟友是坚忍和节制。}他说，\Quote{这些关于主的事，若保持纯洁，便与它们一同喜乐，即智慧、领悟、智识、知识。} 前述的德行既是知识的元素，结果便是信德更为基础，对于诺斯替者如同呼吸对于现世生命一样必要。正如没有四大元素就不能生活，同样没有信德也不能获得知识。因此它是真理的支持。\par

% structure: p0031 source=h2 target=subsection reason=relative-heading-rank; base=h2

\subsection{第七章 恐惧的益处。对异议的回应。}

那些谴责恐惧的人，攻击法律；如果攻击法律，也显然攻击赐下法律的天主。因为这三个要素必然呈现在所讨论的主题中：统治者、他的管理以及被统治者。那么，如果按照假设，他们废除了法律，那么必然的结果是，每个被情欲引导、追求享乐的人，必定忽视正义、轻视神性，并毫无恐惧地放纵不敬和不义，从而脱离了真理。\par

他们说，恐惧是非理性的错乱和心灵的扰动。你说什么？既然命令是由圣言赐给我的，这个定义又怎能成立？但命令禁止，将恐惧悬在那些因训诫而受管教之人的头上。\par

恐惧并非不合理。因此它是合理的。怎能不呢？因为它这样劝诫：\Emph{不可杀人，不可奸淫，不可偷盗，不可作假见证。}但如果他们愿意在名称上争论，就让哲学家们把对法律的恐惧称为\OriginalTerm{谨慎的畏怖}{\GreekText{εὐλάβεια}}，即合理的\OriginalTerm{回避}{\GreekText{ἔκκλισις}}。法塞利的\Person{克里托劳斯}不无恰当地称他们为\Quote{名战者}（\OriginalTerm{名战者}{\GreekText{ὀνοματομάκοι}}）。那么，命令在名称改变后，便显得极其美好可爱。因此\OriginalTerm{谨慎的畏怖}{\GreekText{εὐλάβεια}}被证明是合理的，即回避有害之物；由此产生对先前罪过的悔改。\Quote{敬畏主是智慧的开端；良好的领悟属于遵行它的人。} 他称智慧为一种实践，即敬畏主为智慧铺路。但如果法律产生恐惧，那么对法律的知识就是智慧的开端；人没有法律就不是智慧的。因此拒绝法律的人是不智慧的；结果他们被视为\OriginalTerm{无神者}{\GreekText{ἄθεοι}}。而今指导是智慧的开端。\BibleQuote{但无神之人藐视智慧和指导，} 见：\BibleRef{箴 1:7} 经上说。\par

让我们看看法律宣布了何种恐怖。如果这些是介于德行与恶行之间的事物，诸如贫穷、疾病、卑微出身等，那么世俗法律就列出这些，并因此而受称赞。逍遥学派的人引入三种善的事物，并认为其对立面是恶的，这种观点与此相符。但赐予我们的法律命令我们逃避真正邪恶的事物——奸淫、不洁、鸡奸、无知、邪恶、灵魂之病、死亡（不是那使灵魂与身体分离的死亡，而是那使灵魂与真理分离的死亡）。因为这些才是真正的恶行，由它们生发出来的工作是可怕而可怖的。神谕说：\Quote{并非不义地，网罗为飞鸟张设；因为同谋流血的，为自己积蓄灾祸。}那么，为何某些异端大声引用宗徒的话说：\Quote{因为法律带来对罪的认识？}\BibleRef{罗 3:20}而说法律不是善的呢？我们对他们说，法律并没有引起罪，而是指出了罪。因为它命令当做之事，谴责不当做之事。善者的本分是教导有益之事，指明有害之事，劝勉践行其一，命令回避另一。宗徒（他们不理解他）说，因着法律，对罪的认识显明出来，并不是说罪从法律获得了存在。而法律，它训练人、作为\OriginalTerm{训导者}{\GreekText{παιδαγωγ}\GreekText{ος}}被赐下为引向基督，\BibleRef{迦 3:24}借着恐惧加以纠正，以管教的方式，为达到因基督而来的成全，又如何不是善的呢？经上说：\Quote{我不愿恶人死亡，只愿他悔改。}现在，诫命产生悔改，因为它阻止不当做的事，命令善行。依我之见，他藉无知意指死亡。\Quote{亲近主的人，充满鞭伤。}\BibleRef{友 8:27}显然，亲近知识的人，因渴慕真理，经历了危险、恐惧、烦恼、痛苦而获益。\Quote{受教的儿子成为智者，聪明的儿子免于火烧。聪明的儿子必接受诫命。}宗徒巴尔纳伯说：\Quote{祸哉，那些自视为聪明，在自己眼中看为精明的人}\BibleRef{依 5:21}，接着又说：\Quote{让我们成为属神的人，成为天主完美的殿宇；让我们尽己所能，实践对天主的敬畏，努力遵守他的命令，好使我们在他的审判中喜乐。}因此，\Quote{对天主的敬畏}被神性地称为智慧的开端。\BibleRef{箴 1:7}\par

% structure: p0036 source=h2 target=subsection reason=relative-heading-rank; base=h2

\subsection{第八章 巴西里德斯与华伦提努关于恐惧为事物原因之奇谈}

这里，巴西里德斯的追随者解释这句话说：\Quote{那位执政者，听见那被服侍的圣神的发言，对声音和景象都感到惊愕，因为听到了超出他所希望的好消息；他的惊愕被称为恐惧，而恐惧成为智慧的开端，这智慧区分各类，辨别、完善、并恢复。因为那超越万有者不仅分离并派遣了世界，而且连选民也一并分离和派遣了。}\par

华伦提努似乎在一封书信中也采纳了这种观点。因为他用以下的话写道：\Quote{正如天使们对这个受造物感到恐惧，因为他所说的话超越了他的构造，由于在他内的那一位曾无形地传递了至上本质的胚芽，并自由地说话；同样，在世上的人类族群里，人的作品成为他们制造者的恐惧——例如，偶像和雕像。并且所有人的手都塑造事物以承受天主的名号：因为亚当被塑造为人的名号，激发了与那先在的人相关的恐惧，仿佛那人在他内；于是他们惊惶失措，迅速毁坏了那作品。}\par

但是，既然只有一个第一因，如以后将表明的，这些人将被证明是喋喋不休和唧唧喳喳的发明者。然而，既然天主认为有益的是，人应从法律和先知那里通过主接受预备性的惩戒，因此对主的敬畏被称为智慧的开端，是由主通过梅瑟赐给悖逆和心硬的人的。因为理性不能说服的人，恐惧能驯服他们；那训导之圣言从一开始就预见到这一点，并用每种方法加以净化，使工具适合于敬虔。惊愕是因奇怪的显现或意料之外的表象（例如，一个消息）而产生的恐惧；而恐惧则是对已存在或将要发生之事的过度惊奇。于是他们不明白，他们用惊奇来描绘那至高的、被他们颂扬的天主，仿佛他受扰动、在惊奇之前处于无知。如果无知先于惊奇；如果这惊奇和恐惧，即智慧的开端，是对天主的敬畏，那么很可能，无知作为原因先于天主的智慧以及一切创造之工，不仅这些，甚至先于恢复，乃至先于选民本身。那么，这无知是对善的无知，还是对恶的无知呢？\par

嗯，如果（恐惧）是好事，为什么它会因惊愕而止息？那么，对他们而言，服务、宣讲和圣洗就是多余的。如果（恐惧）是坏事，那坏的东西怎么能成为最好的东西的原因？因为如果无知没有先存在，服务者就不会降临，也不会如他们所说，惊愕会抓住\Quote{王子}；他也不会从恐惧中获得智慧的开端，以便区分被选者和世俗者。如果对先存之人的恐惧使天使们共谋反对自己的作品，认为超然本质的一个无形种子寄居在那受造物中，或者由于无根据的怀疑激起嫉妒——这是不可信的——那么，天使们就成了受托付给他们的受造物的谋杀者，如同对待一个孩子一样，他们被证明有最严重的无知。或者假设他们因被卷入预知而受影响。但他们不会在他们的攻击中预知而共谋反对他们预知的事；他们也不会因预知而对自己作品感到恐惧，当他们感知到超然种子时。或者，最后，假设他们信赖自己的知识而胆敢（但这对他们也是不可能的），在得知普累若麻中的卓越之后，共谋反对人。此外，他们也动手攻击那按肖像而造的，其中也有原型，且与保持不变的知识一起，是不可毁灭的。\par

因此，针对这些人以及其他某些人，特别是马西翁派，圣经尽管他们不听，仍然呼喊说：\Quote{听从我的人必安然居住，得享安静，不怕灾祸。}\BibleRef{箴 1:33}\par

那么，他们将把法律看作什么呢？他们不会称它为邪恶，而是公义的；将善与公义区分开来。但是主命令我们畏惧邪恶时，并不是用一种恶替换另一种恶，而是用相反者废除相反者。现在恶是善的对立面，正如公义是不义的对立面。因此，如果那种由敬畏主所产生的无畏被称为善的开始，那么敬畏就是善事。而来自法律的敬畏不仅是公义的，而且是善的，因为它除去了恶。但是通过恐惧引入无畏，它并不是通过心理混乱产生冷淡，而是通过纪律产生情感的节制。因此，当我们听到：\Quote{尊敬主，要刚强；但不要惧怕除祂以外的，}\BibleRef{箴 7:2}我们理解这是指畏惧犯罪，并遵行天主所赐的命令，这就是来自天主的尊敬。因为对天主的敬畏是\OriginalTerm{畏惧}{\GreekText{Δέος}}。但如果恐惧是心灵的混乱——如有些人所认为的那样，恐惧是心灵的混乱——但并非所有的恐惧都是混乱。迷信确实是心灵的混乱，是对魔鬼的恐惧，这些魔鬼产生并受情感激动的支配。另一方面，因此，对天主的敬畏，祂不受混乱的影响，是没有混乱的。因为人害怕的不是天主，而是远离天主。而害怕这个——即陷入罪恶——的人，害怕并畏惧那些罪恶。而害怕堕落的人，希望自己免于败坏和混乱。\Quote{智慧人惧怕，就远离恶事；愚昧人却狂傲自恃，与它混杂，}经上说；又说：\Quote{敬畏主是力量的望德。}\par

% structure: p0043 source=h2 target=subsection reason=relative-heading-rank; base=h2

\subsection{第九章 基督徒诸德的关系}

因此，这样的敬畏导致悔改和望德。如今望德是对善事的期待，或是对不在场的善事的乐观期待；有利的环境是为了善的望德而被假定的，我们已学到这望德引领到爱。而爱被证明是对理性、生活和行为的赞同，或简言之，生活中的共融，或者是在享受同伴时，与正确理性结合的友谊和情感的强度。而\OriginalTerm{同伴}{\GreekText{ἑταῖρος}}是另一个自我；正如我们称那些由同一话语重生的人为弟兄。与爱相近的是款待，这是一门致力于款待陌生人的相宜艺术。而那些是陌生人，对他们而言，世界上的事是陌生的。因为我们视为世俗的，是那些寄望于地上和肉欲的人。\Quote{不要效法这个世界，只要心意更新而变化，叫你们察验何为天主的善良、纯全、可喜悦的旨意。}\BibleRef{罗 12:2}\par

因此，款待致力于对陌生人有益的事；\OriginalTerm{访客}{\GreekText{ἐπίξενοι}}是\OriginalTerm{陌生人}{\GreekText{ξένοι}}；朋友是访客；弟兄是朋友。\Quote{亲爱的兄弟，}\Person{荷马}说。\par

仁爱，也是自然情感，是对人的慈爱；自然情感是与生俱来的习性，在爱友人或家人中实践，它们都随爱而来。若我们内在真实的人是属灵的，仁爱便是向同享同一圣神之人的弟兄之爱。自然情感则是善意或情爱的保守；情爱是其完美彰显；被爱则是在行为中取悦，通过吸引与感召。人们因共识而达成一致，这共识就是对共同享受的美善的认识。因为\OriginalTerm{同心}{\GreekText{ὁμογνωμοσύνη}}是\OriginalTerm{见解的和谐}{\GreekText{συμφωνία} \GreekText{γνωμῶν}}。经上说：\Quote{你们的爱不可虚伪}，\Quote{要憎恶邪恶，依附良善，实行弟兄友爱}，一直到\Quote{若是可能，总要尽力与众人和睦}。然后说：\Quote{不可为恶所胜}，\Quote{反要以善胜恶}。同一位宗徒承认他为犹太人作证，\Quote{他们对天主有热情，但不按着真知识。因为他们不认识天主的正义，企图建立自己的正义，就不服从天主的正义。}\BibleRef{罗 10:2}因为他们不知道也不遵行法律的旨意，而自以为法律所希望的是什么。他们不相信法律是预言，只相信字面；他们出于恐惧而追随，并非出于心意和信德。\Quote{法律的目标是基督，使凡信的人都获得正义。}\BibleRef{罗 10:4}法律曾预言他，为所有相信的人。梅瑟曾对他们说：\Quote{我要用那不成子民的，激起你们的妒火；用愚昧的国民，惹起你们的怒气。}\BibleRef{罗 10:19}依撒意亚又说：\Quote{我没有寻找我的，让他们找到了我；没有访问我的，我向他们显现。}\BibleRef{依 45:2}——显然是在主来临之前；此后向以色列所说的预言现在正合宜：\Quote{我整天向悖逆和反抗的子民伸出双手。}你看见先知清楚宣告，召叫外邦人的原因就是子民的悖逆与反抗吗？然而天主的仁慈也在他们身上显明。因为宗徒说：\Quote{因着他们的过犯，救恩临到了外邦人，为激起他们的妒火}\BibleRef{罗 11:11}并使之愿意悔改。\WorkTitle{牧者}明确谈论了那些安眠的人，承认在基督显现之前，甚至在法律之前，就有一些义人出自外邦人和犹太人，因着在天主面前的蒙悦纳——如亚伯尔、诺厄，以及其他任何义人。他接着说：\Quote{宗徒和导师们，曾以权能和信德宣讲天主圣子的名，他们安眠后，向那些先前安眠的人宣讲。}然后他补充说：\Quote{他们给了他们宣讲的印记。于是他们同他们一起下到水中，又上来。但这些人活着下去，又活着上来。而那些先前安眠的人，死着下去，却活着上来。因此，他们藉着这些人得以复活，并知道了天主圣子的名。所以他们也同他们一起上来，被嵌入塔的建筑中，未经雕琢而一同被建造；他们在义德和极大的纯洁中安眠，只缺少这个印记。}宗徒说：\Quote{没有法律的外邦人，若顺着本性遵行法律的事，他们虽然没有法律，自己就是自己的法律。}\BibleRef{罗 2:14}\par

既然德行彼此相连，何必再说已经证明的事：信德藉着悔改而希望，畏惧藉着信德；忍耐和实践这些，连同学习，终结于爱，而爱因知识而完美。但必须注意，唯有天主应被视为本性智慧。因此，教导真理的智慧就是天主的德能，而知识的完美包含在其中。哲学家爱慕真理，如今被视为朋友，因其爱而成为真仆人。知识的开端是惊异于对象，正如柏拉图在其\Emph{泰阿泰德}中所说；玛窦在\Emph{传统}中劝告说：\Quote{惊异于你面前的事}，以此作为进一步知识的基础。同样在\WorkTitle{希伯来人福音}中写道：\Quote{惊异者将统治，统治后得以安息。}因此，无知的人，若仍无知，不可能进行哲学思考，因为尚未把握智慧的理念；因为哲学是努力把握真实存在及其相关学问。它并非使人养成良好行为习惯，而是知道如何运用、行动和劳作，以肖似天主。我是指救主天主，通过大祭司圣言服事宇宙的天主，藉着圣言，真正美善和正义得以看见。虔诚是与天主相应和相称的行为。\par

% structure: p0048 source=h2 target=subsection reason=relative-heading-rank; base=h2

\subsection{第十章　哲学家致力于何事}

因此，我们的哲学家投身于这三件事：第一，默观；第二，遵守诫命；第三，培育善人。三者和合，便形成诺斯替者。其中任何一项缺失，知识的要素便跛行。圣经因此神圣地说道：\Quote{主对梅瑟说：你告诉以色列子民说：我是上主，你们的天主。你们不要照你们曾居住的埃及地的风俗行事，也不要照我领你们去的迦南地的风俗行事，不要随从他们的习惯。你们应遵行我的判断，遵守我的法令，按此而行：我是上主，你们的天主。你们应遵守我的一切诫命，并遵行它们；遵行的人必因此生活。我是上主，你们的天主。}\BibleRef{肋 18:1} 那么，埃及和迦南地或是世界和欺骗的象征，或是苦难和困厄的象征；神谕指示我们须避免什么，又指示什么既是神圣而非世俗的，必须遵守。当它说：\Quote{遵行的人必因此生活，}\BibleRef{迦 3:12} 它就宣告了希伯来人自身的纠正，以及我们这些亲近之人的训练和进步；它同时宣告了他们的生命和我们的生命。因为\Quote{那些死在过犯中的，已与基督一同活过来，}\BibleRef{弗 2:5} 藉着我们的盟约。因为圣经多次重复\Quote{我是上主，你们的天主}这句话，以最能劝诫的方式使人羞愧，教导我们追随那颁赐诫命的天主，并温柔地劝勉我们寻求天主，尽力认识祂；这是最高的默观，探究最伟大的奥迹，真正的知识，成为理性无法驳倒的。这唯独是智慧的知识，行为的正直从未与它分离。\par

% structure: p0050 source=h2 target=subsection reason=relative-heading-rank; base=h2

\subsection{第十一章　藉着信德而来的知识最为确实}

然而，那些自以为有智慧的人的知识，无论是蛮族教派还是希腊哲学家，按照宗徒所说，\Quote{使人骄傲。}\BibleRef{格前 8:1} 但那按着真正哲学所传授的、科学性的证明——即知识——却基于信德。我们可以说，那是理性的过程，从已接受的前提，为所争论之事获得信德。既然信德是双重的——知识的信德和意见的信德——那么证明也可称为双重的，一种基于知识，另一种基于意见；同样，知识和预知也被指定为双重的，一种是本质上准确的，一种是有缺陷的。我们所拥有的证明，岂非唯一的真证明，因为它来自圣经、神圣著作和宗徒所说的\Quote{天主所教导的智慧}？学习也是对诫命的服从，即信德在于天主。信德是天主的能力，是真理的力量。例如，经上说：\Quote{如果你们有信德像芥菜子那样大，你们就能移山。}\BibleRef{玛 17:20} 又说：\Quote{照你们的信德给你们成就吧。}\BibleRef{玛 9:29} 一人因信德得医治，获得痊愈；死者因相信会复活之人的力量而复活。然而，基于意见的证明是人的证明，是修辞论证或辩证三段论的结果。因为我们所提到的最高证明，通过将圣经引述并向渴望学习者的灵魂开启，产生理智的信德；其结果是\OriginalTerm{真知}{gnosis}。因为如果为了证明争议点而被引用的东西，被视为真实的，因为是神圣和先知的，那么显然，据此推理得出的结论也将真实地推论出来；证明的合法结果将是知识。因此，当命令将天上神粮的纪念物存放在金罐中时，经上说：\Quote{敖默尔是十分之一厄法。}因为在我们内，这三个量度表示三种准则：对可感物的感觉，对名称和话语的言语，以及对理智对象的理智。因此，诺斯替者将避免言语、思想、感觉和行动中的错误，因为他已听见\Quote{凡注视妇女而有意贪恋她的，已经犯了奸淫，}\BibleRef{玛 5:28} 也默想\Quote{心里洁净的人是有福的，因为他们要看见天主，}并知道\Quote{不是入口的使人污秽，而是出口的才使人污秽。因为从心里发出来的是恶念。}\BibleRef{玛 5:8} 我认为，这就是按照天主的真正而公正的尺度，用以度量可度量之物——即涵盖人的\Quote{十}，上述三个量度概括了它。有身体和灵魂，五种官能，言语，生殖能力——理智或精神官能，或你爱怎么称呼都行。总之，我们必须超越其他一切，止于理智；正如在宇宙中，越过九重：第一重由四大元素组成，置于一处平等互换；然后七颗游星和一颗不游动的行星，第九重，到达完备之数，即超越九重之上的第十层，我们必须达到对天主的认识，简言之，在受造物之后渴慕造物主。因此，厄法中的十分之一和祭献的十分之一都呈献给天主；逾越节从第十天开始，是从一切烦恼和一切可感之物中的过渡。\par

因此，诺斯替者因信德而坚定；但自认为有智慧的人，受不稳定且摇摆不定的冲动驱使，不会触及真理之事。所以经上合理地写道：\Quote{加音从天主面前离去，住在伊甸东方的诺得地方。}如今\OriginalTerm{诺得}{Naid}被解释为\OriginalTerm{骚动}{commotion}，\OriginalTerm{伊甸}{Eden}为\OriginalTerm{愉悦}{delight}；信德、知识和平安都是愉悦，那悖逆者从此被逐出。但自认为有智慧的人，连神圣诫命的开端都不愿听从；反而如同自己的老师，脱缰，自愿投入翻腾的骚动中，从未受造的知识沉沦到必朽的和受造之物，在不同时间持有不同意见。\Quote{那些没有引导的人，如叶子般坠落。}\par

理性，这主导性原则，保持不动并引导灵魂，被称为其舵手。因为通往不变者的途径，是通过真正不变的方式获得的。因此，亚巴郎站在上主面前，走近说话。\BibleRef{创 18:22} 又对梅瑟说：\Quote{但你站在我这里。}\BibleRef{出 34:2} 而西蒙的追随者希望被同化，在品行上与他们所崇拜的站立的形象相似。所以，信德和真理的知识，使灵魂——它选择它们——始终如一、平稳均衡。因为对于虚伪之人，变幻、改变和背离是相宜的，正如对于诺斯替者，平静、安息和平安是相宜的。如同哲学因骄傲和自负而声名狼藉，同样，\OriginalTerm{知识（诺斯）}{gnosis}也因冒用同一名称的假知识而受损；关于此，宗徒写信说：\Quote{啊，弟茂德，你要保管所交托给你的，避开亵渎的空谈和假冒知识的反对；有些人自称有这知识，却偏离了信德。}\BibleRef{弟前 6:20}\par

异端者被这句话定了罪，就拒绝致弟茂德的书信。那么，如果主是真理、智慧和天主的能力——事实上祂正是——就表明真正的诺斯替者是认识祂、并藉着祂认识祂的父的人。因为他的看法与说\Quote{义人的嘴唇知道崇高的事}的那位相同。\par

% structure: p0055 source=h2 target=subsection reason=relative-heading-rank; base=h2

\subsection{第十二章 双重的信德}

信德如同时间都是双重的，我们会发现美德成对共居。因为记忆与过去相关，望德与未来相关。我们相信过去的事确实发生过，未来的事将要发生。另一方面，我们爱，因信德而确信过去如此，因望德而期待未来。因为在一切事上，爱伴随着认识唯一天主的诺斯替者。\Quote{看，祂所创造的一切都很好。}\BibleRef{创 1:31} 祂既认识又赞叹。孝爱增长寿命；\OriginalTerm{敬畏上主}{fear of the Lord}增添年日。因此，如同日子是生命历程的一部分，同样，敬畏是爱的开端，通过发展而成为信德，然后是爱。但我敬畏父，同时既怕又爱，并非像我怕并恨一头野兽（因为敬畏是双重的）。再者，因怕受罚而敬畏，我是在借着采取敬畏而爱自己。怕得罪父亲的人，爱自己。所以，那被发现拥有信德的人是有福的，因为他是由爱和敬畏组成的。信德是获救的能力，是通往永生的力量。再者，预言是预知；知识是预言的领悟，即对那揭示万物的主所预知之事的知识。\par

因此，对已预言之事的知识显示三重结果：或是很久以前已发生的事，或是现在存在的事，或是将要发生的事。那么，已完成或所盼望的两端都归属于信德；而当前的行为为确认两端提供了有说服力的论证。因为，如果预言只有一个，一部分正在实现，另一部分已经应验；由此，真理——无论是所盼望的还是已过的——都得到证实。因为它首先成为现在，然后对我们成为过去；所以，对已过之事的信念是对过去事件的把握，而对未来的望德是对未来事件的把握。\par

不仅柏拉图主义者，而且斯多葛主义者都说，同意在我们自己的力量之内。因此，一切意见、判断、推测和知识——我们借此生活并与人不断交往——都是一种同意；这无非就是信德。不信是信德的缺失，表明同意和信德拥有力量；因为不存在的东西不能被称为缺失。如果你考虑真理，你会发现人天生容易被误导而同意虚假之事，尽管他拥有相信真理所需的资源。\Quote{因此，环绕教会的那个德行，}正如\Person{牧者}所说，\Quote{就是信德，天主的选民借此得救；而那使人刚毅的是节制。随后是淳朴、知识、纯洁、端方、爱德，}所有这些皆是信德的女儿。再者，\Quote{信德领路，敬畏建造，爱德成全。}因此他说道，主应被敬畏以建树，而非魔鬼被畏惧以致毁灭。再者，主的作为——即祂的诫命——应被爱和遵行；但魔鬼的作为应被畏惧而不行。因为对天主的敬畏训练并恢复至爱；而对魔鬼作为的畏惧则与憎恨同住。同一经卷也说\Quote{悔改是高明的智慧。因为那悔改自己所行的人，不再像以前那样行事或说话。但他因自己的罪而折磨自己，却使灵魂获益。因此，罪的赦免与悔改不同；但两者都表明那在我们的能力之下。}\par

% structure: p0059 source=h2 target=subsection reason=relative-heading-rank; base=h2

\subsection{第十三章 论第一次和第二次悔改}

那么，凡得蒙赦罪的人，就不应再犯罪。因为，除了从罪中第一次且唯一的悔改（这是指最初在外教生活中的先前罪恶——我是指在无知中所犯的）之外，对那些蒙召的人，立刻提供了一种悔改，这悔改涤除灵魂宝座上的过犯，以使信德得以建立。主认识人心，预知未来，从起初、从太初就预见了人的善变和魔鬼的诡诈；他嫉妒人得蒙赦罪，就向天主的仆人提供罪的诱因，巧妙地作恶，使他们与自己一同堕落。因此，极其仁慈的主，对那些虽然已有信德却陷于任何过犯的人，赐予第二次悔改，为的是若有人在蒙召后受诱惑，被强力或欺骗所胜，仍可领受一种无需后悔的悔改。\Quote{因为如果我们明知真道之后，仍然故意犯罪，就再没有赎罪的祭了，只有恐惧地等待审判和烧灭众敌人的烈火。} \BibleRef{希 10:26} 然而，屡次连续不断的悔罪，与根本不信的人毫无区别，唯一不同的是他们意识到自己在犯罪。我不知道二者中哪个更糟：一个人是明知故犯，还是悔改后再次犯罪。因为在证明过程中，两方面都显明有罪：一方面是作恶者犯罪时自己谴责的罪，另一方面是预知即将行恶却仍动手去行的人的罪。那在忿怒和快乐中放纵自己的人，是在他明知之事中放纵自己；而那悔改之后又再次冲入快乐的人，与起初故意犯罪的人相差无几。因为一个人再次行他所悔改之事，且谴责自己所行，却仍心甘情愿地去作。\par

那么，凡从外邦人和旧生活中投奔信德的人，已一次获得了罪的赦免。但此后犯罪的人，虽因悔改而得宽恕，也当畏惧，因为他不再被洗涤以得罪赦。因为那\Quote{不是由血，也不是由肉欲的意愿}\BibleRef{若 1:13}所重生，而是由圣神所生的人，不仅必须抛弃他原先奉为神明的偶像，也必须抛弃他先前生活的工作；这重生在于悔改，不再犯同样的过失。因为屡次悔改和因缺乏训练而轻易改变，正是再次犯罪的实践。因此，为自己屡次过犯而频繁求赦，只是悔改的表象，而非悔改本身。但经上说：\Quote{正直人的正义修平道路。}\BibleRef{箴 11:5}又说：\Quote{无辜者的正义必使他的道路正直。}\BibleRef{箴 13:6}不但如此，\Quote{如父亲怜爱子女，主也怜爱敬畏祂的人。}大卫写道：\Quote{含泪播种的，必含笑收割；}即那些在忏悔中告明的人。\Quote{凡敬畏主的人都是蒙福的。}你在福音中看到相应的祝福。\Quote{不要害怕，}经上说，\Quote{当人富足，当他家产增加时：因为他死时什么都不能带走，他的光荣不会随他下去。}\Quote{但我因你的仁慈进入你的居所；我要在敬畏中向你的圣殿俯伏朝拜你；主，求你以你的正义领我进入。}欲念是灵魂朝向或离开某物的运动。激情是超过理性尺度的过度欲念，或是放纵不羁、违逆圣言的欲念。因此，激情是灵魂违反本性的扰乱，违逆理性。然而，反抗、分心和不服，一如服从，都在我们的能力之内。因此，自愿的行为受审判。但若有人逐一检视每一种激情，便会发现它们是不理性的冲动。\par

% structure: p0062 source=h2 target=subsection reason=relative-heading-rank; base=h2

\subsection{第十四章 事物如何可能是不自愿的}

凡非自愿之事，不属审判之列。然而非自愿有两种：因无知而作，及因不得已而作。因为对那被称为以非自愿方式犯罪的人，你将如何审判？或是不认识自己，如发狂的克莱奥梅尼与阿塔玛斯；或是不认识他所行之事，如埃斯库罗斯，在舞台上泄漏秘仪，经阿勒奥珀格斯审判，他因表明从未受过启蒙而被开释。或是不认识所行何事，如有人本欲击敌，却误杀家人；或是不认识所用之工具，如有人以装有圆头的标枪操练，因标枪脱头而致人死亡；或是不认识方式，如在运动场上击倒对手而死，因其所争取的是胜利而非死亡；或是不认识行事之目的，如医生施予解毒良药却致人死亡，因他给药并非为此，而是为救人。当时法律惩罚非自愿杀人者，例如非自愿遗精者，但与自愿者惩罚不同。然而若将情欲转向真理，此人亦当作为自愿行为受罚。因为事实上，不能克制生育之言的人应受惩罚，这是因为灵魂的理性冲动成了类似于喋喋不休的东西。\Quote{忠信的人选择将事隐藏在自己心中。}\BibleRef{箴 11:13}\newline{}因此，取决于选择之事乃审判的对象。\Quote{因为主鉴察人心肺腑。}\Quote{凡注视而动了淫念的}\BibleRef{玛 5:28}必受审判。故而经上说：\Quote{不可贪恋。}\BibleRef{出 20:17}\Quote{这百姓用嘴唇尊敬我，}经上说，\Quote{心却远离我。}因为天主垂顾心思本身，既然罗特的妻子只是自愿转向世俗之恶，主便将她留作无知觉的块状，使她成为盐柱，定住她不再前行，并非作为愚钝无用的雕像，而是为使具有属灵感知能力的人得着调味和盐渍。\par

% structure: p0064 source=h2 target=subsection reason=relative-heading-rank; base=h2

\subsection{第15章 论各种自愿行为及由此产生的罪}

凡自愿之事，或出于欲望，或出于选择，或出于意图。彼此相近者有罪、过失、罪行。例如，生活奢侈放荡是罪；误将朋友当作敌人而伤害是过失；掘墓或亵渎圣物是罪行。犯罪源于不能确定当作何事，或因无力行善；如同某人跌入沟壑，或因不知，或因身体软弱不能跃过。然而致力自我训练、服从诫命，乃在于我们自己的能力；若我们不愿与此事有关，完全放纵于情欲，我们便犯罪，甚而伤害自己的灵魂。因为著名的拉伊奥斯在悲剧中说：\par

\Quote{你所劝告我的这一切，我并非不知；\newline{}但我虽清醒，本性却强迫我；}\par

即他放纵于情欲。同样，美狄亚也在舞台上呼喊：\par

\Quote{我知道我将犯下何等罪恶，\newline{}但情欲强于我的决心。}\par

再者，连埃阿斯也并非缄默；当他要自杀时，喊道：\par

\Quote{没有什么痛苦像耻辱那样啃噬自由人的灵魂。\newline{}我如此受苦；灾难的深痕\newline{}时时从深处搅动我，\newline{}被愤怒的苦刺所激动。}\par

愤怒使这些人成为悲剧的主题，情欲则使成千上万的人——淮德拉、安提亚、厄里菲勒——\par

\Quote{她为她亲爱的丈夫取了宝贵的金子。}\par

另一部戏剧中，喜剧的特拉索尼德斯说道：\par

\Quote{一个下贱的女子使我成了她的奴隶。}\par

过失是违反理性之罪；自愿之罪是罪行（\OriginalTerm{不义}{\GreekText{ἀδικία}}）；罪行是自愿的邪恶。所以罪在我方面是自愿的。因此使徒说：\Quote{罪不得管辖你们，因为你们不在法律之下，而在恩宠之下。}\BibleRef{罗 4:7}他对相信的人说：\Quote{因祂的创伤我们得了痊愈。}\BibleRef{伯前 2:24}过失是他人对我的非自愿行为，而只有罪行（\OriginalTerm{不义}{\GreekText{ἀδικία}}）是自愿的，无论是我自己的行为还是他人的。圣咏作者暗示了这些罪的区分，当他称那些\OriginalTerm{过犯}{\GreekText{ἀνομίας}}被天主涂抹、罪恶（\OriginalTerm{罪}{\GreekText{ἁμαρτίας}}）被遮盖的人为有福时。另一些罪主不计算，其余的祂赦免。因为经上记着：\Quote{过犯得赦免、罪恶被遮盖的人，是有福的。主不计算罪愆、口中没有诡诈的人，是有福的。}这福分临到那些藉我主耶稣基督被天主拣选的人。因为\Quote{爱遮掩众多的罪。}\BibleRef{伯前 4:8}它们被那位\Quote{喜悦罪人悔改胜过死亡的主}\BibleRef{则 33:11}所涂抹。那些不因选择而生的罪不被计算；\Quote{因为凡动淫念的，已经心里犯奸淫了，}\BibleRef{玛 5:28}经上说。光照人的圣言赦免罪过：\Quote{那时，上主说，他们必寻找以色列的过犯，却寻不着；寻找犹大的罪恶，也找不到。}\BibleRef{耶 1:20}因为谁能像我？谁能站在我面前？\BibleRef{耶 49:19}你看那唯一的天主被宣称为善，按行为施行报应，并赦免罪恶。若望在他的长篇书信中也清楚教导罪的区分，他说：\Quote{若有人看见弟兄犯了不至于死的罪，就当祈求，并要将生命赐给他：就是给那些犯不至于死之罪的人。}因为\Quote{有至于死的罪，我不说当为那罪祈求。一切不义都是罪，但也有不至于死的罪。}\BibleRef{若一 5:16}\par

\Person{达味}，以及他之前的\Person{梅瑟}，以下列话语显示出对三条诫命的认识：\Quote{不随从恶人的计谋，便有福了；}正如鱼在黑暗中沉入深处；因为那些没有鳞的——\Person{梅瑟}禁止触摸的——是在海底觅食。\Quote{不站罪人的道路，}如同那些表面敬畏主、实际却犯罪的人，好比母猪，饥饿时嚎叫，饱足时却不认其主人。\Quote{不坐亵慢者的座位，}如同伺机捕食的飞鸟。\Person{梅瑟}吩咐不可吃母猪、鹰、隼、乌鸦，以及任何无鳞的鱼。\Person{巴尔纳伯}如是说。而我曾听一位精通此类事的人说：\Quote{恶人的计谋}指外邦人，\Quote{罪人的道路}指犹太人的信仰，而\Quote{亵慢者的座位}则指异端。另一个人更恰当地说：第一福分归于那些没有追随悖逆天主的邪恶思想的人；第二福分归于那些不留在宽大道路上的人，无论是受法律培育者，还是已悔改的外邦人。而\Quote{亵慢者的座位}则是剧场和审判台，或更确切地说，是对邪恶致命权力的顺从及其行为的共谋。\Quote{但他喜爱主的法律。} \Person{伯多禄}在他的\WorkTitle{宣讲}中称主为法律与圣言。立法者似乎以不同方式教导三种罪的解释——以哑鱼象征言语之罪，因为有时\Emph{沉默胜于言语，因为沉默有安全的赏报；}以贪婪食肉的飞鸟象征行为之罪。母猪喜爱污秽和粪土；我们不应有被\Quote{玷污}的\Quote{良心。} \BibleRef{格前 8:7}\par

因此，先知恰当地说：\Quote{恶人却不如此，乃像糠秕，被风吹散。因此，恶人在审判中必站立不住}（因为他们已被定罪，因为\Quote{不信的人已被定罪} \BibleRef{若 3:18}），\Quote{罪人在义人的会中也必如此，}因为他们已被定罪，以致不能与那些未跌倒的人联合。\Quote{因为主知道义人的道路；恶人的道路却必灭亡。}\par

同样，主清楚表明罪恶和过犯是在我们自己的能力之内，祂通过规定与疾病相应的治疗方式，显示祂愿意我们被牧者们纠正，正如在\BibleBook{厄则克耳}{先知书}中，我认为祂责备其中一些人不遵守诫命：\Quote{瘦弱的，你们没有养壮，}等等，直到\Quote{也没有一人寻回，或是挽回。} \BibleRef{则 34:4}\par

因为\Quote{当一个罪人得救时，在父面前有极大的喜乐，}主说。因此\Person{亚巴郎}极受赞美，因为\Quote{他照主对他所说的而行。}受此启发，希腊一位智者说出了格言：\Quote{跟随天主。} \Person{依撒意亚}说：\Quote{虔诚人筹划智慧的事。}因为筹划是寻求当前情境下的正确行动方式，而良好的筹划乃是筹划中的智慧。再者，天主在赦免\Person{加音}之后，不是很快引入了悔改的\Person{哈诺客}吗？这表明悔改的本性产生宽恕；但宽恕不在于豁免，而在于医治。同一例证是百姓在\Person{亚郎}之前铸造牛犊。受此启发，希腊一位智者说出了格言：\Quote{宽恕胜于惩罚；}同样，\Quote{作保，祸患就在眼前，}出自\Person{撒罗满}的话语：\Quote{我儿，你若为朋友作保，就会将手交给敌人；因为人的嘴唇是自己坚固的圈套，他被自己口中的言语缠住。} \BibleRef{箴 6:1}而\Quote{认识你自己}这句格言更神秘地取自：\Quote{你看见了你的兄弟，你看见了你的天主。}同样，\Quote{你要全心爱主你的天主，并爱邻人如己；}因为经上说：\Quote{全法律和先知都系于这两条诫命。}以下经文也与之相符：\Quote{我对你们说了这些事，好使我的喜乐得以圆满；这是我的诫命，叫你们彼此相爱，如同我爱了你们一样。} \BibleRef{若 15:11} \Quote{因为主仁慈而宽恕；主恩待一切。} \Quote{认识你自己}更明确且经常地由\Person{梅瑟}表达，当他命令说：\Quote{你要留心对自己。} \Quote{因此，施舍和信德的行为洁净罪恶。} \Quote{因敬畏主，人人远离恶事。} \BibleRef{箴 3:7} \Quote{敬畏主是教导和智慧。} \BibleRef{德 1:27}\par

% structure: p0080 source=h2 target=subsection reason=relative-heading-rank; base=h2

\subsection{第十六章 如何解释圣经中赋予天主人类情感的那些段落}

在此又出现吹毛求疵者，他们说喜乐与痛苦是灵魂的情欲：因为他们把喜乐定义为理性的兴奋和欢跃，即因善事而欢喜；把怜悯定义为对无辜受苦者的痛苦；并说此类情感是灵魂的情绪和情欲。但我们，看来，在这些事上仍未停止按血肉理解圣经；并从我们自身的情感出发，按我们的冲动来解释不动情的天主之旨意；且按我们所能听闻的方式，假定全能者也是如此，因而陷入不敬的谬误。因为天主的神性不能被如实地宣示；而是如同我们这些被肉体束缚的人所能听闻的那样，先知们对我们说话；主为了拯救而迁就人的软弱。既然天主的旨意是，那服从诫命并悔改自己罪恶的人应当得救恩，而我们因自己的救恩而喜乐，主便藉先知说话，将我们的喜乐归于自己；正如在福音中慈爱地说：\Quote{我饿了，你们给了我吃；我渴了，你们给了我喝。因为你们对我这些最小兄弟中的一个所做的，就是对我做的。}\BibleRef{玛 25:35} 因此，祂虽非亲自受滋养，却因祂愿意被滋养之人的滋养而受滋养；同样，祂毫无变化地，因悔改者如祂所愿地喜乐而喜乐。又因天主富于怜悯，祂是善的，藉法律和先知颁布命令，更亲近地藉祂圣子的显现，拯救并怜悯那些蒙受仁慈的人，如经上所说；并且按常理，大的怜悯小的；人不能大于人，因人生性为人；但天主在一切事上都大于人；因此，若大的怜悯小的，则唯有天主会怜悯我们。因为人被造是为藉义德而交流，并将从天主所领受的施予他人，这是由于他天生的仁爱与关系，以及他所服从的命令。但天主与我们并无天然关系，如异端作者们所主张的那样；无论是假定祂从无中造了我们，还是假定祂从物质中造了我们；因为前者根本不存在，后者则与天主全然相异——除非我们竟敢说我们是祂的一部分，与天主同体。我不知道一个认识天主的人如何能忍受听到此说，当他看到我们的生活，并见到我们陷于何等邪恶之中。因为这样就会得出，——说出此话便是亵渎——天主在\OriginalTerm{某些}{certain}部分中犯了罪，如果这些部分是整体的一部分且补足整体；若不能补足，则也不能是部分。但天主因祂自己的良善，本性上富于怜悯，眷顾我们，尽管我们既非祂的部分，也非祂本性的儿女。这正是天主良善的最大明证：尽管我们与祂的关系如此，且本性上完全疏远，祂仍然眷顾我们。因为动物对其后代的感情是天然的，而亲近心灵之间的友谊是亲密的结果。但天主的怜悯对我们来说是丰富的，我们与祂毫无关联；我说，无论在我们的本质或本性上，还是在我们本质的特殊运作上，而仅仅在于我们是祂旨意的化工。祂召唤那自愿、藉着纪律与教导接受真理知识的人，归于义子之名，这是一切进步中最大的。\BibleQuote{过犯捉住人；各人被自己罪恶的绳索捆绑。}\BibleRef{箴 5:22} 而天主无可指责。事实上，\BibleQuote{因虔敬而时刻战栗的人是有福的。}\BibleRef{箴 28:14}\par

% structure: p0082 source=h2 target=subsection reason=relative-heading-rank; base=h2

\subsection{第十七章  论各种知识}

因此，就如\OriginalTerm{知识}{\GreekText{ἐπιστήμη}}是一种理智状态，由此产生认识的活动，并成为理性不可反驳的把握；同样，无知是一种消退的印象，能被理性驱散。而能被理性推翻或精心构建的东西，都在我们的能力之内。与知识相近的是\Emph{经验}、\OriginalTerm{认知}{\GreekText{εἴδησις}}、\OriginalTerm{理解}{\GreekText{σύνεσις}}、知觉和\OriginalTerm{科学}{\GreekText{γνῶσις}}。\OriginalTerm{认知}{\GreekText{εἴδησις}}是按种属对普遍者的知识；经验是全面的知识，它探究每一事物的本性。\OriginalTerm{知觉}{\GreekText{νόησις}}是对理智对象的知识；\OriginalTerm{理解}{\GreekText{σύνεσις}}是对所比较之物的知识，或一种不可废止的比较，或比较判断与知识所涉及的对象的能力，无论是单一对象、每个对象，还是构成一个理性的一切。而\OriginalTerm{科学}{\GreekText{γνῶσις}}是对事物本身的知识，或与所发生之事相和谐的知识。真理是对真实之物的知识；真理的心智习惯是对真实之物的知识。知识由理性构成，不能被另一种理性推翻。我们不做某事，要么是因为不能，要么是不愿——或两者皆有。因此我们不飞，因为我们既不能也不愿；我们目前不游泳，例如，因为我们确实能，但选择不做；我们不像主，因为我们愿意，但不能做到：\Quote{因为没有门徒大过师傅的，能像师傅一样就够了：}\BibleRef{玛 10:24}并非在本质上（因为那因义子之名而有的，不可能在实体上等同于那因本性而有的）；而只在于我们已获得不朽，并真实地参与对实相的默观，并藉着属于祂的事物看见圣父。因此，意愿先于一切；因为理智能力是意志的仆役。经上说：\Quote{你愿意，就能做到。}在诺斯替者那里，意愿、判断和努力是一致的。因为若决断相同，则意见和判断也相同；因此他的言语、生活和行为都合乎规范。\BibleQuote{正直的心寻求知识，并听从它。}\BibleQuote{天主教导了我智慧，我认识了圣者的知识。}\BibleRef{箴 30:3}\par

因此，意愿先于一切；因为理智能力是意志的仆役。经上说：\Quote{你愿意，就能做到。}在诺斯替者那里，意愿、判断和努力是一致的。因为若决断相同，则意见和判断也相同；因此他的言语、生活和行为都合乎规范。\BibleQuote{正直的心寻求知识，并听从它。}\BibleQuote{天主教导了我智慧，我认识了圣者的知识。}\BibleRef{箴 30:3}\par

% structure: p0085 source=h2 target=subsection reason=relative-heading-rank; base=h2

\subsection{第十八章 梅瑟法律乃一切伦理之源泉，希腊人从中汲取其伦理思想。}

由此可见，所有其他德行，即梅瑟所描绘的勇敢、节制、智慧、正义、坚忍、忍耐、端庄、自制，以及此外的虔敬，都为希腊人提供了整个伦理领域的初步知识。\par

但人人皆知，教导敬拜与尊崇的虔敬，是最崇高且最古老的缘由；而法律本身彰显正义，兼教导智慧——藉由远离可感知的偶像，并引人归向宇宙的创造者与圣父。由此情感，犹如从泉源，一切智慧得以增长。\Quote{恶人的祭献为上主所憎恶，正直人的祈祷却蒙祂悦纳。} \BibleRef{箴 15:8} 因为\Quote{正义比祭献更蒙天主悦纳。}在依撒意亚书中我们也读到类似的话：\Quote{你们奉献许多牺牲，于我有何益处？上主说；}以及整段经文。\Quote{解除一切邪恶的束缚；因为这是悦纳上主的祭献：一颗痛悔的心寻求它的创造者。} \BibleRef{依 58:6} \Quote{欺诈的天秤为上主所憎恶，公平的砝码却蒙祂悦纳。} \BibleRef{箴 11:1} 于是\Person{毕达哥拉斯}劝诫\Quote{不要跨越天秤；}异端派别的宣称被称为欺诈的正义；\Quote{不义之人的舌头必被剪除，义人的口唇却涌流智慧。} \BibleRef{箴 10:31} \Quote{因为他们称智者和明达者为无用之人。}然而要逐一引证关于这些德行的经文未免冗长，因为整部圣经都称颂它们。既然他们定义勇敢为对可惧、不可惧以及中间之事的认识；节制为一种经由选择与回避而持守智慧判断的心境；与勇敢相连的是忍耐，即坚忍，是对可忍受与不可忍受之事的认识；而宽宏是超越环境的知识；与节制相连的是谨慎，即按理性回避；遵守诫命，即无过失地持守，可获致安全的生活。没有勇敢就没有忍耐，没有节制也没有自制。这些德行彼此相随；谁有德行的系列，谁就有救恩，即持守幸福状态。因此，在探讨这些德行时，我们应当正确地对它们全部加以探究；因为一个人若以灵知的方式拥有一种德行，由于它们彼此相随，便拥有全部。自制是一种不超越似乎合乎正确理性的品质。自制的人遏制违背正确理性的冲动，或克制自己不沉溺于违背正确理性的欲望。节制也离不开勇敢；因为从诫命生发出智慧——它遵行命令的天主——以及那效法天主性情的，即正义；凭此，在自制中，我们纯洁地致力于虔敬及由此而来的行为方式，与天主一致；尽可能如同我们这些本性会死的人一样，相似于主。这就是以智慧为正义与圣洁；因为神性一无所缺，一无所受；因此严格地说，它不能自制，因为它从未受到需要抑制的扰乱；而我们的本性可能受扰乱，需要自制，藉此修养自己，使需用变少，力图在品格上接近神性。因为善人作为不朽与有死本性之间的边界，需用甚少；由于身体及出生本身而有需求，但被理性自制教导只需求极少。\par

法律禁止男人\BibleQuote{穿戴妇女的服饰？}这有何理由？见：\BibleRef{申 22:5}。岂不是要我们成为男子汉，不要在个人外表与行为，也不要在思想与言语上显得柔弱吗？因为它愿意那献身于真理的人，在忍耐与坚韧的行为中，在生活、品行、言语和日夜的操练中，都显出刚毅；即使有必要以自己的流血作证。此外，经上说：\BibleQuote{若有人新建了房屋，尚未入住；或栽种了新葡萄园，尚未尝过果子；或聘定了处女，尚未迎娶她}——仁慈的法律命令这样的人免服兵役：首先出于军事理由，免得他们心系私欲，在战争中变得懈怠；因为只有不受情欲羁绊的人才能勇敢面对危险；也出于人道，因为考虑到战争的不确定性，法律认为一个人不应不能享受自己的劳动成果，而另一个人不劳而获，夺取劳动所得。法律似乎也指出灵魂的刚毅，它规定播种者应收获果实，建造者应居住，聘定者应迎娶：因为它为劳动者提供的不是虚妄的希望。按照有真知者的言论：\BibleQuote{善人已死或活着的希望不会消失，}智慧说；\BibleQuote{我爱爱我的人；寻求我的人必得平安，}见：\BibleRef{箴 8:17}，以及其它。那么怎样？米德扬人的妇女不是曾以自己的美貌，通过放纵情欲，在希伯来人攻打她们时，诱惑他们离开智慧而陷入不敬虔吗？因为她们引诱他们脱离庄重的生活，以美貌引他们陷入放荡的享乐，使他们沉迷于偶像祭献和异族妇女；同时被妇女和快乐所胜，他们背叛了天主，也背叛了法律。全体百姓差一点因妇女的诡计而落入敌人之手，直到危难之时，恐惧的警告将他们拉回。于是，幸存者勇敢地投入敬虔的争战，最终胜过了敌人。\BibleQuote{敬畏上主是智慧的开端，认识至圣者是明智；明悟法律是美好明悟的特征。}见：\BibleRef{箴 9:10}。因此，那些认为法律产生令人不安的恐惧的人，既不善于理解法律，也未曾真正领悟它；因为\BibleQuote{敬畏上主使人得生命，但错误的人必受知识所不见的苦痛。}见：\BibleRef{箴 19:23}。因此，\Person{巴尔纳伯}神秘地说：\Quote{愿统御万有的天主也赐给你们智慧、明悟、学问、认识祂的诫命和忍耐。因此，要受天主教导，寻求主向你们要求什么，好使祂在审判之日发现你们警醒等候这些事。}他称他们为\Quote{爱与和平的子女}，以真知的方式。\par

关于施舍与通财，虽可论述颇多，但只需指出：法律禁止一个\OriginalTerm{弟兄}{brother}收取利息；将\OriginalTerm{弟兄}{brother}不仅指同父母所生者，亦指同种族、同情感、同享\OriginalTerm{圣言}{Word}者；认为不应为金钱收取利息，而应敞开双手和心怀，馈赠给有需要的人。因为天主，这恩宠的创始者和施与者，以人类中最宝贵的事物——温和、温柔、宽宏、名誉、声望——作为合适的利息。你难道不认为这条诫命充满了仁爱吗？同样地，它说：\Quote{每天支付穷人的工钱}，教训人不可延迟支付应得的工钱；因为，我认为，穷人由于贫困，对于未来的热情便瘫痪了。此外，经上说：\Quote{债权人不得强行进入债务人的家中取走抵押品。} 而应当让前者请求取出抵押品，后者若有，便不可迟疑。\BibleRef{申 24:10} 在收获时，法律禁止主人占有从手中掉落的庄稼；同样，在收割时，法律命人留下一部分不收割；以此显著地训练拥有者去分享和宽宏大量，放弃自己的所有给那些有需要的人，从而为穷人提供生计。你看，法律如何同时宣告了天主的公义和良善，祂毫不吝啬地赐予众人食物。在采摘葡萄时，法律禁止摘葡萄者再回去摘取留下的葡萄，也不可拾取落下的葡萄；对橄榄采摘者也同样告诫。\BibleRef{肋 19:10} 此外，果实和牲畜的什一奉献教导我们对天主的虔诚，以及不可贪婪地攫取一切，而要将美善的恩赐分给邻人。我想，\OriginalTerm{司铎}{priest}们正是靠这些和初熟之果维持生计的。因此我们现在明白，法律教导我们虔诚、慷慨、公义和仁爱。因为它不是命令在第七年让土地休耕，并吩咐穷人无畏地享用那由神力生长的果实，大自然为了众人之益而耕种土地吗？\BibleRef{出 33:10} 那么，怎能说法律不是人道的、公义的导师呢？同样，在第五十年，法律命人实行与第七年相同的事；此外，还要将各人自己的田地归还原主，若他因某种原因在此期间放弃了它；它通过限制享乐的期限来限制那些贪图占有之人的欲望，并选择那些因长期贫困而付出代价的人不应遭受终生的惩罚。\Quote{然而，施舍和信德的行为是君王的护卫，祝福临于施舍者的头上；怜悯穷人的人有福了。} 因为他因爱人类的创造主，而向与自己相似的人显示了爱。上述细节关于安息和继承收回方面，还有更自然的解释，但目前不加以讨论。\par

如今爱以多种方式被理解，以温良、温柔、忍耐、慷慨、不嫉妒、无仇恨、忘怀伤害的形式而存在。在这一切中，爱是不可分割、不可区分的：它的本性就是分施。此外，经上说：\Quote{如果你看见你亲戚、朋友，或一般你所认识的任何人的牲畜在荒野中迷路，要将其牵回并归还；若主人远离，就把它养在你自己的牲畜中，直到他回来再归还。} 它教导一种自然的交流，即所找到的应视为寄存物，并且我们不得对敌人怀恨。\BibleQuote{主的诫命是生命的泉源} 确实 \BibleQuote{使人远离死亡的圈套}。\BibleRef{箴 14:27} 还有呢？它难道没有命令我们\Quote{不仅要爱陌生人如同朋友和亲属，而且要如同我们自己，无论在身体上还是在灵魂上吗？} 不仅如此，它还尊敬万民，不对那些作恶的人怀恨。因此，经上明确说：\BibleQuote{你不可憎恨埃及人，因为你在埃及地曾为外方人；} \BibleRef{申 23:7} 以\OriginalTerm{埃及人}{Egyptian}一词指代那个民族的人，或世界上的任何人。对敌人，即使他们列阵城下企图夺取城池，也不应视之为敌人，直到他们被传令官的声音召来和平。\BibleRef{申 20:10}\par

此外，法律禁止与女俘交合以羞辱她。它说：\BibleQuote{却要容她按自己的意愿哀哭三十天，并给她更换衣服，然后与她结合，作为你的合法妻子。} \BibleRef{申 21:10} 因为它认为此事不应出于放荡或像娼妓那样为钱而行，而应只为生育儿女。你看到人性与节制相结合吗？主人爱上了他的女俘，法律不允许他满足自己的私欲，而是通过规定一段时间来遏制他的情欲；此外，它剃去女俘的头发，以羞辱那可耻的爱情：因为若促使他结婚的是理智，那么即使她变得丑陋，他也会依恋她。然后，若有人在情欲之后，不再愿意与女俘同居，法律规定不许可出卖她，也不可再将她当作奴婢，而应使她自由，从服役中释放，免得在娶另一妻子之后，她遭受由嫉妒引起的无法忍受的痛苦。\par

还有什么呢？主吩咐要减轻并扶起敌人负重下的牲畜；间接教导我们不要因邻人的不幸而得意，也不要幸灾乐祸于敌人；以此来训练那些受教的人为敌人祈祷。因为祂不允许我们因邻人的福乐而忧伤，也不允许我们从邻人的灾祸中取乐。若你发现任何敌人的牲畜走失，你要超越分歧的刺激，领回并归还它。因为忘怀伤害会带来良善，良善又带来仇恨的消解。由此我们适于和睦，而这通向幸福。你若认为某人习于敌对，并发现他因私欲或愤怒而错误地执迷不悟，就要引导他走向良善。那么，这引人归向基督的律法岂不显得仁慈而温和？难道同一位置的天主，从起初到末了都以正义为特征，不是适当地运用各种方式以成就救恩吗？主说：\Quote{你们要仁慈，好蒙受仁慈；你们要宽恕，好被宽恕。你们怎样待人，也必怎样被对待；你们怎样给予，也必怎样被给予；你们怎样判断，也必怎样被判断；你们怎样施恩，也必怎样蒙恩；你们用什么量器量给人，也必用什么量器量给你们。} 此外，[律法]禁止那些为糊口而作奴仆的人被烙上耻辱，并将那些因借贷而沦为奴隶的人在第七年完全释放。它还禁止将求庇者交予惩罚。因此，那句神谕是千真万确的：\Quote{金银在炉中试炼，主却选择人心。仁慈的人宽容忍耐；每有忧虑的人便有智慧。因为忧虑将落在智慧人身上；他思虑求索，必寻得生命；寻求天主的，必找到知识与正义。凡正确寻求祂的，必得平安。} 在我看来，\Person{毕达哥拉斯}似乎是效法律法而对无理性的生物怀有温和。例如，他禁止立即使用羊群、山羊群和牛群中刚出生的幼崽；甚至不允许以祭献为借口，既为幼崽也为母兽着想；借此以低于人的生物，通过无理性的动物，训练人走向温和。他说：\Quote{因此，将幼崽交给母畜，至少头七天。} 因为若没有原因，一切都不会发生，而奶水如阵雨般赐予产子动物以养育后代，那将刚生下的幼崽从奶水的供应中夺走的人，便是侮辱自然。那么，让希腊人以及任何诋毁律法的人感到羞愧吧！因为律法在对无理性的生物时显出温和，而他们却遗弃人子；尽管律法早已在预言性的诫命中遏制了他们的残忍。因为若律法禁止将无理性生物的幼崽在哺乳前与母畜分离，那么，在人的情形中，它更是预先为残忍和野蛮的性情提供了疗方；这样，即使他们轻视自然，也不至于轻视教导。他们被允许以羊羔和羔羊饱食，也许还有理由将幼崽与母畜分离。但遗弃孩子又有什么理由呢？那人当初若不愿生儿育女，就不该结婚；当然更不应因放纵淫欲而成为杀子凶手。再者，人道的律法禁止在同一天内宰杀幼崽和母畜。由此，罗马人对于被判处死刑的孕妇，不允许在她分娩前执行刑罚。律法也明确禁止宰杀怀孕的动物，直到它们产仔，间接抑制了人对他人的不义倾向。这样，也将其仁慈延伸至无理性的生物；以便我们通过对异类生物的人道实践，在同族之间也能充分实践人道。此外，那些在动物产仔前踢其腹部，以便享用混合奶的肉食的人，使得为\Emph{胎儿}出生而创造的子宫成为其坟墓，而律法明确命令：\BibleQuote{但也不可用羊羔母的奶煮羊羔。} 因为活物的营养不可成为死物的佐料，生命的源泉不可助长肉身的消耗。同一条律法命令：\BibleQuote{不可笼住牛嘴，因为牛在打谷；工人应得他的粮食。} \BibleRef{申 25:4}\par

并且它禁止将牛和驴同轭耕地；\BibleRef{申 22:10} 这或许是指示动物之间的不一致；同时教导我们不可因种族差异而冤枉任何异族之人，并让他负轭，若除了种族差异外别无其他原因——这种差异根本不是原因，既非邪恶，也非邪恶的结果。在我看来，这寓言也似乎表明：\OriginalTerm{圣言}{Word}的耕作不该同等分配给洁净与不洁净的人、信徒与不信者；因为牛是洁净的，而驴被算作不洁净的动物。但慈祥的圣言，充满人性，教导我们不可为害而砍伐栽培的树木，也不可在收割前割下庄稼；并且根本不可毁坏栽培的果实——无论是土地的果实还是灵魂的果实：因为它不允许将敌人的土地变为荒芜。\par

此外，农夫在这些事情上从法律中获益。因为它命令新栽种的树木要连续三年受到培育，剪去多余的枝条，以免它们因负载过重而被压弯；又命令在树木周围耕种和挖掘，免得因养分被消耗而使得力量因匮乏而耗尽，这样树木就不会因长出吸枝而阻碍其生长。它也不允许从未成熟的树上摘取不完美的果实，而是在三年之后，即第四年，将初熟的果实奉献给天主，此时树木已成熟。\par

这种耕作方式可作为一种教导，教我们必须剪除罪的生长以及心灵周围无用的杂草，这些杂草环绕着生命的果实生长，直到信德的嫩芽完全并变得强壮。因为在第四年，对于正在牢固地接受教理讲授的人，需要时间，四枢德被奉献给天主，只有第三德已经与第四德结合，即主的位格。赞颂的祭献高于全燔祭：因为经上说：\BibleQuote{祂赐予力量以获得能力。}\BibleRef{申 8:18} 如果你的事务沐浴在繁荣的阳光下，要获得并保持力量，在知识中获取能力。因为这些例子表明，一切美好事物和恩赐都由天主供给；我们作为天主恩宠的仆役，应当播种天主的恩惠，使我们接近的人变得高贵和良善；这样，尽己所能，使节制者使他人节制，勇敢者使他人高尚，智慧者使他人明智，正义者使他人正义。\par

% structure: p0096 source=h2 target=subsection reason=relative-heading-rank; base=h2

\subsection{第十九章 真知者是天主的模仿者，尤其在施恩上}

那按天主的肖像和模样而有的人，才是真知者。他尽可能模仿天主，在促成相似的一切事情上无所欠缺，只要合宜。他实践自制和忍耐，正义地生活，驾驭私欲，尽己所能施舍所有，并以言行行善。经上说：\BibleQuote{在天国里，谁若实行并教导，将是最伟大的。}\BibleRef{玛 5:19} 他通过赐予类似的恩惠来模仿天主。因为天主的恩赐是为了共同的益处。经上说：\BibleQuote{凡胆敢妄为行事者，便是挑衅天主。}\BibleRef{户 15:30} 因为傲慢是灵魂的恶习，对于它以及其他罪，祂命令我们悔改；通过调整我们的生活，从混乱状态转向这三件事的改善——口、心、手。这些是标记：手代表行动，心代表意愿，口代表言语。因此，关于忏悔者，这句话说得很好：\BibleQuote{你今天选择了天主作你的天主；天主也选择了你作祂的子民。}\BibleRef{申 26:17} 因为那急于事奉自存者、身为恳求者的人，天主收纳他归于自己；纵使他仅一人之数，也与全体子民同受尊荣。因为他是子民的一部分，从原来的状态恢复后，成为其补充；全体由部分而得名。\par

然而，高贵本身表现在选择并实践最好的事物。因为对\Person{亚当}而言，他所拥有的那种高贵有何益处？没有凡人作他的父亲；因为他本人是所有出生之人的父亲。他轻易选择了卑劣的事物，跟随他的妻子，忽视了真实和良善；为此，他以不朽的生命换取了有死的生命，但并非永远。而\Person{诺厄}，他的起源与\Person{亚当}不同，因天主的眷顾而得救。因为他将自己奉献给了天主。至于\Person{亚巴郎}，他由三位妻子生了子女，并非为了满足快乐，而是如我所认为的，起初为了繁衍种族，结果只有一个儿子承继了他父亲的祝福，而其余的孩子脱离了家族；从他而生的双胞胎中，弟弟赢得了父亲的宠爱并获得了他的祈祷，成为继承人，而哥哥服侍了他。因为对一个恶人来说，最大的恩赐就是不能自主。\par

这个安排是预言性和预像性的。圣经清楚地表明一切属于智慧人，当它说：\BibleQuote{因为天主对我施了仁慈，我拥有了一切。}\BibleRef{创 33:11} 因为它教导我们只应愿望一件事，一切事物都藉此而有，所应许的也归给了配得的人。因此，那成为天国继承人的善人，圣经也藉着天主的智慧，将他登记为与古代正义之人的同乡，那些人或在法律之下，或在法律之前，都按法律生活，他们的行为成了我们的法律；此外，为教导智慧人是君王，圣经引入了另一个民族，对他说：\BibleQuote{你在我们中间是天主的君王；}\BibleRef{创 23:6} 那些被治理的人出于对善人德行的钦佩，自愿服从他。\par

现在哲学家柏拉图定义幸福的目的，说它是尽可能肖似天主；这或许是赞同法律的训诫（因为伟大的天性摆脱了情欲者，有时能击中真理的标记，正如毕达哥拉斯学派的斐洛在叙述梅瑟历史时所说的），或者因为受到当时某些神谕的教导，他始终渴求教导。因为法律说：\Quote{你们当跟随主你们的天主，遵守我的诫命。}\BibleRef{申 13:4} 因为法律称肖似为跟随；这样的跟随尽其所能地使人肖似。主说：\Quote{你们应当慈悲仁慈，如同你们的天上的圣父是仁慈的一样。}\BibleRef{路 6:36} 由此，斯多葛派也提出学说，认为符合自然的生活是目的，恰当地将天主的名称改为自然；因为自然也延伸到植物、种子、树木和石头。因此明确地说：\Quote{恶人不明白法律；但爱法律的人以法律为护卫。}\BibleRef{箴 28:4} \Quote{明智人的智慧明白自己的道路；愚人的愚昧却在错误中。}\BibleRef{箴 14:8} 因为预言说：\Quote{我将垂顾谁呢？岂不是那温良和善、敬畏我言语的人吗？}\par

我们受教导说，友谊有三种：其中第一种和最好的源于德行，因为基于理性的爱是坚固的；第二种和居间的是作为回报，是社交性的、慷慨的、对生活有益的；因为基于恩惠的友谊是相互的。第三种和最后的，我们断言是建立在亲密关系上的；另一些人则认为它是那种多变不定的、建立在快乐上的形式。毕达哥拉斯学派的希波达默斯似乎最出色地描述了友谊：\Quote{建立在认识神之上的友谊，建立在人类恩赐上的友谊，以及建立在动物快乐上的友谊。}有哲学家的友谊——人的友谊——和动物的友谊。因为天主的形象实际上是行善的人，他也在行善中获得善；如同舵手既拯救也被拯救。因此，当一个人获得所求时，他不对施予者说：\Quote{你施予得好}，而是说：\Quote{你领受得好}。所以施予者领受，领受者施予。\Quote{义人怜悯并施予仁慈。}\BibleRef{箴 21:26} \Quote{但温良的人将居住在地上，无辜的人将存留在其中；但悖逆者必从其中被铲除。}\BibleRef{箴 2:21} 荷马似乎也预言性地对信徒说：\Quote{施予你的朋友。}甚至必须帮助敌人，使他不再为敌。因为通过帮助，善意得以凝聚，敌意得以消解。\Quote{若已有乐意的心，按人所有的蒙悦纳，不是按人所没有的；因为这不是要别人轻松，你们受困苦，而是出于此刻的均等}等等。\Quote{他施舍，给予穷人；他的仁义永远长存}，圣经说。因为符合肖像和样式，不是指身体而言（因为可朽的与不朽的相似是错误的），而是指心智和理性，主适当地在这上面印上样式的印记，既是关于行善，也是关于行使治理。因为政府不是由身体的特质来管理，而是由心智的判断来管理。因为圣人的谋略，国家得以良好管理，家庭也是如此。\par

% structure: p0102 source=h2 target=subsection reason=relative-heading-rank; base=h2

\subsection{真正的诺斯替者操练忍耐与节制}

忍耐本身也迫使自己走向神圣的肖似，通过耐心收获其果实——不动情，如果记住关于阿纳尼亚斯的事迹；他属于一个群体，其中达尼尔先知充满神圣信德，也是其中之一。达尼尔住在巴比伦，如同罗特住在索多玛，以及亚巴郎——他后来成为天主的朋友——住在加色丁地。巴比伦王将达尼尔投入充满野兽的坑中；万王之王，忠实的主，将他平安地提上来。这样的忍耐，诺斯替者作为诺斯替者将会拥有。他在试炼中要祝福，如同高贵的约伯；如同约纳被鲸鱼吞下，他要祈祷，信德将使他复原，向尼尼微人预言；即使与狮子一同被关，他要驯服野兽；即使被投入火中，他会被露水沾湿，却不被烧毁。他要夜间作证；他要白天作证；以言语、以生活、以品行作证。与主同住，他要继续做他的密友，按照圣神共享同一炉灶；肉身纯洁，心地纯洁，言语圣化。据说：\Quote{世界被他钉在十字架上，他也被世界钉在十字架上。}他背负救主的十字架，将跟随主的足迹，如同天主，成为至圣所中的至圣。\par

因此，神圣的法律在记念一切美德的同时，特别训练人节制，将此作为德行的基础；并事先通过禁止我们食用那些天然肥腻的东西，如猪的品种——肉质丰满，来预备我们达到节制。因为这样的用途是为饕客而设的。据说有一位哲学家在给出\OriginalTerm{猪}{\GreekText{ὗς}}的词源时，说它是\Quote{祭杀}（\OriginalTerm{祭杀}{\GreekText{θύς}}），因为它只适合宰杀（\OriginalTerm{牺牲}{\GreekText{θύσιν}}）和杀害；因为这种动物的生命没有其他目的，只是为了在肉中膨胀。同样，法律压制我们的欲望，禁止食用没有鳍和鳞的鱼；因为这些鱼在肉质和肥腻上超过其他鱼。据我看来，基于此，秘仪不仅禁止触摸某些动物，而且还在祭杀的牺牲中移走某些部分，其原因只有入教者知道。那么，如果我们要控制肚腹和肚腹以下的部分，显然我们早已从主那里听到，我们要通过法律来抑制情欲。\par

如果我們真誠地譴責那作為情欲燃料的東西——我指的是享樂，那麼這將完全成就。他們說，享樂的理念是一種溫和而平淡的興奮，伴隨著某種感覺。他們說，\Person{墨涅拉俄斯}被此俘獲，在攻陷\Place{特洛伊}後，衝去殺死\Person{海倫}，視她為如此災難的原因，但卻未能得逞，因為被她的美貌所征服，使他想到享樂。因此，悲劇詩人們嘲笑地對他發出侮辱性的呼喊：——\par

\Quote{但是，當你看見她的胸懷時，你的劍\newline{}丟棄，並以一個親吻擁抱了那背叛者、\newline{}永遠美麗的可憐蟲。}\par

又說：——\par

\Quote{難道劍就這樣被美貌挫鈍了？}\par

我同意\Person{安提斯泰尼}他所說的，\Quote{如果能抓住\Person{阿佛洛狄忒}，我會射殺她；因為她毀滅了我們許多美麗而良善的女人。}他又說，\Quote{愛情是自然的缺點，而那些落入其權力之下的可憐蟲稱這種疾病為神明。}因為這些話表明，愚蠢的人因對享樂的無知而被征服，而我們不應允許享樂進入，即使它被稱為一個神明，也就是說，即使它被天主賜予為了生育的必要。\Person{色諾芬}明確稱享樂為缺點，說：\Quote{可憐蟲，你知道什麼好處，或者你有什麼可敬的目標？它甚至不等待對甜食的慾望，在飢餓前就吃，在口渴前就喝；為了你能愉快地吃，尋找好的廚師；為了你能愉快地喝，購買昂貴的酒；在夏天奔跑尋找雪；為了你能愉快地睡覺，不僅提供軟床，還提供床架。}因此，正如\Person{阿里斯頓}所說，\Quote{對抗享樂、痛苦、恐懼和情欲的整個四分音階，需要大量的練習和搏鬥。}\par

\Quote{因為是這些，正是這些穿過我們的內臟，\newline{}擾亂人們的心靈。}\par

因為即使是被認為嚴肅的人，享樂也使他們的心變得蠟一般，據\Person{柏拉圖}說；因為\Quote{每一種快樂和痛苦都把人的靈魂釘在身體上}，而那個人並沒有將自己從激情中分離並釘死自己。\Quote{那喪失自己生命的}，主說，\Quote{必要拯救它}；或為了主的緣故冒險捨棄它，如同祂為我們所做的，或使它脫離與其習慣生命的聯繫。因為你若願意解開、抽離、分開（這就是十字架的意義）你的靈魂脫離此生命中的歡樂和享樂，你將擁有它，在盼望的希望中找到並安息。而這將是死亡的操練，如果我們滿足於那些僅按本性衡量的願望，這些願望不越過那些合乎本性的界限——藉著過度或違反本性——在其中犯罪的潛能出現。\BibleQuote{因此我們必須穿上天主的全副武裝，使我們能抵擋魔鬼的詭計；因為我們的戰爭武器不是屬血肉的，而是藉著天主有能力摧毀營壘，推翻推理和一切自高自大阻礙對天主的認識的事物，並將每一個思想俘獲，使之服從基督}，\BibleRef{弗 6:11}神聖的宗徒說。\par

因為在被運動的物體中，一些由衝動和表象而運動，如動物；另一些由位移而運動，如無生物。而在無生命之物中，植物，他們說，是為了生長而由位移運動，如果我們承認植物是無生命的。那麼，石頭屬於靜止狀態。植物有本性；無理性的動物擁有衝動和知覺，以及前述兩個特徵。但推理官能，作為人類靈魂所特有的，不應像無理性的動物那樣被驅動，而應分辨表象，不被它們帶走。我們所說的那些能力，在易受影響的靈魂面前展示美景、榮譽、通姦、享樂以及諸如此類誘人的幻象；如同驅趕牲畜的人向它們伸出樹枝。然後，欺騙了那些無法分辨真與假的享樂、以及消逝俗豔與神聖之美的人後，他們將他們引向奴役。每種欺騙，藉著不斷向精神施加壓力，在它上面印下自己的形像；而靈魂不知不覺地背上激情的形像，這源於誘餌和我們的同意。\par

巴西里德的追随者习惯称情欲为\OriginalTerm{附加灵}{appendages}，他们说这些本质上是一些依附于理性灵魂的灵，源于某种原始的扰动与混乱；又有其他混杂而异质的灵性本性附着于它们，如狼、猿、狮、山羊的本性，这些动物的特性显现在灵魂周围，他们说，使灵魂的欲望与动物的形象相似。因为他们模仿他们所承载特性的那些动物的行为。这些附加灵不仅与无理性动物的冲动和感知相关，还影响植物运动和美感，因为植物特性也附着于它们。它们还具有特定状态（如钢的硬度）的特性。但我们将在此后讨论灵魂时反驳这一教义。目前只需指出：根据巴西里德，人如同木马（按诗意的神话），在一个身体内包含了如此众多不同的灵。因此，巴西里德之子依西多禄在其著作\WorkTitle{关于附于我们的灵魂}中，虽同意这一教义，却似自责地写道：\Quote{若我说服任何人灵魂不可分割，恶人的情欲由附加灵的暴力引起，那么恶人便有极大借口说：\InnerQuote{我被强迫，我被带走，我非自愿地做了，我是不情愿的}——尽管他本人引领了对恶欲的渴望，并未抵抗附加灵的袭击。但我们必须在理性部分取得优势，表明自己是我们内低级受造的支配者。}因为他也假设我们内有两个灵魂，如同毕达哥拉斯学派，我们后将提及。\newline{}\par

瓦伦提努也在致某些人的信中，就附加灵确切写道：\Quote{有一位善者，其临在即是通过圣子所彰显的启示，且唯独藉祂，心才能因驱逐心中每一恶灵而变得纯洁；因为住在其内的众多灵不使心纯洁；它们各自施行己事，屡以不当的情欲凌辱心。心似乎被对待如同客栈：后者有洞有沟，时常充满污秽；人们肮脏地住在其内，并不顾惜这场所为他人之物。心亦然，只要不被顾惜，便是不洁且众多恶魔的居所。但当唯一善的圣父来访时，它便被圣化，光照闪耀。拥有如此之心的人有福了，以致\Quote{他会看见天主。}}\par

那么，让他们告诉我们，为何这样的灵魂从起初就不被顾惜？要么它不配（于是顾惜它如同出于悔改），要么它是得救的本性（如他所说），而这本性必然从起初就因亲近而被顾惜，不给不洁之灵任何入口，除非被强迫并发现软弱。因为若他承认因悔改而选择了更好，他将不情愿地说出那正是我们所持真理教导的：即救恩来自因服从而有的转变，而非出于本性。正如从大地和沼泽升起的雾气聚成云雾和阴云，肉身情欲的蒸汽将邪恶状态带给灵魂，在灵魂前散布快乐的偶像。因此，它们将黑暗笼罩在理智的光明之上，灵吸引从情欲升起的雾气，并因执迷于快乐而加厚情欲的云团。金子不是从土中整块取出，而是通过冶炼得以纯净；然后，当变得纯净时，它被称为金子，而土被净化。因为\BibleQuote{祈求，就得给你们，}\BibleRef{玛 7:7}这是对那些能自行选择最佳者说的。至于我们如何说魔鬼的权势和不洁之灵将罪播种到罪人的灵魂中，无需我多言，只引用宗徒的巴尔纳伯（他是七十门徒之一，保禄的同工）为证，他如此说：\Quote{在我们相信天主之前，我们心的居所是不稳定的，确实是座人手所造的殿。因为它充满了偶像崇拜，是恶魔的房屋，因为做了相反天主的事。}\par

他因此说罪人行使与恶魔相称的活动，但并未说灵本身住在不信者的灵魂中。因此，他又加上：\Quote{要使主的殿荣耀地被建造。学习，既然获得了罪的赦免；并将我们的希望寄于那名，让我们成为新的，从起初被重新创造。}因为他所说的不是恶魔被从我们内驱逐，而是我们相信之前所犯的像恶魔的罪被赦免。因此他正确地以如下内容对比：\Quote{所以天主确实住在我们的家中。祂住在我们内。如何？祂信德之言、祂应许的召唤、祂法令的智慧、祂训导的诫命（住在我们内）。}\par

我知道我遇到了一种异端；其首领常说以快乐战胜快乐，这位可敬的诺斯替者以假装战斗的方式推进快乐，因为他自称是诺斯替者；他说，一个未尝试快乐的人禁戒快乐不算大事，但一个涉足其中而不被胜过的才算大事；因此他藉快乐训练自己。这可怜人不知他正被快乐的艺术所欺骗。显然，昔勒尼的亚里斯提卜附和了这一观点——那位夸耀真理的诡辩家的观点。因此，当被指责与科林斯娼妓持续同居时，他说：\Quote{我拥有莱斯，但我不被她拥有。}\par

此类人也包括那些自诩追随尼苛劳（Nicolaus）、曲解其箴言者。箴言说：\Quote{肉体必须被滥用。}但这位可敬的人表明，必须抑制快乐与私欲，通过这样的训练消磨肉体的冲动与倾向。而他们却像山羊一样纵情享乐，仿佛在侮辱身体，过着放纵的生活；不知道肉体本会朽坏，终将消解；而他们的灵魂却埋葬在恶习的泥沼中；他们追随的实是快乐本身的教导，而非宗徒之人的教导。因为他们与萨尔达纳帕卢斯（Sardanapalus）有何区别？其生活见于铭文中：\par

\Quote{我所吃的，我所纵欲享受的，\newline{}及我在爱情中感受的快乐，皆为我所有。\newline{}但那些众多的幸福之物都已遗落，\newline{}因为我亦为尘土，曾统治伟大的尼努斯（Ninus）。}\par

快乐的感觉绝非必然，而只是某些自然需求——饥饿、口渴、寒冷、婚姻——的伴随物。如果可能，无需快乐即可饮水、进食或生育子女，那么快乐的其他需要便无从证明。因为快乐既非功能，亦非状态，也非我们的一部分；它作为辅助物被引入生命，如同人们所说，盐用于调味食物。但当它挣脱约束、主宰家庭时，首先产生\OriginalTerm{私欲偏情}{concupiscence}，即一种非理性的朝向满足之物的倾向与冲动；这导致厄璧鸠鲁（Epicurus）将快乐定为哲学家的目标。因此，他将身体的健康状态及对它的确定希望神化。因为除了放纵于自我享乐者之饕餮与过剩，奢侈又是什么？狄奥革涅斯（Diogenes）在一部悲剧中意味深长地写道：\par

\Quote{谁心中依恋于\newline{}柔弱污秽的奢侈之乐，\newline{}便不愿承受丝毫劳苦。}\par

接下来的话语，虽措辞粗鄙，却恰如其分地描绘了那些纵欲者。\par

因此，在我看来，神圣的法律必然以恐惧来警告，使哲学家通过警醒与谨慎，获得并保持无焦虑的状态，在一切事上持续无过犯、无罪愆。因为平安与自由，若非通过与我们情欲不断的、不屈的斗争，便无法获得。因为那些顽强的、如奥林匹克竞赛般的对手，比黄蜂更敏锐，可谓如此；尤其快乐，不仅白日，连夜间也在梦中以巫术般诱人、诡计多端地叮咬。那么，希腊人自己教导快乐是恐惧的奴隶，又怎能正确贬斥法律？索克拉特斯（Socrates）因此嘱咐\Quote{人们要提防在不饥饿时进食、不口渴时饮水的诱惑，以及美人的秋波与亲吻，因为这些比蝎子和蜘蛛的毒液更致命。}而安提斯泰奈斯（Antisthenes）宁愿\Emph{疯癫，也不愿快乐}。特拜的克拉特斯（Crates）说：\par

\Quote{驾驭这些，因灵魂的性情而欢欣，\newline{}既不被黄金征服，也不被憔悴的爱情征服，\newline{}他们也不再是纵欲者的仆从。}\par

最后他推论道：\par

\Quote{这些人，不被奴役的快乐所束缚或屈服，\newline{}热爱不朽的王国与自由。}\par

他用其他话语明确写道：\Quote{对情欲无节制的倾向，其止息乃是饥饿或绞索。}\par

喜剧诗人们在贬低斯多葛派芝诺（Zeno）的教导时，也证实了以下内容：\par

\Quote{他宣扬一种虚妄的哲学：\newline{}他教人节食，收徒收一个面包，\newline{}以干无花果为调味，饮水为饮。}\par

所有这些人，毫不羞愧地坦承谨慎所带来的益处。而那真实且不违理性的智慧，不依赖空言与神谕，而依赖无敌的防御甲胄与强有力的奥迹，致力于神圣诫命、锻炼与实践，按照从圣言所受的灵感，领受神圣的能力。\par

已经，诗人们描述宙斯的埃癸斯为\par

\Quote{可畏的，四周环绕着恐怖，\newline{}其上刻有纷争、勇力与令人胆寒的溃败；\newline{}其上还有戈耳工之首，可怖的怪物，\newline{}可怕、凶恶，是持埃癸斯之宙斯的标志。}\par

但对于那些能够正确理解救恩的人，我不知道还有什么比法律的严肃及其女儿——敬畏——更为珍贵。因为当有人说音调过高时，正如主在某些人身上所说的那样，致使那些向往祂的人中，有些可能唱走调，我并不认为这真是音调过高，而只是针对那些不愿负天主之轭的人而言。因为对于松弛软弱的人，中等的似乎过高；对于不义的人，所遭遇的似乎过于严厉。因为那些因对罪的好感而倾向于宽恕的人，认为真理是严苛，严厉是野蛮，而那个不与它们一同犯罪、不被它们拖累的人，则被视作冷酷无情。悲剧因此对普鲁托（Pluto）写得很好：\par

\Quote{你问，你将来到何种神明面前？\newline{}祂不知怜悯也不施恩惠，\newline{}只爱赤裸的正义。}\par

因为虽然你尚且不能遵行法律所吩咐的事，然而，考虑到法律中为我们树立了最高贵的榜样，我们就能滋养和增长对自由的热爱；这样，我们将尽其所能更热切地获益，邀请一些事，效法一些事，畏惧另一些事。因为古时那些按照法律生活的义人，\Quote{并非来自传说的橡树，或来自岩石}；因为他们愿意真正地从事哲学，便将自己完全奉献给天主，并被归于信德之下。\Person{芝诺}曾论及印度人说，他宁愿看见一个印度人被烤，也不愿学会所有关于忍受痛苦的论证。但我们每天在眼前展示出大量的殉道者，他们被烧死、刺死、斩首。这一切都是法律所激起的敬畏——作为引导人归向基督的\OriginalTerm{导师}{pædagogue}——训练他们用鲜血彰显自己的虔敬。\Quote{天主屹立在诸神的会中；祂在诸神中施行审判。}他们是谁？那些超越快乐、超越情欲、认识自己所行的人——\OriginalTerm{诺斯替者}{Gnostics}，他们大于世界。\Quote{我曾说：你们是神；都是至高者的儿子。}主在对谁说话？对那些尽量拒绝一切属于人之事的人。宗徒又说：\Quote{你们不再属于血肉，而是属于圣神。}\BibleRef{罗 8:9}他又说：\Quote{我们虽然在血肉中行事，却不凭血肉作战。}\BibleRef{格后 10:3}\Quote{因为血肉之身不能继承天主的国，可朽坏也不能继承不朽坏。}\BibleRef{格前 15:50}\Quote{看，你们要像人一样死去。}圣神这样说，为责备我们。\par

我们必须操练自己，谨慎对待那些易受情欲支配的事物，如同真正的哲学家那样逃避那些激发贪欲的食物，以及在床笫和奢侈中放荡不羁的行为；而那些倾向于奢侈的感觉，对他人是实在的酬报，对我们却不再如此。因为天主最大的恩赐是自制。祂亲自说过：\Quote{我决不离开你，也决不弃舍你。}\BibleRef{希 13:5}因为祂已按真正的拣选判定你为配得。这样，当我们虔诚地努力前进时，我们将从信德到信德，戴上主温和的轭，一位御者驱策我们每人迈向救恩，好赢得有福的适当果实。\Quote{\Emph{锻炼是}}，根据科斯岛的\Person{希波克拉底斯}，\Quote{\Emph{不仅是身体的健康，也是灵魂的健康——不畏劳苦——对食物有贪婪的食欲。}}\par

% structure: p0137 source=h2 target=subsection reason=relative-heading-rank; base=h2

\subsection{第二十一章 诸哲学家论至善}

\Person{伊壁鸠鲁}将幸福置于不饥、不渴、不冷，说出了那句貌似神明的话，不虔地宣称即使在这一点上他也要与父神朱庇特争战；他教导说，幸福就是胜利，仿佛这是活在污秽中的猪群，而非理性的哲学家。因为受快乐统治的是\OriginalTerm{昔兰尼学派}{Cyrenaics}和\Person{伊壁鸠鲁}；这些人明确地说，快乐地生活是终极目的，快乐是唯一的圆满之善。\Person{伊壁鸠鲁}还说，痛苦的消除就是快乐，并且说应当优先选择那从自身起初就吸引人的东西，即完全处于运动中的东西。\Person{狄诺玛库}和\Person{卡利丰}说，终极目的是尽其所能达到并享受快乐；\Person{逍遥学派的希罗尼穆斯}说，大目的是无忧无虑地生活，唯一的终极之善是幸福；同样属于该学派的\Person{狄奥多罗斯}也宣称，目的是无扰而善地生活。实际上，\Person{伊壁鸠鲁}和\OriginalTerm{昔兰尼学派}{Cyrenaics}说快乐是第一要务；因为为了快乐，他们说，美德被引入，并产生了快乐。按照\Person{卡利丰}的追随者，美德是为了快乐而被引入，但后来，看到自身的美，它使自己与第一原则即快乐同样受到珍视。\par

但\Person{亚里士多德学派}规定，按照美德生活是目的，但并非每个有美德的人都达到幸福或目的。因为智者，为不自愿的厄运所烦恼和牵累，并因此愿意逃离生命，既不幸也不快乐。因为美德需要时间；因为那存在于完美之人的[美德]不是一天就能获得的，因为他们说，孩子从不快乐。但人的生命是一个完美的时间，因此幸福由三类善事物完成。因此，穷人、卑贱者、甚至患病者、奴隶，都不能成为其中之一。\par

另一方面，斯多葛派的芝诺认为终向是依照美德生活；克莱安西斯认为终向是合乎自然地正确运用理性，他将其理解为选择合乎自然的事物。他的朋友安提帕特鲁斯认为终向在于持续且坚定不移地选择合乎自然的事物，并排斥那些违反自然的事物。阿凯德摩斯则将终向解释为：在选择合乎自然的最大和首要事物时，不可能超越它。此外，帕奈蒂乌斯宣称终向是依照自然赋予我们的手段而生活。最后，波西多尼乌斯说终向是致力于沉思真理和宇宙秩序，并尽己所能塑造自己，丝毫不受灵魂非理性部分的影响。一些后来的斯多葛派将大终向定义为合乎人的构成而生活。何必提及阿里斯通？他说终向是漠然；但漠然只是放弃了漠然之物。我要提出赫里鲁斯的意见吗？赫里鲁斯规定终向是依照知识而生活。因为有些人认为较晚近的学园派门徒将终向定义为心灵对其自身印象的稳定抽象。此外，逍遥学派的吕库斯常说最终终向是灵魂的真喜乐；而勒乌奇穆斯则说是在善中的喜乐。同为逍遥学派的克里托劳斯说终向是合乎自然地正确流淌的生命之完美，所指的完美由传统的三类达成。\par

然而，我们不能满足于此，而是应尽己所能引述自然论者就此观点的教导；因为他们说克拉佐美奈的阿那克萨戈拉确认沉思及由此而来的自由是生命的终向；以弗所的赫拉克利特认为是满足。本都的赫拉克利德斯记载，毕达哥拉斯教导说，对数字完美的认知是灵魂的幸福。阿布德拉学派也教导终向的存在。德谟克利特在其著作\WorkTitle{论至善}中说终向是愉快，他称之为福祉，并常高呼：\Quote{因为快乐与无快乐是达到成熟年龄者的界限；}赫卡泰乌斯认为终向是自足；库齐库斯的阿波罗多图斯认为是愉悦；瑙西芬尼则认为是无畏，他说这正是德谟克利特所谓的不动心。除此之外，狄奥提摩斯宣布终向是善的完美，他称之为福祉。再者，安提斯泰尼认为是谦卑。至于被称为安尼刻里派、属于昔兰尼学派的人，并没有为整体生命规定明确的终向；而是说每个行动都有其自身的终向，即从该行动产生的快乐。这些昔兰尼派拒绝伊壁鸠鲁对快乐的定义，即痛苦的消除，称其为死人的状态；因为我们不仅因快乐而喜悦，也因同伴关系和荣誉；而伊壁鸠鲁认为灵魂的一切喜乐都源于先前肉体的感觉。美特罗多鲁斯在其著作\WorkTitle{论我们自身幸福的源泉大于来自对象的幸福}中说：\Quote{灵魂的善若非肉体的健康状态及其持续的确切希望，又是什么呢？}\par

% structure: p0142 source=h2 target=subsection reason=relative-heading-rank; base=h2

\subsection{第二十二章 柏拉图的意见：至善在于肖似天主，及其与圣经的一致}

此外，哲学家\Person{柏拉图}说，目的是双重的：一种是可传达的，首先存在于理念形式本身，他称之为\Quote{善}；另一种是分有它并从它获得相似性的，例如那些获得德性和真正哲学的人。因此，\Person{克莱安特斯}在第二卷\WorkTitle{论快乐}中也说，\Person{苏格拉底}到处教导说，正义的人与幸福的人是同一的，并咒骂那第一个将正义与有用分开的人，认为他做了不虔敬的事。因为那些将有用与根据法律正当的事分开的人确实是不虔敬的。\Person{柏拉图}自己说，\OriginalTerm{幸福}{\GreekText{εὐδαιμονία}}是正确地拥有精灵，而灵魂的统治部分被称为精灵；他称\OriginalTerm{幸福}{\GreekText{εὐδαιμονία}}为最完美、最完全的善。有时他称它为一致和谐的生活，有时称它为符合德性的最高完善；他将此置于对善的知识和对天主的相似上，表明相似就是正义和圣洁加上智慧。因为，难道不正是这样，我们的一些作者理解到，人一受造就接受了\Quote{按肖像的}，而\Quote{按模样的}要在以后完满时接受吗？现在，\Person{柏拉图}教导说，有德之人将拥有这种伴随谦卑的相似，他解释如下：\BibleQuote{那谦卑自己的，必被高举。}\BibleRef{路 14:11}他因此在\WorkTitle{法律篇}中说：\Quote{天主，正如古语所说，占据万物的开始、中间和终结，在绕行圆周时直行而过。正义始终伴随祂，惩罚那些背离神圣法律的人。}你看他如何将敬畏与神圣法律联系起来。他因此加上：\Quote{那愿幸福的人，紧抓住它，将谦卑而美好地跟随。}然后，接续这些话，并以敬畏告诫，他加上：\Quote{那么，什么行为是蒙天主喜悦并符合天主的？那用一个古词概括的：\InnerQuote{\Emph{相似者必喜爱相似者}}，如同合乎比例者；但不合比例者既不相爱，也不与合乎比例者相爱。因此，那愿蒙天主喜爱的人，必须尽己所能成为像祂一样的。基于同样理由，我们的自制之人蒙天主喜爱。但没有自制的人是不像且不同的。}他说这是一个古老的教义，指出从他而来的教导源自法律。在\WorkTitle{泰阿泰德篇}中承认邪恶围绕必死之本性及此地之后，他加上：\Quote{因此我们必须尽快逃离此处。因为逃离就是尽己所能相似天主。而相似就是成为圣洁、正义且具有智慧。}\Person{柏拉图}的侄子\Person{斯彪西波}说，幸福是那些按本性生活之人的完美状态，或是善的状态；所有人都渴望这种状态，但只有善人达到宁静；因此德性是幸福的创造者。而\Person{卡尔西登的色诺克拉底}定义幸福为严格意义上的德性及附属能力的拥有。然后他清楚地说，其所在处是灵魂；其实现者是德性；这些德性的部分形成可称赞的行为、良好的习性与倾向、运动及关系；而身体及外在之物并非没有这些。因为\Person{色诺克拉底}的门徒\Person{波勒摩}似乎认为幸福是对一切善物、或大多数及最大善物的充足。他因此主张，幸福从不离开德性；德性脱离身体及外在之物便足以幸福。这些就如此吧。对这些意见的反驳将在适当时候提出。但我们有责任达到那未完成的目的，服从命令——即天主——并按照它们生活，无可指摘且明智地，通过神圣旨意的知识；并尽可能地按正确理性相似，就是目的，并恢复为圣子完全接纳的义子身份，那永远通过大司祭荣耀圣父，祂屈尊称我们为兄弟和同承继者。宗徒简洁地描述这目的，在\BibleBook{致罗马人}{书}中写道：\BibleQuote{现在，但你们从罪中得了释放，成为天主的奴仆，你们有果实归于圣洁，结局就是永生。}\BibleRef{罗 6:22} 又视望德为双重的——所期望的和已领受的——他现在教导说目的是望德的重获。\BibleQuote{因为忍耐，}他说，\BibleQuote{生老练，老练生望德；望德不叫人羞耻，因为天主的爱藉着所赐给我们的圣神，已倾注在我们心中了。}\BibleRef{罗 5:4} 因着这爱和望德的恢复，他说，在另一处：\BibleQuote{\Emph{那安息}已为我们存留。}你会在\BibleBook{}{厄则克耳}中找到类似的，如下：\BibleQuote{犯罪的人，必丧亡。但若一个人是正义的，且行正义与公道，没有在山上吃过，也没有举目向以色列家的偶像，没有玷污邻人的妻子，没有亲近在月经期的妇人（因为他不愿人的种子被辱），也不欺压人；他归还债务人的抵押，不取利息；他转手不作恶；在人与人之间施行真实审判；遵行我的律例，谨守我的诫命，而行真理；他是正义的，必得生存——这是上主阿多奈说的。}\BibleRef{则 18:4} \Person{依撒意亚}也这样劝诫未信者要生活严肃，并劝诫诺斯替者要专注，证明人的德性与天主的德性不相同，他这样说：\BibleQuote{你们应寻求上主，找到祂时就呼求祂。当祂临近你们时，愿恶人离弃自己的道路，不义之人离弃自己的思念；愿他归向上主，必得怜悯，}直到\BibleQuote{你们的思念与我的思念不同。} 那么，我们，按照高贵的宗徒，\BibleQuote{靠着圣神，基于信德，等候正义的望德。因为在基督内，受割损不受割损都无关紧要，唯有关乎运行爱德的信德。}\BibleRef{迦 5:5} 我们愿你们每一位都显出同样的勤奋，以达至望德的圆满确信，直到\BibleQuote{永远按默基瑟德的品级为大司祭。}\BibleRef{希 6:11} 同样，与\Person{保禄}一样，\Quote{全德之智慧}说：\BibleQuote{那听我的，必安然住在望德中。}\BibleRef{箴 1:33} 因为望德的重获也被同一术语称为\Quote{望德}。对于\Quote{必安居}一词，它最美地加上了\Quote{安然}，表明这样的人已得安息，因为他已领受了所期望的望德。因此也加上：\BibleQuote{并得安静，不怕任何灾祸。}宗徒在\BibleBook{致格林多人}{第一封书信}中公开而明确地说：\BibleQuote{你们该效法我，如我效法基督一样，}\BibleRef{格前 11:1} 为使这事成就。如果你们属于我，而我属于基督，那么你们就是基督的效法者，而基督是天主的效法者。\Emph{因此，相似天主，使人尽可能成为正义、圣洁且有智慧，}他以此作为信德的目标，而目的就是藉信德所实现的应许的重获。从这些学说涌出上面我们所指定的那些关于\Quote{目的}之论者的泉源。但这些已足够。\par

% structure: p0144 source=h2 target=subsection reason=relative-heading-rank; base=h2

\subsection{第二十三章  论婚姻}

既然快乐与情欲似乎属于婚姻的范畴，那么也当论及婚姻。婚姻是男女为了生育合法子女而进行的首次结合。因此，喜剧诗人\Person{默南德}如此写道：\par

\Quote{为了生育合法子女，\newline{}我将我的女儿赐予你。}\par

我们要问：是否应当结婚？这是所谓相对的问题之一。因为有些人必须结婚，人必须处于某种状态，而且他必须与处于某种状态的某人结婚。并非人人都要结婚，也非总是如此。而是有适宜的时间、适宜的人和适宜的年龄。既不是每个人都应娶妻，也不是任何女人都该娶，也不是总是、也不是任何方式、也不是轻率。只有处于特定情况的人、这样的人、在必要的时候、为了子女的缘故、并且双方在各方面相配、妻子是出于爱情而非强逼或胁迫而爱那爱她的丈夫。因此，\Person{亚巴郎}视自己的妻子为妹妹，说道：\Quote{她实在是我的妹妹，是我父亲的女儿，却不是我的母亲的女儿；她做了我的妻子。}\BibleRef{创 20:12}教导我们，同母的子女不应结婚。让我们简要地追溯历史。\Person{柏拉图}将婚姻列在外在好事之中，为了种族的延续，并将某种永续性如同火炬般传递给子子孙孙。\Person{德谟克利特}拒绝婚姻和生育子女，理由是从中产生的许多烦恼和更必要的事务的抽离。\Person{伊壁鸠鲁}以及那些将善置于快乐和没有烦恼痛苦中的人也同意。按照斯多葛派的意见，婚姻和养育子女是无所谓的事；而按照逍遥派，则是善。简言之，这些人循着他们的教条言词，沦为快乐的奴隶；有的用妾，有的有情妇，大多数用少年。那在花园里与情妇在一起的所谓四贤哲，以行动尊崇了快乐。那么，那些判断某些事不适合自己，却命令别人去做，或相反的人，不能逃脱用牛与驴同轭的诅咒。圣经简短地指明了这一点，说：\Quote{你所憎恶的事，不可对别人做。}\BibleRef{多 4:15}\par

但赞同婚姻的人说，大自然使我们适于婚姻，这从我们身体的构造（男性和女性）可以明显看出。他们不断宣告那条命令：\Quote{你们要生育繁殖。}\BibleRef{创 1:28}尽管情况如此，但对他们来说，似乎可耻的是，天主所造的人竟比无理性的受造物更放荡，那些动物并不与许多滥交，而是与同种的一个交配，例如鸽子、斑鸠以及类似的动物。此外，他们说：\Quote{无子女的人缺乏按照自然的完美，因为没有将自己的适当继任者置于自己的位置。因为完美的人是从自己产生出与自己相似者，或者更确切地说，当他看到自己产生了同样的东西时；即，当那被生者达到了与生者相同的本性。}因此，我们必须无论如何结婚，既为了我们的国家，为了子女的延续，也为了世界的完美（就我们而言）；因为诗人们也怜悯那半完美而无子女的婚姻，却宣称多产者是幸福的。但是主要是身体的疾病显示婚姻的必要性。因为妻子的照料和持续的忠诚似乎超过所有其他亲戚朋友的忍耐力，正如在同情心方面超越他们；而最重要的是，她乐于耐心看护。的确，按照圣经，她是必要的帮助者。那么，喜剧诗人\Person{默南德}，一面贬低婚姻，一面又陈述其益处，如此回答一个说过以下话的人：\par

\Quote{我厌恶这事，\newline{}因为你处理得笨拙。}\par

然后他补充道：\par

\Quote{你看到其中的困难和令你烦恼的事，\newline{}但你没有看到益处。}\par

如此等等。\par

现在，婚姻对年长者有帮助，提供一个配偶来照顾自己，并养育子女来赡养自己的晚年。\par

\Quote{因为对一个人死后，他的儿女带来名声，\newline{}正如软木浮子托住渔网，\newline{}从深处救起钓线。}\par

据悲剧诗人\Person{索福克勒斯}所说。\par

此外，立法者不允许未婚者担任最高行政官职。例如，斯巴达的立法者不仅对独身处以罚款，而且对一妻制、晚婚和独身生活也处以罚款。而著名的\Person{柏拉图}命令未结婚的人向公共金库缴纳妻子的赡养费，并向官员支付一笔适当的费用。因为，如果他们不结婚、不生育子女，他们就会尽其所能造成人的匮乏，并瓦解国家和由之构成的世界，不敬地取消了神圣的生育。逃避与妻子和儿女共同生活也是怯懦和软弱的。因为，失去某物是坏事，那么拥有它一定是好事；其他事物也是如此。但失去子女，他们说，是最大的坏事之一；因此拥有子女是好事；如果是这样，婚姻也是好事。有话说：\par

\Quote{没有父亲，就永远不会有孩子，\newline{}没有母亲，孩子的受孕就不能发生。\newline{}婚姻造就父亲，如同丈夫造就母亲。}\par

因此，\Person{荷马}使一件事成为恳切祈求的对象：\par

\Quote{丈夫和房屋；}\par

但并非单纯如此，而是伴随着良好的和谐。因为其他人的婚姻是放纵的协议；而哲学家的婚姻则导致合乎理性的和谐，命令妻子不要以外表装饰自己，而要以品德；并嘱咐丈夫不要将他们的婚配妻子当作情妇，以肉身纵欲为目的；而要利用婚姻作为整个生活的帮助，并达到最好的自制。\par

依我之见，远比按季播种的麦粒和大麦种子更为卓越的，是那被播种的人，万物为他而生长；而节制有度的农夫常播撒那等种子。一切污秽和玷污的行为，必须从婚姻中清除；这样，非理性动物的交合就不会被用来指责我们，说它比人类的生育结合更合乎自然。必须承认，其中有些动物在受命之时便停止，将受造之物交托于天主眷顾。\par

悲剧作家笔下，\Person{波吕克塞娜}虽被杀害，却仍被描述为临终时极其注意体面地倒下——\par

\Quote{隐藏那本应从人眼前遮掩之物。}\par

婚姻对她而言是一场灾祸。受制于情欲并屈服于它们，是极端的奴役；而克制情欲则是唯一的自由。因而神圣的\WorkTitle{圣经}说，那些违犯诫命的人被卖给陌生人，即违背本性的罪，直到他们回头并悔改。所以婚姻作为神圣的形象，必须保持纯洁，免受那些玷污它的事物侵扰。我们应当与主一同苏醒，并怀着感恩和祈祷入睡——\par

\Quote{无论在你入睡时，还是当圣光来临，}\par

在全部生活中宣认主；灵魂中拥有虔诚，身体上实行节制。因为引导言谈举止得当，是悦乐天主的。污秽的言语是通往厚颜之路；而两者的终局都是污秽的行为。\par

\WorkTitle{圣经}劝诫婚姻，且不允许解除婚姻关系，这明确记载在法律中：\BibleQuote{除非为奸淫的缘故，你不可休妻；}并且它视在配偶尚在时离异另娶为奸淫。妻子若不过分装饰打扮，便免于诽谤猜疑，同时她专心于祈祷和恳求；避免频繁离家，尽量闭门不接见所有非亲属，并视家务琐事胜过不当的轻浮。经上说：\BibleQuote{娶被休之妻的人犯奸淫；若有人休妻，便使她成为淫妇，}即迫使她犯奸淫。休妻者不仅自己犯此罪，娶她的人也有份，因为给了那妇人犯罪的机会；倘若他不娶她，她便会回到丈夫身边。那么法律是什么呢？\BibleRef{肋 20:10}为遏止情欲的冲动，它命令淫妇一经定罪，便要处死；若属司祭家族，则投入火中。\BibleRef{肋 21:9}奸夫也被用石头打死，但不在同一地点，甚至他们的死亡也不得相同。法律与福音并不相悖，而是彼此一致。否则怎能如此？因为一位主是两者的作者。行淫的女子活在罪中，对诫命而言是死的；但悔改的女子，因生活的转变而仿佛重生，获得生命的再生；旧日的娼妓已死，而藉悔改重生者又回到了生命。圣神为厄则克耳所说的话作证，宣告：\BibleQuote{我不喜欢罪人之死，而喜欢他回头。}\BibleRef{则 33:11}他们被石头打死，这是由于心硬，对不信的法律如同死人。但司祭之女所受的刑罚加重，因为\BibleQuote{凡多给谁的，就向谁多要。}\BibleRef{路 12:48}\par

鉴于章节的篇幅和数量，我们在此结束\WorkTitle{杂记}第二卷。\par

\SourceRecord{Translated by William Wilson.；Ante-Nicene Fathers；Vol. 2.；Edited by Alexander Roberts, James Donaldson, and A. Cleveland Coxe.；Buffalo, NY: Christian Literature Publishing Co.,；1885.；<http://www.newadvent.org/fathers/02102.htm>.}

% structure: p0001 source=h1 target=section reason=page-title

\section{杂集（第三卷）}

% structure: p0002 source=h2 target=subsection reason=relative-heading-rank; base=h2

\subsection{第一章——驳斥\OriginalTerm{巴西里德斯}{Basilidis}关于\OriginalTerm{节欲}{continentia}与\OriginalTerm{婚姻}{nuptiis}的主张}

\OriginalTerm{华伦提尼派}{Valentiniani}因从上由神圣流溢引出结合，故认肯婚姻；而\OriginalTerm{巴西里德斯派}{Basilidis}的追随者却说：\Quote{当宗徒问及是否最好不娶妻时，主回答：\BibleQuote{这话不是人人能领受的。因为有些阉人是从母胎生的，有些是被人使然的。}} \BibleRef{玛 19:11}他们这样解释这句话：\Quote{有些人从出生起就自然厌恶女性，他们若使用这自然的自制，不娶妻便是做得好。这些人，他们说，是从出生时的阉人。至于从需要而来的，则是戏院中的操练者，他们因贪图荣誉而自制。此外，因意外被阉割的，是因需要而成了阉人。因此，因需要而成为阉人的，并非按\OriginalTerm{理性}{logos}而成阉人。但为天国的缘故自阉的人，他们说，是为逃避婚姻的烦恼，因他们害怕管理家业的重担和忧虑。至于那句话：\BibleQuote{结婚比被火烧更好，} \BibleRef{玛 19:11}他们说，宗徒的意思是：不要把你的灵魂投进火里，日夜反抗，害怕从节欲中坠落。因为当灵魂忙于反抗时，它就与希望分离了。}——因此，\Person{伊西多尔}在\WorkTitle{道德论}中如此说：\Quote{你当耐心忍受好争吵的妇人，免得你被天主的恩宠夺去；并且当你在种子中燃起欲火时，你当存着良善的良心祈祷。但是，当你的感恩祈祷堕入祈求，而后你站立，却仍不免跌倒和恐惧时，你当娶妻。若有一个青年，或贫穷，或软弱，按\OriginalTerm{理性}{logos}不适合娶妻，他不应离开弟兄；他应说：我进入了圣所，我什么也不能忍受。若他起了什么疑虑，他应说：弟兄，求你按手在我身上，使我不犯罪；他立刻会在心里和身体上体验帮助。只愿他愿意行善，便能达到。有时我们口头说：我们不愿犯罪；但我们的心却趋向罪恶。这样的人因恐惧而不作所愿的事，免得刑罚加于他身上。但人类有些东西是必需且自然的。需要衣服既是必需品又是自然的：而性欲是自然的，但不是必需的。}我引这些话是为了责备巴西里德斯派，他们生活不正，仿佛因完美而拥有犯罪的能力，或全然因本性而得救，即使现在犯罪，因他们被本性尊贵所选。但最初的教义建筑师并没有给他们行这些事的权力。因此，愿他们不要仰望基督的名，并比外邦人中最放纵者更放纵地生活，使这名蒙受咒骂。\BibleQuote{因为这样的人是假宗徒，是欺诈的工人，}直到那句：\BibleQuote{他们的结局必按他们的行为而定。}所以，\OriginalTerm{节欲}{continentia}是根据对天主的宣认而轻视身体；因为节欲不仅涉及性事，也涉及灵魂不正当地贪求、不满足于必需的其他事物。节欲也涉及舌头、获取、使用和欲望。但它并不只教导我们有节制，因为它赐予我们\OriginalTerm{节制}{temperantia}，作为神圣的能力和恩宠。因此，我们必须说明我们对于所讨论之事的意见。我们确实称贞洁以及那些蒙天主赐予此恩的人为有福的；我们赞叹一夫一妻制，即仅在一次婚姻中的正直；但我们却说应当怜悯他人，并\BibleQuote{彼此分担重担} \BibleRef{迦 6:2}，免得\BibleQuote{谁若自以为站得稳} \BibleRef{格后 10:12}自己也会\BibleQuote{跌倒}。关于第二次婚姻：\BibleQuote{若你被火烧，}宗徒说，\BibleQuote{就当结婚。} \BibleRef{格前 7:9}\par

% structure: p0004 source=h2 target=subsection reason=relative-heading-rank; base=h2

\subsection{第二章——驳斥\OriginalTerm{卡尔波克拉特斯}{Carpocratis}与\OriginalTerm{埃皮法尼斯}{Epiphanis}关于\OriginalTerm{妇女共有}{Feminarum Communitate}的主张}

\Person{卡尔波克拉特斯}的那些后裔和\Person{埃皮凡}，认为妻子应当公有；由于他们，对基督之名的最大的诽谤得以散发。这位\Person{埃皮凡}，也流传有其著作，是\Person{卡尔波克拉特斯}的儿子，母亲是\Place{亚历山大里亚}人，名某某；父亲是\Place{亚历山大里亚}人，母亲是\Place{刻法勒尼亚}人。他只活了十七年，在\Place{刻法勒尼亚}的\Place{萨摩}城，被当作神而受敬拜。在那里，用巨大石头建造了圣殿、祭台、神庙、博物馆，并祝圣；每当新月，\Place{刻法勒尼亚}人聚集，祭祀、奠酒、宴饮、唱赞美诗，庆祝\Person{埃皮凡}被列入神祇的生日。他从父亲学习了全部学科和\Person{柏拉图}的哲学。他是\OriginalTerm{一元知识}{monadicæ cognitionis}的首脑。现今的\Person{卡尔波克拉特斯}派(\OriginalTerm{卡尔波克拉特派}{Carpocratiani})异端也是从他而出。他在\WorkTitle{论正义}一书中说：\Quote{天主的正义是某种与平等共融的共融。确实，天空均匀地四面延伸，环绕整个大地。黑夜平等地展示所有星辰；天主从高天平等地倾注了白日之主和光明之父——太阳，给所有能看见的人（他们都共同观看），因为他不区分富人和穷人、民众的首领、愚人和智者、女人和男人、自由人和奴隶。对牲畜也是如此。他平等地倾注了这共同的东西给所有生物。他藉着自己的正义确认了善人和恶人，因为无人能多得，也不能从邻人夺取，好使自己有双倍的光。太阳给所有生物提供共同的滋养，以共同的正义平等地赐予一切：对于这类事物，牛类像牛一样，猪类像猪一样，羊类像羊一样，其余一切也是如此。因为正义在这些事上显现为共融。然后藉着共融，一切按照它们的种类同样地繁殖，共同的滋养由所有在地上吃草的牲畜平等地食用；如此，没有任何法律限制，而是出于那赐予者——命令提供者的供应，适宜而正义地临于一切。但对生育也没有设立法律，否则早已废除：它们平等地繁殖和生产，拥有来自正义的内在共融：万物的创造者和父亲，藉着他的正义颁布法律，平等地共同赐给所有生物眼睛以观看，不区分女性和男性，也不区分理性的和非理性的，简言之，不区分任何事物；而是以平等和共融同样地分配视力，凭一个命令赐予一切。但人类的律法，他说，由于不能惩罚无知，反而教导人违法：因为律法的私有权分割并啃噬了神圣律法的共融；不明白宗徒的话说：\InnerQuote{我藉着法律认识了罪。}他说，\Quote{我的}和\Quote{你的}这些概念是通过律法产生的，以致人们不再共同享用（因为原本是共有的）土地、财产，甚至婚姻。因为神明为大家共同创造了所有事物，不拒绝麻雀或小偷；谷物和其他果实也是如此。一旦共融和平等被破坏，就产生了偷牲畜和果实的小偷。既然天主原为人类创造了万有共同享有，并使女人与男人共同结合，也使所有生物同样地连结，他就宣布了正义，即共融与平等。但那些这样出生的人，否认了那促成他们生育的共融。他说，如果一个人娶了一个妻子，而所有人都可以分享，就像其他正直的动物一样。}他说了这些话后，又补充说：\Quote{他给男性注入了强烈而有力的欲望以延续后代，法律、习俗或其他任何事物都不能消除：这是天主的法令。}既然他通过这些话明显破坏了法律和福音，我们怎能进一步在我们的论述中审查他呢？因为一方面说：\BibleQuote{不可奸淫。} \BibleRef{出 20:13} 另一方面说：\BibleQuote{凡注视妇女而生贪念的，已经犯了奸淫。} \BibleRef{玛 5:28} 而法律说的\BibleQuote{不可贪恋，} \BibleRef{出 20:17} 表明惟一天主，即藉着法律、先知和福音所预告的那一位。因为他说：\BibleQuote{不可贪恋你邻人的妻子。} 但邻人对犹太人来说不是另一个犹太人：因为他是兄弟，有同一圣神；因此，他所说的\Quote{邻人}是指外邦人。那能分享圣神的人，怎能不是邻人呢？因为\Person{亚巴郎}不仅是希伯来人的父亲，也是外邦人的父亲。若一个妇女通奸，与她行淫的人处死刑，\BibleRef{申 22:22}那么显然，那条诫命\BibleQuote{不可贪恋你邻人的妻子}是对外邦人说的：好使人按照法律，既远离邻人的妻子，也远离姐妹，然后清楚听见主说：\BibleQuote{我却对你们说：不可贪恋。} 这个\Quote{我}字的附加，显示了诫命的更大效力。但\Person{卡尔波克拉特斯}和\Person{埃皮凡}与天主作战，甚至在那广为流传的\WorkTitle{论正义}一书中，也从他接着说的话显明：\Quote{从此，谁若认为立法者的话\InnerQuote{不可贪恋}是可笑，就应听到更可笑的话：\InnerQuote{你邻人的财物。} 因为那赐下欲望以维持生育者，却命令人除去它，尽管他不从任何动物那里除去它。而那句\InnerQuote{我邻人的妻子}，更可笑地强迫共融成为私有。}这些就是杰出的\Person{卡尔波克拉特斯}派的教义。他们说，这些人和一些其他类似邪恶的效仿者，聚会用餐（我不愿称他们的聚集为\OriginalTerm{爱宴}{agape}），男人和女人一起，在吃足激发情欲的食物后，熄灯（因为灯光会为他们的淫秽正义带来羞耻），背光而行，任意与什么人交合；并在这种\OriginalTerm{爱宴}{agape}的共融中操练后，白天就随意向妇女要求\Person{卡尔波克拉特斯}的法律（因为学习神圣的事是亵渎）的服从。依我看来，这些法律本该加在\Person{卡尔波克拉特斯}身上，即狗、猪和山羊的淫欲。我认为他也不正确理解了\Person{柏拉图}在\WorkTitle{共和国}中说的应使众人妻子共有：他所说的是指尚未出嫁的女子应由求婚者共有，正如剧场也由观众共有；但每个女子一旦被先占者占有，就不再共有，即已出嫁者。但\Person{克桑托斯}在\WorkTitle{魔法}中写道：\Quote{巫师与母亲和女儿交合，他们说与姐妹交合是允许的，妻子是共有的，不是强暴或秘密地，而是双方同意，当一方愿意娶另一方的妻子时。}关于这些和类似的异端，我认为\Person{犹达}在书信中预言般地说：\Quote{同样地，这些做梦的人（因为他们不警醒地追求真理），直到\InnerQuote{他们的口说夸大的话。}} \BibleRef{犹 1:8}\par

% structure: p0006 source=h2 target=subsection reason=relative-heading-rank; base=h2

\subsection{第三章——柏拉图与其他古人在何种程度上预先教导了马尔基翁派与其他异端者，他们因认为受造物是恶的且人出生是为受罚而禁欲}

的确，如果连柏拉图本人和毕达哥拉斯派，以及后来的马尔基翁派，都认为生育是恶的，那么他们远远不会设立共妻制。但马尔基翁派确实说本性是恶的，来自恶的物质，并由一位公义的工匠和创造者所造。因此，他们不愿充满由创造者所造的世界，而想禁绝婚姻，以此抵抗自己的创造者，并奔向那召叫他们的善者：但并非奔向那位他们说具有不同性情的天主。因此，既然他们不愿在此留下任何属于自己的东西，他们并非出于坚定的意向而节制，而是出于对创造者的仇恨，不愿使用他所造之物。然而，这些因与天主进行不虔敬的战争而脱离自然之思想的人，藐视天主的宽容和良善，尽管他们不愿娶妻，却仍使用受造的食物，呼吸创造者的空气，好成为他的工具，留在他所有的事物中，并聆听他们所传报的一种闻所未闻的新认知，即使他们也因在此接受了福音而应向世界之主感恩。但当我们论述原则时，我们将极其精确地反驳他们。至于我们提及的哲学家，马尔基翁派从他们那里学到生育是恶的，并以此为自家的教义夸耀，他们并不认为本性是恶的，而是那散布真理的灵魂是恶的。因为他们承认灵魂是神圣的，却将其带入这个世界作为惩罚之地。按照他们的看法，灵魂必须进入身体以作补赎。但这教义并不符合马尔基翁派，而符合那些认为灵魂被投入、被捆绑、仿佛从一个器皿倒入另一个器皿的人。反驳他们将是另一个适当的时机，当我们论及灵魂时。因此，赫拉克利特似乎以咒骂来谴责生育，他说：\Quote{因为人出生后愿意活着，并拥有死亡，或者更确切地说，愿意安息；他们也留下儿女，好让死亡发生。}显然，恩培多克勒也同意他，说：——\par

\Quote{我哀悼并痛哭，看见这不幸的陌生世界。}\par

以及更进一步：——\par

\Quote{因为他从活物中造出死物，改变形态。}\par

然后又说道：——\par

\Quote{哀哉！人类是多么不幸而可怜，\newline{}你们是从多少纷争和哀哭中出生的？}\par

西彼拉也说：——\par

\Quote{有死的人们，你们只是肉体，什么都不是；}\par

同样，那位诗人写道：——\par

\Quote{滋养的大地没有生出比人更不幸的东西。}\par

甚至泰奥格尼斯也表明生育是恶的，这样说：——\par

\Quote{对有死的病弱之人来说，最好是不出生，\newline{}也未见明媚阳光的闪耀，\newline{}或者一出生就跨入冥府之门。}\par

悲剧诗人欧里庇得斯接着这些也写道：——\par

\Quote{我们本当聚集开会，\newline{}为出生者哀哭，因为他要遭遇多少苦难；\newline{}但对于已得安息的死者，\newline{}则应带着喜乐和祝贺送葬。}\par

他又说了类似的话：——\par

\Quote{谁知道活着是否就是死去，\newline{}而死去是否就是活着？}\par

看来希罗多德也引述梭伦的话说：\Quote{啊，克罗伊斯，每个人无非是灾祸。}至于他关于克勒奥比斯和比同的故事，无非是谴责生育而赞美死亡。\par

荷马说：\Quote{人类世代如树叶。}但柏拉图在\WorkTitle{克拉底鲁篇}中，将以下言论归于奥斐斯：灵魂在身体内受罚；他说：\Quote{确实，这身体是灵魂的\OriginalTerm{坟墓}{\GreekText{σῆμα}}，}有人认为这是\OriginalTerm{纪念碑}{monumentum}，仿佛它当下就被埋葬；而且因为灵魂通过身体\OriginalTerm{指示}{\GreekText{σημαίνει}}，\Quote{凡它能指示的一切，因此称它为\OriginalTerm{坟墓}{\GreekText{σῆμα}}是恰当的。此外，我认为奥斐斯之所以特别赋予这名称，是因为灵魂在这身体中偿还罪过的惩罚。}值得记住菲洛劳斯所说的。这位毕达哥拉斯派学者这样说：\Quote{古代的圣贤和先知也证明，灵魂与身体结合是为了偿还惩罚，并如同埋葬在\Emph{坟墓}中一般。}甚至品达在谈及厄琉息斯秘仪时也说：\Quote{有福的人，当他看见那地下的共同之物时，便知晓生命的终结，也知晓宙斯所赐的统治权。}柏拉图在\WorkTitle{斐多篇}中同样不避讳这样写：\Quote{此外，那些为我们设立这些秘仪的人，不是别的，}直到\Quote{与诸神共居。}然而，当他说：\Quote{只要我们拥有身体，我们的灵魂与这种恶混合，我们就永远无法充分获得我们所渴望的？}他岂非指明出生是最大恶的起因？而且，在\WorkTitle{斐多篇}中他也证实：\Quote{事实上，那些真正从事哲学的人，别人看不出他们专注于任何其他事，而只是去死和成为死人。}又说：\Quote{同样，在这里，哲学家的灵魂最轻视身体，逃离它，而寻求独处。}这岂非与神圣的宗徒相一致？宗徒说：\BibleQuote{我这个人真不幸！谁能救我脱离这死亡的身体呢？}\BibleRef{罗 7:24}除非他比喻地称那些被拖入罪恶之人的合作为\BibleQuote{死亡的身体}。而且，柏拉图甚至早于马尔基昂就厌恶作为生育之始的交合，他在\WorkTitle{理想国}第一卷中，在赞美老年之后，接着说：\Quote{我希望你知道，当其他（身体的）快感越来越离开我时，我对交谈的渴望和从中获得的快乐却与日俱增。}当提及性事时，他说：\Quote{请说好话！我实在乐意从那如同疯狂粗暴的主人那里逃离。}又在\WorkTitle{斐多篇}中，他谴责出生，说：\Quote{因此，关于这些事在奥秘中所说的，就是这样一个说法：我们人类处于某种牢狱之中。}又说：\Quote{那些被发现比别人活得更虔诚的人，他们就是从这些尘世之地如同从监狱中解脱释放，上升到更高处纯净的居所。}然而，虽然如此，他仍认为天主正确地管理世界；因此他说：\Quote{人不应该自行解脱或逃走。}简言之，他并未给马尔基昂提供视物质为恶的机会，因为他虔诚地论及世界：\Quote{因那位建造者，它拥有一切美善；但因先前的丑陋，世界自身承受了一切在天穹之内产生的不便和不义，并将它们植入生物之中。}他更清楚地补充道：\Quote{它的缺陷是物质性的构造，伴随原始本性而来；因为在它被当前的装饰美化之前，它是极其畸形和无序的。}尽管如此，在\WorkTitle{法律篇}中他也哀叹人类，这样说：\Quote{诸神怜悯被本性劳苦所压迫的人类，为他们设立了劳苦的缓解，即节庆的交替。}在\WorkTitle{厄庇诺米斯篇}中，他进一步追究诸神怜悯的原因，这样说：\Quote{从一开始，被生育本身对任何生物都是沉重的：首先，因为它们参与了在子宫内构成的过程；其次，出生本身，以及之后的抚养和教育，通过无数的劳苦，这一切都发生了，正如我们所有人所说。}那么，赫拉克利特岂非也说出生是死亡？毕达哥拉斯同样，如同苏格拉底在\WorkTitle{高尔吉亚篇}中说：\Quote{我们醒时所见的乃是死亡；我们睡时所见的乃是梦。}不过，关于这些已足够。当我们讨论本原时，我们将考察这些矛盾——也是哲学家所暗示、马尔基昂派以其教义所断言的。然而，我认为已经足够清楚地表明：马尔基昂是从柏拉图那里忘恩负义且无知地接受了外邦异端教义的机缘。现在，让我们继续关于节制的论述。我们曾说，希腊人针对生育子女说了许多不适之处，他们着眼于常伴随生育的种种不便。马尔基昂派不敬地采纳了这些，并对造物主忘恩负义。悲剧这样说：\par

\Quote{人生而不出生更好。\newline{}然后以痛苦的阵痛生下儿女，\newline{}但若生下后，若愚笨的人知道，\newline{}我受苦，因我看出我养活恶人，\newline{}而失去善人。若我养活善人，}\par

\Quote{我混杂的心仍然因恐惧而融化。\newline{}这有何好处？难道仅仅\newline{}吐出灵魂还不够，除非你更多受苦？}\par

而且又同样：\par

\Quote{我心中存有古老的信念：\newline{}人决不应生育子女，\newline{}因他看到我们生育儿女所招致的灾祸。}\par

在接下来的部分，他显然也将灾祸的起因追溯到本原，这样说：\par

\Quote{啊！可怜的人，生来就遭受灾祸，\newline{}人出生，你的悲惨生涯\newline{}从此开始：以太开始向所有\newline{}有气息的凡人传授生命的气息；\newline{}凡人啊，不要恼怒于必死之事。}\par

然后他又说了类似的话：\par

\Quote{所有凡人中，没有一个是有福的，\newline{}也没有一个是没有烦恼的。}\par

接着又：\par

\Quote{唉！何等繁多、何等众多的灾祸临到人，\newline{}何等虚无，其终点无人知晓。}\par

并且同样：\par

\Quote{没有一个凡人永远有福。}\par

依此理由，他们也说毕达哥拉斯派禁绝性事。但我看来，他们倒娶妻以生子，但生子后愿意节制性欲。因此他们神秘地禁食豆子，并非因为豆子产气、难消化、引起纷乱的梦，也非如那诗句所说：\Quote{吃豆子等于啃父亲的头颅；}\par

\Quote{因为啃豆子和啃父亲的头颅是一样的；}\par

但是更确切地说，如果吃了豆子，会使妇女不育。\Person{泰奥弗拉斯托斯}在\WorkTitle{论植物的原因}第五卷中记载，如果把豆荚放在新栽树木的根部，会使树木枯死。而且，家养的母鸡如果经常吃它们，也会变得不育。\par

% structure: p0040 source=h2 target=subsection reason=relative-heading-rank; base=h2

\subsection{第四章：异端者用哪些借口来放纵各类可允许的和情欲的行为。}

在那些受异端引导的人中，我们曾提到本都的马尔西翁，他因与造物主争战而拒绝使用世俗之物。他的节制（如果这可以称为节制）之因正是造物主本身；这个与天主搏斗的巨人自以为与造物主为敌，却勉强保持节制，同时攻击受造物和天主的工程。倘若他们引用主对斐理伯说的话：\BibleQuote{任凭死人埋葬他们的死人，你要跟随我。}\BibleRef{玛 8:22}那么他们应当考虑，斐理伯也带着同样的形像，并没有污秽的尸体。那么，他既然有肉身，怎么没有尸体呢？因为他从坟墓中复活了，主以死亡制伏他的恶习，他却为基督而活。我们也记得卡尔波克拉特斯的观点中那邪恶的妇女共融。当我们谈到尼苛劳的言论时，我们省略了这一点：他们说，他有一个美貌的妻子，在救主升天后，因嫉妒受到宗徒们的责备，他便将妻子带到众人中间，允许她嫁给任何她愿意的人。他们说，这一行为与那句\Quote{应当滥用肉身}的话相符。因此，那些盲目而轻率地追随他的言行、并继续他的异端的人，无耻地放纵淫乱。但我听说，尼苛劳除了娶过的那位妻子外，从未用过其他女人；他的子女中，女儿们以童贞之身终老，儿子则保持纯洁。既然这样，将妻子带到宗徒们中间（据说他因她而犯嫉妒）是对恶习的驱除和净化；而靠努力追求来的对快乐的节制教导了\Quote{滥用肉身}，即操练肉身。因为，我认为，他们不想按主的诫命\BibleQuote{事奉两个主人}\BibleRef{玛 6:24}，即快乐和天主。因此他们说，玛弟亚也曾这样教导：\Quote{要与肉身争战，运用它，不给予任何不洁之物以快乐；要通过信仰和认识增进灵魂。}有些人甚至公开宣称两性结合是神秘的共融，从而污辱了这个名称。正如我们说\Quote{行事}，既指做坏事，也指做好事，将同一名称赋予两者；同样，\Quote{共融}一词也常被使用：因为好的共融在于分享金钱、食物和衣服；而他们却不敬地将任何两性结合称为\Quote{共融。}因此，他们中有人说，当接近我们一位面容姣好的贞女时，说：经上记载：\BibleQuote{凡求你的，就给他。}\BibleRef{玛 5:24}那女子却十分正派地回答，仿佛不明白那人的轻浮：你去与母亲商议婚姻吧。啊，不敬！这些不节制的共融者，情欲的兄弟，不仅玷污哲学，也玷污整个生命，他们篡改并败坏真理，甚至将其掩埋；这些最不幸的人将肉体的交合共融神圣化，并认为这能带领他们进入天主的国。那么，这种共融就引向妓院，当猪和山羊与他们共融后，妓院中那些随叫随到、接纳所有人的娼妓便成为他们最大的希望。\BibleQuote{但你们从基督所学的不是这样。如果你们听过他，并在内中受过教导，如同真理在耶稣基督内，就该脱去你们从前的生活，旧人，那因私欲的迷惑而败坏的。但你们应在心思的灵里更新，并穿上新人，这新人是按照天主在正义和真理的圣洁中创造的。}\BibleRef{弗 4:20}，以肖似天主。\BibleQuote{所以，你们该效法天主，如同蒙爱的儿女；并在爱中行走，就如基督爱了我们，为我们舍弃了自己，作为馨香的祭品和牺牲，献给天主。至于淫乱、一切不洁或贪财，在你们中连提也不可，如同圣徒相称；也不可有猥亵和戏谑的话。}\BibleRef{弗 5:1}因为宗徒教导你们即使言语也要保持贞洁，他写道：\BibleQuote{你们要知道这一点，凡是淫乱的，}等等，直到\BibleQuote{倒要责备。}\BibleRef{弗 5:5}他们的教义出自某本伪经。因此我要引用那作为他们不节制的根源的言辞：不论他们自认是这本书的作者（看他们的疯狂，尽管他们受情欲的驱动而将之归于天主），还是从别人那里曲解地接受了这卓越的教义，那么其言辞如下：\Quote{万物原为一体；当它的一体性不愿再独自存留时，从它发出一种气息，并与它结合，生出了爱子。此后，从那气息又生出另一气息，与之结合，产生了不能看见也不能听见的权能，}直到\Quote{每一个以自己的名字。}因为如果这些人也像瓦伦提尼安派那样设立属灵的共融，或许有人会接受他们的观点；但将肉体的情欲共融引入神圣的预言，是绝望于救恩的人所为。普罗迪科斯的追随者也如此主张，他们以虚假的名称自称\Quote{诺斯替者}：他们自称天生是首生天主的儿子；却滥用这种高贵与自由，随心所欲地生活；他们以情欲为所欲为，自以为不受任何约束，如同\Quote{安息日的主}，并在各方面超越常人，是王室之子。他们说，王没有成文法。首先，他们并不能为所欲为：许多事会阻止他们，即使他们渴望并尝试。而且，他们行事不像君王，倒像鞭刑犯：他们暗中行奸淫，害怕被捉住，避免被判罪，恐惧受惩罚。自由怎会与不节制和污秽的言语相关？\BibleQuote{因为每个犯罪的人都是奴隶，}\BibleRef{罗 6:16}宗徒说。但那将自己交付给一切欲望的人，又如何能建立属天主的生命？主说：\Quote{我却对你们说：不要贪恋。}有人愿意自愿犯罪，并断定奸淫可行，放纵情欲，破坏他人的婚姻，而我们却怜悯那些被迫犯罪的人吗？如果他们在别人的世界到来时，在别人的东西上不忠信，他们就不会拥有真的。一位客人侮辱市民并伤害他们，而不是像外乡人那样使用必需品生活，不冒犯市民？他们行与外邦人憎恨之事相同的事，即不服从法律，行不义、不节制、贪婪、奸淫，却自称只有他们认识天主？他们应当在他人之地正直生活，以真正显示王族的品性。然而，他们与人类立法者和神圣法律为敌，因为他们决心不义地、违法地生活。那位在《户籍纪》中\Quote{刺穿}淫乱者的人，被显示为更属于天主。\BibleQuote{如果我们说，}若望在书信中说，\BibleQuote{我们与他（即天主）相通，却在黑暗中行走，我们就是说谎，不行真理。如果我们像他在光明中一样在光明中行走，我们就彼此相通，他儿子耶稣的血洗净我们的一切罪过。}\BibleRef{若一 1:6-7}那么，这些行如此之事、甚至与世上最坏的人相似的人，如何比世上的人好？我认为，行为相似的人，本性也相似。那些自以为在品级上超越他人的人，也应当在道德上超越，以避免被关入监牢。事实上，如主说：\BibleQuote{除非你们的义德超过经师和法利塞人的义德，你们决进不了天国。}\BibleRef{玛 5:20}关于饮食的节制，由达尼尔证明。\BibleRef{达 1:1}简单说，关于服从，咏唱的达味说：\Quote{青年人如何洁净自己的行径？}他随即听到：\Quote{由谨守你的言语，全心遵守。}耶肋米亚说：\BibleQuote{上主这样说：你们不要学外邦人的道路。}\BibleRef{耶 10:2}另有一些微不足道的人因此说，人是由不同的权能塑造的：上至肚脐的部分属于更神圣的技艺，下至下部则属于低等的；因此他们也追求交合。但他们没有注意到，上部的部分也追求营养，并以某种方式放纵情欲。他们甚至反对基督，他曾对法利塞人说，同一位天主造了人的\BibleQuote{内在}和\BibleQuote{外在}。\BibleRef{路 11:40}而且，欲望并非属于身体，即使通过身体发生。另有一些人，我们称之为\Quote{反律者}，即\Quote{对抗者}和敌对者，他们说，万有的天主是我们的自然之父，他所造的一切都是好的；但他所造的一个存在散布了莠子，生出了恶的本性；这恶将我们所有人都牵连其中，使我们成为父的仇敌。因此，我们也要对抗这恶，以报复父，违背第二者的意愿。既然他说：\Quote{不可奸淫，}我们（他们说）就行奸淫，以破坏他的诫命。对此我们要回答：我们得知，假先知和那些伪装真理的人，凭行为可以认出；如果你们的行为名声败坏，你们怎能还说你们持守真理？要么没有恶，那么你们指控的那位（作为天主的对抗者）就不该受责备，他没有造成任何恶；因为恶树会被一并砍倒；要么恶存在并且持续，请告诉我们，那些赐下的诫命——\Emph{关于正义}、节制、忍耐、含忍等——是善还是恶？如果诫命是恶的，禁止行许多无耻之事，那么恶就会针对自身立法以消灭自身，这不可能；如果诫命是善的，那么他们对抗善的诫命，就是承认自己与善为敌，行恶。而且，救主本人（他们认为只该服从他）禁止仇恨和咒骂，他说：\Quote{当你同对头在路上时，要设法与他友好分手。}那么，他们要么也否认基督的劝告，对抗对头；要么如果他们是朋友，就不愿与他争战。怎么？难道你们不知道，杰出的人们（我仿佛面对你们说话），你们对抗正确的诫命，就是抵抗自己的救恩？你们不是推翻有益的训令，而是推翻自己。主说：\BibleQuote{你们的光也当这样照在人前，叫他们看见你们的善行。}\BibleRef{玛 5:16}而你们却暴露自己的情欲和不节制。此外，如果你们想破坏立法者的诫命，为什么你们只试图破坏那些涉及节制的，如\Quote{不可奸淫}和\Quote{不可与幼童行淫}等，却由于你们的不节制，不破坏由他造成的冬天，使你们在隆冬制造夏天？你们也不能使陆地可航行，海可行走，如同史书所说野蛮的泽尔克西曾想做的？为什么你们不反对所有的诫命？因为他说：\Quote{你们要生育繁殖，}你们这些对抗者应当完全不用交合。当他说：\Quote{我赐给你们一切可吃用的，}你们应当什么也不享用。甚至当他说：\BibleQuote{以眼还眼，}\BibleRef{出 21:24}你们应当不以对抗还对抗。当他命令窃贼偿还\BibleQuote{四倍，}\BibleRef{出 22:1}你们应当给窃贼添上一些。同样，当你们对抗诫命\BibleQuote{你当全心爱你的天主，}\BibleRef{申 6:5}时，你们应当连万有的天主也不爱。又如，当他说：\BibleQuote{你不可铸造神像或偶像，}\BibleRef{申 27:15}你们就应崇拜偶像。那么，你们这些声称反抗造物主、却追求与娼妓和奸夫相似之事的人，怎么不是不敬？你们怎么没有意识到，你们使你们视为软弱的那位显得更伟大？因为成就的是他的意愿，而不是善者的意愿？相反，你们自己倒显出你们称之为父的那位是软弱的。他们还从一些先知书中摘取文句，胡乱拼凑，将其中寓意之言当作直接而字面的解释。他们说，经上记载：\BibleQuote{他们对抗了天主，却得了平安。}\BibleRef{拉 3:15}他们加上\Quote{无羞耻的天主}，并将这句话当作命令性的劝告；他们认为对抗造物主有助于救恩。但\Quote{无羞耻的}并没有写在\Quote{天主}前。即使如此，你们要理解，那被称为魔鬼的才是无羞耻的：或因为他诽谤人，或因为他控告罪人，或因为他是背逆者。因此，所说到的百姓，因自己的罪恶受惩罚时，痛苦呻吟，用这些话发怨言：别的民族行不义却不被惩罚，他们自己却在每件事上受折磨；以致耶肋米亚也说：\BibleQuote{为什么恶人的道路顺利？}\BibleRef{耶 12:1}这与马拉基亚先前提到的相似：\Quote{他们对抗了天主，却得了平安。}因为先知们受天主默感，不仅宣告从天主听到的话，也常以反驳的方式说出百姓中流传的话，并像是报告人们提出的问题；比如刚才提到的这句话。难道宗徒不是针对这些人，在《罗马人书》中写道：\BibleQuote{也不像我们被毁谤、有人谓我们说：让我们作恶以成善，这些人受审判是公正的。}\BibleRef{罗 3:8}这些人，在阅读时以语气歪曲圣经，以迎合自己的欲望，并通过改变某些重音和标点，将那些明智而有益的诫命引向自己的享乐。\BibleQuote{你们以言语激怒了上主，}马拉基亚说，\BibleQuote{却说：我们在什么事上激怒了你？因你们说：凡行恶的，在上主眼中是善的，并且他喜悦他们；公义的天主在哪里？}\BibleRef{拉 2:17}\par

% structure: p0042 source=h2 target=subsection reason=relative-heading-rank; base=h2

\subsection{第五章——记述两类异端者：第一类声称万物皆可随意，并予以驳斥。}

因此，为避免我们更深入地挖掘此点而记起更多荒谬的异端，亦免得我们因需要逐一反驳每一异端而感到羞愧，并使得这些注疏过于冗长，让我们将一切异端分为两类来回应。他们要么教导无分别地生活，要么过度节制，以不敬和仇恨自称持守克制。先处理第一类。如果任何一种生活方式都可随意选择，那么当然节制的生活方式也可选择；如果随意选择任何一种生活方式都是安全的，那么显然，以节制和按德行生活的方式是最安全的。因为既然\BibleQuote{安息日的主}，即使他生活无节制，也无需向谁交代，那么那过适度与节制生活的人，更无需向谁负责。\BibleQuote{凡事都可行，但不都有益处}，宗徒说。如果凡事都可行，那么显然适度与节制也可行。因此，正如那运用自由以行德之人当受称赞，那赐予我们自主权并容许我们按己意生活者更当受尊敬和朝拜，因他未让我们的选择与回避必然服务于任何人。若两者同样安全——选择无节制者与选择节制者——然而并非同等高尚体面。那纵情于逸乐者，讨好肉体；而节制者释放灵魂——肉身的主宰——脱离扰动。若他们说我们\BibleQuote{蒙召得自由，只是不可将自由当作放纵情欲的机会}\BibleRef{迦 5:13}，依宗徒之意。如果必须顺从欲望，且可耻卑劣的生活如同无分别可以选取（如他们所说），那么要么完全顺从欲望——若如此，则须行一切最无耻最邪恶之事，跟随那劝诱我们者；要么有些欲望应回避，便不再无分别地生活，也不应无耻地服侍我们最低贱的部分——肚腹与私处——为欲望所驱使讨好自己的尸身。欲望因享乐得滋养而存活；反之，若受阻挠干扰，则萎靡。那被肉体享乐征服的人，怎能相似于主，或认识天主呢？一切享乐的开端是欲望；欲望是因缺乏而渴望某物带来的烦恼与焦虑。因此，我看那些接受这种生活方式的人，无非正如所言，\par

\Quote{在耻辱之外感受痛苦；}\par

他们为自己招致祸患，现在和将来都选择之。因此，若\BibleQuote{凡事都可行}，且不必因恶行而害怕失去希望，他们或许有借口过卑劣可怜的生活。然而，幸福生活已藉诫命向我们显明，所有人都当跟随它，不曲解任何所说之言，不轻忽任何适宜之事（即使最微小的），跟随圣言的引导；因为若我们偏离了它，必然陷入不朽之恶；但若我们遵循圣经（信者藉之进入），为尽可能地相似于主，便不应无分别地生活，而应尽力远离享乐与欲望，并关心灵魂，在唯独天主那里持守。因为洁净无瑕的心灵，在某种程度上适于接受天主的权能，因天主的肖像在其中升起：\BibleQuote{凡向主怀着这希望的，就洁净自己，如同他洁净一样}\BibleRef{若一 3:3}。那些仍被情感牵引的人，绝不可能获得对天主的认识；因此也不能达到终向，因他们毫无对天主的认识。不认识天主似乎可归咎于那未达终向者；而生活方式导致人不认识天主。因为绝不可能有人既有知识，又不以谄媚肉体为羞。享乐是善，与唯独美善是善，或与唯独美丽的主、唯独善的天主、唯独可受爱慕者，永远不能相容。\BibleQuote{你们在基督内受了割损，不是人手所行的割损，而是脱去肉欲之身的基督割损}\BibleRef{哥 2:11}。\BibleQuote{因此，你们既然与基督一同复活了，就该追求天上的事，思念天上的事，而不是地上的事；因为你们已经死了，你们的生命与基督一同藏在天主内；}而非他们所行的淫乱。\BibleQuote{所以要致死属于地上的肢体，即淫乱、污秽、邪情、恶欲，因这些事天主的义怒才来到。}他们也担当了忿怒、恼怒、恶毒、毁谤、口中污秽之言语，\BibleQuote{脱去了旧人和他的行为，穿上了新人，这新人按创造者的肖像而更新，进入真知}\BibleRef{哥 3:4}。因为生活方式明显揭露那些知道诫命的人：言语如何，生活也如何。树由果实认识，而非由花、叶和枝。因此知识由果实和生活方式而来，而非由言语和花。我们不说赤裸的言语就是知识，而是某种神性的科学，那光因遵守诫命而生在灵魂内，显明一切由生成而来的事物，教导人认识自己，并教导如何达成。因为知识之于心灵，如同眼睛之于身体。他们也不可称那服侍享乐的自由为自由，如同那些说胆汁甜的人。我们学到，我们的主使我们从享乐、欲望和其他扰动中解脱的自由才是自由。\BibleQuote{那说\InnerQuote{我认识主}，却不遵守他诫命的，是说谎的，真理不在他内}\BibleRef{若一 2:4}，若望说。\par

% structure: p0046 source=h2 target=subsection reason=relative-heading-rank; base=h2

\subsection{第六章——攻击第二类异端者：即那些因对万有创造者天主的亵渎看法而实行节制的异端者。}

反之，对另一类异端者，他们外表借着\OriginalTerm{节制}{continentiam}而行不虔敬之事，针对受造物和圣造物主——唯一全能的天主——他们说不应准许婚姻和生养子女，也不应把其他不幸的人带到世上，也不应为死亡提供滋养，对此必须提出反对：首先，若望的话：\BibleQuote{现今有许多假基督出现了，由此我们知道这是最后的时刻。他们从我们中间出去，却不属于我们；如果他们属于我们，必会存留在我们中间。}\BibleRef{若一 2:18}其次，也要驳斥并推翻他们所提出的，即：\Quote{\Person{撒罗默}问：死亡要存续到何时？}——并非因为生命是恶的，受造物是恶的——\Quote{主说：直到你们女人生育为止。}而是教导自然的连续：因为出生后必然跟随死亡。律法确实要使我们远离逸乐和一切羞辱与恶行。其目的就是引导我们从不义走向正义，选择正当的婚姻、生养子女和良好的生活方式。主却\Quote{不是来废除律法，而是来成全}\BibleRef{玛 5:17}。成全，并非由于律法有所欠缺，而是因为律法和先知的预言通过祂的来临得以完成。因为正确的生活方式，甚至对律法之前生活正义的人，也是通过\OriginalTerm{圣言}{Logon}预言的。因此，那些不认识节制的大众，依肉身生活，而非依圣神；没有圣神的肉身不过是尘土和灰烬。主从思想判断奸淫。因为难道不能节制地使用婚姻，而不试图拆散\Quote{天主所结合的}\BibleRef{玛 19:6}吗？离婚的鼓吹者正是教导这些，因此他们的名字在外邦人中受到辱骂和诅咒。而这些人说\OriginalTerm{交合}{coitum}是邪恶的，但他们自己也从交合中获得他们的本质，他们怎么不是邪恶的呢？至于那些被圣化的人，我认为他们的\OriginalTerm{种子}{semen}也是圣的。我们应当被圣化的，不仅是灵魂，还有行为、生活以及身体。因为宗徒\Person{保禄}说\Quote{不信主的丈夫因妻子成了圣洁的}，或者\Quote{妻子因丈夫成了圣洁的}\BibleRef{格前 7:14}，这是什么意思呢？主对那些询问离婚问题的人所说的话又是什么意思：\Quote{是否可以休妻，既然\Person{梅瑟}允许了？}\Quote{祂说：\Person{梅瑟}因你们的心硬才写了这诫命。你们没有读过天主对\OriginalTerm{原祖}{protoplasto}说的话吗？\InnerQuote{二人成为一体}？所以，凡休妻的，除非因奸淫的缘故，就是使她犯奸淫。}\BibleRef{玛 19:3}\Quote{但在复活后，也不娶也不嫁。}关于肚腹和食物也说：\Quote{食物是为肚腹，肚腹是为食物；但天主将使两者都消灭。}\BibleRef{格前 6:13}这是针对那些认为应当像公山羊和公羊一样生活的人，好使他们无忧无虑、毫不畏惧地吃喝和交合。\par

如果他们接受复活（如他们所说），因而削弱并废除婚姻，他们就不吃也不喝；因为宗徒在复活的事上说：\BibleQuote{毁灭了肚腹和食物。} 那么，他们怎么会饥饿、口渴、经历肉体的情感以及其他事情，而一个藉着基督接受了所希望的完美复活的人，是不会经历这些的？就连那些敬拜偶像的人，也禁绝食物和情欲。然而，他说：\BibleQuote{天主的国不在于吃喝。} \BibleRef{罗 14:17} 当然，那些敬拜天使和魔鬼的人更关心同时戒绝酒、肉和情欲之事。正如谦卑是温良，而不是身体的苦行；同样，节制是灵魂的德行，不在于外表，而在于内心。但也有一些人公开说婚姻就是淫乱，并断定它是魔鬼传授的。这些自夸的吹嘘者说他们效法主，主既没有娶妻，也没有拥有世间任何东西；他们夸耀自己比别人更理解福音。但圣经对他们说：\BibleQuote{天主拒绝骄傲的人，却赐恩宠给谦卑的人。} \BibleRef{雅 4:6} 其次，他们不知道主为什么不娶妻。首先，祂有自己的净配教会；其次，祂不是普通人，需要按肉身有一个助手；祂也不需要生儿育女，因为祂永远存在，并且是唯一的天主子。这位主自己说：\BibleQuote{天主所结合的，人不可分开。} \BibleRef{玛 19:6} 又说：\BibleQuote{就如诺厄的日子怎样，人子来临也要怎样：那时人们婚嫁、建造、种植；又如在罗特的日子怎样。} \BibleRef{玛 24:37} 祂表明这话不是对万民说的，因为祂接着说：\BibleQuote{人子来临时，能在世上找到信德吗？} \BibleRef{路 18:8} 又说：\BibleQuote{那日子怀孕的和哺乳的有祸了。} 虽然这些话也是以寓意方式说的。因此，祂也没有确定那些'时候和日期'，\BibleQuote{就是父凭自己的权柄所定的，} \BibleRef{宗 1:7} 为使世界藉着世代延续。他们不明白那句：\BibleQuote{不是所有人都能领悟这话：有些阉人，生来就是如此；有些阉人，被人阉割；有些阉人，为了天国而自阉。能领悟的，就领悟吧。} \BibleRef{玛 19:11} 他们不知道，在论及离婚后，有人问：\BibleQuote{如果夫妻的情形是这样，倒不如不娶。} 那时主说：\BibleQuote{不是所有人都能领悟这话，只有那些得了恩赐的人才能。} \BibleRef{玛 19:10} 因为那些问的人，想要从祂那里知道，如果妻子因淫乱被定罪和休弃，是否允许另娶一个。他们说不少运动员也为了训练身体而节制情欲：例如克罗顿的\Person{阿斯堤卢斯}和希梅拉的\Person{克里松}。还有\Person{阿莫埃布斯}，一位琴师，刚结婚后也抑制自己远离新娘；昔兰尼的\Person{亚里士多德}独自藐视他的情人\Person{莱丝}。他曾向一位妓女发誓，如果她在与对手争斗时帮助他，他就带她回故乡；在实现后，他巧妙地履行了誓言，让人画了一幅她的肖像，放在昔兰尼，正如\Person{伊斯特}在\WorkTitle{De proprietate certaminum}一书中所写。所以，贞洁也并非善，除非是为了爱天主而做的。关于那些厌恶婚姻的人，真福\Person{保禄}说：\BibleQuote{在最后的日子里，有些人将背弃信德，听从欺诈的神和魔鬼的教义，禁止婚嫁，禁戒食物。} 他又说：\BibleQuote{不要让任何人藉着自愿的谦卑和苦身克己来欺骗你们。} 同样，他也写了：\BibleQuote{你已娶了妻子吗？不要寻求分离。你没有妻子吗？不要寻求妻子。} \BibleRef{格前 7:27} 又说：\BibleQuote{每人要拥有自己的妻子，免得撒殚诱惑你们。} 怎么样？古代的义人不是也以感恩之心享用受造物吗？有些人甚至在婚姻中保持了节制并养育了子女。乌鸦给\Person{厄里亚}带来食物，饼和肉。甚至\Person{撒慕尔}先知也把吃剩的羊腿给\Person{撒乌耳}吃。但这些自称在生活方式上卓越的人，甚至不能与他们的行为相比。因此，\BibleQuote{那不吃的人，不要轻视吃的人；那吃的人，也不要判断不吃的人，因为天主已经接纳了他。} \BibleRef{罗 14:3} 甚至主论到自己说：\BibleQuote{若翰来了，不吃不喝，他们就说他附了魔；人子来了，也吃也喝，他们就说：看，一个贪吃好酒的人，税吏和罪人的朋友。} \BibleRef{玛 11:18} 难道他们也谴责宗徒们吗？\Person{伯多禄}和\Person{斐理伯}都生了子女：\Person{斐理伯}还把女儿们嫁人。\Person{保禄}也毫不羞愧地在某封书信中称一位妇女为他的\BibleQuote{配偶}，但他没有带她同行，因为他不需要这样的服务。所以他在一封信中说：\BibleQuote{难道我们没有权力带一位姐妹、妻子同行，如同其余的宗徒一样吗？} \BibleRef{格前 9:5} 但他们，正如合理的，专注于不可分离的职务，即宣讲，所带的妇女不是作为妻子，而是作为姐妹，她们要一起服务那些管家务的妇女；藉着她们，主的教导也能进入内室，没有任何指责或恶意猜疑。因为我们知道，最卓越的\Person{保禄}在给\Person{弟茂德}的第二封信中教导了关于女执事的一切。然而这位\Person{保禄}呼喊说：\BibleQuote{天主的国不在于饮食，}也不是戒酒和肉；\BibleQuote{而在于正义、平安和在圣神内的喜乐。} \BibleRef{罗 14:17} 他们中谁披着羊皮，束着皮带，像\Person{厄里亚}一样周游？谁穿着苦衣，其余赤身，赤脚，像\Person{依撒意亚}？或者只穿着一条麻布腰带，像\Person{耶肋米亚}？谁又能效法\Person{若翰}的诺斯替生活方式？但即使是那样生活，真福先知们也向造物主感恩。至于\Person{卡尔波克拉特}的义德，以及那些像他一样追求无耻共享的人，这样就被驳倒了；因为他说：\BibleQuote{向你求的，就给他；}接着又说：\BibleQuote{愿意向你借贷的，不要拒绝；} \BibleRef{玛 5:42} 这是教导共享，而不是那种乱伦和淫秽的共享。\par

然而，如果没有人拥有并能够出借，那借的和借出的又怎能存在呢？此外，当主说：\BibleQuote{我饿了，你们给了我吃的；我渴了，你们给了我喝的；我作客，你们收留了我；我赤身露体，你们给了我穿的；}\BibleRef{玛 25:35} 随后说：\BibleQuote{你们对我这些最小兄弟中的一个所做的，就是对我做的。}\BibleRef{玛 25:40} 难道在旧约中祂没有颁布同样的法律吗？\BibleQuote{周济穷人的，是借贷给天主。}\BibleRef{箴 19:17} 又说：\BibleQuote{不要拒绝向穷人施恩。}\BibleRef{箴 3:27} 又说：\BibleQuote{不可让仁慈和忠诚离弃你。}\BibleRef{箴 3:3} \BibleQuote{贫穷使人谦卑，富足的手使人致富。}\BibleRef{箴 10:4} 接着又说：\BibleQuote{不放贷取利的人，必蒙悦纳。}以及：\BibleQuote{灵魂的赎价是自己的财富。}\BibleRef{箴 13:8} 这不是清楚表明，正如世界由对立之物组成，即冷与热、湿与干，同样也由施与受的人组成吗？再者，当祂说：\BibleQuote{你若愿意成为成全的，去变卖你所有的，分给穷人，}祂驳斥了那个夸口\BibleQuote{从小遵守了一切诫命}的人；因为他没有做到\BibleQuote{你应爱近人如你自己}；那时，当他被主成全时，他正学习藉着爱德去分享和施与。所以，主并没有正当地禁止人富有，而是禁止人不义和贪得无厌地富有。\BibleQuote{以不义急速获得的财产，必将减少。}\BibleRef{箴 13:11} \BibleQuote{有人播种增多，有人收获反少。}\BibleRef{箴 11:23} 经上关于这些人说：\BibleQuote{他施舍，给了穷人，他的正义永远常存。}因为那\BibleQuote{播种并收获更多的}人，正是藉着尘世和暂时的分享与分配，获取天上的和永恒的事物。另有一种人，他不施与任何人，而枉然\BibleQuote{在地上积蓄财宝，那里有锈和虫腐蚀。}\BibleRef{玛 6:19} 关于此人经上写道：\BibleQuote{那聚集工资的，是将工资装入漏袋。}\BibleRef{盖 1:6} 主在福音中说到这个人的\BibleQuote{田地肥沃}\BibleRef{路 12:16}；然后，当他想要存放果实，并\BibleQuote{建造更大的粮仓}时，他便藉着拟人法对自己说：\BibleQuote{你有许多财物，为你存留多年，吃、喝、享乐吧！}\BibleQuote{愚昧的人哪！今夜就要索回你的灵魂；你所准备的，将归谁呢？}\par

% structure: p0050 source=h2 target=subsection reason=relative-heading-rank; base=h2

\subsection{第七章——基督徒的节制在哪方面胜过哲学家所自称的节制。}

因此，按希腊哲学家的意见，属人的节制——我指的是那种与贪欲争斗，并在行为上不侍奉贪欲的节制。但按我们的意见，节制是不贪恋；不是一个人在贪恋时勇敢行事，而是连贪恋本身也加以克制。若非天主的恩宠，这种节制无法获得。为此祂说：\BibleQuote{你们求，就给你们。}\BibleRef{玛 7:7}梅瑟也领受了这恩宠，他身穿粗糙的躯体，却能四十天不饥不渴。然而，健康胜过患病时谈论健康；同样，光明胜过谈论光明；存在于真理中的节制胜过哲人所教导的节制。因为哪里有光，哪里就没有黑暗；但贪恋独自盘踞时，即使它停止肉体活动，却仍借着记忆与不在场的事物交往。总之，我们关于婚姻、饮食及其它事上的论述，是要我们不做任何出于贪欲的事，只愿那些必要的事。因为我们不是贪欲之子，而是意志之子；那为生育子女而娶妻的人，应当操练节制，甚至不贪恋自己的妻子，他应以正直和节制的意愿关爱妻子、致力于子女。因为我们所学的不是\BibleQuote{为肉体安排，去放纵私欲}，而是\BibleQuote{行事为人要端正，好像行在白昼}，跟随基督和主的光明生活规范，\BibleQuote{不可荒宴醉酒，不可好色邪荡，不可争竞嫉妒}。但实在说来，不应只在一种事情上考虑节制，即情欲之事，也应在一切其它事上考虑，那些灵魂贪求奢侈、不满足于必需、而追求享乐的事。节制就是轻视钱财；以伟大崇高的精神藐视逸乐、财产、观赏；管住口舌，以理性战胜邪恶。事实上，有些天使也曾因不节制、被贪欲征服而从天上坠落到此。瓦伦提诺在致阿加托波德的书信中写道：\Quote{耶稣忍受到一切，祂是节制的，为自己挣得了神性；祂以独特的方式吃喝，却不消化食物；祂内里存有如此大的节制之力，以致营养在祂内也不朽坏，因为祂本身没有朽坏。}因此，我们出于对主的爱，并为了正直本身，拥抱节制，圣化圣神的宫殿。因为正直就是\Quote{为天国而自阉}\BibleRef{玛 19:12}，断绝一切贪欲，并\Quote{洁净良心，除去死行，以事奉永生天主}\BibleRef{希 9:14}。而那些出于对肉体的憎恨而渴望摆脱婚姻结合和正当饮食的人，是无知、不敬且不合理的节制，正如许多其他民族一样。婆罗门确实不吃有生命之物，也不饮酒；但他们中有些人每日如我们一样进食；另有一些人则每三天进食一次，如亚历山大·波利希斯托在\WorkTitle{印度志}中所说的；他们轻视死亡，视活着为无物，因为他们相信再生；有些人敬奉赫拉克勒斯和潘。印度人中被称为\OriginalTerm{可敬者}{\GreekText{Σεμνοί}}的那些人，即\OriginalTerm{可敬者}{venerandi}，终生赤身露体：他们操练真理，预言未来，敬拜一座金字塔，他们以为某位神的骨头安放在其下。裸行智者和所谓\OriginalTerm{可敬者}{\GreekText{Σεμνοί}}都不与女子交合，因为他们认为这违反自然且不义；因此他们保持自身贞洁。也有些女子是贞女，被称为\OriginalTerm{可敬者}{\GreekText{Σεμναί}}，即\OriginalTerm{可敬者}{venerandæ}。他们似乎观察天象，并借此预兆预言某些未来之事。\par

% structure: p0052 source=h2 target=subsection reason=relative-heading-rank; base=h2

\subsection{第八章——解释异端者为贬抑婚姻而引用的圣经章节；首先驳斥异端者对宗徒在\BibleRef{罗 6:14}的话的曲解。}

既然那些倡导无差别论的人曲解了几处圣经，以为它们支持他们挑逗性的享乐，尤其是这一节：\BibleQuote{罪必不能作你们的主，因为你们不在法律之下，乃在恩宠之下。}\BibleRef{罗 6:14}以及一些其它类似的经文，此后我不宜再提（因我不教海盗行径），那么让我们稍加驳斥他们的论点。因为这位杰出的宗徒在紧接着上述话语的字句中，给出了对控罪的解答：\BibleQuote{这却怎么样呢？我们在恩宠之下，不在法律之下，就可以犯罪吗？断乎不可。}\BibleRef{罗 6:15}他如此神圣而先知般地随即瓦解了享乐的诡辩之术。因此，看来他们不明白：\BibleQuote{我们众人必须出现在基督的审判台前，好使每人借着身体，照他所做的，或善或恶，领取应得的报应}\BibleRef{格后 5:10}，即每人通过身体所做的，他要领受。\Quote{所以谁若在基督内，他就是新受造物}，不再沉溺于罪恶：\Quote{旧事已过}，我们脱去了旧生活：\BibleQuote{看，一切都成了新的}\BibleRef{格后 5:16}，从淫乱到贞洁，从不节制到节制，从不义到正义。\BibleQuote{正义与不义有什么相通？光明与黑暗有什么联系？基督与贝里雅耳有什么和谐？信者与不信者有什么相干？天主的圣殿与偶像有什么一致？所以，我们既有这些恩许，就当洁净自己，除去肉体和心灵的一切玷污，以敬畏天主而成全圣德。}\BibleRef{格后 7:1}\par

% structure: p0054 source=h2 target=subsection reason=relative-heading-rank; base=h2

\subsection{第九章——解释基督对撒罗默所说的话，异端者将其作为贬抑婚姻的言论而引用。}

但是那些借着那美好的节制而抵抗天主的受造物的人，也引用那些对\Person{撒罗默}所说的话，我们之前曾提及；这些话，按我看来，是记载在\WorkTitle{埃及人福音}中的。因为他们说救主亲自说过：\Quote{我来是为瓦解妇人的工作。}妇人，指的是情欲；工作，则指生育和朽坏。那么他们会说什么呢？难道这事工停止了吗？他们不会这么说：因为世界仍然处于同一治理之下。但主并没有说谎；他实际上瓦解了情欲的工作：贪婪、纷争、虚荣、对女人疯狂的爱、娈童、暴食、奢侈和挥霍，以及诸如此类的事。这些事的起源就是灵魂的死亡：因为我们确实\BibleQuote{死于过犯}。\BibleRef{弗 2:5}而那个妇人就是不节制。受造物的产生和朽坏因它们的本性而必然发生，直到被选者的完全分别和恢复，借此那些曾与世俗混合和混乱的本质得以恢复其本有特性。因此，当\OriginalTerm{圣言}{Logos}曾谈论终结时，撒罗默问道：\Quote{人要到几时才死亡呢？}圣经用双重意义称人：既指那显现的人，也指灵魂；又指那得救的人和那不得救的人。而灵魂的死亡被称为罪。因此主谨慎而深思熟虑地回答：\Quote{直到妇女生育之时，}即直到情欲活动为止。宗徒说：\BibleQuote{因此，就如罪恶藉着一人进入了世界，死亡藉着罪恶也进入了世界，如此死亡便殃及了众人，因为众人都犯了罪；死亡从亚当直到梅瑟统治了人。}\BibleRef{罗 5:12}然而，按照天主圣意治理的自然必然性，死亡紧随生育而来；身体的结合之后就是它们的分离。如果生育是为了教训和认识，那么分离将是恢复的原因。正如妇人因生育而被视为死亡的原因，同样地，因同一原因她也将被称为生命的引领者。因此，那个最先开始违命的女人被称为\BibleQuote{生命}\BibleRef{创 3:20}，因为承继的缘故：她是一切受生者和犯罪者、义人和不义之人的母亲，因为我们每个人称自己为义，或相反使自己成为悖逆者。因此，我不认为宗徒厌恶肉身的生命，当他说：\BibleQuote{但无论在什么环境中，我满怀信心，现今和往常一样，基督必在我的身体上受显扬，或借生命，或借死亡。因为在我看来，生活就是基督，死亡就是利益。但如果活在肉身内，这对我也是工作的果实，那么我该选择什么，我不知道。我被夹在两者之间：我渴望解脱而与基督同在，那实在好得多；但为你们，我留在肉身内更为必要。}\BibleRef{斐 1:20}通过这些，我以为他清楚地表明：离开身体进入完美，就是对天主的爱；而他心存感恩地容忍留在肉身的现在，是为了那些需要救恩的人。不仅如此，那些宁愿追随任何事物而不愿追随真理的福音准则的人，不是也在关于对撒罗默所说的话之后，附加了以下内容吗？因为当她说：\Quote{那么我没有生育是对的。}意即生育没有被恰当地接受；主回答说：\Quote{你吃各种蔬菜，但不要吃有苦味的。}借此他表明，节制或婚姻都取决于我们的能力，而不是出于诫命的禁止；并且婚姻给受造界带来一些帮助，此外他还加以解释。因此，没有人应认为按圣言结婚的人犯了罪，除非他认为养育儿女是苦事；然而，对许多人来说，没有孩子却是极其烦恼的。也不应有人以为生育儿女是苦事，因为它使人陷入事务，从神圣事物中分心。因为有些人无法轻易忍受独身生活，就寻求婚姻：毕竟婚姻是一件愉快的事，人可节制地享受且无害；而且我们每个人都如此自主，可以选择是否愿意生育儿女。但我明白，有些人以婚姻的困难为借口而禁绝婚姻，却不合神圣知识而堕入了非人道和憎恨人类；在他们中间爱德消失了；另一些人则结婚，却沉迷于奢侈和享乐，法律在某种程度上伴随他们，他们就如先知所说，\Quote{成了像畜牲一样。}\par

% structure: p0056 source=h2 target=subsection reason=relative-heading-rank; base=h2

\subsection{第十章 基督的话\BibleRef{玛 18:20}的神秘解释}

那\BibleQuote{两个或三个人，因我的名聚在一起}，在他们\BibleQuote{中间}有主的是谁呢？\BibleRef{玛 18:20}难道他不是说男人、女人和孩子是三个人，因为女人与男人因天主结合吗？但如果有人想要束腰并轻装，不愿生儿育女，因为生儿育女有烦恼和牵挂，宗徒说，\BibleQuote{让他没有妻子，如同我一样。}\BibleRef{格前 7:7}但他们解释主的说话，好像他说过：对于多数人，天主是创造者和生育的主宰；但对于一个人，即被选者，救主是另一位神、即善神的儿子。但事实并非如此：其实，对于在婚姻中诚实适度行事并生育子女的人，天主也藉着圣子与他们同在；对于按照\OriginalTerm{理性}{Logon}（即理性）而节制的人，同一位天主也与他同在。此外，还有其他三种：忿怒、贪欲和理性；肉体、灵魂和灵，则是另一种方式。或许，之前提到的三人组也暗示了召叫、第二类选民和第三类被置于首等尊荣的种类，与他们同在的是考虑万物的天主的能力，不分裂却落入分裂之中。因此，凡按应当的方式运用灵魂的自然运作的人，确实渴望适宜之事，恨恶有害之事，正如诫命所命：\BibleQuote{你要祝福那祝福你的，诅咒那诅咒你的。}但当人超越了忿怒和贪欲，因着万有的创造者天主而对受造物怀有真正的情爱，以\OriginalTerm{诺斯替方式}{gnostice}安排生活，变得相似救主，获得了容易的节制习惯，并将认知、信德和爱德结合起来，在这部分使用简单的判断，真正成为属神的人，完全不能容纳从忿怒和贪欲产生的意念，如此，被造者亲自塑造的人成为完美的人，他确实配得被称为主的兄弟，同时是朋友和儿子。因此，\BibleQuote{两个或三个人}聚集在同一个地方，即在诺斯替者身上。由三个人所估算的多人的和谐，主与他们同在，也可以象征一个教会、一个人、一个民族。难道不是这样吗：当主颁布律法时，他与一个犹太人同在；但当他藉预言并派遣\Person{耶肋米亚}到\Place{巴比伦}，又藉预言召叫外邦人时，他聚集了两个民族；第三个是一个，由两个\BibleQuote{被造成一个新的人}，主在其中行走并居住于教会本身？律法和先知，连同\BibleBook{}{福音}，都因基督的名聚集在同一个认知中。因此，那些因憎恨而不娶妻，或出于贪欲而无分别地滥用肉身的人，不在那些被拯救、主与他们同在的人之列。\par

% structure: p0058 source=h2 target=subsection reason=relative-heading-rank; base=h2

\subsection{第十一章。——解释法律和基督关于不贪恋的诫命。}

既然这些已经显明，来吧，让我们引用那些反对异端诡辩者的圣经，并阐明按照理性应遵守的节欲法则。凡能理解哪段圣经与哪个异端相反的人，便适时运用它来反驳那些捏造与诫命相悖教义的人。此外，为从高处重新回顾，法律正如我们先前所说，先宣告了：\BibleQuote{不可贪图你近人的妻子}\BibleRef{出 20:17}；随后又在\WorkTitle{新约}中与主的声音相连，祂亲自说：\BibleQuote{你们听见了法律吩咐：不可奸淫。我却对你们说：不可贪恋。}\BibleRef{玛 5:27}。法律愿意男子节制地使用妻子，仅为生子，这由以下事实清楚表明：它禁止无妻之人立刻与女俘交合；若他一心想娶，便允许她剃发，哀哭三十天；如果即便如此情欲仍不消减，那时才可生儿育女，因为支配的冲动在指定时间后已被证明是合情理的。因此，你找不到任何古人在圣经中与孕妇交合，而是等怀胎、生产、断奶后，才再与妻同房。你会发现梅瑟的父亲遵守了这一目的和规则：在生下亚郎后隔了三年才生了梅瑟。同样，肋未支派遵守了天主所传的这条自然法律，以少于其他支派的人数进入了预许之地；因为当八人依法婚姻播种种时，不易繁衍，他们不仅等待怀胎，还等待断奶。因此梅瑟也循序渐进地引导犹太人走向节欲，当他们连续三天后禁绝性欲时，命令他们听天主的话：\BibleQuote{三天}\BibleRef{出 19:20}。\BibleQuote{我们原是天主的殿，正如先知所说：我要在他们内居住，在他们内往来；我要作他们的天主，他们要作我的子民。}若我们按诫命生活，无论个人还是整个教会。\BibleQuote{因此，你们要从他们中间出来，分别出来，主说：不要触摸不洁之物；我就接纳你们，我要作你们的父亲，你们要作我的子女，全能的主说。}这先知性的话命令我们分离，不是如他们所说从娶妻者中分离，而是从仍生活在淫乱中的外邦人中，也从先前所说的异端中，如同不洁不敬者。为此，保禄也向那些与所说之人类似的人发言：\BibleQuote{所以，亲爱的，我们既有这些恩许，就当洁净自己，除去肉体和心灵的一切玷污，在敬畏天主中成全圣德。}\BibleRef{格后 7:1}\BibleQuote{我以天主的妒爱妒爱你们，因为我曾把你们许配给一个丈夫，如同把贞洁的童女献给基督。}\BibleRef{格后 11:2}教会不再与其他婚姻结合，因为她已有新郎；但我们每人都有权力娶任何合法的女子，尤其是在初婚中。\BibleQuote{但我害怕，正如蛇以诡诈诱惑了厄娃，你们的心志也会从基督内的纯朴中败坏。}\BibleRef{格后 11:3}宗徒以敬虔和导师的风范说了这话。因此，可敬的伯多禄也说：\BibleQuote{亲爱的，我劝你们作客旅和寄居的，要戒绝肉体情欲，这些情欲与灵魂争战；在外邦人中要品行端正，好使那些毁谤你们是作恶的人，因见到你们的善行，归光荣于天主。你们是自由的人，却不可用自由作邪恶的掩饰，但要作天主的仆人。}同样，保禄在\WorkTitle{罗马书}中也写道：\BibleQuote{我们已死于罪恶，怎能还在罪恶中活着？因为我们旧人已与基督同钉十字架，使罪恶的身体被毁坏，}直到那话：\BibleQuote{也不要将你们的肢体献给罪恶作不义的武器。}\BibleRef{罗 6:13}既然到了这里，我似乎不该忽略指出：宗徒通过法律、先知和\WorkTitle{福音}宣告了同一天主。因为在\WorkTitle{罗马书}中，他把写在\WorkTitle{福音}中的\BibleQuote{不可贪恋}归于法律，知道通过法律和先知宣告的是同一位父，也是通过祂所宣告的。他说：\BibleQuote{那么，我们说什么呢？法律是罪吗？绝对不是！但我若不藉法律，就不认识罪；因为若不是法律说\InnerQuote{不可贪恋}，我就不知道贪恋。}\BibleRef{罗 7:7}如果持不同意见者反对，认为保禄的话是针对造物主，那么请他们读接下来所说的：\BibleQuote{我知道，善不在我内，即不在我肉体内，}\BibleRef{罗 7:18}并读此前和此后的话。因为先说了：\BibleQuote{而是住在我内的罪；}因此按理说：\BibleQuote{善不在我肉体内。}\BibleRef{罗 7:17}随后接着说：\BibleQuote{但如果我所作的，是我所不愿意的，那么作那事的已不是我，而是住在我内的罪。}那\BibleQuote{敌对}天主的法律和\BibleQuote{我的理智}，\BibleQuote{把我掳去，使我隶属于那在我肢体中作王的罪的法律。我真不幸！谁能救我脱离这死亡的身体？}他永不厌倦地用各种方式帮助，毫不畏惧地总结：\BibleQuote{因为圣神内的法律已使我脱离罪和死的法律了；}因为\BibleQuote{天主藉着圣子，在肉体中定了罪案，使法律的要求成就在我们身上，因为我们不随从肉体，而随从圣神。}此外，他进一步阐明先前的话，喊道：\BibleQuote{身体因罪而死，}指明它不是圣殿，而是灵魂的坟墓。当身体被圣化归于天主时，他说：\BibleQuote{那使耶稣从死者中复活的圣神，如果住在你们内，祂必使你们会死的身体藉着住在你们内的圣神而活。}\BibleRef{罗 8:10}他又责备耽于享乐的人，附加说：\BibleQuote{随从肉体的心思是死亡；因为随从肉体的人，思念肉体的事；随从肉体的心思是与天主为敌，因为它不服从天主的法律。凡属于肉体的人，}不像某些人所说，\BibleQuote{不能中悦天主，}而是如我们先前所说。然后为区分他们，他对教会说：\BibleQuote{但你们已不属于肉体，而是属于圣神，只要天主的圣神住在你们内。谁若没有基督的圣神，就不属于基督。如果基督在你们内，身体固然因罪而死，精神却因正义而活。所以我们弟兄们，并不欠肉体的债，而随从肉体生活。如果你们随从肉体生活，必要死；但如果藉圣神，治死肉体的行为，必能活。因为凡被天主的圣神引导的，都是天主的子女。}反对可憎的贵族和自由（由持不同见解者引入，以淫欲为荣），他补充说：\BibleQuote{因为你们没有领受奴役的圣神，再陷于恐惧；而是领受了义子的圣神，使我们呼号\InnerQuote{阿爸，父！}}即为此我们领受了圣神，为认识我们所祈求的那一位，祂是真正的父，万有的独一之父，祂如父亲般为救恩教导和惩戒，并施加恐惧的威胁。\par

% structure: p0060 source=h2 target=subsection reason=relative-heading-rank; base=h2

\subsection{第十二章——解释宗徒的话\BibleRef{格前 7:5,39-40}以及其他相关的圣经章节。}

但是，\BibleQuote{经双方同意，暂离是为了专务祈祷}的婚姻，是节制的教导。他加上了\BibleQuote{经双方同意}，以免有人拆散婚姻；\BibleQuote{但只在暂时间内}，\BibleRef{格前7:5}以免娶了妻子的人因不得已而操练节制而滑入罪恶，并且因顾惜自己的配偶而贪恋他人的。因此，他认为那因供养处女而举止不当的人，正确地说这婚姻是有损害的婚姻。然而，每个人——无论是选择贞洁的，还是为了生育子女而结婚的——都必须坚定地持守自己的志向，不可转向更坏的。因为如果能增进并加强自己生活的目标，他就在天主面前获得更大的尊荣，由于纯洁而合理的节制。但如果他超越了自己选择的规则，他就会坠入更大的、关乎希望的荣耀。因为正如贞洁有其职分，婚姻也有其属于主的特殊职务和侍奉，即照顾子女和妻子。因为那在婚姻中成全的人所正当地主张的，是婚姻的亲密关系，使他在共同的家中表现出对一切的关心和照顾。因此，他说，应当设立\BibleQuote{主教}，他们从自己的家操练来管理整个教会。所以，\BibleQuote{每个人}当在\BibleQuote{他所蒙召的事上} \BibleRef{格前7:24}以工作完成服务，好使在基督内成为自由，并得到他服务应得的报酬。他接着论述法律，使用比喻说：\BibleQuote{因为有了丈夫的妇人，}他说，\BibleQuote{当丈夫活着时，她被法律束缚，} \BibleRef{罗8:2}等等。他又说：\BibleQuote{妇人被束缚，只要她的丈夫活着；但如果丈夫死了，她就自由了，可以再嫁，但只可在主内。但如果她这样守下去，就更蒙福——这是我的意见。} \BibleRef{格前7:39}但在前一段中，他说：\BibleQuote{你们已被杀死，} \BibleQuote{脱离了法律，}不是婚姻，\BibleQuote{好使你们归属另一位，就是从死者中复活的那一位，} \BibleRef{罗7:4}就是新娘和教会；她必须保持贞洁，远离内在的、与真理为敌的思虑，以及外在试探她的人，即那些追随异端的人，他们说服你们离开独一丈夫，即全能天主：\BibleQuote{免得像蛇诱惑了厄娃，} \BibleRef{格后11:3}那被称为\BibleQuote{生命}的，我们也被异端狡猾的诱惑引入歧途，违犯诫命。第二段则确立了一夫一妻制：因为，不像某些人所认为的，不应理解为妇人受丈夫束缚、肉体受朽坏联结；他是在指责那些不敬的人，他们公然将婚姻的发明归于魔鬼，因此立法者面临被咒骂的危险。我认为\Place{叙利亚}人\Person{塔提安}敢于教导这样的道理。他在\WorkTitle{依照救世主的完善}一书中确实写道：同意确实结合于祈祷；但朽坏的共融却解除了祈祷的效力。他确实非常谨慎地通过让步加以限制。因为当他再次允许\BibleQuote{为了撒殚和不节制而再次相聚} \BibleRef{格前7:5}时，他宣布那将要服从的人\BibleQuote{将事奉两个主人：} \BibleRef{玛6:24}通过同意事奉天主，通过不同意事奉不节制、奸淫和魔鬼。他说这些话是在解释宗徒。但他以诡辩歪曲真理，用真事来确证虚假：因为我们也承认不节制和奸淫是魔鬼的恶习和情欲；但是，节制婚姻的同意介于其间，它既引导人持续祈祷，又使人为了生育子女而体面地结合。\BibleQuote{认识}一词确实被圣经用来指生育子女的时期，当它说：\BibleQuote{亚当认识了自己的妻子厄娃，她怀了孕，生了儿子，给他起名叫舍特：因为天主给我立了另一个儿子替代亚伯尔。} \BibleRef{创4:25}你看见他们咒骂的是什么人？他们攻击诚实而有节制的播种，并将生育归于魔鬼。因为他没有简单地称\BibleQuote{神}，而是加上了定冠词，即\OriginalTerm{天主}{\GreekText{ὁ} \GreekText{Θεός}}，意指那全能者。接着宗徒所说的话：\BibleQuote{你们再相聚是为了撒殚，} \BibleRef{格前7:5}，其目的是为了消除转向其他欲望的机会。因为暂时的同意并没有彻底排斥自然的欲求：通过这同意，宗徒再次引入婚姻的结合，不是为了不节制、奸淫和魔鬼的行为，而是为了不被不节制、奸淫和魔鬼所制服。塔提安也区分旧人和新人，但不像我们所说的：他称\BibleQuote{旧人}为法律，\BibleQuote{新人}为福音。我们同样同意他，但不是像他那样，他因废除法律而视之为另一位天主；而是同一位主和主人，当祂更新旧物时，不再允许多妻制（因为当人们需要增长繁衍时，天主曾需要多妻制），但引入了一夫一妻制，为了生育子女和料理家务，为此女人被赐予作助手；如果宗徒为了不节制的\BibleQuote{火}而允许第二次婚姻的宽免，那么这人在旧约下确实不犯罪（因为法律没有禁止），但他没有达到福音所要求的圆满生活。但那守在自己身边、通过死亡解除而保持纯洁结合的人，为自己赢得了天上的荣耀，并怀着感恩和乐意的心服从那神圣的安排，藉此他不能从主的侍奉中分离。但现今通过主的天主眷顾，并不吩咐那从婚榻起身的人像从前一样被洗涤：因为主通过一次圣洗洗净信徒，使之达到完全的习俗，祂没有必然地阻止生育子女，祂也将梅瑟的多次圣洗包含在一次圣洗中。同样，法律通过肉体的生育预示重生，曾将圣洗应用于生殖能力，而不是因为它厌恶人的生育。因为显现为人所生，这等同于精液的排泄。所以，并非多次性交是生殖的，而是子宫的接受承认了生育，当在自然的工厂中种子形成胚胎。怎么可能是，婚姻只是旧的和法律的发明，而从主而来的婚姻是陌生的，既然我们承认同一天主？因为\BibleQuote{天主所结合的，人}不可\BibleQuote{拆散}；\BibleRef{玛19:6}何况父所命令的，子也会遵守。如果同一位同时是立法者和福音作者，祂绝不会自相矛盾。因为法律存活着，既是属神的，并以真知被人理解：我们却\BibleQuote{藉着基督的身体，已死于法律，好使我们归属于另一位，即从死者中复活的那一位，}就是法律所预言的，\BibleQuote{为使天主结果实。} \BibleRef{罗7:4}因此，\BibleQuote{法律是圣的，诫命也是圣的、正义的、善的。} \BibleRef{罗7:12}所以我们死于法律，即死于法律所意指的罪，法律通过命令当做的事和禁止不当做的事而指出罪，但并不产生罪，斥责内在的罪，\BibleQuote{好显露出罪来。}但如果按照法律缔结的婚姻是罪，我不知道一个人如何能说自己认识天主，竟然说天主的命令是罪。但如果\BibleQuote{法律是圣的}，婚姻就是圣的。因此宗徒将此奥秘联系到基督和教会：就如\BibleQuote{由肉生的属于肉，由神生的属于神，} \BibleRef{若3:6}不仅在于生育，也在于学习。如今\BibleQuote{儿女是圣的}，\BibleRef{格前7:14}是主的话语所喜悦的喜乐，它们与灵魂订婚。所以奸淫和婚姻是分开的，因为魔鬼远离天主。\BibleQuote{因此你们也藉着基督的身体死于法律，好使你们归属于那从死者中复活的一位。} \BibleRef{罗7:4}同时，如果你们顺从，就立即蒙垂听，因为我们同样从法律的真理服从同一主，祂从远处命令。难道关于这样的人，圣神岂不公开\BibleQuote{说：在末后的时期，有些人将离弃信德，听从欺诈的神和魔鬼的教训，这教训出于假冒伪善的撒谎者，他们的良心已被烙伤，他们禁止嫁娶，禁戒食物，就是天主所造，让有信德并认识真理的人以感谢的心领受的。因为天主的一切受造物都是好的，若以感谢的心领受，就没有一样是可舍弃的。因为藉着天主的话和祈祷得以圣化。}所以，绝不应禁止结婚，也不应禁止吃肉或喝酒。因为经上记载：\BibleQuote{不可吃肉，也不可喝酒，凡使你弟兄跌绊的事，一概不作才好。} \BibleRef{罗14:21}又说：\BibleQuote{我愿众人像我一样。} \BibleRef{格前7:8}但使用的人\BibleQuote{以感谢的心}，\BibleRef{罗14:19}而不使用的人也同样\BibleQuote{以感谢的心}，并以节制和适度的享用生活，合乎理性。简言之，所有宗徒的书信，教导节制和自律，同时包含无数关于婚姻、生育子女和家务管理的诫命，从未禁止或废除正当而节制的婚姻；而是保持法律与福音的和谐，两者都接纳：那以感谢天主而节制地使用婚姻的人，和那按照主的意愿生活在贞洁中的人，如同\BibleQuote{各人蒙召的时候}无过犯且完全地选择。\BibleQuote{雅各布的帐幕将超越全地，}先知说，那器皿本身因神的光荣而荣耀。但有人攻击生育，说生育会朽坏、灭亡；又有人将救主的话\BibleQuote{不要在地上积蓄财宝，那里有虫蛀和锈蚀} \BibleRef{玛6:19}曲解为针对生育子女，并且无耻地加上先知的话：\BibleQuote{你们众人如同衣服会变旧，蛀虫将吃掉你们。} \BibleRef{依50:9}但我们并不反对圣经，也不否认我们的身体会朽坏、是易逝的。也许先知在那里对作为罪人的人预言了灭亡。但救主关于生育子女的话什么也没说，而是劝告那些只富有财物却不愿帮助穷人的人要施舍和分享。因此祂说：\BibleQuote{你们不要为那可朽坏的食粮操劳，而要为那存留到永生的食粮操劳。} \BibleRef{若6:27}他们也同样引用关于死人复活的话：\BibleQuote{那世界的儿子们不娶也不嫁。} \BibleRef{路20:35}但若有人思考这个问题和那些提问的人，就会发现主不是谴责婚姻，而是在复活中为肉欲的期待提供补救。那句\BibleQuote{今世的儿子们}，\BibleRef{路20:34}不是说为了区分另一个世界，而是如同说：在这世界出生的人，由于生育而为儿子，他们生养儿女；因为没有生育，没有人能经过此生；但这一生育，接受类似的朽坏，不再属于那与今生分离的人。\BibleQuote{所以，我们只有一位父，就是那在天上的。} \BibleRef{玛23:9}然而祂也是所有受造物的父。\BibleQuote{所以你们不要称地上的任何人为父。} \BibleRef{玛23:9}如同说：不要认为那以肉体撒种生育你们的人是你们本性的作者和原因，而是生育的辅助原因，甚或只是一个工具。因此，祂要我们回转变成小孩子，承认那真正的父，通过水重生，而这是创造中的另一种撒种。但他说：\BibleQuote{那独身的人，挂虑主的事；那娶了妻子的人，想怎样悦乐妻子。} 怎么？难道那些按天主旨意悦乐妻子的人，不能感谢天主吗？难道那娶了妻子的人，不能同时与婚姻一起挂虑主的事吗？但是，正如\BibleQuote{那未出嫁的，挂虑主的事，一心使身体和灵魂圣洁；}同样，那出嫁的，也挂虑丈夫的事和主的事，在主内挂虑，使身体和灵魂圣洁。因为两者在主内都是圣的：一个作为妻子，另一个作为处女。为了羞辱和制止那些倾向再婚的人，宗徒特意用严厉的口吻说话：他说：\BibleQuote{看，各种罪都在身体之外，惟有行淫的人，是得罪自己的身体。} \BibleRef{格前6:18}如果有人胆敢说婚姻是奸淫，那他就是再次攻击法律和主，并加以诅咒。因为正如贪婪和贪多被称为奸淫，因为它与知足相悖；又如偶像崇拜是将独一的天主分散到众多，同样，奸淫是从一个婚姻滑向多个。因为，如我们所说，在宗徒那里，奸淫和通奸有三种含义。论到这些，先知说：\BibleQuote{你们因自己的罪恶被出卖。}又说：\BibleQuote{你在异地受了玷污：} \BibleRef{依50:1}认为那与外人身体而不是与婚姻中为生育子女所给的身体的结合是邪恶的。因此宗徒也说：\BibleQuote{我愿意年轻的寡妇再嫁，生养儿女，治理家务，不给敌人以辱骂的机会。因为有些已转去随从撒殚。} \BibleRef{弟前5:14}而且他确实接受一妻的丈夫；无论是司铎、执事或平信徒，只要使用婚姻无可指摘：\BibleQuote{但他将因生育子女而得救。} \BibleRef{弟前3:15}再者，救主称犹太人为\BibleQuote{邪恶淫乱的世代}，教导他们不认识法律，却\BibleQuote{随从长老的传授和人的吩咐，} \BibleQuote{玷污了法律，}好像独身的主和主人没有赐给法律一样。也许祂也暗示他们被外来的欲望所奴役，因不断服侍罪恶而被卖给外邦人。因为在犹太人中间，公共妇女是不允许的：相反，通奸是被禁止的。那说\BibleQuote{我娶了妻子，不能来} \BibleRef{路19:20}赴天主的筵席的，便是一个例子，用以指责那些因享乐而背离神圣命令的人；否则，在来临之前和之后娶妻的义人，即使是宗徒，也不能得救。如果他们引用先知的话：\BibleQuote{我在一切仇敌中间衰老了，}他们应将仇敌理解为罪恶。但有一种罪，不是婚姻，而是奸淫；否则他们也会说生育是罪，并且生育的创造主也是罪。\par

% structure: p0062 source=h2 target=subsection reason=relative-heading-rank; base=h2

\subsection{第十三章——回应异端者尤利乌斯·卡西安的言论；以及同一人所引伪福音之文。}

\Person{尤利乌斯·卡西安}，幻影论派的首领，也使用这样的论点。在其著作\WorkTitle{论节欲}或\WorkTitle{论贞洁}中，他这样说道：\Quote{没有人应该说，因为我们有这样的肢体，女性以不同方式被塑造，男性又以另一种方式：女性是为了接纳，男性是为了播种，因此天主允许了这种交合。因为如果此构造来自我们所趋向的天主，祂就不会称阉人是有福的；先知也不会说他们\InnerQuote{不是不结果子的树}\BibleRef{依 56:3}，将树喻为人，那人出于自愿和自主选择了如此自阉的意念。}为了支持其不虔敬的意见，他进一步争辩道：\Quote{然而，人岂非能正当地责备救主，如果祂改变了我们，将我们从错误中解放出来，并脱离了肢体、附加物和羞耻部位的联系？}在此他与\Person{塔提安}意见相同；后者出自\Person{瓦伦廷}学派。因此卡西安说：\Quote{当\Person{撒罗默}询问那些她所问之事何时会显明时，主说：当你们践踏了羞耻的衣饰，当二者成为一体，且男性与女性既非男性亦非女性之时。}首先，在我们所传的四个福音书中并没有这句话，而是在那部所谓\WorkTitle{埃及福音}中。其次，在我看来，他似乎不知道：愤怒指的是男性的欲求，而女性指的是贪欲；对其运作的后果是悔改和羞耻。因此，当一个人既不顺从愤怒也不顺从贪欲（贪欲因习惯和不良教育而增强，遮蔽并掩盖理性），而是除去由此产生的黑暗，因悔改而感到羞耻，使灵魂与精神在服从\OriginalTerm{圣言}{Logos}或理性中合一；那么，如\Person{保禄}所说，\Quote{在我们内没有男性和女性之分}。因为灵魂脱离那区分男性和女性的形态，便被引向合一，因为那时候它两者都不是。然而，这位如此著名的人过多地按柏拉图主义认为，灵魂原本是神圣的，被欲望软化，来到此地进入生成和腐朽。\par

% structure: p0064 source=h2 target=subsection reason=relative-heading-rank; base=h2

\subsection{第十四章——解释\BibleRef{格后 11:3}和\BibleRef{弗 4:24}。}

事实上，他几乎勉强地迫使\Person{保禄}将生育归于欺骗，当他说：\BibleQuote{但我怕你们的心意，像蛇曾用诡诈迷惑了厄娃一样，会从基督的单纯中败坏。}\BibleRef{格后 11:3}然而，主确实也\BibleQuote{来}到那些\BibleQuote{迷失了}的人那里\BibleRef{玛 18:11}。但他们并不是从原先更高来源坠入此地生成（因为生成是全能的造物主的创造，祂从未将灵魂从更好的带向更坏的）；而是对那些在思想或意念上迷失了的人，救主来到我们这里——这些人在遵守诫命上因我们的不服从而败坏，因为我们过分贪求快乐；尤其我们的原祖提前了时间，在婚期之前就贪求婚姻的恩宠并迷失了：因为\BibleQuote{凡注视妇女而贪恋她的，就已经在心上犯奸淫了}\BibleRef{玛 5:28}，如同那等不及意愿的时间的人。因此，正是主，祂当时也在谴责那先于婚姻的贪欲。所以，当\Person{宗徒}说：\BibleQuote{穿上新人，它是按照天主的肖像创造的}\BibleRef{弗 4:24}，祂是对我们说的，我们是由全能者的意愿所造，正如我们所被造的那样。但祂说\Quote{旧人}，并非指生与重生，而是指不服从的生活与服从的重生。至于\BibleQuote{皮衣}\BibleRef{创 3:21}，\Person{卡西安}认为是指身体：在这一点上，以后我们将证明他以及与他同感的人都错了，当我们关于人的起源，在适当处理了先前必须说明的问题之后，展开阐述。因为他说：\BibleQuote{那些被地上事物支配的，生育并出生；但\Emph{我们的家乡是在天上，我们也从那里等待救主。}\BibleRef{斐 3:20}}因此，我们正确知道这些话也是对我们需要说的：我们应该以客旅和寄居者的身份过生活；有妻子的要像没有妻子；拥有财产的要像没有财产；生育儿女的要像在生必死的人，像要留下产业，甚至必要时过没有妻子的生活；不是带着过分的行为，而是怀着高远的心志。\par

% structure: p0066 source=h2 target=subsection reason=relative-heading-rank; base=h2

\subsection{第十五章——解释\BibleRef{格前 7:1}，\BibleRef{路 14:26}，及\BibleRef{依 56:2-3}。}

主又说：\BibleQuote{男人不触摸妻子是好的，但因着淫乱，各人当有自己的妻子；}\BibleRef{格前 7:1} 仿佛解释这一点，他又说：\BibleQuote{免得撒殚诱惑你们。}\BibleRef{格前 7:5} 因为他不是对那因着只为生儿育女而节制地使用婚姻的人说\Quote{因着放纵}，而是对那些想要越过生儿育女的目标的人说的：免得我们的仇敌过多地点头同意，激起对他人欢愉的欲望。或许因为那仇敌因嫉妒而抵抗那些正义生活的人，并与他们争战，想要把他们引到自己行列中，藉着艰难的节制为他们制造机会。所以他正当地说：\BibleQuote{结婚比被火烧更好}\BibleRef{格前 7:9}，为要\Quote{丈夫对妻子尽义务，妻子对丈夫亦然，且不要互相亏负}这本是赐予生育的神圣帮助。\Quote{谁不恨自己的父亲、母亲、妻子、儿女，不能做我的门徒。}\BibleRef{路 14:26} 他并非命令恨自己的亲属：因为他说：\Quote{应孝敬父亲和母亲，为你好}\BibleRef{出 20:12}，而是说，不要被与理性相悖的欲望带走，也不要顺从世俗的风俗。因为家由家族构成，城邦由家构成；同样，\Person{保禄}也说那些从事婚姻的人\Quote{取悦世界}。\BibleRef{格前 7:33} 主又说：\BibleQuote{娶了妻子的，不要休妻；没有娶的，不要娶；}\BibleRef{格前 7:10} 那因着贞洁的意愿而明认不娶妻子的，应保持独身。因此，同一主藉着先知\Person{依撒意亚}给予两方合适的应许，这样说：\Quote{阉者不要说：我是枯木；}因为主对阉者说：\Quote{如果你们遵守我的安息日，并遵行我所吩咐的一切，我必赐给你们比儿女更好的地方。}因为贞洁不能独自使人成义，阉者的安息日也不能，除非他遵行诫命。他又对娶了妻子的人说：\BibleQuote{我的选民不会徒劳劳动，也不会在咒骂中生儿育女，因为种子是主所祝福的。}\BibleRef{依 65:23} 因为那按照圣言生养儿女、并在主内教育他们的，如同那藉着真正的教理讲授和培育而生育的，也有某种赏报存留，正如那蒙选的种子一样。但有些人将\Quote{咒骂}理解为生儿育女本身，而不明白圣经是在说反对那些人本身。因为那些真正是主选民的人，不制定教义，也不生育在咒骂和异端中的儿女。因此，所谓阉者，不是指那被强行阉割了身体部分的人，也不是指那独身的人，而是指那不生真理的人。这里的木头先前是枯干的；但如果他顺从圣言，并遵守安息日（藉着戒除罪恶），并遵行诫命，他将比那些没有正确生活培育、仅凭言语受教的人更为尊贵。师傅说：\BibleQuote{孩子们，我跟你们在一起的时候不多了。}\BibleRef{若 13:33} 因此，\Person{保禄}写给\Place{迦拉达}人时也说：\BibleQuote{我的孩子们，我再次为你们受生产之苦，直到基督在你们内形成。}\BibleRef{迦 4:19} 他又写信给\Place{格林多}人：\BibleQuote{即使你们在基督里有千万导师，但父亲却不多；因为我在基督里藉着福音生了你们。}\BibleRef{格前 4:15} 因此，\Quote{阉者不得进入天主的教会}\BibleRef{申 23:1}，因为他是荒芜的，不结果子，既无生活培育，也无言语。但那些\Quote{自阉}一切罪恶，\Quote{为天国}的人\BibleRef{玛 19:12}，他们是有福的，他们在世俗中斋戒。\par

% structure: p0068 source=h2 target=subsection reason=relative-heading-rank; base=h2

\subsection{第十六章。——\BibleRef{耶 20:14}；\BibleRef{约 14:3}；\BibleRef{咏 50:5}；\BibleRef{格前 9:27}，解释。}

\Quote{可咒骂}的\Quote{是我出生的那一日，愿它不被祝愿，} \BibleRef{耶 20:14} 耶肋米亚说：他并非绝对地咒骂生育，而是痛恨人民的罪恶和悖逆。因此他接着说：\Quote{我为何出生，为看到劳苦和痛苦，使我的岁月在永远的羞辱中？} \BibleRef{耶 20:18} 不但如此，凡宣讲真理的人，都因听者的悖逆而似乎陷入惩罚，并处于危险中。\Quote{为何我母亲的子宫不成为我的坟墓，使我不见雅各伯的苦难和以色列家族的辛劳？} 厄斯德拉先知说。\Quote{无人能洁净污秽，} 约伯说，\Quote{即使他的一生只有一天。} \BibleRef{约 14:4} 那么，请告诉我们，初生的婴儿在哪里行淫？或者，未行任何事的婴儿如何落在亚当的咒骂之下？因此，他们似乎必然要得出结论，说生育是恶的，不仅是身体的，也是灵魂的，因为灵魂使身体存在。当达味说：\Quote{我在罪恶中受孕，我母亲在邪恶中怀了我：} 他先知性地提到了母亲厄娃；但厄娃曾是\Quote{众生之母}；而如果他是\Quote{在罪恶中受孕}，但他本身并不在罪恶中，也不是罪恶。至于那些从罪恶归向信仰的人，是否可以说他们是从罪的习惯如同从\Quote{母亲}归向\Quote{生命}，十二先知中的一位会为我作证，他说：\Quote{若我为不虔敬献上长子，我腹中的果实为我灵魂的罪恶。} \BibleRef{米 6:7} 他并不责备那位说\Quote{你们要生育繁殖} \BibleRef{创 1:28} 的天主，而是指出生后的最初冲动，那时我们还不认识天主，他称之为\Quote{不虔敬}。如果有人以此理由说生育是恶的，那么他也应当说生育是善的，因为通过生育我们认识了真理。\Quote{你们要公正地洗净自己，不要再犯罪。因为有些人对天主无知，} 即那些犯罪的人。\Quote{因为我们的战斗不是对抗血和肉，而是对抗属神的邪恶。} \BibleRef{弗 6:12} 而这世界的\Quote{黑暗的掌权者}有力量试探，因此赐予宽恕。为此保禄也说：\Quote{我克制我的身体，使它服从；因为竞赛的人，在一切事上节制，} 也就是说，在一切事上节制，并非戒绝一切，而是节制地使用那些他判断为可用的东西，\Quote{他们为获得可朽坏的花冠，我们却是为获得不可朽坏的花冠，} 在搏斗中得胜，不是不经过努力就获得花冠。有些人甚至认为寡妇优于贞女，因为她以大勇蔑视了已体验过的快乐。\par

% structure: p0070 source=h2 target=subsection reason=relative-heading-rank; base=h2

\subsection{第十七章——那些主张婚姻和生育为恶的人，也辱骂天主的创造和福音的救恩计划。}

但若生育是恶的，那么亵渎者就要说曾参与生育的主是恶的，生育的童贞玛利亚是恶的。我何其不幸！多少、多大的恶！他们以咒骂攻击天主的旨意和创造的奥秘，因为他们攻击生育。由此卡西安捏造了\OriginalTerm{幻影说}{Docesin}；由此马尔基昂和瓦伦提诺也认为身体是属魂的，因为他们说，人若致力于情欲，\Quote{就与畜类相似了}。然而，当人确实因淫欲而疯狂，想要侵犯他人时，那时那人才真是如野兽一般：\Quote{他们成了发情的公马，各向邻人的妻子嘶鸣。} \BibleRef{耶 5:8} 假如他说蛇从无理性的动物那里取得了它计谋的理据，说服亚当同意与厄娃同房，仿佛否则原始父母就不会使用这种本性（如某些人所认为的），那么受造物再次被辱骂，因为它使人类比无理性的动物更软弱，而最初由天主形成的人追随了它们的榜样。但如果本性确实像畜类一样引导他们生育子女，而他们因受欺骗过早行动，当时还年轻，那么天主的公义审判针对那些没有等候他旨意的人；但生育是神圣的，世界因它而存在，本质、本性、天使、权能、灵魂、诫命、法律、福音、对天主的认识，都因它而存在。\Quote{凡有血肉的都如草，他的一切荣耀都像草上的花；草必枯萎，花必凋谢，但上主的话永远存留，} 那话傅油于灵魂，并使它与圣神结合。此外，如果我们教会中的救恩计划怎能在没有身体的情况下完成呢？因为连作为教会元首的那一位，也在肉身中度过了一生，虽无形无像，为教导我们注目于神圣原因的、无形无体的本性。\Quote{生命树，} 先知说，\Quote{是在善的愿望中，} \BibleRef{箴 13:12} 教导我们良善纯洁的愿望，那是在永生之主内的。现在，有些人认为丈夫与妻子在婚姻中的结合（称为\Quote{认识}）是罪：他们指出这由吃\Quote{知善恶树的果子} \BibleRef{创 3:5} 预示，通过\Quote{认识}一词的暗示\BibleRef{创 4:1}，这词标志着诫命的违犯。但若如此，那么对真理的认识也是吃生命树的果子。因此，正直节制的婚姻也能分享那树的果实。但我们先前已说过，可以善用、也可以滥用婚姻；这就是\Quote{知识树}，如果我们不违犯婚姻的法律。那么，我们的救主岂非既拯救灵魂，也拯救身体脱离情欲？因为如果肉身是灵魂的仇敌，他决不会通过恢复健康来加强肉身以对抗灵魂。\Quote{弟兄们，我说这话：肉和血不能承受天主的国，可朽坏的也不能承受不可朽坏的。} \BibleRef{格前 15:50} 罪恶既然是\Quote{可朽坏的}，就不能与作为正义的\Quote{不可朽坏的}有分。\Quote{你们竟这样愚蠢吗？你们既以圣神开始，如今却要以肉身结束吗？} \BibleRef{迦 3:3}\par

% structure: p0072 source=h2 target=subsection reason=relative-heading-rank; base=h2

\subsection{第十八章——应当避免两种极端意见：第一种是那些因憎恨造物主而禁绝婚姻的人；第二种是那些以此为契机放纵无耻情欲的人。}

因此，论到义德与救恩的和谐——这和谐是可敬而稳固的——有些人，如我们已指出的，过分强调它，以亵渎和咒骂的态度，连同各种不敬虔，接受节制的实践；其实他们本可虔诚地选择那符合健全规则的贞洁；他们应为所赐的恩宠而感恩，不该憎恨受造之物，也不该鄙视那些缔结婚姻的人。因为世界是被创造的，贞洁也是被创造的；双方都应感恩于自己所处的境遇，只要他们也了解自己所处的境遇。另一些人则放纵无度，狂妄无礼，正如圣经所说：\BibleQuote{他们成了向女人发情的公马，向邻舍的妻子嘶鸣。}\BibleRef{耶 5:8}以致他们自己无法自制，还说服邻人也追求快感；他们不幸地听从了圣经的那句话：\BibleQuote{你与我们同分，我们共用同一个钱袋，让我们共有一个钱囊。}\BibleRef{箴 1:14}因着这些人，同一位先知劝告我们说：\BibleQuote{不要与他们同行，要约束你的脚离开他们的路径。因为他们不是徒然设网罗捕鸟；他们既流血分赃，就为自己积蓄恶财。}这就是说，他们为自己追求不洁，也教导邻人同样的事，如同用尾巴击打的战士\BibleRef{默 9:10}，先知说，希腊人称之为\OriginalTerm{尾巴}{\GreekText{κέρκους}}。这些人可能就是先知所指的纵欲无度的人，他们用尾巴作战，是黑暗的可怒之子\BibleRef{弗 2:3}，在信德上玷污，自相残杀，杀害亲属。\BibleQuote{所以，你们要把旧酵除净，好使你们成为新和的面团。}\BibleRef{格前 5:7}宗徒向我们呼喊。他又因类似的人而愤怒，命令：\BibleQuote{若有人称为弟兄，却是奸淫的、贪婪的、拜偶像的、辱骂的、醉酒的、勒索的，这样的人不可与他交往，甚至与他一同吃饭都不可。}\BibleRef{格前 5:11}\BibleQuote{因为我已借着法律而死于法律，为叫我能为天主而活，我已同基督被钉在十字架上；我活着，已不是我活着，而是基督在我内活着；现今我在肉身内生活，是靠着对天主子的信德而生活。}\BibleRef{迦 2:19}贞洁而蒙福地借着服从诫命而生活。因此，那时我是在肉体中属肉体地活着；现今我在肉身内生活，是靠着对天主子的信德而生活。\BibleQuote{外邦人的路你们不要走，撒玛黎雅城你们不要进。}\BibleRef{玛 10:5}主如此说，劝我们离开相反的生活方式；因为\BibleQuote{邪恶的人的道路是恶的，这就是所有行恶之人的生活。}\BibleRef{箴 1:18}主说：\BibleQuote{那人有祸了，他若没有出生，比他绊倒我的一个选民更好。倒不如把一块磨石套在他脖子上，沉入海里，比他使我的一个人跌倒更好。因为天主的名因他们而受亵渎。}（\BibleRef{玛 26:24}；\BibleRef{罗 2:24}）因此，宗徒清楚地写道：\BibleQuote{我先前在信上写过，不可与奸淫的人交往。}\BibleRef{格前 5:11}直到此处：\BibleQuote{身体不是为奸淫，而是为主，主也是为身体。}\BibleRef{格前 6:13}他并不是将婚姻称为奸淫，而是通过下文表明：\BibleQuote{你们不知道那与娼妓结合的，便是与她成为一体吗？}\BibleRef{格前 6:16}难道有人能在婚前称处女为娼妓吗？\BibleQuote{你们不要彼此亏负，除非两相情愿，暂时分房。}\BibleRef{格前 7:5}他用\Quote{亏负}一词表明婚姻的债务是生养子女；这一点他在前文已经表明，说：\BibleQuote{丈夫对妻子要尽其所欠，妻子对丈夫也要如此。}\BibleRef{格前 7:3}此债务偿还后，妻子在持家和在基督内的信德上成为帮助者。他还更明确地说：\BibleQuote{对于已婚者，我——不是主——命令：妻子不可离开丈夫；若是离开了，就当守节，或与丈夫和好；丈夫也不可休妻。至于其余的人，是我说，不是主说：若某弟兄……直到如今她是圣洁的了。}\BibleRef{格前 7:14}对此，那些攻击法律和婚姻、仿佛婚姻仅是法律所允许而非在新约中也允许的人，又能说什么呢？对于这些法律，那些厌恶播种和生育的人能说什么呢？当宗徒设立\BibleQuote{善于管理自己家庭的}\BibleRef{弟前 3:4}主教作为教会领袖时，他岂不是以\BibleQuote{一个妻子的丈夫}\BibleRef{弟前 3:2}之家庭作为主的家庭吗？\BibleQuote{在洁净的人，一切都是洁净的；但在污秽和不信的人，什么都不洁净，连他们的理智和良心都污秽了。}\BibleRef{铎 1:15}至于逾越规则的快感，他说：\BibleQuote{你们不要受骗：奸淫的、拜偶像的、通奸的、作娈童的、好男色的、贪婪的、偷窃的、醉酒的、辱骂的、勒索的，都不能承继天主的国。我们这些曾陷于这些罪中的人，也已经洗净了。}\BibleRef{格前 6:9-11}而那沉浸于这种无度行为的人，却是从节制洗入奸淫，主张应纵容情欲和私情，教导人从有节制变为无节制，将自己的希望锚定在肢体无耻的快乐上；他们使前往他们那里的人被摈弃于天主的国之外，而非被登记在内，以伪称的知识为名，走上那导向外面黑暗的道路。\BibleQuote{此外，弟兄们，凡是真实的、凡是高尚的、凡是正义的、凡是纯洁的、凡是可爱的、凡是荣誉的，若有任何德行，若有任何可称赞的，你们要思念这些；你们在我身上所学到的、所领受的、所听到的、所看到的，你们要实践这些；那么和平的天主必与你们同在。}\BibleRef{斐 4:8}伯多禄在书信中也同样说：\BibleQuote{好使你们的信德和望德在于天主，你们既因服从真理而洁净了灵魂，}\BibleRef{伯前 1:21}如同服从的儿女，不要效法从前无知时的私欲，而要效法那召你们的圣者，在一切行为上也成为圣的。因为经上记载：\BibleQuote{你们应当是圣的，因为我是圣的。}\BibleRef{伯前 1:16}无论如何，我们这针对那些伪称知识之人的必要辩驳，已经将我们带得比适当更远，使我们的论述过于冗长。因此，\WorkTitle{杂论}第三卷，即关于真正哲学的注释，至此结束。\par

\SourceRecord{Translated by William Wilson.；Ante-Nicene Fathers；Vol. 2.；Edited by Alexander Roberts, James Donaldson, and A. Cleveland Coxe.；Buffalo, NY: Christian Literature Publishing Co.,；1885.；<http://www.newadvent.org/fathers/02103.htm>.}

% structure: p0001 source=h1 target=section reason=page-title

\section{\WorkTitle{杂文集}（第四卷）}

% structure: p0002 source=h2 target=subsection reason=relative-heading-rank; base=h2

\subsection{第一章 内容次序}

依我看来，接下来应论述殉道与何者为成全之人。随此将包括依所论之点要求而须述说者，以及奴仆与自由人，无论男女，皆应同等修习哲学。在后续部分，完成关于信德与探究的讨论后，我们将阐述象征的领域。如此，在简要结束伦理论述后，我们将展示蛮族哲学带给希腊人的益处。在此概要之后，将呈现对圣经的简要解释，既针对希腊人也针对犹太人，以及在前几卷\WorkTitle{杂文集}中（因不得不顾及众多事宜）未能包括的一切要点，依照该作品的起始，旨在以一篇注释完成。除此之外，在完成概要后，我们将尽我们所提议的范围，论述希腊人与蛮族关于本原的物理学说，凡其意见传至我们者，并反驳哲学家所构想的主要观点。\par

在这些之后，自然应继以神学俯览，然后讨论关于预言的传承意见。如此，既证明我们所信之圣经因其全能的权威而有效，我们便能连续地研读它们，并由此向所有异端显示：惟一天主与全能之主，确实由法律、先知以及蒙福的福音所宣讲。当我们试图在写作中消除他们所提主张的力量，并用圣经本身说服他们（尽管非其所愿）时，许多反对异端者的反驳正等待着我们。\par

完成我们在注释中所提议的全部内容后（若圣神愿意，我们服务于迫切的需要——因为在达至真理之前，非常必要以序言方式述说当言之物），我们将着手真正的\OriginalTerm{真知}{\GreekText{γνωστικός}}自然的科学，在进入大奥秘之前先接受小奥秘的传授。如此，在预先清除并勾勒出需要初步说明和陈述的主题后，便无任何障碍妨碍对神圣之物的真正神圣宣告。自然之学，或更恰当地说，观察，即按照真理准则包含于真知传统中者，依赖于关于宇宙生成的讨论，由此上升至神学领域。因此，我们将从先知所传的创世开始阐述传承的内容，同时引入异端者的主张，并尽力驳斥之。但若天主愿意，且如祂所默感，这将被书写；如今我们必须进行所提议之事，完成伦理论述。\par

% structure: p0006 source=h2 target=subsection reason=relative-heading-rank; base=h2

\subsection{第二章 \WorkTitle{杂文集}或\WorkTitle{缀集}一名之含义}

我们这些笔记，如我们为粗心不熟之读者多次所说，应具备多样特征——正如其名所示，拼凑而成——不断从一件事转到另一件事，在讨论序列中暗示一事而论证另一事。\Quote{寻找金子的人，}\Person{赫拉克利特}说，\Quote{挖许多土，只找到少量金子。}但真正金族之人，在开采与他们相近之物时，将在少量中发现大量。因为圣言将找到理解它的人。\WorkTitle{杂文集}的笔记有助于记得并表达真理，对那能以理性探究者而言。此外，你必须进行其他劳作与研究。因为对于踏上陌生道路之人，仅仅指出方向便已足够；此后他们必须自己行走并发现其余。据说，有一个奴隶曾问神谕他应做什么才能取悦主人，皮提亚女祭司回答说：\Quote{你若寻找，必将寻见。}确实，如所表明，发现隐藏的善是困难之事，因为\par

\Quote{在德行之前设置努力，\newline{}通往它的路又长又陡，\newline{}起初崎岖；但到达顶峰后，\newline{}便变容易，虽然先前困难。}\par

\Quote{因为门是窄的，路是狭的}，诚然，是主的道路。\Quote{天国是以猛力夺取的。}\par

因此，\Quote{你们求，必要给你们}，持守真正王道，不偏左右。所以，毫不奇怪，本论著所包含之教义种子的生殖能力在小空间中巨大，正如圣经所说，\Quote{田间的百草}\BibleRef{约 5:25}。因此，\WorkTitle{杂文集}的笔记有其恰当标题，奇妙地类似\Person{索福克勒斯}所描述的那古老而采自各种什物的祭品：\par

\Quote{那里有羊的卷毛，也有葡萄，\newline{}奠祭与贮藏的葡萄，\newline{}并掺有各类水果，\newline{}橄榄的油脂，与最巧妙的\newline{}黄色蜜蜂的蜡制作品。}\par

同样，我们这\WorkTitle{杂文集}，正如喜剧诗人\Person{提摩克勒斯}庄稼汉所说，出产\Quote{无花果、橄榄、无花果干、蜂蜜，如同一切果实丰盛的田地}；因这丰饶，他补充说：\par

\Quote{你所说的不是农作的收获冠，而是一种丰饶之冠。}\par

因为雅典人惯常呼喊：\par

\Quote{收获冠带来无花果和肥美的面包，\newline{}以及杯中的蜂蜜和橄榄油供你涂抹。}\par

我们必须常常摇晃并扬起这大量种子的混合物，如同在簸筛中，以便分离麦子。\par

% structure: p0017 source=h2 target=subsection reason=relative-heading-rank; base=h2

\subsection{第三章 人的真正卓越}

大多数人有不稳定和漫不经心的性情，如同暴风雨的本性。\Quote{缺乏信德做了许多好事，而信德做了坏事。} 埃皮卡摩斯说：\Quote{不要忘记练习不信；因为它是灵魂的肌腱。} 现在，不信真理带来死亡，正如相信带来生命；反之，相信谎言和不信真理则加速毁灭。自制和放纵也是如此。克制自己不去行善是恶行的工作；但远离恶是救恩的开端。因此，安息日通过禁绝邪恶，似乎表明自制。并且，请问，人在哪方面与野兽不同，而天主的天使在哪方面比人更明智？\BibleQuote{你使他稍微逊于天使。} 因为有些人并不将这段圣经解释为论及主，虽然祂也取了肉身，而是论及完人和\OriginalTerm{真知者}{gnostic}，在时间上比天使低下，且因身体的遮盖[而如此]。那么，我称智慧不外乎科学，因为生命与生命并无不同。因为活着是必死本性的共同之事，即对人类和对那被赐予不死者而言；同样，沉思和自制的能力也是如此，二者中一个更为卓越。基于此，毕达哥拉斯在我看来曾说过：唯有天主是明智的，因为宗徒在致罗马人书也写道：\Quote{为了所有外邦人的信德的服从，通过耶稣基督，使独一明智的天主被认识；} \BibleRef{罗 16:26} 而他本人是哲学家，因他与天主的友谊。因此，经上记着说：\BibleQuote{天主与梅瑟交谈，如同朋友与朋友交谈。} \BibleRef{出 33:11} 那么，真理既然对天主是清晰的，就立即产生真理。而真知者热爱真理。\Quote{去，} 经上说，\Quote{到蚂蚁那里，懒惰的人，并作蜜蜂的门徒；} 撒罗满如此说。\BibleRef{箴 6:6} 因为如果每个受造物本性的特有功能只有一个，无论是牛、马还是狗，那么人的特有功能是什么呢？在我看来，他像是半人马，一个帖撒利的虚构物，由理智和非理智部分，即灵魂和身体组成。身体耕种土地，并趋向它；但灵魂被提升到天主：经过真哲学的训练，它奔向上面它的同类，远离身体的贪欲，此外还有辛劳和恐惧，尽管我们已经表明忍耐和恐惧属于善人。因为如果\Quote{通过法律是罪的认知，} \BibleRef{罗 3:20} 如那些诋毁法律的人所声称，且\Quote{在法律之前，罪已经在世界上，} \BibleRef{罗 5:13} 然而\Quote{没有法律，罪是死的，} \BibleRef{罗 7:6} 我们反对他们。因为当你除去恐惧的原因（罪），你就除去了恐惧；更不用说惩罚，当你除去了引起贪欲的东西。\Quote{法律不是为义人设立的，} \BibleRef{弟前 1:9} 圣经说。那么，赫拉克利特说：\Quote{如果这些事没有发生，他们就不会知道正义的名字。} 苏格拉底说：\Quote{法律不是为了好人的缘故而设立的。} 但那些吹毛求疵的人甚至不知道这一点，正如宗徒所说：\Quote{那爱弟兄的人不行恶；} 因为这一点：\Quote{不可杀人，不可奸淫，不可偷盗；若有任何别的诫命，都包含在这句话里：你应当爱你的近人如你自己。} \BibleRef{罗 13:8} 同样，经上也说：\Quote{你应当全心、全灵、全意爱主你的天主，并爱你的近人如你自己。} \BibleRef{路 10:27} 并且\Quote{如果那爱近人的人不行恶，} 并且如果\Quote{每条诫命都包含在这句话里，就是爱我们的近人，} 那么诫命通过用恐惧威胁，产生爱，而不是仇恨。因此，法律产生了恐惧的情感。\Quote{所以，法律是圣洁的，} 并且实际上是\Quote{属神的，} 按照宗徒所言。那么，我们必须恰当地研究身体的本性和灵魂的本质，把握每一者的终向，并且不将死亡视为邪恶。\Quote{因为你们曾是罪的奴仆，} 宗徒说，\Quote{就脱离了正义。在那些你们现在以为耻的事上，你们得了什么果实？因为那些事的结局就是死亡。但现在，你们从罪中得了释放，成为天主的奴仆，你们有果实归于圣洁，结局是永生。因为罪的工价是死亡；但天主的恩赐是永生，在我们的主耶稣基督内。} \BibleRef{罗 6:20} 那么，可以大胆断言，已经表明：死亡是在罪的状态下灵魂与身体的结合；而生命是与罪的分离。有许多贪欲的木桩和沟渠阻碍我们，必须跳过愤怒和怒气的坑，并且我们必须避免那些阴谋对付我们之人的一切诡计——他们不再愿意\Quote{通过镜子}看见天主的认识。\par

\Quote{远见的宙斯从人身上取走一半的美德，\newline{}当他将人贬为奴隶状态时。}\par

圣经视那些\Quote{在罪之下}和\Quote{被卖给罪}的人为奴隶，他们是享乐和身体的爱好者；并且是野兽而非人，\Quote{那些变得像牛犊、马匹，嘶叫着追求邻舍的妻子的人。} 放荡者是\Quote{好色的驴，} 贪婪者是\Quote{残暴的狼，} 欺骗者是\Quote{蛇。} 因此，灵魂与身体的分离，作为终生的学习，在哲学家\OriginalTerm{真知者}{gnostic}中产生敏捷，使他能轻松承受自然的死亡，即解除灵魂与身体捆绑的锁链。\Quote{世界已被钉在十字架上与我，我也被钉在十字架上于世界，} 宗徒说；\Quote{现今我活着，虽在肉身内，却以在天上生活为念。} \BibleRef{迦 6:14}\par

% structure: p0021 source=h2 target=subsection reason=relative-heading-rank; base=h2

\subsection{第四章 殉道之赞颂}

因此，合情合理地，\OriginalTerm{诺斯替者}{gnostic}当被激怒时，轻易地服从，并将自己的身体交给那要求的人；既已事先脱去这躯壳的情感，不是侮辱那诱惑者，反而，依我看来，是在训练他说服他——\par

\Quote{从何等尊荣与何等广大的财富坠落，}\par

他实在向自己见证他对天主忠诚忠信；向诱惑者见证他白白嫉妒那因爱而忠诚的人；向主见证那因他的教训而受感动的确信，他不会因怕死而背离这教训；此外，他也以自己的行动证实了宣讲的真理性，表明他所奔赴的天主是有能力的。你会惊叹于他的爱，他满心感激地显明这爱，与那与他相亲的联合，并借着他宝贵的血羞辱了不信的人。他因戒命而避免因害怕否认基督；也不因盼望备好的赏赐而出卖信德，却因爱主而欣然离此生命；也许既感谢那赐给他离世缘由的人，也感谢那设计陷害他的人，因为他获得了一个他自己未曾提供的可敬理由，借着他的忍耐向那人表明他是怎样的人，并借着他的爱向主表明，甚至在他出生前，那位知道殉道者抉择的主已向他显明。于是他勇敢地走向主，他的朋友，他自愿为此朋友交出身体，并按审判官所期望的，交出了灵魂，从我主救主那里听到诗歌般的话语：\Quote{亲爱的兄弟}，因他生命的相似。我们称殉道为完美，不是因为人像他人一样结束生命，而是因为他展现了爱的完美工作。古人称赞希腊人中战死者的死亡，不是说他们劝人暴死，而是因为结束生命于战争的人，脱离了对死亡的恐惧，与身体分离时未经历先前的痛苦或灵魂衰弱，如同患病之人。因为他们死时身心软弱且贪生；因此他们没有交出纯洁的灵魂，而是带着他们的欲望如同铅坠；只有那些在美德上显赫的人除外。有些人带着欲望战死，这些人与那些因疾病而耗竭的人毫无区别。\par

如果向天主宣认是殉道，那么每个在认识天主中纯洁生活、遵守诫命的灵魂，无论以何种方式脱离身体，都是藉生活和言语的见证人——一生直至离去都在流信仰如血。例如，主在福音中说：\BibleQuote{凡为福音和我的名舍弃父亲、母亲、弟兄等等}，\BibleRef{玛 19:29} 他是有福的；这不是指简单的殉道，而是指诺斯替式的殉道，如同那按福音规则行事的人，因爱主（因为对圣名的认识和对福音的理解指出真知，而不是裸的名称），以致舍弃世俗的亲属、财富和一切所有，为度不受激情的生活。\Quote{母亲}比喻指国家和生计；\Quote{父亲}是公民政体的法律：高尚的义人必须怀着感恩轻视它们；为成为天主的友人和获得圣地中的右手之位，如宗徒们所做的一样。\par

赫拉克利特说：\Quote{神和人都尊敬战死沙场的人；} 柏拉图在\WorkTitle{理想国}第五卷写道：\Quote{在军中服役而死的人中，凡死后赢得声誉的，难道我们不会说他是金族之首吗？当然。} 但金族与诸神在一起，他们在天上，在固定的天球上，主要掌管对人类行使眷顾。现在有些异端误解了主，既邪恶又怯懦地贪生；他们说真正的殉道是对唯一真神的知识（我们也承认这一点），而那个借着死亡宣认的人是一个自杀者；并引证其他类似怯懦的诡辩。我们将在适当时候答复他们；因为他们在第一原则上与我们不同。现在，我们也说那些冲向死亡的人（因为有一些不属于我们，只是同名的，急于交出自己，这些可怜的人因仇恨造物主而死），我们说这些人驱逐自己而不成为殉道者，即使他们公开受罚。因为他们没有保存信仰殉道的特征标记，因为他们不认识唯一真神，却将自己交付给虚空的死亡，如同印度\OriginalTerm{裸形智者}{Gymnosophists}投身于无用的火焰。\par

但既然这些假名者诽谤身体，让他们知道身体的和谐构造有助于理解，这理解引导向善的本性。因此，在\WorkTitle{理想国}第三卷中，柏拉图（他们大声引用他的权威来贬低生育）说：\Quote{为了灵魂的和谐，必须照顾身体}，借着身体，那宣布真理的人可能活着，且活得好。因为我们是通过生命与健康之路学习真知。但那人若不能在没有被必要之事占据的情况下达到高处，而是通过它们做那趋向知识的事，他不是选择好好地活着吗？因此，在活着中，好的生活得以确保。那在身体中献身于好生活的人，正被送往不朽的境界。\par

% structure: p0028 source=h2 target=subsection reason=relative-heading-rank; base=h2

\subsection{第五章 论轻看痛苦、贫穷及其他外在事物}

斯多葛学派堪称令人敬慕的对象，他们断言灵魂不受身体影响，既不因疾病而趋于恶，也不因健康而趋于善；他们声称这两者都是无关紧要的。的确，约伯因超乎寻常的节制和卓越的信德，从富足变为贫穷，从受尊敬变得被轻视，从俊美变得丑陋，从健康变为疾病，他被描绘为一个良好的榜样，使试探者蒙羞，祝福他的造物主；他将后来发生的事视为最初的事，并最清楚地教导：真知者有可能善用一切境遇。宗徒表明，古代的事迹被提出作为我们纠正的模范，他说：\Quote{以致我所带的锁链在御营全军和其余的人中已显明是为基督的缘故；并且那些在主内的弟兄，许多因我所带的锁链而笃信不疑，更敢无所畏惧地宣讲天主的圣言。}\BibleRef{斐 1:13}——因为殉道者的证言是光荣圣化的皈依的榜样。\Quote{圣经上所写的，都是为教训我们写的，使我们因圣经所生的忍耐和安慰，可得着安慰的盼望。}\BibleRef{罗 15:4}当痛苦临到时，灵魂似乎逃避它，并认为脱离当前痛苦是可贵的。在那时，灵魂从学习中松懈下来，其他德行也被忽视。然而我们并不说是德行本身受苦，因为德行不受疾病影响。但那同时拥有德行和疾病的人，因后者的压力而受苦；倘若尚未达到自制习惯的人不是高尚的人，他就会心烦意乱；不能忍受痛苦便被视同逃避痛苦。\par

贫穷也是如此。因为它迫使灵魂脱离必要的事物，我指的是默观和纯全的无罪，迫使那尚未全心热爱天主的人专注于生计；反之，健康和充足的必需品使灵魂自由无碍，并能够善用手边的一切。\Quote{因为，}宗徒说，\Quote{这样的人在肉身必有苦难。但我愿意宽免你们。我愿意你们无所挂虑，以合宜和专务主的事，没有分心。}\par

这些事物之所以应避免，并非为它们本身，而是为身体的缘故；照顾身体是为灵魂的缘故，因为灵魂是身体的归宿。因此，过真知者生活的人有必要知道什么是适宜的。既然享乐不是善事这一事实从某些享乐是恶的这一点得以承认，那么根据这一推论，善显得是恶，恶显得是善。此外，如果我们选择某些享乐而避开另一些，那么并非每一种享乐都是善事。\par

同样，同样的原则也适用于痛苦：我们忍受某些痛苦，避开另一些。但选择与避开是依据知识来进行的；所以善事不是享乐，而是知识，我们凭知识在特定时间选择某种享乐。殉道者选择透过当前痛苦而存在于前景中的享乐。如果痛苦被视为存在于口渴中，享乐存在于饮水中，那么之前的痛苦便成为享乐的成因。但恶不能成为善的成因。因此，两者都不是恶。西蒙尼德斯（亚里士多德亦然）写道：\Quote{健康是第一好事，美貌是第二好事，第三好事是不靠欺骗致富。}\par

墨伽拉的泰奥格尼斯说：——\par

\Quote{为了摆脱贫穷，你必须，\newline{}啊，库尔诺斯，从陡峭的岩石上\newline{}纵身跳入深海。}\par

另一方面，喜剧诗人安提法奈斯说：\Quote{普路托斯（财富）一旦抓住那些比别人看得更清楚的人，就使他们变瞎。}现在，诗人们宣告他从出生就瞎：——\par

\Quote{便带出他，生来瞎眼，未见阳光。}\par

卡尔基斯的欧福里翁说：——\par

\Quote{财富和奢侈的盛宴，\newline{}对人而言是男子气概最糟的训练。}\par

欧里庇得斯在\WorkTitle{亚历山大}中写道：——\par

\Quote{据说，\newline{}贫穷通过不幸获得了智慧；\newline{}但巨大的财富不仅将攻陷\newline{}斯巴达，而且将攻陷每一座城。}\par

\Quote{因此，凡人拥有的不仅是白银或黄金的货币，而且还有德行，}如索福克勒斯所说。\par

% structure: p0042 source=h2 target=subsection reason=relative-heading-rank; base=h2

\subsection{第六章 真福八端中的几点}

我们的至圣救主将贫穷与财富以及类似之事，应用于属灵之物和感官之物上。因为当祂说：\BibleQuote{为义而受迫害的人是有福的，}\BibleRef{玛 5:10}祂清楚教导我们在任何情况下都要寻求殉道者，若有人为义而贫穷，就证明他所爱的义是善的；若他\Quote{为义而饥渴，}就证明义是最好的。同样，为义而哭泣哀恸的人，证明这至善的律法是美好的。因此，正如\Quote{受迫害的人}一样，那些为义而\Quote{饥渴的人}，被那认可这真正渴望（连饥荒也无法阻止的渴望）的祂称为\Quote{有福的}。若他们\Quote{对义本身感到饥渴，}他们便是有福的。\Quote{贫穷的人有福了，}无论是在\Quote{神贫}还是在境遇中——即若为义之故。祂称有福的不是单纯的穷人，而是那些愿意为义成为穷人的人——那些为获得\Quote{美善}而轻视世俗尊荣的人；同样，那些因洁德而形貌与品格优美的人，那些出身高贵且可敬的人，因义而获得义子名分，因此\BibleQuote{获得了成为天主子女的权能，}\BibleRef{若 1:12}能\BibleQuote{践踏蛇和蝎子，}并统治恶魔和\BibleQuote{仇敌的军队。}\BibleRef{路 10:19}总之，主的训诫欣然将灵魂从肉身中引开，即使肉身在分离时挣扎。\Quote{因为那爱惜自己生命的，必要丧失生命；那丧失自己生命的，必要获得生命，}只要我们将自己可朽的部分与天主的不朽结合。天主的旨意[是要我们]获得天主的认识，即不朽的传达。因此，那依悔改之言认识自己生命是罪恶的，将丧失它——从罪恶（即生命被夺走之处）中丧失；但丧失它，将找到它，依据那向信德复活、向罪恶而死的服从。这便是所谓的\Quote{找到生命}，\Quote{认识自己}。\par

然而，斯多葛派说，那导向神圣事物的皈依，是通过一种改变而实现的，即灵魂被改变为智慧。而柏拉图说：\Quote{论灵魂转向更美好的事物，以及从某种黑夜般的白昼中转变。}如今，哲学家们也容许善人按照理性退出生命，当有人剥夺了他的积极活动，以至于不再留有行动的希望时。而那强迫我们否认我们所爱之主的审判者，我认为他显明了谁是、谁不是天主的朋友。在这种情况下，甚至没有留下余地来考察一个人更偏爱什么——人的威胁还是天主的爱。而戒除恶行，在某种程度上，似乎导致恶习倾向的减弱和消失，其活力因消极而被摧毁。这就是\Quote{变卖你所有的，分给穷人，然后来跟随我}\BibleRef{玛 19:21}的含义——即，跟随主所说的话。有人说，通过\Quote{你所有的}，祂指的是灵魂中那些与本性质不相似的事物，尽管他们无法说这些如何施舍给穷人。因为天主按照各人的功过施予一切，祂的分配是公义的。因此，轻看天主在你慷慨中分给你的财产，服从我所说的话；赶紧攀登圣神的高峰，不仅因戒除邪恶而成义，此外还因基督般的慈善而得以成全。在这件事上，祂定了那自夸已履行了法律诫命的人的罪，指出他不爱近人；正是通过慈善，那按诺斯替上升阶梯是安息日的主的爱，彰显了自己。那么，按照我的看法，我们必须投靠救恩之言，既不是出于对惩罚的恐惧，也不是出于对恩赐的应许，而是为了善本身。如此而行的人，站在圣所的右边；而那些认为借着施舍易朽之物，就能换取不朽之物的人，在关于两兄弟的比喻中被称为\Quote{雇工}。在此，对于\Quote{相似和肖像}这一说法，难道不是有了些光亮吗？因为有些人按照基督的相似而生活，而那些站在左边的人则按照自己的肖像而生活。那么，有两样东西出自真理，一个根在两者之下——然而选择并不相等，或者更确切地说，选择中的差别并不相等。在我看来，通过模仿而选择，与那些按照知识而选择的人的选择不同，正如被点燃的东西与被照亮的东西不同一样。那么，以色列就是那按照圣经的相似之光。但肖像则是另一回事。拉匝禄的比喻，通过显示富人和穷人的肖像，意味着什么？而\Quote{没有人能事奉两个主人：天主和玛门}这句话又是什么意思？——主如此称呼贪爱钱财。例如，那些贪婪的人被邀请赴宴，却没有回应邀请，不是因为他们拥有财产，而是因为他们对所拥有的过度爱恋。那么，\Quote{狐狸}有洞。祂称那些专心于从地中开采挖掘的财富的邪恶和属世的人为狐狸。同样，关于黑落德：\Quote{你们去告诉那只狐狸：看，我今天明天驱魔治病，第三天就要完成了。}\BibleRef{路 13:32}因为祂将\Quote{天空的飞鸟}这个名称应用于那些与其他飞鸟不同的人——那些真正纯洁的人，那些有能力飞向天上圣言知识的人。因为不仅财富，而且荣誉、婚姻和贫穷，对于那些不适当的人来说，都有无数的忧虑。祂在四种种子的比喻中指出了这些忧虑，当时祂说\Quote{落在荆棘中的圣言种子}和篱笆中被它们所窒息，不能结果实。因此，必须学习如何利用每一件事，以便按照知识，通过善生，为永生之境受训练。因为经上说：\Quote{我见过恶人高升，耸立如黎巴嫩的香柏木；我经过，}圣经说，\Quote{看，他已不在了；我寻找他，他的地方也找不着了。保持纯真，注视正直：因为和平的人终有结局。}这样的人将是那用全心真诚相信，且全魂平静的人。\Quote{因为这百姓用嘴唇尊敬我，他们的心却远离主。}\Quote{他们用口祝福，心里却咒骂。}\Quote{他们用口爱祂，用舌头向祂说谎；但他们的心不向祂正直，他们也忠于祂的盟约。}因此\Quote{愿说谎的嘴唇成为无声，愿主毁掉夸夸其谈的舌头：那些说，我们要夸大我们的舌头，我们的嘴唇是我们自己的；谁能作我们的主？因为困苦人的受迫害和穷乏人的叹息，我现在要起来，说：我要把他安置在稳妥之处；我要为他辩白。}因为基督属于谦卑的人，他们不抬高自己反对祂的羊群。\Quote{因此，不要在地上为自己积蓄财宝，那里有虫蛀、锈蚀，也有贼挖洞偷窃。}\BibleRef{玛 6:19}主说，这或许是在责备贪婪的人，或许也责备那些只是忧虑和充满挂虑的人，以及那些放纵身体的人。因为情欲、疾病和邪恶的念头\Quote{偷窃}心灵和整个人。但我们真正的\Quote{财宝}是在与我们心灵相连的地方，因为它赐予正义的传递能力，表明我们必须将借此获得的东西归给我们旧日生活的习惯，并投靠天主，恳求仁慈。祂实在就是\Quote{那不会朽坏的钱袋}，永生的储备，\Quote{在天上不枯竭的财宝。}\BibleRef{路 12:33}\Quote{因为我要怜悯谁，就怜悯谁。}\BibleRef{罗 9:15}主说。他们向那些为了正义而愿贫穷的人说这些话。因为他们曾在诫命中听见\Quote{那宽阔广大的路通向灭亡，有许多人从那里进去。}\BibleRef{玛 7:13}这个断言不是针对别的东西，而是针对放荡、爱女人、爱荣誉、野心和类似的激情。因为祂这样说：\Quote{糊涂人哪！今夜就要索回你的灵魂；你所预备的一切，将归谁呢？}\BibleRef{路 12:20}而诫命正是用这些话表达的：\Quote{因此，你们要谨慎，躲避一切贪婪。因为一个人的生命不在于他拥有丰盛的财产。因为人纵然赚得了全世界，却赔上了自己的灵魂，为他有什么益处？或者，人还能拿什么作为自己灵魂的代价？}\BibleRef{玛 16:26}\Quote{所以，我告诉你们：不要为你们的生命忧虑吃什么，也不要为你们的身体忧虑穿什么。因为生命贵于食物，身体贵于衣服。}\BibleRef{玛 6:31}又说：\Quote{因为你们的父知道你们需要这一切。}\Quote{但你们先该寻求天主的国和它的义德，}因为这些都是大事，而那些属于此世的小事\Quote{自会加给你们。}\BibleRef{玛 6:32}那么，祂不是明明劝我们跟随诺斯替生活，并吩咐我们在言行上寻求真理吗？因此，训练灵魂的基督衡量一个人是否富有，不是根据他的恩赐，而是根据他的选择。因此，据说匝凯（或按有些人说是玛窦），税吏长，一听到主屈尊到他那里去，就说：\Quote{主，如果我讹诈了谁什么，我就偿还他四倍；}于是救主说：\Quote{今天人子来了，找到了那失落的。}再者，看见富人按照他们的财富向银库投钱，而寡妇投了两个小钱，祂说\Quote{那寡妇投的比众人还多，}因为\Quote{他们都由他们的富余中捐献，而她却由她的匮乏中捐献。}因为祂使一切为灵魂的操练服务，祂说：\Quote{温良的人是有福的，因为他们要承受土地。}\BibleRef{玛 5:5}温良的人就是那些在灵魂中平息了不信的战争、忿怒和情欲以及其他从属形式的战争的人。祂称赞那些出于选择而非出于必要而温良的人。因为在主那里有酬报和许多\Quote{住所}，与人们的生活相应。\Quote{谁若，}祂说，\Quote{因先知的名接纳一位先知，必得先知的赏报；谁若因义人的名接纳一个义人，必得义人的赏报；谁若接纳我这门徒中最小的一个，必不失他的赏报。}\BibleRef{玛 10:41}再者，祂用数目不等的时辰指明了美德按功绩的差别和崇高的赏报；此外，通过给每个工人同等的赏报——即救恩，由银币所象征——祂指明了正义的平等；而祂又用那些工作了不等时间的人暗示了蒙召者的不同。因此，他们将按照他们认为配得的合适住所而工作，作为赏报，成为不可言喻的管理和服务中的同工。\Quote{那么，}柏拉图说，\Quote{那些似乎被召向圣洁生活的人，就是那些摆脱并释放了如监狱般的地上场所，而抵达了高处纯洁居所的人。}他用更清晰的词句又表达了同样的事：\Quote{那些藉哲学充分净化了这些事物的人，完全无身体地生活直到永远。尽管他们被包裹在某些形状中；有些是空气的形状，有些是火的形状。}他进一步补充：\Quote{他们到达比那些更美的居所，这不容易描述，现在也没有足够的时间。}因此有理地，\Quote{哀恸的人是有福的，因为他们要受安慰。}\BibleRef{玛 5:4}因为那些悔改了先前邪恶生活的人将获得\OriginalTerm{蒙召}{\GreekText{κλῆσιν}}，因为\OriginalTerm{受安慰}{\GreekText{παρακληθῆναι}}就是这个意思。有两种风格的忏悔者。较常见的是因所做的事而恐惧；但另一种较特殊的，是心灵因良心而产生的羞愧感。那么，不论在此处或别处（因为没有一处地方缺乏天主的恩惠），祂又说：\Quote{怜悯人的人是有福的，因为他们要受怜悯。}怜悯并不像一些哲学家所想象的那样，是因别人的灾祸而感到痛苦，而是某种善，正如先知们所说的。因为经上说：\Quote{我要的是怜悯，而不是祭献。}祂所说的怜悯人，不仅指那些实行怜悯行为的人，也包括那些愿意实行却无能为力的人；他们在意愿上已经做了。因为有时我们愿意通过施舍金钱或亲自努力来施行怜悯，比如帮助一个匮乏的人，或帮助一个生病的人，或陪伴一个处在紧急情况中的人；但因贫穷、疾病或年老（这也是一种自然的疾病）而不能实现我们的意图，对于我们所受激励之事，无法将之引导至我们所愿的结局。那些怀有同等意愿的人，与有能力的人分享同样的荣誉，尽管后者在资源方面有优势。既然达到救恩成全有两条路，即工作和知识，祂称\Quote{内心纯洁的人是有福的，因为他们要看见天主。}如果我们真看事情的真相，知识就是灵魂主导官能的净化，并且是一种善的活动。因此，有些事物本身是善的，另一些则通过参与善而是善的，正如我们说善行是善的。但没有占据质料地位的中介事物，既不能构成善行也不能构成恶行，我指的是生命、健康以及其他必要的或环境性的事物。那么，就肉体的情欲而言是纯洁的，就神圣思想而言是纯洁的，他指的是那些获得天主知识的人，此时灵魂的主导官能没有任何虚假之物阻碍其力量。因此，当那以诺斯替方式分享这神圣品质的人投身于默观，在纯洁中与神圣者相通时，他就更接近进入无感情的本然状态，以至于不再具有科学和拥有知识，而是成为科学和知识本身。\par

\Quote{缔造和平的人是有福的}，因为他们制服并驯服了那作战反对理智倾向的法律、愤怒的威胁、情欲的诱饵以及其他对抗理性的激情；他们生活在善行与真知之中，将得以恢复为更宝贵的义子身份。因此，完美的缔造和平是在一切境况下保持平安不变；称天主的眷顾为神圣而良善；其存在在于对神圣与人事的认识，由此视世间的对立为受造界最和谐的秩序。那些教导与罪的诡计作战的人转向信德与和平的人，也是缔造和平者。在我看来，当主教导我们为了爱天主必须灵知地藐视死亡时，这便是一切德行的总纲。祂说：\BibleQuote{为义而受迫害的人是有福的，因为他们将称为天主的儿子；} \BibleRef{玛 5:10} 或者，如一些改编福音者所说：\Quote{被义所迫害的人是有福的，因为他们将是完美的。} 又说：\Quote{为我的缘故受迫害的人是有福的，因为他们将有一个不受迫害的地方。} 又说：\BibleQuote{几时人们因人子的缘故恨你们、断绝你们、辱骂你们、弃绝你们的名字如同可憎的，你们是有福的。} \BibleRef{路 6:22} 只要我们不以憎恨对待迫害我们的人，不因受他们惩罚而怀恨他们，好像我们比预期的更迟受到考验；而是也知道每一次考验都是作证的机会。\par

% structure: p0046 source=h2 target=subsection reason=relative-heading-rank; base=h2

\subsection{第七章 殉道者的真福}

那说谎、不忠、投奔魔鬼阵营的人，我们以为他处于何等邪恶之中？因此，他不信天主，就是否定主，或者说他欺骗了自己的希望；不信主的人，就是不遵行祂命令的人。\par

再者，否认主的人，岂不也否认自己？因为剥夺自己与主的关系，岂不是抢夺了主人的权威？所以，否认救主就是否认生命；因为\BibleQuote{这生命就是光}。\BibleRef{若 1:4} 祂称那些有小信德的人不是无信，而是伪善者，\BibleRef{玛 6:30} 他们身上刻着名字，却否认自己真正的信仰。但忠信者被称为仆人和朋友。因此，若人爱自己，他就爱主，并为得救而宣认，为拯救自己的灵魂。即使你因爱为邻人而死，并把救主视为我们的邻人（因为拯救的天主居于被拯救者近旁），你是因生命而选择死亡，为了自己而不是为了邻人受苦。岂不因此祂称他为弟兄？那因爱主而受苦的人，是为自己的救恩受苦；反之，那为自己的救恩而死的人，是为爱主而忍受。因为祂就是生命，在祂所受的苦难中，祂愿意受苦，好让我们因祂的苦难而活。\par

\BibleQuote{你们为什么称呼我主啊主啊，却不遵行我的话呢？} \BibleRef{路 6:46} 因为\BibleQuote{这百姓用嘴唇尊敬我，心却远离我}，\BibleRef{依 29:15} 他们是另一族人，依赖另一位，又自愿卖身给另一位；但那些遵行主命令的人，在每一行动中作证，做祂所愿，并恒常呼求主名；且以行动作证他们所信靠的那一位，他们是已把肉身同邪情私欲钉在十字架上的人。\BibleQuote{如果我们依靠圣神而生活，就也应依靠圣神而行事。} \BibleRef{迦 5:24} \BibleQuote{那随从肉身撒种的，必由肉身收获败坏；那随从圣神撒种的，必由圣神收获永生。} \BibleRef{迦 6:8}\par

但那些可怜的人，以为为主流血作证是最残暴的死亡，却不知道这样的死亡之门是真生命的开始；他们既不理解死后属于圣洁生活者的荣耀，也不理解对不义和不洁生活者的惩罚。我不仅依据我们的圣经这样说（因为几乎所有诫命都指示这些），他们甚至连自己的言论也听不进去。因为毕达哥拉斯学派的\Person{特阿诺}写道：\Quote{对恶人而言，生命确实是盛宴，他们行恶后便死去；若非灵魂不死，死亡将是天赐。}\Person{柏拉图}在他的\WorkTitle{斐多篇}中也说：\Quote{如果死亡是摆脱一切，}等等。因此，我们不应像\Person{埃斯库罗斯}的\WorkTitle{泰勒福斯}中所想的那样，\Quote{只有一条路通往冥府。}道路很多，引向那里的罪也很多。像那些不忠信者这样深度迷失的人，\Person{阿里斯托芬}恰当地将他们作为喜剧题材。\Quote{来吧，}他说，\Quote{你们这些隐晦生命的人，你们像叶子的族类，软弱，蜡像，阴影的族群，短暂，易逝，朝生暮死。}还有\Person{埃庇卡摩斯}：\Quote{人的这种本性是吹胀的皮囊。}救主也曾对我们说：\BibleQuote{心神固然切愿，但肉体却软弱。}\BibleRef{玛 26:41}\BibleQuote{因为随从肉性的思念，是与天主为仇，}宗徒解释说：\BibleQuote{它不服从天主的法律，而且也不能服从。凡属于肉性的人，都不能中悦天主。}接着进一步解释，免得有人像\Person{马西昂}那样认为受造物是邪恶的。\BibleQuote{但如果基督在你们内，身体固然因罪而死亡，但圣神却因正义而生活。}又说：\BibleQuote{如果你们随从肉性生活，必要死亡。因为我断定现时的苦楚，与将来在我们身上要显示的光荣，是不能较量的。如果我们与祂同受苦楚，也必与祂同受光荣，作为基督的同承继者。而且我们知道，天主使一切协助那些爱祂的人，就是那些按祂的旨意蒙召的人，获得益处。因为祂所预知的人，也预定他们与祂儿子的肖像相同，使祂在众多弟兄中作长子。祂所预定的，也召叫了他们；祂所召叫的，也使他们成义；祂所成义的，也使他们受光荣。}\par

你们看，为爱之故的殉道是受教导的。如果你愿意为得利的酬报而成为殉道者，你将再听到：\BibleQuote{因为我们得救是在乎盼望；只是所见的盼望不是盼望，谁还盼望他所见的呢？如果我们盼望那看不见的，就要耐心等待。}\BibleRef{罗 7:24}\BibleQuote{如果我们为义而受苦，}伯多禄说，\BibleQuote{我们是有福的。不要怕人的威吓，也不要惊慌；但要在你们心中尊主天主为圣；常常准备答复任何询问你们心中盼望理由的人，但要以温柔和敬畏的心，保持良心的纯洁，这样，在你们被毁谤的事上，那些诽谤你们在基督内的善行的人，必要惭愧。因为如果天主的旨意要你们因行善而受苦，总比因作恶而受苦更好。}但如果有人强辩说，软弱的肉身如何能抵抗那些邪恶权势的精力与鬼神呢？那么，让他知道：依靠全能的主，我们对抗黑暗的权势和死亡。\BibleQuote{你还在说话的时候，}祂说，\BibleQuote{看，我在这里。}请看那无往不胜的护佑者保护我们。\BibleQuote{所以，对于为试炼你们而燃起的烈火，不要感到惊奇，好像发生了异常的事；但既然你们有分于基督的苦难，就要欢喜，好使在祂光荣显现时，你们也能欢喜踊跃。如果你们因基督的名受辱骂，就有福了，因为光荣的神，即天主的神，安息在你们身上。}正如经上所记载：\BibleQuote{为了你的缘故，我们终日被杀死，被视为待宰的羊。然而，靠着爱我们的那位，我们在这一切事上，大获全胜。}\BibleRef{罗 8:36}\par

\Quote{你想从我心中探知的事，\newline{}你必不能探知，即使你从我的头顶到脚底\newline{}施加可怕的锯子，\newline{}即使你用锁链捆绑我。}\par

这是一位女子在悲剧中表现英勇地说的。而\Person{安提戈涅}，藐视\Person{克瑞翁}的宣告，勇敢地说：——\par

\Quote{发出这宣告的不是\Person{宙斯}。}\par

但是向我们发出宣告的是天主，祂必须被相信。因为心里相信，可使人成义；口里承认，可使人得救。因此圣经说：\BibleQuote{凡信祂的人，不致于羞耻。}\BibleRef{罗 10:10}相应地，\Person{西摩尼得斯}公正地写道：\Quote{据说美德居住在接近不可攀登的岩石之间，但她迅速穿越纯净之地。她也并非所有凡人的眼睛都能看见。不被心中烦恼的汗水浸透的人，无法攀上勇敢的顶峰。}\Person{品达}说：——\par

\Quote{青年们的忧虑思虑，伴随着辛劳旋转，\newline{}终将找到荣耀：及时地，他们的功绩\newline{}将在辉煌的以太中灿烂闪耀。}\par

\Person{埃斯库罗斯}也领会了这一思想，说：——\par

\Quote{辛劳者应得，\newline{}作为他劳苦的产物，来自神明的荣耀。}\par

\Quote{因为伟大的命运达到伟大的结局，}根据\Person{赫拉克利特}说：——\par

\Quote{有什么奴隶，不在乎死亡呢？}\par

\BibleQuote{因为天主赐给我们的，不是怯懦之神，而是大能、爱德和健全之神。所以不要以给我们的主作证为耻，也不要以我这为主被囚的人为耻，}他写信给\Person{弟茂德}。\BibleRef{弟后 1:7}这样的人将\BibleQuote{持守那善的，}按照宗徒的话，\BibleRef{罗 12:9}\BibleQuote{憎恨邪恶，拥有无伪的爱德；因为爱别人的就满全了法律。}\BibleRef{罗 13:8}那么，如果这位我们为之作证的天主，正如祂是希望的天主，我们承认我们的希望，奔向希望，\Quote{充满良善，满溢一切知识。}\par

印度智者对马其顿的亚历山大说：\Quote{你将人的身体从一个地方运到另一个地方。但你绝不能强迫我们的灵魂去做我们不愿做的事。火对人是最可怕的折磨，我们却蔑视它。}因此赫拉克利特宁愿选择一件事，即荣耀，胜过一切；并宣称\Quote{他容许民众像牲畜一样填饱自己。}\par

\Quote{因身体而有诸多劳苦，\newline{}我们为它建造了有顶的房屋，\newline{}并发现了如何挖掘白银、耕种土地，\newline{}以及所有其余我们以名称知晓的事物。}\par

对群众而言，这种虚妄的劳动是可欲的。但对我们，宗徒说：\Quote{我们知道，我们的旧人已与他同钉十字架，使罪身丧失机能，免得我们再作罪的奴隶。}\BibleRef{罗 6:6}宗徒岂非明白地接着这样说，以显示对群众在信德方面的轻视？\Quote{我以为天主把我们宗徒列在最后，好像定了死罪的人；因为我们成了一台戏，给世界、天使和世人观看。直到此时此刻，我们仍是又饥又渴，衣不蔽体，受人拳打，居无定所，并且亲手劳苦做工。被人咒骂，我们就祝福；被人迫害，我们就忍受；被人毁谤，我们就劝勉；我们成了世界上的污秽，万物中的渣滓。}同样，柏拉图在\Emph{\WorkTitle{理想国}}中也有此话：\Quote{义人纵然被拉上刑架，被挖去双眼，也会幸福。}\OriginalTerm{灵知者}{Gnostic}绝不会将终极目的置于现世生命，而在于永远快乐和蒙福，并作天主君王般的朋友。即使蒙受羞辱、流放、抄家，尤其死亡，也决不会失去自由和对天主显著的爱。\Quote{爱德凡事包容，凡事忍耐，}\BibleRef{格前 13:7}确信天主眷顾安排一切妥当。\Quote{我劝勉你们，}因此经上说，\Quote{你们要效法我。}救恩的第一步是伴随敬畏的训诲，由此我们避免错误；第二步是望德，因之我们渴望最善的事；但爱德，按合适的方式，藉由如今按知识的培育而达至成全。因为希腊人，我不明所以，将事件归因于无理性的必然，承认自己不情愿地屈服于它们。因此厄里庇得斯说：——\par

\Quote{夫人，请听我所宣告：\newline{}没有凡人不受劳苦；\newline{}他埋葬儿女，又生育其他，\newline{}他自己也会死去。凡人的苦难就是这样。}\par

然后他补充道：——\par

\Quote{我们必须承受本性中不可避免的事，并经历它们：\newline{}对凡人来说，没有一件必然之事是可怕的。}\par

对那些追求成全的人，提出了理性的\OriginalTerm{灵知}{gnosis}，其基础是\Quote{\OriginalTerm{神圣三德}{sacred Triad}}。\Quote{信德、望德、爱德；其中最大的是爱德。}的确，\Quote{凡事都可行，但不全有益；}宗徒说：\Quote{凡事我都可行，但不全造就人。}\BibleRef{格前 10:23}并且，\Quote{各人不可只求自己的益处，也该求别人的益处，}\BibleRef{格前 10:24}以便能同时实行和教导，建立并建造。因为\Quote{大地和其中的万物属于主，}这是公认的；但软弱者的良心被扶持。\Quote{我说良心，不是自己的，而是他人的；为什么我的自由被他人的良心论断？若我因恩宠而分享，为什么因我感谢的事而受毁谤？所以，你们或吃或喝，或无论做什么，一切都要为光荣天主而做。}\Quote{因为我们虽然在血肉中行走，却不按血气作战；因为我们作战的武器不是属血肉的，而是藉天主有力摧毁堡垒，摧毁思想及一切高举起来反对认识基督的傲慢事物。}配备这些武器，\OriginalTerm{灵知者}{Gnostic}说：主啊，赐予机会，接受证明；让这件可怕的事过去；因我对祢的爱，我藐视危险。\par

\Quote{因为唯独在人类事物中，\newline{}德行不接受外来的报酬，\newline{}而是以自身作为其劳苦的奖赏。}\par

\Quote{所以，你们既是天主的选民，圣洁而蒙爱的人，要穿上怜悯的心肠、仁慈、谦卑、温柔、忍耐。在这一切之上，要有爱，它是成全的纽带。并要让天主的平安在你们心中作主；你们蒙召也是为此，成为一个身体；且要心存感谢。}你们这些仍在身内的人，像古时的义人一样，享受灵魂的无情与安宁。\par

% structure: p0071 source=h2 target=subsection reason=relative-heading-rank; base=h2

\subsection{第八章 妇女与男子一样，奴隶与自由人一样，皆为殉道荣冠的候选者}

既然不仅伊索尼亚人、马其顿人和拉刻代蒙人在遭受酷刑时忍受了，正如埃拉托斯特尼在其著作\WorkTitle{论善与恶}中所说；而且埃利亚的芝诺，在被迫泄露秘密时，抗拒酷刑，什么也没有承认；他在临终时咬断舌头，吐向暴君，有人称其为尼阿尔科斯，有人称其为德慕洛斯。毕达哥拉斯学派的特奥多图斯也做了类似的事，还有拉克得斯的朋友保禄，正如佩加蒙的提摩太在其著作\WorkTitle{论哲学家的刚毅}中和阿凯库斯在\WorkTitle{伦理学}中所说。还有罗马人波斯图穆斯，被佩乌克提翁俘虏时，没有泄露一个秘密；而是把手放在火上，像放在铜块上一样，一动不动，面部肌肉毫无变化。我略去阿那克萨耳科斯的事例，他喊道：\Quote{朝装着阿那克萨耳科斯的口袋捣吧，因为你捣的不是阿那克萨耳科斯，}当时他正被暴君命令用铁杵捣碎。因此，无论是幸福的希望还是对天主的爱，都不会遭受厄运，而是保持自由，即使被掷入最凶猛的野兽中或吞灭一切的烈火中；即使被暴君的酷刑折磨。由于它依赖于天主的恩宠，它高高升起，不受奴役，将身体交给那些只能触及其身体的人。据说有一个蛮族，不为哲学所累，每年选一位使者到英雄匝莫尔克息斯那里。匝莫尔克息斯是毕达哥拉斯的一位门徒。那么，被认为最有价值的人被处死，使那些实践了哲学却未被选中的人感到悲痛，因为他们被认为不配得到一种幸福的效劳。\par

因此，教会充满了这样的人，既有贞洁的妇女，也有男子，他们一生都在默想那唤醒人归向基督的死亡。因为一个人的生活若按我们的方式塑造，即使没有学问，也可以从事哲学思考，无论是蛮族人，还是希腊人，无论是奴隶，还是老人、少年或妇女。因为自制是选择它的所有人类所共有的。我们承认，每个种族都有相同的本性，相同的德行。就人性而言，女人并不拥有一种本性，男人展示另一种，而是相同的：德行也是如此。因此，如果节制和正义以及被视为随之而来的任何品质是男人的德行，那么只有男人才有德，而女人则为放纵和不义。但这样说甚至令人反感。因此，女人应当像男人、奴隶和自由人一样，实践节制、正义和每一种其他德行；因为同一本性拥有同一德行，这是恰当的结论。我们不是说女人的本性等同于男人的本性，因为她是女人。因为无疑，理当认为他们之间应该存在某种差异，基于此，一个是男性，另一个是女性。因此，我们说话和分娩属于女人，因为她是女人，而不是因为她是人。但若男女之间没有差异，两者就会做同样的事，受同样的苦。既然在灵魂方面是相同的，她将达到同样的德行；但由于身体构造的特殊差异，她被注定生育和操持家务。\BibleQuote{因为我要你们知道，}宗徒说，\BibleQuote{一切男人之首是基督，女人之首是男人：因为男人不是出于女人，而是女人出于男人。并且女人没有男人，男人也没有女人，在主内。}因为我们说男人应当节制，胜于快乐；同样也认为女人应当节制，并练于对抗快乐。\BibleQuote{但我却说：你们随从圣神而行，就绝不会满足肉身的私欲，}宗徒的命令劝诫说，\BibleQuote{因为肉身私欲相反圣神，圣神也相反肉身；这些都是彼此相敌}，（不是善对恶，而是有益的斗争），他因此加上，以致你们不能做你们愿意做的事。\BibleQuote{肉身的行为是明显的：就是淫乱、不洁、放荡、拜偶像、邪术、仇恨、纷争、嫉妒、忿怒、争吵、分党、异端、妒忌、醉酒、宴乐，以及类似的事；我预先告诉你们，如同以前说过的一样：行这样事的人，决不能继承天主的国。然而圣神的果实是：仁爱、喜乐、平安、忍耐、良善、温和、忠信、柔和。}他认为，罪人为\Quote{肉身}，义人为\Quote{圣神}。此外，应当采取刚毅，以产生信心和坚忍，以致\BibleQuote{有人打你的右脸，连左脸也转过来给他；有人要拿你的里衣，连外衣也让他拿去，}强力抑制愤怒。因为我们不像亚马孙人那样训练妇女在战争中刚毅；因为我们希望甚至男子也要和平。我听说萨尔马提亚妇女作战不亚于男子；此外，萨卡妇女也向后射箭，像男子一样佯装逃跑。我也知道，伊比利亚附近的妇女从事男子的工作和劳动，即使接近分娩也不停下手头的任务；甚至在产痛挣扎中，生下孩子后，抱起婴儿带回家。此外，女性与男性一样管理家庭、狩猎和放牧：——\par

\Quote{母犬克瑞萨在公鹿的踪迹上疾跑。}\par

因此，妇女应当与男子同样地从事哲学思考，尽管男子在各方面更优越，除非他们变得女性化。那么，对于全人类来说，纪律和德行是必需的，如果他们追求幸福的话。欧里庇得斯有时这样写，有时那样写，多么轻率！有一次：\Quote{因为每个妻子都在丈夫之下，即使最优秀的男人娶了名声好的女子。}另一次：——\par

\Quote{贞洁的妇女是她丈夫的奴仆，\newline{}而不贞洁的妇女因愚妄藐视她的配偶。\newline{}……因为没有比这更好、更卓越的事了，\newline{}即是当夫妻二人操持家务时，\newline{}心意和谐。}\par

所以统治权是头。如果\BibleQuote{主是男人的头，而男人是女人的头，} 男人\BibleQuote{既是天主的肖像和光荣，就是女人之主。} 因此，在\WorkTitle{厄弗所书}中也记载：\BibleQuote{你们要怀着敬畏基督的心，彼此服从。妻子要服从自己的丈夫，如同服从主，因为丈夫是妻子的头，如同基督是教会的头，又是这身体的救主。丈夫要爱妻子，如同基督爱了教会。所以丈夫也应当爱妻子如同自己的身体；爱妻子就是爱自己。因为从来没有人恨过自己的身体。} \BibleRef{弗 5:21} 在\WorkTitle{哥罗森书}中也说：\BibleQuote{妻子要服从自己的丈夫，这在主内是合宜的。丈夫要爱妻子，不要对她们苛刻。子女要凡事听从父母，因为这是主所喜悦的。父亲不要惹子女发怒，免得他们灰心。仆人要凡事服从肉身的主人，不要只在眼前服事，像讨人喜欢的，而要心存诚实，敬畏主。无论做什么，都要从心里做，像是给主做的，不是给人做的，因为你们知道，从主那里必得基业作为赏报；你们服事的是主基督。行不义的，必受不义的报应，主并不偏待人。主人要给你们仆人正义和公道，因为你们知道你们也有一位主在天上，在那里既没有希腊人，也没有犹太人，没有受割礼的，也没有未受割礼的，没有化外人，没有西古提人，没有为奴的，也没有自主的；只有基督是一切，也在一切之内。} 地上的教会是天上的教会的肖像，正如我们也祈祷\BibleQuote{愿天主的旨意行在地上如同行在天上。} \BibleRef{玛 6:10} \BibleQuote{所以，你们要穿上怜悯的心肠、仁慈、谦卑、柔和、忍耐；彼此担待，互相宽恕，若有人与谁有争执，如同基督宽恕了你们，你们也要这样。在这一切之上，还要加上爱德，因为爱德是完美的联系。并要让天主的平安在你们心中做主，你们被召同归一体，也要存感激的心。} 因为屡次引用同一段圣经来使\Person{马西昂}羞愧，并没有妨碍；或许他能被说服而悔改，因为他会认识到，信友应当感谢造物主\Person{天主}，祂召叫了我们，并在肉身中宣讲了福音。从这些考虑，信德的合一就清楚了，也显示出谁是成全的人。因此，即使有些人勉强，并尽其所能地抵抗，即使丈夫或主人以惩罚相威胁，家仆和妻子都会修习哲学。此外，自由人，即使受到暴君以死亡相威胁，被带到法庭前，所有财产都濒临危险，也绝不会放弃虔诚。与恶丈夫同住的妻子，或有了坏父亲的儿子，或有了坏主人的家仆，也绝不会不能高尚地持守德行。正如人为德行、为自由、为自己而死是高尚的，同样，女人这样做也是高尚的。因为这不专属于男性的本性，而是属于善的本性。因此，老人、青年和仆人都会忠信地活着，如果需要，也忠信地死去；这将是因死亡而得生命。所以我们知道，孩子、妇女和仆人常常违背他们的父亲、主人和丈夫的意愿，达到了最高的卓越境界。因此，那些立志虔诚生活的人，当有人似乎对他们施加强制时，仍然应当表现出乐意。但我认为，他们更应当表现出热切，并以非凡的精力奋斗，以免被胜过而放弃了最好且最不可或缺的劝告。因为，我认为，这无法比较：是作全能者的追随者更好，还是选择魔鬼的黑暗更好。因为凡我们因他人而做的事，应当常常去做，努力尊重那些为他们的缘故而应当做这些事的人，将给予他们的满足视为我们的准则；但那些为我们自己而非为他人而做的事，应当同样认真去做，无论它们是否可能取悦某些人。如果某些无关紧要的事获得了如此尊荣，以至于似乎值得采纳，尽管违背某些人的意愿；那么，德行更应当被我们视为值得争取的，同时只看什么是可正确地做的事，无论别人看起来是否好。然后，\Person{厄比鸠鲁}写信给\Person{梅诺厄斯}说：\Quote{不要让青年人延迟修习哲学，也不要让老年人对修习哲学感到厌倦；因为没有一个人是未到年龄或超过年龄而不再关心他灵魂的健康。那说修习哲学的时候未到或已过的人，就像说幸福的时候未到或已过的人一样。因此，青年和老年都应当修习哲学：前者，为了在变老时，因过去所积累的惠益而在好事上变得年轻；后者，为了他因对未来无所恐惧而同时既年轻又老迈。}\par

% structure: p0078 source=h2 target=subsection reason=relative-heading-rank; base=h2

\subsection{第九章 基督关于殉道的言论}

关于殉道，主已明确说过，我们将散见于各处的记载收集起来。\BibleQuote{但我告诉你们：凡在人面前承认我的，人子在天使面前也要承认他；凡在人面前否认我的，我在天使面前也必否认他。}\BibleRef{路 12:8} 凡在这淫乱罪恶的世代以我和我的话为耻的，人子在他父的光荣中同圣天使来临的时候，也要以他为耻。所以，凡在人面前承认我的，我在我天上的父面前也必承认他。\BibleRef{玛 10:32} \BibleQuote{当人押送你们到会堂、官长和有权柄者面前时，不要预先思虑怎样申辩，或说什么话；因为圣神要在那一刻教给你们该说的话。}\BibleRef{路 12:11} 在解释这段话时，华伦提努派中最杰出的\Person{赫拉克利翁}明确地说：\Quote{有一种宣认是藉着信德和行为，另一种是藉着言语。用言语、在官长面前所作的宣认，多数人以为这是唯一的宣认。这并不正确：伪善者也能以此宣认。但这句话也并非普遍适用，因为所有得救者都曾用言语宣认并离世，其中有\Person{玛窦}、\Person{斐理伯}、\Person{多默}、\Person{肋未}和许多其他人。口唇的宣认不是普遍的，而是个别的。但祂现在所指的乃是普遍的宣认，即与信仰相称的行为和行动。在这宣认之后，若有必要且合理，才有个别的宣认，即当着官长。因为凡先以心志宣认的人，才能以言语正确宣认。祂恰当地对宣认者用了\InnerQuote{在我内}的表述，对否认者用了\InnerQuote{我}的表述。因为那些人虽以言语承认祂，却在行为上否认祂，并未在行为上承认祂。惟独那些按祂的宣认和行为而生活的人，才\InnerQuote{在祂内}宣认；在他们内，祂也宣认，因为祂被包含在他们内、被他们持守。因此\InnerQuote{祂决不能否认自己}。而不在祂内的人则否认祂。因为祂没有说\InnerQuote{凡否认我}在我内，而说\InnerQuote{我}。因为没有一个在祂内的人会否认祂。\InnerQuote{在人面前}的表述同样适用于得救者和外邦人：以行为面对前者，以言语面对后者。因此他们绝不能否认祂。而不在祂内的人则否认祂。} 以上是\Person{赫拉克利翁}的话。在其他方面，他似乎与我们在本节中的观点一致；但他没有注意到这一点：如果有些人未以行为和生命\Quote{在人面前承认基督}，却显示他们心里相信，在法庭上用口承认祂，并在遭受酷刑至死时不否认祂。心志既被承认，尤其至死不变，便除去了由肉身欲望产生的一切情欲。因为可以说，在生命终结时，突然有了行动上的悔改，并在言语的见证中产生了对基督的真切宣认。但如果父的圣神在我们内作证，我们怎能还是伪善者呢？那些人被认为只用言语作证。但若合适，有些人会得到申辩的机会，为藉着他们的见证与宣认，使所有人获益——教会内的人得以坚定，外邦人中那些致力于寻求救恩的人感到惊奇并被引向信仰，其余人则惊讶不已。因此，宣认是绝对必要的，因为它在我们能力之内。但为我们的信仰申辩并非普遍必要，因为这不取决于我们。\Quote{唯有坚忍到底的，才可得救。}因为智慧人中谁不愿在天主内为王，甚至服务呢？所以有些人\Quote{承认认识天主}，如宗徒所说，\Quote{但在行为上否认祂，是可憎的、悖逆的，并在各样善事上被摒弃。}\BibleRef{铎 1:16} 这些人尽管只这样承认，最终将做一件善事：他们的见证似乎是光荣地洗净罪恶。例如，\WorkTitle{牧者}说：\Quote{如果你的心变得纯洁无瑕，你必逃脱野兽的力量。}主自己也说：\Quote{撒但想要筛你们像筛麦子一样，但我已为你祈祷了。}\BibleRef{路 22:31} 因此，唯独主为洁净那些谋害祂、不信祂的人，\Quote{饮了这杯}。宗徒效法祂，为使成为真正的诺斯替信徒和成全的人，为他们所建立的教会受苦。同样，追随宗徒足迹的诺斯替信徒也应当无罪，并因爱主而爱自己的弟兄，以致若时机召唤，他们能毫无绊倒地为教会忍受苦难，\Quote{饮这杯}。那些在生活中以行为、在法庭上以言语作见证的人，无论怀着希望还是猜疑恐惧，都胜过那些仅以口宣认救恩的人。但若有人上升到爱德，他便是真正有福且真实的殉道者，已藉着主完美地宣认了诫命和天主；他因爱主而承认主为兄弟，为天主完全舍弃自己，当人被索要时，愉快而慈爱地将自己交出，如同交还一笔寄存物。\par

% structure: p0080 source=h2 target=subsection reason=relative-heading-rank; base=h2

\subsection{第十章 那些自献于殉道者的人受谴责}

再者，当祂说：\BibleQuote{他们在这一城迫害你们，你们就逃到另一城去}\BibleRef{玛 10:23}，祂并非劝人逃跑，好像迫害是邪恶的事；也不是命令他们藉逃跑躲避死亡，仿佛害怕死亡，而是愿意我们既不作任何恶事的主动者，也不作任何恶事的同谋，不论是对我们自己，还是对迫害者和凶手。因为祂在某种程度上命令我们照顾自己。但违背者乃是鲁莽和冒失的。如果杀害天主的人得罪天主，那么主动站到审判台前的人也使自己成为死亡的责任者。同样，那不躲避迫害，却因胆敢而自投罗网的人也是如此。这样的人，在他能力所及的范围内，成了迫害者罪行的同伙。如果他还加以挑衅，他就是完全有罪的，是在挑衅野兽。同样，如果他提供任何冲突、惩罚、报复或敌意的起因，他就给了迫害的机会。因此，我们被命令不要依恋此生任何事物，反而要给那拿走我们外衣的人\BibleQuote{连我们的内衣也给他}，不仅使我们能继续保持超然、不受过度情感束缚，而且使我们不致因报复而激怒迫害者反对我们自己，并煽动他们亵渎那名。\par

% structure: p0082 source=h2 target=subsection reason=relative-heading-rank; base=h2

\subsection{第十一章　反驳\Quote{既然天主眷顾你们，为何你们仍受苦}的异议}

但是，他们说：如果天主眷顾你们，你们为何被迫害和处死？祂把你们交给这个了吗？不，我们不认为主愿意我们卷入灾祸，而是祂先知性地预言了将要发生的事——我们将为祂的名被迫害、被屠杀、被钉十字架。所以，并非祂愿意我们被迫害，而是祂藉着对要发生之事的预言预先告诉我们将要遭受什么，训练我们忍耐，并应许给继承产业，尽管我们不是独自受罚，而是与许多人一同受罚。但那些人，据说，因为是作恶者，所以公义地受罚。相应地，他们不情愿地见证了我们的正义，因为我们为了正义而受不义之罚。但法官的不义并不影响天主的眷顾。因为法官必须是他自己意见的主人——不是像无生命的机器被线牵着，只靠外部原因才运动。因此，他根据他的判断被审判，我们也根据我们对可欲之事的选择和我们的忍耐而被审判。虽然我们没有行不义，法官却看我们为行不义者，因为他既不认识也不愿认识我们，而是被无根据的偏见所影响，因此他也要受审判。所以，他们迫害我们，并不是认为我们是作恶者，而是以为仅仅因我们是基督徒，以这样的方式生活和劝勉他人采纳同样生活，就犯了悖逆生命的罪。\par

但你们为何被迫害时不被帮助？他们说。对我们来说，被死亡释放到主那里去，从而经历生命的改变，如同从一个生命时期到另一个时期，这有什么损害呢？如果我们正确思考，我们应当感谢那些提供迅速离世途径的人，如果我们是出于爱而作见证；否则，我们在众人眼中就成了卑鄙的人。倘若他们也认识真理，所有人都会奔向那条路，就不会有选择。但我们这信仰，作为世界的光，谴责不信。\Quote{即使阿尼图斯和梅勒图斯杀了我，他们丝毫不能伤害我；因为我认为正确者被错误者伤害是不合宜的。}（苏格拉底语）因此，我们每个人都可以有信心地说：\BibleQuote{主是我的助佑；我不畏惧：人能对我做什么？} \BibleQuote{义人的灵魂在主手中，灾祸不能触及他们。}\par

% structure: p0085 source=h2 target=subsection reason=relative-heading-rank; base=h2

\subsection{第十二章　驳斥巴西里德对殉道的观念}

\Person{巴西里德}在《\OriginalTerm{释经书}{Exegetics}》第二十三卷中，关于那些因殉道而受惩罚的人，这样表达自己的观点：\Quote{因为我说这话，凡无意中在其他事上违反诫命而遭受上述苦难的人，是被那带领他们的善心引至这善终，但在其他方面被控告；因此他们不是作为公认的恶行而受罚，也不是像奸夫或杀人犯那样受谴责，而是因为他们是基督徒；这会使他们得到安慰，以致他们不认为自己在受苦。而若有一个从未犯罪的人遭受苦难——这是罕见的情形——他也不会因权势的阴谋而受苦，而是像那似乎未犯罪的孩童那样受苦。}接着他又补充：\Quote{因此，正如那先前未犯罪、自身也无实际过犯的孩童，只拥有那犯罪的潜能，在遭受苦难时却获益，从许多困难中得到好处；同样，即使一个成全的人实际未犯罪，在他忍受苦难时，也像孩童一样受苦。由于他内有罪恶的原则，但不利用犯罪的机会，所以他不犯罪；因此不可将他视为未犯罪。因为那图谋奸淫的人，即使未成功，也是奸夫；那图谋杀人的人，即使未能行凶，也是凶手；同样，若我看见一个我所指认的无罪之人受苦，即使他没有作恶，我也应因他存心犯罪而称他为恶人。因为无论什么，我都宁可说，也不愿称\OriginalTerm{上智安排}{Providence}为恶。}然后，他接着明确论及主，如同论及人：若你离开这一切观察，而用（或许代表某些人的话）来羞辱我说：\InnerQuote{那么这人犯了罪；因为这人受苦了}——若你允许，我会说：\InnerQuote{他没有犯罪；而是像一位受苦的孩童。}若你更加逼问，我会说：\InnerQuote{你所指的那人是人，但天主是公义的：}\Quote{因为没有一人是洁净的，}正如某人所说的，\InnerQuote{脱离污秽。}\BibleRef{约 14:4}但\Person{巴西里德}的假设说，灵魂在另一生命中已犯了罪，在此生受惩罚——被选的灵魂因殉道而得尊荣，其他的则受适当的惩罚净化。这怎能成立，既然宣认并受罚与否取决于我们自己？因为对于将来否认的人，按\Person{巴西里德}所说的\OriginalTerm{上智安排}{Providence}就不存在了。那么我要问他，对于一个被捕的宣认者，他是否会按\OriginalTerm{上智安排}{Providence}而宣认并受罚？因为若否认，他就不会受罚。但如果为了逃避和规避惩罚此人之必要，他会说那些将否认之人的毁灭是出于\OriginalTerm{上智安排}{Providence}，那么他就是违背自己的意愿成为殉道者。那么那为见证者所存留于天上的极荣耀的赏报又如何是见证的赏报呢？若\OriginalTerm{上智安排}{Providence}不允许罪人犯罪之极致，则在两种情况中都是不义的：既不拯救那为正义而受罚的人，又拯救那意图作恶（他已在意志上完成了恶行，但\OriginalTerm{上智安排}{Providence}阻止了行为）并偏袒罪人。而又何等亵渎，竟将魔鬼神化，胆敢称主为有罪的人！因为魔鬼试探我们，他知道我们是什么，却不知道我们能否坚持，他企图使我们背弃信仰，并企图使我们服从他。这一切之所以允许他，部分是由于拯救我们的必要——我们因诫命而有机会自救；部分是为使那试探而未成功的魔鬼蒙羞；并为巩固教会成员和钦佩所显毅力的良心。但如果殉道是作为惩罚的报应，那么信仰和教义——殉道由此而来——便是惩罚的共犯——还有什么比这更荒谬的呢？至于这些教义——灵魂是否转换到另一身体，以及魔鬼的事——将在适当时候提及。但现在，对已说过的话，让我们补充以下内容：如此，对殉道前所犯罪过的报应信仰何在？为真理而受迫害、忍受的爱主之情何在？对宣认者的赞美和对否认者的谴责何在？正当的行为、克制私欲、不恨任何受造物又有何用？但如\Person{巴西里德}自己所说，假定天主所启示旨意的一部分是爱一切事物，因为一切事物都与整体相关；另一部分是\Quote{不贪恋任何事物}，第三部分是\Quote{不恨任何事物}；那么按照天主的旨意，这些也将成为惩罚，而这是亵渎的想法。因为主并不是因圣父的旨意而受苦，那些受迫害的人也不是因天主的旨意而受迫害；因为二者必居其一：要么因天主旨意的迫害是好事，要么那些下令并施加迫害的人是无罪的。但没有什么能离开宇宙之主的旨意。因此只能说，这些事的发生未受天主的阻止；只有这样才能保全天主的眷顾和良善。因此我们不可认为祂主动产生苦难（切勿如此想！）；但必须确信祂不阻止那些制造苦难者，而是将仇敌的罪行导向善：\Quote{我必将拆毁这墙，使它成为被人践踏的。}\BibleRef{依 5:5}\OriginalTerm{上智安排}{Providence}是一门训练的技艺；对其他人是为各人自己的罪，对主和祂的宗徒则是为我们的罪。对此，神圣的宗徒说：\BibleQuote{因为天主的旨意就是要你们成圣，远离淫乱；要你们每一个人晓得怎样用圣洁和尊贵保守自己的身体；不放纵邪淫的情欲，像那些不认识主的外邦人；谁也不可在这事上越分，欺压自己的弟兄；因为主在一切这样的事上必报应，正如我们预先告诉你们、并且证明过的。因为天主召我们不是为不洁，而是为圣洁。所以那轻视的，不是轻视人，而是轻视那将圣神赐给你们的天主。}\BibleRef{得前 4:3}因此，主没有从我们的这种圣化中被禁止。所以，若他们中有人回答说，殉道者是为本次降生之前所犯的罪受罚，并且他将在此生再次收获其行为的果实，因为这是（神圣管理）的安排，我们将问他，报应是否由\OriginalTerm{上智安排}{Providence}而来。若它不是出于神圣管理，那么净化之安排就瓦解了，他们的假设也倒塌了；但若净化出于\OriginalTerm{上智安排}{Providence}，则惩罚也出于\OriginalTerm{上智安排}{Providence}。然而\OriginalTerm{上智安排}{Providence}，虽可说始于统治者，却由宇宙之神植入于实体同其起源之中。既然如此，他们必须承认：要么惩罚不公正，那些谴责和迫害殉道者的人做得对；要么迫害也是出于天主的旨意。那么劳苦和恐惧，正如他们所说，并非如铁锈之于铁那样偶然发生，而是出于灵魂自身的意志降临。关于这些点还有很多可说的，将在适当时候留待以后探讨。\par

% structure: p0087 source=h2 target=subsection reason=relative-heading-rank; base=h2

\subsection{第十三章 驳斥\Person{华伦提努}关于废除死亡的奇谈}

\Person{华伦提努}在一篇讲道中写道：\Quote{你们原是不死不灭的，是永生的子女，死亡要分配给你们，好让你们挥霍浪费它，好让死亡在你们内并藉着你们而死去；因为当我们消融世界而你们自己并未消融时，你们就掌管一切受造物和一切朽坏。} 因为他也和\Person{巴西里德}一样，认为有一类人是生来就得救的，这个不同的种族从上方来到我们这里是为废除死亡，而死亡的根源是这世界的创造者所为。因此他也这样解释那圣经：\BibleQuote{谁也不会看见天主的面而仍能活着}，仿佛祂是死亡的原因。关于这位天主，他在以下表达中作了那些暗示：\Quote{形象如何不如生面，世界也如何不如活的\OriginalTerm{永世}{Æon}。那么，形象的原因是什么？是面容的威严，它向画师展示形象，以他的名字受尊敬；因为形式并不完全符合生命，但名字弥补了肖像中所缺的。天主的不见性也有助于被造之物的信德。} 因为那被称为天主和圣父的创造者，他称之为\Quote{画师}和\Quote{智慧}，所形成的万物就是祂的肖像，是为荣耀那看不见的一位；因为从一对出来的事物是互补的，而从单一出来的事物是肖像。但是，既然所见到的并非祂的一部分，灵魂来自中间体，这中间体是不同的；这就是不同的灵的灵感，一般来说，就是吹入灵魂的东西，即灵的肖像。一般来说，关于那按肖像受造的创造者，他们说是由\WorkTitle{创世纪}中关于人类起源的一个可感知的形象所预言的；而他们将相似性归于自己，教导说那不同的灵被加入了，而这创造者并不知晓。那么，当我们论述那在律法、先知和福音中所宣告的天主的统一性时，我们也会讨论这个；因为这是个至高的主题。但我们必须前进到急切的问题。如果为了消除死亡，那特殊的种族已经来到，那么就不是\Person{基督}消除了死亡，除非祂也被说成是与他们同质。如果祂消除死亡是为了不触及那特殊的种族，那么就不是这些创造者的对手——他们按自己教义的原则，将来自上方的生命吹入他们中间之灵的肖像——消除了死亡。但若他们说这是经由祂的母亲发生的，或若他们说他们与\Person{基督}一起对抗死亡，就让他们承认他们的秘密教义：他们竟敢攻击创造者的神能，通过纠正祂的创造，仿佛他们更优越，努力拯救那祂无法从朽坏中救出的生命肖像。那么，主就会高于天主创造者；因为儿子绝不会与父亲争斗，尤其是在众神之间。但是，关于万物的创造者、全能的主，就是圣子的圣父这一点，我们推迟到这些点的讨论中，在那里我们承担了驳斥异端的任务，表明独有祂是祂所宣告的天主。\par

但是，宗徒写信给我们论及忍受苦难时，说：\BibleQuote{这是出于天主，因为祂恩赐你们为\Person{基督}的缘故，不但信祂，而且为祂受苦；有同样的斗争，就如你们在我身上所见的，现在也在我身上所听到的。所以，如果在\Person{基督}内有任何劝勉，有任何爱心的安慰，有任何圣神的相通，有任何慈悲怜悯，你们就要使我的喜乐满足，要同心合意，有同样的爱，同心一致，思念同一件事。如果他在信德的祭献和服务上被奠祭，喜乐并欢庆} 与\Person{斐理伯人}同乐，宗徒向他们说话，称他们为\InnerQuote{同乐的同享}，\BibleRef{斐 1:7} 他怎么说他们同心并有同一个灵魂呢？同样，论及\Person{弟茂德}和他自己时，他写道：\BibleQuote{我没有别人像他那样同心，会真诚地关心你们的事。因为各人寻求自己的事，不寻求\Person{耶稣基督}的事。} \BibleRef{斐 2:20}\par

那么，上述之人不要以责备的方式称我们为\Quote{自然的人}（\OriginalTerm{自然的人}{\GreekText{ψυκικοί}}），弗里吉亚人也不可；因为他们现在也称那些不投身于新预言的人为\Quote{自然的人}（\OriginalTerm{自然的人}{\GreekText{ψυκικοί}}），我们将在关于\Quote{预言}的论述中与他们讨论。因此，完美的人应当操练爱德，并由此急赴神圣的友谊，出于爱而满全诫命。爱仇敌并非意味着爱邪恶、不敬、奸淫或偷窃；而是爱那窃贼、不敬者、奸淫者，并非就他犯罪而言，也不是就他玷污人之名的行为而言，而是就他是一个人、是天主的化工而言。的确，罪是一种活动，而非存有；因此它不是天主的化工。罪人被称为天主的仇敌——即，是那些他们不服从的诫命的仇敌，正如服从的人成为朋友，前者因共融而得名，后者因疏离而得名，这是自由意志的结果；因为没有仇敌和罪人，就没有仇视和罪。而\Quote{一无所贪}的诫命，并非如同那些可贪求之物不属于我们，也不是教导我们不要有欲望，正如那些教导造物主不同于第一位天主的人所认为的，也并非如同受造物是可憎和恶的（因为此类意见是不敬的）。但我们说，世界之物\Emph{不是我们自己的}，并非如同它们是怪物，也并非如同它们不属于天主、宇宙的主，而是因为我们不永远存留于其中；从拥有方面说，它们不是我们的，且相继从一人传递到另一人；但就使用而言，它们属于我们这些为此受造的人，只要有必要继续与之相伴。因此，依照自然的欲望，不被允许的事物要正当使用，避免一切过度和不当的情欲。\par

% structure: p0091 source=h2 target=subsection reason=relative-heading-rank; base=h2

\subsection{第十四章 对众人，甚至对仇敌的爱}

仁慈何其伟大！经上说：\BibleQuote{爱你们的仇敌，}\BibleQuote{祝福那些诅咒你们的人，为那些恶意对待你们的人祈祷，}\BibleRef{玛 5:44}以及类似的话；又加上：\BibleQuote{好使你们成为你们在天上之父的子女}，暗示与天主的相似。又说：\BibleQuote{赶快与你的对头和好，当你还在路上和他一起时。}\BibleRef{玛 5:25}对头并非如某些人所想的而是身体，而是魔鬼，以及那些模仿他的人，他们以人的身份与我们同行，在现世生活中效仿他的行为。那么，那些承认自己属于基督，却发现自己身处魔鬼的作为之中的人，遭受最敌意的对待是不可避免的。因为经上记载：\Quote{免得他把你交给审判官，审判官把你交给撒殚国度的差役。}\BibleQuote{因为我深信，无论是死，}通过迫害者的攻击，\BibleQuote{是生，}在这世界，\BibleQuote{是天使，}那背叛的天使，\BibleQuote{是掌权者}（而撒殚的权能是他所选择的生命，因为属于他的黑暗的掌权者和权势就是如此），\BibleQuote{是现在的事，}我们在有生之年存于其中，如同士兵所怀的希望和商人的盈利，\BibleQuote{是高天，是深渊，或是任何其他受造之物，}由于人所特有的能力——都不能反对那按自由意志行事之人的信德。\Quote{受造之物}与活动同义，是我们的工作，而这种活动\BibleQuote{不能使我们与天主在基督耶稣我们的主内的爱相隔绝。}\BibleRef{罗 8:38}你得到了关于真知者殉道的概要。\par

% structure: p0093 source=h2 target=subsection reason=relative-heading-rank; base=h2

\subsection{第十五章 论避免冒犯}

\Quote{我们知道我们都有知识}——在普通事物上的普通知识，以及知道有一位天主的知识。因为他写信给信徒；因此他接着说，\Quote{但\OriginalTerm{知识}{\GreekText{γνῶσις}}并不在众人内，}只传递给少数人。还有人说，关于祭偶像之物的知识并没有在众人中传播，\Quote{以免我们的自由成为软弱人的绊脚石。因为你的知识，那软弱的就要丧亡。}若他们说：\Quote{凡在市场上出卖的，是否应当购买？}又以诘问的口气加上：\Quote{不必追问问题，}\BibleRef{格前 10:25}仿佛等价于\Quote{追问问题，}他们给出了可笑的解释。因为宗徒说：\Quote{凡在市场上出卖的其他东西，都不必追问问题而买，}但宗徒们的公函中提及的例外，是\Quote{在圣神的同意下，}这写在宗徒大事录中，并由保禄亲手传给信徒。因为宗徒们表明\Quote{他们必须戒避祭偶像之物、血、勒死的牲畜和奸淫；戒避这些，他们便能行得好。}因此，宗徒所表达的则是另一回事：\Quote{我们难道没有权力吃、喝吗？我们难道没有权力带领一位姊妹、一位妻子，如同其余的宗徒、主的弟兄和\OriginalTerm{刻法}{Cephas}一样吗？但我们没有使用这权力，}他说，\Quote{反而忍受一切，以免我们对基督的福音造成阻碍。}即，通过背负重担，当需要为一切事不受约束时；或者成为那些想实行节制的人的榜样，不鼓励彼此贪婪地吃摆在我们面前的食物，也不贸然与女人来往。尤其对于受委托执行如此管理的人来说，必须向门徒展示纯洁的榜样。\Quote{因为我虽从众人中得了自由，却使自己成为众人的仆人，}经上说，\Quote{为要赢得众人。凡在竞赛中争取优胜的，在一切事上都要节制。}\BibleRef{格前 9:19}\Quote{然而大地是属于主的，并其中的一切。}\BibleRef{格前 10:26}因此，为了良心的缘故，我们应当戒避所应戒避的。\Quote{我说的是良心，不是他自己的，}因为它具备知识，\Quote{而是别人的，}免得他受到不良训练，并在无知中模仿他所不知道的，而成为藐视者而非强心之人。\Quote{因为我的自由为什么被人凭别人的良心来判断呢？我若因恩宠是分享者，为什么我因感谢的事而受毁谤呢？你们无论做什么，都要为光荣天主而做，}\BibleRef{格前 10:28}——即你们被信德的规则所命令去做的事。\par

% structure: p0095 source=h2 target=subsection reason=relative-heading-rank; base=h2

\subsection{第十六章 关于殉道者的恒心、忍耐和爱德的圣经段落}

\BibleQuote{\InnerQuote{人心里相信，就可成义；口里宣认，就可救恩。}因此圣经说：\InnerQuote{凡信祂的人，必不至于蒙羞。}这就是我们所传的信德之言：因为你若口里宣认耶稣为主，心里相信天主使祂从死者中复活，你必得救。}这里清楚地描述了完美的义，在实践和默观中得以成全。因此我们要\BibleQuote{祝福那些迫害我们的人；祝福，而不可咒骂。}\BibleRef{罗 12:14}\BibleQuote{因为我们的夸耀是这点：良心的见证，即凭着神圣和真诚，我们认识天主，}借着这不显著的实例显出爱德的工作，即\BibleQuote{不是凭属血肉的智慧，而是凭天主的恩宠，我们在世上为人处事。}\BibleRef{格后 1:12}宗徒至此论及知识；在致格林多人后书中，他将普通的\Quote{信德的教导}称为知识的馨香。\BibleQuote{直到今日，在诵读旧约时，同样的帕子还留在许多人身上，}\BibleRef{格后 3:14}因为没有转向主而未被揭开。因此，他也向那些能领悟的人显示了复活，即那仍属肉身的生命，用肚腹爬行的生命。因此，他也将\Quote{毒蛇的种类}之名应用于纵欲者，他们服侍肚腹和私处，为了世俗的享乐而彼此砍头。\BibleQuote{孩子们，我们爱人，不可只在言语和口舌上，}若望说，教导他们要完美，\BibleQuote{而要在行动和真理上；由此我们就知道我们是出于真理的。}\BibleRef{若一 3:18}如果\BibleQuote{天主是爱，}那么虔诚也是爱：\BibleQuote{爱里没有恐惧；但完美的爱驱除恐惧。}\BibleQuote{我们遵守祂的诫命，这就是爱天主了。}\BibleRef{若一 5:3}再者，对于渴望成为\OriginalTerm{诺斯替者}{Gnostic}的人，经上写道：\BibleQuote{但你应当在言语、行为、爱德、信德、纯洁上，成为信徒的榜样。}\BibleRef{弟前 4:12}因为完美的信德，我想，不同于普通的信德。而神圣的宗徒以这些话为诺斯替者提供了准则，写道：\BibleQuote{因为我已经学会了，无论在什么境况，都可以知足。我知道怎样处卑贱，也知道怎样处丰富；或饱足，或饥饿，或有余，或缺乏，随事随在，我都得了秘决。我靠着那加给我力量的，凡事都能做。}\BibleRef{斐 4:11}此外，当与别人讨论以使他们羞愧时，他也不怕说：\BibleQuote{你们要追忆往昔，那时你们蒙光照之后，忍受了极大的痛苦斗争：一面被辱骂和磨难所观看；一面成为那些受这样对待的人的同伴。因为你们曾怜悯了我的锁链，并欣然接受了你们财产的掠夺，知道你们有更好更持久的产业。所以不要丢弃你们的勇气，它必得大赏报。原来你们需要忍耐，好使你们行完了天主的旨意，而获得所应许的。因为还有一点点时候，那要来的就来，必不迟延。但义人必因信德而生活；如果他退缩，我的心就不喜悦他。但我们不是退缩以致丧亡的人，而是有信德以致灵魂得救的人。}\BibleRef{希 10:32}然后他提出一连串神圣的榜样。他说：难道不是\BibleQuote{因着信德，}这样忍耐，他们才显贵吗？他们\BibleQuote{忍受了戏弄和鞭打，甚至锁链和监禁；他们被石头砸死，被试探，被刀杀死；他们披着绵羊、山羊的皮，到处流窜，受贫乏、患难、虐待，这世界配不上他们；他们漂流在旷野、山岭、山洞和地穴。所有这些人都因信德得了美好的见证，却未得到天主的应许，}(用\OriginalTerm{中断省略}{parasiopesis}所表达的部分留给读者理解，即\Quote{单独})。他接着加上：\BibleQuote{天主为我们预备了更好的事（因为祂是良善的），叫他们若不与我们同得，就不能完全。因此，既有这样一大片云彩，}神圣而透明，\BibleQuote{围绕我们的见证人，我们就该放下各样的重担，和容易缠累我们的罪，以忍耐奔那摆在我们前头的路程，仰望我们信德的创始者和完成者耶稣。}既然他指明了义人在基督里的一种救恩，并且对于我们的救恩，他已明确表达了前者，关于梅瑟也说了同样的话，他加上：\BibleQuote{他以基督的凌辱比埃及的财宝更富裕，因为他注视的是赏报。因着信德，他离开了埃及，不怕王怒；因为他坚持如同看见那看不见的主。}神圣的智慧论及殉道者说：\BibleQuote{在愚人眼中，他们似乎是死了，他们的离世被视为灾祸，他们从我们中间离去是痛苦。但他们处于平安中。因为虽然在别人看来，他们受了惩罚，但他们的希望充满不死。}他接着加上，教导殉道是一种光荣的净炼：\BibleQuote{他们受了少许管教，却要大大蒙福；因为天主考验了他们，}即容许他们受试炼，为考验他们，并使那试炼的作者蒙羞，\BibleQuote{并发现他们配得上祂自己，}显然被称为儿子。\BibleQuote{祂试验了他们，有如炉中的黄金；又悦纳了他们，有如全燔祭的牺牲。在他们蒙眷顾的时候，他们要光芒四射，有如禾秸间跑过的火星。他们要审判万民，治理众民，上主要永远统辖他们。}\par

% structure: p0097 source=h2 target=subsection reason=relative-heading-rank; base=h2

\subsection{第十七章 摘自克莱孟致格林多人书论殉道}

此外，在\WorkTitle{致格林多书}中，宗徒克肋孟也描绘了诺斯替信徒的形象，说：\Quote{因为你们中间有谁曾寄居在你们那里，却没有证明你们那完美而坚定的信德？谁不赞赏你们那健全而温和的虔诚？谁不歌颂你们那慷慨好客的风范？谁不祝贺你们那完全而确切的知识？因为你们行事不偏不倚，遵行天主的诫命；}等等。\par

然后更清楚地说：\Quote{让我们把目光投向那些为祂辉煌的荣耀献上完美服务的人。让我们以哈诺客为例，他因服从而被视为义人，遂被接升天；还有诺厄，他因相信而得救；以及亚巴郎，他因信德和好客被称为天主的朋友，并成为依撒格的父亲。}\Quote{因好客和虔诚，罗特从索多玛得救。}\Quote{因信德和好客，妓女辣哈布得救。}\Quote{因忍耐和信德，他们披着山羊皮、绵羊皮和骆驼毛的褶衣行走，宣讲基督的王国。我们提到祂的先知：厄里亚、厄里叟、厄则克耳和若翰。}\par

\Quote{因为亚巴郎，因他自由的信德被称为\Quote{天主的朋友}，并没有因荣耀而自高自大，而是谦逊地说：\BibleQuote{我是灰尘和灰烬。}}\BibleRef{创 18:27}\Quote{关于约伯，经上这样写道：\BibleQuote{约伯是公正、无可指摘、真实、虔诚的人，远离一切邪恶。}}\BibleRef{约 1:1}他藉忍耐胜过了试探者，同时为天主作证并蒙天主作证；他保持谦卑，说：\Quote{没有一个人是纯洁无玷的，即使他的生命只有一天。}\Quote{梅瑟，\Quote{在他的全家尽忠的仆人}，对那位从荆棘中发出神谕的天主说：\BibleQuote{我是谁，祢竟打发我去？我拙口笨舌，说话结巴。}以人的言语传达天主的声音。又说：\BibleQuote{我是从锅中的烟。}}\Quote{因为天主拒绝骄傲的人，却赐恩宠给谦卑的人。}\BibleRef{雅 4:6}\par

\Quote{达味也是如此，主为他作证说：\BibleQuote{我找到了一个合我心意的人，达味，叶瑟的儿子。我用圣油膏抹了他。}但他也对天主说：\BibleQuote{天主，求祢按祢的仁慈怜悯我；按祢丰厚的慈悲涂抹我的过犯。求祢彻底洗净我的不义，清除我的罪愆。因为我认我的过犯，我的罪愆常在我面前。}}\Quote{然后，他以真知识的节制暗示不受法律约束的罪，补充说：\BibleQuote{我唯独得罪了祢，在祢眼前行了恶。}}\Quote{因为圣经某处说：\BibleQuote{上主的神是灯，探查腹中的隐密。}}\BibleRef{箴 20:27}一个人越因行善而成为诺斯替信徒，光照的圣神就越接近他。\Quote{因此，主亲近义人，我们所生的任何心思和意念都逃不过祂——我是指主耶稣，}祂以全能旨意审察我们的心，\Quote{祂的血为我们而祝圣。因此，让我们尊敬在上者，敬重长老；让我们尊重年轻人，教导天主的纪律。}因为那相称地实行并教导主命令的人是有福的；他具有宽宏的心和默观真理的心。\Quote{让我们引导我们的妻子行善；让她们表现出，}他说，\Quote{贞洁的可爱情操；让她们显示温良的无邪意愿；让她们以沉默表达舌头的柔和；让她们的爱不依私好，而是以圣洁对一切敬畏天主的人施予同等的爱。让我们的儿女分享在基督内的纪律；让他们学习谦卑在天主面前有何益处，圣洁之爱在天主面前有何力量，敬畏上主有多可爱和伟大，它拯救所有以纯洁之心在其内行走的人：因为祂是思想和情感的探索者，祂的气息在我们内，祂若愿意，就会收回。}\par

\Quote{如今这一切都因在基督内的信德而得以确认。\BibleQuote{孩子们，来吧，}主说，\BibleQuote{听我，我要教导你们敬畏上主。谁渴望生命，喜爱看到好日子？}然后祂附上关于数字七和八的诺斯替奥秘。\BibleQuote{禁止你的舌头出恶言，你的嘴唇不说诡诈的话。远离恶，行善。寻求和平，并追随它。}因为这些话暗示\OriginalTerm{知识}{gnosis}，连同远离恶和行善，教导我们要藉言语和行为得以成全。\BibleQuote{上主的眼目看顾义人，祂的耳朵垂听他们的祈祷。上主的脸却敌对作恶的人，要从地上剪除他们的纪念。义人呼求，上主垂听，救他脱离一切患难。} \BibleQuote{罪人的鞭伤甚多；但那盼望上主的人，必有慈爱环绕。} }\Quote{丰厚的仁慈，}他高洁地说，\Quote{环绕那信靠上主的人。}\par

因为在\WorkTitle{致格林多书}中写道：\Quote{藉着耶稣基督，我们愚昧昏暗的心智得以照亮。藉着祂，上主愿意我们品尝不朽的知识。}为了更明确地表明知识的独特性质，他补充说：\Quote{既然这些事对我们而言是清楚的，我们要深入默观神圣的知识，就应当按着秩序遵行上主命令我们在指定时刻所行的一切。因此，智慧人不仅要以言语，更要以善行显出自己的智慧。谦卑的人不要自证，而应容许别人为他作证。在肉身纯洁的人不要自夸，要知道是另一个人赐给他节德。弟兄们，你们看，我们越陷入危险，就越配得知识。}\par

% structure: p0104 source=h2 target=subsection reason=relative-heading-rank; base=h2

\subsection{第十八章　论爱德与克制欲望}

\Quote{因此，我们慈善的有节制倾向}，根据\Person{克肋孟}，\Quote{寻求公共益处}；无论是通过遭受殉道，还是通过言行教导——后者是双重的，口传与书写。这就是爱，爱天主并爱近人。 \Quote{这导向那不可言说的高峰。\InnerQuote{爱遮盖许多罪过。\BibleRef{雅 5:20} 爱凡事包容，凡事忍耐。} \BibleRef{格前 13:7} 爱使我们与天主结合，凡事和谐。在爱内，天主所有的选民都得以成全。离开爱，没有什么是天主所喜悦的。} \Quote{论及它的完美，无可言表。}据说。\Quote{谁能适合在其中被发现，除非天主认为谁相称？} 关于这点，宗徒\Person{保禄}说：\Quote{我若舍身，却没有爱，就成了鸣的锣，响的钹。}他的意思是，若我作证不是出于由\OriginalTerm{诺斯替者}{gnostic}之爱决定的意向，而是出于恐惧和期待报酬，为了向主作证而启动双唇，好承认主，那么我只是一个普通人，空喊主的名，却不认识祂。 \Quote{因为有一族人用嘴唇爱；另有一族人将身体交付火焚。} \Quote{并且我若将所有财物施舍}，他说，不是出于爱的通传原则，而是为了回报，无论是从受惠者还是从应许的主那里；\Quote{并且我若有全备的信德，能够移山}，并除去蒙蔽的情欲，却不因爱而对主忠实，\Quote{我就算不得什么}，相比于那作为\OriginalTerm{诺斯替者}{Gnostic}作证的人，相比于群众，算不得更好。\par

\Quote{从亚当直到今日，所有世代都已逝去。但那些因天主的恩宠而在爱内得以成全的人，将占有虔诚者的位置，这些虔诚者将在基督王国来临的巡视中显现。}爱不允许犯罪；但若因仇敌的干扰而陷入此类情况，效法\Person{达味}，它将歌唱：\Quote{我要向上主承认，这比献上有角有蹄的公牛犊更令他喜悦。愿贫穷者看见而欢乐。}因为他说：\Quote{以赞颂的祭献祭献天主，向主偿还你的誓愿，并在患难之日呼求我，我必拯救你，你必荣耀我。}\Quote{因为天主所悦纳的祭献是破碎的心灵。}\par

\Quote{天主}，既善，\Quote{是爱}，据说。\BibleRef{若一 4:8} 祂的\Quote{爱不加害于人}，\BibleRef{罗 13:10}既不伤害，也不报复，总之，按照天主的肖像，对众人行善。\Quote{爱}因此\Quote{是法律的满全}；\BibleRef{罗 13:10}正如基督，即爱我们的主的临在，以及我们爱的教导和按照基督的纪律。因此，靠着爱，不可奸淫、不可贪恋邻人妻子的诫命得以满全，这些罪曾受恐惧禁止。\par

\Quote{给\OriginalTerm{诺斯替者}{Gnostic}准备了眼所未见，耳所未闻，人心所未想到的}但给那操练单纯信德的人，祂作证以百倍回报他所舍弃的——这应许已落入人的理解范围之内。\par

此时，我记起一个自称为\OriginalTerm{诺斯替者}{Gnostic}的人。他解释这句话：\Quote{但我告诉你们：凡注视妇女而起贪念的，已经犯了奸淫。}\BibleRef{玛 5:28}他认为被谴责的不是单纯的欲望，而是如果通过欲望，那由此产生的行为，超越欲望而实现在欲望中。因为梦境运用幻想和身体。因此，历史学家叙述了正义者\Person{博科里斯}的以下判决。一位青年爱上了一个妓女，说服女孩以约定报酬次日来他处。但他的欲望意外地通过梦见抓住女孩而提前满足，当他的意中人按约定来时，他阻止她进入。但她得知后，要求报酬，说她已经这样满足了情人的欲望。他们于是来到法官面前。法官命令青年在太阳下举起装有报酬的钱袋，并叫妓女攫取影子；诙谐地令他支付影子的报酬，以回报拥抱的影子。\par

因此，当灵魂赞同所见之景时，人便做梦。但那睁着眼做梦的，是观看以致贪恋；不仅如那位诺斯底派所言，若在看见妇人时心中幻想交合——这已是贪欲之行为，因为是贪欲——而且，若人看见身体之美（圣言说），且肉身在他眼中因贪欲而显得美丽，以血肉和罪恶的方式观看，他便因赞赏而被定罪。因为另一方面，那以贞洁之爱观看美丽的人，并不认为肉体是美的，而是精神；在我看来，他欣赏身体如同一个形象，借其美丽将自己提升至艺术家和那真正的美丽；向上等待上升的天使展示神圣的象征，即正义的明印；我是指蒙悦纳的傅油，即因圣神的相通而喜乐的灵魂中居住的性情品质。这荣耀曾照耀在梅瑟的面容上，百姓不能直视。因此他为那荣耀蒙上一层面纱，给那些以血肉观看的人。因为那些收税的人，扣留带来任何世俗之物、被自己情欲所累的人。但那摆脱一切应纳税之物、充满知识及行为正义的人，他们便以美好的祝愿放行，祝福那人及其工作。\Quote{他的生命必不衰败}——那被\Quote{水渠}滋养的活树之叶。如今义人被比作结果实的树，而不只是那些高大性质的树。在按律法的祭献供物中，有那些在祭品中寻找瑕疵的人。精通此类事物的人区分倾向（\OriginalTerm{倾向（propension）}{\GreekText{ὄρεξις}}）与贪欲（\OriginalTerm{贪欲（lust）}{\GreekText{ἐπιθυμία}}）；他们将后者归为不合理性的，属于享乐和放纵；而倾向，作为理性的运动，他们归为自然之需。\par

% structure: p0111 source=h2 target=subsection reason=relative-heading-rank; base=h2

\subsection{第十九章——女人与男人同样能达到成全}

在这成全中，男人和女人同样可以分享。因此，不仅是梅瑟听见天主说：\Quote{我曾一次、两次对你们说：我看见了这百姓，看，他们是硬着颈项的。你容我灭绝他们，从天下涂抹他们的名；我要使你成为一个比这更大、更奇妙的邦国。}他回答时并非顾念自己，而是众人的救恩：\Quote{决不，上主！求你赦免这百姓的罪，不然就从生命册上涂抹我的名吧。}他的成全何其大——情愿与百姓同死，而不愿独自得救！\par

但友弟德，她在妇女中成为成全者，在城市被困时，应长老们的恳求，前往外邦人的营寨，为了国家不计一切危险，凭着对天主的信德将自己交在敌人手中；她立刻得到了信德的赏报——虽是一介女子，却战胜了信德的仇敌，取得了敖罗斐乃的首级。再者，艾斯德尔凭着信德成为成全者，她将以色列从国王和总督的残酷中拯救出来：一个孤单的女子，以禁食苦行，挡住了万千武装之手，以她的信德废除了暴君的谕令；她平息了王，抑制了哈曼，并借着对天主的完美祈祷，保全了以色列无损。我略过苏撒纳和梅瑟的姐妹不提，因为后者是先知率领军队的同伴，超越希伯来妇女中所有以智慧著称的人；而前者以超然的谦逊，甘愿被放荡的仰慕者判处死刑，至死是贞洁的不动摇的殉道者。\par

哲学家狄昂也讲述，某位吕西迪卡因过度谦逊，穿着衣服沐浴；而斐洛特拉，当她进入浴池时，随着水覆盖赤裸部分，她逐渐后退长袍；然后起身时又慢慢穿上。阿提卡的莱埃娜不是勇敢地忍受了酷刑吗？她参与了哈尔摩狄乌斯和阿里斯托革顿反对希帕尔科斯的密谋，虽受重刑，却一言不发。他们说，阿尔戈斯的妇女在诗人忒勒西拉带领下，仅仅现身便击退了强悍的斯巴达人；而她使她们视死如归。同样，撰写《达那伊得斯》的人关于达那俄斯的女儿们说：——\par

\Quote{于是达那俄斯的女儿们迅速武装，\newline{}在美丽的河流之前，庄严的尼罗河，}\par

等等。\par

其他诗人歌颂阿塔兰忒狩猎的敏捷、安提克勒亚对子女的爱、阿尔刻斯提斯对丈夫的爱、马卡里亚和许阿金托斯们的勇气。我还说什么呢？毕达哥拉斯派特阿诺不是如此精进于哲学，以致对那注视她并说\Quote{你的手臂很美}的人，她回答\Quote{是的，但它不属于公众}。同样得体的，还有以下记载的回答。当被问及女人与丈夫同房后何时参拜地母节时，她说：\Quote{从自己丈夫那里立刻就去，从陌生人那里永远不去。}兰普萨库斯的特弥斯托，佐伊罗斯的女儿，勒昂特斯的妻子，研习了伊壁鸠鲁哲学；同样，毕达哥拉斯派特阿诺的女儿弥亚，以及撰写了狄奥尼修斯历史书的阿里格诺忒，也是如此。\par

被称为克洛诺斯的狄奥多若的女儿们都成了辩证学家，正如辩证学家斐洛在\WorkTitle{梅涅克塞诺斯}中所说的，她们的名字如下——梅涅克塞涅、阿尔吉亚、特奥格尼斯、阿耳忒弥西亚、潘塔克莱亚。我还记得一位女犬儒派——她叫做希帕尔基亚，马罗尼人，克拉特斯的妻子——在她身上，那所谓的犬婚是在佩西勒柱廊举行的。还有昔兰尼的阿瑞忒，阿里斯提波斯的女儿，她教育了她的儿子阿里斯提波斯，后者被称为\Quote{母亲教成的}。阿耳基斯的拉斯忒涅亚和弗利乌斯的阿克西奥忒亚曾与柏拉图一起研究哲学。此外，米利都的阿斯帕西娅，喜剧作家们对她多有记述，她受苏格拉底训练哲学，受伯里克利训练修辞。其余的我因篇幅省略，既不提及女诗人科琳娜、忒勒西拉、弥亚和萨福，也不提及女画家，如克拉提努斯的女儿伊瑞涅和尼尔刻斯的女儿阿那克珊德拉，据狄迪莫斯在\WorkTitle{会饮录}中的记述。林狄人智者兼君王克勒奥布洛斯的女儿，不羞于为父亲的客人洗脚。亚巴郎的妻子、有福的撒辣，亲自为天使们准备了炭火烤的饼。希伯来人的王族少女牧羊。荷马笔下的瑙西卡也因此去了洗衣槽。\par

明智的女子首先会选择说服她的丈夫与她一同追求幸福之事。如果那不可行，就让她独自热切地追求德行，在一切事上获得丈夫的同意，从而永不违背他的意愿，除非是那些被认为有助于德行和救恩的事。但如果有人阻止妻子或女仆过这样一种生活，而她的心志已坚定于此，那么这个人所做的事就显然无非是决心将她从正义和节制中赶走，并选择使自己的家庭变得邪恶和放荡。\par

因此，男子或女子不可能在任何事情上有所成就而不借助教育、实践和训练的好处；我们说过，德行不在于他人，而首先在于我们自己。其他事情可以借着对抗来抑制；但关于自身的事，即使最坚持不懈的努力也完全无能为力。因为恩赐是来自天主的，而非任何他人的能力。因此，放纵应被视为除了放纵者本人之外没有任何他人之恶；而节制，则被视为能够实践它的人之善。\par

% structure: p0121 source=h2 target=subsection reason=relative-heading-rank; base=h2

\subsection{第二十章 贤妻}

欧里庇得斯在劝诫时描述了那位得体地爱着自己丈夫的女子：\par

\Quote{当丈夫说什么时，\newline{}她若沉默，就该认为他说得好；\newline{}若她想说什么，就应尽力取悦丈夫。}\par

他接着补充道：\par

\Quote{妻子应与丈夫甜蜜地同忧，\newline{}若有什么不幸降临到他身上，\newline{}并应共同分担忧伤和喜乐。}\par

然后，描述她在不幸中仍温柔和善良，他补充道：\par

\Quote{而我，当你生病时，将分担你的痛苦承受它；\newline{}我也将分担你的不幸。}\par

以及：\par

\Quote{没有什么对我是苦涩的，\newline{}因为对朋友应当快乐，\newline{}否则友谊是什么？}\par

那么，按照圣言成就的婚姻，是圣化的，如果这结合是在天主的权下，并以\BibleQuote{真诚的心，完全的信德，良心洗净，身体用清水洗净，并坚持望德的宣认；因为那应许的是信实的}来实行。婚姻的幸福绝不应以财富或美貌来衡量，而应以德行来衡量。\par

悲剧说：\par

\Quote{美貌对妻子帮助不了丈夫；\newline{}但德行帮助了许多；因为每一个好妻子\newline{}若依恋她的丈夫，就知道如何实行节制。}\par

然后，作为劝诫，他说：\par

\Quote{首先，这在于有智识的她，\newline{}即使丈夫丑陋，他也必须显得好看；\newline{}因为判断在于心灵，而非眼睛。}\par

等等。\par

圣经极合宜地说，女人被天主赐给男人作为\Quote{一个助手}。那么，在我看来，显然她将负起责任，以明智和劝勉来弥补丈夫在家政中引起的各种烦恼。如果他不顺从，她将尽力——就人性之可能——度一个无罪的生活；无论是按照理性必须死，还是必须活；她认为天主在这条道路上是她帮助者和伴侣，是她当前及未来的真正捍卫者和救主；使祂成为她一切行为的领袖和向导，视清醒与正义为自己的工作，并以天主的恩宠为终向。因此，宗徒在\WorkTitle{弟铎书}中优美地写道：\BibleQuote{年长妇女应举止虔诚，不说谗言，不沉溺于酒；她们应劝诫年轻妇女爱丈夫、爱子女、明智、贞洁、持家、良善、服从自己的丈夫，以免天主的话受亵渎。}\BibleRef{铎 2:3} 但他又说：\BibleQuote{你们要与众人追求和平与圣洁，非圣洁没有人能见主；又要谨慎，免得有人成为淫乱或亵圣的人，如同厄撒乌为了一碗食物出卖了长子的名分；免得有苦根长出来扰乱你们，且使许多人因此被玷污。}\BibleRef{希 13:14} 然后，作为对婚姻问题的最后收笔，他补充说：\BibleQuote{婚姻在各方面都应受尊重，床榻也无玷污；但淫乱者和奸淫者，天主必要审判。}\BibleRef{希 13:4} 而且，既然在成全方面，男人与女人拥有同一个目标、同一个结局，伯多禄在他的书信中说：\BibleQuote{虽然如今在暂时的考验中，你们可能因各种试探而忧愁；但你们信德的考验，比那被火试验而朽坏的金子更为珍贵，必在耶稣基督显现时，得着称赞、荣耀和尊敬；你们虽然没有见过祂，却爱祂；如今虽然看不见祂，却因信祂而怀着不可言传、充满光荣的喜乐，并获得你们信德的成果，即灵魂的救恩。}\BibleRef{伯前 1:6} 因此，保禄也为了基督的缘故而喜乐，因为\BibleQuote{他更劳苦，更多次受鞭打，更常常冒死。}\BibleRef{格后 11:23}\par

% structure: p0137 source=h2 target=subsection reason=relative-heading-rank; base=h2

\subsection{第二十一章 完人之描述，或称诺斯替者}

此處我發現，成全之事按照與那在每一德行上超越者的關係，以多種方式被理解。因此，一個人被成全為虔誠者、忍耐者、節制者、工作者、殉道者、靈知者。但我不知道在仍為人時，有任何人在一切事上同時是成全的——僅就律法的字面而論——除了那為我們披上人性的獨一者。那麼，誰是成全的呢？是那宣稱遠離罪惡的人。是的，這是通往福音和行善的道路。但就律法之人而言，靈知者的成全，是接納福音，好使那在律法之下的人得以成全。因為如此，那在律法之後的梅瑟曾預言，必須聽從，好使我們能按宗徒所說，接受基督——法律的圓滿。\BibleRef{申 18:15}但如今在福音中，靈知者不僅藉著將律法用作階梯而進步，更藉著理解和領悟它——正如那賜予盟約的主將它傳授給宗徒們。並且，若他行事端正（因為藉著惡行獲得\OriginalTerm{靈知}{gnosis}，當然是不可能的）；而且，若他更進一步，作了極其正確的宣認，出於愛而成為殉道者，在人間獲得相當的聲譽——即便如此，他在肉體內也不能被稱為成全；因為是生命的終結才配得這個稱號，此時靈知者殉道者首先顯示了成全的工作，正確地展現了它，並感恩地流了自己的血，交付了靈魂：那時他將是有福的，並真正被宣告為成全的，\Quote{這卓越的能力是出於天主，而不是出於我們}，正如宗徒所說。只讓我們保存自由意志和愛：\Quote{事事憂苦，卻不絕望；遭迫害，卻不被遺棄；被打倒，卻不被毀滅。}\BibleRef{格後 4:8}因為那些追求成全的人，按同一宗徒，必須\Quote{在一切事上不給人任何絆腳石，反而在各樣事上表明自己是天主的僕人}。並且，作為結果，他們也應當向人讓步；因為為了辱罵的毀謗，這是合理的。以下是具體的描述：\Quote{以極大的忍耐，在患難、匱乏、困苦中，在拷打、監禁、騷亂、勞苦、不眠、禁食中，以純潔、知識、堅忍、仁慈、聖神、無偽的愛德、真理之言、天主的德能}\BibleRef{格後 6:3}，好使我們成為天主的聖殿，\Quote{從一切肉體和靈魂的污穢中}潔淨。\Quote{並且，我——祂說——要接納你們；我要作你們的父親，你們要作我的兒女，這是全能的主說的。}\Quote{所以，我們——他說——當在敬畏天主中成全聖德。}因為雖然恐懼產生痛苦，\Quote{我歡喜——他說——不是因為你們憂愁，而是因為你們曾憂愁以致悔改。因為你們照著天主的意思憂愁，凡事就不至於因我們受虧損。因為依著天主的意思憂愁，就生出沒有後悔的悔改，以致於救恩；但世俗的憂愁叫人死。你看，你們照著天主的意思憂愁，從此就生出何等的殷勤、自訴、自恨、恐懼、想念、熱心、責罰！在一切事上，你們都表明自己是潔淨的。}\BibleRef{格後 7:1}這些是靈知修練的預備練習。既然全能的天主自己\Quote{賜予一些宗徒，一些先知，一些傳福音者，一些牧者及教師，為成全聖徒，盡職事之工，建立基督的身體，直到我們眾人達到信仰的合一，和對天主子圓滿的認識，成為成全的人，達到基督圓滿年齡的程度}；那麼我們就當努力達到與靈知者相稱的成年，並在肉體中儘可能成全，以完美的和諧在此世與天主的旨意一致，為恢復那真正成全的高貴和親屬關係，達到基督的圓滿——這圓滿完全依賴於我們的成全。\par

如今我們看到神聖的宗徒在何處、如何、何時提到成全的人，以及他如何指出成全者的差別。另一方面：\Quote{聖神的顯示是為眾人的益處。有人藉著聖神得了智慧的言語；有人藉著同一聖神得了知識的言語；有人藉著同一聖神得了信德；有人藉著同一聖神得了治病的恩賜；有人行了奇蹟；有人得了先知之恩；有人得了辨別神恩；有人得了各種語言；有人得了解釋語言：這一切都是同一聖神所運作，隨心所欲分給各人。}\BibleRef{格前 12:7}既是如此，先知們在預言上成全，義人在義德上成全，殉道者在宣認上成全，其他人在宣講上成全——並非他們不分享共同的德行，而是在他們被委派的事上精通。因為哪個有理智的人會說先知不義呢？難道義人如亞巴郎沒有說過預言嗎？\par

\Quote{因為天主賜給一人戰功，\newline{}給另一人舞蹈的成就，\newline{}給另一人豎琴和歌曲。}\par

荷马说。\Quote{但每人都有天主所赐的独特恩赐}\BibleRef{格前 7:7}——一人是这样，另一人是那样。但宗徒们在一切事上都是完备的。如果你愿意，你将在他们的言行和著作中找到知识、生命、宣讲、正义、纯洁、预言。因此，我们必须知道，虽然保禄在时间上是年轻的——他紧接主升天之后便兴盛起来——但他的著作依赖于旧约，呼吸并言说旧约。因为对基督的信德和福音的知识是法律的解释和成全；因此曾对希伯来人说：\Quote{如果你们不相信，你们必不能明了；}\BibleRef{依 7:9}意思是，除非你们相信法律所预言并由法律以神谕方式传达的，你们就不能明了旧约，那是主藉祂的来临所阐释的。\par

% structure: p0142 source=h2 target=subsection reason=relative-heading-rank; base=h2

\subsection{第二十二章 真知者行善，非因畏惧惩罚或盼望回报，而只为善本身}

那么，具有理解和洞察力的人便是真知者。他的事业不是戒除邪恶（因为这是达到最高成全的一步），也不是因畏惧而行善。因为经上记载：\Quote{我往哪里去躲避祢的面？我若升到天上，祢在那里；我若逃到海极，祢的右手也在那里；我若下到深渊，祢的圣神也在那里。}他也不因盼望所应许的报酬而行善。因为经上说：\Quote{看哪，主，祂的赏报在祂面前，要按各人的行为赐给各人；天主为爱祂的人所预备的，是眼所未见，耳所未闻，人心所未想到的。}但唯有出于爱、并因善本身的卓越而行善，才是真知者的选择。如今，在天主位格中，对主说：\Quote{你向我求，我将外邦人赐给你为基业；}教导祂求一个真正君王的请求——即无价地拯救人类，使我们继承并拥有主。相反，为任何实用目的而渴望关于天主的知识，以便做这事或不做那事，不适合真知者；知识本身作为沉思的理由便足够了。因为我敢断言，那为神圣知识本身而献身于知识的人，选择知识并非因为他希望得救。因为心智的操练通过练习延长为永恒的操练。而心智的永恒操练是理智存在的本质，它源于不间断的混合过程，并保持永恒的沉思，一种活生生的实体。那么，我们能否设想有人向真知者提出：他宁愿选择关于天主的知识还是永久的救恩？如果这两者是完全相同的，并且可以分开，他将毫不犹豫地选择关于天主的知识，认为那从爱上升到知识的信德特质因其本身而值得追求。那么，这便是成全之人的第一种行善形式：行善不是为任何与己相关的好处，而是因为他判断行善是正当的；而能力在万有中有力地运作，在行动本身中成为善；不是在某些事上善，在另一些事上不善；而是由行善的习惯构成，既不为了荣耀，也不像哲学家所说的为了名誉，也不为了来自人或天主的报酬；而是为了按主的肖像和模样度日。\par

如果他在行善时遇到任何逆境，他将毫无怨言地任凭补偿过去，就如同那是善的，因为他公义而良善的，\Quote{对义人和不义的人。}主对这样的人说：\Quote{你们应当如同你们的天父一样成全。}\par

对他而言，肉身是死的；只有他自己活着，将坟墓祝圣为献给主的圣殿，将那旧有罪过的灵魂转向天主。\par

这样的人不再节制，而已经达到无动于衷的境界，等待披上神圣的肖像。经上说：\Quote{你若施舍，不要让任何人知道；你若禁食，要梳头擦脸，惟独让天主知道，}不让人知道。甚至那施恩的人也不该知道自己在施恩；因为这样他有时会仁慈，有时不仁慈。当他因习惯而行善时，他将模仿善的本性，他的性情将成为他的本性和实践。没有必要挪移那些被高举的人，但那些行走的人有必要跨越整个窄路，达到必需的目标。因为这就是被父所吸引，成为配得从天主领受恩宠的能力，以便毫无阻碍地奔跑。如果有人恨恶被选者，这样的人知道他们的无知，并怜悯他们愚昧的理智。\par

因此，知识本身合理地爱慕并教导无知者，教导整个受造界荣耀全能的天主。如果这样的人教导爱天主，他不会在任何情况下失去美德，无论是清醒、在梦中或在任何异象中；因为习惯永远不会因脱离习惯而失去自身。那么，无论是说知识是习惯还是性情，由于不同的意见从未获得进入，主导的官能保持不变，不容许通过白天运动在梦中形成幻想的景象来改变外观。因此，主也命令\Quote{要警醒}，使我们的灵魂即使在梦中也不被激情扰乱；并且要使夜间的生命如同在白天一样纯洁无瑕。因为尽可能地肖似天主，就是保持心志与同样事物的关系。而这就是心志作为心志的关系。\par

但性情的差异源于对物质事物的过分贪恋。正因如此，按我看来，他们称黑夜为\OriginalTerm{欧弗洛涅}{Euphrone}；因为那时灵魂摆脱了感官知觉，返回自身，更真实地把握\OriginalTerm{智慧}{\GreekText{φρόνησις}}。因此，奥秘仪式大多在夜间举行，表示灵魂在夜间离开身体。\BibleQuote{所以我们不要像别人一样睡觉，而要保持警醒和清醒。因为睡觉的人是在夜里睡觉，醉酒的人是在夜里醉酒；但我们这些属于白昼的人，应当清醒，穿上信德和爱德的胸甲，戴上救恩的望德作头盔。} \BibleRef{得前5:6} 再者，他们对睡眠所说的一切，同样也适用于死亡。因为二者都显示灵魂的离开，只是程度不同；我们也可从\Person{赫拉克利特}得到这点：\Quote{人在自身中触摸到黑夜，当他死亡、其光明熄灭时；而活着时，当他睡觉，他触摸到死者；醒着时，当他闭上眼睛，他触摸到睡着的人。} \Quote{看见主的人是有福的}，按宗徒所说；\Quote{因为现在是时候从睡眠中醒来了。因为我们的救恩现在比我们初信时更近了。黑夜已深，白昼将近。所以我们应抛弃黑暗的行为，穿上光明的武器。} 他用白昼和光明比喻性地指圣子，用光明的武器喻指恩许。\par

所以据说我们应当洗净身体去参加祭献和祈祷，保持洁净明亮；这种外在的装饰和净化是为作标记。而纯洁在于思念神圣的思想。此外，还有圣洗圣事的象征，这也是从\Person{梅瑟}传给诗人的，如下：—\par

\Quote{她打了水，身上穿着洁净的衣服。}\par

这是\Person{佩涅洛佩}要去祈祷：—\par

\Quote{而\Person{忒勒马科斯}，\newline{}在苍茫的大海中洗了手，向\Person{雅典娜}祈祷。}\par

犹太人的习俗是起床后经常清洗。于是便有了那句名言：—\par

\Quote{要纯洁，不是用水洗身，而是用心。}\par

因为在我看来，圣德就是心灵、行为、思想乃至言语的完全纯洁，最终连梦境中也无罪。\par

我认为，对于人而言，充分的净化就是彻底而坚固的悔改。如果我们谴责自己先前的行为，在往后前进时深思熟虑，使心灵摆脱那些通过感官取悦我们的事物以及先前的过犯。\par

那么，如果我们为\OriginalTerm{知识}{\GreekText{ἐπιστήμη}}追溯词源，其含义来源于\OriginalTerm{放置}{\GreekText{στάσις}}；因为我们的灵魂从前飘荡不定，时而这样，时而那样，如今它安顿在对象上。同样，信德也可作词源解释，即我们灵魂在面对存有之物时的安顿（\OriginalTerm{放置}{\GreekText{στάσις}}）。\par

但我们渴望了解那一位始终并在一切实事上正义的人；他既不畏惧法律所施加的惩罚，也不害怕因为与那些同他一起生活且迫害受害者的人为伍而怀恨邪恶，也不惧怕来自作恶者的危险，却仍然保持正义。因为一个人若因这些考虑而避免任何错事，便不是自愿良善，而是出于恐惧才为善。就连\Person{伊壁鸠鲁}也说，在他眼中聪明的人\Quote{不会为了利益而伤害任何人；因为他无法说服自己能逃过察觉。}因此，如果他明知不会被察觉，按他的说法，他就会作恶。这些都是黑暗的教义。再者，若有人因指望天主因义行而赐予的赏报而避免作恶，依此假设，他也不是自发地良善。因为如同恐惧使那人正义，同样赏报使这人正义；或者更准确地说，使他显得正义。但对于死后的希望——好人得善望，坏人则相反——不仅追随蛮族智慧的人，而且毕达哥拉斯学派也知晓。因为后者也为那些追求哲学的人设立了希望作为目的。而\Person{苏格拉底}在\WorkTitle{斐多篇}中说\Quote{善灵带着善望离开此处；}接着，在谴责邪恶时，他又有相反的说法：\Quote{因为他们怀着恶望而活。}\Person{赫拉克利特}在他的论人文章中也明显与他一致：\Quote{人在死后将面对他们既未希望也不曾想到之事。}因此，\Person{保禄}神圣地写道：\BibleQuote{患难生坚忍，坚忍生经验，经验生希望；而希望不会使人羞惭。}\BibleRef{罗5:3}因为坚忍是由于对未来的希望。现在，希望与赏报和希望的恢复同义；这希望不会使人羞惭，不再受辱。\par

但是那仅仅因为被召唤而服从召唤的人，既非出于恐惧，也非为了享乐，正走在通往\OriginalTerm{知识}{\GreekText{γνῶσις}}的道路上。因为他并不考虑是否有任何外在的有利可图的收益或享乐随之而来；而是被那真正可爱之爱的对象所吸引，并被引导到必要的事物上，实践虔诚。因此，即使我们假设他从天主那里获准去做禁止的事而不受惩罚，即使他得到应许，将获得作为奖赏的蒙福者的美善事物；此外，即使他能说服自己天主会被他所作所为蒙蔽（这是不可能的），他也绝不会愿意做任何违背正确理性的事，因为他已经选择了那真正美善且因其自身而值得选择、因而值得被爱的事物。因为我们听说，善并不在于腹中的食物。但他已听到\BibleQuote{食物不能使我们得悦纳}，\BibleRef{格前 8:8}也非婚姻，亦非出于无知的禁婚；而是有德行的知识性的品行。因为狗，作为一种无理性的动物，可以说是节制的，因为它害怕举起的棍子，因此远离肉。但如果将所预言的应许取消，将威胁的恐惧消除，将逼近的危险移除，这些人的内心就会显露出来。\par

% structure: p0160 source=h2 target=subsection reason=relative-heading-rank; base=h2

\subsection{第二十三章 同一主题续}

因为这与事物本身的性质不符，他们不能以真正知识者的方式领悟真理：即所有为我们受造之物都是善的，例如婚姻和生育，若能适度使用；并且变得无欲无求，通过肖似天主而修德，比善更好。但在外部事物——愉快或不愉快——上，他们有些节制，有些不节制。而在他们因厌恶而节制的事物上，他们显然在责备受造物和造物主；虽然表面上他们忠信地行走，但他们所持的观点是不虔敬的。那条诫命\BibleQuote{不可贪恋}，既不需要恐惧的必要性——这迫使远离愉快之事；也不需要那因应许而劝说克制情欲冲动的奖赏。\par

而那些因应许而服从天主的人，被享乐的诱饵所捕，选择服从不是为了诫命本身，而是为了应许。脱离感官对象，作为必然结果，并不会产生对理智对象的依恋。相反，对理智对象的依恋自然而然地成为对知识者一种抽离感官对象的影响力；因为他凭借对美善的选择，\OriginalTerm{按照知识}{\GreekText{γνωστικῶς}}选择了美善，赞赏生育，并通过圣化造物主来圣化肖似天主。但我要摆脱贪欲，让他说，主啊，是为了与祢结合。因为受造的治理是善的，万物都管理得当；没有一件事是无因而起的。我必须在你之内，全能者。如果我就在那里，我就接近祢。我要摆脱恐惧，以便能够接近祢，并满足于少量，实践祢在善与似善之间的公正选择。\par

宗徒以奥秘和神圣的方式，教导我们那真正恩宠的选择——不是以拒绝其他事物为恶的方式，而是以做比善更好的事——他说：\BibleQuote{所以，那把自己女儿出嫁的，做得好；那不把她出嫁的，做得更好；这是就适宜和专心事主而言的。}\BibleRef{格前 7:38}\par

现在我们知晓，困难的事并非必要；而必要的事已被天主仁慈地使其易于获得。因此\Person{德谟克利特}说得对，\Quote{本性与教导}彼此相似。我们已简要说明了原因。因为教导使人和谐，并通过和谐使人合乎本性；一个人是天性如此，还是通过时间和教育转变而成，都无关紧要。主提供了两者：即通过受造，以及通过盟约而再创造和更新。那有益于更优越者的是更可取的；而比一切都优越的是心智。因此，真正美善的事被视为最令人愉悦，并自行产生所渴望的果实——灵魂的安宁。经上说：\BibleQuote{听从我的人，必在平安中休息，坦然无惧，不怕任何灾祸}\BibleRef{箴 1:33}\BibleQuote{全心、全意依靠天主}\BibleRef{箴 3:5}。\par

如此，知识者已经有可能成为天主。\BibleQuote{我曾说：你们是神，都是至高者的儿子。}\Person{恩培多克勒}也说智者的灵魂成为神，他写道：\par

\Quote{终为先知、诗人和医生，\newline{}以及凡人中的佼佼者，\newline{}从此生出至尊的神明。}\par

那么，从基因上考虑，人是依照与同生之灵的理念而形成的。因为他并非在自然的工场中被造为无形无状，在那里，人的产生奥秘地完成，艺术和本质都是共通的。但是个体的人是根据他选择对象在灵魂中留下的印痕而被镌刻的。因此我们说\Person{亚当}就其形成而言是完美的；因为人的理念和形态的任何特有特征都不缺少；但在进入存在时他获得了完美。他因服从而成义；这就是达到成年，就他而言是如此。其原因在于他的选择，尤其是他选择了禁止的事。天主并非原因。\par

因为产生是双重的——被产生之物的产生和生长之物的产生。而人的刚毅，按照他们的说法，使人——他是易受扰动的——在本质上无所畏惧且不可战胜；愤怒是心智在忍耐、坚忍等事物上的卫星；自制和有益的理智则管辖欲望。但天主是不受激情的，无愤怒，无欲望。祂也不是因避免可怕之事而无恐惧，也不是因控制欲望而有节制。因为天主的本性既不会遇到任何可怕之事，祂也不逃避恐惧；正如祂不会感到欲望，从而统治欲望一样。因此，那毕达哥拉斯式的说法神秘地针对我们而说出：\Quote{人应当成为一}，因为大司祭自身就是一位，天主是一位，在永恒涌流美善的不变状态中。如今，救主已除去了愤怒连同情欲，愤怒是复仇的欲望。因为普遍而言，易受感受属于每一种欲望；而人当被纯粹神化达到无感状态时，就成为一。因此，正如在海中被锚定的人，他们拉锚，但并非把锚拖向自己，而是把自己拖向锚；同样，那些按照灵知生活将天主拉向自己的人，不知不觉地把自己带到天主面前：因为敬畏天主的人，敬畏自己。于是在默观生活中，人在敬拜天主时关注自己，并通过自己无瑕的净化而神圣地瞻仰至圣的天主；因为自制，当其临在、不断审视和默观自身时，就尽可能地被同化于天主。\par

% structure: p0169 source=h2 target=subsection reason=relative-heading-rank; base=h2

\subsection{第24章. 天主惩罚的理由与目的}

现在，那在我们的能力范围之内的事，是我们与其对立面同样能主宰的——例如，从事哲学或不从事，相信或不相信。因此，由于我们同样主宰每一对立面，取决于我们的事就被认为是可能的。如今，诫命可由我们去行或不行，我们因此合理地受到称赞或责备。而那些因自己所犯的罪而受惩罚的人，只为他们所犯的罪受惩罚；因为已做的事是过去的事，已做的事永远不能撤销。因此，在信德之前所犯的罪被主赦免，并非因为那些罪被撤销了，而是仿佛它们未曾发生过。但\Person{巴西里得}说：\Quote{并非所有罪，而只是非自愿和出于无知的罪被赦免}，仿佛这样的恩惠是由人而不是由天主所赐。圣经对这样的人说：\Quote{你以为我像你一样。}但如果我们因自愿的罪而受惩罚，我们受惩罚并非为了使已犯的罪被撤销，而是因为它们已被犯下。但惩罚对犯罪者无补于撤销其罪，而是使他不再犯罪，也使他不再陷入同样的事。因此，良善的天主出于这三个原因而纠正：第一，使被纠正者比从前更好；第二，使那些能被榜样挽救的人受到警告而退避；第三，使那受伤害的人不易被轻视，且易接受伤害。纠正有两种方式：教导性的和惩罚性的，我们称之为惩戒。应当知道，那些在圣洗后陷入罪的人正是受惩戒的人；因为之前的行为被赦免，之后的行为被净化。正是针对不信者才说：\Quote{他们被算作被风吹散在地面上的糠秕，以及从器皿中滴落的水滴。}\par

% structure: p0171 source=h2 target=subsection reason=relative-heading-rank; base=h2

\subsection{第25章. 真正的成全在于认识与爱慕天主}

\Quote{拥有知识之文化的人有福了，他不被驱向伤害公民或行不义之事，而是默观不朽本性的不衰秩序，它如何、以何种方式和样式存在。对于这样的人，卑劣的行为不会附着。}因此，\Person{柏拉图}正确地说：\Quote{那投身于默观理型的人将作为神生活在人间；心智是理型的处所，而天主是心智。}他说默观看不见的天主的人将作为神生活在人间。在\WorkTitle{智者篇}中，\Person{苏格拉底}称那位来自\Place{埃利亚}的辩证者客人为\Quote{神}：\Quote{这样的神如同异乡客人，常常访问城市。}因为当灵魂超越生成领域，独自存留，居住在理型中间时，如同\WorkTitle{泰阿泰德篇}中的歌队长，现在成为了天使，它将在默观中与基督同在，始终注视天主的旨意；实际上\par

\Quote{惟独有智慧，而这些如同影子般飘过。}\par

\Quote{让死人埋葬他们的死人。}因此\Person{耶肋米亚}说：\BibleQuote{我要用我愤怒所击打的尘世死人充满它。}\BibleRef{耶 33:5}\par

因此，天主既然不是论证的对象，便不能成为科学的知识对象。但圣子是智慧、知识、真理，以及一切与此相关者。祂也是可以论证和描述的。圣神的所有能力，共同成为一体，都终止于同一点——即圣子。但祂就每一种能力的理念而言是无法言说的。圣子既不是单纯的一个事物，也不是许多部分，而是作为一切的一个事物；因此祂也是一切。因为祂是所有能力卷起并统一成一个统一体的圆。为此，圣言被称为\Quote{阿耳法}和\Quote{敖默加}，唯独在祂那里，终结成为开始，又毫无间断地在原初的起始处终结。因此，相信祂并借着祂，就成为一，在祂内不可分解地合一；而不相信则是分离、分开、分裂。\par

\BibleQuote{因此主这样说：每个外邦之子，在心灵上未受割损，在肉体上也未受割损}（即，在身体和灵魂上不洁）：\BibleQuote{没有一个外方人进入以色列家中，唯有肋未人。}\BibleRef{则 44:9} 祂称那些不愿相信、反倒不信的人为外方人。只有那些生活纯洁、成为天主真正司铎的人。因此，在一切受割损的支派中，那些被傅油为大司祭、君王和先知的人，被算为更圣洁。因此祂命令他们不可触摸尸体，或接近死者；这并不是因为身体被玷污，而是因为罪恶与悖逆成了肉身，成了具体，且是死的，因此可憎。于是，只有当父亲、母亲、儿子、女儿死亡时，司铎才能进入，因为这些人仅与血肉和种子相关，司铎因他们而获得进入生命的直接原因。他们洁净自己七天，这是创造完成的时期。因为在第七日庆祝安息；在第八日他带来赎罪祭，如厄则克耳所记载的，根据这赎罪祭，恩许将被领受。\BibleRef{则 44:27} 而完美的赎罪祭，依我看来，就是那在福音中的仁慈信德，它出自法律和先知，以及那在普遍服从中显示出来的纯洁，连同放弃世物；为的是那感恩地将帐幕交出，这帐幕来自灵魂的享乐。那么，无论这时间是通过所列举的七个时期回归于至高的安息，或是有些人所算的层层上升的七重天；又或那靠近理智力世界的固定球体被称为第八层，这一表述都表明\OriginalTerm{诺斯替者}{Gnostic}应当从受造界和罪恶的领域上升。在这七天之后，为罪献祭。因为仍有变化的恐惧，且它触及第七圈。义人\Person{约伯}说：\BibleQuote{我赤身出于母胎，也必赤身归回那里。}\BibleRef{约 1:21} 不是赤身无有财物，因为那是微不足道且常见的事；而是作为义人，他赤身离开邪恶与罪恶，以及那跟随行为败坏者的丑陋形状。因为这是所说的：\BibleQuote{你们若不悔改，变成如同小孩子一样，}\BibleRef{玛 18:3} 肉身纯洁，灵魂因禁戒恶行而圣洁；表明祂愿意我们成为如同祂从我们的母亲——水中——生出我们时一样。因为一代接一代的目的，是通过进步而成为不朽。\Quote{恶人的灯必要熄灭。} 诺斯替者所分有的身体与灵魂的纯洁，全智的\Person{梅瑟}藉着重复描述\Person{黎贝加}身上身体和灵魂的不朽而表明，如此说：\BibleQuote{那少女容貌俊美，且没有男子认识她。}\BibleRef{创 24:16} \Person{黎贝加}，解释起来，意思是\Quote{天主的荣耀}；而天主的荣耀是不朽。这实在是义，不是贪恋别物，而是完全成为被祝圣的主的殿堂。义是生命的平安和良好的状态，主在向她说\BibleQuote{去吧，平安地去吧}时，将她释放到这种状态中。\BibleRef{谷 5:34} 因为\Place{撒冷}，按解释，是平安；我们的救主被登记为撒冷王，正如梅瑟所说，\Person{默基瑟德}是撒冷王、至高天主的司祭，他拿出饼和酒，献上祝圣的食物，作为圣体圣事的预像。而\Person{默基瑟德}被解释为\Quote{正义的王}；这名字是义与平安的同义词。然而，\Person{巴西里德}却认为，正义和她的女儿平安驻留在第八层天球上。\par

但我们必须从物理学转向伦理学，后者更为清晰；因为关于这些的论述将在本论文之后接着讨论。救主亲自明确地引领我们进入奥秘，根据那悲剧的话：——\par

\Quote{看见那些看见者，他也赐予秘仪。}\par

若你问，\par

\Quote{这些秘仪，其性质是什么？}\par

你将再次听见：——\par

\Quote{未受入教礼的凡人不得知晓。}\par

若有任何人好奇地询问它们是什么，让他听吧：——\par

\Quote{你不可听，但值得知道；\newline{}神的礼仪憎恶行不敬的人。}\par

如今，那无始的天主，是宇宙的完美起始，也是起始的产生者。因此，正如祂是存有，祂是行动领域的第一原理；正如祂是善，祂是道德的第一原理；作为理智，另一方面，祂是推理和判断的第一原理。因此，也只有祂是导师，祂是至高圣父的独生子，人类的训导者。\par

% structure: p0186 source=h2 target=subsection reason=relative-heading-rank; base=h2

\subsection{第二十六章 完美之人如何对待身体和世物}

那麼，那些贬低受造的存在、诋毁身体的人是错误的；他们未曾考虑，人的身体被造成直立是为了默观上天，感官的结构是为了知识，而肢体和器官的安排是为了善，而非为了享乐。因此，这居所成为接纳那对天主至为宝贵的灵魂的容器，并因灵魂与身体的圣化而受圣神的尊崇，得以圆满救主的完美。在灵知者身上，可以找到三种德性的顺序：他道德地、物理地并逻辑地致力于天主。因为智慧是对神圣与人性事物的知识；义德是灵魂诸部分的和谐；圣德是事奉天主。但如果有人借引\Person{依撒意亚}的话说：\BibleQuote{凡有血肉的，都如草，人的一切光荣，如草上的花；草枯了，花也谢了；但主的话永远存留}，\BibleRef{依 40:6} 就让他听听圣神通过\Person{耶肋米亚}对此事的解释：\BibleQuote{我将他们像被风吹往旷野的干枯枝叶一样四散。这就是你不顺从的份和报应，主说。因为你忘记了我，信赖了谎言，所以我要将你的后裔暴露在你的面前；你的耻辱必被看见，你的奸淫，你的嘶鸣}，等等。\BibleRef{耶 13:24} 因为\Quote{草上的花}、\Quote{顺从肉身生活}和\Quote{属血肉的}（如宗徒所说）指的是那些陷于罪恶中的人。人的灵魂无疑是人的更好部分，身体是较低的部分。但灵魂并非本性为善，身体也并非本性为恶。并非不善的东西就是恶。因为有些事物处于中间状态，其中有可取的，也有可弃的。因此，人的结构（处于可感事物之间）必然由相异而非对立的事物——身体和灵魂——组成。\par

因此，善行总是属于更好、支配的灵，而纵欲和罪恶的行为则归于更坏、有罪的那个。\par

智者和灵知者的灵魂，犹如在身体中作客，严肃而恭敬地对待身体，不怀有过度的情感，好似若离别时刻召唤，就要离开这帐幕。经上说：\Quote{我是世上的客旅，是与你们同住的寄居者。}因此\Person{巴西里得}认为，被他理解的人（选民）是世界的异乡人，本性超尘世。但情况并非如此。因为万有都出于同一的天主。没有人在本性上是世界的异乡人，因为万有的本质是一，天主也是一。但那被选者作为寄居者生活，知道万有都是被拥有和管理的；他使用毕达哥拉斯派认为是三种善物的事物。身体，如同被派往远方朝圣者，在途中使用客栈和住所，关切世界的事物，他停留之处；但离开其居所和财产时不过分动情；欣然跟随那引他离开生命者；绝不、在任何时候都不回头；为他的寄居而感恩，为他的离去而祝福（天主），拥抱天上的宅邸。\BibleQuote{因为我们知道，如果我们地上的帐棚被拆毁了，我们必由天主获得一所房舍，一所非人手所造的、永远在天上的房舍。因为我们在这帐棚里呻吟，切望穿上我们那从天上来的住所，只要穿上，就不至于赤身了。因为我们因信德行走，并非因目睹}，如宗徒所说；\BibleQuote{我们更愿意离开身体，与主同在。}这\Quote{更}是比较而言。比较存在于相似的事物中，正如在勇士中更勇敢的人比较勇敢，在懦夫中则最为勇敢。因此他接着说：\BibleQuote{所以，我们无论住在身内，或离开身外，都努力使他喜悦} \BibleRef{格后 5:9}即天主，万有（包括世界和超世界的事物）都是他的工作和创造。我钦佩\Person{埃庇卡摩斯}，他明确地说：—\par

\Quote{怀有虔诚心智的人，在临终时，\newline{} 不会遭受任何邪恶。灵魂将居于天上；}\par

以及那位歌者所唱：—\par

\Quote{恶人的灵魂在地球上天空下飘荡，\newline{} 在不可避免的邪恶轭下忍受凶杀般的痛苦；\newline{} 但虔诚者的灵魂居于天上，\newline{} 在歌声中赞美那伟大、有福者。}\par

灵魂不是被从天上遣送到更坏的地方。因为天主使一切趋于更好。但选择了最佳生命——即来自天主和正义的生命——的灵魂，以地交换天。因此，有知识的约伯明智地说：\BibleQuote{现在我知道祢能做一切事；对祢没有不可能的事。因为谁告诉我我不知道的事，伟大而奇妙的事，我之前不熟悉？我自觉卑贱，视自己为尘土和灰烬。}因为那处于无知状态而有罪的人，\BibleQuote{是尘土和灰烬}；而那处于知识状态，尽可能肖似天主的人，已经是属神的，因此是被选者。而圣经称无感觉和悖逆的人为\BibleQuote{地}，将由耶肋米亚先知说明，他论及约雅金和他的弟兄说：\BibleQuote{大地，大地，请听上主的话；记下这人，当作被逐的人。} \BibleRef{耶 22:29} 另一位先知又说：\BibleQuote{天啊，请听；地啊，请侧耳倾听。} \BibleRef{依 1:2} 他把理解称为\BibleQuote{耳}，把\OriginalTerm{真知者}{Gnostic}的灵魂——即那致力于默观天界和神圣事物，并因此成为\OriginalTerm{以色列人}{Israelite}的人——称为\BibleQuote{天}。相反，他把那选择无知和心硬的人称为\BibleQuote{地}。而\BibleQuote{侧耳}一词源自\Quote{听觉器官}，\Quote{耳朵}，将属血肉的事归于那些依附感官事物的人。米该亚先知说到这样的人：\BibleQuote{请听上主的话，你们这些生活在痛苦中的百姓。}亚巴郎说：\BibleQuote{绝不做这事。上主是审判全地的主；} \BibleRef{创 18:25} \BibleQuote{因为不信的人，}按救主的话，\BibleQuote{已被定罪了。} \BibleRef{若 3:18} 在\WorkTitle{列王纪}中写着上主的审判和判决，如此说：\BibleQuote{上主听从义人，但祂不拯救恶人，因为他们不愿认识天主。}因为全能者不会做不合理的事。面对这话，异端们作何解释？既然圣经宣告全能的天主是善的，不是邪恶和错误的作者，如果无知确实源于人不认识，那么天主不做任何荒谬的事。\BibleQuote{因为这神，}经上说，\BibleQuote{是我们的神，除祂以外没有拯救者。} \BibleRef{依 45:21} \BibleQuote{因为在天主前没有不义，} \BibleRef{罗 9:14} 按宗徒所说。先知又清楚教导天主的旨意和真知的成熟，说：\BibleQuote{现今，以色列，上主你的天主向你要求什么？不过是敬畏上主你的天主，行走在祂的一切道路上，爱祂，单独事奉祂？} \BibleRef{申 10:12} 祂向你要求，你拥有选择救恩的能力。那么，\OriginalTerm{毕达哥拉斯学派}{Pythagoreans}要我们\Quote{出声祈祷}是什么意思？在我看来，并非他们认为神性不能听见沉默言语者，而是因为他们希望祈祷是正确的，没有人会在众人知晓的情况下羞愧地祈祷。然而，我们将在恰当的时候适时讨论祈祷。但我们应该有高声呼喊的善行，如同配得上\BibleQuote{在白天行走的人}。 \BibleRef{罗 13:13} \BibleQuote{让你们的善行照耀，} \BibleRef{玛 5:16} 看哪，一个人及其善行都在他面前。\BibleQuote{因为看哪，天主和祂的作为。} \BibleRef{依 62:11} 因为真知者必须尽可能效法天主。而诗人们在他们的作品中称被选者为神般的人、神、与神同等、有与\OriginalTerm{宙斯}{Zeus}同等的智慧、有神的谋略、肖似神——在我看来，这略为触及了\BibleQuote{按天主的肖像和模样}的说法。 \BibleRef{创 1:26}\par

\Person{欧里庇得斯}据此说：\Quote{黄金之翼绕我背，我穿\OriginalTerm{塞壬}{Sirens}的翼履；我将飞升入广阔苍穹，与\OriginalTerm{宙斯}{Zeus}相会。}\par

但我要祈求基督的圣神使我飞到我的耶路撒冷。因为\OriginalTerm{斯多葛学派}{Stoics}说，天正是城市，但地上的地方不是城市；因为它们被称为城市，其实不是。因为城市是重要的东西，人民是整齐的团体，是受法律约束的人群，如同教会受圣言约束—一座地上不可攻破的城市—不受暴政统治；是天上人间神圣意志的产物。诗人们用笔墨创造这城市的形象。因为\OriginalTerm{极北族人}{Hyperboreans}、\OriginalTerm{阿里玛斯比亚人}{Arimaspian}的城市、以及\OriginalTerm{厄吕西翁原野}{Elysian plains}，都是义人的共和国。我们知道\OriginalTerm{柏拉图}{Plato}的城市被放置在天空中作为模范。\par

\SourceRecord{Translated by William Wilson.；Ante-Nicene Fathers；Vol. 2.；Edited by Alexander Roberts, James Donaldson, and A. Cleveland Coxe.；Buffalo, NY: Christian Literature Publishing Co.,；1885.；<http://www.newadvent.org/fathers/02104.htm>.}

% structure: p0001 source=h1 target=section reason=page-title

\section{\WorkTitle{杂集（第五卷）}}

% structure: p0002 source=h2 target=subsection reason=relative-heading-rank; base=h2

\subsection{第一章  论信德}

关于真知者，我们已略有所述。现在继续后续，必须再次思考信德；因为有些人作此区分：信德关乎圣子，知识关乎圣神。但他们没有注意到，为了真正相信圣子，我们必须相信祂是圣子，相信祂来了，以及如何来，为何来，并相信祂的苦难；我们还必须知道谁是圣子。知识没有信德不行，信德没有知识也不行。圣父没有圣子也不行，因为圣子与圣父同在。圣子是教导有关圣父的真教师；而为了相信圣子，我们必须认识圣父，圣子也与圣父同在。再者，为了认识圣父，我们必须相信圣子，相信是圣子教导；因为从信德到知识，经由圣子，便是圣父。而按照真知准则——那真正是真知的——对圣子和圣父的认识，就是藉着真理达成并领悟真理。\par

我们，就是那些相信不被相信之事的人，是那些认识不被认识之事的人；也就是说，是那些认识众人不信也不认识、但少数人相信并认识之事的人；是真知者，不是以言语描述行为，而是在默观中操练的真知者。在听众耳中说话的人是有福的。信德是灵魂的耳朵。主暗示信德即是如此，当祂说：\Quote{有耳能听的，听吧！}\BibleRef{玛 11:15}；好叫人藉着相信能领悟祂所说的，如祂所说。最古老的诗人荷马也使用\Quote{听}字代替\Quote{领悟}——以属指代类——写道：\par

\Quote{他们最听他。}\par

总而言之，双方信德的一致与和谐趋向同一目标——救恩。我们在宗徒那里有一个无疑的见证：\Quote{因为我切愿见你们，把一些属神的恩赐分给你们，为使你们坚固；也就是说，我在你们中间，因你我所共有的信德，同得安慰。}\BibleRef{罗 1:11}接着他又补充说：\Quote{天主的正义从信德到信德显示出来。}\BibleRef{罗 1:17}那么，宗徒显然宣布了一种双重的信德，或者不如说是一种能够成长与完善的信德；因为普通的信德埋藏在底下作为根基。因此，对那些愿意被治愈、并被信德感动的人，祂加上说：\Quote{你的信德救了你。}\BibleRef{玛 9:22}但那被优美地建造起来的信德，在信徒身上得以成全，又藉由教训与圣言所产生的信德而得以完善，为遵守诫命。宗徒们便是如此，在他们身上，据说\Quote{信德移山倒树}。因此，他们领悟到其力量的伟大，便求\Quote{增加他们的信德}\BibleRef{路 17:5}；这信德有益地扎根于土壤，\Quote{像一粒芥菜子}，并在此中优美地生长，以至于崇高事物的理由都寄托其上。因为如果人本性上认识天主，如巴西里得所认为的——他称高级的理智同时为信德和王权，以及相称于造物主本质的受造；并解释说，在祂近旁不是能力，而是本质、本性和实体；且说信德不是灵魂行使自由意志的理性认同，而是直接属于受造物的未定义之美——那么，旧约和新约的诫命便是多余的，如果人像瓦伦提努所主张的那样，是因本性而得救，又像巴西里得所认为的，是因本性而成为信徒和选民；那么本性总有一天能够无需救主的显现而闪耀。但若他们说救主的到访是必要的，那么他们的本性特质便丧失了，选民藉由教训、净化及善工而得救。亚巴郎听见那在玛默勒橡树下应许的声音：\Quote{我要把这地赐给你和你的后裔}，他便藉着听而相信了；他或者是选民，或者不是。如果他不是，他怎能立刻相信，仿佛出于本性？如果他是选民，他们的假设便被废弃，因为连在主来临之前也已发现有选民，且得救了：\Quote{就以此算为他的正义。}\BibleRef{创 15:6}因为如果有人追随玛尔强，敢说\OriginalTerm{造物主}{\GreekText{Δημιουργόν}}在主来临之前就已拯救了相信祂的人（选民以他们自己的救恩得救）；那么善者的能力便黯然失色；因为它只在后来，且晚于他们所说给予吉祥词语的造物主，才试图去拯救，并且藉着教训和模仿他。但如果善者如此拯救，按他们的说法，他拯救的不是自己的人，也不是在创造万物的那一位同意之下尝试救恩，而是藉着强迫或欺骗。那么他既然后来才行动，怎能仍然是善的？但如果地域不同，全能者的住所远离善天主的住所，然而那位拯救者先开始行动的意志，并不亚于善天主。从先前已证明的来看，不信的人被证明是无知的：\Quote{他们的道路弯曲，他们不知平安，}先知说。\BibleRef{依 59:8}\Quote{至于愚拙无学问的辩论，}神圣的保禄劝勉\Quote{要避免，因为能生争论。}\BibleRef{弟后 2:23}埃斯库罗斯也宣告：\par

\Quote{在无益的事上，不要徒劳。}\par

因为那与信德相符的探究，它在信德的根基上建立起真理的庄严知识，我们知道这是最好的。现在我们知道，既不是那些清楚的事物成为探究的对象（例如问是否白天，而当时正是白天），也不是那些未知且永远不会变得清楚的事物（例如星星的数目是偶数还是奇数），也不是那些可转换的事物；所谓可转换，是指那些可以由持相反意见的人同样说出的事情，例如子宫中的东西是否是活物。第四种情况是，当从双方中的任何一方提出一个无法回答且不可反驳的论据时。那么，如果探究的根据根据所有这些模式都被移除，信德就确定了。因为我们向他们提出一个无法回答的考虑：是天主在说话，并以文字前来帮助我们，涉及我所探究的每一个点。那么，谁如此不敬虔，竟不相信天主，并要求天主像人一样提供证据？再者，有些问题需要感官的证据，例如问火是否是热的，雪是否是白的；有些则需要劝诫和责备，例如是否应该孝敬父母的问题。还有一些问题是该受惩罚的，例如要求提供天主眷顾存在的证据。既然有天主眷顾，那么认为全部预言和关乎救主的经济并非按照天主眷顾而发生，就是亵渎的。或许甚至不应该试图证明这些点，因为神圣的天主眷顾从其所有可见的巧妙而智慧的作品中显而易见，其中一些按秩序发生，一些在秩序中显现。那赐予我们存在和生命者，也赐予我们理性，愿意我们理性而正直地生活。因为宇宙之圣父的圣言不是\OriginalTerm{发出的言语}{\GreekText{λόγος} \GreekText{προφορικός}}，而是天主的智慧和最显明的仁慈，也是祂的大能，这大能是全能的、真正神圣的，并且并非不能被那些不宣认的人所理解——那全能的意志。但由于有些人不信，有些人好辩，并非所有人都达到美善的圆满。因为若不运用自由意志，就不可能达到它；但也不完全取决于我们自己的意图；例如，什么是注定要发生的。\Quote{因为我们因恩宠得救：}确实不是没有善行；但我们必须通过为善而被塑造，获得对善的倾向。我们还必须拥有那专注于追求美善的健康心智；为此我们极其需要天主的恩宠、正确的教导、神圣的感召力以及圣父的吸引。因为，被束缚在这个属地的身体中，我们通过身体感知感官对象；但我们通过逻辑官能本身把握理智对象。但如果有人期望通过感官把握一切，他就远离了真理。因此，属灵地，宗徒论及认识天主写道：\BibleQuote{因为现在我们如同透过镜子观看，但那时就要面对面了。}\BibleRef{格前 13:12}因为真理的异象只赐给少数人。因此，柏拉图在\WorkTitle{爱庇诺米斯}中说道：\Quote{我不说所有人都可能蒙福和快乐；只有少数。当我们活着时，我宣布这是实情。但死后我有美好的盼望，我将获得一切。}梅瑟中所载与此相同：\BibleQuote{没有人能见到我的面而存活。}\BibleRef{出 33:20}因为显然，在生命期间没有人能够清楚地领会天主。但\BibleQuote{心地纯洁的人要看见天主。}\BibleRef{玛 5:8}当他们达到最终的圆满时。因为灵魂变得过于衰弱，无法把握现实，我们需要一位神圣的老师。救主被派遣下来——一位获取美善的老师与领袖——伟大天主眷顾的秘密而神圣的标志。\Quote{那么，文士在哪里？今世的探究者在哪里？难道天主不是使今世的智慧成为愚妄吗？}\BibleRef{格前 1:20}圣经如此说。又说：\Quote{我要摧毁智慧人的智慧，废弃明哲人的明哲。}\BibleRef{格前 1:19}这显然是指那些自以为聪明、好辩的人。因此耶肋米亚说得极好：\Quote{主这样说：站在路上，探问永恒的道路，哪一条是善道，就在其上行走，你们必为你们的灵魂找到赎罪。}\BibleRef{耶 6:16}他说，询问、请教那些知道的人，不要争论和争辩。在学到真理之路后，让我们行在正道上，不转离，直到我们达到所渴望的。因此，罗马国王（名叫努马），是一位毕达哥拉斯主义者，在所有人中首先为信德与和平建造了一座圣殿，这是有理由的。\Quote{亚巴郎因信德，被算为正义。}他，钻研高深的大气现象哲学和天上运动的神圣哲学，被称为亚巴郎，解释为\Quote{崇高的父亲}。但后来，他仰望天空，无论是他在圣神中看见了圣子（如一些人解释），或是一位荣耀的天使，或以其他方式认出天主超越受造物及其中的一切秩序，他另外接受了\Quote{阿尔法}，即对唯一真天主的认识，被称为亚巴辣罕，不再是自然哲学家，而是成为智慧者和爱天主者。因为这解释为\Quote{声音的蒙选之父}。因为声音是指发出的言语：心灵是它的父亲；善人的心灵是蒙选的。我忍不住极其赞美阿格里根特的诗人，他这样歌颂信德：——\par

\Quote{朋友们，我知道，在我将要讲述的神话中有真理。\newline{}我将讲述它们。但对人来说，信德的尝试非常困难，\newline{}而且令心灵厌烦。}\par

为此，宗徒也劝勉说：\BibleQuote{你们的信德不该在于人的智慧，}那些以说服为业的人，\BibleQuote{而应在于天主的能力，}\BibleRef{格前 2:5}这能力无需证据，仅凭信德就能拯救。\Quote{那些自命不凡的人中最有见识的知道如何保持警惕。公义将捉拿说谎的伪造者和见证人，}厄弗所人说。因为他从外邦哲学中获得知识，熟悉对行为败坏之人的火炼，斯多葛派后来称之为\OriginalTerm{大火}{\GreekText{ἐκπύρωσις}}，他们也教导每个人将完全如其原样复活，如此论及复活；而柏拉图则如下说，大地在特定时期被火和水净化：\Quote{人类曾有许多种毁灭方式，且将有极大的毁灭由火和水而来；其他较短暂者则有无数原因。}稍后他补充说：\Quote{事实上，围绕大地和天体的诸物也有变化；在漫长周期的进程中，大地上的物体因一场大火而毁灭。}接着他论及洪水说：\Quote{然而，当天神们再次以洪水淹没大地以用水净化时，山上那些牧人和牧羊人被拯救；而你们城中的那些人则被河流冲入海中。}我们在第一本杂集中已指出，希腊哲学家被称为窃贼，因为他们未加承认地从梅瑟和先知那里窃取了主要教义。对此我们还要补充：那些获得更高职位的天使，因沉溺于享乐，将所知的秘密告诉了妇女；而其余的天使则隐藏了它们，或者说，将它们保留至主的来临。由此产生了眷顾的教义和高深之事的启示；预言既已赋予希腊哲学家，关于教义的研究便在哲学家中兴起，有时正确（当他们击中目标时），有时错误（当他们不理解预言之比喻的奥秘时）。这一点，我们打算在概述需要提及的事项时简要指出。那么，我们说，信德不应是停滞不前的，而应伴随探究。因为我并非说我们根本不应探究。因为经上说：\BibleQuote{寻找，就会找到。}\BibleRef{玛 7:7}\par

\Quote{所求之物可得捕获，\newline{}但被忽视者则逃逸。}\par

据\Person{索福克勒斯}所言。\par

喜剧诗人\Person{梅南德}也说类似的话：——\par

\Quote{一切所寻，\newline{}最智者说，需焦虑思虑。}\par

但我们应当正确引导灵魂的视觉官能去发现，并清除障碍；完全抛弃争辩、嫉妒和纷争，这些注定要在人类中悲惨地消亡。\par

\Place{弗利乌斯}的\Person{提蒙}写得非常优美：——\par

\Quote{纷争，凡人的灾祸，徒劳地尖叫行进，\newline{}是凶杀争吵与不睦的姊妹，\newline{}它盲目地翻卷一切。但随后\newline{}它将头转向人类，将他们抛向希望。}\par

稍后他补充道：——\par

\Quote{因为谁使这些陷入致命的斗争？\newline{}一群与回声同步的乌合之众；因对沉默者发怒，\newline{}它兴起一场恶疾攻击人类，许多人丧亡；}\par

至于否认虚假的言论、两难之境、隐藏之物、\OriginalTerm{连锁推理}{Sorites}和\OriginalTerm{鳄鱼推理}{Crocodilean}、显露之物、以及含混和诡辩。因此，探究天主，若不趋向争辩而趋向发现，则是有益的。因为达味经上记载：\BibleQuote{贫乏的人吃了，且得饱饫；寻求主的人将赞美祂。你们的心将永远生活。}因为那些以真正寻求寻求祂的人，赞美主，将被来自天主的恩赐所充满，即知识。他们的灵魂将生活；因为灵魂比喻性地称为心，它赋予生命：因为通过子，圣父被认识。\par

我们不应将自己的耳朵交托给所有轻率说话和写作的人。因为许多被耳朵抓住的杯子也会变脏并失去耳朵；此外，当它们跌落时就破碎了。同样，那些以许多琐事玷污了信德的纯净听觉的人，最终对真理变成聋子，变得无用并跌落地上。因此，我们吩咐孩童在亲吻亲人时，要抓住他们的耳朵，这并非没有理由；表明爱的情感是由听觉产生的。而\BibleQuote{天主就是爱}\BibleRef{若一 4:16}，正如\Quote{天主是信实的}；我们必须藉天主的爱而与祂结合：好能以相似看见相似，天真纯洁地聆听真理之言，如同顺服我们的孩童。而这就是那位在\Place{埃皮达鲁斯}神庙入口处刻下如下铭文的人（无论他是谁）所指示的：——\par

\Quote{进入那馨香圣殿者，\newline{}必须纯洁。}\par

纯洁就是\Quote{怀着圣洁的思想}。\BibleQuote{除非你们变成如同这些小孩子，你们不能进入，}经上说，\BibleQuote{天国。}因为在那里，天主的圣殿被看见建立在三个基础上——信德、望德和爱德。\par

% structure: p0024 source=h2 target=subsection reason=relative-heading-rank; base=h2

\subsection{第二章——论望德}

关于信德，我们已在希腊人的作品中引用了充分的证据。但为了不因热心收集大量关于望德和爱德的证据而越出界限，只需提及\WorkTitle{克里托篇}中\Person{苏格拉底}的话就足够了，他更偏爱善生善死而非生命本身，认为我们死后有另一个生命的希望。\par

他在《\WorkTitle{费德若篇}》中亦说：\Quote{只有在分离的状态下，灵魂才能分沾那真实且超越人类智慧的智慧；当藉着哲学之爱达到望德的终点，渴望将把它吹向天堂时，}他说，\Quote{它便接受了另一种不朽生命的开端。}在《\WorkTitle{会饮篇}》中他说：\Quote{一切天赋中皆注入了对生成相似者的自然之爱，在人类中只是生成人，在善人中则生成他自己的对应者。但善人若不拥有完美的德行，便不可能做到这一点，他将在那些投奔他的青年中培育这些德行。}正如他在《\WorkTitle{泰阿泰德篇}》中所说：\Quote{他将生育并成全人。因为有的藉身体生育，有的藉灵魂生育；}因为在异邦哲学家那里，教导和启迪被称为再生；并且\Quote{我在耶稣基督内生了你们}\BibleRef{格前 4:15}，那位良善的宗徒在某处说道。\par

恩培多克勒也将友爱列为元素之一，视之为一种结合的爱：——\par

\Quote{你用心灵去看它，不要张着嘴呆看。}\par

巴门尼德也在他的诗中暗示望德，这样说：——\par

\Quote{然而你定要用心灵去看那不在场的如同在场，\newline{}因为它不会将那存在者从其把握中分离，\newline{}即使它散遍世界各处或被结合起来。}\par

% structure: p0031 source=h2 target=subsection reason=relative-heading-rank; base=h2

\subsection{第三章 信德与望德的对象唯以心灵感知}

因为那怀有望德的人，如同那相信的人，是以心灵感知理智的对象和未来的事物。那么，如果我们肯定某物是正义的，肯定它是善的，并且我们也说真理是某种东西，但我们从未用眼睛见过这些对象中的任何一个，而是仅以心灵。如今天主的圣言说：\Quote{我是真理。}\BibleRef{若 14:6}圣言因此必须被心灵所默观。\Quote{你断言，}有人说，\Quote{有真正的哲学家吗？}\Quote{是的，}我说，\Quote{那些热爱默观真理的人。}柏拉图在《\WorkTitle{费德若篇}》中谈及真理时，也将它作为理念来展示。理念是天主的观念；异邦人称之为天主的圣言。原文如下：\Quote{因为那时必须敢于说出真理，尤其在谈论真理时。因为灵魂的本质是无色、无形、不可触摸的，唯有其向导天主才能看见。}现在，发出的圣言是创造的原因；然后祂也生了自己，\Quote{圣言成了血肉}\BibleRef{若 1:14}，为能被人看见。义人将寻求那源自爱德的发现，若他急忙奔赴，便会成功。因为经上说：\Quote{给敲门的就开门：你们求，就给你们。}\BibleRef{玛 7:7}\Quote{因为用强暴夺取天国的人}\BibleRef{玛 11:12}并非在争辩的言谈中如此；而是藉着持续的正直生活和不间断的祈祷，被说成\Quote{用强夺取}，擦去他们先前罪过留下的污点。\par

\Quote{你可以获得邪恶，甚至大量地获得。\newline{}而劳作的人，天主帮助他；\newline{}因为缪斯的礼物，难以赢得，\newline{}并不摆在任何人面前，任他拿走。}\par

因此，对无知的认识，是依照圣言行走的第一步。无知的人寻求了，寻求便找到了老师；找到了便相信了，相信了便怀有望德；从此以后爱了，便与所爱的相似——努力成为他最初所爱的。苏格拉底教导阿尔喀比亚德的方法即是如此，后者这样问道：\Quote{你不认为我将以别的方式知道什么是正确的吗？}\Quote{是的，如果你找到了。}\Quote{但你不认为我找到了吗？}\Quote{当然，如果你寻求了。}\par

\Quote{那么你不认为我寻求过吗？}\Quote{是的，如果你认为你不知道。}同样，明智童女们的灯在无知的大黑暗中点燃，圣经用\Quote{黑夜}来象征。智慧的灵魂，纯洁如童女，明白自己处于世界的无知中，便点燃灯，唤醒心灵，照亮黑暗，驱散无知，寻求真理，等候老师的出现。\par

\Quote{那么，群众，}我说，\Quote{不能成为哲学家。}\par

\Quote{持杖者众多，但巴克斯信徒却少，}依柏拉图所说。\Quote{因为蒙召的人多，但被选的人少。}\BibleRef{玛 20:16}\Quote{知识并非人人都有}\BibleRef{格前 8:7}，宗徒说。\Quote{并请为我们祈祷，使我们脱离无理和邪恶的人，因为不是人人都有信德。}\BibleRef{得后 3:1}斯多葛派的克莱安塞斯在他的《\WorkTitle{诗学}》中写道：——\par

\Quote{不要追求荣耀，想要突然变得智慧，\newline{}也不可畏惧多数人无辨识力和轻率的意见；\newline{}因为大众没有理智、智慧或正确的判断，\newline{}你只能在少数人中发现这些。}\par

而喜剧诗人更精辟地简短说道：——\par

\Quote{凭喧闹来判断正确之事是可耻的。}\par

因为他们，我想，听到了那卓越的智慧，它对我们说：\Quote{在愚昧人当中要把握你的机会，在明智人当中要持续。}\BibleRef{德 27:12}又说：\Quote{智慧人隐藏自己的知识。}\BibleRef{箴 10:14}因为多数人要求证明作为真理的担保，不满足于单纯的因信德得救。\par

\Quote{但坚决不信那盛行的邪恶，\newline{}正如我们缪斯的保证所吩咐的，\newline{}藉着剖析你胸中的言语来认识。}\par

\Quote{因为这是恶人的习性，}恩培多克勒说，\Quote{想以不信来压倒真理。}而我们的教义是可信的、值得相信的，希腊人将在后续更深入的研究中知道这一点。因为我们由相似者教导相似者。因为撒罗满说：\BibleQuote{要照愚昧人的愚妄回答他。}\BibleRef{箴 26:5}因此，对于那些询问我们智慧的人，我们要提出适合他们的东西，以便他们能极其容易地通过自己的观念而达到对真理的信仰。因为\BibleQuote{我成了众人中之一切，为的是要救得众人。}\BibleRef{格前 9:22}既然神圣恩宠的\BibleQuote{雨}也降在\BibleQuote{义人及不义的人}身上。\BibleRef{玛 5:45}\BibleQuote{难道天主只是犹太人的天主吗？不也是外邦人的天主吗？是的，也是外邦人的天主，因为天主只有一个，}\BibleRef{罗 3:29}——尊贵的宗徒如此宣告。\par

% structure: p0044 source=h2 target=subsection reason=relative-heading-rank; base=h2

\subsection{第四章 神圣事物在圣经与外教作者中皆以象征隐藏}

但是既然他们既不信正直的善，也不信得救的知识，我们认定他们声称属于我们的东西，因为一切都是天主的；尤其因为从我们传到希腊人的善，让我们按他们所能听的方式处理这些事。因为大众不是以真理，而是以他们喜欢的东西来判断智慧和正义。他们不会比与自己相似的事物更高兴。因为那仍然盲目和哑口的人，没有理解力，也没有救主所赐的沉思灵魂的无眩目光明，就像未参与奥秘的人或未参与舞蹈的不懂音乐者，尚未纯洁、不配纯洁的真理，仍然不协和、无序且属物质，必须站在神圣的合唱团之外。\BibleQuote{因为我们以属神的事解释属神的事。}\BibleRef{格前 2:13}因此，按照隐藏的方式，真正神圣的圣言——真正神圣且对我们最为必要——存放在真理的圣所中，被埃及人用他们所谓的\OriginalTerm{圣所}{adyta}表示，被希伯来用帷幔表示。只有祝圣过的人——也就是那些奉献给天主，为了对唯一神圣者的爱而在情欲的欲望上受割损的人——才能接近它们。因为柏拉图也认为\Quote{不洁者接触纯洁者}是不合法的。\par

由此，预言和神谕以隐语说出，奥秘不是立即展示给所有人，而是经过一定的净化和前期指导之后。\par

\Quote{因为缪斯那时并非\newline{}贪图利益或做买卖；\newline{}忒耳普西科瑞的甜蜜、\newline{}蜂蜜调子、银铃般柔声的\newline{}旋律也没有被商品化。}\par

在埃及受教导的人首先学习那种称为书信体的埃及字母；其次是僧侣体，为神圣抄写员所用；最后是象形文字，其中一种通过基本元素的是直译体，另一种是象征体。在象征体中，一种是通过模仿直接说，另一种是像形似地书写；还有一种完全是寓言式的，使用某些隐语。\par

他们想用文字表达太阳时，画一个圆圈；表达月亮时，画一个像月亮的形状，类似其自身形状。但在使用比喻风格时，他们通过转换、转移、改变和多种变形来绘制字符。在神学神话中叙述国王的赞美时，他们用阳刻浮雕书写。让以下作为第三种类型——隐语体的标本。至于其他星辰，由于其倾斜的轨迹，他们将其画成蛇的身体；而太阳则像甲虫，因为它制造一个牛粪的圆球，并滚在面前。他们说这种生物在地下生活六个月，另一半年份在地上生活，并将种子射入球中，然后生育；而且没有雌性甲虫。总之，所有谈到神圣事物的人，无论是蛮族还是希腊人，都隐藏了万物的首要原理，并以隐语、象征、寓言、比喻和类似的修辞手法传达真理。这样的例子还有希腊人的神谕。皮提亚的阿波罗被称为\OriginalTerm{洛西阿斯}{Loxias}。此外，希腊人中被称为智者的格言，用简短的话语指示了重要内容的展开。当然，有这样的格言：\Quote{珍惜时间：}要么是因为生命短暂，我们不应虚度光阴；要么它命令你节省个人开支，以便即使你活很多年，必需品也不会缺乏。同样地，格言\Quote{认识你自己}也表明了许多事：你是有死的，你是人类出生的；而且，在生活的其他长处比较中，你无足轻重，因为你说你富有或出名；或者另一方面，即使你富有或出名，你也不会仅仅因为你的优势而受到尊敬。它说：要知道你为何而生，你是谁的肖像；你的本质是什么，你的受造是什么，你与天主的关系如何，以及诸如此类。圣神藉着先知依撒意亚说：\BibleQuote{我要将隐藏的宝物、黑暗中的宝藏赐给你。}\BibleRef{依 45:3}智慧是难以追寻的，是天主的宝藏和永不匮乏的财富。但那些受这些先知教导神学的人——诗人——也以隐义的方式大量创作哲学。我指的是俄耳甫斯、利诺斯、穆萨埃乌斯、荷马、赫西奥德，以及这些类型的智者。诗歌的说服风格对他们来说是蒙蔽众人的面纱。梦和征兆对人多多少少是模糊的，不是出于嫉妒（因为认为天主有情感是错误的），而是为了研究引入对隐语的理解，加速发现真理。因此索福克勒斯这位悲剧诗人在某处说：\par

\Quote{我知道天主是这样的：\newline{}永远向智者启示隐语，\newline{}但对悖逆的恶者，即使话语简短作为导师，}\par

以坏代替简单。然后明确地关于我们全部的圣经，如同用比喻说出，在\WorkTitle{圣咏}中写道：\BibleQuote{我的百姓阿，你们要听我的法律；请侧耳听我口中的言语。我要开口说比喻，我要从起初发出我的谜语。}同样，高贵的宗徒也这样说：\BibleQuote{然而，在成全的人中，我们讲论智慧，却不是这世上的智慧，也不是这世上将要消灭的掌权者的智慧；我们讲的，是天主隐藏的、奥秘的智慧，这智慧是这世上的掌权者中没有一个知道的，因为如果他们知道了，就不会把荣耀之主钉在十字架上。}\BibleRef{格前 2:6}\par

哲学家们并没有努力轻视主的显现。因此，宗徒所斥责的是犹太人中的智者之见解。所以他补充说：\BibleQuote{然而，正如经上所载：天主为爱祂的人所预备的，是眼睛所未见，耳朵所未闻，人心所未曾想到的。但天主藉着圣神将这一切启示给我们，因为圣神洞察一切，就连天主的深奥也洞悉。}\BibleRef{格前 2:9}因为他承认属神的人和\OriginalTerm{诺斯替者}{Gnostic}是天主所赋的圣神的门徒，这圣神就是基督的心意。\BibleQuote{然而属血气的人不接受天主圣神的事，因为为他是愚妄。}\BibleRef{格前 2:14}现在，宗徒与诺斯替的成全相对比，称普通的信德为\Emph{根基}，有时又称之为\Emph{奶}，他这样写道：\BibleQuote{弟兄们，我从前对你们说话，不能把你们当作属神的，只得把你们当作属血肉的，当作在基督内的婴孩。我给你们喝奶，没有给你们吃饭，因为那时你们还不能吃，就是如今还是不能，因为你们仍是属血肉的。因为在你们中间有嫉妒、纷争，你们岂不是属血肉的，照着人的样子行事么？}\BibleRef{格前 3:1}这些事是那些罪人的选择。但那戒绝这些事的人，思想神圣的事，分享诺斯替的食物。\BibleQuote{按照所赐给我的恩宠，}宗徒说，\BibleQuote{我好像一个明智的建筑师，奠定了根基，别人在上面建造金、银、宝石。}\BibleRef{格前 3:10}这就是在基督耶稣的信德根基上的\OriginalTerm{诺斯替的}{gnostic}上层建筑。但\Quote{草、木、禾秸}是异端的添加物。\Quote{但火要试验各人的工程是哪种。}暗指诺斯替的大厦，在\WorkTitle{致罗马人书}中他也说：\BibleQuote{因为我切切想见你们，要把属神的恩赐分给你们，使你们得以坚固。}\BibleRef{罗 1:11}这样的恩赐不可能不蒙遮掩而写出来。\par

% structure: p0053 source=h2 target=subsection reason=relative-heading-rank; base=h2

\subsection{第五章 论\Person{毕达哥拉斯}的象征}

现在，\Person{毕达哥拉斯}的象征与蛮族哲学有最隐晦的联系。例如，撒摩斯人劝告\Quote{不要在屋内养燕子}；也就是说，不要接纳多嘴、低声私语、饶舌、不能保守秘密的人。\Quote{因为燕子、斑鸠和田野的麻雀，都认识它们进入的时候，}\BibleRef{耶 8:6}圣经说；而人绝不应与琐碎之事同居。斑鸠咕咕叫，显示抱怨的忘恩负义，理应被赶出屋外。\par

\Quote{不要坐在一处又一处地对我低声怨言。}\par

燕子也使人想起\Person{潘狄翁}的寓言，因为它所记载的事件令人憎恶，其中一些我们听说\Person{特柔斯}所遭受，一些是他所施加的。它也追逐音乐的蝉，因此迫害圣言的人应当被赶走。\par

\Quote{秉执权杖的\Person{赫拉}，其目光遍察\Place{奥林帕斯}，\newline{}我有一个生锈的舌头橱柜，}\par

诗歌这样说。\Person{埃斯库罗斯}也说：—\par

\Quote{但是，我也有一个钥匙作为我舌头的守卫。}\par

再次，\Person{毕达哥拉斯}命令：\Quote{当锅从火上提起时，不要在灰烬中留下它的印记，而要把它抹平；}以及\Quote{人们起床后，要抖一抖床单。}因为他暗示，不仅需要除去印记，而且连愤怒的痕迹也不要留下；当它停止沸腾时，要平息，一切受伤害的记忆都要擦去。圣经说：\BibleQuote{不要让太阳在你们的愤怒中落下。}\BibleRef{弗 4:26}那说过\BibleQuote{不可贪恋}\BibleRef{出 20:17}的，除去了所有错误的记忆；因为愤怒在温和的灵魂中是贪欲的冲动，尤其寻求非理性的报复。同样，\Quote{床单要被抖开}，以便没有在睡眠中流泄的回忆，也没有白天的睡眠；此外，也没有夜间的快乐。他暗示，黑暗的幻觉应该被真理的光迅速驱散。\Person{大卫}说：\BibleQuote{你们愤怒却不要犯罪。}教导我们不可赞同冲动，也不可行动而追随它，从而巩固愤怒。\par

再次，\Quote{不要在地上航行}是\Person{毕达哥拉斯}的谚语，表明应该拒绝捐税和类似的契约，因为它们麻烦且不稳定。所以圣言也说，收税的人很难得救。\par

再者，\Quote{勿戴戒指，也不可在其上雕刻神祇的像}，这是毕达哥拉斯所吩咐的；正如摩西早在多时以前就明文规定，不可制造雕刻的、铸成的、模制的或绘制的像，好叫我们不执着于可感觉的事物，而转向理智的对象：因为习于目视会贬损对神性事物的敬畏；而藉着物质去崇拜那非物质的，便是以感官羞辱祂。为此，埃及最有智慧的祭司决定将雅典娜的神庙建成\OriginalTerm{露天的}{hypæthral}，正如希伯来人建造了没有偶像的圣殿。有些人在崇拜天主时，制作一个包含星辰的天象图，并以此朝拜，尽管圣经说：\Quote{让我们按我们的肖像，按我们的模样造人。}\BibleRef{创 1:26} 我以为也值得引述毕达哥拉斯学派的欧律苏斯的话如下，他在其著作\WorkTitle{论命运}中说，\Quote{造物主在造人时，以自己为典范}，接着又说：\Quote{身体与其他事物相似，因为它由同样的材料制成，并由最卓越的工匠塑造，他以自己为\OriginalTerm{原型}{archetype}而制作了它。}总之，毕达哥拉斯及其追随者，还有柏拉图以及大多数其他哲学家，最熟悉那位立法者，这可以从他们的学说中得出结论。藉着一种神圣的占卜之言，并非没有天主的助力，他们赞同某些先知性的宣告，并在部分和方面中把握真理，以并非隐晦的措辞，且不超越对事物的解释，当他们确认了与真理的关系时，便尊崇它。因此，希腊哲学就像一支人们点燃的灯芯，从太阳那里盗取了光。但自\OriginalTerm{圣言}{Word}宣告以来，那全部圣洁的光辉便照耀出来。夜间在房屋里，盗取的光有用；但白昼火焰燃烧，整夜被这样一颗理智之光的太阳照亮。\par

毕达哥拉斯在他的一句话中概括了摩西关于正义的论述，他说：\Quote{不要跨过天平；}即不要逾越分配上的平等，如此尊崇正义。\par

\Quote{那将朋友永远连结于朋友，\newline{}城邦于城邦——盟邦于盟邦，\newline{}因平等对人而言即是正义；\newline{}但较少的总是与较大的为敌，\newline{}仇恨的日子便开始，}\par

正如以诗意的优雅所说。\par

因此，主说：\Quote{你们背起我的轭，因为它是柔和的，是轻省的。}\BibleRef{玛 11:29} 当门徒争大时，祂以简朴的方式命令平等，说\Quote{他们必须变成像小孩子一样。}\BibleRef{玛 18:3} 同样，宗徒也写道：\Quote{在基督内，没有为奴的或自由的，也没有希腊人或犹太人。因为在基督耶稣内的受造乃是新的，是平等、无争、不贪求，是正义的。}因为嫉妒、怨恨和苦毒，都站在神圣的合唱团之外。\par

同样，那些精通奥秘的人禁止\Quote{吃心脏}，教导我们不应因违背心愿之事而藉着懈怠和烦恼啃噬并消耗灵魂。因此，荷马所说的那人是可悲的，他独自游荡，\Quote{吃自己的心。}但是，既然福音设想了两条路——宗徒们也是如此，与所有先知一样——并称那一条为\Quote{狭窄而拥挤的}，那是按照诫命和禁令而限定的；而相反的那条通向丧亡的路是\Quote{宽阔而宽敞的}，向享乐和忿怒敞开。他们又说：\Quote{不随从恶人的计谋，不站在罪人的道路上的人是有福的。}由此也产生了凯乌斯的普罗狄库斯关于德行与恶行的寓言。毕达哥拉斯也不避讳禁止在公共大道上行走，命令不要追随多数人的粗陋且矛盾的见解。阿里斯托克利托斯在其著作\WorkTitle{驳赫拉克利奥多罗斯}的第一卷中，提到一封信如此说：\Quote{阿泰阿斯，司基提亚人的王，致拜占庭人：不要损害我的收入，免得我的母马饮你们的水；}因为这位蛮族人象征性地表示他要向他们开战。同样，诗人欧福里翁引入涅斯托尔说：\par

\Quote{我们尚未在西摩伊斯河中浸湿阿开亚人的战马。}\par

因此，埃及人在他们的神庙前放置斯芬克斯像，以表明关于天主的道理是隐晦难解的；也许也表明我们应当既爱又敬畏神圣者：爱祂，因祂对虔敬者温柔仁慈；敬畏祂，因祂对不敬者不偏不倚、公正；因为斯芬克斯同时呈现了野兽和人的形象。\par

% structure: p0070 source=h2 target=subsection reason=relative-heading-rank; base=h2

\subsection{第六章 会幕及其器皿的奥秘含义}

逐一列举所有先知和法律，详述以隐语说出的话，将是冗长的；因为几乎全部圣经都以这种方式发言。我想，对于任何有智慧的人来说，为证明当前的观点，选取几个例子便已足够。\par

隐藏表现在围绕圣殿的七重环行中，这在希伯来人那里有提及；还有长袍上的饰物，以各种象征指示那从天上延伸到地上的和谐，这些象征关涉可见之物。幔子和帷幔用蓝色、紫色、朱红色和细麻织成。如此便表明，元素的本性包含了天主的启示。因为紫色来自水，细麻来自地；蓝色是暗的，如同空气；朱红色如同火。\par

在遮盖物和帐幔的中间，就是司铎可以进入的地方，设有香坛，它是置于这宇宙中央的大地的象征；从那里升起馨烟。他们说，那介于内帐幔（只有大司祭在规定的日子才可以进入）与围绕它的外院（所有希伯来人皆可自由进出）之间的地方，是天与地的中点。但其他人说，它是理智世界和感觉世界的象征。于是，遮盖物——大众不信的屏障——在五根柱子前拉开，阻挡着周围空间的人。\par

如此奥秘地，救主擘开五饼，喂饱了听众的人群。因为人群是庞大的，他们固守于感觉之物，仿佛那是唯一存在的东西。柏拉图说：\Quote{请环顾四周，看看}，\Quote{不要让任何未入道者听见}。这样的人认为，除了他们能用双手紧握的东西之外，别无存在；却不承认行动、生成过程以及整个不可见者属于存在的领域。因为这样的人固守于五种感官。但天主的认识是这类人的耳朵及类似器官所无法触及的。因此，圣子被称为圣父的面容，祂通过披上血肉，向五种感官启示圣父的特性。\Quote{但如果我们生活在圣神内，就让我们也随从圣神而行。}\BibleRef{迦 5:25}\Quote{因为我们行走凭信德，并非凭眼目，}\BibleRef{格后 5:7}尊贵的宗徒说。于是，帐幔内隐藏着司祭的职务，并使从事它的人远离外面的人。\par

再者，进入至圣所的入口也有帐幔。有四根柱子，这是古老盟约中神圣四元组的标志。此外，那个只有能进入至圣所的人才知道的四字母奥秘之名，称为\OriginalTerm{雅威}{Jave}，解释为\Quote{是者且将存在者}。在希腊人中，天主的名字也包含四个字母。\par

如今，主独自进入理智世界，通过苦难进入对不可言喻者的认识，升越一切凭声音可知之名。灯台也被放置在香坛的南边；通过它，显示出七星的运行，它们向南旋转。因为灯台两边伸出三根枝子，上面有灯；同样，太阳犹如灯台，置于所有行星的中央，以一种神圣的乐音向上下各方分施光明。\par

金灯台传达另一个象征基督的奥迹，不仅在于形状，更在于祂光照那些相信祂并怀有望德的人，他们藉着首生者的职务而看见，这光照是\Quote{多次并以多种方式，}\BibleRef{希 1:1}。他们说，主的七目\Quote{是那停留在从叶瑟的根发出的枝子上的七神。}\par

在香坛的北边放置了一张桌子，上面有\Quote{饼的陈列}；因为最滋养的风是北风。因此，这象征了教会中的某些座位，它们共同组成一个身体和一个集会。\par

关于圣约柜所记载的事情，象征了思想世界的特性，它被隐藏并对众人关闭。\par

那些金像，每个有六翼，或象征双熊（如有些人所愿），或更确切地说，象征两个半球。名称为\OriginalTerm{革鲁宾}{cherubim}，意为\Quote{丰富的知识}。但两者共有十二翼，通过黄道带和时间（黄道带在其上运行）指出感觉世界。我想，正是关于它们，论及自然的悲剧说道：—\par

\Quote{不倦的时光在永流中圆满旋转，\newline{}自身产生自身。而双熊\newline{}凭借其翅膀的迅速游走运动，\newline{}守护着\Person{阿特拉斯}的天极。}\par

而\Person{阿特拉斯}，那不受苦的天极，可能意指固定的天球，或者更恰当地说，不动的永恒。但我认为更好的方式是考虑约柜，它从希伯来语词\OriginalTerm{特波塔}{Thebotha}得名，象征其他事物。它被解释为\Quote{各处之一代一}。那么，究竟它是指第八区域和思想世界，还是指包罗万有、无形且不可见的天主，我们暂时可以不作定论。但它象征了与朝拜神灵（即革鲁宾所意指的）同在的安息。\par

因为那位禁止制造雕像者，祂自己绝不可能按圣物的样式制造任何图像。天上也根本没有任何复合之物或具有感觉的受造之物。但面容是理性灵魂的象征，翅膀是左右权势的高贵服务者和能力；而声音是在不断默观中的荣耀之悦。让这奥秘的解释进行到这里就够了。\par

大司祭的长袍是感觉世界的象征。七颗行星由五块宝石和两颗红宝石代表，对应土星和月亮。前者属南方、湿润、属土、沉重；后者属大气，因而被一些人称为\Person{阿尔忒弥斯}，仿佛\OriginalTerm{艾若托摩斯}{Ærotomos}（切割空气）；空气是多云的。由于它们协作生产下界的事物，那些按天主眷顾掌管行星者，被正确地置于胸部和肩膀上；藉着它们，是创造之工，即第一周。而胸膛是心和灵魂的所在。\par

另一种解释是，这些石头可能是救恩的不同阶段；有些占据整个得救身体的上部，有些则占据下部。悬挂在礼服上的三百六十个铃铛，代表一年的时间，\Quote{主悦纳之年}，宣扬并响彻救主的奇妙显现。再者，宽大的金冠象征主的王权，因为\Quote{教会之首}是救主\BibleRef{弗 5:23}。头上的冠冕是最高的王权标记；此外我们也曾听说：\Quote{基督之首是天主和我們主耶穌基督的父}\BibleRef{格前 11:3}。此外，还有胸牌，包括厄弗得（工作的象征）和\OriginalTerm{神谕}{\GreekText{λογίον}}；这神谕指示出\OriginalTerm{圣言}{\GreekText{λόγος}}，万物由圣言构造而成，且是天堂的象征，由圣言所造，并且服从于基督——万有之首，因为它以同样方式并类似地运动。因此，厄弗得上发光的绿宝石，象征太阳和月亮——自然的辅助者。我认为，肩膀是手臂的起点。\par

胸前的十二块宝石，排成四行，为我们描述了黄道带的圆圈，对应一年四季的变化。法律和先知本应放在主的头下，因为两约中都提到了义人。因为若我们说宗徒们同时是先知和义人，我们说得很好，\Quote{因为同一圣神在所有人内工作}\BibleRef{格前 12:11}。正如主超出整个世界，甚至超出思想世界，同样，刻在板上的名字被认为象征超乎一切统权和大能；它的铭文涉及成文的法律和感官的显现。所表达的正是天主的名；因为，正如圣子看见圣父的良善，救主天主运行，被称为万物的首要原理，它从不可见的天主首先被形塑，在万世之前，并塑造了其后出现的一切事物。不仅如此，神谕展示了由圣言呼喊和宣讲的预言，以及将要来的审判；因为正是同一圣言预言、审判并辨别一切事物。\par

他们又说，礼服预示了在肉身内的服务，藉此祂被更密切地看见与世界的关系。于是，大司铎脱下他祝圣的礼服（宇宙以及宇宙中的受造物，因祂同意所造的都很好而被祝圣），洗自己，穿上另一件长袍——可以这么说，是一种至圣所的长袍——它将伴他进入至圣所；在我看来，这展示了肋未人和诺斯替者，作为其他司铎（那些在水中沐浴、仅以信德为衣、并期待自己个别居所的人）的首领，他自己区分理智的对象与感官的事物，超越其他司铎，急忙走向理念世界的入口，以洗去自己从下界沾染的事物，不是用水，如从前加入肋未支派时被洁净那样。但是，已经藉着诺斯替的圣言在内心被净化，并彻底规整，并将从司铎领受的生活方式提升到最高点，在言语和生活上都完全圣化，穿上荣耀的明亮礼服，接受了那属神的、完美之人的不可言喻的遗产，\Quote{眼所未见，耳所未闻，人心未曾想到的}；并成为儿子和朋友，他如今在面对面中，以无法满足的默观而满溢。因为没有什么比亲耳听见圣言本身更好的了，祂藉着圣经启发更完备的理智。因为经上这样说：\Quote{他应脱下他进入圣所时所穿的亚麻长袍；把它放在那里，在圣所内用水洗自己的身体，然后穿上自己的礼服}\BibleRef{肋 16:23}。但在我看来，主以某种方式脱下和穿上，即降至感官领域；另一种方式，透过祂而信的人脱下和穿上，如宗徒所暗示的，祝圣的祭披。从此，按主的肖像，从神圣支派中选出了最配得的人作大司铎，而那些被选为君王和先知职位的人被傅了油。\par

% structure: p0088 source=h2 target=subsection reason=relative-heading-rank; base=h2

\subsection{第七章 埃及象征与神圣事物的谜语}

因此，埃及人也没有将他们的奥秘交付给每一个人，也没有将神圣事物的知识泄漏给世俗之人；只传给那些注定要登上王位的人，以及那些从教养、文化和出身而被认为最配得上的司铎。因此，在隐藏方面，埃及人的谜语与希伯来人的谜语相似。埃及人中，有些将太阳画在船上，有些则画在鳄鱼上。他们由此表示，太阳穿越优美湿润的空气，产生时间；这时间在其他一些祭司的记述中以鳄鱼为象征。此外，在埃及的\OriginalTerm{底奥斯波里斯}{Diospolis}，在名为\OriginalTerm{皮隆}{Pylon}的神庙上，画有一个男孩作为生产的象征，一个老人作为衰朽的象征。而鹰是天主的象征，鱼是仇恨的象征；按照另一种象征，鳄鱼是厚颜无耻的象征。那么，整体的象征结合起来，似乎教导说：\Quote{哦，你们这些生而会死的人，天主仇恨厚颜无耻。}\par

另有一些人用贵重材料制作耳朵和眼睛，并祝圣它们，将它们供奉在神庙中给诸神——这清楚地表明天主看见并听见一切。此外，狮子在他们那里是力量和勇力的象征，如同牛明显是大地本身、农耕、食物，马则是坚毅和信心的象征；而斯芬克斯则是力量与智慧的结合——它完全拥有狮子的身体和人的面容。类似地，为了表示智慧、记忆、能力和技艺，人在神庙中被雕刻成雕像。在他们所谓的诸神Komasiæ中，他们抬着金像——两只狗、一只鹰和一只朱鹭；这四尊雕像被称为四个字母。狗是两个半球的象征，它们仿佛环绕并守护；鹰是太阳的象征，因它炽热且具有破坏性（所以他们将瘟疫疾病归因于太阳）；朱鹭是月亮的象征，将阴暗的部分比作深色羽毛，将明亮的部分比作光。有些人认为狗指的是回归线，它们守护并监视太阳向南和向北的路径。鹰代表赤道线，它高而炎热干燥，朱鹭代表黄道线。因为朱鹭似乎比其他动物更早向埃及人提供了发明数字和度量的最初原理，正如斜线提供了圆的原理一样。\par

% structure: p0091 source=h2 target=subsection reason=relative-heading-rank; base=h2

\subsection{第八章. 诗人与哲学家使用象征风格}

然而，不仅是埃及人中具有最高智慧的人，而且其他从事哲学的野蛮人也偏爱象征风格。据说，斯基泰人的国王伊丹图里斯，正如叙罗斯的斐瑞库德斯所记载的那样，当大流士渡过伊斯特河进行战争威胁时，他送去了一件象征物代替信件，包括一只老鼠、一只青蛙、一只鸟、一支标枪、一把犁。对这些象征物有疑问——正如可以预料的那样，千夫长奥龙托帕加斯说，他们必须放弃王国：老鼠意味着住所，青蛙意味着水，鸟意味着空气，犁意味着土地，标枪意味着武器。但克西弗罗德斯解释相反；他说：\Quote{如果我们不像鸟一样飞走，或像老鼠一样钻入地下，或像青蛙一样潜入水底，我们就逃不过他们的箭；因为我们不是这领土的主人。}\par

据说斯基泰人阿纳卡西斯在睡觉时用左手遮盖私处，用右手遮盖嘴巴，以表明两者都应当被克制，但克制舌头比克制情欲更为重要。\par

我何必在野蛮人身上多费唇舌，既然我可以举出希腊人极其热衷于使用隐藏方法呢？毕达哥拉斯派学者安德罗库德斯说，著名的所谓以弗所字母属于象征一类。因为他说\OriginalTerm{无影}{\GreekText{ἄσκιον}}意味着黑暗，因为它没有影子；\OriginalTerm{有影}{\GreekText{κατάσκιον}}意味着光，因为它用光线投下影子；\OriginalTerm{地}{\GreekText{λίξ}}是大地，根据古代的名称；\OriginalTerm{四}{\GreekText{τετράς}}是一年，与四季相关；\OriginalTerm{征服者}{\GreekText{δαμναμενεύς}}是太阳，它征服一切（\OriginalTerm{征服}{\GreekText{δαμάζων}}）；而\OriginalTerm{吉祥}{\GreekText{τὰ} \GreekText{αἴσια}}是真实的声音。于是象征表明：神圣的事物已经和谐地安排——黑暗对光明，太阳对一年，大地对自然万物的生产过程。此外，语法家狄奥尼修斯·特拉克斯在其著作\WorkTitle{论圆圈中象征意义的阐释}中明确说：有些人不仅通过言语，而且通过符号来传达行为：通过言语，如所谓的德尔斐箴言\Quote{凡事勿过度}、\Quote{认识你自己}等；通过符号，如诸神神庙中旋转的轮子（源于埃及人），以及给崇拜者的树枝。因为色雷斯人俄耳甫斯说：——\par

\Quote{凡世间人们看顾的树枝之功业，\newline{}无一在心目中有一命运，而万事万物\newline{}绕转；不可停留于一点，\newline{}但各保持其初始的一份赛程。}\par

树枝或作为最初食物的象征，或者是为了让大众知道果实普遍生长和萌发，存留很长时间；但分配给他们的生命期限是短暂的。正因如此，他们认为树枝被赐予；或许也是让他们知道，正如这些树枝另一方面被焚烧，他们自己也迅速离开此生，将成为火的燃料。\par

因此，象征性解释的方式对许多目的非常有用；它有助于正确的神学和虔诚，有助于显示智慧、实践简洁和展示智慧。\Quote{使用象征性言语是智者的特征，}语法家狄迪莫斯恰当地评论道，\Quote{并对其中所指示的进行解释。}事实上，儿童最初步的教育就包含对四种元素的解释；据说弗里吉亚人称水为Bedu，正如俄耳甫斯所说：——\par

\Quote{明澈的水倾泻而下，是宁芙们的Bedu。}\par

迪翁·泰特斯似乎也写道：——\par

拿取Bedu，浇在你手上，然后转向占卜。\par

另一方面，喜剧诗人菲利狄斯理解Bedu为空气，作为生命赐予者（Biodoros），在以下诗句中：——\par

\Quote{我祈祷我能吸入有益的Bedu，\newline{}它是健康最重要的部分；\newline{}吸入纯净、无污的空气。}\par

库齐库斯的尼安忒斯同意同一观点，他写道马其顿的祭司们祈求Bedu，他们解释为\Emph{空气}，使之对他们和他们的子女仁慈。有些人无知地认为Zaps是火（来自\OriginalTerm{沸腾}{\GreekText{ζέσιν}}）；因为海如此称呼，如欧福里翁在给塞奥里达斯的答复中：——\par

\Quote{而Zaps，船舶的毁灭者，把它撞碎在岩石上。}\par

狄奥尼修斯·伊安布斯类似地说：——\par

\Quote{咸涩的Zaps在狂怒的深海上呜咽。}\par

同样地，年轻喜剧家克拉提努斯：——\par

\Quote{札普斯吐出虾和小鱼。}\par

罗得岛的西米亚斯：—\par

\Quote{伊格内特斯和忒尔喀涅斯的咸涩之父札普斯诞生了。}\par

而\OriginalTerm{地}{\GreekText{χθ}\GreekText{ων}}是散布至广大的大地（\OriginalTerm{\GreekText{κεχυμ}\GreekText{ενη}}{}）。普莱克特隆，据一些人说，是天极（\OriginalTerm{\GreekText{π}\GreekText{ολος}}{}），据另一些人说，是空气，它击打（\OriginalTerm{\GreekText{πλ}\GreekText{ησσοντα}}{}）并促动自然与生长，且充满万有。但这些人没有读过哲学家克莱安特斯，他明确称普莱克特隆为太阳；因为它在东方投出光线，犹如击打世界，引导光行其和谐之路。而太阳也意指其余星辰。\par

而斯芬克斯，依照诗人阿拉托斯，并非对宇宙的理解和世界的旋转；而或许是那渗透并维系宇宙的精神音调。但更宜视之为以太，它维系并压迫一切；正如恩培多克勒所说：—\par

\Quote{但来吧，首先我将述说太阳，万物的第一本原，\newline{} 我们所见的一切都从中生出，\newline{} 大地、波涛深渊与湿润空气；\newline{} 提坦与以太，它束缚万物于周遭。}\par

科孚岛的阿波罗多洛斯说，这些诗句是先知布朗库斯在为米利都人净化瘟疫时所吟诵的；因为他用月桂树枝洒向众人，领唱赞歌，大致如下：—\par

\Quote{歌唱吧，少年们，赫卡厄尔戈斯和赫卡厄尔加。}\par

民众应和他，说道：\Quote{贝杜，札普斯，地，普莱克特隆，斯芬克斯，克纳克兹比，克堤普特斯，弗莱格摩斯，德罗普斯。}卡利马科斯在抑扬格诗中记载了此事。克纳克兹比，按词源，是瘟疫，因其\OriginalTerm{啃噬}{\GreekText{κνα}\GreekText{ιειν}}和\OriginalTerm{毁灭}{\GreekText{διαφθε}\GreekText{ιρειν}}，而\OriginalTerm{霹雳}{\GreekText{θ}\GreekText{υψαι}}是用雷霆吞噬。悲剧诗人特斯庇斯认为这些词所指另有他物，他这样写道：\Quote{看，我向你献上白色的克纳克兹比奠酒，从黄色的乳母中压取。看，啊，双角潘，将克堤普特斯奶酪与红蜜调和，我置于你神圣的祭台上。看，我为你倾洒布罗米乌斯闪亮的光芒。}依我之见，他意指灵魂最初如奶般的滋养，来自二十四元素，随后以凝固之乳为食物。最后，他教导圣言之葡萄的血，那闪亮的美酒，即教导之圆满喜乐。而德罗普斯是运作的圣言，始于基础训练，进至人的成长，燃烧并光照人直到成熟的尺度。\par

第三段据说是儿童的习字帖——\OriginalTerm{马尔普特斯}{\GreekText{μ}\GreekText{αρπτες}}、\OriginalTerm{斯芬克斯}{\GreekText{σφ}\GreekText{ιγξ}}、\OriginalTerm{克洛普斯}{\GreekText{κλ}\GreekText{ωψ}}、\OriginalTerm{津克特顿}{\GreekText{ζυνχθηδ}\GreekText{ον}}。依我之见，它意指通过元素和世界的安排，我们必须进至更完善的知识，因为永恒救恩靠力量和辛劳获得；因为\OriginalTerm{马尔普特斯}{\GreekText{μ}\GreekText{αρψαι}}是抓取。而世界和谐由斯芬克斯意指；\OriginalTerm{津克特顿}{\GreekText{ζυνχθηδ}\GreekText{ον}}意指困难；而\OriginalTerm{克洛普斯}{\GreekText{κλ}\GreekText{ωψ}}同时意指主和日的隐秘知识。噫！厄庇革涅斯在其论\WorkTitle{奥菲斯的诗}一书中，展示奥菲斯诗中的特征时，难道不是说\Quote{曲杖}（\OriginalTerm{\GreekText{κερα}\GreekText{ισι}}{}）意指\Quote{犁}；经线（\OriginalTerm{\GreekText{στ}\GreekText{ημοσι}}{}）意指犁沟；纬线（\OriginalTerm{\GreekText{μ}\GreekText{ιτος}}{}）是种子的比喻；宙斯之泪意指阵雨；而\Quote{部分}（\OriginalTerm{\GreekText{μο}\GreekText{ιραι}}{}）又是月亮的相位，即三十日、十五日和新月，奥菲斯因此称它们为\Quote{白衣者}，作为光的部分？同样，春天因本性而被称为\Quote{繁花}；夜因休息而被称为\Quote{寂静}；月亮因其中的面孔而被称为\Quote{戈尔贡}；而\Quote{神学家}称必须播种的时节为阿佛洛狄忒。同样，毕达哥拉斯派也比喻地称行星为\Quote{珀耳塞福涅的狗}；并将海的隐喻名称应用于\Quote{克洛诺斯之泪}。在诗人和哲学家中，可发现无数谜语般的表述；也有整卷书隐晦地呈现作者的思想，如赫拉克利特的\WorkTitle{论自然}，他正是因此被称为\Quote{晦涩者}。类似此书的是叙罗斯的斐瑞居德斯的\WorkTitle{神学}；因为诗人幼佛里翁、卡利马科斯的\WorkTitle{原因}、吕科弗隆的\WorkTitle{亚历山德拉}等，被提给所有文法家作为解释的练习。\par

因此，我们所要论述的蛮族哲学，也应当像所显示的那样，以隐晦和象征的方式预言。摩西的诫命即是如此：\Quote{这些不洁之物，即猪、隼、鹰和乌鸦，不可吃。}因为猪是肉欲和不洁的饮食贪欲，以及淫荡污秽的纵欲的象征，总是发痒、属物质、躺在泥中，为宰杀和毁灭而被养肥。\par

他又命令吃那些分蹄和反刍的动物；\Quote{示意}，\Person{巴尔纳伯}说，\Quote{我们应当依附那些敬畏主的人，他们在心中默想所领受的那部分圣言，依附那些宣说并遵守主法令的人，依附那些以默想为乐事的人，他们反刍主的圣言。分蹄又是什么？那就是义人在此世行走，并期待圣洁的永恒来临。}然后他补充说，\Quote{看\Person{梅瑟}立法多么好！但他们从何处能明白或领悟这些事呢？我们正确领悟的人，按主所愿的方式宣示诫命；因此祂割损了我们的耳朵和心，使我们能领会这些事。当他引用\InnerQuote{你不可吃鹰、隼、鸢和乌鸦}时，他说：\InnerQuote{你不可依附或效法那些不知以劳作和汗水谋生，却以掠夺和非法方式生活的人。}因为鹰象征掠夺，隼象征不义，乌鸦象征贪婪。经上也记着：\InnerQuote{与纯洁的人，你将成为纯洁；与正直的人，你将成为正直；与乖僻的人，你将成为乖僻。}我们有责任依附圣人，因为依附他们的人将被圣化。}\par

因此\Person{特奥格尼斯}写道：——\par

\Quote{因为从善者你将学到善事；\newline{}但你若与恶者混杂，就会毁掉你或许拥有的任何心智。}\par

再者，当诗歌中说，\Quote{因为祂大获全胜：将马和骑马者投进了海}，\BibleRef{出 15:1}。那多肢体的野蛮情欲、贪欲，连同骑手，他放松缰绳纵情享乐，\Quote{祂投进了海}，将他们抛入世界的混乱之中。同样，\Person{柏拉图}在他的著作\WorkTitle{论灵魂}中说，驾战车者和跑开的马——即分为愤怒和贪欲两部分的非理性部分——坠落；因此这个神话暗示，\Person{法厄同}是因马匹放荡而被抛出的。同样，关于\Person{若瑟}：兄弟们嫉妒这位青年，他因知识而拥有非凡的远见，便剥下他的彩衣，把他扔进一个坑里（坑是空的，没有水），拒绝这位善人源于爱慕教诲的多样知识；或者，他们凭藉根据法律的朴素信德，将他扔进无水之坑，卖他到没有天主圣言的\Place{埃及}。那坑缺乏知识；被扔进去并剥去知识后，那个在不知不觉中变得有智慧的人，失去了知识，看起来像他的兄弟们。另一种解释，彩衣是贪欲，它通入一个张开的坑。\Quote{如果有人挖开或凿出一个坑}，经上说，\Quote{而不盖好，若有牛或驴掉进去，坑主应赔偿金钱，交给邻人；死兽归他。} \BibleRef{出 21:33}。这里加上那预言：\Quote{牛认识主人，驴认识主人的槽；以色列却不认识我。} \BibleRef{依 1:3}。为了不让那些接触了你所教导的知识的人，变得不能持守真理，而不顺从而跌倒，据说，要在对待圣言上谨慎，将活泉封闭在深处，远离那些非理性靠近的人，却要给渴求真理的人递上饮品。那么，要向那些不配接受知识深处的人隐藏它，并这样盖住坑。那么，坑主，即\OriginalTerm{诺斯替者}{Gnostic}，自己将受惩罚，因他人绊倒而受责备，并因圣言的伟大而被压倒，他自己能力微小；或者将工人转移到思辨领域，并因此使他脱离即刻的信德。\Quote{将付钱}，呈上清算，并将他的账目提交给\OriginalTerm{全能意志}{the omnipotent Will}。\par

那么，这就是\Quote{法律和先知直到\Person{若翰}为止}的预象，\BibleRef{玛 11:13}。而他虽然说话更加清晰，不再预言，而是指出那从起初就被象征性地宣告的祂现在已临在，却仍然说：\Quote{我是不配解主的鞋带的。}因为他承认自己不配给如此伟大的权能施行洗礼；因为那些净化他人的人，理应使灵魂脱离身体及其罪恶，如同脚脱离鞋带。或许这也意味着\Person{救主}对我们能力的最终运用——我是指直接的能力——即祂的临在，隐藏在预言的谜语中，既然他（\Person{若翰}）通过指出那曾被预言的那一位，并指示已经来到的临在，步入光明，解开了\OriginalTerm{旧约经济学}{the [old] economy}学问的鞋带，揭示了象征的含义。\par

\Person{罗马人}在遗嘱中实行的惯例在此适用；那些天平和硬币象征公义，解放奴隶，以及搓耳朵。因为这些惯例是为了公义地处理事务；那些是为了分配荣誉；最后一项是，偶然在旁的人，如同被强加了一项负担，应站立倾听并充当中介。\par

% structure: p0125 source=h2 target=subsection reason=relative-heading-rank; base=h2

\subsection{第九章 用象征掩盖真理的原因}

但看来，我在急于证明自己的论点时，不知不觉超出了必要范围。因为若要引证那些以象征方式从事哲学思考的大量例子，我的生命将不够用。因此，为了纪念、简洁和吸引人归向真理，\OriginalTerm{异邦哲学的圣经}{the Scriptures of the Barbarian philosophy}就是这样。\par

因为只有那些经常接近它们，并借着信德和整个生活试炼它们的人，才能得到真正的哲学和真正的神学。他们也希望我们要求一位解释者和向导。因为他们如此认为，从深知真理的人手中接受真理，我们会更加认真，不易受骗，而那些相称的人将获益。此外，一切透过面纱照耀的事物，都使真理显得更伟大、更庄严；就像透过水闪耀的果实，以及透过面纱的形体，这些都给予它们附加的反映。因为，除了未被隐藏的事物以一种方式被感知外，环绕的光线也揭示了缺陷。既然如此，我们可以从隐藏的表达中得出多种含义，正如我们所做的，无知和未受教育的人就会失败。但\OriginalTerm{诺斯替者}{Gnostior}却能领会。现在，并不希望所有事物都无差别地暴露给所有人，也不希望将智慧的益处传授给那些甚至在梦中都未净化灵魂的人（因为不允许将经过如此艰苦努力获得的东西随意传给每一个偶然前来的人）；也不应将圣言的奥秘解释给亵渎者。\par

他们说，毕达哥拉斯学派的希帕尔科斯因用明白的语言写下毕达哥拉斯的教义，被逐出学派，并为他立了一根柱子，仿佛他已死了一般。因此，在蛮族哲学中，他们称那些背离了教义、使心灵屈从于肉体情欲的人为死人。\BibleQuote{不义与正义有什么交谊？}按照神圣的宗徒的话。\BibleQuote{光明与黑暗有什么相通？基督与贝里雅耳有什么和谐？信者与不信者有什么份？}\BibleRef{格后 6:14} 因为奥林匹斯神和必死者的荣誉是分开的。\BibleQuote{因此，你们要从他们中间出来，与他们分离，主说，不要触摸不洁之物；我要收纳你们，我要作你们的父亲，你们要作我的儿女。}\BibleRef{格后 6:17}\par

不仅毕达哥拉斯学派和柏拉图隐藏了许多事物，伊壁鸠鲁学派也说他们有不可言说的事物，不允许所有人阅读那些著作。斯多葛学派也说，第一位芝诺写下了一些东西，他们不轻易让门徒阅读，除非他们先证明自己是否是真正的哲学家。亚里士多德的弟子们说，他们的论著有些是祕传的，有些是普通的、公开的。此外，那些创立奥秘的人，作为哲学家，将他们的教义埋藏在神话中，以免为所有人知晓。那么，他们是否通过隐藏人类的意见来阻止无知者触碰它们呢？而对实在的神圣和神圣的沉思予以隐藏，不是更有益吗？但不仅是蛮族哲学的教义或毕达哥拉斯的神话。甚至柏拉图著作中的那些神话（在\WorkTitle{共和篇}中，关于亚美尼亚的希罗的神话；在\WorkTitle{高尔吉亚篇}中，关于埃阿科斯和拉达曼提斯的神话；在\WorkTitle{斐多篇}中，关于塔耳塔洛斯的神话；在\WorkTitle{普罗泰戈拉篇}中，关于普罗米修斯和厄庇墨透斯的神话；此外，在\WorkTitle{大西洋篇}中，关于亚特兰蒂斯人和雅典人之间的战争的神话）都必须以寓言方式解释，并非在一切表达上绝对地，而是在那些表达总体意义的表达上。我们将发现这些是以象征方式在寓言的帷幕下指出的。此外，毕达哥拉斯社团和与同伴的双重交往，其中大多数被称为\OriginalTerm{听讲者}{\GreekText{ἀκουσματικοί}}，而另一些真正热爱哲学的人被称为\OriginalTerm{学习者}{\GreekText{μαθηματικοί}}，然而也表明有些话是对大众说的，有些则向他们隐藏。也许，逍遥学派教学的双重种类——一种称为或然的，另一种称为可知的——非常接近意见与荣耀和真理之间的区分。\par

\Quote{为赢得人间美誉之花，\newline{}不要被诱导说出任何超出正义的话。}\par

因此，伊奥尼亚的缪斯明确说：\Quote{大多数自以为聪明的人跟随吟游诗人并利用法律，知道许多人是坏的，少数是好的；但最好的人追求荣耀：因为最好的人选择人间的永恒荣耀胜过一切。但大众像野兽一样填塞自己，以肚腹和私处以及我们中最卑贱的事物来衡量幸福。}而伟大的埃利亚的巴门尼德被引入如此描述两条道路的教导：——\par

\Quote{一条是令人信服的真理的无畏之心；\newline{}另一条在于人的意见，其中没有真信德。}\par

% structure: p0133 source=h2 target=subsection reason=relative-heading-rank; base=h2

\subsection{第十章 宗徒们关于隐藏信仰奥秘的意见}

因此，神圣的宗徒有道理地说：\BibleQuote{这奥秘藉启示已向我彰显（如我先前简要所写，你们照着读，便能明白我对基督奥秘的认识），这奥秘在以前的世代没有对人子显示过，如今却已启示给祂的圣宗徒和先知。}\BibleRef{弗 3:3} 因为有一种对成全者的教导，关于此，他写信给哥罗森人说：\BibleQuote{我们不断为你们祈祷，恳求你们在一切智慧和属神的领悟上，充满对祂旨意的认识，好使你们行事为人对得起主，事事蒙祂喜悦，在一切善事上结出果实，不断增加对天主的认识，照着祂光荣的德能，以一切大能坚强起来。}\BibleRef{哥 1:9} 他又说：\BibleQuote{按照天主恩宠所托付我的职分，为使你们完成天主的圣言；这奥秘从万世以来、从历代以来隐藏着，如今却显示给祂的圣徒们：天主愿意他们知道，这奥秘在外邦人中是何等光荣的丰富。}\BibleRef{哥 1:25} 因此，一方面，这些奥秘隐藏直到宗徒的时代，并由他们从主那里领受而传授，且隐藏在旧约中，显明给圣徒们。另一方面，有\BibleQuote{这奥秘在外邦人中光荣的丰富}，就是对基督的信德和望德；他在另一处称之为\Quote{根基}。\BibleRef{哥 1:27} 再者，他仿佛急于揭示这知识，如此写道：\Quote{用一切智慧警戒每个人，好使我们将每个人（整个的人）在基督内成全地呈献。}不是泛指每个人，因为没有人会是不信的。他也并非称每个信基督的人为成全者；而是说整个的人，仿佛是说全人，如同在身体和灵魂上得到净化。因为知识并不属于所有人，他明确补充说：\BibleQuote{在爱中彼此联结，充满确信的知识的一切丰富，以认识天主在基督内的奥秘，因为在祂内蕴藏着一切智慧和知识的宝藏。}\BibleRef{哥 2:2} \BibleQuote{你们要恒心祈祷，警醒着，怀着感恩的心。}\BibleRef{哥 4:2} 感恩不仅为灵魂和属神的祝福，也为身体和身体上的美物。他更清楚地揭示知识不属于所有人，补充说：\BibleQuote{同时也为你们祈祷，求天主给我们打开传讲基督奥秘的门，我就是为此被捆绑，好使我按所当说的将这奥秘显明出来。}\BibleRef{哥 4:3} 因为在希伯来人当中，确实有些传统是口传的。\Quote{按时间说，你们本该做导师了，}这是说，他们仿佛已在旧约中老去，\Quote{却还需要有人再将天主圣言的基础教导你们；你们成了需要奶，而不是坚硬食物的人。因为凡吃奶的，对正义的话还没有经验，因为他是婴孩，尚在学习初阶课程。但坚硬的食物属于成年人，他们因运用而锻炼了感官，能够分辨善恶。因此，我们应离开基督道理的基础，向成全者迈进。}\par

\Person{巴拿巴}也曾亲自与宗徒在外邦人的职务中宣讲圣言，他说：\Quote{我非常简单地给你们写信，为使你们明白。}然后下面，已展现出更清楚的诺斯替传统痕迹，他说：\Quote{另一位先知\Person{梅瑟}对他们说什么？看，上主天主这样说：进入那美好之地，就是上主天主、\Person{亚巴郎}、\Person{依撒格}和\Person{雅各伯}的天主所誓许的；你们承受了那流奶流蜜的地为产业。}知识说什么？学习，盼望，它说，在\Person{耶稣}内，祂将要在肉身中显示给你们。因为人是受苦之地；因为\Person{亚当}是由地土而形成的。那么，关于流奶流蜜的良好土地，它说了什么？弟兄们，愿我们的主受赞美，祂将智慧和对祂奥秘的理解放在我们心中。因为先知说：\Quote{谁明白上主的比喻呢？唯有智慧且明达的人，以及爱慕他主人的人。}只有少数人能领会这些事。因为主并非出于嫉妒而宣布福音说：\Quote{我的奥秘属于我和我家中的儿子们。}这样使选民安全无虞，超越焦虑；使属于所选择和所领受的事物超越嫉妒的范围。因为不拥有善的知识的人是邪恶的：因为只有一位是善的，就是父；不认识父就是死亡，认识祂就是永生，通过参与那不朽者的德能。而不朽就是参与神性；但背离对天主的认识就带来败坏。先知又说：\BibleQuote{我将赐给你隐藏、幽暗、看不见的宝藏，使他们知道我是上主。}\BibleRef{依 45:3} 同样，\Person{达味}歌唱说：\Quote{因为，看，你喜爱真理；你将智慧的隐晦和隐藏的事显明给我。}\Quote{日复一日吐露言语}（指清楚写出的），\Quote{夜复一夜传达知识}（指隐藏于奥秘面纱中的）；\Quote{没有话语，没有言词，它们的声音不会被天主听到，}祂说：\Quote{人做隐秘的事，我岂能看不见他？}\par

因此，揭示隐藏之事的教导被称为\OriginalTerm{光照}{illumination}，因为唯有教师才打开约柜的盖子，这与诗人们所说的相反：\Quote{宙斯堵住了好运之罐，却打开了恶运之罐。} \Quote{因为我知道，}宗徒说，\Quote{当我到你们那里时，我必带着基督圆满的祝福而来；}\BibleRef{罗 15:29} 指明属灵的恩赐和\OriginalTerm{真知者}{gnostic}的交流，他既在场，便渴望将之传达给在场的他们，作为\Quote{基督的圆满，照着从永世以来所隐藏的\OriginalTerm{奥迹}{mystery}的启示，如今却借着先知的圣经，按照永生天主的命令，显明给万民，使他们因信德而服从，}即那些相信如此的外邦人。但只有少数人才被显明那\OriginalTerm{奥迹}{mystery}中所包含的事。\par

因此，柏拉图在\WorkTitle{书信}中论及天主时说得对：\Quote{我们必须用\OriginalTerm{谜语}{enigmas}说话；以免书板无论在海路或陆路上偶然失落，阅读的人仍一无所知。}因为宇宙的天主超越一切言语、一切概念、一切思想，绝不能付诸笔墨，即使以祂自己的能力也无法言说。柏拉图也表明了这一点，他说：\Quote{既然如此，当心免得有一天你因现今之事而懊悔，以不配的方式离去。最大的保障不是写下来，而是学习；因为写下来的东西不可能不消失。}\par

与此相似的是圣宗徒保禄所说的，他保存了那预言性的、真正古老的秘密，希腊人从中汲取了美好的教导：\Quote{然而在成全的人中，我们讲论智慧，却不是今世的智慧，也不是今世掌权者——即将消灭者——的智慧；我们讲论的是天主隐藏于\OriginalTerm{奥迹}{mystery}中的智慧。}\BibleRef{格前 2:6} 接着，他以下面的话劝诫不可将他的话泄露给大众：\Quote{弟兄们，我从前对你们说话，不能把你们当作属神的人，只能当作属血肉的人，当作在基督内的婴孩。我给你们\OriginalTerm{奶}{milk}喝，没有给\OriginalTerm{肉}{meat}吃，因为你们那时还不能吃，就是如今还是不能，因为你们仍然是属血肉的。}\BibleRef{格前 3:1}\par

如果\Quote{\OriginalTerm{奶}{milk}}被宗徒说成属于婴孩，而\Quote{\OriginalTerm{肉}{meat}}是成年人的食物，那么奶将被理解为教理讲授——仿佛是灵魂的第一食物。而肉是\OriginalTerm{奥秘默观}{mystic contemplation}；因为这是圣言的肉和血，即对天主性体德能与本质的理解。经上说：\Quote{你们尝尝，并看看主是基督。}因为祂就这样将祂自己赐予那些以更属灵的方式分享这种食物的人；这时灵魂按照爱真理的柏拉图所说，滋养自己。因为对天主本质的认识是圣言的肉和饮料。因此柏拉图在\WorkTitle{理想国}第二卷中也说：\Quote{那些不是献上一头母猪，而是某种伟大而艰难的祭品的人}才应当探究天主。宗徒写道：\Quote{我们的逾越节基督已被祭献，为我们做了牺牲；}\BibleRef{格前 5:7}——这确实是一种难以获得的祭品，天主子为我们而奉献。\par

% structure: p0140 source=h2 target=subsection reason=relative-heading-rank; base=h2

\subsection{第十一章　脱离物质事物为获得对天主的真知识所必需}

如今，天主所悦纳的祭献是从身体及其情欲中毫不退缩的\OriginalTerm{抽象}{abstraction}。这才是真正的虔诚。因此，苏格拉底不正是正确地将哲学称为\OriginalTerm{死亡操练}{practice of Death}吗？因为那在思辨中不使用眼睛，也不从其他感官中获取任何东西，而是以纯洁的心灵本身投向对象的人，才实践真正的哲学。这就是毕达哥拉斯所规定的\OriginalTerm{五年静默}{silence of five years}的意义，他命令门徒这样做；好使他们从感官对象中抽象出来，仅以心灵默观神性。希腊人的领袖正是从梅瑟那里汲取了这些哲学信条。因为他命令将\OriginalTerm{全燔祭}{holocausts}剥皮并分成块。因为\OriginalTerm{真知者的灵魂}{gnostic soul}必须被祝圣于光明，剥去\OriginalTerm{物质的覆盖}{integuments of matter}，脱离身体的轻浮和一切情欲——这些是通过虚妄和欺骗的意见获得的——并脱去肉体的贪欲。但大多数人，披着易朽之物，如同蛤蜊，在过度的欲望中缩成一团，如同刺猬，对那有福且不朽的天主怀有与自身相同的观念。他们没有注意到，尽管天主临近我们，却将万般事物赐予我们，而祂自己并不分享这些：出生，祂自己非受生；食物，祂一无所缺；成长，祂永远不变；长寿和永生，祂不朽且永不衰老。因此，不要让任何人认为手、脚、口、眼、来去、忿怒和威胁，是希伯来人所言的天主的属性。绝非如此；而是这些称谓中的某些在\OriginalTerm{寓意义}{allegorical sense}中被更神圣地使用，随着论述的推进，我们将在适当的时候解释。\par

\Quote{万药之智是万能药}，卡利马科斯在\WorkTitle{铭辞}中写道。\Quote{一人从另一人获得智慧，无论过去还是现在}，巴库里德斯在\WorkTitle{颂歌}中说；\Quote{因为要找到不可言说之言的门扉并非易事。}因此，伊索克拉底在\WorkTitle{泛雅典娜颂}中优美地写道，他提出疑问：\Quote{那么，谁是受过良好训练的人？}接着补充：\Quote{第一，那些善于处理日常事务的人，其见解顺应时机，且多半能击中要害；第二，那些以合宜正直的态度对待接近他们的人，对他人的烦扰和侵扰轻描淡写、从容应对，而与交往之人尽可能以极度审慎和节制行事；此外，那些能驾驭自己的享乐，不为不幸所过度压倒，而是在其中以配得我们共有本性的男子气概行事的人；第四——也是最重要的——那些不被顺境腐蚀、不迷失自我或变得高傲，而是继续留在明智人群之列的人。}然后他给论述加上顶石：\Quote{那些灵魂的倾向不仅适合其中一件事，而是适合所有事的人——我断言他们是智慧和完美的人，拥有所有德行。}\par

你看到希腊人如何神化\OriginalTerm{知识的生活}{gnostic life}（却不知道如何认识它）了吗？而知识是什么，他们连做梦也不知道。那么，如果在我们中间一致认为知识是理性的食粮，那么根据圣经，\Quote{真正有福的是那些饥渴慕真理的人：因为他们必得饱足}，以永久的食物。与这些话最奇妙地和谐一致，戏剧哲学家欧里庇得斯在以下诗句中被发现——不知何故，同时暗示圣父和圣子：——\par

\Quote{向祢，万有之主，我献上\newline{}饼和奠祭，哦，宙斯，\newline{}或者哈得斯是祢愿被称呼之名；\newline{}请接受我全果的供奉，\newline{}稀有、满溢、倾注。}\par

因为全燔祭和稀有的祭献为我们是基督。而他不知不觉地提及救主，他将明确，当他补充说：——\par

\Quote{因为祢，在天神之中，\newline{}执掌宙斯的权杖，也分享\newline{}对地上之人的统治。}\par

然后他明确地说：——\par

\Quote{请将光明送入渴望知道\newline{}冲突的源起和罪恶的根源的人类灵魂，\newline{}以及我们必须向有福之神献上\newline{}应得祭礼的仪式，以便\newline{}我们可以从烦恼中获得安息。}\par

因此，在希腊人中盛行的奥秘中，洁净占据首位；正如在野蛮人中洗濯一样。此后是小奥秘，它们有一些教导基础和预备，为接下来的事做准备；而大奥秘，在其中关于宇宙已无他事可学，只需沉思和理解本性及事物。\par

我们将通过告解理解净化方式，通过分析理解沉思方式，从基础属性开始，通过分析上升到第一概念；从身体中剥离其物理属性，除去深度的维度，然后宽度的维度，然后长度的维度。因为剩下的点是一个单位，可以说，具有位置；如果我们从中抽离位置，就剩下统一的概念。\par

那么，如果我们抽离所有属于身体和所谓无形体的事物，将自己投入基督的伟大中，然后通过圣洁进入无限，我们或许能到达全能者的概念，知道祂不是什么，而不是祂是什么。形状和运动，或站立，或宝座，或地方，或右手或左手，完全不应被认为属于宇宙之父，尽管经文如此书写。但这些各是什么意思，将在适当的地方展示。第一因不在空间内，而是超越空间、时间、名称和概念。\par

因此梅瑟也说：\BibleQuote{求祢将祢自己显示给我}\BibleRef{出 33:18}——最清楚地表明天主不能被人教导，也无法用言语表达，而只能通过祂自己的大能来认识。因为探究是模糊昏暗的；但知识之恩宠是由祂通过圣子而来的。撒罗满将最清楚地为我们作证，如此说：\BibleQuote{人的谨慎不在我内；但天主赐我智慧，我知晓神圣之事。}\BibleRef{箴 30:2}现在梅瑟以寓言描述神圣的谨慎，称之为种植在乐园中的生命树；这乐园或许是世界，其中所有出于创造的事物生长。在其中，圣言亦开花结果，\Quote{成了血肉}，并给那些\Quote{尝过祂的恩慈}的人带来生命；因为祂来对我们的认识并非没有木料之树。因为我们的生命被挂在其上，以便我们相信。撒罗满又说：\BibleQuote{她是一棵不朽的树给那些抓住她的人。}\BibleRef{箴 3:18}\Quote{看，我将生命与死亡摆在你们面前，要爱主你们的天主，行在祂的道路上，听祂的声音，信托于生命。但如果你们违犯我所赐给你们的法令和判断，你们将彻底毁灭。因为这就是生命和你们日子的长久：爱主你们的天主。}\par

再者：\BibleQuote{亚巴郎在第三天来到天主所指示的地方，举目一望，远远看见了那地方。}\BibleRef{创 22:3}因为第一天是由看见美善事物所构成的；第二天是灵魂的最佳渴望；第三天，理智感知属神的事物，理解的双眼被第三天复活的那位师傅打开。三天可能是印记的奥迹，在其中天主真正被相信。因此他远远看见那地方。因为天主的领域是难以达到的；柏拉图称之为理念的领域，从梅瑟那里学到这是一个包含一切万有普遍性的地方。但它被亚巴郎远远看见，正确地，因为他处于受造界，并立刻被天使引入。因此宗徒说：\Quote{现在我们藉着镜子观看，但那时要面对面，}藉着理智那些唯一纯洁和非物质的接触。在推理中，有可能推论关于天主，如果有人尝试不用任何感官，藉着理智达到那独特的；并且不退出存在的领域，直到上升到超越它的事物，他藉着理智本身领悟那善，在思想世界的边界处移动，按柏拉图所说。\par

再者，梅瑟不允许在许多地方建造祭台和圣殿，而只建造一个天主的圣殿，宣告世界是独生的，如巴西里德所说，并且天主是唯一的，但这在巴西里德看来似乎尚未显明。既然诸知的梅瑟不把那不能被限定者限定在空间内，他在圣殿中没有设立任何像以供敬拜；显示天主是无形无像的，无法被限定；并且以某种方式藉着对圣殿中祂名字的尊崇引导希伯来人认识天主。此外，圣言禁止建造圣殿和一切祭献，暗示全能者不被任何事物所包含，藉着祂所说的：\BibleQuote{你们要为我建筑什么殿宇？\Person{主}说。天是我的宝座，}\BibleRef{依 66:1}等等。关于祭献也同样：\BibleQuote{我不喜欢公牛的血和羔羊的脂肪，}以及圣神藉着先知在后续所禁止的。\par

因此，\Person{欧里庇得斯}与这些观点极其吻合，当他写道：——\par

\Quote{什么房屋由工匠之手建造，\newline{}用墙壁的褶皱能遮盖天主的形状？}\par

论到祭献他这样说：——\par

\Quote{因为天主一无所缺，如果祂真是天主。\newline{}这些是吟游诗人的可怜神话。}\par

\Quote{因为天主创造世界并非出于需要；而是为了从人和其他的神祇与魔鬼那里获得尊荣，从受造界收取一种贡赋，从我们这里收取烟雾，从神祇和魔鬼那里收取他们适当的供奉，}柏拉图说。因此，极其有教益地，保禄在\BibleBook{宗徒}{大事录}中说：\BibleQuote{创造宇宙及其中的万物的天主，既是天地的主宰，并不住人手所造的殿宇，也不受人手的侍奉，好像祂需要什么似的，因祂自己将生命、气息和一切赐给众人。}\BibleRef{宗 17:24}而斯多噶学派的创始人\Person{芝诺}，在他的\WorkTitle{共和国}一书中说，\Quote{我们不应该建造殿宇或偶像；因为没有一件作品是配得上众神的。}他毫不畏惧地写下这些话：\Quote{不需要建造殿宇。因为殿宇不值得珍贵，也不应被视为神圣。因为没有任何东西是值得珍贵和神圣的，如果它是建造者和工匠的作品。}因此，柏拉图正确地认识到世界是天主的圣殿，向公民指出城中一个地方，他们的偶像应该被放置在那里。他说：\Quote{那么，不要让任何人再向神祇奉献殿宇。因为在其他城邦中，金和银，在私人手中和在殿宇中，是一种引人嫉妒的财产；象牙，一种已经离开生命的物体，不是神圣的还愿供物；钢铁和黄铜是战争的工具；但无论谁想要奉献，就让他用同一棵树的木材，以及用于普通殿宇的石头。}因此，在伟大的书信中他正确地说：\Quote{因为它不像其他学科那样能够言说。但由于与这个主题极大的亲密以及与之共处，一束突然的光芒，如同由闪烁的火点燃的，在灵魂中产生，并滋养自身。}这些陈述难道不像\Person{索福尼亚}先知所说的吗？\Quote{主的神把我提去，带我上到第五层天，我看见称为主的天使；他们的冠冕安放在圣神中；每一个都有宝座，比初升太阳的光亮七倍；他们住在救恩的殿宇中，歌颂那不可言说的至高的天主。}\par

% structure: p0160 source=h2 target=subsection reason=relative-heading-rank; base=h2

\subsection{第十二章　天主无法用言语或心智把握}

\Quote{因为发现这个宇宙的父和创造者既是一桩难以完成的任务；而一旦找到祂，也不可能向所有人宣告祂。因为这绝不像其他教导科目那样能够言说，}爱好真理的柏拉图说。因为那正确听到全智的\Person{梅瑟}登上高山作神圣默观，达到理智对象的顶峰的人，必然命令全体人民不要跟随他。当圣经说\BibleQuote{梅瑟进入了天主所在的浓云黑暗}时，这向那些能理解的人显示，天主是无形无像且不可言说的。而\Quote{黑暗}——实际上是群众的不信和无知——阻碍了真理的光芒。再者，神学家\Person{俄耳甫斯}，从这方面得到帮助，说：——\par

\Quote{一是完美的在自己内，万有都是一的产物，}\par

或\Quote{被生出}；因为也是这样写的。他补充说：——\par

\Quote{祂\newline{}没有凡人见过，但祂看见一切。}\par

他更清楚地说：——\par

\Quote{我看不见祂，因为周围有一片云\newline{}笼罩；因为凡人的眼睛太小，\newline{}凡人的瞳孔——那里只长着肉和骨。}\par

宗徒将为这些陈述作证：\BibleQuote{我认得一个在基督内的人，被提到第三层天上，又从那里被带到乐园里，听到了不可言传的话，那是人不可说的。}——以此暗示不可能表达天主，并指出神圣之事非人力可述；若他确实开始在第三层天之上发言，正如在那里可以合法地将奥秘启示给蒙选的灵魂一样。因为我知道\Person{柏拉图}著作中的内容（因为来自蛮族哲学的例子很多，现在根据先前应许，等待适当时机的著作提供了这些例子）。因为在\WorkTitle{蒂迈欧篇}中，他怀疑我们是否应当把多个世界理解为多个天，还是这一个，他在名称上没有区分，用同一个名称称呼世界和天。但陈述的话如下：\Quote{那么，我们说得对，是有一个天，还是有许多无限的天？更正确的说法是只有一个，如果它确实是按照原型创造的。}此外，在\WorkTitle{罗马人写给格林多人的书信}中写道：\Quote{一个人类无法测度的海洋及其后的世界。}因此，高贵的宗徒呼喊：\BibleQuote{天主的智慧和知识何其深奥！}\BibleRef{罗 11:33}\par

先知命令做无酵饼，岂不正是这个意思，暗示真正神圣的奥秘之言，关于非受生者和祂的权能，应当隐藏？为证实这些事，宗徒在\WorkTitle{致格林多人书}中明言：\BibleQuote{然而，在成全的人中，我们讲论智慧，但不是这世代的智慧，也不是这世代将要消灭的执政者的智慧；我们讲论的是隐藏在天主奥秘中的智慧。}\BibleRef{格前 2:6}又在另一处说：\BibleQuote{为充分认识天主的奥秘，就是基督，在祂内蕴藏着一切智慧和知识的宝藏。}\BibleRef{哥 2:2}这些事救主亲自封印，祂说：\BibleQuote{天国的奥秘是给你们知道的。}福音又说，救主以奥秘向宗徒说话。因为论到祂，预言说：\BibleQuote{祂要开口说比喻，要说出从创世以来隐藏的事。}现在，主用酵母的比喻显示隐藏；因为祂说：\BibleQuote{天国好像酵母，妇人取了，藏在三斗面里，直到全部发了酵。}\BibleRef{玛 13:33}因为三元灵魂因服从而得救，透过由信德隐藏在其内的属灵能力；或者因为赐给我们的圣言的力量，强大有力，秘密且无形地将每位接受它的人吸引到自己内，并把它保存在自己内，使其整个系统归于统一。\par

因此，\Person{梭伦}论及天主时，最有智慧地写道：—\par

\Quote{最难的是把握那不可见的心智尺度\newline{}它独自持守万物的界限。}\par

因为\Place{阿格里琴托}的诗人说：\Quote{神性，}—\par

\Quote{不能以我们的眼睛接近\newline{}也不能用手掌握；但说服的大道\newline{}至高无上，引人进入心灵。}\par

宗徒若望说：\BibleQuote{从来没有人见过天主；只有那在父怀里的独生子天主，祂给我们详述了。}\BibleRef{若 1:18}——称不可见性和不可言说性为天主的怀抱。因此有人称它为深渊，因为它包含并怀抱万有，不可接近且无界限。\par

因为如何能表达那既不是类，也不是种差，也不是种，也不是个体，也不是数；不仅如此，既不是事件，也不是事件承受者的东西呢？没有人能正确完全地表达祂。因为由于祂的伟大，祂被列为万有，是宇宙的圣父。也不能论及祂的任何部分。因为太一是不可分的；因此祂也是无限的，不是从不可测度的意义上，而是从无维度和无限制的意义上。因此祂是无形无名的。如果我们命名祂，我们这样做并不恰当，称之为太一、善、理智、绝对存在、圣父、天主、造物主、主。我们这样说并非提供祂的名字；而是由于缺乏，我们使用好的名称，以使心灵有这些作为支撑点，以免在其他方面犯错。因为每一个名称本身并不表达天主；但全部合起来表示全能者的权能。因为谓词要么从事物本身所属的方面表达，要么从它们的相互关系中表达。但这些都不适用于天主。祂也不能被论证的科学所把握。因为它依赖于首要和更清楚的原则。但没有任何东西先于非受生者。因此，我们只得藉着神圣恩宠和从祂而来的圣言来认识未知者；正如路加在\WorkTitle{宗徒大事录}中记载，保禄说：\par

\BibleQuote{诸位雅典人，我看你们在各方面都更敬畏神明。因为我巡行各处，观看你们所敬拜的，曾见一座祭坛，上面写着\InnerQuote{未识之神}。你们所敬拜而不认识的，我向你们宣告。}\BibleRef{宗 17:22}\par

% structure: p0176 source=h2 target=subsection reason=relative-heading-rank; base=h2

\subsection{第十三章 天主的知識是神聖的恩賜，按哲學家所言}

所以，凡落在名称之下的，无论他们愿意与否，都是受造的。那么，是圣父亲自吸引每一位度纯洁生活并已达到对真福不朽本性认识的人；还是我们内在的自由意志，通过认识善，如同体操运动员所说的那样，跃过障碍；然而，灵魂没有卓越的恩宠就不能生翼、翱翔，并被提升到更高的领域，脱去一切沉重之物，将自己交付给它的同类元素。\par

\Person{柏拉图}也在《\WorkTitle{美诺篇}》中说，美德是天主所赐的，如以下语句所示：\Quote{由这个论证，哦，\Person{美诺}，美德对于拥有它的人而言，是由天主眷顾而来的。}那么，那普遍来到的\OriginalTerm{真知倾向}{the gnostic disposition}岂不正是以隐晦方式被称为\Quote{天主眷顾}吗？而他又更明确地补充说：\Quote{那么，如果在这整篇论述中我们探究得当，结果就是：美德既非出于本性，也非可教，而是由天主眷顾、并非没有理智地产生于拥有它的人身上。}那属于天主的上智，作为圣父的能力，确实激发我们的自由意志，接纳信德，并以它的极峰共融回报选民的努力。\par

现在我将引述\Person{柏拉图}本人，他显然认为应当相信天主的子女。因为他在《\WorkTitle{蒂迈欧篇}》中论述可见而生之神时说道：\Quote{但说到其他神灵，要知道他们的诞生，对我们来说就太难了。我们必须相信那些早先说过的人，他们自称是诸神的后裔，并且熟知自己的祖先，尽管他们没有用可能和必然的证明来说话。}我认为希腊人不可能作出比这更清晰的见证，证明我们的救主以及那些受傅为先知的人（后者被称为天主的儿子，而主则是祂自己的圣子）在神圣事物上是真实的见证者。因此他补充说，他们既受灵感，就应当被相信。若有人以更悲壮的语调说，我们不应相信，\par

\Quote{因为告诉我这些事的并非宙斯，}\par

却让他知道，是天主自己借着祂的圣子颁布了圣经。而宣报自己事物的人是应当被相信的。\BibleQuote{除了圣子，没有人认识圣父；除非圣子愿意启示给谁，}主这样说。\BibleRef{玛 11:27}那么，根据\Person{柏拉图}，这是应当被相信的，尽管它是在旧约和新约中宣报和述说的，\Quote{没有可能和必然的证明}。\BibleQuote{除非你们相信，}主说，\BibleQuote{你们必要死在你们的罪中。}\BibleRef{若 8:24}又说：\BibleQuote{信的人有永生。}\Quote{凡投靠祂的人都是有福的。}因为投靠是超越信德的。当一个人相信天主子是我们的导师时，他就相信祂的教导是真实的。正如按照\Person{恩培多克勒}，\Quote{教导使心智增长}，同样，对主的信赖使信德增长。\par

那么，我们说：诋毁哲学、贬抑信德、赞美不义并以放荡生活为乐，这是同一类人的特征。但现在信德，如果它是灵魂的自愿同意，它仍然是善行的实行者、正直行为的基础；如果\Person{亚里士多德}严格地教导说，\OriginalTerm{制作}{\GreekText{ποιεῖν}}适用于理性动物和无生命之物，而\OriginalTerm{行动}{\GreekText{πράττειν}}只适用于人，那么让他纠正那些说天主是宇宙的\OriginalTerm{造物主}{\GreekText{ποιητής}}的人。而\OriginalTerm{所做之事}{\GreekText{πρακτόν}}，他说，是善的或必然的。那么，行恶不是善，因为没有人为了别的什么而行恶；凡是必然的都不是自愿的。那么行恶是自愿的，所以它不是必然的。善人与恶人的区别尤其在于意向和善欲。因为灵魂的一切堕落都伴随着缺乏约束；而由激情行事的人，是出于缺乏约束和堕落而行事的。\par

我不得不从各方面佩服那句神圣的话：\BibleQuote{我实实在在告诉你们：凡不从门进入羊栈，而从别处爬进去的，便是贼，是强盗。但从门进去的，才是羊的牧人。看门的人给他开门。}随后主解释说：\BibleQuote{我就是羊的门。}那么，人必须借着基督学习真理而得救，即使他们获得了哲学。因为如今清楚地显明了\BibleQuote{那在别的世代没有被知道的，如今却被启示给人类的子孙。}\BibleRef{弗 3:5}因为在所有思想正直的人中，对唯一全能之天主始终有一种自然的显示；而大多数没有完全摆脱对真理的羞愧的人，也在天主眷顾中领悟到永恒的恩惠。总之，\Person{查尔西顿的色诺克拉底}并非毫无希望，认为神性的观念甚至存在于非理性动物中。而\Person{德谟克利特}，尽管不情愿，也会通过他的教条的结果而承认这一点，因为他认为同样的形象从神圣本体发出，既对人发出，也对非理性动物发出。人并不缺乏神圣观念，按《创世纪》所载，人分享了灵感，被赋予了比其他有生命之物更纯净的本质。因此，\Person{毕达哥拉斯学派}说，理智由天主眷顾来到人身上，正如\Person{柏拉图}和\Person{亚里士多德}所承认的；但我们断言，圣神启示那已经相信的人。\Person{柏拉图主义者}认为，理智是神圣安排在灵魂中的流出，并把灵魂置于身体中。因为十二先知之一的\Person{岳厄尔}明确地说：\BibleQuote{此后，我要将我的圣神倾注在所有人身上，你们的儿女们要讲预言。}\BibleRef{岳 2:28}但圣神并非作为天主的一部分而在我们每人心中。这安排如何发生，以及圣神是什么，我们将在论预言和论灵魂的书中说明。但按照\Person{赫拉克利特}，\Quote{不信任善于隐藏知识的深度，因为不信任逃避无知。}\par

% structure: p0184 source=h2 target=subsection reason=relative-heading-rank; base=h2

\subsection{第十四章 希腊人从希伯来人的剽窃}

让我们完整地添加以下内容，现在更清楚地展示希腊人从外方哲学的剽窃。\par

现在斯多噶派认为天主如同灵魂，本质上是身体和灵。这些在他们的著作中随处可见。暂且不要将他们那些寓言当作诺斯替真理来审视；他们是否像灵巧的竞技者一样，表面上是一回事，实则意指另一样。他们声称天主遍及一切存在；而我们只称祂为造物主，且是藉着圣言而创造的造物主。他们被《智慧篇》中的话误导了：\Quote{祂因自己的纯洁，遍及并渗透万有}；\BibleRef{智 7:24} 因为他们不明白这话是指着智慧说的，智慧是天主创造中的首生者。\par

就依他们所说吧。但哲学家们，斯多噶派、柏拉图、毕达哥拉斯，甚至逍遥派的亚里士多德，都认为物质是原初原理之一，而非只有一个原初原理。那么让他们知道，他们所谓的物质，被他们说成是无性质、无形的，而柏拉图更大胆地称之为非存在。难道他不是非常奥秘地，知道真正实在的第一原因是唯一的，用这样的话说：\Quote{那么，现在让我们持这样的意见。关于宇宙的第一原理或诸原理，或者我们应该对这一切持有何种意见，我们现在不必谈论，没有别的原因，只因为按照现在的论述形式，难以解释我们的思想。}但无疑那先知式的表述：\Quote{大地是混沌空虚的}为他们提供了物质本质的基础。\par

而\Quote{偶然}的引入由此启发伊壁鸠鲁，他误解了\Quote{虚而又虚，万事皆虚}这句话。而亚里士多德从这篇圣咏中提出天主眷顾仅延伸到月球：\Quote{主，祢的慈爱上达高天；祢的信实远及云霄。}因为在主来临之前，先知奥秘的解释还未显明。\par

另一方面，死后的惩罚和火的刑罚，被所有诗学缪斯和希腊哲学从蛮族哲学中剽窃了。因此，柏拉图在《共和国》最后一卷中用这样的话明确说：\Quote{那时，这些人相貌凶猛如火，站在旁边，听到声音，就抓住一些人，带到一边；将阿里代乌斯等人的手脚和头捆起来，扔下去，剥皮，沿着道路拖曳，用荆棘撕裂他们的肉。}因为火人象征那些抓住并惩罚恶人的天使。经上说：\Quote{祂以风为使者，以火焰为仆役。}这由此证明灵魂是不朽的。因为受折磨或受矫正的，既然有感觉，就是活着，尽管说是受苦。难道柏拉图不知道火河、地渊和塔尔塔罗斯——被蛮族人称为革赫纳——他还预言性地命名了科库图斯、阿克戎和皮里弗勒格同，并引入这样的惩戒性折磨来训练吗？\par

但指明\Quote{天使}，正如经上所说\Quote{这些小子中的、最小的，他们看见天主}，以及守护天使对我们的眷顾，他不避讳地写道：\Quote{当所有灵魂都选择了几种生命，按所抽的签，他们依次走向拉刻西斯；她给每个灵魂派去他所选择的神祇作为生活的向导和共同完成目的者。}也许苏格拉底的神祇也为他想到了类似的东西。\par

再者，哲学家们从梅瑟那里听到后，就教导世界是被造的。所以柏拉图明确说：\Quote{世界究竟是没有开始，还是从某个开始获得了起源？因为它是可看见的，所以是可触摸的；因为可触摸，所以有身体。}他接着说：\Quote{要找到这个宇宙的造物主和父亲是困难的。}他不仅表明世界是被造的，而且指出它是作为儿子由祂所生，并被称作其父，因为它从祂独一而来，从无有中生出。斯多噶派也持有世界是被造的观点。\par

而蛮族哲学所讲的那个魔鬼，即众魔之首，是一个邪恶的灵，柏拉图在《法律篇》第十卷中用这些话断言：\Quote{我们必须说，那遍及一切运动之物的灵，也遍及上天吗？那么，是一个还是多个？我说多个，为你回答。让我们假设不少于两个——一个是仁慈的，另一个有能力成就相反的事。}同样在《费德若斯篇》中他写道：\Quote{现在还有其他邪恶。但某个神祇将快乐与当前大多数事物混合了。}此外，在《法律篇》第十卷中，他明确说出了那宗徒性的思想：\Quote{我们的战斗不是对抗血和肉，而是对抗率领者，对抗掌权者，对抗天界中的邪恶之神；}他这样写道：\Quote{因为我们一致认为，天上充满许多善的存在；但也充满这些的对立面，而且更多；我们断言这样的战斗是不死的，需要非凡的警惕。}\par

外邦哲学认识可理解世界和可感知世界——前者是原型，后者是所谓模型的肖像；并将前者归于\OriginalTerm{太一}{Monad}，因其由理智所感知，将可感知世界归于数目六。因为六被毕达哥拉斯派称为婚姻，作为生育之数；它将无形天、圣地和理智之光置于太一中。因为\Quote{在起初，}据说，\Quote{天主创造了天地，地是看不见的。}接着又说：\Quote{天主说：有光！就有了光。}\BibleRef{创 1:1}而在物质的宇宙论中，他创造了一个固体天（固体是可感知的），一个可见的地和一个可见的光。\Person{柏拉图}岂非由此显得把生物的理念留在理智世界中，并按照它们的种类把理智对象变成可感知的种？所以\Person{梅瑟}说得对，\Person{柏拉图}称为\Quote{地上的帐幕}的身体是由泥土形成的，而理性的灵魂是由天主吹入人面上的。因为他们说，那主宰官能位于那里；把感官进入第一个人解释为灵魂的增加。\par

因此，人也说\Quote{是按天主的肖像和模样造的。}因为天主的肖像是神圣而尊贵的圣言，不可受苦的人；肖像的肖像乃是人的理智。你若想以另一名称认识模样，便会在\Person{梅瑟}那里发现它的名称：神圣的对应。因为他说：\Quote{跟随你们的主天主，遵守他的诫命。}\BibleRef{申 13:4}我认为所有有德之人都是天主的仆人和追随者。因此斯多葛派说哲学的目的是按自然生活；\Person{柏拉图}说肖似天主，正如我们在\WorkTitle{杂录}第二部中所表明的。而斯多葛派的\Person{齐诺}，从\Person{柏拉图}那里借用，\Person{柏拉图}又从外邦哲学借用，说所有善人彼此都是朋友。因为\Person{苏格拉底}在\WorkTitle{费德罗篇}中说：\Quote{恶人做恶人的朋友，善人不做善人的朋友，这并没有被规定；}他在\WorkTitle{吕西斯篇}中也充分证明，友谊从不在邪恶与恶习中保持。雅典异乡人同样说：\Quote{有一种取悦天主并符合天主的品行，基于一个古老的基本原则，即同类喜爱同类，只要在度量之内。但超出度量的事物既与度量内的事物不相容，也与度量外的事物不相容。现在，天主是我们万物的度量。}接着，\Person{柏拉图}补充说：\Quote{因为每一个善人像每一个其他善人；同样，由于像天主，他便被每一个善人和天主所喜爱。}于此，我恰想起以下的话。在\WorkTitle{蒂迈欧篇}的结尾，他说：\Quote{你必须按照其原始本性，将那感知者同化于那被感知者；正是通过这种同化，你才能达到诸神为人设定的最高生活的目的，无论是现在还是将来。}因为那些与这些有同等的能力。那寻求的，必不停止，直到寻见；寻见了，必要惊叹；惊叹了，必要为王；为王了，必要安息。还有什么？\Person{泰勒斯}那些话不也是从这些而来吗？天主永远受到荣耀，且被我们称为鉴察人心者，他解释这一事实。当\Person{泰勒斯}被问及：神性是什么？他回答说：那既无始也无终的。当另一人问：\Quote{一个人能在做任何事时逃避对神的认知吗？}他回答说：\Quote{连思想都不能逃避，如何能做呢？}\par

此外，外邦哲学承认唯有善是可敬的，美德足以获得幸福，因为它说：\Quote{看，我已将善与恶、生命与死亡摆在你眼前，为使你选择生命。}因为它称善为\Quote{生命，}选择生命是可敬的，选择反面是\Quote{恶。}善与生命的目的是成为爱天主者：\Quote{因为这是你的生命和你的年日，}即爱那趋向真理的事。这些要点更加清楚。因为救主在命令爱天主和我们的近人时说：\Quote{全部法律和先知都基于这两条诫命。}这些是斯多葛派所宣扬的教义；在此之前，\Person{苏格拉底}在\WorkTitle{费德罗篇}中祈祷说：\Quote{啊，\Person{潘}，以及你们其他的神，赐我内在的美。}而在\WorkTitle{泰阿泰德篇}中他明确说：\Quote{因为那说得好的\OriginalTerm{好}{\GreekText{καλῶς}}也是美和善的。}而在\WorkTitle{普罗泰戈拉篇}中，他向\Person{普罗泰戈拉}的同伴宣称，他遇见了一个比\Person{阿尔西比亚德}更美的人，如果智慧的东西最美的话。因为他曾说美德是灵魂的美，反之，恶是灵魂的畸形。因此，斯多葛派的\Person{安提帕特}就此主题写了三本书，题为\Quote{根据\Person{柏拉图}，唯有美是善，}表明按他的意思，美德足以获得幸福；并援引了几个与斯多葛派一致的教义。而且，生活在\Person{托勒密·费拉德尔弗斯}时代的\Person{阿里斯托布鲁}，被\BibleBook{玛加伯}{书}摘要的编纂者提及，著有大量书籍表明逍遥学派哲学源自\Person{梅瑟}法律及其他先知。情况就是这样。\par

柏拉图清楚地称我们为弟兄，作为同一天主和同一位导师的人，他在以下话语中这样说：\newline{}\Quote{因为你们这些在城邦中的人完全是弟兄（我们将继续我们的故事对他们说）。但那位形成你们的天主，在你们出生时，在你们当中那些适合统治的人身上混合了金子，因此你们受到最高的尊荣；在那些辅助者身上混合了银子；在农夫和其他工人身上混合了铁和铜。}因此，必然地，有些人拥抱和热爱那些属于知识的事物；另一些人则拥抱意见的事物。也许他是在预言那倾向于知识的蒙选本性，如果他通过假设三种本性，并不像有些人所想的那样描述三种政体：犹太人的政体（银子）、希腊人的政体（第三种）和基督徒的政体（与基督徒混合了王者的金子，即圣神，金子般的）。\par

在展示基督徒生活时，他在\WorkTitle{泰阿泰德}中这样写道：\newline{}\Quote{现在让我们谈论最高的原则。因为我们为什么要谈论那些滥用哲学的人？他们既不知道去市场的路，也不知道法庭或元老院，或城邦的任何其他公共集会。至于口头或书面的法律和法令，他们既看不见也听不见。但政治结社和公众集会的党派情绪，以及与乐师的狂欢[占据他们]；他们甚至从未梦想参与事务。有人在城邦中行为良好或恶劣吗？或者有什么邪恶从祖先传给人吗？——这些对他们来说如同海沙一样不为所知。这个人甚至不知道他不知道这一切；但实际上，他的身体只是安放并居住在城邦中，而这个人自己，根据品达的说法，飞向地下和天上，研究天文学，探索所有方面的整个自然。}\par

再次，与主的这句话\Quote{你们的话该当是：是就说是，非就说非}可以比较以下的话：\Quote{但接纳虚假、毁灭真理，在任何方面都是不合法的。}与禁止起誓相一致的，还有\WorkTitle{法律篇}第十卷中的话：\Quote{让赞美和誓言在一切事物中都不存在。}\par

总之，毕达哥拉斯、苏格拉底和柏拉图说，他们在密切注视由天主创造并不断保存的宇宙结构时，听到了天主的声音。因为他们听到了梅瑟说：\BibleQuote{他说了，就成就了}，把天主的圣言描述为一个行动。\par

并且基于人从尘土形成的说法，哲学家们经常称身体为尘土的。荷马也不犹豫地以下面的话作为诅咒：\par

\Quote{但你们全都变成土和水。}\par

正如依撒意亚所说：\Quote{并践踏他们如同泥土。}而卡利马科斯清楚地写道：\par

\Quote{那是这样的年份，\newline{}\newline{}鸟儿、游鱼、四足动物，\newline{}\newline{}都像普罗米修斯的泥土那样说话。}\par

而同一作者又说：\newline{}\par

\Quote{如果你是被普罗米修斯塑造的，\newline{}\newline{}而你不是来自别的泥土。}\par

赫西俄德关于潘多拉这样说：\newline{}\par

\Quote{并吩咐著名的赫淮斯托斯，尽他最快的速度，\newline{}\newline{}把土和水揉合，并注入人的声音和心灵。}\par

因此，斯多葛派将自然定义为一种有艺术的火焰，系统地进展到生成。而天主和祂的圣言，按照圣经，被象征性地称为火和光。但是，荷马本人难道不是在解释水从陆地退去、旱地清晰显露时，关于特堤斯和俄刻阿诺斯这样说吗：\newline{}\par

\Quote{因为现在他们长时间避免\newline{}彼此的床笫和爱情？}\par

再次，在一切事物中的能力，被希腊人中最有智慧的人归给天主；埃庇卡摩斯——他是毕达哥拉斯派——说：\newline{}\par

\Quote{没有事物逃过神明的眼目。这你应该知道。\newline{}祂是我们的监察者。对天主而言，没有不可能的事。}\par

而抒情诗人说：\newline{}\par

\Quote{天主能从阴暗的夜晚\newline{}举起无瑕的光芒，\newline{}并能在阴暗的幽暗中\newline{}遮蔽白日的纯净辉煌。}\par

唯独祂能够在白昼期间制造夜晚，祂就是天主。\par

在\WorkTitle{天文现象}中，阿拉托斯这样写道：\par

\Quote{让我们从宙斯开始；我们作为人，\newline{}决不能将他未表达而留下。一切都充满了宙斯，\newline{}街道和人群，大海也充满了，\newline{}海岸，处处我们都享受宙斯。}\par

他又说：\newline{}\par

\Quote{因为我们也是\newline{}他的后代；……}\par

即藉由受造。\par

\Quote{他，对人类温和，\newline{}显示慈爱的征兆，并唤醒他们从事工作。\newline{}因为他在天空设立了这些征兆，\newline{}散布星辰，并以星星为年份加冕；\newline{}向人显示季节的工作，\newline{}好使一切事物按确定秩序进行。\newline{}他们永远先敬拜他，最后也敬拜他：\newline{}万福，父啊，伟大的奇迹——对人类伟大的恩赐。}\par

而在他之前，荷马根据梅瑟的说法，在赫淮斯托斯打造的盾牌上塑造世界，说：\par

\Quote{上面他塑造了大地、天空和海洋，\newline{}以及所有装饰天空的星辰。}\par

因为诗歌和散文中歌颂的宙斯引导心灵走向天主。而且，可以说，德谟克利特已经写道：\Quote{少数人在光明中，他们向那个地方伸出双手，我们现在希腊人称其为空气。宙斯说一切，听一切，分配和收回一切，他是万王之王。}而更神秘地，彼奥提亚的品达，作为一位毕达哥拉斯派，说：\newline{}\par

\Quote{神和人是同一族类，\newline{}两者都从同一位母亲得到气息；}\par

即来自物质；并称万物的唯一创造者为父、首要的工匠，祂根据功绩提供通往神性的进步途径。\par

我且不提柏拉图；在致\Person{埃拉斯托}和\Person{科里斯库}的书信中，他明明从希伯来圣经中展示了某种意义上的父与子，并如此劝诫道：\Quote{在宣誓时，以并非无知的庄重与教养，邀请庄重之姊妹——万有之天主，并呼祂为主，即元首之父与创造者；你们若以真正哲学的精神研习，就必认识祂。}而在《\WorkTitle{蒂迈欧篇}》中的致辞称创造者为父，说：\Quote{你们众神中之神，我为你们之父，也是你们作品的创造者。}因此，当他说：\Quote{万物围绕着万王之王，因祂而有万物；祂是万善之因；第二位周围的是第二序列的事物；第三位周围的是第三序列的}时，我理解这无非是指天主圣三；因为第三位是圣神，第二位是圣子，万物藉着祂照圣父的旨意受造。\par

同一位柏拉图在《\WorkTitle{共和国}》第十卷中提到了\Person{厄洛斯}，他是\Person{阿尔墨纽斯}之子，即\Person{琐罗亚斯德}。\Person{琐罗亚斯德}写道：\Quote{这些是由阿尔墨纽斯之子、潘菲利亚人琐罗亚斯德所作：我在战斗中死去，到了冥府，从神明那里学到这些。}柏拉图说，这个琐罗亚斯德被置于火葬柴堆上，十二天后复活了。他也许暗指复活，也许暗指灵魂上升的道路经过黄道十二宫；他本人说，下降至生的道路也相同。同样，我们也要理解\Person{赫拉克勒斯}的十二项劳作，之后灵魂便从整个尘世中获得解脱。\par

我不略过\Person{恩培多克勒}，他这样物理地谈论万物的更新，即在于转化为火的本质，这将要发生。而最清楚的是\Person{以弗所的赫拉克利特}，他以为有一个永恒的世界，并承认一个可朽的世界——即在其秩序方面，并非与前一个不同，只是从某方面看。但他知道由普遍本质构成的不可朽的世界永远具有特定本性，这一点他清楚地说：\Quote{这个万物同一的世界，既非任何神，也非任何人所创造。它过去是、现在是、将来永远是可燃的火，按尺度点燃，按尺度熄灭。}至于他教导世界是受生且可朽的，由以下所言显示：\Quote{火的转化：首先是海；海的一半是陆地，一半是炽热蒸汽。}他说这些是能力的效果。因为火借统治万物的天主圣言，通过空气转化为湿气，这湿气仿佛是宇宙变化的种子；他称之为海。从海中又产生地、天及其所含万物。至于它们如何再次恢复并点燃，他清楚地在这些言词中表明：\Quote{海扩散并按相同的规则度量，这规则在它成为地之前就已存在。}类似地，对其他元素也应作相同理解。最著名的斯多亚派学者与他教导类似的学说，论及大火与世界治理，以及严格意义上的世界与人，并我们灵魂的持续。\par

柏拉图在《\WorkTitle{共和国}》第七卷中，称\Quote{此地的白昼为黑夜}，我想，是因为\Quote{这黑暗世界的主宰者}；\BibleRef{弗 6:12} 而灵魂进入身体，睡眠与死亡，与\Person{赫拉克利特}相似。圣神籍着大卫岂没有谕示般地预言了救主，说：\Quote{我睡了，又安眠；我醒来，因为主必扶持我}？因为祂不仅象征性地将基督的复活称为从睡眠中醒来；也将主降生成人喻称为睡眠。救主亲自吩咐：\Quote{警醒}，这等于说：\Quote{学习如何生活，并努力将灵魂与身体分离}。\par

柏拉图在《\WorkTitle{共和国}》第十卷中预言性地谈到主日，说：\Quote{当他们在草地上各自度过了七天，第八天就要出发，并在四天内到达。}草地应理解为固定球体，因其温和宜人，是虔敬者的居所；七天指七个行星各自的运行，以及整个通向安息终点的实践技艺。在漫游的天体之后，旅程引向天堂，即第八种运行和第八日。他说灵魂在第四天离开，指出经过四个元素的通道。但第七日不仅被希伯来人，也被希腊人尊为圣日；根据它，整个动物和植物的世界都运转。\Person{赫西俄德}论及它说：—\par

\Quote{第一日、第四日、第七日被视为圣日。}\par

并且再次：—\par

\Quote{而在第七日，太阳灿烂的圆轮。}\par

\Person{荷马}说：—\par

\Quote{而在第七日，圣日就来临了。}\par

并且：—\par

\Quote{第七日是圣日。}\par

并且再次：—\par

\Quote{那是第七日，一切事物都完成了。}\par

并且再次：—\par

\Quote{而在第七日早晨，我们离开\Place{阿刻戎河}。}\par

\Person{卡利马科斯}也写道：—\par

\Quote{那是第七日早晨，他们完成了所有事情。}\par

并且再次：—\par

\Quote{在好日子中，第七日和第七赛程。}\par

并且：—\par

\Quote{第七日属于首要，第七日是完美的。}\par

并且：—\par

\Quote{如今所有七颗星在星空的天上造成，\newline{}在圆圈中闪耀，如同年份出现。}\par

\Person{梭伦}的挽歌也极度神化第七日。\par

又怎样呢？这岂不与圣经所说相似，\Quote{让我们移除义人，因他令我们不安？}\newline{}\BibleRef{智 2:12} 那时柏拉图几乎预言了救恩的计划，在\WorkTitle{理想国}第二卷中说道：\Quote{那被认定为义的人将受鞭打，被拉上刑架，被捆绑，被挖去眼睛；最后，受尽一切苦楚后，将被钉死。}\par

苏格拉底派的安提斯泰尼，转述那先知的话：\Quote{你们把我比作谁呢？上主说，} 说：\Quote{天主不像任何人；因此无人能藉着形象认识祂。}\par

雅典人色诺芬也以如下话语表达了类似的情怀：\Quote{那震动万物、自身却不动者，显然是一位伟大而强有力的。但祂形像如何，却显不出。太阳也一样，愿普照四方，却不许任何人直视自己。若有人胆敢凝视，便会失去视力。}\par

\Quote{什么肉身能用肉眼看见天上、真实的\newline{}不朽的天主，祂的居所在天极？\newline{}甚至连在太阳的明亮光芒前\newline{}凡人，作为可朽者，也不配站立，\newline{}}—\par

女先知此前已说过。于是科洛丰的色诺芬尼正确地教导说天主是独一且非物质的，他补充道：—\par

\Quote{在神与人中有一位至高的天主；\newline{}在形体、心智上与凡人不同。\newline{}}\par

再者：—\par

\Quote{但人们以为神是出生的，\newline{}穿着衣服，且有声音和形状。\newline{}}\par

再者：—\par

\Quote{倘若牛或狮子有手\newline{}或能用手像人一样描绘作品\newline{}若野兽画出神的形像\newline{}马就会把它们画得像马\newline{}对牛则画成牛，并把它们的身体造\newline{}成它们自己所拥有的形状。\newline{}}\par

那么，让我们听听抒情诗人巴库利德斯谈论神圣者：—\par

\Quote{他们从不屈服于可怕的疾病，\newline{}且无可指责；丝毫不像凡人。\newline{}}\par

还有斯多葛派的克莱安西斯，他在一首论神的诗中这样写道：—\par

\Quote{如果你问善的本性是什么，请听——\newline{}那是规律的、公正的、神圣的、虔诚的，\newline{}自治的、有用的、公平的、适当的，\newline{}严肃的、独立的、永远有益的，\newline{}不感到恐惧或悲伤、有利的、无痛苦的，\newline{}有益的、愉快的、安全的、友善的，\newline{}受尊敬的、与自身一致：可敬的，\newline{}谦卑的、谨慎的、温良的、热心的，\newline{}永久的、无可指责的、永远持续的。\newline{}}\par

同一位，暗地里贬低群众的偶像崇拜，补充道：—\par

\Quote{每个注重意见的人都是卑贱的，\newline{}因为他想从中获取任何好处。\newline{}}\par

那么我们不应按照群众的意见来思想天主。\par

\Quote{因为我不认为他会偷偷地，\newline{}模仿恶棍的伪装，\newline{}作为男人上你的床，\newline{}}\par

安菲翁对安提俄珀说道。索福克勒斯也清楚地写道：—\par

\Quote{宙斯娶了他的母亲，\newline{}不是以黄金的形像，也不是覆盖着\newline{}天鹅的羽毛，如同他使普琉戎女子怀孕那样；\newline{}而是一个彻头彻尾的男人。\newline{}}\par

他继续下去，补充道：—\par

\Quote{很快那奸夫站在了新娘的阶梯上。\newline{}}\par

然后他更清楚地详述了传说中的宙斯的淫荡：—\par

\Quote{但他既不碰食物也不碰洁净的水，\newline{}而是心中刺痛地上了床，整夜\newline{}纵欲。\newline{}}\par

但让这些留给剧院的愚行吧。\par

赫拉克利特清楚地说：\Quote{但对于那永恒的圣言，人们既在听见之前，也在初次听见之时，都不能理解。} 而里尔琴手梅拉尼皮德斯在歌中唱道：—\newline{}\Quote{听我，噢父，人的奇迹，\newline{}永活灵魂的统治者。\newline{}}\par

\Quote{听我，噢父，人的奇迹，\newline{}永活灵魂的统治者。\newline{}}\par

伟大的巴门尼德，正如柏拉图在\WorkTitle{智者篇}中所说，这样论及天主：—\par

\Quote{非常，因为祂是非受生且不可毁灭的，\newline{}完整的、独生的、不动的、无始的。\newline{}}\par

赫西俄德也说：—\par

\Quote{因为祂是众神中所有的王与主。\newline{}无人能与众神在权能上争胜。\newline{}}\par

更有甚者，悲剧远离偶像，教导人仰望上天。索福克勒斯，如同赫卡泰乌斯（他撰写了关于亚伯拉罕和埃及人的史书）所说，在舞台上清楚地呼喊：—\par

\Quote{事实上只有一位，天主是独一的，\newline{}祂创造了天和广阔的大地，\newline{}深渊的蓝色波涛，以及风的力量。\newline{}但我们凡人，许多人心灵迷途甚远，\newline{}为了安慰我们的苦难，立起了\newline{}神的形像——或用石头，或用黄铜，\newline{}或用黄金或象牙制成的雕像；\newline{}并为这些献上祭品和虚浮的节庆，\newline{}自以为是虔诚的。\newline{}}\par

欧里庇得斯在舞台上的悲剧中说：—\par

\Quote{你看见这崇高无边的以太吗？\newline{}它用湿润的臂膀环抱着大地。\newline{}这就是宙斯，\newline{}并视之为天主。\newline{}}\par

而在\WorkTitle{珀里托俄斯}剧中，同一人以悲剧风格写下这些诗句：—\par

\Quote{你，自生者，在以太的轮上\newline{}转动了万物的本性，\newline{}在你周围，光明和暗黑的繁星之夜，\newline{}以及无数的星辰之群，也不停地舞蹈。\newline{}}\par

因为在那里他说创造性的理智是自生的。接下来适用于宇宙，其中有光明和黑暗的对立。\par

欧福里翁之子埃斯库罗斯，也非常庄严地论及天主：—\par

\Quote{以太是宙斯，宙斯是大地，宙斯是天；\newline{}宇宙是宙斯，以及所有之上。\newline{}}\par

我知道柏拉图赞同赫拉克利特，后者写道：\Quote{惟有那智慧的东西不愿被表达，就是指宙斯的名字。} 再者，\Quote{法律就是服从一人的意志。} 你若想引用那句格言：\Quote{有耳能听的，就让他听吧。} 你会发现以弗所人表达如下：\Quote{那些听了却不理解的人就像聋子。谚语反对他们，说他们在场却不在场。}\par

但你想从希腊人那里明确地听到一个第一本原吗？洛克里的提迈欧，在\WorkTitle{论自然}中，将以如下话语作证：\Quote{一切事物有一个非受生的第一本原。因为如果它受生，它就不再是第一本原；而第一本原将是从它受生的那个。} 因为这正确的意见源自以下经文：\Quote{听，} 它说，\Quote{噢以色列，上主你的天主是独一的，你只应事奉祂。}\newline{}\BibleRef{申 6:4}\par

\Quote{看，祂全然可靠、全然无误。}\par

\Person{西彼尔}说。\par

\Person{荷马}也凭藉一种幸福的占卜，显然在以下诗句中提及了\Person{圣父}和\Person{圣子}：—\par

\Quote{若乌提斯独自一人向你施暴，\newline{}且无法逃脱宙斯降下的疾病——\newline{}因为独眼巨人并不理会持埃癸斯的宙斯。}\par

而在他之前，\Person{俄耳甫斯}谈到这一点时说道：—\par

\Quote{伟大宙斯之子，持埃癸斯的宙斯之父。}\par

\Person{卡尔西顿人色诺克拉底}提及至高的宙斯和次等的宙斯，留下了圣父和圣子的提示。\Person{荷马}虽然将众神描绘成受人类情欲支配，却似乎认识神圣的存在，而这正是\Person{伊壁鸠鲁}所不敬畏的。于是他说道：—\par

\Quote{为何，珀琉斯之子，你虽是凡人，\newline{}却以快脚追赶我这不朽的天神？\newline{}你还不曾知道\newline{}我是一位神吗？}\par

因为他表明，神性既不能被凡人抓住，也不能用脚、手、眼或任何身体来领悟。\BibleQuote{你们将主比作谁？或将他比成什么样式？}圣经说。\BibleQuote{工匠不是造了偶像吗？金匠熔金镀金，}以及后面的内容。\par

喜剧诗人\Person{埃庇卡摩斯}在他的\WorkTitle{共和国}中清楚地说到\Person{圣言}，其词如下：—\par

\Quote{人类的生活只需要计算和数字，\newline{}我们靠数字和计算而活，因为它们拯救凡人。}\par

然后他明确地补充道：—\par

\Quote{理性统治凡人，且独自保存习俗。}\par

接着：—\par

\Quote{人里面有推理，也有神圣理性。\newline{}理性植入人内，以供给生命和生计，\newline{}但神圣理性在所有情况下照料技艺，\newline{}始终教导它们做有益之事。\newline{}因为人并非发现技艺，而是天主带来；\newline{}人的理性源自神圣理性。}\par

圣神也通过\Person{依撒意亚}呐喊：\BibleQuote{你们的众多祭献有何用？主说。公羊的全燔祭我已饱足，肥犊的脂油和公牛的血我不喜爱；}稍后又加上：\BibleQuote{你们要洗濯，自洁。从你们灵魂中除掉邪恶，}等等。\par

喜剧诗人\Person{米南德}用这些词写道：—\par

\Quote{若有人通过献祭，一群\newline{}公牛或山羊，啊潘菲鲁斯，凭宙斯之名，\newline{}或类似之物；通过制作工艺品，\newline{}金衣或紫袍，象牙\newline{}或翡翠的雕像，以为借此\newline{}能使天主息怒，他就错了，\newline{}心思空虚。因为人必须证明\newline{}自己是个有价值的人，既不玷污处女，\newline{}也不是奸夫，不偷窃，不杀人\newline{}因贪恋世间财富，啊潘菲鲁斯。\newline{}不，连一根针线也不要贪图。因为天主\newline{}看见你，就在你身旁。}\par

\BibleQuote{我是近处的天主，}\BibleRef{耶 23:23}经上说，\BibleQuote{不是远处的天主。人能在隐秘处行事，而我看不见他吗？}\par

\Person{米南德}再次用那圣经的话进行释义：\BibleQuote{献上正义的祭献，并信赖主。}他这样写道：—\par

\Quote{甚至不要觊觎别人\newline{}的一根针线，最亲爱的朋友；\newline{}因为天主喜悦正义的作为，并因此\newline{}允许那昼夜耕田\newline{}劳碌的人增加世间财富。\newline{}但要向天主献祭，并成为正义之人，\newline{}闪耀的不是华服，而是你的内心；\newline{}当你听到雷声时，不要逃跑，\newline{}自知无一过错，\newline{}善先生，因为天主永远在旁看见你。}\par

\BibleQuote{你还在说话时，}圣经说，\BibleQuote{我要说：我在这里。}\BibleRef{依 65:24}\par

喜剧诗人\Person{狄菲鲁斯}再次论及审判如下：—\par

\Quote{你以为，啊尼克拉图斯，那些\newline{}生前享尽各种奢华的死者\newline{}能逃脱神，如同被遗忘？\newline{}有一只正义之眼，看见一切。\newline{}因为我们认为，通往哈得斯有两条路——\newline{}一条为善人，另一条为恶人。\newline{}但如果大地永远掩埋两者，那么\newline{}就抢劫、偷窃、掠夺、作乱吧。\newline{}但你们不要犯错。因为在哈得斯里有审判，\newline{}那是万有之主天主\newline{}要执行的，他的名我畏惧不敢提，\newline{}他赐罪人地上寿命延长。\newline{}若任何凡人认为，日复一日\newline{}作恶时能避开众神的锐眼，\newline{}他的思想邪恶；当正义有了\newline{}闲暇，他必被察觉如此\newline{}思想。看，无论你是谁，若你说\newline{}没有神。有，确实有。\newline{}若有人本性邪恶，行恶，\newline{}让他赎罪吧；因为这样的人\newline{}不久将受到应得的惩罚。}\par

与此一致，悲剧也在以下诗行中说道：—\par

\Quote{因为那时刻必将来临，必将来临，\newline{}那时金眼的以太将打开它\newline{}储藏的火宝；吞噬的火焰\newline{}将猛烈烧尽地上一切\newline{}和天上一切。}\par

稍后他又补充道：—\par

\Quote{当整个世界消失，\newline{}海洋深渊的波浪尽都退去，\newline{}大地无树；火焰包围中，\newline{}空气不再生出有翼的族类；\newline{}那毁灭一切的，将恢复一切。}\par

我们在\Person{俄耳甫斯}的赞美诗中也能找到与此类似的表达，写如下：—\par

\Quote{因为隐藏一切后，又将它们带回到\newline{}喜悦的光明中，从他神圣的心\newline{}发出，满有眷顾。}\par

如果我们终其一生圣洁而正义地生活，我们在此世是有福的，在离开之后将更为有福；不是暂时拥有幸福，而是得以安息于永恒。\par

\Quote{在同一炉灶和餐桌旁，与\newline{}其余不朽的神明同坐，我们全然自由\newline{}免于人间的患难，不受伤害，}\par

\Person{恩培多克勒}的哲学诗歌说。所以，根据希腊人，没有谁大到可以超越审判，没有谁小到可以逃脱其注意。\par

同一位\Person{俄耳甫斯}这样说：—\par

\Quote{但你要注视神圣的圣言，留心，\newline{}正确守护心灵这\newline{}理智的容器，并正直地行走，\newline{}只将你的目光\newline{}投向世界不朽的君王。}\par

再次，关于天主，他说祂不可见，只有一位\Person{加色丁}族人——意指\Person{亚巴郎}或其子——认识祂，他这样说：—\par

\Quote{但有一位加色丁族后裔；\newline{}因为他深知太阳的轨道，\newline{}以及天球如何\newline{}绕地运动，在圆周中\newline{}绕轴均匀旋转，风如何\newline{}在空中和海上驾驭战车。}\par

然后，彷佛在释义这句话：\BibleQuote{天是我的宝座，大地是我的脚凳，} \BibleRef{依 66:1}，他补充说：\par

\Quote{但在广阔的天上，\newline{}祂稳固地坐在金座上，\newline{}大地在祂脚下。\newline{}祂的右手伸到海洋的边界；\newline{}山岳的根基因祂的忿怒而颤抖，\newline{}无法承受祂的大能。祂全然属天，\newline{}并在地上完成一切。\newline{}祂是元始、中与终末。\newline{}但我不敢提及祢。\newline{}我四肢和心灵颤抖。祂从高处统治。}\par

等等。因为这些经文指示了这些预言性的话：\Quote{若祢打开天，山岭必因祢的临在而战栗；它们要熔化，如蜡在火前熔化；}以及依撒意亚所说：\Quote{谁曾用手掌量过天，用拳量过大地？}又说：\par

\Quote{统治以太、\newline{}冥府、海洋与陆地的主，\newline{}祢以雷电震撼奥林帕斯的坚固居所。\newline{}众魔惧怕祢，众神之群也敬畏祢。\newline{}命运女神虽无情，\newline{}也顺服祢。不死者，我们母亲的父亲！\newline{}祢的忿怒使万物摇动；祢驱动风，\newline{}以云遮蔽世界，以太被祢的闪电劈开。\newline{}星辰的秩序属于祢，\newline{}它们按祢不变的号令运行。\newline{}在祢燃烧的宝座前，\newline{}天使侍立，多行奇事，\newline{}奉命为人类做一切事。祢的春天光华，饰以紫花；\newline{}祢的冬日以寒云侵袭；\newline{}祢的秋日赐予喧闹的巴克斯。}\par

然后他明确指名全能的天主说：\par

\Quote{不死不朽者，\newline{}唯对不朽者才可述说！来吧，\newline{}最伟大的神，伴随刚强的必然。\newline{}可畏、不可战胜、伟大、不死者，\newline{}以太加冕于祢。}……\par

通过\Quote{我们母亲的父亲}（\OriginalTerm{我们母亲的父亲}{\GreekText{μητροπάτωρ}}）这一措辞，他不仅暗示了从无中的创造，而且给那些引入流溢说的人提供了想象天主有配偶的场合。他释义了那些预言性的圣经——依撒意亚中的：\BibleQuote{我是那建立雷霆、创造风的手奠定了天上的万军；}\BibleRef{亚 4:13}以及梅瑟中的：\BibleQuote{看，看，我就是祂，除我以外没有别的神：我使人死，也使人活；我击伤，也医治；无人能从我的手中救出。}\BibleRef{申 32:39}\par

\Quote{祂从善中，为凡人种植灾祸、\newline{}残酷的战争和流泪的苦难，}\par

据\Person{俄耳甫斯}所说。\par

帕罗斯岛的\Person{阿基洛库斯}也有类似的话。\par

\Quote{哦，宙斯，诸天的权能属于祢，\newline{}祢将暴虐与不义加于人。}\par

再让色雷斯的\Person{俄耳甫斯}为我们歌唱：\par

\Quote{祂的右手伸到海洋的边界；\newline{}大地在祂脚下。}\par

这些显然来自以下经文：\BibleQuote{上主必拯救有居民的城邑，并用手掌握全地如巢穴；}\BibleRef{依 10:14}\BibleQuote{是上主以能力造大地，}正如\Person{耶肋米亚}所说，\BibleQuote{以智慧立定大地。}\BibleRef{耶 10:12}此外，\Person{福库利得斯}称天使为\Quote{神灵}，他在以下诗句中解释说其中一些是善的，另一些是恶的（因为我们也知道有些是背弃者）：\par

\Quote{有神灵——有些在这里，有些在那里——被派来管理人；\newline{}有些在人之初（进入生活时）是为了抵御灾祸。}\par

所以，喜剧诗人\Person{菲勒蒙}也正确地用这些话摧毁偶像崇拜：\par

\Quote{对我们而言，\newline{}命运并非神祇：根本没有这样的神。\newline{}偶然发生、自然而至的事，才被称为命运。}\par

悲剧诗人\Person{索福克勒斯}说：\par

\Quote{诸神也非万事如意，\newline{}只有宙斯例外；因为他是元始，也是终末。}\par

而\Person{俄耳甫斯}：\par

\Quote{同一位大能者，伟大的、\newline{}燃烧的诸天，是唯一的神。万物是同一存在；\newline{}其中一切——火、水、大地——都围绕旋转。}\par

等等。\par

抒情诗人\Person{品达}，彷佛在酒神狂欢中，明确地说：\par

\Quote{神是什么？万有。}\par

又说：\par

\Quote{神，祂造了所有凡人。}\par

当他说——\par

\Quote{你身为凡人，\newline{}何等渺小，竟想为自己求得智慧？\newline{}凡人的心智难以洞察神的谋略；\newline{}你由凡人之母所生，}\par

他思想出自以下经文：\BibleQuote{谁曾知晓上主的心意，或作过祂的谋士？}\BibleRef{依 40:13} \Person{赫西俄德}也同意上文所说的，他写道：\par

\Quote{没有先知，出自地上的人，\newline{}能知晓持羊皮盾的宙斯的心意。}\par

同样地，雅典的\Person{梭伦}在\WorkTitle{哀歌}中，追随\Person{赫西俄德}写道：\par

\Quote{不朽者的心意对人完全未知。}\par

再者，\Person{梅瑟}曾预言女人因过犯将在艰难和痛苦中生产，一位并非无名的诗人写道：\par

\Quote{白昼他们永无安息，\newline{}辛劳与灾祸不断；\newline{}夜间呻吟不绝。\newline{}神将忧苦赐予人。}\par

此外，当\Person{荷马}说——\par

\Quote{天父亲自举起金天平，}\par

他暗示天主是公义的。\par

喜剧诗人\Person{米南德}在启示神时，说：\par

\Quote{每个人出生时，\newline{}都有一位守护神灵被指派为他一生的良善向导。\newline{}因为神灵是恶的并损害善生，\newline{}是不可想象的。}\par

然后他补充说：\par

\Quote{\OriginalTerm{一切善者就是神}{\GreekText{Ἅπαντα} \GreekText{δ᾿} \GreekText{ἁγαθὸν} \GreekText{εἱναι} \GreekText{τὸν} \GreekText{Θεόν}},}\par

意思是\Quote{每一个善者就是神}，或者更恰当地说，\Quote{神在一切事上都是善的}。\par

悲剧诗人\Person{埃斯库罗斯}在展示神的能力时，不犹豫地称祂为至高者，用这些话：\par

\Quote{将天主与凡人分开；\newline{}不要以为祂像你一样是肉身的。\newline{}你不认识祂。\newline{}有时祂显现为火，可畏的力；\newline{}有时为水；有时为黑暗；\newline{}在兽类中隐约映现，在风、云、\newline{}闪电、雷、雨中；\newline{}海洋和岩石侍奉祂，每个泉源和水流、\newline{}江河。群山颤抖，大地、\newline{}浩瀚的海渊和山岳的崇高，\newline{}当君主的可畏之眼注视它们：\newline{}至高全能天主的荣耀是无所不能的。}\par

难道他看起来不是在改写这段话：\Quote{在主的面前大地战栗吗？}除此之外，最具预言能力的\Person{阿波罗}被迫——如此为天主的荣耀作证——论及\Person{雅典娜}，当\Person{玛代人}向\Place{希腊}开战时，她曾为\Place{阿提卡}恳求和祈求\Person{宙斯}。神谕如下：——\par

\Quote{帕拉斯不能平息奥林匹亚的宙斯，\newline{}尽管她以许多言语和明智的劝告祈祷；\newline{}但他将把许多不死者的庙宇交给吞噬之火，\newline{}这些庙宇现在因恐惧而颤抖，浸透在汗水中；}\par

等等。\par

\Person{塞阿里达斯}在他的书\WorkTitle{论自然}中写道：\Quote{那时存在一个真正真实的万有之本原[\OriginalTerm{第一原则}{first principle}]}——一个。因为在起初的那个存有是唯一且单独的。\par

\Quote{除了大王以外没有别的，}\par

\Person{奥尔弗斯}说。与此一致，喜剧诗人\Person{狄菲洛斯}极其格言式地说道，那\par

\Quote{万有之父，\newline{}惟独向祂献上不断的崇敬，\newline{}这等福泽的发明者和作者。}\par

因此，柏拉图正确地\Quote{使最好的本性习惯于达到那项我们先前所说的最高研究，既看见\OriginalTerm{善}{the good}，也完成那上升。而这，显然，不是抛掷陶片；而是灵魂从黑夜般的白昼转向那真正回归真实的存在，我们将断言这便是真哲学。}他认为凡分享此等境界的人属于\OriginalTerm{金族}{the golden race}，当他说：\Quote{你们都是弟兄；而那些属于\OriginalTerm{金族}{the golden race}的人在各方面最能做出最准确的判断。}\par

那么，万有的圣父和造化者，被万物按照万有的方式，通过内在的能力并无需教导而领会——无生命之物与有生命之受造物同情；而在有生命者中，有些已经是不死的，在白昼行动。但那仍然可朽的，有些在恐惧中，仍在母腹中被携带；其他则凭自己的独立理性自行调节。而在人中，全是希腊人和异邦人。但任何地方的农人、游牧者，甚至城市居民，没有一个种族能活着而不被灌输了具有一个至高存有的信仰。因此，每个东方民族，每个触及西岸的民族，或北方，以及向南的每一个，——全都对那设立了统治的那一位怀有同一的先入之见；因为祂最普遍的运行同样遍及一切。更甚的是，希腊人中致力于探究的哲学家们，从异邦哲学出发，将眷顾归给那\Quote{不可见的、唯一的、最有力的、最精巧的、至高的万有最美之原因；}——并不知晓这些真理的推论，除非由我们教导，甚至不知道天主如何能被自然地认知；而只是，如我们已多次说过的，通过真正的\OriginalTerm{迂回说法}{periphrasis}。因此宗徒正确地断言，\BibleQuote{祂只是犹太人的天主，不也是希腊人的吗？}——不仅先知性地说明，在希腊人中信仰的希腊人会认识天主；而且暗示，就能力而言，主是万有的天主，真正万有的君王。因为他们不知道祂是什么，也不知道祂如何是主、圣父和造化者，以及其余真理的系统，除非由其教导。因此，先知的话语也具有与宗徒之言相同的力量。因为\Person{依撒意亚}说，\BibleQuote{若你们说：我们信赖主我们的天主：现在请与我主亚述王结盟。}他又加：\BibleQuote{现在，我们上来攻打这地，难道不是出于主吗？}而\Person{约纳}，他自己也是先知，在他所说的话中暗示了同一回事：\BibleQuote{船长到他那里，对他说：你为什么打鼾？起来，呼求你的天主，好使我们得救，免于灭亡。}因为\BibleQuote{你的天主}这个说法，他仿佛是对一个因知识而认识祂的人说的；而\BibleQuote{愿天主拯救我们}这个说法，则揭示了外教人心中的意识，他们曾用心思考万有的主宰，却尚未相信。同样的话又再次出现：\BibleQuote{他对他们说：我是主的仆人；我敬畏主，天上的天主。}又再次：\BibleQuote{他说：我们切不可为这人的性命而灭亡。}而\Person{玛拉基亚}先知明确地显示天主说：\BibleQuote{我不从你们手中接受祭品。因为从日出到日落，我的名在外邦人中受颂扬；并且在每一处，祭品奉献给我。}又再次：\BibleQuote{因为我是伟大的君王，全能的上主说；我的名在万国中显明。}什么名？圣子在已经信仰的希腊人中显明圣父。\par

柏拉图在以下段落中展示了\OriginalTerm{自由意志}{free-will}：\Quote{德性没有主人；每个人依其尊崇或轻蔑它的程度，而分有它。过错在于自由选择的应用。}但天主是无辜的。因为祂决不是邪恶的作者。\par

\Quote{哦，好战的\Person{特洛伊人}，}抒情诗人说，——\par

\Quote{高高统治的\Person{宙斯}，他洞察万物，\newline{}并不是凡人巨大苦难的原因；\newline{}而是在所有人的能力之中，找到\newline{}正义，神圣，纯洁，\newline{}秩序的伴侣，\newline{}以及明智的\Person{特弥斯}，\newline{}你们是有福者的子孙，\newline{}因找到她作为你们的伴侣。}\par

而\Person{品达}也明确引入\OriginalTerm{救主宙斯}{Zeus Soter}，\Person{特弥斯}的配偶，宣告他为王、救主、公义者，在以下诗句中：——\par

\Quote{首先，谨慎的\Person{特弥斯}，天生神圣，\newline{}乘着金色骏马，在大洋之岩旁，\newline{}命运女神们将她带到崇高的阶梯，\newline{}奥林帕斯的光辉入口，\newline{}作为\OriginalTerm{救主宙斯}{Zeus Soter}永远的配偶，\newline{}而她把时序女神们，金冠、美果、善良，生育出来。}\par

因此，那不顺服真理、因人的教导而自大的人，是可悲而可怜的，据\Person{欧里庇得斯}说：——\par

\Quote{谁看见这些事，却不领悟天主，\newline{}反而口吐崇高论题，远远抛出\newline{}邪恶的欺骗；固执其中，舌头\newline{}射出它的箭矢，关于不可见的事，\newline{}毫无意义。}\par

让那愿意的人，走近真教导，从\Person{巴门尼德}（\Place{埃利亚}人）那里学习，他应许说：——\par

\Quote{以太之性，以及所有星辰，\newline{}你将在以太中知晓，以及效果，\newline{}神圣太阳之明灯的所有无形光辉，\newline{}及其产生之源。圆眼月亮\newline{}旋转的影响与本性，你\newline{}将学会；环绕的天穹你将知晓；\newline{}它从何而生；必然如何将它\newline{}摄取并锁住，以守护星界的边际。}\par

而\Person{梅特罗多罗斯}虽是一位伊壁鸠鲁派信徒，却以神圣感动说了这样一番话：\Quote{要记住，哦，\Person{梅内斯特拉特}，作为拥有有限生命的凡胎，你的灵魂已经上升，直到你看见无尽的时间，和万物的无限；以及将要发生之事，和已经发生之事；}当与蒙福之合唱队一起时，按照\Person{柏拉图}的话，我们将凝视蒙福的景象与神视；我们跟随\Person{宙斯}，而其他人跟随其他神祇，若可如此说，以领受进入最蒙福的奥迹的入门礼：我们将庆祝它，我们自身成为完美的，且不被在我们时间终结时等候我们的恶事所扰；并被引入对完美而宁静的异象的认识，并在纯净的阳光中默观它们；我们自身纯洁，且如今不再被那物所分别——那物，当携带着它时，我们称为身体，如同牡蛎被绑在其壳上一般。\par

毕达哥拉斯学派称天为\OriginalTerm{对地}{Antichthon}。\Person{耶肋米亚}说在这地上，\Quote{我必将在儿女中安置你们，将拣选之地赐给你们作为全能天主的产业；}\BibleRef{耶 3:19} 而那些承受它的人将统治大地。无数万万千千的例子涌上我心头，本可举出。但为了匀称，现在话语必须停止，以免我们印证悲剧家\Person{阿伽松}的话：——\par

\Quote{将我们的副业当作正业，\newline{}并将我们的正业当作副业。}\par

那么，正如我所想，已经清楚地表明了希腊人被主称为盗贼应如何被理解，我乐意地离开哲学家的教理。因为如果我们遍历他们的语录，我们会马上收集如此多的笔记，在表明整个希腊智慧源于蛮族哲学。然而，当我们收集希腊人中间流行的关于第一原理的意见时，我们将再次触及这个思考，按必要性所需。\par

但从刚才所说，它默默地落在我们身上，去思考能穿越其中波涛的人该如何阅读希腊书籍。因此\par

\Quote{拥有神圣理智财富的人有福了。}\par

正如\Person{恩培多克勒}所说的那样。\par

\Quote{但可悲的是，他关心关于众神的黑暗意见。}\par

他神圣地展示了知识与无知是幸福与痛苦的界限。\Quote{因为哲学家应当熟悉许多事物，}按照\Person{赫拉克利特}的话；并且确实必须\par

\Quote{他，寻求成为一个善者，在许多事上犯错。}\par

那么，从刚才所说，对我们来说现在清楚了，天主的恩惠是永恒的，并且从无始的原则，平等的自然正义达到了所有人，按照各个种族的值得——从未有过开始。因为天主并非始于成为主和善，祂常是祂所是。祂也永不停止行善，即使祂使万有终结。我们每个人都按祂的意愿分受祂的恩惠。因为选民的区别是由值得灵魂的选择和实践参与介入而造成的。\par

那么，就这样，让我们符合真正哲学的第五卷诺斯替笔记杂录到此结束。\par

\SourceRecord{Translated by William Wilson.；Ante-Nicene Fathers；Vol. 2.；Edited by Alexander Roberts, James Donaldson, and A. Cleveland Coxe.；Buffalo, NY: Christian Literature Publishing Co.,；1885.；<http://www.newadvent.org/fathers/02105.htm>.}

% structure: p0001 source=h1 target=section reason=page-title

\section{\WorkTitle{杂记（卷六）}}

% structure: p0002 source=h2 target=subsection reason=relative-heading-rank; base=h2

\subsection{第一章 计划。}

第六卷及第七卷杂记，依照真哲学，尽可能描绘其中所传达的伦理论证，并展示了诺斯替者在生活中的样貌，接下来要向哲学家表明，他绝非如他们所想的那样不虔敬，而是唯独他真正虔敬，通过简明地展示诺斯替者的宗教形式，尽可能安全地将其写在一本参考书中。因为主吩咐说：\Quote{要为那存留到永生的食物劳碌。}\BibleRef{若 6:27}先知也说：\Quote{那撒种在各水边，使牛和驴践踏的，是有福的，}\BibleRef{依 32:20}［即］百姓，从法律和外邦人中，聚集在一个信仰里。\par

\Quote{软弱的人只吃蔬菜，}依照尊贵的宗徒。\BibleRef{罗 14:2}我们分为三卷的\WorkTitle{导师}已经展示了从童年状态开始的训练和培育，即从初步教导因着信德而成长的历程；并且在那些被登记在人数之列的人中，为接受诺斯替知识预备了已具有德性的灵魂。希腊人清楚地从我们将在这些篇章中所说的内容学到，他们在不虔敬地逼迫那爱天主之人的同时，自己也在行不虔敬；然后，随着笔记的进展，按照\WorkTitle{杂记}的风格，我们必须解决希腊人和外邦人关于主来临所提出的难题。\par

在草地上，花朵五彩缤纷地绽放；在园中，果树的种植并不按种类与其他种类分开。如果有人采集各式各样，编成了博学的文集、赫利孔山、蜂房和袍服，那么，用随意回忆起来的东西，既不在顺序上也不在表达上加以净化，而是故意散布，\WorkTitle{杂记}的形式就像草地一样混杂多样。既然如此，我的笔记将成为点燃的火花；对于有资格获得知识的人，如果他偶然遇到它们，通过努力进行研究将对他有益和有利。因为劳动不仅应当先于食物，而且更应当先于知识，对于那些借着\Quote{窄门狭路}——真正是主的道路——迈向永恒和蒙福救恩的人。\par

我们的知识和我们属灵的园圃就是救主自己；我们被移植到他里面，从旧生命被转移和移植到好土地中。移植有助于结果实。主，我们被移植到他里面，就是光与真知识。\par

现在，知识还有另一种双重意义上的说法：普遍称呼的、在所有人身上显现的知识（同样地，理解力和领悟力也如此），普遍地，关于个别对象的知识；不仅理性能力，非理性能力也同样分享这种知识，我绝不会称之为知识，因为通过感官掌握事物是自然的。但唯独被称为知识的东西，带有判断和理性的印记，在运用它时，将只有理性的认知，纯粹适用于思想的对象，并来自灵魂的纯粹活动。达味说：\Quote{他是一个善人，他怜悯}（那些因错误而毁灭的人）\Quote{并且借贷}（通过真理之言的交流），不是随意地，因为\Quote{他要在审判中分配他的话语：}以深远的考量\Quote{他施舍，他给予穷人。}\par

% structure: p0008 source=h2 target=subsection reason=relative-heading-rank; base=h2

\subsection{第二章 剽窃主题再续。希腊人互相剽窃}

在处理所提议的论题之前，我们必须先作为序言，补充第五卷结尾所缺少的内容。因为既然我们已经表明象征风格是古老的，不仅被我们的先知使用，也被大多数古代希腊人以及不少其他外邦野蛮人使用，那么就必须继续讲解入教者的奥秘。我推迟阐明这些内容，直到我们进展到驳斥希腊人关于第一原则的言论；因为我们将表明这些奥秘属于同一个思辨分支。并且，通过证明希腊思想的宣告被赐予我们的圣经中的真理所全面照亮——按照意思来理解——我们已经证明，且不说令人反感的话，真理的被盗传到了他们那里。\par

来吧，让我们引用希腊人自身作为见证来反对他们自己的剽窃。因为，既然他们互相偷窃，他们就证实了他们是窃贼；尽管违反他们的意愿，他们被揭露，秘密地将属于我们的真理据为己有，传给同族的人。因为如果他们彼此不肯罢手，他们几乎不会放过我们的作者。关于哲学教义，我什么也不会说，因为那些学派划分的作者本人，在书面中承认，为了不被指控忘恩负义，他们从\Person{苏格拉底}那里获得了最重要的教义。但是，在利用了几个最常被提及、在希腊人中享有声誉之人的少数证词，并揭露了他们剽窃的风格，从不同时期挑选了它们之后，我将转向接下来的内容。\par

因此，\Person{俄耳甫斯}作了这句诗：——\par

\Quote{因为没有什么比女人更无耻、更可鄙，}\par

\Person{荷马}明言：——\par

\Quote{因为没有什么比女人更可怕、更无耻。}\par

\Person{穆赛俄斯}曾写道：——\par

\Quote{因为技艺远胜于力量，}——\par

\Person{荷马}说：——\par

\Quote{靠技艺而非力量，伐木者远胜一筹。}\par

再者，\Person{穆赛俄斯}曾作诗行：——\par

\Quote{就像肥沃的田地长出叶子，\newline{}在梣树上，一些凋零，一些生长，\newline{}人类种族也是如此旋转着它的叶子，}——\par

\Person{荷马}转录：——\par

\Quote{一些叶子被风吹落在地。\newline{}萌芽的树木生出一些；在春天时节，\newline{}它们来临。如此，一代人兴起，一代人消逝。}\par

再者，\Person{荷马}曾说：——\par

\Quote{It is unholy to exult over dead men,}\par

\Person{Archilochus} and \Person{Cratinus} write, the former:—\par

\Quote{It is not noble at dead men to sneer;}\par

and \Person{Cratinus} in the \WorkTitle{Lacones}:—\par

\Quote{For men 'tis dreadful to exult\newline{}Much o'er the stalwart dead.}\par

Again, \Person{Archilochus}, transferring that Homeric line:—\par

\Quote{I erred, nor say I nay: instead of many} —\par

writes thus:—\par

\Quote{I erred, and this mischief has somehow seized another.}\par

As certainly also that line:—\par

\Quote{Even-handed war the slayer slays.}\par

He also, altering, has given forth thus:—\par

\Quote{I will do it.\newline{}For Mars to men in truth is evenhanded.}\par

Also, translating the following:—\par

\Quote{The issues of victory among men depend on the gods,}\par

he openly encourages youth, in the following iambic:—\par

\Quote{Victory's issues on the gods depend.}\par

Again, \Person{Homer} having said:—\par

\Quote{With feet unwashed sleeping on the ground,}\par

\Person{Euripides} writes in \WorkTitle{Erechtheus}:—\par

\Quote{Upon the plain spread with no couch they sleep,\newline{}Nor in the streams of water lave their feet.}\par

\Person{Archilochus} having likewise said:—\par

\Quote{But one with this and one with that\newline{}His heart delights,}—\par

in correspondence with the Homeric line:—\par

\Quote{For one in these deeds, one in those delights,} —\par

\Person{Euripides} says in \WorkTitle{Œneus}:—\par

\Quote{But one in these ways, one in those, has more delight.}\par

And I have heard \Person{Æschylus} saying:—\par

\Quote{He who is happy ought to stay at home;\newline{}There should he also stay, who speeds not well.}\par

And \Person{Euripides}, too, shouting the like on the stage:—\par

\Quote{Happy the man who, prosperous, stays at home.}\par

\Person{Menander}, too, on comedy, saying:—\par

\Quote{He ought at home to stay, and free remain,\newline{}Or be no longer rightly happy.}\par

Again, \Person{Theognis} having said:—\par

\Quote{The exile has no comrade dear and true,}—\par

\Person{Euripides} has written:—\par

\Quote{Far from the poor flies every friend.}\par

And \Person{Epicharmus}, saying:—\par

\Quote{Daughter, woe worth the day!\newline{}You who are old I marry to a youth;}\par

and adding:—\par

\Quote{For the young husband takes some other girl,\newline{}And for another husband longs the wife,}—\par

\Person{Euripides} writes:—\par

\Quote{'Tis bad to yoke an old wife to a youth;\newline{}For he desires to share another's bed,\newline{}And she, by him deserted, mischief plots.}\par

\Person{Euripides} having, besides, said in the \WorkTitle{Medea}:—\par

\Quote{For no good do a bad man's gifts,}—\par

\Person{Sophocles} in \WorkTitle{Ajax Flagellifer} utters this iambic:—\par

\Quote{For foes' gifts are no gifts, nor any boon.}\par

\Person{Solon} having written:—\par

\Quote{For surfeit insolence begets,\newline{}When store of wealth attends.}\par

\Person{Theognis} writes in the same way:—\par

\Quote{For surfeit insolence begets,\newline{}When store of wealth attends the bad.}\par

Whence also \Person{Thucydides}, in the \WorkTitle{Histories}, says: \Quote{Many men, to whom in a great degree, and in a short time, unlooked-for prosperity comes, are wont to turn to insolence.} And \Person{Philistus} likewise imitates the same sentiment, expressing himself thus: \Quote{And the many things which turn out prosperously to men, in accordance with reason, have an incredibly dangerous tendency to misfortune. For those who meet with unlooked success beyond their expectations, are for the most part wont to turn to insolence.} Again, \Person{Euripides} having written:—\par

\Quote{For children sprung of parents who have led\newline{}A hard and toilsome life, superior are;}\par

\Person{Critias} writes: \Quote{For I begin with a man's origin: how far the best and strongest in body will he be, if his father exercises himself, and eats in a hardy way, and subjects his body to toilsome labour; and if the mother of the future child be strong in body, and give herself exercise.}\par

Again, \Person{Homer} having said of the \Person{Hephæstus}-made shield:—\par

\Quote{Upon it earth and heaven and sea he made,\newline{}And Ocean's rivers' mighty strength portrayed,}\par

\Person{Pherecydes} of \Place{Syros} says:— \Quote{\Person{Zas} makes a cloak large and beautiful, and works on it earth and \Person{Ogenus}, and the palace of \Person{Ogenus}.}\par

And \Person{Homer} having said:—\par

\Quote{Shame, which greatly hurts a man or helps,} —\par

\Person{Euripides} writes in \WorkTitle{Erechtheus}:—\par

\Quote{Of shame I find it hard to judge;\newline{}'Tis needed. 'Tis at times a great mischief.}\par

Take, by way of parallel, such plagiarisms as the following, from those who flourished together, and were rivals of each other. From the \WorkTitle{Orestes} of \Person{Euripides}:—\par

\Quote{Dear charm of sleep, aid in disease.}\par

From the \WorkTitle{Eriphyle} of \Person{Sophocles}:—\par

\Quote{Hie you to sleep, healer of that disease.}\par

And from the \WorkTitle{Antigone} of \Person{Sophocles}:—\par

\Quote{Bastardy is opprobrious in name; but the nature is equal;}\par

And from the \WorkTitle{Aleuades} of \Person{Sophocles}:—\par

\Quote{Each good thing has its nature equal.}\par

Again, in the \WorkTitle{Ctimenus} of \Person{Euripides}:—\par

\Quote{For him who toils, God helps;}\par

And in the \WorkTitle{Minos} of \Person{Sophocles};\par

\Quote{To those who act not, fortune is no ally;}\par

And from the \WorkTitle{Alexander} of \Person{Euripides}:—\par

\Quote{But time will show; and learning, by that test,\newline{}I shall know whether you are good or bad;}\par

And from the \WorkTitle{Hipponos} of \Person{Sophocles}:—\par

\Quote{Besides, conceal nothing; since Time,\newline{}That sees all, hears all, all things will unfold.}\par

But let us similarly run over the following; for \Person{Eumelus} having composed the line,\par

\Quote{Of Memory and Olympian Zeus the daughters nine,}\par

\Person{Solon} thus begins the elegy:—\par

\Quote{Of Memory and Olympian Zeus the children bright.}\par

Again, \Person{Euripides}, paraphrasing the Homeric line:—\par

\Quote{What, whence are you? Your city and your parents, where?}\par

employs the following iambics in \WorkTitle{Ægeus}:—\par

\Quote{What country shall we say that you have left\newline{}To roam in exile, what your land— the bound\newline{}Of your own native soil? Who you begot?\newline{}And of what father do you call yourself the son?}\par

And what? \Person{Theognis} having said:—\par

\Quote{Wine largely drunk is bad; but if one use\newline{}It with discretion, 'tis not bad, but good,}—\par

does not \Person{Panyasis} write?\par

\Quote{Above the gods' best gift to men ranks wine,\newline{}In measure drunk; but in excess the worst.}\par

\Person{Hesiod}, too, saying:—\par

\Quote{But for the fire to you I'll give a plague,\newline{}For all men to delight themselves withal,}—\par

\Person{Euripides} writes:—\par

\Quote{And for the fire\newline{}Another fire greater and unconquerable,\newline{}Sprung up in the shape of women}\par

And in addition, \Person{Homer}, saying:—\par

\Quote{There is no satiating the greedy paunch,\newline{}Baneful, which many plagues has caused to men.}\par

\Person{Euripides} says:—\par

\Quote{Dire need and baneful paunch me overcome;\newline{}From which all evils come.}\par

此外，喜剧诗人\Person{卡里亚斯}曾写道：\par

\Quote{With madmen, all men must be mad, they say,}—\par

\Person{米南德}在\WorkTitle{Poloumenoi}中表达了类似的意思，说道：\par

\Quote{The presence of wisdom is not always suitable:\newline{}One sometimes must with others play the fool.}\par

特奥斯的\Person{安提马科斯}曾说道：\par

\Quote{From gifts, to mortals many ills arise,}—\par

\Person{奥吉亚斯}写下了这样的诗句：\par

\Quote{For gifts men's mind and acts deceive.}\par

\Person{赫西俄德}曾说道：\par

\Quote{Than a good wife, no man a better thing\newline{}Ere gained; than a bad wife, a worse,}—\par

\Person{西摩尼得斯}说道：\par

\Quote{A better prize than a good wife no man\newline{}Ere gained, than a bad one nought worse.}\par

同样，\Person{埃庇卡摩斯}曾说道：\par

\Quote{As destined long to live, and yet not long,\newline{}Think of yourself.}—\par

\Person{欧里庇得斯}写道：\par

\Quote{Why? Seeing the wealth we have uncertain is,\newline{}Why don't we live as free from care, as pleasant\newline{}As we may?}\par

同样，喜剧诗人\Person{狄菲罗斯}曾说道：\par

\Quote{The life of men is prone to change,}—\par

\Person{波西狄波斯}说道：\par

\Quote{No man of mortal mould his life has passed\newline{}From suffering free. Nor to the end again\newline{}Has continued prosperous.}\par

\Person{柏拉图}在描写人是一种易变的存在时，也对你说了同样的话。同样，\Person{欧里庇得斯}曾说道：\par

\Quote{Oh life to mortal men of trouble full,\newline{}How slippery in everything are you!\newline{}Now you grow, and now you decay away.\newline{}And there is set no limit, no, not one,\newline{}For mortals of their course to make an end,\newline{}Except when Death's remorseless final end\newline{}Comes, sent from Zeus,}—\par

\Person{狄菲罗斯}写道：\par

\Quote{There is no life which has not its own ills,\newline{}Pains, cares, thefts, and anxieties, disease;\newline{}And Death, as a physician, coming, gives\newline{}Rest to their victims in his quiet sleep.}\par

此外，\Person{欧里庇得斯}曾说道：\par

\Quote{Many are fortune's shapes,\newline{}And many things contrary to expectation the gods perform,}—\par

悲剧诗人\Person{忒奥得克忒斯}同样写道：\par

\Quote{The instability of mortals' fates.}\par

\Person{巴库里得斯}曾说道：\par

\Quote{To few alone of mortals is it given\newline{}To reach hoary age, being prosperous all the while,\newline{}And not meet with calamities,}—\par

喜剧诗人\Person{摩斯基翁}写道：\par

\Quote{But he of all men is most blest,\newline{}Who leads throughout an equal life.}\par

你会发现，\Person{泰奥格尼斯}曾说道：\par

\Quote{For no advantage to a man grown old\newline{}A young wife is, who will not, as a ship\newline{}The helm, obey,}—\par

喜剧诗人\Person{阿里斯托芬}写道：\par

\Quote{An old man to a young wife suits but ill.}\par

因为\Person{阿那克里翁}曾写道：\par

\Quote{Luxurious love I sing,\newline{}With flowery garlands graced,\newline{}He is of gods the king,\newline{}He mortal men subdues,}—\par

\Person{欧里庇得斯}写道：\par

\Quote{For love not only men attacks,\newline{}And women; but disturbs\newline{}The souls of gods above, and to the sea\newline{}Descends.}\par

但是，为了避免冗长，在我们急于展示希腊人在措辞和教义上的剽窃倾向时，请允许我们引用厄勒亚的智者\Person{希庇亚斯}在论述此问题时明确给出的证词，他如此说道：\Quote{Of these things some perchance are said by Orpheus, some briefly by Musæus; some in one place, others in other places; some by Hesiod, some by Homer, some by the rest of the poets; and some in prose compositions, some by Greeks, some by Barbarians. And I from all these, placing together the things of most importance and of kindred character, will make the present discourse new and varied.}\par

为了让我们看到哲学、历史乃至修辞学也免不了同样的责难，有必要从中略举数例。克罗顿的\Person{阿尔克迈翁}曾说道：\Quote{It is easier to guard against a man who is an enemy than a friend,} \Person{索福克勒斯}在\WorkTitle{安提戈涅}中写道：\par

\Quote{For what sore more grievous than a bad friend?}\par

\Person{色诺芬}说道：\Quote{No man can injure enemies in any way other than by appearing to be a friend.}\par

\Person{欧里庇得斯}在\WorkTitle{忒勒福斯}中曾说道：\par

\Quote{Shall we Greeks be slaves to Barbarians?}—\par

\Person{特拉西马库斯}在致拉里萨人的演说中说道：\Quote{Shall we be slaves to Archelaus— Greeks to a Barbarian?}\par

\Person{俄耳甫斯}曾说道：\par

\Quote{Water is the change for soul, and death for water;\newline{}From water is earth, and what comes from earth is again water,\newline{}And from that, soul, which changes the whole ether;}\par

\Person{赫拉克利特}把这些诗句中的表述整理起来，如此写道：\par

\Quote{It is death for souls to become water, and death for\newline{}water to become earth; and from earth comes water,\newline{}and from water soul.}\par

毕达哥拉斯学派的\Person{阿塔马斯}曾说道：\Quote{Thus was produced the beginning of the universe; and there are four roots— fire, water, air, earth: for from these is the origination of what is produced,}—阿格里根特的\Person{恩培多克勒}写道：\par

\Quote{The four roots of all things first do you hear—\newline{}Fire, water, earth, and ether's boundless height:\newline{}For of these all that was, is, shall be, comes.}\par

\Person{柏拉图}曾说道：\Quote{Wherefore also the gods, knowing men, release sooner from life those they value most,} \Person{米南德}写道：\par

\Quote{Whom the gods love, dies young.}\par

\Person{欧里庇得斯}在\WorkTitle{俄诺玛俄斯}中曾写道：\par

\Quote{We judge of things obscure from what we see;}\par

而在\WorkTitle{福尼克斯}中：\par

\Quote{By signs the obscure is fairly grasped,}—\par

\Person{许珀里得斯}说道：\Quote{But we must investigate things unseen by learning from signs and probabilities.} \Person{伊索克拉底}曾说道：\Quote{We must conjecture the future by the past,} \Person{安多基德斯}也不避讳地说道：\Quote{For we must make use of what has happened previously as signs in reference to what is to be.} 此外，\Person{泰奥格尼斯}曾说道：\par

\Quote{The evil of counterfeit silver and gold is not intolerable,\newline{}O Cyrnus, and to a wise man is not difficult of detection;\newline{}But if the mind of a friend is hidden in his breast,\newline{}If he is false, and has a treacherous heart within,\newline{}This is the basest thing for mortals, caused by God,\newline{}And of all things the hardest to detect,}—\par

\Person{欧里庇得斯}写道：\par

\Quote{Oh Zeus, why have you given to men clear tests\newline{}Of spurious gold, while on the body grows\newline{}No mark sufficing to discover clear\newline{}The wicked man?}\par

\Person{许珀里得斯}本人也说道：\Quote{There is no feature of the mind impressed on the countenance of men.}\par

再者，\Person{斯塔西努斯}写下如下诗句：——\par

\Quote{愚人，杀了父亲，留下儿女，}——\par

\Person{色诺芬}说：\Quote{因为我似乎做了类似的事，如同一个人杀了父亲却饶了儿女。} 而\Person{索福克勒斯}在\WorkTitle{安提戈涅}中写道：——\par

\Quote{母亲和父亲如今已在冥府，\newline{}再无兄弟能够为我降生。}——\par

\Person{希罗多德}说：\Quote{母亲和父亲已不在，我不会再有另一个兄弟。} 此外，\Person{泰奥蓬波斯}写道：——\par

\Quote{老人真可谓两次做孩童；}\par

在他之前，\Person{索福克勒斯}在\WorkTitle{佩琉斯}中写道：——\par

\Quote{我是埃阿科斯之子佩琉斯，唯一的管家，\newline{}如今他年迈，我引导他，重新训练他，\newline{}因为老人再次成为孩童，}——\par

\Person{安提丰}演说家说：\Quote{对老人的照料如同对孩童的照料。} 哲学家\Person{柏拉图}也说：\Quote{老人，正如似乎，将两次成为孩童。} 此外，\Person{修昔底德}曾说：\Quote{在马拉松我们独当一面，}——\Person{德摩斯梯尼}说：\Quote{由那些在马拉松独当一面的人。} 我也不会忽略以下。\Person{克拉提努斯}在\WorkTitle{皮提尼}中曾说：——\par

\Quote{准备工作或许你们知道。}\par

\Person{安多基德斯}演说家说：\Quote{准备工作，陪审员们，以及敌人的热忱，几乎你们所有人都知道。} 同样，\Person{尼基亚斯}在关于押金的演说中，反对\Person{吕西亚斯}，说：\Quote{对方当事人的准备工作和热忱，你们看到了，哦陪审员们。} 在他之后，\Person{埃斯基涅斯}说：\Quote{你们看到了准备工作，哦雅典人，以及战线。} 再者，\Person{德摩斯梯尼}曾说：\Quote{在这场竞赛中投入了怎样的热忱和怎样的拉票，哦雅典人，我想几乎你们所有人都知道；} 而\Person{菲利努斯}类似地：\Quote{在这竞赛中发生了怎样的热忱、怎样的列阵，陪审员们，我想你们中无人不知。} 再者，\Person{伊索克拉底}曾说：\Quote{仿佛她与他的财富有亲，而非与他本人，} \Person{吕西亚斯}在\WorkTitle{俄耳甫斯}中说：\Quote{他显然不是与人有亲，而是与钱有亲。} 此后，\Person{荷马}也写道：——\par

\Quote{朋友，如果在这场战争中，通过逃跑，\newline{}我们就能免除衰老和死亡，\newline{}我自己就不会冲在最前线，\newline{}也不会派你去光荣的战斗；\newline{}但如今——因为死亡的无数命运\newline{}无论如何都缠身，人既无法逃脱\newline{}也躲不开——来吧。我们要么给某人带来\newline{}荣耀，要么某人给我们带来，}——\par

\Person{泰奥蓬波斯}写道：\Quote{因为如果避开当前的危险，我们就能在安全中度过余生，那么表现出贪生怕死并不奇怪。但如今，生命遭遇如此多致命的事件，以致战死似乎更可取。} 还有呢？智者\Person{奇利}说出了格言：\Quote{给人作保，灾祸就在眼前，} 难道\Person{埃皮卡莫斯}没有用其他话表达同样的意思吗？他说：\Quote{担保是灾祸之女，损失是担保之女？} 此外，医生\Person{希波克拉底}写道：\Quote{你必须注意时间、地点、年龄和疾病，} \Person{欧里庇得斯}在\Emph{六音步诗}中说：——\par

\Quote{那些要良好实践医术的人，\newline{}必须研究人们的生活方式，注意\newline{}土壤，并这样考虑疾病。}\par

\Person{荷马}再次写道：——\par

\Quote{我说凡人无法逃脱命运，}——\par

\Person{阿尔基努斯}说：\Quote{所有凡人迟早都注定要死；} \Person{德摩斯梯尼}：\Quote{对所有人而言，死亡是生命的终结，即使有人把自己关在笼子里。}\par

再者，\Person{希罗多德}在关于斯巴达人格劳库斯的演说中说，皮提亚曾言：\Quote{在神身上，说和做是等同的，} \Person{阿里斯托芬}说：——\par

\Quote{因为想和做是等同的。}\par

在他之前，\Person{埃利亚的巴门尼德}说：——\par

\Quote{因为思想和存在是同一的。}\par

\Person{柏拉图}曾说：\Quote{并且我们将说明，也许并不荒谬，爱的起点是视觉；希望减弱激情，记忆滋养它，交往保存它；} 喜剧诗人\Person{菲勒蒙}不也写道：——\par

\Quote{先是众人看见，然后倾慕；\newline{}继而凝视，继而产生希望；\newline{}爱情就这样产生？}\par

再者，\Person{德摩斯梯尼}曾说：\Quote{因为对我们所有人而言，死亡是一笔债务，}等等，\Person{法诺克勒斯}在\WorkTitle{爱}或\WorkTitle{美}中写道：——\par

\Quote{但从命运不间断的线中逃脱\newline{}对地上养育之人是不可能的。}\par

你还会发现，\Person{柏拉图}曾说：\Quote{每株植物的萌芽，若按自身本性的德性获得良好的开端，对引导其适切的终点最为有力；} 这位历史学家写道：\Quote{此外，野生植物在度过早期生长期后，要变成栽培植物是不自然的；} 以及\Person{恩培多克勒}的如下诗句：——\par

\Quote{因为我已经曾是男孩和女孩，\newline{}以及灌木、鸟和海中哑鱼，}——\par

\Person{欧里庇得斯}在\WorkTitle{克律西波斯}中转录：——\par

\Quote{但存在的万物中\newline{}无一物死亡；而是分解开来，\newline{}从另一物中分解，\newline{}显现另一种形式。}\par

\Person{柏拉图}曾在\WorkTitle{理想国}中说，女人是共有的，\Person{欧里庇得斯}在\WorkTitle{普罗忒西拉奥斯}中写道：——\par

\Quote{所以，女人的床是共有的。}\par

再者，\Person{欧里庇得斯}曾写道：——\par

\Quote{对有节制者而言，满足就够了}——\par

\Person{伊壁鸠鲁}明确地说：\Quote{满足是最大的财富。}\par

再者，\Person{阿里斯托芬}曾写道：——\par

\Quote{你将安全地享受生活，行正义，\newline{}远离纷扰，免于恐惧而安居，}——\par

\Person{伊壁鸠鲁}说：\Quote{正义的最大果实是宁静。}\par

那么，就让这些希腊人剽窃思想的样本，如此这般，对能领会的人而言，足以作为清晰的例证。\par

他们不仅被查出剽窃和改写希腊人的思想和措辞，正如将要表明的那样；而且他们也将被证实拥有完全偷来的东西。因为他们完全偷窃他人的作品，将其当作自己的出版；如从穆赛俄斯窃取整本\WorkTitle{论特斯普罗提亚人}的库瑞涅人欧加蒙，从林都斯的庇西努斯\WorkTitle{赫拉克勒亚}的卡米罗斯人庇桑德，以及从萨摩斯的克勒奥斐卢斯\WorkTitle{奥卡利亚的陷落}的哈利卡纳苏斯人潘亚西斯。\par

你也会发现，伟大的诗人荷马从俄耳甫斯的\WorkTitle{狄俄尼索斯的消失}中逐字逐句地拿走了这些词句：——\par

\Quote{如同一个人培育一棵茂盛的橄榄幼苗。}\par

在\WorkTitle{神谱}中，俄耳甫斯论克洛诺斯说：——\par

\Quote{他躺着，粗壮的脖子歪向一边；\newline{}全能睡眠已将他攫住。}\par

荷马将这些转用到独眼巨人身上。赫西俄德论墨兰波斯写道：——\par

\Quote{欣然听闻，不朽者所分配予\newline{}凡人的，勇者与懦夫昭然有别；}\par

等等，逐字逐句取自诗人穆赛俄斯。\par

喜剧诗人阿里斯托芬在\WorkTitle{地母节妇女}第一部中，从克拉提努斯的\WorkTitle{焚烧}中转移了词句。喜剧诗人柏拉图（喜剧诗人）和阿里斯托芬在\WorkTitle{代达罗斯}中互相抄袭。阿里斯托芬之子阿拉罗斯所作的\WorkTitle{科卡路斯}，被喜剧诗人菲勒蒙改编，并制成名为\WorkTitle{假子}的喜剧。\par

编年史家欧墨鲁斯和阿库西拉俄斯将赫西俄德的内容改写成散文，当作自己的发表。勒昂提翁的戈尔吉亚和那克索斯的欧得摩斯这两位史家，从墨勒萨戈拉斯那里偷窃。此外，有普罗孔涅索斯的彼翁，他摘要并抄写了古代卡德摩斯、阿尔基洛科斯、亚里士多德、勒安得罗斯、赫兰尼科斯、赫卡泰俄斯、安德罗提翁和菲洛科洛斯的作品。墨伽拉的狄埃乌基达斯将他的论文开头从赫兰尼科斯的\WorkTitle{丢卡利翁}中转用。我略过以弗所的赫拉克利特，他从俄耳甫斯那里取了极多内容。\par

灵魂不死之说，柏拉图得自毕达哥拉斯，而毕达哥拉斯得自埃及人。许多柏拉图主义者著书证明，正如我们开头所说，斯多葛派和亚里士多德从柏拉图那里取得了他们教义中最主要的部分。伊壁鸠鲁也从德谟克利特那里窃取了他的主要教义。这些事就暂且如此吧。因为如果我打算详细讨论这个问题，揭露希腊人自私的剽窃，以及他们如何声称发现了他们从我们这里获得的最好的学说，那么生命将不足以完成。\par

% structure: p0234 source=h2 target=subsection reason=relative-heading-rank; base=h2

\subsection{第三章　希腊人剽窃希伯来圣书中的奇迹}

现在他们不仅被证实从蛮族人那里借用了学说，而且还把在我们记录中、由那些度圣洁生活、关注我们的人通过来自上天的神圣力量所行的奇迹，当作希腊神话中的奇事来讲述。我们要问他们：他们所讲的事是真还是假？他们不会说假的，因为故意编造谎言会使他们自责极其愚蠢，所以他们必然承认是真的。那么，由梅瑟和其他先知所行的奇迹，在他们看来又怎么会是不可信的呢？因为全能的天主眷顾所有人，用诫命、用威胁、用奇迹、用温和的许愿，使一些人得救。\par

且说，希腊人曾遭遇长期干旱，土地歉收。据说幸存者因饥荒而到德尔斐恳求皮提亚女祭司，问他们如何才能摆脱这场灾难。女祭司宣告：解除困境的唯一办法是借助埃阿科斯的祈祷。他们说服了埃阿科斯，他登上希腊山丘，向天伸出洁净的双手，呼求共同的天主，恳求祂怜悯荒芜的希腊。他祈祷时，雷声异常响起，整个周围的天空乌云密布。暴雨倾盆而下，充满了整个地区。结果是通过埃阿科斯祈祷的耕耘，带来了丰饶的丰收。\par

经上记载：\Quote{撒慕尔呼求了上主，} \Quote{上主就发出声音，并在收割的日子降下雨水。} \BibleRef{撒上 11:18} 你看见没有？那\Quote{降雨给义人和不义的人} \BibleRef{玛 5:45} 的，是同一的天主。我们的全部圣经充满了这样的例子：天主因义人的祈祷，垂听并成就他们的每一个请求。\par

再者，希腊人讲述说，有一次夏季风没有吹起，阿里斯泰俄斯在凯奥斯向伊斯米亚的宙斯献祭。那时农作物被灼热烤焦，因往常使土地润泽的风没有吹拂，造成巨大的破坏，而他轻易地召回了风。\par

在薛西斯远征希腊时，德尔斐的女祭司回答说：——\par

\Quote{德尔斐人啊，向风祈祷，这会更有利。}\par

——于是他们筑起祭坛，向风献祭，便得到了风的帮助。因为强风猛烈吹袭塞皮阿斯角，粉碎了波斯远征的整个装备。阿格里真托的恩培多克勒被称为\Quote{止风者}。据说有一次，从阿格里真托的山上吹来一阵强烈而有害的风，导致居民患病，妻子不育，他使风停止了。因此他本人写道：——\par

\Quote{你要使不知疲倦的风的力量平息，\newline{}它们冲向大地，毁坏凡人的庄稼，\newline{}并按你意愿带回复仇的狂风。}\par

他们说，有些人跟随他，这些人是使用占卜的，有些是长期被重病折磨的。那么，他们显然相信了我们的圣经中所记载的医治、标记和奇事。因为如果某些能力能移动风并分配降雨，就让他们听听圣咏作者的话：\BibleQuote{万军的上主，你的居所是多么可爱！} 这是上主，是众掌权者、宰制者和统治者的主，梅瑟所谈论的那位；好使我们能与祂同在。\BibleQuote{你要割损你坚硬的心，不要再硬着颈项。因为祂是万主之主，万神之神，伟大的天主，大能者，} \BibleRef{申 10:16} 等等。依撒意亚也说：\BibleQuote{请你们举目向上，看看谁创造了这一切？} \BibleRef{依 40:26}\par

有些人说，瘟疫、冰雹、暴风雨以及类似的事，不仅由于物质紊乱而发生，也由于魔鬼和恶天使的愤怒。例如，他们说，在\Place{克勒奥内}的术士们，观察天空的现象，当云层即将降下冰雹时，他们通过咒语和祭献来避免愤怒的威胁。若有时缺少动物，他们就用自己手指的血来献祭。女先知\Person{狄欧提玛}，在雅典人于瘟疫之前献祭时，使瘟疫延迟了十年。\Person{克里特的厄庇墨尼得斯}的祭献，使波斯战争推迟了同样的时间。而人们认为，称这些神灵为神明或是天使，并没有什么区别。那些精通雕像祝圣的人，在许多庙宇中为死者修建了坟墓，将这些灵魂称为\OriginalTerm{邪灵}{Dæmons}，并教导人们敬拜它们；因为这些灵魂因纯洁的生活，藉天主的预知，获得了在环绕大地的空间中游走以服务于人的能力。因为他们知道有些灵魂本质上是保持在身体内的。但关于这些，随着工作的进展，在\WorkTitle{关于天使的论著}中，我们将进行讨论。\par

\Person{德谟克利特}，通过观察天体现象预测许多事情，被称为\Quote{智慧}（\OriginalTerm{智慧}{\GreekText{Σοφία}}）。当他受到他的兄弟\Person{达玛苏斯}的热诚接待时，他根据某些星辰预测将有大雨。因此，有些被他说服的人，就收聚了他们的庄稼；因为当时是夏天，庄稼还在打谷场上。但其他人却失去了一切，因为意外的大暴雨突然降临。\par

那么，希腊人怎能再不相信\Place{西乃山}上那神圣的显现呢？那时火燃烧着，却不烧毁山上生长的任何东西；号角声发出，不需乐器吹奏。因为那被称为天主降临山上的事，乃是天主能力的来临，遍及整个世界，并宣告\BibleQuote{那不可接近的光} \BibleRef{弟前 6:16}。\par

因为按照圣经，这就是寓意。然而，正如\Person{阿黎斯托步罗}所说，火被看见了，当时全体群众——不少于一百万人，还不算未成年人——聚集在山周围，山麓周长不少于五天的路程。在整个显现的地方，他们看见燃烧的火，仿佛在四周安营；因此，降临不是地方性的。因为天主无所不在。\par

编撰故事的人说，在\Place{不列颠}岛上，有一座山，山下有一个洞穴，山顶上有一个裂缝；因此，当风进入洞穴，冲进裂缝的深处时，会听到一种声音，好像铙钹在音乐般地撞击。在树林中，当树叶被一阵突如其来的风刮动时，常常会发出一种类似鸟鸣的声音。\par

那些编写\WorkTitle{波斯史}的人也叙述说，在高地，在术士的国度里，有三座山坐落在一片广阔的平原上，那些经过该地的人，来到第一座山时，听到一种混乱的声音，好像几万人在呐喊，如同打仗的阵势；到了中间那座山，他们听到更响亮、更清晰的喧嚷；最后听到人们唱着凯歌，仿佛胜利者。在我看来，整个声音的原因是地形的平坦和多洞穴的性质；空气进入后，被弹回并汇聚到同一点，发出相当有力的声音。就让这些如此吧。但全能的天主，即使没有媒介，也能通过耳朵产生声音和异象，表明祂的伟大具有超越常规的自然秩序，以便使至今不信的灵魂皈依，并接受所赐的诫命。但既然有云和高的山，怎么不可能听到不同的声音，风因主动因而动？因此先知也说：\BibleQuote{你们听见了言语的声音，却没有看见什么形象。} \BibleRef{申 4:12} 你看到，主的圣言——没有形状的圣言，圣言的能力，主的光明的圣言，从天上、从上而来的真理，进入教会的集会，是通过光明的直接事工而运作的。\par

% structure: p0250 source=h2 target=subsection reason=relative-heading-rank; base=h2

\subsection{第四章 希腊人从埃及和印度的赤身智者中汲取了许多他们的哲学信条}

我们将发现另一项证实的确证，即优秀的哲学家们从我们这里采纳了他们最卓越的教义，却夸耀说各派系的某些教义是从其他蛮族人，主要是从埃及人那里采集来的——既有其他教义，也有灵魂转生这一特别教义。因为埃及人追求他们自己的哲学。这主要体现在他们的神圣礼仪中。首先出场的是歌咏者，手持音乐标志之一。因为他们说，他必须学习两本赫尔墨斯之书，其中一本包含诸神的颂歌，第二本关于君王生活的规章。歌咏者之后出场的是占星者，手持时辰仪和棕榈枝，这些是占星术的象征。他必须将四本赫尔墨斯占星书经常挂在嘴边。其中一本是关于可见恒星之秩序，另一本关于太阳月亮的合相与发光现象，其余关于它们的升起。按顺序，随后出场的是神圣书记，头上有翅膀，手中拿着书和尺子，其中有墨水和书写的芦苇笔。他必须熟悉所谓的象形文字，并通晓宇宙志和地理学、太阳月亮的方位、以及五大行星；此外还有埃及的描述和尼罗河地图；以及祭司装备的描述和祝圣给他们的场所，还有度量衡和神圣礼仪中使用的物品。接着，袍服执事紧随前面的人，手持正义肘尺和奠酒杯。他通晓所有被称为\OriginalTerm{培迪提克}{Pædeutic}（训练相关）和\OriginalTerm{莫斯科普提克}{Moschophatic}（祭献相关）的知识。还有十本书涉及他们对诸神的敬拜，包括埃及的崇拜；如关于祭献、初果、颂歌、祈祷、游行、节庆等。最后行的是先知，手臂中公开抱着水瓶；其后跟着携带饼食的人。他作为圣殿的管理者，学习那十本被称为\Quote{\OriginalTerm{希耶拉提克}{Hieratic}}的书；它们包含所有关于法律、诸神以及司铎的全部训练。因为在埃及人中，先知也负责收入的分配。因此有四十二本不可或缺的赫尔墨斯之书；其中三十六本包含埃及人的全部哲学，由上述人物学习；另外六本是医学书，由\OriginalTerm{持像者}{Pastophoroi}学习——论及身体结构、疾病、器械、药物、眼科、最后关于妇女。简而言之，埃及人的习俗就是如此。\par

印度人的哲学也受到赞誉。\Person{马其顿的亚历山大}选出了十位印度\OriginalTerm{裸智派}{Gymnosophists}中看起来最出色和最善辩的，向他们提出难题，并威胁要处死那些回答不当者；他命令其中最年长的一人来裁决。\par

于是第一个人被问及他是否认为活人比死人更多，他说：活人；因为死人不存在。第二个人被问及海洋还是陆地养育了更大的动物，他说：陆地；因为海洋是陆地的一部分。第三个人被问哪种动物最狡猾？答道：迄今为止尚未认识的那个——人。第四个人被问及他们为什么使他们的王子\Person{萨巴}造反，答道：因为他们希望他活得好，而不是死得惨。第五个人被问及他认为白天还是黑夜在先，他说：一天。因为令人困惑的问题必须有令人困惑的答案。第六个人被问及：人怎样才会最受爱戴？答道：最为强大，这样他就不会胆怯。第七个人被问及一个人怎样才能成为神？答道：如果他做人所不能做的事。第八个人被问及生命与死亡哪个更强？答道：生命，因为它承载如此多的苦难。第九个人被问及人活到什么时候为好？答道：直到他认为死不比生更好为止。当\Person{亚历山大}命令第十个人发言（因为他是裁判）时，他说：\Quote{一人说得比另一人更糟。}当\Person{亚历山大}说，那么，你给出这样的判决，难道不该先死吗？他答道：大王，您说过要先杀死第一个回答最差的人，但您又如何证明自己言行一致呢？\par

而希腊人被称为各种著作的剽窃者，我认为已有大量证据充分证明了这一点。\par

% structure: p0255 source=h2 target=subsection reason=relative-heading-rank; base=h2

\subsection{第5章 希腊人对真天主有一些认识}

而希腊人中最有声望的人认识天主，不是通过确切的知识，而是通过间接的表达，\Person{伯多禄}在\WorkTitle{宣讲}中说：\Quote{你们要知道，只有一位天主，祂造了万物的起始，并掌握终结的权力；祂是不可见的，却看见万物；祂不能被包含，却包含万物；祂一无所缺，万物都需要祂，并且依靠祂而存在；祂是不可理解的、永恒的、非受造的，祂藉着\InnerQuote{祂能力的圣言}造就了万物，也就是说，按照\OriginalTerm{诺斯替}{gnostic}经文，就是祂的圣子。}\par

接着他补充道：\Quote{不要像希腊人那样敬拜这位天主}——清楚表明，希腊人中的卓越者与我们敬拜同一天主，但他们并未通过圆满的知识学会圣子所传授的真理。\Quote{所以不要敬拜}，他并非说不要敬拜希腊人所敬拜的天主，而是说\Quote{不要像希腊人那样}——改变敬拜天主的方式，而非宣告另一位天主。那么，\Quote{不要像希腊人那样}这句话是什么意思，伯多禄自己会解释，他补充道：\Quote{因为他们被无知所驱使，不认识天主（像我们按照圆满知识那样认识）；反而将天主赐予他们使用能力的东西——木石、铜铁、金银——物质——塑造成形，又将那些供使用和占有的奴隶竖立起来，敬拜它们。至于天主赐予他们作食物的——空中的飞鸟、海中的鱼、地上的爬虫、田野的四足牲畜连同走兽、鼬鼠、老鼠、猫、狗、猿猴，以及他们自己的食物——他们却当作祭品献给会死的人；将死物献给死者，如同献给神，对天主忘恩负义，借此否认祂的存在。}至于我们和希腊人认识同一天主，虽方式不同，他将这样推论：\Quote{也不要像犹太人那样敬拜；因为他们以为只有自己认识天主，其实并不认识，竟敬拜天使、总领天使、月亮和月朔。若月亮不显现，他们就不守安息日（即所谓的第一日）；也不守月朔、无酵节、节日和大日。}然后他对这个问题作了总结：\Quote{所以你们也要这样行，圣洁而正义地学习我们所传授给你们的；持守它们，以一种新的方式，藉着基督敬拜天主。}因为我们在圣经中看到，正如主所说：\Quote{看，我要与你们订立新约，不像我与你们祖先在曷勒布山所订立的那约。}\BibleRef{耶 31:31} 祂与我们订立了新约；因为属于希腊人和犹太人的是旧的。但我们以一种新的方式，在第三种形式中敬拜祂的，是基督徒。因为，我想，他清楚地表明：同一的独一天主，被希腊人以外邦人的方式认识，被犹太人以犹太人的方式认识，而被我们以一种新的、属神的方式认识。\par

此外，他还表明，同一位天主既赐予了两种盟约，也将希腊哲学赐予了希腊人，全能者藉此在希腊人中得到光荣。这一点是清楚的。因此，来自希腊的教化和来自法律的教化，一同被聚集到得救之人的同一民族中——那些接受信德的人：并非这三个民族被时间分隔，以致有人会推测有三种本性，而是在同一位主的不同的盟约中，受同一位主的话语所训练。因为，正如天主愿意藉着赐予犹太人先知来拯救他们，同样也藉着兴起他们自己语言中的先知（按他们能领受天主恩惠的能力），祂将最卓越的希腊人从普通人群中分别出来。此外，宗徒保禄在\WorkTitle{伯多禄的宣讲}中也会说明，他说：\Quote{也取用希腊书籍，读读西彼耳，看如何表明天主是独一的，以及未来如何被预示。再取用希斯达斯培，读一读，你将发现更清晰、更明白地描述了天主子，以及多少君王将纠集军队攻击基督，恼恨祂和那些持守祂名的人，以及祂的忠信者、祂的忍耐和祂的来临。}然后他一言以蔽之问我们：\Quote{世界和世上的一切属于谁？岂不属于天主吗？}因此，伯多禄说，主对宗徒们说：\Quote{那时，如果以色列人中有人愿意悔改，并因我的名而信从天主，他的罪将得以赦免，过了十二年之后。你们要出去到世界中去，免得有人说：我们没有听见。}\par

% structure: p0259 source=h2 target=subsection reason=relative-heading-rank; base=h2

\subsection{第六章 福音曾在阴间向犹太人和外邦人宣讲。}

但正如福音的宣讲现在按适当的时候来到，同样，法律和先知也是在适当的时候赐予异邦人的，哲学赐予希腊人，以预备他们的耳朵接受福音。拯救以色列的主说：\Quote{因此，在悦纳的时候，我应允了你；在救恩的日子，我帮助了你。我立了你作为万民的盟约，使你复兴大地，承受荒芜之地为产业；对被囚的说：出来吧！对在黑暗中的人说：显现吧！}因为如果\Quote{被囚的}是犹太人，主曾对他们说：\Quote{你们这些愿意的，从你们的捆绑中出来吧！}——意指自愿被捆绑的人，他们背负着人所吩咐的\Quote{难挑的重担}\BibleRef{玛 23:4}——那么显然，\Quote{在黑暗中的人}就是那些灵魂的宰制能力被埋葬在偶像崇拜中的人。\par

因为对于按法律成为义人的人来说，信德是欠缺的。因此主在治愈他们时也说：\Quote{你的信德救了你。}但对于按哲学成为义人的人来说，不仅需要对主的信德，还需要放弃偶像崇拜。一旦真理显现，他们也立即懊悔从前的行为。\par

因此主向阴间的人宣讲福音。相应地，经上记载：\Quote{阴间对毁灭说：我们没有看见祂的形状，却听见了祂的声音。}这并非指那个地方本身，如上面的话所说，听见了声音，而是指那些被安置在阴间、自甘毁灭的人，如同那些自愿从船上跳入海中的人。他们，就是那些听见天主能力和声音的人。因为，哪个有理智的人会设想义人的灵魂和罪人的灵魂处于同一惩罚中，从而指控天主眷顾不公呢？\par

然而，圣经岂不是显示主曾向洪水灭亡的人，更确切地说，向被锁链捆住、被囚禁在\Quote{看守和警卫}中的人传了福音？\BibleRef{伯前 3:19}此外，在\WorkTitle{杂文集}第二卷中已经证明，宗徒们效法主，向阴府中的人传了福音。因为我认为，正如在此世一样，在那里，最好的门徒也应当效法老师；这样，主使那些属于希伯来人的人悔改，而宗徒们则使外邦人悔改，即那些按法律和哲学生活在正义中、却未能以完美方式、而是以有罪方式结束生命的人。因为合乎天主眷顾的安排是：那些在正义上更有价值、生活卓越的人，在悔改自己的过犯之后（虽然他们身处另一处所，但明确属于全能天主子民之列），应各按各人的认识而得救。\par

而且，我认为，救主也施展祂的大能，因为拯救是祂的工作；因此祂也借着福音宣讲，吸引那些无论身处何处、愿意相信祂的人归向救恩。那么，如果主降到阴府，除了传福音之外别无目的，正如祂确实降下了一样，那么祂要么向所有人传福音，要么只向希伯来人。如果向所有人，那么所有相信的人都将得救，即使他们是外邦人，只要在那里宣认信仰；因为天主的惩罚是拯救性的、训诫性的，引人皈依，祂宁愿罪人悔改，而不愿他死亡；尤其因为灵魂虽然被私欲所蒙蔽，但当脱离身体后，由于不再被卑劣的肉身所阻碍，便能更清楚地感知。\par

那么，如果祂只向犹太人传福音（他们缺少对救主的认识和信德），那么显然，既然天主不偏待人，宗徒们也必须像在此世一样，在阴府向那些准备好皈依的外邦人传福音。\WorkTitle{牧者}说得很好：\Quote{他们于是同他们一起下到水里，又升上来。但这些是活着下去，活着升上来的；而那些已经安眠的人，是死了下去，却活着升上来的。}此外，福音书\BibleRef{玛 27:52}说：\Quote{许多长眠者的身体竟起来}——显然是被转移到了更好的状态。于是，通过救主的恩宠安排，发生了一场普遍的移动和转移。\par

那么，一个义人，作为义人，无论他属于法律还是属于希腊人，与另一个义人并无区别。因为天主不但是犹太人的主，也是所有人的主，更是那些认识祂的人的父亲。因为如果按法律生活得好就是生活，那么按理性在法律中生活也是生活；那些在法律之前生活正直的人，被归于信德，并被判断为义人——那么显然，那些在法律之外、由于声音的特殊性质而生活正直的人，虽然身处阴府和监牢\BibleRef{伯前 3:19}，一旦听到主的声音（无论是祂本人的声音，还是通过祂的宗徒发出的声音），便立刻回头并相信了。因为我们记起主是\Quote{天主的德能}\BibleRef{格前 1:24}，而德能绝不软弱。\par

因此，我认为已经证明：天主是良善的，主是有大能的，他们以正义和平等拯救所有转向祂的人，无论在此世还是别处。因为天主的主动德能不仅在这里，而是无处不在、永远工作。因此，在\WorkTitle{伯多禄宣讲}中，主在复活后对门徒们说：\Quote{我拣选了你们十二个门徒，认为你们配得上我，}主愿意他们成为宗徒，并判断他们是忠信的，派遣他们到世上去，到地上的人中间，使他们知道只有一位天主，清楚地预示将要借着对基督的信德发生的事：那些听见并相信的将得救；而那些听见后不信的，将作见证，不能以\Quote{我们没有听见}为借口。\par

那么，在阴府中岂不也实行了同样的恩宠安排，以至于在那里，所有灵魂一听到宣告，或者表示悔改，或者承认他们的惩罚是公义的，因为他们不信？若使那些在主来临之前离世的人（他们没有听到福音，并且由于相信或不相信而没有提供任何理由）要么得救，要么受惩罚，那将是极大的不公。因为让这些人未经审判就被定罪，而只有那些在主来临之后生活的人才能享有天主的正义，这是不合理的。但对所有理性的灵魂，从上头曾说：\Quote{你们中无论谁在无知中行事，没有清楚认识天主，若在醒悟之后悔改，他的一切罪都将被赦免。}\Quote{因为，看，}经上说：\Quote{我已将生命与死亡摆在你们面前，使你们可以拣选生命。}天主说祂摆出，而不是祂制造两者，是为了比较选择。在另一处圣经中祂说：\Quote{你们若听从我，愿意听从，就可以吃地上的美物。但你们若不听从我，不愿意听从，刀剑必吞灭你们：因为上主的口说了这些话。}\BibleRef{依 1:19}\par

再者，达味明言（或者更确切地说，是主在圣人的位格中，从世界奠基起，每个在不同时期因信德而得救并将要得救的人），说：\BibleQuote{我的心欢喜，我的舌头快乐，我的肉身也要安居在指望中。因为祢必不将我的灵魂撇在阴间，也不叫祢的圣者见朽坏。祢已将生命的道路指示我，必叫我因祢的临在充满喜乐。}因此，正如百姓在主眼中为宝贵，整个圣洁的百姓也是如此；连同犹太人，那从外邦人中归化的人，也是以归化者的名号被预言的。因为经上说得对：\BibleQuote{牛必与熊同食。}\BibleRef{依 11:7}因为犹太人以牛为象征，牛是负轭的洁净动物，按法律既分蹄又反刍。而外邦人以熊为象征，熊是不洁的野兽。这动物生出一团无定形的肉，仅靠舌头舔成兽的形状。因为那从外邦人中归化的人，藉着圣言从野兽般的生活被塑造成温良；一旦被驯服，就变得洁净，如牛一般。例如，先知说：\BibleQuote{野狗、鸵鸟和田野的走兽都必尊崇我。}\BibleRef{依 43:20}在不洁的动物之列，田野的走兽是众所周知的，即属于世界的；因为那些在信德上野性、在生活上污秽、未按法律之正义成义的人，被称为走兽。但藉着主的信德，他们从走兽转变为属天主的人，从改变的愿望进到事实。因为主劝诫一些人，对那些已经尝试的人，祂伸出援手，将他们拉起。\BibleQuote{因为主不看人的情面，也不重视尊贵；祂造了大人和小人，对一切人同样关怀。}\BibleRef{智 6:7}达味又说：\Quote{因为外邦人陷在自己所掘的坑中，他们的脚被自己隐藏的网罗缠住。}\Quote{但主作了穷人的避难所，在患难中作及时的帮助。}所以，那些在患难中的人，适时地听到了福音。因此经上说：\Quote{在列邦中宣扬祂的作为，}使他们不致受不公义的审判。\par

如果祂向肉身中的人传福音，使他们不致受不公义的审判，那么怎能设想祂没有因此向那些在祂来临之前已离世的人传福音呢？\Quote{因为公义的主喜爱公义，祂的面容垂顾正直。}\Quote{但喜爱邪恶的人，憎恨自己的灵魂。}\par

如果在洪水中，一切有罪的内身都灭亡了，惩罚加于他们是为纠正，我们必须首先相信：天主的旨意是训诲且慈惠的，祂拯救那些归向祂的人。再者，那更精微的实体——灵魂，绝不会因更粗糙的水元素而受伤害，因其精微单纯的本性，它是不可触的，被称为无形体。然而，因罪而变得粗重的一切，连同那与灵魂相争的肉性之神，都会被抛弃。\par

现在，瓦伦提诺——那些宣扬共融之人的领袖，在他的著作\WorkTitle{论朋友的交往}中写道：\Quote{所写的许多东西，虽然是在普通书籍中，却发现在天主的教会里写有。因为那些从心里发出的言论是虚妄的。因为写在心里的法律，就是那被爱者的人民——他们爱祂，也为祂所爱。}无论他将犹太人的著作还是哲学家的著作称为\Quote{普通书籍}，他都将真理视为公共的。而巴西里德的儿子兼弟子依西多禄，在\WorkTitle{先知帕尔乔释义}第一卷中也写道：\Quote{阿提卡人说，有某些事情曾通过伴随苏格拉底的一个灵魔暗示给他。亚里士多德说，所有人都有灵魔陪伴他们，在他们处于肉身期间——他将这一预言性的教导加以摄取，并移入自己的著作中，却未承认他从何处摘取了这一论述。}此外，在他的著作第二卷中，他这样写道：\Quote{不要有人认为，我们所说的只属于选民的那些话，曾被任何哲学家先说过。因为这不是他们的发现。他们从我们的先知那里借用之后，将其归之于他们眼中的智者。}在同一处又说：\Quote{因为在我看来，那些自称为哲学家的人，这样做是为了学习什么是长翅膀的橡树及其上五彩斑斓的披风，所有这些都被斐瑞居德斯用作神学寓言，取自查姆的预言。}\par

% structure: p0273 source=h2 target=subsection reason=relative-heading-rank; base=h2

\subsection{第七章 真正的哲学是什么及其名称由何处来}

正如我们早已指出的，我们提出的主题并非各派中的学问，而是真正的哲学，即严格系统的智慧，它提供对涉及生命之事的认识。我们定义智慧为确定的知识，是对神与人之事的确凿无误的把握，涵盖现在、过去和未来，这是主藉着祂的来临和先知教导我们的。而且，它因已被传达而不可驳倒。因此，它全然真实，符合[天主的]旨意，因为是通过圣子而得知的。在一方面它是永恒的，在另一方面它随时间而成为有用的。部分地，它是同一的；部分地，它是多样的且无关紧要的；部分地，它毫无情感冲动；部分地，它带有热切的愿望；部分地，它是完美的；部分地，它是不完备的。\par

因此，这种智慧——灵魂与理性的正直，以及生活的纯洁——是哲学所渴望的对象，哲学慈爱地、爱慕地朝向智慧，并倾尽全力以获得它。\par

现在，在我们当中，那些爱慕智慧——万有的创造者与导师，即圣子之知识——的人被称为哲学家；而在希腊人当中，那些从事德性论辩的人也被称为哲学家。哲学由各派（我指的是哲学学派）中不可辩驳的教义与相应的生活方式汇集而成；这些教义从外邦人那里窃取了天主赐予的恩宠，并用希腊言辞加以装饰。因为他们有些是借用的，有些是误解的。至于其他方面，他们因其感动而说的话，尚未完全推敲；另一些则凭人的臆测和推理，其中他们也跌倒。他们认为自己已经完美地把握了真理，但按我们的理解，只是部分而已。因此，他们除了这个世界之外一无所知。这就像几何学，通过平面上的图形处理度量、大小和形状；又像绘画，在场景中似乎包含了全部视野。但它按照艺术的规则，利用视线路径所产生的符号，对视野作了虚假的描述。通过这种方式，视野中的高点、低点和中间点得以保留；有些物体似乎出现在前景，有些在背景，另一些则以其他方式出现在平滑的平面上。同样，哲学家们仿效绘画的方式复制真理。在他们每一个人身上，自爱始终是他们所有错误的根源。因此，人不应当为了那止于人的荣耀而受自爱所驱使；而应当爱天主，以智慧成为真正圣洁。所以，如果有人把特殊的当作普遍的，把仆役当作主，他就错失了真理，不明白达味在忏悔中所说的话：\Quote{我像吃饼一样吃了土〔灰〕。} 在他看来，自爱和自负就是土和错误。但如果是这样，科学和知识就来自教导。如果有教导，你就必须寻求师傅。克雷安德斯认齐诺为师，梅特罗多洛斯认伊壁鸠鲁为师，泰奥弗拉斯托斯认亚里士多德为师，柏拉图认苏格拉底为师。但若我来到毕达哥拉斯、斐瑞居德斯、泰勒斯和最初的贤者们那里，我就无法找到他们的师傅。即使你说是埃及人、印度人、巴比伦人乃至玛吉人，我也不会停止追问他们的师傅。我将你引向最初的人类，并从那里开始探究谁是他们的师傅。不是任何人，因为他们还没有学习。也不是任何天使，因为天使以天使的方式说话，人听不见；我们也无法指望他们像我们一样有耳朵，有相应的舌头；也没有人会把发声器官——即嘴唇以及邻近的部分、喉咙、气管、胸膛、呼吸和振动空气——归于天使。天主远在不可接近的圣洁中，甚至与总领天使分离，祂不会大声呼喊。\par

我们也已经听说，天使们及其统治者——因为他们有开端——学习了真理。因此，我们必须继续上升，寻找他们的师傅。既然无始的存在是独一的全能天主；那么首生者也是独一的，\BibleQuote{万物是藉着祂而造的，凡受造的，没有一样不是藉着祂而造的。} \BibleRef{若 1:3} \Quote{因为独一的天主，祂形成了万物的开端；} 指出\Quote{首生的圣子}，彼得写道，准确地理解这句话，\BibleQuote{在起初，天主创造了天地。} \BibleRef{创 1:1} 祂被所有先知称为智慧。这就是那一切受造物的导师、天主的同谋者，祂预知一切；祂从高处，从世界最初的奠基开始，\BibleQuote{多次并以多种方式，} \BibleRef{希 1:1} 训练并成全；因此，\BibleQuote{不要在地上称人为你们的师傅。} \BibleRef{玛 23:8}\par

你看，真正的哲学从何而来，尽管法律是真理的影像和影子；因为法律是真理的影子。但希腊人的自爱宣称某些人是他们的师傅。因此，正如整个家族追溯到天主创造者\BibleRef{弗 3:14}，同样，一切成义的善事的教导也归向主，并引领和促成这一点。\par

但如果任何受造物以任何方式接受了真理的种子，却没有培育它们，反而将它们投在贫瘠无雨的土壤中，以杂草窒息它们，如同法利塞人违背法律，引入人的教导——这些的原因不在于那位教师，而在于那些选择不顺服的人。但他们中相信主的来临和圣经简明教导的人，就达到了对法律的认识；同样，那些热爱哲学的人，藉着主的教导，被引入对真正哲学的认识：\Quote{因为主的言语是纯净的言语，在火中炼过，在地上试过，七次净化。} 正如银子多次被炼净，义人也受到考验，成为主的钱币，接受君王的肖像。或者，既然\Person{撒罗满}也称\Quote{义人的舌头是经过火炼的金子}，\BibleRef{箴 10:20} 表明那经过考验且明智的教训，应受赞美和接纳，当它在地上充分试炼时，也就是说，当\OriginalTerm{真知者}{Gnostic}的灵魂以多种方式圣化，藉着脱离地上的火。它所居住的身体也被净化，被用于圣殿的纯洁。但首先发生在身体上的净化，灵魂居先，是戒除恶事，有些人认为这就是成全，事实上，这是普通信徒——犹太人和希腊人——的成全。但在\OriginalTerm{真知者}{Gnostic}方面，在被认为是他人成全的事之后，他的正义进升到行善的行动。并且，无论谁，若正义增长的力量进升到行善，在他身上，成全便存留在效法天主的固定行善习惯中。因为那些是\Person{亚巴郎}的子孙，更是天主的仆人，就是\Quote{蒙召者}；而\Person{雅各伯}的子孙就是被选者——他们绊倒了邪恶的势力。\par

如果我们断言基督自己就是智慧，并且是祂的工程在先知们中显现，藉此可以学到真知传统，如同祂亲自在祂临在时教导宗徒们；那么结论就是：那\Emph{\OriginalTerm{真知}{gnosis}}，即对现在、未来和过去之事的认识和理解，是确实可靠的，因为是由天主之子传授和启示的，就是智慧。\par

如果智慧人的目标是默观，那么那些仍是哲学家的人的目标也指向它，但从未达到，除非在学习的过程中接受了已启示的先知性话语，藉此它把握现在、未来和过去——它们如何存在、曾经存在和将要存在。\par

这\Emph{\OriginalTerm{真知}{gnosis}}本身是由宗徒们以口传方式传授下来，传给少数人的。因此，知识或智慧应当被操练，直到达到默观的永恒不变习惯。\par

% structure: p0283 source=h2 target=subsection reason=relative-heading-rank; base=h2

\subsection{\Emph{第八章 哲学是天主赐予的知识}}

因为\Person{保禄}在书信中并未明显贬低哲学；但他认为，那已达到\OriginalTerm{真知者}{Gnostic}高度的人，再回到希腊的\Quote{哲学}是不相称的，他形象地称之为\Quote{世俗小学}，认为它是最初级的东西，是为真理所作的预备训练。因此，他写给那些再次从信德退回到法律的希伯来人时说：\Quote{你们还需要有人教导你们天主圣言的初步道理吗？你们成了需要奶，而不是坚硬食物的人。}\BibleRef{希 5:12} 同样，写给那些是希腊归化者的哥罗森人：\Quote{你们要谨慎，免得有人用哲学和虚空的欺骗，按人的传统，按世俗小学，而不是按基督，把你们掳去。}\BibleRef{哥 2:8}——引诱他们再次回到哲学这初级的教义。\par

若有人说，哲学是由希腊人通过人类智慧发现的，但我却发现圣经说明悟是由天主赐予的。因此，圣咏作者将明悟视为最大的恩赐，并且恳求说：\Quote{我是祢的仆人，求祢赐我明悟。} \Person{达味}在祈求丰富的知识经验时，岂不也写道：\Quote{求祢教导我温和、训练和知识，因为我信从了祢的诫命？} 他承认盟约具有最高权威，且是赐给更卓越者的。因此，圣咏再次论及天主说：\Quote{祂没有这样对待任何民族，也没有将祂的判决显示给他们。} 这说法\Quote{祂没有这样做}表明祂做了，但不是\Quote{这样。} 那么，\Quote{这样}是比较而言的，指向优越性，这在我们身上实现。先知本可以简单地说\Quote{祂没有做，}而不用\Quote{这样。}\par

此外，\Person{伯多禄}在\BibleBook{宗徒}{大事录}中说：\Quote{我确实看出天主是不偏待人的；但每个民族中，凡敬畏祂并履行正义的，都蒙祂接纳。}\BibleRef{宗 10:34}\par

天主不偏待人，这并非在时间中，而是从永恒开始。祂的恩惠也没有起点，也不限于地方或人群；因为祂的恩惠不受局部限制。\Quote{敞开正义之门}，经上说；\Quote{进入其中，我要向主谢恩。这是主的门，义人要从它进入。}\Person{巴拿巴}解释先知的话说：\Quote{有许多门开着，那在正义中的门，就是在基督内的门，凡进入的都蒙福。}同一意义的还有以下先知的话：\Quote{主在多水之上；}不仅指不同的盟约，也指教导的方式，无论是在希腊人中还是在异邦人中，都有助于正义。\Person{达味}清楚地为真理作证，歌唱说：\Quote{愿罪人转向\OriginalTerm{阴府}{Hades}，所有忘记天主的民族。}他们显然忘记了祂，他们从前曾记得祂，并在忘记祂之前就认识祂，却将祂抛弃。所以，在万民中也有对天主模糊的认识。关于这些，就说到这里。\par

现在，\OriginalTerm{真知者}{Gnostic}必须有学识。既然希腊人说，由\Person{普罗泰戈拉}带头，发明了以论证反对论证的方法，那么适当地谈谈这类论证是有必要的。因为圣经说：\Quote{多言者，也当听对方的话。}\BibleRef{约 11:2} 然而谁能明白主的比喻呢？惟有智者、明达者、以及爱主的人。这样的人应当忠实，能宣讲其知识，在辨别话语上有智慧，在行动上灵巧，并保持纯洁。\Quote{他显得越大，就越当谦卑，}——这是\Person{克莱孟}在\WorkTitle{致格林多人书}中所说的——\Quote{这样的人能够遵行命令：\InnerQuote{有些人要从火中救出，另有些人要怀着怜悯，有所区分。}}\par

修剪镰刀当然主要是为修剪；但我们也用它来分开缠绕的嫩枝，割除与葡萄藤一同生长的荆棘，这些荆棘不太容易够到。所有这些都与修剪有关。同样，人主要是为认识天主而造的；但人还丈量土地，从事农耕，进行哲学思辨；这些追求中，一样有助于生活，一样有助于生活美好，第三样则有助于研究可以证明的事物。此外，让那些说哲学源于魔鬼的人知道，圣经说：\Quote{魔鬼伪装成光明的天使。}\BibleRef{格后 11:14} 当它要做什么？显然是要预言。但如果它作为光明的天使预言，它就会说真话。如果它预言属于天使和光明的事，那么当它按照相似的行为方式转化时，它就会预言有益的事，尽管它在背信方面是不同性质的。因为如果不先吸引慕知者与之共融，然后继而将他引向谎言，它怎能欺骗任何人呢？尤其它会被发现知道真理，即使不是完全理解，至少也不是全然不熟悉。\par

哲学并不因此是虚假的，尽管窃贼和说谎者因运作方式的转化而说出真理。也不应因说话的人而愚昧地对所说的话宣判（这也应在那些如今被声称预言的人身上留意）；而应审查所说的话，看它是否持守真理。\par

一般而言，我们断言一切为生活所必需和有益的事物都来自天主，并且哲学尤其赐给了希腊人，作为他们特有的盟约——它本身是通往基督哲学的阶梯——尽管那些从事希腊哲学的人自愿对真理塞耳不听，轻蔑异邦人的声音，或者也因国家法律而惧怕悬在信徒头上的危险。\par

正如在蛮族哲学中，同样在希腊哲学中，那位莠子的农夫也播种了\Quote{莠子}；因此，异端与茁壮的麦子一同在我们中间生长；而那些在希腊哲学中宣扬伊壁鸠鲁的不敬与纵欲，以及其他违背正确理性传播的学说的人，在希腊人中如同天主所赐农艺的劣果。这种纵欲自私的哲学，宗徒称之为\Quote{今世的智慧}；因为它只教导今世及关乎今世的事，并因此在统治权方面，服从于今世的掌权者。故此，这种片面的哲学甚为初浅，而真正完美的科学则处理理智的对象，这些对象超越世界的领域，并且比那些\BibleQuote{眼所未见，耳所未闻，人心所未想到的}更为属灵，直到那位导师向我们揭示它们的奥秘；揭开至圣所；并且按上升的次序，向那些真正而非虚假地承受主嗣业的人显示更圣洁的事。因为我们现今敢断言（因为这是以知识为特征的信德），这样一个人知道万事，并以确定的领悟理解万事，关乎那些对我们困难、真正属于真知的事——如雅各伯、伯多禄、若望、保禄以及其他宗徒。因为预言充满\OriginalTerm{知识}{gnosis}，因为它是主所赐，又由主向宗徒解释。\OriginalTerm{知识}{gnosis}不正是理性灵魂的一种属性吗？它为此锻炼自己，以便借着知识成为不朽的。因为知识和冲动都是灵魂的能力。冲动被发现是同意后的运动。因为对行动有冲动的人，首先接受行动的知识，其次才是冲动。让我们进一步留意这一点。因为学习先于行动（因为自然地，做其所愿之事的人首先学习它；知识来自学习，冲动跟随知识，之后是行动）；知识成为一切理性行动的开端和起因。所以，理性灵魂的独特本性正当地唯独由此特征所界定；因为实际上，冲动如同知识，由存在的事物所激发。而\OriginalTerm{知识}{gnosis}本质上是灵魂对存在物的默观，或对某一物或某些物，当它达至成全时，则是对所有物的默观。虽然有人说，智者相信有些事不可理解，以至于对它们有某种理解，因为他理解不可理解之事是不可理解的；这是普遍的，属于那些只能感知微小事物的人。因为这样的人肯定有些事不可理解。\par

但我所说的那位诺斯替者，他自己理解别人似乎不可理解的事；相信没有什么是天主之子不可理解的，因此没有什么不可教导的。因为那为爱我们而受苦的那一位，不会隐瞒任何为我们教导所需的知识元素。因此，这信德成为确实的证据；因为真理跟随天主所传授的。但若有人渴望广博的知识，\BibleQuote{他懂得往古，推测未来；明白言词的曲折，谜语的解答。智慧的弟子预知征兆和奇迹，以及时代和时期的结局。}\BibleRef{智 7:17}\par

% structure: p0294 source=h2 target=subsection reason=relative-heading-rank; base=h2

\subsection{第九章　诺斯替者摆脱一切灵魂的扰动}

诺斯替者是这样的，他仅受维持身体所需的情绪影响，如饥、渴等。但在救主的情况中，[以为]身体作为身体需要必要的帮助以维持其存在，那是可笑的。因为祂吃东西，不是为了身体——身体是由神圣的能力维持的——而是为了不让那些与祂在一起的人对祂产生不同的想法；如同后来一些人确实以为祂是\OriginalTerm{幻影}{\GreekText{δοκήσει}}显现的。但祂是完全\OriginalTerm{不动情的}{\GreekText{ἀπαθής}}；对任何情感运动——无论是快乐还是痛苦——都无动于衷。而宗徒们，通过主的教导，以最诺斯替的方式掌握了愤怒、恐惧和贪欲，因此他们甚至不受那些看似好的情感运动的影响，如勇敢、热忱、喜乐、欲望，因为他们心境稳定，丝毫不改变；在主复活之后，持续不变地处于训练状态。\par

即使承认上述情感在理性产生时是好的，但在完人身上仍不可接受，因为他不能行使勇敢：因为他不会遇到令人恐惧的事，因为他认为生活中发生的事没有一件是可畏惧的；也没有任何事能使他脱离——他对天主的爱。他也不需要喜乐的心情；因为他不会陷入痛苦，确信一切事都进行顺利。他不愤怒；因为没有事能激起他的愤怒，因为他永远爱天主，全心转向祂，因此不恨天主任何的受造物。他也不嫉妒；因为他什么也不缺，为达致相似，使他成为卓越和良善所需的一切。因此，他也不以这种普通的情感爱任何人，而是爱造物主于受造物中。因此，他也不陷入任何欲望和渴求；就灵魂而言，他也不需要任何属于他人的东西，因为他通过爱与那被爱者结合，由自由选择以及由训练而来的习惯，更接近祂，并因美好事物的丰盈而得福。\par

因此，基于这些理由，他被迫在\OriginalTerm{不动心}{impassibility}上变得像他的导师。因为天主的圣言是理智的，正如心灵的肖像唯在人类身上得见。同样，善人在灵魂的形式和模样上也像天主。另一方面，天主也像人。因为每个人独特的形式是那界定我们的心灵。因此，那些得罪人的人也是不圣洁和不虔敬的。因为说\OriginalTerm{真知者}{gnostic}和完美的人必须根除愤怒和勇气，那是可笑的——若没有这些，他将无法与环境抗争或忍受可怕之事。但若我们从他那里夺走\OriginalTerm{欲望}{desire}，他将完全被困扰所压倒，因而极为卑劣地离开此生。除非如某些人所认为，他拥有欲望，否则他将不会对类似卓越和善的事物产生欲望。那么，如果所有与善的联合都伴随着欲望，据说，渴望卓越的人如何能保持不动心呢？\par

但这些人似乎不知道\OriginalTerm{爱德}{love}的神性。因为爱不是爱者一方的欲望，而是一种情感的关系，将\OriginalTerm{真知者}{Gnostic}恢复到\OriginalTerm{信德}{faith}的合一——超越时间和地点。但藉着爱已经身处于他注定所处的境界中、并藉着知识预先把握了希望的人，并不渴望任何事物，因为他已尽可能拥有了所渴望的事物本身。因此，可以预期的是，他持续操练真知的爱，处于那不变的状态中。\par

因此，他也不会急切地渴望与优美相似，因为他已藉着爱拥有优美。对于那藉着爱获得与\OriginalTerm{不动心的天主}{impassible God}的亲近、并藉着爱将自己登记在天主朋友中的人，勇气和欲望对他还有什么必要呢？\par

因此，我们必须将\OriginalTerm{真知者}{gnostic}和完美的人从灵魂的一切情欲中解救出来。因为\OriginalTerm{知识}{gnosis}产生实践，实践产生习惯或性情；而这种状态产生\OriginalTerm{不动心}{impassibility}，而非情欲的节制。欲望的完全根除所收获的果实是不动心。但真知者既不分享那些通常被赞颂为善的情感——即与情欲相关的良好情感，如喜悦（与快乐相关）、沮丧（与痛苦相连）、谨慎（隶属于恐惧）——也不分享高昂，因为它与愤怒为伍；尽管有些人说这些不再是恶，而已经是善。因为那一次藉着爱被成全、永恒且不满足地享受默观无限喜乐的人，不可能以渺小卑下的事物为乐。对于获得了\Quote{不可接近的光}\BibleRef{弟前 6:16}的人，还有什么理性原因促使他回到世界的美善呢？虽然就时间和地点而言尚未真实，但藉着那真知的爱——藉此继承和完美复归随之而来——酬报的施予者以行动实现了真知者藉着真知的抉择、藉着爱预先把握的事物。\par

因为藉着为爱\OriginalTerm{主}{Lord}而走向主，虽然他的帐篷可见于地上，他并没有从生命中退出。因为那是不容许他的。但他已将灵魂从\OriginalTerm{情欲}{passions}中退出。因为那是赐予他的。另一方面，他活着，已治死了自己的私欲，不再使用身体，而是容许身体使用必需品，以免给出解体的理由。\par

那么，一个不在危险中的人，既不在现场，而已经完全与所爱者同在，他还有什么需要\OriginalTerm{刚毅}{fortitude}呢？对于一个不需要\OriginalTerm{自制}{self-restraint}的人，自制又有何必要？因为拥有需要自制来控制的欲望，是尚未洁净而受情欲支配之人的特征。现在，刚毅因恐惧和怯懦而被取用。因为天主的朋友——\Quote{天主在创造世界以前已拣选了他}\BibleRef{弗 1:4}，使他被登记在至高的\Quote{\OriginalTerm{义子}{adoption}}中——再陷入快乐或恐惧、忙于压制情欲，就不再合宜。因为我敢断言，正如他藉着将要做的事和将获得的恩赐而被预定，同样他也因他所知和所爱者而预定了自己，不像群众那样模糊地猜测未来，而是藉着真知之信德把握了他人所隐藏的。藉着爱，未来已为他成为现在。因为他藉着预言和基督的来临，相信了不说谎的天主。他所相信的，他已拥有，并持守着许诺。那许诺者是真理。藉着许诺者的可信性，他藉着知识牢牢抓住了对许诺结局的确信。那知晓在所处环境中对未来的确实把握的人，藉着爱迎接未来。因此，那确信将获得真正美善的人，不会祈求获得现世之物，而是祈求始终持守那正中靶心的信德，并获得成功。此外，他将祈祷尽可能多的人变得像他一样，为光荣天主，这光荣藉着知识得以成全。因为那变得像救主的人也致力于拯救，忠实地执行诫命，直至人性所能容纳的肖像。这就是藉着行动和对真正正义的知识来敬拜天主。主不会等待这人在祈祷中发声。祂说：\Quote{求，我必做；想，我必给。}\par

简言之，不变者在可变者中取得坚定与稳固是不可能的。但主导官能若处于永恒变化中，因而也不稳定，则习惯的力量无法维持。因为一个被外在事件与偶然不断改变的人，怎能拥有习惯、意向，一言以蔽之，拥有\OriginalTerm{科学知识}{\GreekText{ἐπιστήμη}}的把握？再者，哲学家也将德性视为习惯、意向和科学。正如\OriginalTerm{知识}{gnosis}并非与生俱来，而是后天获得的，其要素的获得需要专注、训练与进步；然后经由不断的实践，它便成为习惯；同样，当它在奥秘的习惯中得以成全时，它便恒存，因爱而绝无谬误。因为他不仅领悟了第一原因及其所产生的原因，并对此确信无疑，牢牢持有坚定、无可反驳且不可动摇的道理；而且对于善恶、一切受造物，总而言之，对于主所谈论的一切，他从真理本身学会了从世界的根基到终结的最精确的真理。他不把看似可信的、或按希腊推理必然的东西置于真理本身之上；而是将主所说的话视为清晰而明显的，尽管对他人是隐藏的；并且他已经领受了一切事物的知识。我们所拥有的神谕陈述现存事物如其本然，未来事物如其将然，过去事物如其已然。\par

在科学之事上，作为唯独拥有科学知识的人，他将居首位，并会论述关于善的讨论，始终专注于理智对象，从上面的原型中推演他在人类事务中的行为；如同航海者依星辰指引船只；随时准备从事一切适当的行动；惯于藐视一切困难和危险，当必须忍受它们时；从不做任何仓促或与自身或公共利益不相称的事；有先见之明；无论醒时还是梦中都不为快乐所动摇。因为习惯了节俭的生活与简朴，他适度、活跃而庄重；只需少量生活必需品；不从事任何多余之事。他并不以这些为首要追求，而是出于生活的共融，作为他旅居世间的必需品，仅取所需。\par

% structure: p0305 source=h2 target=subsection reason=relative-heading-rank; base=h2

\subsection{第十章 知识者运用一切人类知识的帮助}

因为对他而言，\OriginalTerm{知识}{gnosis}是首要之事。因此，他应用那些训练知识的科目，从每一门学问中汲取其对真理的贡献。于是，在音乐中研究和谐的比例；在算术中注意数字的增加与减少及其相互关系，以及多数事物如何归属于某种数字比例；研究几何学这一抽象本质，他感知到持续的距离和一种不同于这些物体的不变本质；又通过天文学，在思想上从大地升起，与天一同高升，并随其运转而运转；永远研究神圣的事物及其彼此间的和谐；\Person{亚巴郎}正是由此出发，上升到对创造它们的那位的认识。此外，知识者将运用辩证法，区分种类中的类别，并掌握存在物的区分，直至达于首要而单纯之物。\par

但群众畏惧希腊哲学，如同孩童畏惧面具，担心它会使他们误入歧途。然而，如果他们拥有的信德（因为我不能称之为知识）会被动听的言辞所瓦解，那么不妨任其瓦解，并让他们承认自己不会持守真理。因为真理是不可动摇的；而虚假的意见则会瓦解。举例来说，我们将一种紫色与另一种紫色比较。所以，如果有人承认自己没有一颗被矫正的心，他就没有兑换银钱的桌子或检验言辞的标准。那么，一个不能当场验证并区分伪币的人，又如何能继续当兑换银钱者呢？\par

\Person{达味}曾呼喊道：\Quote{义人永远不动摇；}因而既不会被欺骗性的言辞，也不会被错误的快乐所动摇。因此他永远不会从自己的产业中被摇动。\Quote{他不怕凶信；}因此既不怕无根据的毁谤，也不怕周围的错误意见。他更不会畏惧诡辩的言辞，因为他能够分辨它们，或正确回答所问的问题。辩证法就是这样一道屏障，使真理不致被诡辩家践踏。\Quote{因为那些以主圣名赞美的人，}按先知所说，\InnerQuote{应当心中喜乐，寻求上主。你们当寻求祂，并要刚强。时时在一切道路上寻求祂的慈颜。}\BibleQuote{因为，天主曾多次并以多种方式说话，}\BibleRef{希 1:1}祂的认识并不限于一种方式。\par

因此，我们的知识者并非利用这些作为德性而成为博学之人。而是借助它们作为工具来区分共同与特殊，从而接纳真理。因为一切错误与虚假意见的原因，在于无法区分事物在哪些方面是共同的，在哪些方面是不同的。除非在区分事物时，人密切留意言辞，否则他便会无意中混淆共同与特殊之处。这种情况一旦发生，他必定会陷入无路可走的迷途与错误。\par

圣经本身对名称和事物的区分也在人的灵魂中产生极大的光明。因为有必要理解表达多种事物的词语，以及表达同一事物的多种词语。其结果是准确的答复。但必须避免那种在无关问题上耗费精力的极大无益；因为诺斯替者利用各门学问作为辅助性的预备练习，以便尽可能不受干扰地、准确地传达真理，并抵御那些企图消灭真理的推理。因此，他不会在有助于学程和希腊哲学的精进方面有所欠缺；但这并非主要的，而是必要的、次要的，且出于环境的原因。因为那些在异端中劳苦之人所邪恶使用的，诺斯替者将正确使用。\par

因此，出现在希腊哲学中的真理是部分的，而真正的真理，正如太阳照耀在白色和黑色之上，显明它们各自的性质。同样，它也暴露了一切似是而非的巧辩。所以，希腊人也正确地宣告了：\Quote{真理女王是伟大德行的开端。}\par

% structure: p0312 source=h2 target=subsection reason=relative-heading-rank; base=h2

\subsection{第十一章　数字比例、几何比及音乐中的奥秘含义}

如同在天文学中我们有亚巴郎为例，同样在算术中我们也有同一个亚巴郎。\Quote{因为，听说罗特被掳，他数点自己家里生养的仆人，共\OriginalTerm{三百一十八}{\GreekText{τιή}}人，}他击败了极多的敌人。\par

他们于是说，代表三百的字符，就其形状而言，是主印记的样式；而\Emph{\OriginalTerm{约塔}{Iota}}和\Emph{\OriginalTerm{厄塔}{Eta}}指示救主的名字；因此，亚巴郎的家人是在救恩中，他们逃向印记和名字，成为被掳者的主，以及跟随他们的许多不信民族的主。\par

如今，数字三百是三乘一百。十被认为是完满的数字。而八是第一个立方，它在长、宽、高各维度上相等。\Quote{人的日子将是，}经上记着，\Quote{\OriginalTerm{一百二十}{\GreekText{ρκ}}岁。}\par

\SourceRecord{Translated by William Wilson.；Ante-Nicene Fathers；Vol. 2.；Edited by Alexander Roberts, James Donaldson, and A. Cleveland Coxe.；Buffalo, NY: Christian Literature Publishing Co.,；1885.；<http://www.newadvent.org/fathers/02106.htm>.}

% structure: p0001 source=h1 target=section reason=page-title

\section{《杂记》（第七卷）}

% structure: p0002 source=h2 target=subsection reason=relative-heading-rank; base=h2

\subsection{第一章 真正的诺斯替者是敬拜天主的人，被不信者不公正地毁谤为无神论者}

现在是时候向希腊人证明，唯有诺斯替者是真正虔诚的；这样，哲学家们得知真正基督徒的品性后，便会谴责自己的愚昧，因他们轻率而不加考虑地迫害基督徒的名号，并且毫无理由地称那些认识真天主的人为不虔诚者。我认为，必须向哲学家们使用更清晰的论证，使他们能凭借自己通过训练已获得的练习得以理解，尽管他们尚未证明自己配得上分享信德的能力。\par

目前我们不拟论及先知之言，因为我们随后会在适当之处引用圣经。但我们要扼要指出其中所指示的要点，在我们对基督信仰的描述中，以免立即引用圣经（尤其因为他们尚未理解其中的话语）而打断论述的连续性。但在指出所指示的事之后，要向那些已相信的人提供丰富的证明。\par

但如果我们的断言在某些群众看来与主的圣经不同，须知这些断言正是从那源头得到气息和生命；并且它们从圣经出发，声称只引出含义，而非字句。因为更深入的处理，既然不合时宜，将理当显得多余。因此，不关注紧急之事将是极度怠惰和有缺陷的；\Quote{真实有福的，是那些查考主的证言、并以全心寻求祂的人。}法律和先知都为天主作证。\par

因此，我们的目的是证明，唯有诺斯替者才是圣洁和虔诚的，并以相称的方式敬拜真天主；而相称于天主的敬拜随之而来的就是被天主所爱并爱天主。他因此按着价值判断一切卓越皆为可敬的；并判断，在我们感官感知的对象中，我们应当尊敬统治者、父母以及每一位年长者；在学问的对象中，应当尊敬最古老的哲学和原始的预言；在理智的观念中，应当尊敬起源于最古老的、无时间且无始的第一本原，即万有的起始——圣子——从他那里我们得以认识更远的原因，即万有的圣父，最古老且最仁慈的一位；非言语所能表达，而要以敬畏、静默和神圣的惊叹来尊崇，并至高崇敬；由主启示给那些学习者所能理解的，并被主所选召者所承认；宗徒说，\Quote{他们的官能因练习而得到了锻炼。}\BibleRef{希 5:14}\par

那么，对于诺斯替者而言，事奉天主就是他灵魂持续不断的研究和专注，以不止息的爱献于天主。因为对人事奉的种类，一种以改进为目的，另一种是服务性的。身体的改进是医术的目标，灵魂的改进是哲学的目标。服务性的侍奉是子女对父母、臣民对统治者所做的。\par

同样，在教会中，长老负责改进的部门，执事负责服务性的部门。在这两种工作中，天使在天主管理地面事务时事奉天主；诺斯替者本人也事奉天主，并向人们展示改进的计划，如同他被指派以纪律训练人使之改过的方式。因为唯有在人间事务中正确而无可指责地事奉天主的人才是虔诚的。正如通过农学的知识和技能，植物的最佳处理方式使其果实得以产生和收获，从而为人们带来益处；同样，诺斯替者的虔诚，收取那些因他而相信的人所结的果实，当不少的人借着他达到知识并因而得救时，便通过他的技能获得最好的收成。正如\OriginalTerm{虔敬}{\GreekText{θεοπρέπεια}}是保持相称于天主的习惯，虔诚的人乃是唯一爱天主的人；而这样的人将知道什么是在知识和生活中相称于那注定要成为\OriginalTerm{神圣}{\GreekText{θεῷ}}并已与天主相似的人。因此，他首先是爱天主的人。因为正如尊敬父亲的人是爱父亲的人，同样，尊敬天主的人是爱天主的人。\par

如此，在我看来，诺斯替能力有三种效果：第一，认识事物；第二，实践圣言所建议的一切；第三，以相称于天主的方式传达隐藏在真理中的奥秘。\par

那么，一个确信天主是全能的，并从祂的唯一圣子学习了神圣奥秘的人，怎能是\OriginalTerm{无神论者}{\GreekText{ἄθπεος}}呢？因为无神论者是认为天主不存在的人。而迷信者是畏惧魔鬼的人；他将万物神化，包括木头和石头；并将精神以及拥有理性生命的人置于束缚之下。\par

% structure: p0011 source=h2 target=subsection reason=relative-heading-rank; base=h2

\subsection{第二章 圣子为万有的统治者和救主}

因此，认识天主是信德的第一步；然后，通过信赖救主的训诲，认为以任何方式行恶都不适于对天主的认识。\par

因此，地上最好的事物是最虔敬的人；天上最好、更近且更纯净的事物是天使，他参与永恒而蒙福的生命。然而，那最亲近那唯独全能者的圣子的本性，是最完美、最神圣、最有权能、最尊贵、最具王权、最仁慈的。这是最高的卓越，它依照圣父的旨意安排万有，以不倦不疲的权能最佳地执掌宇宙之舵，运作它所运行的一切，同时注视其隐秘的计划。因为从祂自身的角度而言，天主子从不被移置；祂不被分割、不被分离、不从此处移到彼处；祂始终无所不在，却不为任何空间所容纳；祂是圆满的心智，圆满的父之光；全是眼睛，看见一切，听见一切，知道一切，以祂的权能洞察诸权能。所有天使和诸神的队伍都服从于祂；祂，圣父的圣言，为那将万有置于祂权下者展示神圣的治理。\par

因此，所有的人也都属于祂：有些通过知识，有些尚未如此；有些作为朋友，有些作为忠信的仆人，有些仅仅作为仆人。祂是导师，通过奥秘训导诺斯替者\OriginalTerm{诺斯替者}{Gnostic}，通过美好的希望训导信徒，通过可感知的惩戒训导心硬之人。因此，祂的眷顾在私下、在公开、在一切地方。\par

而我们称为救主和主的这一位，就是天主子，预言性的圣经明确证明了这一点。所以，万有的主，希腊人和外邦人的主，劝勉那些愿意的人。因为祂并不强迫那能（通过选择和完成来自祂的、与把握希望相关的事）从祂领受救恩的人。\par

祂也是那藉着次等天使将哲学赐予希腊人的那一位。因为按照古代神圣的秩序，天使被分配在各民族中。但相信的人的荣耀是\Quote{主的产业}。因为主要么不关心所有的人——这或许是因祂不能（这不应被设想，因为那将是软弱的证明），或因祂不愿（这不是善者的属性）；而为我们的缘故取了能受苦的肉身的那一位，远非沉溺于安逸。要么祂关心所有的人——这适合于那成为万有之主的那一位。因为祂是救主；不是某些人的救主，而非其他人的。但祂按照每人所拥有的适应程度，将祂的恩惠分施给希腊人和外邦人，甚至给他们中预定、并在适当时候蒙召的信徒和选民。那同等呼召众人、并给以特别卓越方式相信者分配特殊尊荣的那一位，永不会嫉妒任何人。那为万有之主、且超越一切服事良善而全能的圣父旨意的那一位，绝不会被他人阻碍。但嫉妒并不触及那从无始以来就不受苦难的主；人的事也不至于到被主嫉妒的地步。而是另一个被情感所触动的\Quote{他}嫉妒。也不能说主因无知而不愿拯救人类，因为祂不知道如何照顾每个人。因为无知不适用于那在世界奠基之前就是圣父谋士的天主。因为祂就是智慧，统治的天主\Quote{在其中喜悦}\BibleRef{箴 8:30}。因为圣子是天主的德能，作为万有被造之前圣父最古老的圣言，也是祂的智慧。因此，祂被恰当地称为由祂所造之物的导师。祂也从不因被享乐所牵引而放弃对人类的关怀，祂取了那本性能受苦的肉身，将其训练至不受苦难\OriginalTerm{不受苦难}{impassibility}的状态。\par

祂如何是救主和主，如果不是所有人的救主和主？但祂是那些因愿意认识而相信之人的救主，也是那些尚未相信之人的主，直到他们得以承认祂，并藉着祂获得那独特而合适的恩赐。\par

如今，主的动力\OriginalTerm{主的动力}{energy of the Lord}指向全能者；圣子，可以说，是圣父的一种动力。因此，救主绝不可能是憎恨人类者；祂因对人性肉身的极大爱，不轻视其能受苦的特性，反而穿上它，来为人类的共同得救；因为那些选择相信之人的信仰是共同的。不仅如此，祂绝不会忽略自己的工程，因为人比其他所有受造活物，在其受造时就被赋予了天主的观念。也不可能有任何其他更好、更适合于人的治理，比那由天主所设立的了。\par

因此，那本性上卓越者常适宜居于低劣者之上，那能妥善管理某事的人适合被指派管理它。如今，真正统治和主持的是圣言及其眷顾，它鉴察万有，并不轻忽任何属于它的事务的关怀。\par

所以，那些选择属于祂的人，就是那些通过信德而得以成全的人。圣子，凭着全能圣父的旨意，是万善的原因，是第一运动动力因\OriginalTerm{第一动力因}{first efficient cause}——一种感官无法把握的德能。因为祂所是的，未被那些因肉身的软弱而无法领会实况的人看见。但祂取了可感觉的肉身，来向人展示通过服从诫命而可能之事。因此，作为圣父的德能，祂轻易胜于祂所愿的，不忽略祂治理中哪怕最细微的一点。否则，万有就不会被祂完美地执行了。\par

但我认为，至高权能的特征是对所有部分进行精确审察，直至最细微之处，终止于宇宙的首要治理者，祂凭圣父的意愿指导所有人的救恩；有些监督者隶属于他人之下，而他人又置于他们之上，直到你来到伟大的大司祭面前。因为在一个原初的第一原则（祂按圣父的意愿行事）上，第一、第二和第三都依附于祂。然后在可见世界的最高端，是蒙福的天使团；而直到我们自身，那些从一位、通过一位而得救并救人者，层层排列，一些在另一些之下。\par

因此，正如铁的最小微粒被赫拉克勒斯石的精神所移动，当它散布在许多铁环上时；同样，被圣神所吸引，有德行的人通过亲和力被添加到第一居所，而其他的人则依次直到最后。但那些因软弱而邪恶的人，从邪恶的贪得无厌中堕落成败坏的状态，既不能控制也不受控制，四处乱转，被激情所卷，坠落于地。\par

因为从一开始的法律就是：德行应是自愿选择的对象。故此，按法律和先于法律所颁布的诫命，并非为义人设立的（\BibleQuote{因为法律不是为义人设立的}\BibleRef{弟前 1:9}），而是规定那选择它们的人应得永生和蒙福的奖赏。\par

但另一方面，他们允许那以恶为乐者与他所选择的对象为伍；而另一方面，那在获得德行和增加义德上不断进步的靈魂，应在宇宙中获致更好的位置，如同在每一步进阶中趋向于\OriginalTerm{不动心}{impassibility}的习惯，直到\BibleQuote{达到圆满成年}\BibleRef{弗 4:13}，即知识和继承的双重卓越。\par

这些有益的变化，按照变化的顺序，通过时间、地点、尊荣、认知、继承和职务来区分，根据每个变化的特定秩序，直到永恒中至超越且不断的对主的默观。\par

那值得爱慕者引导每一个出于对知识的热爱而全心投入默观者进入对自身的默观。因此，主从同一泉源中汲取祂所赐的第一和第二诫命，既不让法律之前的人没有法律，也不让不熟悉\OriginalTerm{蛮族哲学}{Barbarian philosophy}原理的人不受约束。因为祂以诫命赐予前者，以哲学赐予后者，将不信封闭至主的来临。因此，凡不信者都无可推诿。因为祂通过不同的进步过程，无论是希腊人还是蛮族人，都引向因信德而得的成全。\par

若有一个希腊人，越过\OriginalTerm{希腊哲学}{Hellenic philosophy}的初步训练，直接进入真道，他虽不识字，却因选择了由信德而至成全的捷径，而超越他人。\par

所以，凡不阻碍人自由选择的一切事物，祂都使之成为德行的辅助，为的是让那些只能隐约看见的人，也能以某种方式显明那唯一全能、良善的天主——从永远到永远藉着祂的圣子施行救恩。\par

而另一方面，祂在任何方面都不是恶的原因。因为宇宙的主宰为了宇宙的救恩，以普遍和特别的方式安排了一切。救恩之义的职能便是尽可能改善一切。因为即使次要之事，也是为着那更好者的救恩和适合人品性的居所而安排的。凡有德行的事物都向更好转变；其转变的适当原因是知识自由的选择，这选择是靈魂在自己能力之内的。但必要的矫正，藉着伟大的监察审判官的美善，通过随行的天使、各种预判性的审判和最终的审判，迫使大罪人悔改。\par

% structure: p0030 source=h2 target=subsection reason=relative-heading-rank; base=h2

\subsection{第三章：\OriginalTerm{诺斯替者}{Gnostic}力求尽可能肖似天主及其圣子}

现在，我暂且不提其余的事，只光荣主。但我断言：\OriginalTerm{诺斯替者}{Gnostic}的靈魂以默观的宏伟超越每一位圣品级的生活方式——其中诸神蒙福的居所按分配而定——被视为圣者中的圣者，从全体中转移至全体，到达优于更好之处的深处，拥抱神圣的异象，不是在镜中或藉着镜子，而是以超越清澈、绝对无餍足的异象，这是热切爱主之靈魂的特权，在无尽的世代中欢庆，因所有卓绝的同一性而受尊荣。这样的异象是\BibleQuote{心地纯洁的人}\BibleRef{玛 5:8}所能见到的。这正是成全的\OriginalTerm{诺斯替者}{Gnostic}的功能：通过伟大的大司祭与天主交往，按自己的容量肖似主，全心服务于天主，这服务趋向于人的救恩——通过关怀那以我们为对象的恩惠；另一方面，通过敬拜、教导和行动上的恩惠。\OriginalTerm{诺斯替者}{Gnostic}甚至塑造并创造自身；此外，他也如同天主，装点那些听他教导的人，尽可能将那从实践而来、趋向\OriginalTerm{不动心}{impassibility}的节制，肖似于那本性拥有\OriginalTerm{不动心}{impassibility}的祂；尤其是不断与主交往和共融。我认为，温和、爱人和卓越的虔诚，是\OriginalTerm{诺斯替者}{Gnostic}相似的法则。我断言这些美德\BibleQuote{是天主所悦纳的祭献}\BibleRef{斐 4:18}；圣经声称\BibleQuote{谦卑的心与正确的知识是天主的全燔祭}；每个被接纳进入圣洁的人都被光照，以达至不可分解的合一。\par

因为福音和宗徒都命令人\Quote{将自己掳去}，杀死\Quote{那因情欲而败坏的旧人}，放弃私欲，脱离罪恶，将新的人从死亡中、从\Quote{旧的生活}中复活。\par

因此，法律暗示了这一点：它命令将罪人除掉，从死亡迁到生命，达到因信德而来的不动情。法律的教师们不理解这一点，因为他们视法律为好争辩的，就给那些徒然诽谤法律的人提供了把柄。因此，我们正确地不向天主献祭，因为祂一无所缺，供给万人一切；但我们荣耀那为我们祭献了自己的主，我们也祭献自己；从无所需要到无所需要，从不动情到不动情。因为天主唯独在我们的救恩中喜悦。因此，我们有理由不向那不被快乐征服者献祭，因为烟雾的蒸气远远停在下面，连最厚的云也达不到；而它们所达到的离他甚远。神性既不缺乏任何东西，也不喜爱快乐、利益或金钱，祂充满一切，供给一切有需要的有受造物。神性既不能被祭品和供物，也不能被光荣和尊荣所赢得；祂也不被任何这类事情影响；但祂只向卓越和良善的人显现，这些人绝不会因威胁的恐惧而背叛正义，也不会因丰厚礼物的许诺而动摇。\par

但那些没有认识到人类灵魂的自决性，以及在生命选择上不能被奴役的人，因粗鲁的不义所行之事感到厌恶，就认为没有天主。与这些人意见相等的是那些沉溺于享乐、痛苦和意外灾祸，并蔑视事件，说没有天主，或者说祂存在但不照管一切的人。还有其他人相信他们所认为的神明能被祭品和礼物所说服，似乎偏爱他们的放荡；他们不相信那唯一真实的天主，祂存在于正义之善的不变性中。\par

因此，诺斯替者是虔诚的，他先关心自己，然后关心邻人，使他们成为至善。因为儿子通过显示自己良善并肖似父亲来取悦好父亲；同样，臣民取悦统治者。因为相信和服从是在我们自己的能力之内。\par

但若有人以为罪恶之因是物质的软弱、无知的不由自主冲动以及（愚昧中）非理性的必然性，那已成为诺斯替者的人通过教导胜过了这些，如同胜过了野兽；他效法天主的计划，尽其所能向愿意的人行善。若被置于权威之中，如同梅瑟，他将为被统治者的救恩而治理；他将通过记录对最优秀者的尊荣和对恶人的惩罚，依照理性为纪律而驯服野性和无信。\par

因为义人的灵魂是卓越的天主肖像，相似于天主；在其中，通过服从诫命，如同在圣地，包围并供奉着人类和不朽者的领袖、善的君王和父，祂是真正的法律、正义和永恒圣言，是每个人个别和众人共同的唯一救主。\par

祂是真正的\OriginalTerm{独生子}{Only-begotten}，是宇宙君王和全能圣父光荣的肖像，祂按照自己的肖像将完美静观的印记印在诺斯替者身上；因此，现在有第三个神的肖像，尽可能相似于\OriginalTerm{第二原因}{Second Cause}，即本质生命，我们藉此生活真正的生命；我们所认为的诺斯替者被描述为在确定和完全不变的事物中行动。\par

因此，他治理自己和属于自己的一切，并把握住神圣科学，真诚地接近真理。因为对可理解之物的认知和掌握必然被称为确定的科学知识，其关于神圣事物的功能是思考什么是\OriginalTerm{第一原因}{First Cause}，什么是\BibleQuote{万有是藉着祂而造的，没有祂什么也不能造}；\BibleRef{若 1:3}以及另一方面，什么是渗透性的，什么是包容性的；什么是结合的，什么是分离的；每一个处于什么位置，贡献什么能力和什么服务。又在人类事物中，人本身是什么，他本性或超本性拥有什么；他应该如何行动或承受；他的德行和恶行是什么；关于善、恶和无关紧要的事；也关于刚毅、明智、节制以及各方面完全的德行，即正义。\par

此外，他在获取智慧时运用明智和正义，运用刚毅不仅在于忍受环境，也在于抑制快乐和欲望、忧愁和愤怒；一般而言，要抵抗一切以强力或欺诈诱惑我们的事物。因为不需要忍受恶习和德行，而是要忍受令人恐惧之事。\par

因此，在医术、训练和惩罚中，疼痛被认为是有益的；藉着它，人的品行得以改善而于他们有益。刚毅的形式有忍耐、宽宏、高昂、慷慨与宏大。为此，他不受大众的指责或恶评，也不屈从于意见或奉承。但在忍受劳苦并同时履行任何义务时，以及在他对一切境遇所展现的男子气概的超越中，他在其余人类中真正显出是一个\OriginalTerm{男人}{\GreekText{ἀνήρ}}。另一方面，他保持智德，在灵魂的平静中实行节制；接受所命令之事，如同接受属于自己之物，厌恶卑下之事，如同与自己相异之物；他变得端庄而超尘脱俗，端庄有序地行事，在各方面、任何事上都不越矩。他极其富有，无所贪求，因为少有需要；藉着对善的知识，他处于一切美善的丰盛之中。因为他正义的首要效果，是喜爱与自己的同类——无论在地上还是天上——共度时光和交往。同样，他也慷慨地施与自己所拥有的一切。他既是爱人的，也就憎恨恶人，对一切邪恶怀有完全的厌恶。因此他必须学习对自己和邻人忠实，并服从诫命。因为真正天主的仆人，是那自动服从祂命令的人。而那已不是藉着诫命，而是藉着知识本身而心地纯洁的人，是天主的朋友。因为我们并非生而自然拥有德行，也非出生后它如身体的某些部分一样自然成长；因为那样它既非出于自愿，也无可赞许。德行也不像言语，藉由日常生活事件的实践而完善（因为恶习恰恰多半由此产生）。因为知识不是藉由任何技艺——无论是牟利之技或与身体保健相关的技艺——而获得的。也不是来自教学课程。因为那只要能预备和磨练灵魂，就已足够了。因为国家的法律或许能够约束恶行；但仅触及表面的劝导之辞，无法产生真理的科学性恒久。\par

如今，希腊哲学在某种意义上净化灵魂，并预先预备它接受信德，真理在此基础上建造知识的殿堂。\par

这才是真正的运动员——他在那巨大的竞技场，即这美丽的世界中，因对一切情欲的真正胜利而获得桂冠。因为制定比赛的是全能的天主，而颁发奖品的是天主的独生子。天使和诸神是观众；这比赛，包含各种操练，是\BibleQuote{不是对抗血和肉}\BibleRef{厄 6:12}，而是对抗藉肉体而运作的无序情欲的属灵势力。谁在这些挣扎中获得制胜，并推翻那如同以某些比赛相威胁的诱惑者，谁就赢得永生。因为天主至公义判决的判决是无误的。观众被召集来观看比赛，运动员在竞技场上较量；那服从教练指导的人获胜。因为天主为所有人提出的赏报都是同等的；祂自己无可指责。那拥有能力者得到仁慈，那运用意志者是有力的。\par

同样，我们接受了理智，好使我们知道我们所作之事。而\InnerQuote{认识你自己}这句格言在此意为认识我们为何而生。我们生来是为了服从诫命，如果我们愿意选择得救。这便是那无法逃避天主的报应。那么人的本分就是服从天主，祂已藉着诫命宣告了丰富的救恩。而宣认就是感恩。因为行善者首先开始行善。那在适当考虑下立即接受并持守诫命的人，是\OriginalTerm{忠信的}{\GreekText{πιστός}}；那以爱尽可能回报恩惠的人，已成为朋友。人方面一个最重要的回报，就是做天主喜悦的事。作为祂自己的作品和与祂相似的成果，导师与救主将人方面的援助与促进视为个人的恩惠与光荣；同样，祂也将那些加诸相信祂之人的伤害，视为对自己忘恩与羞辱。因为还有什么羞辱能触及天主呢？因此，根本不可能给予与从主所获恩惠等价的回报。\par

正如虐待财产者侮辱物主，虐待士兵者侮辱统帅一样，对他所祝圣之人的虐待，就是对主的藐视。\par

因为，正如太阳不仅照耀天空和全世界，普照陆地和海洋，而且透过窗户和小缝隙将光线送入房屋最隐秘的角落，同样，圣言遍及各处，将祂的眼目投向生活行动中最细微的细节。\par

% structure: p0047 source=h2 target=subsection reason=relative-heading-rank; base=h2

\subsection{第四章 外邦人按自己模样造神，一切迷信由此而生}

如今，正如希腊人把诸神描绘成具有人形，同样也把他们描绘成具有人的情欲。而正如他们各自描绘的神像与自己相似，如\Person{色诺芬尼}所说：\Quote{埃塞俄比亚人是黑皮肤和猿鼻，色雷斯人是红润和黄褐}；同样，他们也使诸神的灵魂与塑造他们的人相似：例如，野蛮人把神塑造成野蛮粗野，而希腊人则塑造成更文明却仍受情欲支配。\par

因此，合理的是，恶人对天主的观念必定是恶劣的，而善人的观念则是最卓越的。所以，那在灵魂上真正具有君王气度和知识的人，既虔诚又远离迷信，他坚信唯一的天主是尊贵的、可敬畏的、庄严的、仁慈的、行善的、一切美善的创造者，而非邪恶的根源。至于希腊人的迷信，我想，在我们题为\WorkTitle{劝勉}的书中已经充分论述了，大量利用了与此相关的历史。因此，没有必要再次详细叙述已经清楚说明的内容。但是，就进入主题时必须指出的必要而言，只需从许多证据中引用少数几个，以证明那些将神性比作最坏之人的不虔敬。因为要么他们的那些神被人伤害，并且因被我们伤害而显得不如人；要么，如果不是这样，那么他们为何对那些并未伤害他们的人发怒，就像一位易怒的老妇被激怒一样？\par

据他们说，阿尔忒弥斯因奥纽斯而对埃托利亚人发怒。因为她是一位女神，难道她不明白他没有献祭并非出于轻蔑，而是由于疏忽，或是以为他已经献祭过了吗？\par

而勒托与雅典娜争论，因为后者对她因在神庙中生育而发怒，她说：\par

\Quote{你喜欢看到从死者身上剥下的\newline{}杀人的战利品，这些\newline{}对你并非不洁，但你却认为\newline{}我在此的分娩是可怕之事，\newline{}尽管其他生物在庙中生养\newline{}幼仔并无害处。}\par

因此，由于对这些易怒的神祇怀有迷信的恐惧，他们便认为所有事件都是邪恶的征兆和原因。如果老鼠穿透了泥筑的祭坛，或者因缺乏其他东西而啃破油瓶；如果一只正在被养肥的公鸡在傍晚啼叫，他们就断定这是某种征兆。\par

米南德在\WorkTitle{迷信者}中对此类人物做了喜剧性描述：\par

\Quote{甲：愿好运降临于我，尊贵的神明！\newline{}我系右鞋带时，\newline{}把它弄断了。}\newline{}\Quote{乙：很可能，愚蠢的傻瓜，\newline{}因为它已经腐烂了，而你吝啬鬼，\newline{}不肯买新的。}\par

安提丰有一个巧妙回答：当某人认为母猪吃掉了自己的小猪是不祥之兆时，他看到母猪因饲养者的吝啬而瘦弱不堪，便说：\Quote{你应该为这个征兆而庆幸：她如此饥饿，却没有吃掉你的孩子。}\par

\Quote{这有什么奇怪？}比翁说，\Quote{如果老鼠找不到东西吃，就啃口袋？}因为令人惊奇的是（正如阿尔克西劳开玩笑时所说）\Quote{口袋吃掉了老鼠。}\par

第欧根尼因此对那个因看到一条蛇缠绕在杵上而惊奇的人说得好：\Quote{别惊奇；如果你看到杵缠绕着蛇，而蛇是直挺挺的，那才更令人惊奇呢。}\par

因为无理性的动物必然奔跑、跳跃、争斗、繁殖并死亡；这些事对它们而言是自然的，对我们而言绝不可能是反常的。\par

\Quote{许多鸟儿在阳光下行走。}\par

喜剧诗人菲勒蒙在喜剧中处理了这类问题：\par

\Quote{当我看到一个人留意谁打了喷嚏，\newline{}或谁说了话，或观察谁在走路，\newline{}我立刻在市场上把他卖掉。\newline{}我们每个人走路、说话、打喷嚏，\newline{}都是为了自己，而不是为了市民：\newline{}事情按其本性发生。}\par

然后，通过节制的生活，人们寻求健康；而通过暴饮暴食和在盛宴中沉溺于饮酒，他们招致疾病。\par

也有许多人害怕立起的铭文。第欧根尼非常机智，他在一个恶人的房子里发现铭文写着：\Quote{以胜利闻名的赫拉克勒斯住在这里；勿让任何坏事进入}，便评论道：\Quote{那么这房子的主人怎么进去呢？}\par

同一群人，正如俗语所说，崇拜每一根棍子和油腻的石头，却惧怕淡黄色羊毛的簇束、盐块、火把、海葱和硫磺，被巫师用某些不洁的赎罪仪式所迷惑。但是天主，真正的天主，只承认义人的品性为圣洁——不承认不义和邪恶。\par

你可以看到那些从被净化之人身上取来的鸡蛋，如果给予必要的温暖，就会孵化。但这不可能发生，如果它们被转移了那经受净化之人的罪过。因此，喜剧诗人狄菲洛斯在喜剧中关于巫师风趣地写道：\par

\Quote{净化普洛托斯之女和她们的父\newline{}普洛托斯·阿班塔得斯，还有第五个，一个老妇，\newline{}如此多人的身体用一支火把、一颗海葱，\newline{}用硫磺和巨响大海的沥青，\newline{}来自平静流淌、深流的海洋。\newline{}但福佑的空气穿过云层送来安提库拉，\newline{}好让我把这臭虫变成雄蜂。}\par

因为米南德说得好：\par

\Quote{菲迪亚斯啊，如果你真有病，\newline{}你必须为它寻找真正的治疗；\newline{}但现在并非如此。对于不真实的病\newline{}我找到了不真实的治疗。相信它\newline{}会对你有益。让妇女们围成一圈\newline{}为你擦拭，并从三处泉源取水。\newline{}加上盐和扁豆，然后洒在你身上。\newline{}每个人都是纯洁的，如果自觉无罪。}\par

例如，悲剧说：\par

\Emph{墨涅拉俄斯}。\Quote{俄瑞斯忒斯，是什么疾病在毁灭你？}\newline{}\Emph{俄瑞斯忒斯}。\Quote{良心。因为我做了可怕的事。}\par

因为实际上，没有别的纯洁，只有远离罪恶。因此埃庇卡摩斯说得极好：\par

\Quote{若你有一颗纯洁的心，\newline{}你的整个身体就是纯洁的。}\par

现在我们也说，需要用正确的理性来净化灵魂中腐败和邪恶的教义，然后据此记忆教义的主要要点。因为在传授奥秘之前，他们认为有必要对即将接受入门者进行某些净化；同样，人也必须放弃不敬神的意见，从而转向真正的传统。\par

% structure: p0075 source=h2 target=subsection reason=relative-heading-rank; base=h2

\subsection{第五章　圣洁的灵魂是比任何人造建筑更卓越的圣殿}

我们不是正确地、真实地不将那不可限定的限定在任何地方吗？我们不是不将包含万物的那一位关闭在人手所造的殿宇中吗？建筑者、石匠和机械工艺的作品怎能成为圣的呢？比这些更高明的，岂不是那些认为空气、包围的空间，甚至整个世界和宇宙，都适合天主的卓越的人吗？\par

确如哲学家们自己所说，人，作为天主的玩物，去制造天主，而天主成为艺术的玩物，这实在是荒谬的；因为所造之物与所造之材料相同，如象牙制品是象牙，金制品是金。现在，机械工匠所造的神像和殿宇是用无生命的材料制成的；因此它们也是无生命的、物质的、凡俗的；如果你完善技艺，它们仍带有机械的粗糙。艺术作品因此不能是神圣的、属神的。\par

有什么能被定位呢？既然没有东西不是被定位的？因为万物都在一个地方。而那被定位的，先前是不被定位的，被某物定位。那么，如果天主被人定位，祂曾不被定位，并且根本不存在。因为不存在者就是不被定位的；因为凡不存在的都不被定位。而存在者不能被不存在者定位；也不能被另一实体定位。因为它也是一个实体。由此推出，它必须由自己定位。那么，一物怎能产生自己呢？或者，那存在者怎能将自身置于存在中呢？无论它先前不被定位，是否已经自己定位？但它是没有存在的；因为不存在者是不被定位的。而假定它被定位，它后来怎能使自己成为先前所是呢？\par

但是，万有属于祂的那一位，怎能需要什么东西呢？如果天主具有人的形状，祂就需要与人类同样需要食物、住所、房屋和随之而来的事物。那些在形状和情感上相似的事物，将需要类似的供养。如果\OriginalTerm{神圣}{\GreekText{τὸ} \GreekText{ἱερόν}}有双重应用，既指天主自己，也指为荣耀祂而建造的结构，那么，我们怎能不恰当地称教会为圣的呢？教会借着真知，为天主的荣耀而造，是天主的\OriginalTerm{神圣}{\GreekText{ἱερόν}}之所，具有巨大价值，不是由机械艺术建造，也不是由骗子的手装饰，而是由天主的旨意造型成为一座殿宇。因为我现在所称的教会，不是地方，而是被选者的集合。这座殿宇更适合接纳天主尊严的伟大。因为那具有高价值的活物，被那值得一切，或者更确切地说，没有等价物的那一位所圣化，凭借后者超越的圣洁。如今，这活物就是\OriginalTerm{诺斯替者}{Gnostic}，是极有价值、被天主所尊荣、天主内居于其中的那位，即关于天主的真知被祝圣。在此，我们也将发现神圣的相似和圣洁的肖像在正义的灵魂中，当它因被净化并施行有福的行为而蒙福时。在此，我们将发现那被定位的和那正在被定位的——前者存在于已经是诺斯替者的人中，后者存在于那些能够成为诺斯替者但尚不配接受关于天主的真知的人中。因为每个注定要相信的人，在天主眼中已经是忠信的，并为祂的荣耀而被树立，是一个被赋予德行的肖像，奉献给天主。\par

% structure: p0080 source=h2 target=subsection reason=relative-heading-rank; base=h2

\subsection{第六章 发自纯洁心灵的祈祷和赞美，不断献上，远胜于祭品}

因此，既然天主不被地方所限定，也从不以活物的形状被代表；同样，祂也没有类似的情感，也没有像受造物那样的需要，以致于因饥饿而渴望祭品作为食物。那些被情感影响的受造物都是有死的。而将食物带给不吸收营养者是无用的。\par

那位喜剧诗人\Person{菲勒克拉底}（Pherecrates），在\WorkTitle{《逃亡者》}中，戏谑地描写众神自己指责人们关于他们的神圣礼仪：——\par

\Quote{当你向众神献祭，\newline{}选择我们应得的一份，\newline{}说来可耻，却不拿取，\newline{}并将两条大腿，直到腹股沟刮净，\newline{}腰部完全裸露，还有脊骨，\newline{}也像用锉刀一样刮干净，\newline{}你们吞下它们，把剩余的部分\newline{}给我们，如同给狗？然后，\newline{}仿佛彼此羞愧，\newline{}用成堆的盐渍大麦掩盖。}\par

另一位喜剧诗人\Person{欧布罗斯}（Eubulus）这样论到祭品：——\par

\Quote{但向众神献上的只是尾巴\newline{}和腿，如同向鸡奸者献祭。}\par

并在\WorkTitle{《塞墨勒》}中引入\Person{狄俄尼索斯}（Dionysus），描写他争论道：——\par

\Quote{首先，如果他们向我献祭，有人\newline{}献上血、膀胱，而非心\newline{}或网膜。因为我从不吃任何\newline{}比大腿更甜的肉。}\par

\Person{米南德}（Menander）写道：——\par

\Quote{腰子的末端、\newline{}苦胆、不可食用的骨头，他们摆\newline{}在众神面前；其余的自己享用。}\par

因为野兽不是避开全燔祭的气味吗？如果这气味真是希腊众神的报酬，他们难道不应该首先神化厨师，厨师被赋予同等的幸福，并敬拜烟囱本身，因为它更靠近那备受珍视的气味吗？\par

\Person{赫西俄德}（Hesiod）说，\Person{宙斯}（Zeus）在分肉时被\Person{普罗米修斯}（Prometheus）欺骗，得到了用狡猾技艺藏在闪闪发光的脂肪中的白色牛骨：——\par

\Quote{因此人类的部族\newline{}在祭台上焚烧祭品的白色骨头，献给不朽的众神。}\par

但是，他们决不会说神性因匮乏产生的欲望而衰弱，以致被滋养。因此，他们将祂描绘成如同植物和无洞的野兽一样，没有欲望而被滋养。他们说这些生物无害地生长，要么靠空气中的密度滋养，要么靠自身身体的呼气滋养。尽管神性虽一无所缺，但按照他们所说被滋养，那么祂既一无所缺，又何需食物呢？但如果祂本性一无所缺，却喜悦被尊荣，那么我们在祈祷中尊荣天主并非没有理由；这样，我们就带着正义献上最好最圣洁的祭品，将其作为礼物献给最正义的\OriginalTerm{圣言}{Word}，我们借着祂接受真知，藉着祂为我们所学的而归荣耀于祂。\par

因此，在我们这里，地上的祭台就是那些献身于祈祷的人的聚会，他们仿佛拥有一个共同的声音和一个共同的心思。\par

然而，如果通过鼻孔吸入的营养物比通过口腔摄入的更为神圣，那么它们仍暗示着呼吸。那么，他们关于天主说了什么呢？祂是否像橡树一族那样呼气？或者祂是否像水生动物那样，通过鳃的扩张而只吸气？或者祂是否像昆虫那样，通过翅膀压缩节段而全方位呼吸？但没有一个神志清醒的人会将天主比作这些中的任何一个。\par

那些通过肺向胸腔扩张而呼吸的生物吸入空气。那么，如果他们给天主赋予内脏、动脉、静脉、神经和肢体，他们就会使祂与人毫无区别。\par

现在，\OriginalTerm{同呼吸}{\GreekText{σύμπνοια}} 正确地用于教会。因为教会的祭献是从圣洁灵魂发出的如馨香之气的言语，同时祭献和整个心思都向天主敞开。他们颂扬\Place{提洛}岛上那非常古老的祭台为神圣；据说只有这座祭台，未被屠杀和死亡玷污，\Person{毕达哥拉斯}才靠近它。那么，当我们说正义的灵魂是真正神圣的祭台，从它升起的馨香是圣洁的祈祷时，他们难道不会相信我们吗？但我相信，祭献是人类为了给食肉找借口而发明的。但若没有这种偶像崇拜，凡愿意的人本可以分食肉类。\par

因为律法中的祭献象征性地表达了我们所实践的虔诚，正如为罪献上的斑鸠和雏鸽表明，灵魂非理性部分的净化是蒙天主悦纳的。但如果任何一个义人没有因食肉而加重灵魂的负担，他便拥有了理性理由的优势，而非像\Person{毕达哥拉斯}及其追随者所梦想的灵魂转世之说。\par

现在，\Person{色诺克拉底}在其独自论述\Quote{来自动物的食物}时，以及\Person{波勒莫}在其著作\WorkTitle{论合乎自然的生活}中，似乎明确地说动物性食物是不健康的，因为它已经被加工并同化到非理性生物的灵魂中。\par

同样，犹太人也特别因猪的不洁而禁食猪肉；因为猪比其他动物更会拱地并毁坏地里的出产。但如果他们说动物被赐予了人——我们同意这一点——但并非完全为了食物。也不是所有动物，而是那些不劳作的。因此，喜剧诗人\Person{柏拉图}在戏剧\WorkTitle{宴席}中说得不错：\par

\Quote{因为四足动物，我们今后不应杀害\newline{}任何东西，除了猪。因为这些动物有\newline{}最可口的肉；而猪身上对我们毫无用处，\newline{}除了刚毛、泥巴和噪音。}\par

因此\Person{伊索}说得不错：\Quote{猪被拖走时大声尖叫，因为它们知道自己除了用于祭献以外毫无用处。}\par

因此\Person{克莱安西斯}也说：\Quote{它们有灵魂代替盐，}以免肉质腐烂。有些人吃它们因为它们无用，另一些人则因为它们毁坏果实。还有些人不吃它们，因为这种动物有强烈的感官倾向。\par

因此，律法除了在放逐罪恶的单一情况下，并不祭献山羊；因为快乐是邪恶的根源。据说吃山羊肉会导致癫痫也是贴切的。他们说猪肉能带来最大的增力。因此它对锻炼身体的人有益；但对致力于灵魂发展的人则不然，因为食肉会导致迟钝。或许有些诺斯替主义者会为了操练而禁食肉类，并为了使肉体不致在情欲上放纵而禁食。\Quote{因为酒，} \Person{安德罗基德斯}说，\Quote{和贪食的肉餐使身体强壮，但使灵魂更加迟钝。}因此，为了清明的理解，此类食物应当被拒绝。\par

因此，埃及人在他们中间施行的洁净礼中，不允许祭司食用肉类；但他们食用鸡肉，认为它最轻；他们不碰鱼，因为某些神话，尤其因为这类食物使肉质松软。但现在，陆地动物和鸟类呼吸与我们生命之气相同的空气，拥有与空气同源的生命原则。但据说鱼类不呼吸这种空气，而是呼吸在最初创造时与水和其余元素混合在一起的空气，这也是物质持久性的一个标志。\par

因此，我们应当向天主献上并非昂贵的祭品，而是祂所喜爱的。律法中所提及的那种调和香料，是由许多舌头和声音在祈祷中组成的，或者更确切地说，是由不同的民族和本性，通过在经济中赐予的恩赐为\Quote{信德的合一}而预备的，并在赞美中聚集在一起，带着纯洁的心思、正义和正直的行为，源于圣洁的工作和正义的祈祷。因为用优雅的诗歌语言来说——\par

\Quote{谁是如此大的愚人，在人类中\newline{}如此轻信，竟以为\newline{}众神以无肉的骨头和胆汁的诡计\newline{}全部焚烧，连饥饿的狗都不配吃，\newline{}就会喜悦，并以此作为奖赏，\newline{}向这样对待他们的人施恩惠，}\par

即使他们碰巧是暴君和强盗？\par

但我们要说，火圣化的不是肉，而是有罪的灵魂；我们指的并非那吞噬一切的世俗之火，而是智慧之火，它穿过火焰而渗透灵魂。\par

% structure: p0110 source=h2 target=subsection reason=relative-heading-rank; base=h2

\subsection{第七章　诺斯替主义者采用何种祈祷，以及它如何被天主垂听}

如今我们受命要尊敬并荣耀同一位，深信祂是圣言、救主和领袖，并且藉着祂而尊敬圣父，并非像某些人那样在特定的日子，而是在我们整个生命中、在一切方式上持续这样做。的确，那因诫命而成义的蒙选之族说：\BibleQuote{一日七次我赞美了祢。}因此，诺斯替信徒并非在特定的地点、或被选定的圣殿、或某些节庆和规定的日子，而是在他整个生命中，在每一处，即使他独自一人，并且在他有任何同样信德之人的地方，他都荣耀天主，也就是说，他承认对认识生活之道的感恩。\par

若一个好人的临在，因他所激发的尊敬与敬畏，常常改善与他交往的人，那么，那位藉着知识、生活和感恩而与天主持续不间断交流的人，岂不更是在每一步中，在行为、言语和心性上全面超越自己吗？这样的人深信天主常在他身边，并不认为天主被限制在某些有限的地方；因此他不会以为有时天主不在他身边，而放纵自己于夜间的过度行为。\par

这样，我们整个生命都在庆节中，深信天主全然处处临在，我们耕种田地，赞美；我们航行海上，咏唱；我们其余的一切交往都按规矩行事。因此，诺斯替信徒与天主极为亲近，在一切事上既庄重又喜乐——庄重，因其灵魂倾向于神性；喜乐，因其思虑天主所赐予我们的人性福泽。\par

如今，知识的卓越性显然由先知提出，他说：\BibleQuote{仁慈、训诲和知识教导我，}藉着层层递进，颂扬成全的至高地位。\par

因此，他是真正王者般的人；他是天主的神圣大司祭。即使在今天，在最具智慧的蛮族人中，也可见到将司祭阶层晋升至王权的情况。因此，他从不把自己交给那在剧场中统治一切的乌合之众，甚至在梦中也不接纳那些为取悦感官而说、而行、而看的事物，因此也不接纳视觉的愉悦，也不接纳在其他享乐中发现的种种快乐，如昂贵的香料和气味（它们迷惑鼻孔），或肉食的烹调以及各种美酒的沉溺（它们诱捕味觉），或多种花朵的芬芳花束（它们通过感官使灵魂柔化）。但他常将一切事物的庄重享受溯源至天主，将食物、饮料和香膏的初果献给万有的赐予者，藉着所赐予他的圣言，在礼物和使用中表达感恩。他很少参加各类混杂的宴会，除非朋友般和谐宴会的邀请促使他去。因为他深信天主知道并洞察一切——不仅是言语，还有思想；因为我们的听觉，通过身体的通道，其理解力并非通过肉体的力量，而是通过一种心理的感知和辨别有意义的声音的理智。因此，天主不具有人形去听；也不需要感官，如斯多葛派所断定的，\Quote{特别是听觉和视觉；否则祂永不能理解。}然而，空气的敏感性、天使极其敏锐的知觉，以及那以不可言喻的力量、无需可感听觉而触及灵魂意识的能力，在思想产生的瞬间知道一切。若有人说声音不能到达天主，而是向下滚动在空气中，那么圣人们的思想不仅穿透空气，也穿透整个世界。而神圣的力量，以光的速度，看透整个灵魂。看！难道意志不也向天主说话，发出它们的声音吗？它们岂不是藉着良心传达？祂，按照自己的旨意，在出生之前就已认识蒙选者，知道将要存在的事如同已存在，还要等待什么声音呢？力量之光不是照彻整个灵魂的深处吗？圣经说，\BibleQuote{知识的灯}，探察\BibleQuote{隐秘之处}？天主全是耳朵，全是眼睛，若可允许我们这样表达。\par

总的来说，对天主不配的见解，无论在圣歌、论述、著作或教义中，都不保存任何虔诚，反而引向卑下和不雅的思想与观念。因此，由于大众对真理的无知，他们的赞扬与责备并无区别。那么，愿望和渴望的对象，简言之，心灵的冲动的对象，就是祈祷的主题。因此，没有人渴望一杯饮料，而是渴望喝可饮之物；没有人渴望继承，而是渴望继承。同样，没有人渴望知识，而是渴望知道；没有人渴望正当的治理，而是渴望参与治理。我们祈祷的对象，就是我们所请求的对象，请求的对象就是渴望的对象。因此，祈祷与渴望按次序紧随，旨在拥有所提供的祝福与利益。\par

那么，那真正拥有这般特质的诺斯替信徒，他发出祈祷和请求，为求那真正属于灵魂的美善之物；他祈祷时，自己也贡献努力以达致良善的习性，以致不再将美好事物当作某种属于他的课程，而成为善本身。\par

因此，这类人尤其应当祈祷，因为他们正确认识神性，并拥有与之相称的道德卓越；他们知道什么是真正美好的，什么是应当祈求的，以及在何种情况下、以何种方式为每个个案祈祷。向那些并非神明的存在祈求，好像他们是神明，是极端愚蠢的；或者祈求那些无益之物，在美好事物的外表下为自己乞求祸害。\par

因此，义理宜然，既然只有一位美善的天主，有些美善的事物唯独由祂赐下，有些则保存不变，我们与天使都这样祈祷。但并非同样地。因为祈求恩赐保存不变，与初次努力获得恩赐，并非同一回事。\par

避免灾祸是祈祷的一种；但此种祈祷决不可用于伤害人，除非诺斯替者在专注于正义时，对那麻木不仁者运用此祈求。\par

那么，更勇敢地说，祈祷是与天主的交谈。因此，即使低声细语，不启唇，我们静默而言，却内里呼喊。因为天主不断倾听一切内在的交谈。同样，我们举目、举手向天，并在祈祷结束时移动脚步，随从那朝向理智本质的灵性热忱；竭力将身体与话语一同从大地抽象出来，提升灵魂，以渴望更美善的事物为翼，强制它前进至圣洁之地，慷慨地轻视肉体的锁链。因为我们深知，诺斯替者甘愿跨越整个世界，正如犹太人确实跨越了埃及，首要地清楚表明，他将尽可能接近天主。\par

现在，若有人指定固定的祈祷时辰——例如第三、第六和第九时辰——然而诺斯替者一生都在祈祷，藉祈祷努力与天主共融。简言之，一旦达至此境界，他留下一切无用之物，因为已获得那以爱德行事之人的完美。但将时辰三等分、以同样多次祈祷尊崇之，那些认识圣洁居所之蒙福三元组者明白。\par

至此，我忆起某些异端者——即普罗迪科斯异端的追随者——所传入的关于无需祈祷的教义。为使它们不致因这不虔敬的智慧（仿佛新奇）而自高自大，让它们知道这早已被称作昔兰尼学派的哲学家所采纳。尽管如此，那些假称［诺斯替者］之人的不圣知识（\OriginalTerm{诺斯}{\GreekText{γνῶσις}}）将在适当时机遭到驳斥；这样，对其的攻讦虽不简短，却不会因插入注释而打断当前的论述——我们在此证明，唯一真正圣洁和虔诚的人，是那按照教会规范真正为诺斯替者的人；唯独他，在祈求并思想时，获准按天主意愿所作的祈求。因为正如天主能行一切祂所愿，同样诺斯替者获得一切他所求。因为普遍而言，天主知道谁配得美善之物，谁不配得；因此祂赐予各人适宜之物。故此，对那些不配者，虽屡次祈求，祂也不赐予；但对那些配得者，祂必赐予。\par

祈求并非多余，尽管美善之物无需功劳便赐下。\par

感谢并为邻人的悔改祈求，是诺斯替者的职责；正如主祈祷时，为祂职务的完成而感谢，祈求尽可能多的人获得知识；为使在得救者中，藉着救恩，藉着知识，天主得光荣，那唯一美善、唯一救主者，藉圣子从世世代代被承认。同样，相信将获得，也是一种以诺斯替方式贮藏的祈祷。\par

但是如果任何与天主交谈的机会成为祈祷，那么就不应错过任何亲近天主的时机。无疑，诺斯替者的圣洁，与天主的福佑眷顾结合，在自愿的宣认中彰显天主完美的恩惠。因为诺斯替者的圣洁，与天主之友的相互善意，是眷顾的一种对应运动。因为天主并非不自愿地美善，像火是热的那样；但在祂那里，施予美善之物是自愿的，即使祂事先接受了祈求。被救者也不会违背其意志而得救，因为他不是无生命的；但他将首先自愿并自由选择奔赴救恩。因此，人也接受了诫命，以便他能自我推动，趋向他愿意选择或避免的一切。因此，天主并非出于必然而行善，而是从祂的自由意志恩赐那些自愿回转的人。因为从天主延伸至我们的眷顾并非服侍性的，如同从下级到上级的服务。但因怜悯我们的软弱，持续的眷顾护理运作，如同牧人对羊群的照顾，君王对其臣民的照顾；我们自己也顺从地对待那些上级——他们按照从天主所受的委托管理我们。\par

因此，那些献上最自由、最像君王的服务——这服务出于虔敬之心和知识——就是神性的仆人和侍从。那么，每个我们怀有关于天主之想法的地点和时间，实际上都是神圣的。\par

因此，当那选择正义、同时怀有感恩之心的人在祈祷中提出祈求时，他有助于获得它，在祈祷中欢喜地抓住所渴望的事物。因为当美善之物的施予者察觉到我们这方的接受性时，一切美善之物随即跟随其观念而来。确实，在祈祷中，人的品格被筛选，看其对义务持何态度。\par

但是，若声音和表达赐予我们是为了理解，天主怎能不听见灵魂本身和思想呢？因为灵魂确实听见灵魂，思想听见思想。因此，天主不等待多言的舌头，如同人间的翻译者，而是绝对知道所有人的思想；那声音向我们传达的，就是我们的思想——那在创造之前祂已知道会进入我们心中的思想——向天主说话。因此，祈祷可以不用声音，将整个灵性本性集中于思想的内在表达，以不分心的转向天主。\par

既然黎明是诞生之日的影像，从那时起，最初从黑暗中照耀出来的光逐渐增强，那么，对那些陷于黑暗中的人，也升起了一个认识真理的日子。与太阳升起的方式相应，祈祷面向东方的日出。因此，最古老的庙宇也面向西方，好教导人们在面对图像时转向东方。\BibleQuote{愿我的祈祷如馨香上达于祢，我高举的双手如晚祭，}圣咏如此说。\par

因此，恶人的祈祷最为有害，不仅对他人如此，对他们自己也是如此。倘若他们祈求并得到他们所谓的幸运，这些幸运在得到后反而伤害他们，因为他们不知道如何使用。他们祈求拥有所没有的，所求的似乎是好事，其实不是。但诺斯替者祈求他所拥有之物的持久、对将要发生之事的适应、以及他将要接受之事的永恒。他祈求那些真正美好的事物，即灵魂的事物，能属于他并常与他同在。因此他不渴求任何缺失之物，满足于现有的一切。因为他并不缺乏属于他自己的美善，藉着天主的恩宠和认识，他已自足。既已自足，他便不再需要其他事物。他知道至高的旨意，祈祷时即已拥有，与全能的力量密切相联，因着无限的爱渴慕成为属神的人，与圣神结合。\par

因此，他心胸宽大，藉着认识拥有最宝贵、最卓越的事物，勤于默观，在灵魂中保持默观对象的持久动力，即认识之敏锐的洞察力。他尽最大努力获得这种能力，通过控制所有攻击心灵的影响，通过不间断地致力于默观，通过克己训练、节制享乐、行事端正，此外还在研究和生活中取得丰富经验，因而拥有言说的自由，不是饶舌之舌的能力，而是使用明白言语的能力，不因偏爱或恐惧隐瞒任何值得在适当时间、极有必要说出的话语。他既从真理本身的奥秘队伍中接受了关于天主的事物，便使用与美德价值相符的语言来强调其伟大，并以祈祷之灵感的提升展示其成果，尽可能以诺斯替的方式与理智和属神的事物相通。\par

因此，他总是温和良善，平易近人，和蔼可亲，坚忍不拔，心怀感恩，拥有良好的良心。这样的人刚正不阿，不仅不会腐化，而且不会受诱惑。因为他从不使自己的灵魂屈服于快乐与痛苦，或被它们俘获。若作为审判者的圣言呼召，他变得坚定不移，丝毫不放纵情欲，公正地走向他所该去之处；他深信一切事情都安排得尽善尽美，选择美德之灵魂的进步走向更善，直到达至善本身，即圣父的前庭，靠近伟大的大司祭。这就是我们的诺斯替者：忠信，深信宇宙之事以最美善的方式管理。他尤其对一切所发生之事深感满意。因此，他合情合理地不求生活中所需用的任何事物，深信知晓万物的天主会赐予善人一切有益之物，即使他们不祈求。\par

因为我认为，正如一切事物按技艺的规则供应给有技艺的人，按外邦人的方式供应给外邦人，同样，一切事物也按诺斯替的方式供应给诺斯替者。从外邦人中皈依的人将祈求信德，而上升至认识的人将祈求爱的成全。已达至顶峰的诺斯替者将祈求默观增长并持久，如同普通人祈求持续的健康。\par

不仅如此，他将祈求永不离弃美德，尽最大的努力合作以成为无误者。因为他知道有些天使由于疏忽而被摔到地上，尚未完全达到合一的状态，未能摆脱对双重性的倾向。\par

但对于那由此训练自己达至认识之顶峰和完美之人崇高境界的人，一切时空之事都有助于他，既然他已选择无误地生活，并接受训练以获得认识的稳定。然而，在那些仍有沉重角落向下倾斜的人身上，即使用信德提升的部分也会被拖下。因此，在藉诺斯替训练获得不可丧失的美德者身上，习性成为本性。正如石头中的重量，此人的认识也无可丧失。不是无端地，而是藉着意志的运用、理性的力量、认识和天主的眷顾，才变得不可丧失。藉着谨慎，它变得不可丧失。他将谨慎以避免犯罪，思虑以防失落美德。\par

认识似乎产生思虑，因为它教导人认识有助于美德持久的事物。因此，最高的就是认识天主；也正因此，美德被保存得不可丧失。认识天主的人是圣洁和虔诚的。因此，我们已证明诺斯替者是唯一的虔诚者。\par

他以現有之善而喜樂，並因那應許之事——彷彿已然臨在——而歡欣。因它們並未如尚缺般逃過他的注意，因為他預先知道它們是何等樣式。既因知識而確信每一未來之事將如何，他便擁有它。因為匱乏與欠缺是以與人相關之物來衡量的。那麼，若他擁有智慧，而智慧是神聖之物，分享那無所欠缺者的人自身也將無所欠缺。因為智慧的傳授並非通過主動與被動互相觸動與停止，亦非通過任何事物被抽離或變得欠缺。因此，主動在傳遞中顯為未減損。所以，我們的\Quote{知者}盡可能擁有一切美善，但不在數量上等同；否則他便無法經由適當的靈性進展階段和管理的行動而改變其位置。\par

天主幫助他，以更親近的眷顧使之尊榮。因為難道萬物不是為善人的緣故，為他們的擁有和益處，或者說為救恩而被造的嗎？他不會剝奪那為德行而存在之物，不給予那些為其被造的人。因為顯然，為尊崇他們卓越的本性和神聖的選擇，他激勵那些已選擇善生的人，賜予他們完成其餘救恩的力量；勸勉一些人，幫助另一些自己已成為配得的人。因為一切美善都能在\Quote{知者}身上產生；如果確實他的目標是明智地認識和實行一切。如同醫生為那些與他合作以達健康的人提供健康，同樣，天主為那些合作以求得知識和善行的人賜予永恆的救恩；既然誡命所吩咐的——連同其實行——都在我們的能力之內，那麼應許便得以實現。\par

在我看來，以下由希臘人所說的話似乎極好。一位在古人中頗負盛名的運動員，為了獲得男子氣概的力量，長期使身體受嚴格訓練，之後去奧林匹亞賽會，他注視著比薩的宙斯雕像，說：\Quote{宙斯啊，若我已為比賽做好一切必要準備，請賜我勝利，正如所應當的。}因為同樣，在\Quote{知者}身上，他無可指責地、憑著良心履行了所有依賴於他的事——在學習、訓練、善行和悅樂天主方面——這一切都共同促成救恩達至圓滿。因此，從我們這裡要求的是那些在我們能力之內的事，以及那些與我們相關——無論現在或將來——的選擇、願望、擁有、使用和持守。\par

因此，與天主往來的人必須擁有純潔無瑕、毫無玷污的靈魂，因為使自己完全成為良善是必要的。\par

他也應以溫和之心與善人一同祈禱。因為參與他人的罪是危險的事。因此，\Quote{知者}將與那些較新近相信的人一同祈禱，為那些他們有責任共同行動的事項祈禱。他的整個生命是一場神聖的慶節。他的祭獻是祈禱、讚頌、餐前誦讀聖經、餐間和睡前吟唱聖詠與讚美詩，夜間還有祈禱。藉這一切，他因不斷的紀念而與神聖的歌詠團合一，專注於具有永恆回憶的默觀。\par

再者，他豈不也知道另一種祭獻，即以教義和金錢施予需要者嗎？當然知道。但他不用口舌作冗長的祈禱，因他已學會向主祈求所需。因此，他在每一處祈禱，但並非公開地、有目共睹地對著群眾。然而在走路、交談、沉默、閱讀和按理性工作時，他每種心境都在祈禱。他若只在靈魂隱密處形成思想，並\Quote{以無可言喻的歎息}\Quote{\BibleRef{羅 8:26}}呼求父，祂便在近旁，在他說話時就在他身邊。既然一切行為只有三個目的，他為其卓越和有用而做每一事；但為快樂而做任何事，他留給那些追求世俗生活的人。\par

% structure: p0144 source=h2 target=subsection reason=relative-heading-rank; base=h2

\subsection{第八章　\Quote{知者}如此忠於真理，以致無需使用誓言}

在這般虔敬中經過考驗的人，遠離說謊和起誓的傾向。因為誓言是一種決定性的肯定，帶著呼喚天主之名。因為那已經一次忠信的人，怎能顯得不忠信而需要誓言？又怎能不讓他的生命成為一個確實而決定性的誓言？他生活、行走，並以堅定而確實的生活和言語展示他肯定之可靠。若錯誤在於行為者或說話者的判斷，而不在於受冤者的受苦，他既不會說謊也不會發假誓以致冤枉天主，因為他知道天主的本性是不能受損的。他也不會為鄰人的緣故——儘管他學會了愛鄰人，即使不親密——而說謊或犯任何過犯。更何況，他也不會為自己的緣故說謊或發假誓，因為他絕不會故意對自己做錯事。\par

他甚至不起誓，寧可肯定時用\Quote{是}，否定時用\Quote{否}來表示斷言。因為起誓就是發誓，或從心中提出某事作為誓言式的確認。對他而言，加上\Quote{我實話實說}這句話，作為對那些不能感知他回答之確定性的人的確認，就已經足夠。因為我認為，他應當保持一種能向外人激發信任的生活，以致連誓言也不會被要求；而對他自己和那些與他交往的人，則保持善意，這是自願的正義。\par

\OriginalTerm{诺斯替信徒}{Gnostic}真正地起誓，但不轻易起誓，很少使用誓言，正如我们所说。他因着与真理的相合而说实话。因此，说实话的誓言是在职责正确中产生的。对于按着真理的极致生活的人，哪里还有发誓的必要呢？因此，不起誓的人远未犯虚誓之罪。而不违犯通过契约所确认之事的人，将永不发誓；因为违犯和履行是由行为来确认的；正如在肯定和发誓中撒谎和虚誓是违反职责的。但公正生活、不违犯任何职责的人，当真理的审判被审视时，他藉着自己的行为说实话。因此，对于他而言，舌头的见证是多余的。\par

因此，他始终深信天主无所不在，且不惧怕说真话，并知道说谎与他不相配，他就满足于天主的知觉和他自己的知觉。因此他不说谎，也不做任何违背契约的事。因此，即使被要求发誓，他也不发誓；他也永不否认以致说假话，即使他死于酷刑。\par

% structure: p0149 source=h2 target=subsection reason=relative-heading-rank; base=h2

\subsection{第九章　教导他人者当在德行上超卓}

由那个承担教导他人首位、领受了藉言语和行为实施世间至善之职务的人，诺斯替信徒的尊威得以增强和加增；他凭此中介与神性接触和共融。正如那些敬拜地上之物的人向它们祈祷，仿佛它们能听见，在它们面前确认契约；同样，在活生生的肖像——人身上，\OriginalTerm{圣言}{Word}的真正威严被可信的导师所接纳；而向他们所行的善举被上达到\OriginalTerm{主}{Lord}，那位真正的、藉着教导按祂肖像创造并谐调的人，更新那接受教导的人使之得救。因为正如希腊人因某种关系称钢铁为\Quote{阿瑞斯}，称酒为\Quote{狄俄尼索斯}；同样，诺斯替信徒视邻人的益处为自己的救恩，可被称为主的活肖像，不是就形式的特殊性而言，而是就能力的象征和宣讲的相似性而言。\par

因此，凡他心中所想的，他都放在舌上，向那些堪听的人说出，言语和生活都出于同意和倾向。因为他思想并说出真理；除非有时像医生为了病人的安全，按照\OriginalTerm{诡辩家}{Sophists}的方式，为了医治而欺骗或说不真实的话。\par

例如：高贵的\OriginalTerm{宗徒}{apostle}给\Person{弟茂德}行了割礼，尽管他大声宣告并写道，人手所做的割礼毫无益处。但为了不因一下子将那些勉强听道的希伯来人从法律拉到内心的割礼（藉着信德）而迫使她们脱离会堂，他\Quote{适应犹太人，成了犹太人，为要赢得众人。}因此，那仅仅为了邻人的益处而屈就适应的人，为了那些他为之适应之人的得救，不参与任何因嫉妒者加于义人的危险而产生的虚伪，这样的人绝不行强迫之事。但仅为了邻人的益处，他将做一些他本不会首先去做的事，除非他为了他们的缘故而行。这样的人为教会舍己，为那些他在信德中所生的门徒；作为榜样给那些能接受那富有人情和爱天主的导师之至高经济的人，为他话语真理的确认，为主爱德的操练。这样的人不为恐惧所奴役，言语真实，忍耐劳苦，永不愿以口言说谎，并常在言语中保持无罪；因为谎言带着欺骗说出，不是静止的话语，而是导致祸害。\par

因此，从各方面看，唯有诺斯替信徒在行为和言语上为真理作证。因为他总在一切事上行得正确，无论在言语、行为还是思想本身。\par

那么，粗略地说，基督徒的虔诚就是这样。如果他能按照职责和正确的理性行这些事，他就是虔诚而正义地行。若是如此，唯有诺斯替信徒才是真正虔诚、正义和敬畏天主的。\par

基督徒并非不虔诚。因为这是我们应向哲学家证明的一点；因此他绝不会以任何方式行任何恶或卑鄙之事（那是不义的）。所以，他并非不虔诚；唯独他敬畏天主，以神圣和尽责的方式敬拜真天主，那宇宙的统治者、君王和主宰，怀着真虔诚。\par

% structure: p0156 source=h2 target=subsection reason=relative-heading-rank; base=h2

\subsection{第十章　完善的步骤}

因为\OriginalTerm{知识}{gnosis}，总的说来，是作为人的人的完美化，是通过对神圣事物的认识而达成，在品格、生活和言语上，与自己及神圣的\OriginalTerm{圣言}{Word}一致和相符。因为藉着它，\OriginalTerm{信德}{faith}得以成全，因为唯独藉着它，信徒成为完美的。\OriginalTerm{信德}{faith}是一种内在的善，不寻求天主，却承认祂的存在，并荣耀祂作为存在。因此，从这信德出发，藉着它藉天主的恩宠而发展，关于祂的知识要尽可能被获得。\par

现在我们断言\OriginalTerm{知识}{gnosis}不同于\OriginalTerm{智慧}{\GreekText{σοφία}}（\OriginalTerm{智慧}{wisdom}），后者是教导的结果。因为凡知识之处，它当然是智慧；但凡是智慧之处，它却不一定是知识。因为\Quote{智慧}一词只出现在所说出的话语的知识中。\par

但是，关于天主，不是怀疑，而是相信，才是知识的基础。但基督既是基础又是上层建筑，万有之始与终藉着祂。而两端——即开始和终结——我意指\OriginalTerm{信德}{faith}和\OriginalTerm{爱德}{love}——不是可教的。但\OriginalTerm{知识}{knowledge}，通过传授从那天主恩宠的托付中来，被托付给那些显明自己堪当的人；并从它发出爱德之光，从光明到光明。因为经上说：\Quote{凡有的，还要给他；}\BibleRef{路 19:26} 对信德给予知识；对知识给予爱德；对爱德给予嗣业。\par

而这事发生，每当人藉着信德、知识和爱德悬挂于主，并与祂一同上升，到达我们信德和爱德的天主与守护者那里。因此最终（由于需要极大的准备和先前的训练，既为了听所讲的，也为了生活的安定，并为了明智地进步到超越法律的义的地步），知识才被托付给那些适合并蒙拣选的人。它引领我们达到无尽而完美的终局，预先教导我们将来要过的生活，即按照天主并与神同在的生活；当我们从因罪而受的一切惩罚和刑罚中解脱出来之后，这些惩罚是为救恩的纪律。在此救赎之后，奖赏和尊荣就归于那些已成为完美的人；当他们完成了净化，并停止了所有的服务，即使是神圣的服务和在圣者中的服务。然后，心中纯洁，接近主，等待他们的是恢复永恒的默观；他们被称为神的称号，注定要与那些由救主首先安置在位置上的其他神一同坐在宝座上。\par

因此，知识在净化方面是快速的，并适合于那蒙悦纳的向更好的转变。因此它轻松地[将灵魂]转移到与灵魂相近的、神圣和圣洁的事物，并藉着自己的光明引导人通过神秘的进阶阶段；直到它使心中纯洁的人安息在冠冕般的安息之所；教导人以知识和理解面对面地凝视天主。因为知识者的灵魂的完美就在于与主同在，在那里它直接服从于祂，超越了所有的净化和服务。\par

那么，信德可以说是对本质的全面知识；而知识是对信德所领受之事的坚强而确切的证明，建立在信德之上，藉着主的教导，将[灵魂]引向无误、科学和理解。并且，在我看来，第一个救恩性的转变是从异教到信德，如我之前所说；第二个是从信德到知识。而后者终止于爱德，此后将爱者给予被爱者，将认知者给予被认知者。并且，也许这样的人已经达到了\Quote{与天使相等}的状态。\BibleRef{路 20:36}因此，在肉身中达到最高卓越之后，不断适时地转向更好的，他飞向祖厅，经过神圣的七重[天居所]到主自己的宅邸；成为一束光，稳定而永恒地持续，完全且在每一部分都不变。\par

我们提到的第一种主的工作方式，是展示虔诚所带来的报偿。在极多的见证中，我将引述一个，由达味先知如此概括地说：\Quote{谁能上主的山，谁能站在祂的圣所？那手洁心清，不向虚妄举起灵魂，也不向邻人发假誓的人。他必从上主领受祝福，从天主他的救主领受仁慈。这就是寻求上主的世代，寻求雅各伯天主面容的人。}在我看来，先知简洁地指出了知识者。达味似乎通过称祂为\Quote{雅各伯天主的面容}而简要地证明了救主是天主，祂宣讲并教导关于圣神的事。因此，宗徒也称圣子是\OriginalTerm{肖像}{\GreekText{χαρακτῆρα}}\BibleRef{希 1:3}，祂教导关于天主的真理，并表达了全能者是独一唯一的天主和圣父，\Quote{除了圣子，没有人认识祂，除非圣子愿意启示给谁。}\BibleRef{玛 11:27}那些\Quote{寻求雅各伯天主面容的人}暗示天主是唯一的；我们的救主和天主将这位唯一的天主描述为善父。而\Quote{寻求祂的世代}是特选的族类，致力于探求知识。因此，宗徒也说：\Quote{我若不以启示、或以知识、或以先知话、或以教训向你们说话，我就对你们毫无益处。}\BibleRef{格前 14:6}\par

虽然那些不是知识者的人有时也做正确的事，但不是按照理性；例如在勇敢方面。因为有些人生来就精神高昂，后来没有理由地培养这种倾向，就冲向许多事，表现得像勇敢的人，有时甚至成功地完成同样的事；正如忍耐对于工匠来说是容易的。但并非出于同一原因或有同一目的；即使他们献出整个身体。\Quote{因为他们没有爱德，}按照宗徒的话。\BibleRef{格前 13:3}\par

因此，知识者的所有行动都是正确的行动；而没有知识的人所做的行动是错误的行动，即使他遵循一个计划；因为他不是出于思考而勇敢行动，也不是将他的行动导向那些从美德到美德的事情，以达到任何有益的目的。\par

其他美德也是如此。同样，类比也保留在宗教中。因此，我们的知识者不仅在圣德方面如此；与知识的虔诚相对应的是关于生活其余行为的诫命。因为我们现在的目的是描述知识者的生活，而不是提出教义体系，我们将在适当的时候解释后者，保留主题的顺序。\par

% structure: p0167 source=h2 target=subsection reason=relative-heading-rank; base=h2

\subsection{第十一章 知识者的生活描述}

关于宇宙，他由于接受了神圣教导而真实而宏伟地思考。于是，始于对受造界的赞叹，并亲身证明了自己接受知识的能力，他成为主的一个顺服的门徒。一听到天主和眷顾，他便因所怀的赞叹而相信。藉由此产生的冲动之力，他全力投入学习，做一切使他能获得所渴望知识的事。随着信德前进，融合了探求的渴望便产生。这样，他便配得上默观，具有如此品质和如此重要性。如此，真知者将品尝天主的旨意。因为他不是将耳朵，而是将灵魂交托给言谈所意指的事物。因此，他藉由言语领会本质和事物，恰当地将灵魂带入本质；例如，以向真知者讲述的独特方式领会诫命：\Quote{不可奸淫}、\Quote{不可杀人}，而不是像其他人所理解的那样。于是，他在学术默观中训练自己，进而操练更广泛的概括和更宏大的命题；深知照先知所说的，\Quote{教导人知识的那位}是主，主藉人的口行事。同样，祂也取了肉身。\par

因此，他理所当然从不将愉悦置于有用之上；即使一个美丽的妇人偶然遇到他，放荡地引诱他，也不能动摇他；因为若瑟主人的妻子未能诱使他偏离自己的坚定；当她强行抓住他的外衣时，他脱下了它——赤裸于罪恶，却穿上了品德的得体。因为倘若主人——我是说埃及人——的眼睛没有看见若瑟，全能者的眼睛却看见了。因为我们听见声音，看见形体；但天主审视事物本身，言语和观看皆由此而出。\par

因此，即使疾病、意外、以及最可怕的事——死亡——临到真知者，他的灵魂仍然坚定不移——深知这一切都是受造界的必然，并且藉由天主的大能，它们成为救恩的药方，藉由管教使难以改过者获益；按功过分配，由真正良善的眷顾所安排。\par

因此，当圣言命令时，并在其命令的范围内，以感恩造物主之情使用受造物，他成了享受它们的主人。\par

他从不怀恨或对任何人积怨，尽管对方的行为可憎。因为他敬拜造物主，并爱那分享生命的人，因他的无知而怜悯并为他祈祷。他确实分受身体的感情，身体本质上是易受苦的，他被缚于身体。但在感觉中，他并非其主要主体。\par

因此，在不自主的环境中，他藉从困扰中退至真正属于他的事物，不被外在于他的事带走。他只适应必要之事，以使灵魂完好无损。因为他希望忠实，不是凭借假设或外表，而是凭知识和真理，即确实的行为和有效的言语。因此，他不仅赞美高尚，而且努力使自己高尚；因爱从良善忠信的仆人变为朋友，藉由习惯的成全，这习惯是他从真确教导和伟大纪律中所获得的纯洁。\par

因此，他努力达到知识的顶峰（\OriginalTerm{知识}{gnosis}）；品格庄重；举止沉着；拥有属于真正真知者的所有优势；定睛于美好的典范——许多生活正直的族长、数不胜数的先知和天使，以及超过一切的、教导并显示他能达到那至上生命的主——因此他不爱世间在他掌握中的一切美好，以免滞留于地，而爱那所希望的，或更确切说，已知道的，希望以便把握。\par

于是他忍受劳苦、考验和磨难，不像那些有男子气概的哲学家，希望眼前的苦难停止，并重新享受快乐；而是知识以最坚定的信念激励他接受将来的希望。因此，他不仅藐视今世的痛苦，也藐视其所有快乐。\par

因此，据说有福的伯多禄看见妻子被带去受死，因她的蒙召和被接回家而欢喜，并非常鼓励和安慰地呼唤她的名字说：\Quote{请记住主。}这就是有福之人的婚姻，以及他们对待至亲之人的完美态度。\par

同样，宗徒也说：\Quote{那娶妻的，要像没有娶妻的}，并认为他的婚姻摆脱了不当的私爱，且与对主的爱不可分离；真正的丈夫曾劝诫他的妻子在离开此生前往主那里时要持守此爱。\par

那些即使在受罚最痛苦时仍感谢天主的人，他们的信德不是显明在对死后希望的确信上吗？因为我认为，他们所拥有的信德是坚固的，且伴随有信德的行为。\par

因此，在所有环境中，真知者的灵魂都是坚强的，处于极其健康和有力的状态，如同运动员的身体。\par

因为他在人事上是谨慎的，能判断义人所当做的事；他从高天由天主那里获得了原理，并为了肖似天主，在身体的痛苦与快乐中获得了节制。他勇敢地对抗恐惧，信赖天主。所以，的确，\OriginalTerm{诺斯替者}{Gnostic}的灵魂，被完美的德性所装饰，是地上天主能力的肖像；其发展是本性、训练和理性共同作用的结果。灵魂的这种美丽成为圣神的宫殿，当它获得与福音相称的整个生活的性情时。这样的人因此能抵抗一切可怕之事的恐惧，不仅是死亡，还有贫穷、疾病、耻辱以及类似的事；他不被快乐征服，是理性欲望的主宰。因为他清楚知道什么是该做的，什么是不该做的；完全知道哪些事真正该怕，哪些不该。因此他明智地承受\OriginalTerm{圣言}{Word}向他提示为必要和必须的事；明智地分辨什么真正是安全的（即善），什么只是看起来安全；以及什么该怕，什么只是看起来该怕，如死亡、疾病和贫穷；这些事更多是在意见中如此，而非在真理中。\par

这就是真正的好人，他没有情欲；借着具有德性的灵魂的习性或性情，他超越了整个情欲的生活。他依靠自己的一切来达到目的。因为那些被称为可怕的意外事件，对好人并不可怕，因为它们不是恶。而真正该怕的事，与\OriginalTerm{诺斯替}{Gnostic}基督徒是无关的，因为它们与善截然对立，因为是恶；而相反的东西不可能在同一时间共存于同一个人身上。那么，他无瑕地演完了天主交给他的人生戏剧，就知道该做什么，该忍受什么。\par

那么，怯懦岂不是来自对什么可怕什么不可怕的无知吗？因此，唯一有勇气的人是\OriginalTerm{诺斯替者}{Gnostic}，他既知道现在和将来的善事；也知道（如我所说）实际上并不可怕的事。因为他知道唯独恶是可憎的，是破坏有助于知识的东西，他穿上主的铠甲，与恶作战。\par

因为如果任何事是因愚昧，以及魔鬼的运作或更确切地说是共谋而起的，这并不就直接是魔鬼或愚昧。因为没有任何行动是智慧。因为智慧是一种习性。而没有任何行动是一种习性。那么，从无知而来的行动，并不就是无知，而是出于无知的恶，但并非无知。因为心灵的扰动或罪都不是恶习，尽管它们来自恶习。\par

所以，没有一个非理性的勇敢者是\OriginalTerm{诺斯替者}{Gnostic}；因为人们可能称儿童为勇敢，他们由于不知什么可怕，就经历可怕的事。所以他们甚至触摸火。而冲到长矛尖端的野兽，凭着野蛮的勇气，也可以称为勇敢。这样的人或许还会称变戏法的人勇敢，他们以某种技巧在剑上翻筋斗，为了可怜的报酬操练有害的技艺。但真正勇敢的人，在眼前看到群众恶意带来的危险时，勇敢地等待一切来临。这样，他就与其他被称为殉道者的人区别开来，因为有些人自己制造机会，不知怎地就冲入危险的中心（因为用温和的语言说是对的）；而他们则按照正确的理性保护自己；然后，当天主真正召唤时，立刻交出自己，并确认召唤，因为他们意识到没有轻率，并在真正的理性刚毅的操练中呈现那人被考验。因此，他们既不像其他人那样因害怕更大的危险而忍受较小的危险，也不因害怕同伴和同感者的责备而继续持守他们召叫的宣认；而是出于对天主的爱，他们自愿服从召唤，除了悦乐天主之外别无目的，并非为了辛苦的报酬。\par

因为有些人出于虚荣心受苦，另一些人出于害怕另一种更严厉的惩罚，还有一些人为了死后的享乐和愉悦，他们是信仰上的孩童；确实有福，但尚未成为在天主爱中的成人，如\OriginalTerm{诺斯替者}{Gnostic}那样。因为正如在体育比赛中，在教会中也有为成人和为孩童的冠冕。但爱本身应被选择，不为任何其他东西。因此，在\OriginalTerm{诺斯替者}{Gnostic}中，知识与刚毅的完美从生活的操练中发展出来，他总在学习掌控情欲。\par

因此，爱使它的运动员无所畏惧、毫不胆怯，并在主内自信，傅油并训练他；正如正义为他的整个生活确保诚实。因为说\Quote{你的是就是是，你的不是就是不是}是正义的摘要。\par

对节制也是如此。因为它既不是为了荣誉的爱，如运动员为了冠冕和名声；也不是为了爱钱，如有些人假装操练节制，以可怕的痛苦追求善。也不是为了健康而爱身体。更不是任何因粗野、未尝过快乐而节制的人，是真正有节制的人。当然，那些过着劳苦生活的人，一尝到快乐，就立刻把节制的坚定性瓦解成快乐。这样的人受法律和恐惧的约束。因为一找到有利时机，他们就欺骗法律，放走善。但为自身而可欲的节制，借着知识而达致完美，常存不变，使人成为自己的主人和主宰；这样，\OriginalTerm{诺斯替者}{Gnostic}就是有节制的、没有情欲的，不能被快乐和痛苦所瓦解，就像他们说金刚石不能被火瓦解一样。\par

那么，这些的原因就是爱，它是一切知识中最神圣、最至高的。\par

因为借着对那至善至高的、以合一为特征的事物的服侍，它使\OriginalTerm{诺斯替者}{Gnostic}同时成为朋友和儿子，真正成长为\Quote{完美的人，达到基督圆满年龄的程度} \BibleRef{弗 4:13}。\par

再者，在同样事物上的一致就是同意。但同者即是一。友谊在相似中达至完美；团体在于一体。因此，真知者由于是唯一真神的爱慕者，是真正完美的人，是天主的朋友，并被置于儿子的等级中。因为这些是高贵、知识以及默观天主的完美之名；当真知者的灵魂变得十分纯洁，被认为配得永远看见全能的天主时，它便领受了这进步的顶峰，据说，\Quote{脸,} 据说，\Quote{面对面。} 因为成为了完全属灵的，并且在属灵的教会中已回归到同类的本性，它安息于天主的安息中。\par

% structure: p0191 source=h2 target=subsection reason=relative-heading-rank; base=h2

\subsection{第十二章 真知者是仁爱的、节制的、轻视世俗之事的}

就让这些事情如此吧。真知者对于身体和灵魂——对于他的邻人，无论是家人，或是合法的仇敌，或是任何人——都持这样态度，他显得平等而相似。因为他并没有\Quote{轻视他的弟兄，} 按天主的法律，这位弟兄同出于一父一母。他确实安慰受苦者，以安慰、鼓励和生活的必需品帮助他；给与一切需要的人，虽然并非同等，而是公义地，按应得的；此外，给与那迫害和仇恨他的人，即使他需要；对那些对他说他是出于恐惧而施与的人，他不太在意，如果他这样做并非出于恐惧，而是为了帮助。因为那些对自己的敌人缺乏贪财之心、嫉恶如仇的人，岂不是更以爱来对待那些属于他们自己的人吗？\par

这样的人由此前进，达到精确的知识，知道主要应当给谁、给多少、何时给、如何给。\par

谁若有任何理由，能成为任何方面都不引起敌意的人之敌人呢？这岂不正如在天主的情形中一样？我们说天主不是任何人的对手，也不是任何人的敌人（因为祂是万有的创造者，没有任何存在的事物不是祂所愿意的；但我们断言，那不顺服、不按祂诫命行走的人，是祂的敌人，因为他们敌视祂的盟约）。我们将在真知者身上发现完全相同的情形，因为他绝不可能以任何方式成为任何人的敌人；但那转向相反道路的人则可视为他的敌人。\par

尤其在我们中间盛行的宽宏习惯被称为\Quote{义德；} 但按应得的分辨大小，针对那些适合的对象的能力，则是最崇高义德的一种形式。\par

有些人以一种粗俗的方式实践某些事情，例如对快感的节制。因为，如在外邦人中，有些人是由于无法获得所见之物、出于对人的恐惧，以及为了更大的快感，而戒绝眼前的享乐；同样，在信仰中，有些人实行自制，或是出于对承诺的重视，或是出于对天主的畏惧。好吧，这样的自制是知识的基础，是朝向更好之事的途径，是朝向成全的努力。因为据说：\Quote{敬畏上主是智慧的开始。} \BibleRef{箴 1:7} 但成全之人出于爱，\Quote{凡事包容，凡事忍耐，} \BibleRef{格前 13:7} \Quote{不讨人的喜悦，只讨天主的喜悦。} \BibleRef{得前 2:4} 虽然赞美随之而来，但并非为他自己的益处，而是为了那些赞美他的人的效法和益处。\par

按另一种看法，那仅仅控制自己情欲的人并不被称为节制的人，而是那也获得了对美好事物的掌控，并确实获得了学问的伟大成就，由此结出德行活动果实的人。因此，真知者在紧急情况发生时，绝不会脱离他所特有的习惯。因为对善的科学拥有是稳固而不可变的，是对神与人之事的知识。知识绝不会变成无知，善也绝不会变成恶。因此，他也吃喝、结婚，并非作为生存的首要目的，而是作为必要的事。我甚至提到结婚，如果圣言规定，并且适宜。因为成为成全之人后，他以宗徒们为榜样；一个人并非真正在独身生活的选择中显为男子汉；而是他超越众人，他因结婚、生育子女和照顾家庭所受的磨练，在不为喜乐或痛苦所动的情况下，在他的家庭关怀中，与天主的爱不可分离，并抵挡了来自儿女、妻子、仆人和财产的一切试探。但那没有家庭的人在很大程度上免于试探。因此，只照顾自己的人，就他个人的救恩而言，虽不如那较差的人，但在生活行为上却更优越，在他对真理的关怀中，确实保存了一个微小的肖像。\par

我们必须尽可能地使灵魂接受各种预备训练，使之能领受知识。你们岂不见蜡如何被软化，铜如何被净化，以接受印在上面的印记吗？正如死亡是灵魂与身体的分离，知识仿佛是一种理性的死亡，催促灵魂脱离情欲，引向善行的生活，以便它随后能自信地对天主说：\Quote{我照祢的旨意生活。} 因为那以讨人喜悦为目的的人不能讨天主的喜悦，因为群众选择的不是那有益的，而是那令人愉快的。但在讨天主喜悦时，一个人作为结果也得到善人中的好感。那么，与肉食、饮料和色欲相关的事，怎能令这样的人喜悦呢？因为他甚至对产生快感的话语、愉快的心灵运动和行为都持怀疑态度。\Quote{因为没有人能事奉两个主人：天主和玛门，} \BibleRef{玛 6:24} 它所说；意思不仅是金钱，而且是从金钱用于各种享乐所产生的资源。实际上，那慷慨而真正认识天主的人，不可能事奉敌对的情欲。\par

那么，惟独有一位从一开始就无\OriginalTerm{私欲偏情}{concupiscence}——就是那为我们成为人的\OriginalTerm{仁爱之主}{philanthropic Lord}。凡努力肖似祂所赐印痕的人，都藉操练力图成为无私欲偏情的。因为那曾怀有私欲偏情而后又自制的人，就如一个寡妇藉节制而复为贞女。这便是知识的赏报——即远离邪恶、行善立功——由救主和导师所赐，也是祂自己所求的，藉此获得救恩。\par

正如那些学会技艺的人凭所学谋生，同样，\OriginalTerm{真知者}{Gnostic}也藉他所知获得生命而得救。因为那不愿根除\OriginalTerm{灵魂的情欲}{passion of the soul}的人，便是自杀。然而，无知似乎是灵魂的饥荒，而知识是灵魂的滋养。\par

这些就是\OriginalTerm{真知的灵魂}{gnostic souls}，福音将它们比作等候主的\OriginalTerm{献身贞女}{consecrated virgins}。因为她们是贞女，在于远离邪恶；又因她们出于对主的爱而等候，并点燃她们的光以默观万物，故她们是明智的灵魂，说：\Quote{主，长久以来我们渴望领受祢；我们按祢所吩咐的而生活，没有违犯祢的任何诫命。因此我们也要求应许。我们祈求有益的事，因为向祢求那至为卓越的，原是不必要的。我们将一切视为美好，即使我们所遇的磨炼——那是祢的安排为训练我们的坚忍而带来的——看似邪恶，也是如此。}\par

所以，\OriginalTerm{真知者}{Gnostic}因其超乎寻常的圣德，宁愿在祈求时不得应允，也不愿在未祈求时得到。\par

他的整个生命就是祈祷并与天主交谈。他若清白无罪，必能得其所愿。因为天主对义人说：\Quote{你求，我必给你；你想，我必成就。}若是有益，他必立刻领受；若是有害，他永不会求，因此也不会得到。事必如其所愿。\par

但若有人对我们说，有些罪人甚至按其请求获得了恩赐，我们应说：这事极少发生，是出于天主公义的仁慈。这恩赐是赐给那些能有益于他人的人。因此，这恩赐并非为祈求者的缘故，而是天主的上智安排预见有人会因他而得救，从而使这恩赐再次成为公义的。对于那些堪当的人，即使他们不祈求，真正美好的事物也会赐予他们。\par

因此，当一个人是正义的，不是出于必要或恐惧或希望，而是出于自由选择时，这便称为\OriginalTerm{王道}{royal road}，是\OriginalTerm{王族}{royal race}所行走的。但旁道是滑溜而险峻的。倘若除去恐惧和尊荣，我不知道那些在哲学家中盛名远扬、如此畅所欲言的人，是否还能忍受苦难。\par

现在，情欲和其他罪恶被称为\Quote{荆棘和蒺藜}。因此，真知者在主的葡萄园中劳作，栽种、修剪、浇灌，是栽种于信德中的\OriginalTerm{神圣农夫}{divine husbandman}。那些没有作恶的人，认为理应得到安逸的报酬。但那出于自由选择而行善的人，要求一份应得的酬劳，如同一个好工人。他必定得到双倍的工资——既为他未作恶，也为他所行的善。\par

这样的\OriginalTerm{真知者}{Gnostic}，除非获得天主的许可并且为了与他同在之人的益处，否则无人能诱惑他；他藉勇毅的忍耐鼓励他们，坚固他们的信德。诚然，正是为此目的——为建立和坚固各教会——蒙福的宗徒们才被置于试炼和殉道之中。\par

因此，\OriginalTerm{真知者}{Gnostic}听到耳畔响起的声音说：\Quote{我所击打的，你要怜悯，}就恳求那些恨他的人能悔改。因为对作恶者的惩罚——在通衢大道上执行——是给孩童观看的；真知者既已出于选择修习成卓越和良善，便不可能从这类景象中得到教训或喜悦。这样，他既已变得不为享乐所软化，也从不陷入罪恶，便不因他人受苦的榜样而受纠正。那已高贵地蔑视这世界所呈现的前景——纵然这些前景是神圣的——的人，更是远不以地上的享乐和景象为乐。\par

因此，\Quote{不是凡向我说\InnerQuote{主啊！主啊！}的人，就能进入天国；而是那承行天主旨意的人。}\BibleRef{玛 7:21}这就是\OriginalTerm{真知工人}{gnostic labourer}，他即使仍在肉身内，也拥有对\OriginalTerm{世俗欲望的掌控}{mastery of worldly desires}；对于未来而仍不可见的事物，他因认识而有确切的信念，以致视它们比眼前的事物更为真实。这位能干的工人为他所认识的事而喜乐，但因陷于生活的必须而受束缚，尚未堪当亲身参与他所认识的事。他使用此生如同使用属于别人的——仅到必要的程度。\par

他也知道那些日子禁食的\OriginalTerm{奥秘}{enigmas}——我指的是\OriginalTerm{第四日和预备日}{Fourth and the Preparation}。因为一名源自\Person{赫耳墨斯}，另一名源自\Person{阿佛洛狄忒}。他在生活中禁食，远离贪婪和纵欲，一切恶习皆由此而生。因为我们已多次根据宗徒指出三种私通：爱享乐、爱金钱、崇拜偶像。因此，他按法律禁食，远离恶行；按福音的圆满，远离恶念。试探加于他，不是为洁净他，而是如我们所说，为邻人的益处，若他经受劳苦和痛苦而轻视并越过它们。\par

同样的道理也适用于快乐。因为一个经历过快乐的人，之后能够克制，这是最高的成就。如果一个人在他不认识的事物上自制，这有什么了不起呢？他依照福音的诫命守主日，就是当他抛弃邪恶的性情，而接纳诺斯替者的性情时，在自己身上光荣主的复活。此外，当他领受了科学思辨的理解时，他认为自己看见了主，将目光转向不可见的事物，尽管他似乎在看他不愿看的事物；他惩戒视觉官能，当他觉察到自己因眼睛注视而感到愉悦时；因为他只愿看见和听见与他相关的事物。\par

在默观弟兄们的灵魂时，他也同时看到了肉体的美，连同灵魂本身，这灵魂已习惯于只注视那善的，而没有肉体的快感。他们确实是弟兄；因为他们蒙选的受造、品格的合一以及行为的本质，使他们按照主希望蒙选者被激励的情感，施行、思考和谈论相同的神圣和善工。\par

因为信德表现在他们选择相同的事物上；知识表现在学习和思考相同的事物上；望德表现在渴望相同的事物上。\par

倘若因生活的必需，他在维持生计上花费少许时间，他就觉得自己受到欺骗，因事务而分心。因此，即使在梦中，他也不看任何不适合蒙选者的事物。因为每一个这样的人，住在城中却轻视城中他人所羡慕的事物，住在城中如同住在旷野，以致地方不能强迫他，而是他的生活方式显明他是义人。\par

简言之，这位诺斯替者以其生活的正直、知识的准确、造福亲属、\Quote{移开邻居的山}并消除他们灵魂的混乱，来弥补宗徒的缺席。尽管我们每个人都是自己的葡萄园和工人。\par

他在做最卓越之事的同时，也希望避开人的注意，在主面前劝服自己按照诫命生活，因相信这些事存在而选择它们。\Quote{因为一个人的心在哪里，他的宝藏也在哪里。}\par

他使自己贫穷，为的是因爱德中的成全，永不忽略一个陷入痛苦的弟兄，尤其若他知道自己比那弟兄更容易忍受匮乏。因此，他把别人的痛苦视为自己的忧伤；若因自己的贫困而施善行遭受任何困苦，他不因此而烦恼，反而更加增进他的施舍。因为他拥有在生活事务中操练的信德，并在实践和默观中赞美福音。事实上，他因实行主的教训而得称赞，\Quote{不是从人，而是从天主}，\BibleRef{罗 2:29}。\par

他被自己的望德所吸引，不尝世上的美物，对这里的一切持有高贵的蔑视；怜悯那些死后受罚的人，他们因惩罚而不情愿地认罪；对自己的离世有清白的良心，并常常预备如同\Quote{外方人和旅客}，对于这里的产业；只记得那些属于他自己的，视这里的一切为不属于他的；不仅仰慕主的诫命，而且可以说，因知识本身而参与天主的旨意；因正义而成为主和祂诫命真正蒙选的密友；作为诺斯替者，拥有王子和君王的尊严；藐视地上和地下的所有黄金，以及从海岸到海洋的统治权，以便他能专注于唯独事奉主。因此，在饮食、婚姻（若圣言命令如此），甚至在梦中，他施行和思想圣洁的事。\par

因此，他时常保持纯洁以祈祷。他也与天使的团体一同祈祷，因为他已具有天使的品级；他永不离天主的圣洁守护；即使他独自祈祷，也有圣人的歌侣与他同站。\par

他认识到信德中的双重因素：信者活动，以及所信之物按其价值的卓越；因为义德也有双重：出于爱德的义德和出于畏惧的义德。因此经上说：\Quote{对主的敬畏是纯洁的，永远常存。}因为那些因畏惧而转向信德和义德的人，将永远长存。如今，畏惧能止恶，而爱德则勉励行善，增进到自发行善的程度；这样人可以从主那里听到：\Quote{我不再称你们为仆人，而是朋友。}而今他可以满怀信心地投身于祈祷。\par

他祈祷的形式是为过去、现在和未来（如同已通过信德而临在）感恩。这之前有知识的领受。他祈求以诺斯替者的身份度过在肉身中指定的生命，如同脱离了肉身，以获得最好的事物，逃避更坏的事物。他也祈求在那些我们犯罪的事物上获得宽恕，并悔改承认它们。\BibleRef{路 18:18}\par

当他离世时，他追随那位召唤他的，几乎如行在前面的召唤者一样迅速，因良善的良心而急切感恩；到达那里后，与基督在一起，因他的纯洁而证明自己堪当通过一种融合拥有基督所传达的天主的能力。因为他不想藉分有热力而温暖，或藉分有火焰而发光，而愿完全是光。\par

他准确地知道这句宣告：\Quote{除非你们恨父亲和母亲，甚至自己的性命，除非你们背起[十字架的]记号。}\BibleRef{路 14:26} 因为他恨那肉体的无节制情感，它们拥有快乐的有力魔力；并高尚地鄙视一切属于肉体受造和营养的事物。他也抗拒那属血肉的灵魂，给倔强的非理性精神套上缰绳：\Quote{因为肉体的欲望相反圣神。}\BibleRef{迦 5:17} 而\Quote{背起[十字架的]记号}就是背着死亡，在活着时与一切告别；因为\Quote{撒播了肉体}与为知识塑造了灵魂之间没有同等的爱。\par

他既已养成了行善的习惯，便行施舍，比说话更快；他祈祷能分担弟兄们的罪过，好使他的亲属得以告解并悔改；并渴望将自己所有的善分给最亲近的人。因此，这些人对他是朋友。于是，按主所命令的农事，促使在他内播下的种子生长，他继续无罪，成为节制的，并以精神与那些如同他的圣者歌团中的世人同住，虽仍被拘留于地上。\par

他日夜言行主的诫命，极其喜乐，不仅在早晨和中午起身时，而且在行走、睡觉、穿衣和脱衣时；他教导他的儿子，如果有儿子的话。他不可与诫命和望德分离，不断感谢天主，如依撒意亚以比喻所说的活物，在一切试探中顺服，他说：\Quote{主赐予的，主也拿去了。}\BibleRef{约 1:21} 因为约伯也是如此；他在财产被掠夺之后，连同身体的健康，因爱主而舍弃一切。因为经上说他是\Quote{正义、圣洁，远离一切邪恶}。\BibleRef{约 1:1} 而\Quote{圣洁}一词指出了对天主的一切义务和整个生活过程。他知道这些，因此他是诺斯替者。因为我们既不可过分依恋这类事物，即使它们是好的人性之物，也不可憎恶它们，如果它们是坏的；而必须超脱两者[好与坏]，践踏后者，将前者转给需要的人。但诺斯替者在适应上谨慎，以免被察觉，或适应变成习性。\par

% structure: p0226 source=h2 target=subsection reason=relative-heading-rank; base=h2

\subsection{第十三章  继续描述诺斯替者}

他从不记念那些得罪过他的人，而是宽恕他们。因此他正义地祈祷说：\Quote{宽恕我们，因为我们也宽恕。}\BibleRef{玛 6:12} 因为这也是天主愿意的事情之一：不贪图什么，不恨任何人。因为众人都是一个意志的作品。岂不是救主愿意诺斯替者如同\Quote{天上的父}\BibleRef{玛 5:48} 一样完善，即祂自己，祂说：\Quote{来，孩子们，听我教导敬畏上主？}祂愿意他不再需要借助天使的帮助，而是从祂自己那里得到，因他已成为堪当的，并通过服从从祂自己那里获得保护。\par

这样的人向主索求，而不只是祈求。当他弟兄缺乏时，诺斯替者不会为自己求满满财富来施舍，而会祈祷他们的需要得到供应。因为诺斯替者如此将他的祈祷赐给需要的人，因他的祈祷他们得到供应，而他不知，也无虚荣。\par

贫乏、疾病以及类似的考验，常为训诫而发送，为纠正过去，并照顾将来。这样的人祈祷从中解脱，因拥有知识的特权，并非为虚荣；而正因他是诺斯替者，他施行施舍，成为天主良善的工具。\par

据传统说，玛窦宗徒常言：\Quote{如果一位蒙选者的邻人犯罪，蒙选者便犯了罪。因为若他按圣言所规定而行，他的邻人也会因对他所活出的生活充满敬畏而不犯罪。}\par

那么，对诺斯替者本身，我们该说什么？\Quote{你们不知道}，宗徒说，\Quote{你们是天主的殿宇吗？}\BibleRef{格前 3:16} 诺斯替者因而神圣，已圣洁，承载天主，并为天主所承载。而圣经显明犯罪与他无关，将跌倒者卖给陌生人，说：\Quote{不要观看陌生女人，以致贪恋}，明确宣告犯罪是外来的，相反天主的殿宇的本性。殿宇大，如教会；也小，如保存亚巴郎后裔的人。因此，让天主安息在他内的人不会再渴望其他事物。他立即离开一切阻碍，轻视一切分心的事物，藉知识穿透天空。穿越精神的本体和一切宰制者与掌权者，他触及最高的宝座，只为那唯一值得认识的事物而疾行。\par

于是，将\Quote{蛇与鸽子}\BibleRef{玛 10:16} 调和，他既完善又良心清白地生活，混合信德与望德，以期待未来。因为他意识到他所获得的恩赐，已变得堪当获得它；并从奴隶地位转变为义子地位，作为知识的结果；认识天主，或更确切地说，被祂所认识，因为结局，他发出与恩宠价值相称的能力。因为工作跟随知识，如同影子跟随身体。\par

因此，他正确地对任何发生的事不感困扰；也不怀疑那些借着上智安排而发生的、为善的事。他也不耻于死亡，具有好良心，并足以被掌权者所见。他，可说是净化了灵魂的一切污点，清清楚楚地知道，他离世后会更好。\par

因此，他从不将享乐和利益置于天主的安排之上，因为他以诫命训练自己，使自己在一切事上中悦于主，在世眼中值得称颂，因为万有都依赖唯一至高的天主。据说，天主子来到自己的地方，自己的人却没有接受他。因此，即使在使用世间之物时，他不仅感谢赞美受造界，而且当正当地使用它们时，他本身也受赞美；因为他的视域终点是以符合诫命的\OriginalTerm{诺斯替的}{gnostic}活动进行的默观。\par

于是，藉着知识收集材料作为默观的食粮，高尚地拥抱了知识的广大，他前进到由此世被提去的圣洁报偿。因为他听见了那圣咏所说的话：\Quote{环抱熙雍，环绕它，讲述它的城楼。}因为我认为，这暗示那些崇高拥抱了圣言的人，以致成为高耸的城楼，并坚定地站在信德和知识中。\par

让这些关于诺斯替者的陈述，以尽可能简洁的方式包含主题的萌芽，向希腊人提出。但要知道，如果那位\Emph{单纯的}信徒正确地做了这些事中的一件或两件，他却不会在所有事上以最高的知识这样做，如同\OriginalTerm{诺斯替者}{Gnostic}一样。\par

% structure: p0237 source=h2 target=subsection reason=relative-heading-rank; base=h2

\subsection{第十四章 基于格林多前书六章1节等对诺斯替者的描述}

现在，论及我们所归诸于诺斯替者的所谓不动情（在此，信徒的完美，\BibleQuote{藉着爱德前进，达到完美的人，达到圆满身量的程度} \BibleRef{弗 4:13}，通过肖似天主并真正成为天使般的），我想起许多其他圣经的见证可以引用。但我认为，由于论述的长度，最好将这样的荣誉委托给那些愿意付出辛劳的人，并留给他们通过选择圣经来阐述教义。\par

因此，我将以最简洁的言语提及一段经文，以免这个话题得不到解释。\par

因为在致格林多人的第一封书信中，神圣的宗徒说：\BibleQuote{你们中间有人与别人有了争执，怎么竟敢在不义的人面前诉讼，而不在圣徒面前呢？你们不知道圣徒将要审判世界吗？} \BibleRef{格前 6:1}，等等。\par

由于此段很长，我们将通过使用最贴切的宗徒言论，以最简洁的语言和一种粗略的方式，解释他描述诺斯替者完美的论述，来展示宗徒话语的含义。因为他不仅举例说明诺斯替者的特征是宁愿受害而不愿害人；而且他教导说，他不怀恨在心，甚至不允许他为伤害他的人祈祷。因为他知道主明确命令\BibleQuote{要为仇敌祈祷} \BibleRef{玛 5:44}。\par

因此，说被伤害的人在不行义的人面前诉讼，无异于说他表露出报复的意愿，以及反过来伤害对方的欲望，这同样也是他自己行不义。\par

而他所说的他愿意\Quote{有些人在圣徒面前诉讼}，是指那些通过祈祷要求害人者因其不义遭受报复的人，并暗示后者比前者好；但他们仍未顺服，如果不完全摆脱怨恨，甚至为仇敌祈祷。\par

那么，他们从悔改中获得正确的倾向是好的，悔改导致信德。因为如果真理似乎有怀有恶意的仇敌，但它并不敌视任何人。\BibleQuote{因为天主使太阳升起，照耀义人，也照耀不义的人} \BibleRef{玛 5:45}，并差遣主自己来到义人和不义的人中间。那竭力肖似天主的人，在极大的无怨恨中，宽恕七十个七次，仿佛一生之久，并在此世全部行程中（由七的数目所指示）向每一个人显示仁慈；若有人在他在肉身内的整个生命期间加害于诺斯替者。因为他不仅认为善人应仅仅放弃自己的财产，让给那些曾加害于他的人；而且希望义人应请求那些法官原谅那些加害他的人的过犯。这是有道理的，如果那些试图加害他的人只在外部和关乎身体的事上（即使到了死亡的程度）占他的便宜；这些没有一样真正属于诺斯替者。\par

再者，一个人若自身背弃了那按照福音的不念旧恶，又怎能\Quote{审判}背信的\Quote{天使}呢？他说：\BibleQuote{你们为什么不情愿受冤？为什么不情愿被亏负？你们反倒冤屈人，亏负人} \BibleRef{格前 6:7}，明显是通过祈祷反对那些因无知而犯罪的人，并尽你所能，剥夺你所祈祷反对的人的天主的慈爱和良善，\Quote{而这些人正是你们的弟兄}——不仅指在信德内的，也指归依者。因为那现在敌对的人以后是否会相信，我们尚不知道。由此清楚得出结论：即使所有人尚未成为我们的弟兄，也应当如此看待。如今，唯有诺斯替者认识到所有人都是同一天主的工程，并且在同一个本性中赋予了同一肖像，虽然有些人可能比其他人更混浊；并且在受造物中，他认出天主的作为，借此他又朝拜天主的旨意。\par

\BibleQuote{难道你们不知道不义的人不会承受天主的国吗？} \BibleRef{格前 6:9} 报复的人行不义，无论是通过行为、言语，还是意愿的意念，这些在法律的训练之后，福音拒绝了。\par

\Quote{你们中间也有人从前是这样}——显然就是那些你们不肯原谅的人；\BibleQuote{但你们已经洗净了，} \BibleRef{格前 6:11} 不像其余的人那样简单，而是带着知识；你们已经抛弃了灵魂的情欲，为要尽可能相似于天主眷顾的良善，藉着忍耐和宽恕\Quote{向义人及不义之人，}将温和的光辉倾注在他们身上，如同太阳一样。\par

诺斯替信徒将藉由心灵的伟大，或藉由对更善之事的效法来达成这一点。这是第三个原因。\Quote{你们宽恕，也必蒙宽恕；}这条诫命，可以说是藉着丰富的良善强迫人走向救恩。\par

\Quote{但你们已经成圣了。}因为达到这种状态的人，已处于圣洁的境地，绝不再陷入任何情欲，仿佛已经脱去了肉体，已在这大地之外成圣了。\par

\Quote{因此，}他说，\Quote{你们是在主的名里成义的。}你们可以说是由祂造成，像祂一样成为义人，并尽可能与圣神融合。因为\BibleQuote{凡事我都可行，但我决不受任何事物的辖制，}\BibleRef{格前 6:12}以致做出、思想或说出任何与福音相悖的事。\BibleQuote{食物是为肚腹，肚腹是为食物；但天主都要废掉这两样。}\BibleRef{格前 6:13}——这是指那些想和活得像为吃而生的人，而不是为了活着而吃，并将知识作为首要目的。他岂不是说这些如同圣教会的肉身部分吗？作为身体，主的教会——属灵而神圣的团体——被象征。因此，那些仅被称为（基督徒）却不按圣言生活的人，就是肉身部分。\Quote{现在}这属灵的\Quote{身体}，即圣教会，\Quote{不是为淫乱的。}那些属于外邦生活的事物，不可藉离弃福音而加以采纳。因为谁在教会内以异教方式行事，无论行为、言语或甚至思想，他就是对教会和自己的身体行淫。谁这样\Quote{与娼妓联合}，即与违背盟约的行为联合，就成了另一个\Quote{身体}，非圣洁的，\Quote{成为一体}，并且有了异教徒的生活和另一种希望。\Quote{但与主在灵里联合的人}便藉着另一种结合成为属灵的身体。\par

这样的人完全是儿子，是圣洁的人，无激情的，诺斯替的，成全的，由主的教导所塑造；为使他行为、言语甚至精神本身都与主亲近，从而获得那因达到如此成年之境而应得的居所。\par

对那有耳的人，这例子已足够。因为无需揭示奥秘，只需为那些有分于知识的人指出足以唤起记忆的要素；他们也将领悟主如何说：\Quote{你们应当成全，如同你们的圣父完全成全一样，} 藉着宽恕罪过、不记恨、并养成无激情的习惯。因为正如我们称一位医生为成全的，一位哲学家为成全的，同样，在我看来，我们也称一位诺斯替信徒为成全的。但这些要点中，没有一项——即使最为重要——是为了相似天主而假设的。因为我们不说，如斯多葛派最不敬地主张的，在人与在天主身上的美德是一样的。我们难道不应该成全，如圣父所愿的吗？因为完全不可能有人像天主一样成全。圣父愿意我们藉着按福音的服从、无可指摘地生活而成为成全的。\par

那么，若语句有省略，我们理解了所缺的部分，以补全那段话，使那些无法理解省略内容的人也能明白，我们就既懂得天主的旨意，也将同时虔诚而大度地行走，合乎诫命的尊荣。\par

% structure: p0254 source=h2 target=subsection reason=relative-heading-rank; base=h2

\subsection{第十五章 解答因异端纷争而犹豫加入教会的反对意见}

既然接下来要答复希腊人和犹太人对我们的反对，并且先前讨论的一些问题中，那些持其他教导的派别也提供了帮助，那么最好先清除面前的障碍，然后为解答难题做好准备，推进到下一部杂论集。\par

首先，他们这样向我们提出反对，说他们不该因派别的分歧而相信。因为当一些人教导一套信条，另一些人教导另一套时，真理就被扭曲了。\par

对此我们说：在你们犹太人当中，以及在希腊最著名的哲学家中，也产生了许多派别。然而你们并不说，因为你们当中的派别彼此不一致，就应当犹豫是否要从事哲学或信奉犹太教。再者，异端将撒在真理中，如同\Quote{莠子撒在麦子中}，这已由主预言过；所预言必然发生的事不能不发生。其原因在于，一切美丽的事物之后总会跟着丑陋的污点。那么，如果有人违背承诺，偏离他在我们面前所作的宣认，我们是否要因为他背弃了表白而放弃真理？正如善人不应背信，也不应食言，即使他人违背承诺；同样，我们也绝不可逾越教会的规范。尤其我们在最重要的事上持守我们的宣认，而他们却背弃它。\par

因此，应当相信那些坚守真理的人。我们可以广泛使用这一答复，对他们说：按照各自学派持有相反意见的医生，实际上同样治疗病人。那么，一个人身体患病需要治疗，会因医学学派不同而不去求医吗？同样，灵魂患病、充满偶像的人，也不应以异端为借口，耽误获得健康、归向天主。\par

此外，据说正是为了\Quote{那些蒙悦纳的人}，才存在异端。\BibleRef{格前 11:19}（宗徒）称那些在达到信德时，以某种辨别力应用主的教导的人为\Quote{蒙悦纳的人}（如同那些被称为精明的兑换银钱者，他们凭假印记区分伪币与真币），或者那些已经在生活和知识上都成为蒙悦纳的人。\par

因此，我们需要更大的注意和思虑，以便探究我们应当如何生活，以及什么是真正的虔诚。显然，正是因为真理难以获得且充满艰辛，才产生了问题，从这些问题中又生出了异端，带着那些没有真正学习或领悟，而仅捕获了知识的虚妄观念之人的自爱与虚荣。因此，我们当以更大的谨慎来审查那唯一的真正真理，其唯一的目标就是真天主。这辛劳之后，是甜蜜的发现与回忆。\par

因此，为了异端的缘故，我们必须承担发现的辛劳，但绝不能放弃真理。因为，如果在我们面前摆着水果，一些是真正的、成熟的，一些是用蜡做的、尽可能像真的，我们不应因相似而两者都避开。相反，要借着理解的沉思和最具决定性的推理，分辨真伪。\par

正如有一条皇家大道，却有许多其他道路：一些通向悬崖，一些通向湍急的河流或深海，没有人会因为道路的多样而退缩不前行，反而会使用那条安全、皇家且常有人走的路；同样，尽管人们对真理各说各话，我们也不可放弃真理，而必须寻求关于它最准确的知识。既然在菜园中杂草也与蔬菜一同生长，难道农夫就应该停止耕种吗？\par

既然我们天生有充足的手段来审查那些陈述，我们就应当发现真理的脉络。因此，如果我们不赞同应当服从的事，不区分敌对的、不雅的、不自然的、虚假的与真实的、一致的、恰当的、合乎自然的事，我们确实会被正当地定罪。这些手段必须被用来达到对真正真理的认识。\par

因此，对希腊人而言，这种借口是徒劳的，因为愿意的人可以找到真理。但对那些提出不合理借口的人，对他们的定罪是无可辩驳的。因为他们是否否认或承认有这样一件事作为证明？我认为所有人都会承认，除非那些剥夺感官的人。既然有证明，就必须屈尊俯就问题，并借着圣经本身以证明的方式确定异端如何失败，以及只有在真理和古老的教会中，才有最精确的知识和真正最好的\OriginalTerm{原则}{\GreekText{αἵρεσις}}。\par

那些偏离真理的人中，有些试图欺骗自己，有些则也欺骗邻人。那些被称为\OriginalTerm{自以为有智慧的人}{\GreekText{δοξόσοφοι}}，认为自己已经找到了真理，却没有真正的证明，他们在自认为已到达安息之地的想法中欺骗自己。其中不乏有许多人，他们为了避免反驳而回避探究，为了避免定罪而逃避教导。而那些欺骗前来寻求他们的人则非常狡猾，他们明知自己一无所知，却以似是而非的论证来模糊真理。\par

但依我之见，似是而非的论证的本性是一回事，真正论证的本性则是另一回事。我们知道，异端的称呼必须在与真理的对比中表达；诡辩家们从真理中汲取某些东西来毁灭人，并将它们埋葬在自己发明的人为技艺中，他们以成为学派的领袖为荣，而不是管理教会。\par

% structure: p0267 source=h2 target=subsection reason=relative-heading-rank; base=h2

\subsection{第十六章　圣经是区分真理与异端的准则}

但那些准备好投入最卓越追求的人，不会停止对真理的探求，直到他们从圣经本身获得证明。\par

有一些共同的准则适用于所有人，如感官；另一些属于那些已将意志和精力用于真理的人——即心智和理性所使用的方法，用以区分真假命题。\par

其次，放弃意见是一件非常伟大的事，即置身于准确知识与意见的轻率智慧之间，并知道那盼望永享安息的人也知道进入安息的门是艰辛\Quote{而狭窄的}。那一旦接受了福音的人，甚至在他得到救恩知识的同一时刻，就\Quote{不要像罗特的妻子那样}回头，正如经上所说；他不可回到先前的生命（这生命依附于感官之物），也不可回到异端中去。因为异端塑造人的品格，却不认识真天主。\Quote{凡爱父亲或母亲超过爱我的，}即那真理的圣父和教师，祂重生、再造并滋养那被选的灵魂，\Quote{就不配属于我}——祂的意思是，不配作天主的儿子和天主的门徒，同时也不配作朋友和同类的本性。\Quote{因为没有人回头，手扶着犁，适合于天主的国。}\BibleRef{路 9:62}\par

但是，看来甚至直到我们的时代，许多人仍因\Person{玛利亚}生了孩子而认为她处于产褥状态，尽管她并没有。因为有些人说，她在分娩之后，经检验被发现仍是童贞。\par

对我们来说，这就是主的圣经，它生了真理并继续是童贞，隐含着真理的奥秘。\Quote{她生了孩子，却又没有生孩子，}圣经说，如同她由自己受孕，而非因结合。因此，圣经为诺斯替派孕育了孩子；而异端没有学习它们，就弃绝了它们，如同没有怀孕一样。\par

现在，所有的人都具有同样的判断力，有些人跟随那说话的圣言，为自己构建论据；而另一些人，沉溺于享乐，按自己的私欲曲解圣经。按我的看法，热爱真理的人需要灵魂的力量。因为那些从事最伟大尝试的人，在最要紧的事上必然会失败；除非他们从真理本身接受真理的准则，并紧抓住真理。但这样的人，因偏离了正道，在大多数细节上犯了错误；这正如你所预料的，因为他们没有经过严格训练以辨别真伪、拣选要点的能力。若他们有这能力，他们就会服从圣经。\par

因此，正如一个人若像那些被喀耳刻下药的人一样，变成了畜生；同样，凡轻蔑教会传统、奔向异端者见解的人，就不再是天主的人，也不再对主保持忠信。但若有人从这欺骗中回头，听见圣经，并将其生命转向真理，他就好像从人变成了神。\par

因为我们拥有教导的根源——主，他借着先知、福音和蒙福的宗徒，\Quote{在古时多次以多种方式}\BibleRef{希 1:1}，从知识的开端引领到终末。但如果有人假定需要另一个根源，那么，根源就不再真正得以保存。\par

因此，那亲自相信圣经和主的声音（这声音通过主而运作，以造福人）的人，就正确地被认为是忠信的。确实，我们用它作为发现事物的标准。凡受批判的东西，在受批判之前是不被相信的；因此，需要批判的东西不能成为第一原理。所以，合理地，我们借着信德把握这不可证明的第一原理，并从第一原理本身丰富地接受关于第一原理的证明，我们就被主的声音培养到认识真理。\par

因为我们不可仅凭人的陈述就赞同他们，因为人同样可以陈述相反的事。但如果仅仅陈述意见不够，而必须确认所陈述的内容，我们就不能等待人的见证，而是借着主的声音来确定所讨论的事，这声音是所有证明中最确实的，甚至可以说是唯一的证明；在这知识中，那些仅仅尝过圣经的是信徒；而那些更进一步、成为真理的正确诠释者的人，则是\OriginalTerm{识者}{Gnostics}。同样，在关乎生活的事上，工匠胜于普通人，并塑造出超越一般观念的事物；因此，我们也同样，通过从圣经本身对圣经作完整的展示，借着信德以证明来说服人。\par

如果那些跟随异端的人也敢利用先知性的圣经；首先，他们不会使用全部圣经，其次，他们不会完整引用，也不会按照预言的整体和结构来引用。而是挑选含糊的措辞，将它们曲解为符合自己的意见，这里那里收集一些言辞；不关注含义，只利用字面。因为在他们几乎所有的引文中，你会发现他们只注意名称，而改变意义；既不认识（按他们所说），也不按照引文的真实本性来使用他们所引用的经文。\par

但真理并非通过改变含义而找到（因为这样人们就颠覆了一切真正的教导），而是通过思考什么完全属于并适合那位至高天主，以及通过从相似的圣经再次确立圣经中所证明的每一点。因此，他们既不愿转向真理，因羞于放弃自我之爱的主张；也不能管理他们的意见，因对圣经施加暴力。但他们首先向人散布了虚假的教义；公然与几乎全部圣经为敌，并不断被我们这些反对他们的人所驳斥；至于其余，甚至现在他们部分地坚持反对承认先知性的圣经，部分地贬低我们，说我们本性不同，无法理解他们所特有的东西。有时，当他们的教义被驳倒时，他们甚至否认自己的教义，羞于公开承认他们在私下里所夸耀的教导。这可以在所有异端中看到，当你检视他们教义的不义时。因为当我们清楚地表明他们与圣经相悖而推翻他们时，可以看到那些捍卫教义的人做了两件事之一：他们要么轻视自己教义的一致性，要么轻视预言本身，或者更确切地说，轻视自己的希望。而且他们总是更倾向于在他们看来更明显的东西，而不倾向于主借着先知和福音所说的，并且由宗徒们所见证和确认的话。\par

因此，看到他们所处的危险（不是关于一个教义，而是关于异端的维护）——无法发现真理；因为他们虽然手头有书可读，却视之为无用而轻视，但在渴望超越普通信德的过程中，他们偏离了真理。由于没有学习教会知识的奥秘，也没有能力领略真理的伟大，懒于深究事物的根源，只作表面阅读，他们就丢弃了圣经。然后，被虚浮的意见所抬举，他们不停地争吵，显然更关心\Emph{显得}是哲学家，而不关心\Emph{是}哲学家。他们没有以事物的必要第一原理为基础；受人的意见影响，然后强迫地使目的适应自己；由于被驳斥，他们与那些从事真正哲学研究的人拳斗，忍受一切，并且如他们所说，使尽浑身解数，甚至走向不虔敬，不信圣经，也不愿被剥夺异端的荣誉和他们在教会中自夸的首位；因此他们也急切地拥抱那被称为阿加佩的宴席中荣誉的躺椅，这被称为阿加佩实是虚名。\par

从我们已有的信仰中得到的真理知识，产生了对尚未相信之事的信德；这信德可以说是论证的本质。但显然，没有一种异端有耳朵听有益的事，只对引人快乐之事敞开。因为如果他们中有人愿意服从真理，他就会得到治愈。\par

然而，治愈自负（如同每一种疾病）有三方面：查明病因及去除的方法；第三，训练灵魂，使其习惯于对所达成的判断采取正确态度。因为，就像有毛病的眼睛一样，被非自然的教义蒙蔽的灵魂也不能清晰地感知真理之光，甚至对眼前之物视而不见。\par

他们说，在浑水中鳗鱼因被弄瞎眼睛而被捕。正如顽劣的男孩把老师挡在外面，这些人也将预言从他们的教会中排除，因着指责和劝诫而对其持怀疑态度。事实上，他们拼凑了大量谎言和虚构，以便显得他们不接纳圣经是合乎理性的。因此，他们并不虔诚，因为他们不喜欢神圣的命令，也就是圣神。正如那些内核无价值的杏仁被称为空杏仁，而非那些里面什么都没有的；同样，我们也称那些缺乏天主的计划和基督的传统的人为空虚的异端；实际上，苦如野杏仁，他们的教义源自他们自己，除了那些因证据确凿而不能丢弃和隐藏的真理之外。\par

因此，正如在战争中士兵不得离开指挥官所指定的岗位，我们也不得离开圣言所指定的岗位，我们已接受祂为知识和生命的导师。但大多数人甚至没有查问过，是否有一位我们应该追随的，以及祂是谁，以及应如何追随祂。因为正如圣言那样，信徒的生活也必须如此，以便能够追随天主，祂从开始就按正确的路径引导一切事物达到终点。\par

但当人违抗了圣言，从而违抗了天主；如果是由于突然受到某种印象而变得无力，他应当确保手头有理性的印象。如果他是因被以前掌控他的习惯所胜而变得\Quote{世俗}，正如圣经所说，那么必须彻底停止这些习惯，并训练灵魂反对它们。如果看来冲突的教义使一些人偏离，必须移除这些教义，并求助于那些调和教义的人，通过新旧约的联系来解释真理，以圣经的魅力驯服那些胆怯的未受过教导者。\par

但显然，我们倾向于基于意见的观念，即使它们与真理相悖。因为真理是严峻而庄重的。现在，灵魂有三种状态——无知、意见、知识——处于无知的是外邦人，处于知识的是真教会，处于意见的是异端。那么，再没有比这更清楚的了：那些有知识的人对他们所知道的事做出肯定，而其他人则对他们基于意见所持的观点做出肯定，就未经证实的肯定而言。\par

因此他们彼此轻蔑嘲笑。而且常常发生同一思想被某些人极为推崇，而被另一些人斥为疯狂。事实上，我们知道，应归于外邦人的纵欲是一回事；在异端教派中更受青睐的口角是另一回事；应属于教会的喜乐是另一回事；应归于真诺斯替者的欢愉是另一回事。正如一个人若投身于\Person{依斯科玛科斯}，他会使他成为农夫；投身于\Person{兰匹斯}，成为水手；投身于\Person{卡里德摩斯}，成为军事指挥官；投身于\Person{西蒙}，成为骑手；投身于\Person{佩尔狄克斯}，成为商人；投身于\Person{克罗比鲁斯}，成为厨师；投身于\Person{阿基劳斯}，成为舞者；投身于\Person{荷马}，成为诗人；投身于\Person{皮浪}，成为好争辩者；投身于\Person{德摩斯梯尼}，成为演说家；投身于\Person{克里西普斯}，成为辩证学家；投身于\Person{亚里士多德}，成为自然学家；投身于\Person{柏拉图}，成为哲学家：同样，听从主并跟随祂所赐予预言的人，将完美地成形为师傅的样式——成为行走在肉体中的神。\par

因此，那些不追随天主到祂所引领之处的人，就从这高位跌落。而祂在受默感的圣经中引领我们。\par

虽然人的行为成千上万，但一切罪的根源只有两个：无知和无能。而两者都取决于我们自己；因为我们不愿学习，也不愿抑制欲望。其中，前者是使人不能正确判断的原因，后者是使人不能遵行正确判断的原因。因为心神被迷惑的人不能正确行动，即使完全有能力做他所知道的事；同样，即使能判断什么是必要的，但如果缺乏行动的能力，他也不能免于责罚。因此，针对这两种罪，有两种相应的纠正方法：对于前者，是知识和来自圣经见证的清晰证明；对于后者，是按照圣言的训练，由信德和敬畏的纪律所规范。而两者都发展为完美的爱德。因为在我看来，诺斯替者的目的在此有二重性——一方面是理论性的默观，另一方面是行动。\par

那么，但愿这些异端者能通过这些笔记学习并改正，归向至高的天主！但如果他们像聋蛇一样，不听那被称为新歌、却极为古老的歌曲，愿他们受天主惩戒，在审判之前接受慈父般的训诫，直到他们羞愧悔改，而不致因轻率的无信仰而冲撞，并自坠审判。\par

有些部分的矫正被称为惩戒，许多我们中曾陷入过犯的人，因偏离了主的子民而招致这些惩戒。但正如孩童被他们的老师或父亲惩戒一样，我们也被天主眷顾所惩戒。然而天主并不惩罚，因为惩罚是对邪恶的报复。祂惩戒，却是为了那些受惩戒者的益处，无论是集体还是个人。\par

我援引这些事，是出于一种愿望，要使那些渴望学习的人免于陷入异端的危险，并想阻止他们表面的无知、愚钝或恶劣的性情，或无论该称之为什么。为了说服并引领那并非完全不可救药的人归向真理，我使用了这些话。因为有些人完全不能忍受那些劝他们转向真理的人的话；他们试图嬉闹，倾吐亵渎真理的言语，自诩通晓宇宙中最伟大的事，却未曾学习、探究、劳苦或发现思想的连贯脉络——对这样的人，与其憎恨他们的乖戾，不如怜悯。\par

但若有人是可治愈的，能像火或钢一样承受真理的直言——这直言割除并烧毁他们的虚假意见——就让他献上灵魂的耳朵。除非因懒惰的倾向而推开真理，或因贪求名声而试图发明新奇之物，否则情况便会如此。因为那些懒惰的人，本有能力从圣经本身为圣经获取适当的证明，却只选择有助于自己私欲的部分。而那些贪求名声的人，则自愿通过各种各样的论证，回避受福的宗徒和教师所传授、且与灵感之言结合的事物；以人的教导反对神圣的传统，为要建立异端。因为，说实话，还有什么可说——我是指在教会知识方面——留给这样的人呢？例如，\OriginalTerm{马尔西昂}{Marcion}、\OriginalTerm{普若迪科}{Prodicus}以及类似的人，他们没有行在正道上。因为他们不可能在智慧上超越先辈，以致发现任何超出先辈所言的东西；他们若能学到先前已阐明的事，就该满足了。\par

唯独我们的诺斯替者，既在圣经中老练成熟，又在教义上持守宗徒的与教会的正统，便最正确地按福音生活，并从法律和先知中发现他所寻求的证明（因他是被主所差遣的）。因为在我看来，诺斯替者的生命无非是与主的传统相符的言行。但\Quote{众人并非都有知识。因为弟兄们，我不愿你们不知道，}宗徒说，\Quote{众人都曾在云下，并分吃了属灵的食物和饮料；}他清楚肯定，所有听了道的人并没有在言行上领会知识的伟大。因此他又补充说：\Quote{但祂并不喜欢他们所有的人。}这是谁？祂说：\Quote{你为什么称我为主，却不遵行我父的旨意？}这就是救主的教导，它对我们而言是属灵的食物和不渴之饮，是诺斯替生命的水。此外，据说知识是\Quote{叫人自高自大。}对此我们说：或许表面的知识才叫人自高自大，如果有人认为这话的意思是\Quote{膨胀。}但若宗徒的话更可能是指\Quote{怀抱伟大而真实的思想，}那么难题就解决了。因此，让我们跟随圣经，确立所说的话：\Quote{智慧，}撒罗满说，\Quote{使她的子女膨胀。}因为主并没有通过祂教导的细节制造傲慢；而是产生对真理的信赖和心智的扩展，这扩展在于圣经所传授的知识，以及对那些引人犯罪之事的轻视——这就是\Quote{膨胀}一词的含义。它教导通过教诲而植入其子女内心的智慧之宏伟。如今宗徒说：\Quote{我要知道的，不是那些自高自大之人的言语，而是能力；}\BibleRef{格前 4:19} 若你们以宽广之心理解圣经（这意味着真实；因为没有什么比真理更伟大）。因为那正是智慧子女的能力所在，他们因之自高自大。他似乎在说：我要知道你们是否正确地怀抱关于知识的伟大思想。\Quote{因为天主，}按达味的话，\Quote{在犹大为人所知，}即那些按知识为以色列人的人。因为犹大被解释为\Quote{宣认。}因此宗徒说得对：\Quote{这\InnerQuote{不可奸淫，不可偷窃，不可贪恋；}以及别的任何诫命，都包含在这一句话里，\InnerQuote{你应当爱你的近人如你自己。}}\BibleRef{罗 13:9}\par

因为我们决不可像那些追随异端的人一样，通过满足自己的情欲和虚荣，欺骗我们的近人，来篡改真理或盗窃教会的准则；我们首先有责任在爱德中教导他们持守真理。因此经上明确说：\Quote{在万民中传扬祂的旨意，}好使他们不被定罪，而让那些先前已听闻的人得以悔改。但那些用舌头说话诡诈的人，已有记载的惩罚等待着他们。\par

% structure: p0296 source=h2 target=subsection reason=relative-heading-rank; base=h2

\subsection{第十七章 教会的传统先于异端传统}

因此，那些坚持不虔之言并将之灌输给他人的人，由于他们对神圣之言的使用不正确，而是邪僻的，他们自己既不进入天国，也不允许他们所欺骗的人获得真理。因为他们没有进入的钥匙，而是一把假的（按常言说是）伪造的钥匙（\OriginalTerm{反钥匙}{\GreekText{ἀντικλεῖς}}），他们不像我们那样通过主的传统、拉开帷幕而进入；而是从旁门突入，在教会的墙上秘密挖掘，跨越真理，把自己立为不虔者的灵魂的秘仪导师。\par

若要证明他们后来举行的集会是在天主教会之后，无须多费唇舌。\par

因为我们主在他来临时所教导的，从\Person{奥古斯都}和\Person{提庇留}开始，在\Person{提庇留}的时期中间完成。\par

宗徒们的教导，包括\Person{保禄}的职务，以\Person{尼禄}为止。后来，在\Person{哈德良}王时代，那些发明异端的人才兴起；他们延续到\Person{安东尼}长者时期，例如\Person{巴西理德}，虽然他声称（如他们所夸耀的）他的师傅是\Person{伯多禄}的翻译者\Person{格劳基亚}。\par

同样，他们声称\Person{瓦伦廷}是\Person{德奥达}的听众。而\Person{德奥达}是\Person{保禄}的门徒。因为与他们同时代的\Person{玛尔基翁}，年老时与较年轻的异端者们同住。在他之后，\Person{西满}曾短暂听过\Person{伯多禄}的宣讲。\par

既然如此，从教会的高度古老性和完全的真理，显然这些后来的异端，以及其后更晚的，都是从真理中伪造出来的新发明。\par

从以上所说，我认为真正的教会，即真正古老的教会，是唯一的，而那些按照天主计划为义的人，都登记在其中。因为既然天主是唯一的，主也是唯一的，那最受尊崇的便因其单一性而受赞美，是唯一第一原则的摹仿。因此，在唯一者的本性中，唯一的教会结为共同的产业，他们却试图把它分裂成许多教派。\par

因此，在实体和理念上，在起源、卓越性上，我们说那古老而唯一的天主教会，把那些已经预定的人——天主在创立世界以前就预知他们会成为义人——聚集到唯一信仰的统一中，这统一来自不同的盟约，或更确切地说，来自唯一的天主通过唯一的主在不同时期按他的旨意所立的唯一盟约。\par

但教会的卓越性，作为合一的原则，在于它的唯一性超越一切，没有相似或相等之物。但关于这点，容后再说。\par

在异端中，有些以一个人的名字命名，如称为\OriginalTerm{瓦伦廷派}{Valentinus}、\OriginalTerm{玛尔基翁派}{Marcion}和\OriginalTerm{巴西理德派}{Basilides}——尽管他们夸耀引述了玛窦的意见，但并不真实；因为教导如此，宗徒的传统也是唯一的。有些以地方命名，如\OriginalTerm{贝拉提克人}{Peratici}；有些以民族命名，如\OriginalTerm{弗吕家派}{Phrygians}；有些以行为命名，如\OriginalTerm{恩克拉提特派}{Encratites}；有些以特殊教义命名，如\OriginalTerm{幻身派}{Docetae}和\OriginalTerm{血质派}{Haermatites}；有些以假设和他们所尊敬的个人命名，如\OriginalTerm{加因派}{Cainists}和\OriginalTerm{奥非特派}{Ophians}；有些以邪恶的实践和暴行命名，如被称为\OriginalTerm{恩提奇特派}{Entychites}的\OriginalTerm{西满派}{Simonians}。\par

% structure: p0307 source=h2 target=subsection reason=relative-heading-rank; base=h2

\subsection{第18章 法律中洁净与不洁净动物的区别是教会、犹太人与异端者之间的区别的象征}

在给那些喜爱默观教会的人打开一扇小窗后——从祭献法律中关于洁净与不洁净的动物（因为这样，普通的犹太人和异端者被神秘地从神圣教会中区分出来）——让我们结束这篇论述。\par

因为那些分蹄又反刍的祭品，圣经视为洁净，为天主所悦纳；因为义人借着信德得以接近圣父和圣子。这就是分蹄者的稳固，他们昼夜研读天主的圣言，并在灵魂的接受教导之处反刍这些圣言；法律以洁净动物反刍的形像表达了这种诺斯替的操练。而那些不具备两者之一的，则被分开视为不洁。\par

那些反刍却不分蹄的，指示大多数犹太人——他们确有天主的圣言，却没有信德和那立足于真理、经由圣子通向圣父的阶梯。因此这类牲畜容易滑倒，因为它们的蹄没有分开，也不依靠信德的双重支持。因为经上说：\BibleQuote{除了那蒙圣子启示的人，没有人认识父。}\BibleRef{路 10:22}\par

同样，那些分蹄却不反刍的，也是不洁的。因为这些指示异端者——他们确实行走在圣父和圣子的名号上，却不能碾磨和细察圣言的清晰宣告，而且他们粗糙地、不精确地行正义的事，即使他们行的话。主对他们说：\BibleQuote{你们为什么称呼我主啊，主啊，却不遵行我所说的话呢？}\BibleRef{路 6:46}\par

那些既不反刍也不分蹄的，完全是污秽的。\par

\Quote{但是你们\Person{麦加拉人}，}\Person{特奥格尼斯}说，\Quote{既非第三，也非第四，\newline{}也非第十二，既不在计算中，也不在数目里。}\par

\BibleQuote{却像糠秕，被风吹离地面，也像从器皿滴下的水珠。}\BibleRef{依 40:15}\par

这些要点既已在前文详细论述，伦理学的大纲也按我们承诺的，以片段的方式概要勾画，并且我们在这里和那里穿插了作为真知识种子的教义，以便神圣传统的发现对任何未入教者都不容易，那么让我们转向我们承诺的内容。\par

然而，\WorkTitle{杂记}不像整齐排列、赏心悦目的部分，倒更像一座荫蔽而杂乱的丘陵，种植着月桂、常春藤、苹果、橄榄和无花果；这种植刻意混合了结果实和不结果实的树木，因为作品的目的是隐蔽，以防那些胆敢偷窃成熟果实的人；但农夫们通过移植嫩枝和植物，将装饰出一座美丽的公园和宜人的小树林。\par

因此，\WorkTitle{杂记}既不追求布局，也不追求文辞；因为甚至在希腊人中，有时也故意不让华丽的文辞出现，而暗中播下教义的种子，不是按照真理，使阅读者勤勉而敏于发现。因为各种鱼类有各种不同的诱饵。\par

现在，在我们这第七部\WorkTitle{杂集}之后，我们将从另一个开端按顺序记述接下来的内容。\par

\SourceRecord{Translated by William Wilson.；Ante-Nicene Fathers；Vol. 2.；Edited by Alexander Roberts, James Donaldson, and A. Cleveland Coxe.；Buffalo, NY: Christian Literature Publishing Co.,；1885.；<http://www.newadvent.org/fathers/02107.htm>.}

% structure: p0001 source=h1 target=section reason=page-title

\section{《杂记》（第八卷）}

% structure: p0002 source=h2 target=subsection reason=relative-heading-rank; base=h2

\subsection{第一章 哲学与神学探究的对象——真理的发现}

但是，最古老的哲学家并未陷入争辩和怀疑，更何况我们这些追随真正哲学的人，圣经吩咐我们要考察和探究。因为较晚近的希腊哲学家，由于空虚而徒劳的虚荣心，被引入无用的空谈、反驳和争吵。相反，蛮族哲学驱除了一切争论，说：\Quote{你们寻求，就会找到；你们敲，就会给你们开；你们求，就会给你们。}\par

因此，通过探究，所提议的疑问和回答就敲开了真理之门，按照所显现的。而当在探究过程中因障碍而打开通道时，就产生了科学的沉思。根据我的看法，对于如此敲门的，所探究的议题就被打开了。\par

对于那些在圣经中如此提问的，天主赐予他们（他们所追求的）天主所赐知识的恩赐，作为理解，通过逻辑探究的真正启迪。因为不先寻求，就不可能找到；不先考察，就不可能寻求；不先通过询问来展开和开启问题以产生清晰性，就不可能考察；再者，不完成整个探究，以后就不可能获得所疑问之点的知识作为奖赏。\par

但寻求者应找到，而寻求者则事先认为他自己不知道。因此，被渴望发现善的事物所吸引，他深思熟虑地寻求，不争竞、不虚荣，询问、回答，并考虑所做的陈述。因为不仅当我们研读神圣圣经时，而且当我们涉及普遍概念时，我们应当进行探究，发现止于某个有用的目的。\par

因为另一个地方和人群等待那些好斗的人以及法庭式的诡辩。但是，对于既爱真理又作真理门徒的人来说，即使在探究中也应保持和平，通过科学的论证前进，不贪爱自己，而爱真理，以达到全面的知识。\par

% structure: p0008 source=h2 target=subsection reason=relative-heading-rank; base=h2

\subsection{第二章 明确界定的必要性}

对于这种性质的教导的开始，还有什么比清晰讨论所提出的术语更好的或更清楚的方法呢？这方法能使所有使用同一语言的人都理解它。难道论证这个术语就像\OriginalTerm{Blityri}{Blityri}这个词一样吗？它只是一个声音，毫无意义。但哲学家、演说家乃至法官，怎么都不会把论证作为一个毫无意义的术语提出；也没有任何争论者不知道这个事实，即含义并不存在？\par

事实上，哲学家们以不同的方式提出论证具有实质存在。因此，如果一个人要正确地处理每个问题，他不能将论述引回到另一个比该术语所表明的意义更普遍接受的基本原则，而这个意义是所有同民族和同语言的人所承认的。\par

然后，从这一点出发，必须探究命题是否具有这个含义。其次，如果证明它具有，就必须精确地考察其本质，它属于哪种，以及它是否曾越出所指定的类别。如果仅仅说自己的想法是不够的（因为对方同样可以随意提出相反的主张），那么所陈述的必须得到确认。如果它的裁决被引回到另一个同样有争议的问题，并且那个问题的裁决又被引回到另一个有争议的问题，那么这将\OriginalTerm{ad infinitum}{ad infinitum}永无止境，将无法证明。但是，如果对一个未被承认的点的信念被引回到一个被所有人承认的点，那么这就要作为教导的起点。因此，每个供讨论的术语都要转换成一个为讨论各方所承认的表达，以形成教学的出发点，引导到所探究点的发现。例如，假设所讨论的术语是\Quote{太阳}。斯多葛派说它是\Quote{从海水点燃的理智之火}。那么，定义岂不比术语更晦涩，需要另一个证明来检验它是否真实吗？因此，最好以普通而清晰的说话方式说：\Quote{最亮的天体被称为太阳。}因为这种表达更可信、更清晰，也为所有人所承认。\par

% structure: p0012 source=h2 target=subsection reason=relative-heading-rank; base=h2

\subsection{第三章 论证的定义}

同样，所有人都会承认，论证是符合理性的论述，从被承认的点出发，在争议点上产生信念。\par

现在，不仅论证、信念和知识，而且预知，也有两种用法。一种是在科学上和确定无疑的，另一种只是基于希望。\par

那么，严格地说，那在学习者灵魂中产生科学信念的被称为论证。另一种只是导致意见。正如，真正的人，拥有通常的判断，和野蛮、兽性的人——都是人。因此喜剧诗人说：\Quote{人只要是人，就是优雅的。}牛、马、狗也是如此，根据动物的优劣。因为通过着眼于类的完善，我们来到严格意义上的含义。例如，我们设想一个在医治能力上毫无欠缺的医生，以及一个在科学知识上毫无欠缺的诺斯替者。\par

现在，证明与三段论不同；因为被证明的要点指示出某一件事，是单一且同一的；正如我们说，怀有身孕便是已非童贞的证明。但由三段论所把握的，虽是一件东西，却从若干前提推论出来；例如，关于\Person{皮托}背叛了\Place{拜占庭人}（如果事实确是如此），所举出的不止一个证据。从公认的前提推出结论叫做\Emph{进行三段论推理}；而从真实前提推出结论叫做\Emph{证明}。\par

可见证明有一种复合的优越性：一方面它为了证明所争论的要点而假定真实的前提，另一方面它推导出从前提而来的结论。如果第一个条件不存在，而第二个条件从第一个条件而来，那么一个人就不是证明，而只是进行了三段论推理。因为从前提推出适当的结论，不过是进行三段论推理。但要使每个前提都是真实的，则不仅是进行了三段论推理，而且也是进行了证明。\par

而得出结论（如此词所示）就是将推理引向终结。在每个推理过程中，所要确定的目标就是终点，也称为结论。但任何简单而基本的陈述（即使是真实的）都不被称为三段论；三段论至少由三个这样的陈述复合而成——两个作为前提，一个作为结论。\par

现在，要么一切事物都需要证明，要么其中有些是自明的。如果是前者，那么因要求每个证明的证明，我们将无穷递进；于是证明便被推翻了。如果是后者，那么这些自明的事物将成为证明的起点[和基本根据]。\par

事实上，哲学家们承认万物的第一原理是不可证明的。因此，如果存在证明，就绝对必须存在着某种自明的事物，这被称为首要且不可证明的。\par

因此，一切证明都追溯到不可证明的信德。\par

而且，在源于信德的源头之后，还会有其他的证明起点——即那些清楚呈现在感官和理智面前的事物。因为感官的现象是单纯的，不能被分解；而理智的现象是单纯的、理性的、基本的。但由它们所产生的事物则是复合的，却同样清楚可靠，并且比前者更多涉及推理能力。因此，我们一切人自然具备的理性特有的本土能力，处理的是相符与不相符。那么，如果发现某个论证属于这样一种类型，即能够从已经相信的要点出发，对尚未相信之事产生信念，我们将断言这正是证明的本质。\par

现在，他们断言证明的本性（正如信念的本性一样）有两种：一种是在听者的灵魂中仅仅产生说服，另一种是产生知识。\par

那么，如果一个人从感官和理智所明白的事物出发，然后推出适当的结论，他就是在真正地证明。但如果他从仅仅是可能的面非基本的事物出发（这些事物既非感官也非理智所明白的），并且他推出了正确的结论，他确实会进行三段论推理，但不会产生科学的证明；但如果他没有推出正确的结论，他就根本不是在进行三段论推理。\par

现在，证明与分析不同。因为被证明的每一点都是通过已被证明的要点来证明的；那些已被证明的要点又是由其他要点先前证明的；直到我们追溯到那些自明的、或感官和理智所明白的要点；这被称为分析。而证明则是，当所讨论的要点通过所有中间步骤到达我们这里。因此，从事证明的人应当非常关注真理，同时忽略前提的术语，不论你称它们为公理、前提或假设。同样地，还必须特别关注结论是基于哪些假设；而至于别人是否称它是一个终结性的或三段论式的命题，他完全可以不在意。\par

因为我断言，想要证明的人必须注意这两件事：假定真实的前提，并从它们推出正确的结论（有些人也称此为\Quote{推论}，因为它是从前提推断出来的）。\par

现在，在每个关于一个问题的命题中，必须有不同的前提，但这些前提必须与所提出的命题相关；而且所提出的东西必须还原为定义。这个定义必须得到所有人的承认。但如果假定了与要建立的命题无关的前提，就不可能得出任何正确的结果；因为整个命题（也称为其本性的问题）被忽略了。\par

因此，在所有问题中，都有某种先前已知的东西——那就是自明的、无需证明即可相信的东西；它必须被用作探究的起点，以及明显结果的判据。\par

% structure: p0029 source=h2 target=subsection reason=relative-heading-rank; base=h2

\subsection{第四章　为避免含糊，我们必须从清晰的定义开始}

因为每个问题都是从已有知识出发来解决的。而对每个研究对象的已有知识，有时仅仅是关于本质的，而其功能却未知（例如石头、植物和动物，我们不知道它们的运作），或者是关于属性、能力，或者（可以这样说）关于对象本身所固有的品质。有时我们可能知道这些能力或属性中的某一个或多个（例如灵魂的欲望和情感），却不知道本质，并将其作为研究对象。但在许多情况下，我们的理智假定了所有这一切，问题就在于，它们（即这些属性等）内在于哪一种本质之中；因为我们在心中形成了对两者（即本质和运作）的概念之后，才着手处理问题。还有一些对象，我们虽然知道它们的运作及其本质，却不知道它们的变型。\par

此即发现\OriginalTerm{真理}{truth}的方法。因为我们必须从认识待讨论的问题开始。因为表达的形式常常欺骗、混淆和扰乱心智，以致不易发现事物应归于何类；例如，\OriginalTerm{胎儿}{fœtus}是否为\OriginalTerm{动物}{animal}。因为我们对动物和胎儿拥有一个概念，便询问胎儿是否是动物，亦即，处于胎儿状态的主体是否具有运动能力，此外还有感觉能力。因此，探究涉及已知主体中的功能和感觉。所以，应首先询问提出问题的人，他称什么为动物。当我们发现同一术语用于多种目的时，尤其应当这样做；并且我们必须考察，该术语所指的含义是有争议的，还是众所公认的。因为若有人说，他把一切生长并受养育的东西称为动物，我们就必须再次进一步询问，他是否认为植物是动物；然后，在他表明这一立场后，他必须说明，处于胎儿状态并被养育的东西是什么。\par

因为\Person{柏拉图}称植物为动物，认为它们仅参与第三种生命，即欲望的生命。但\Person{亚里士多德}虽然认为植物具有生长和营养的生命，却不认为适合称它们为动物；因为他认为，只有拥有另一种生命——感觉生命——的东西才可合理地称为动物。\Person{斯多葛学派}则不称营养能力为生命。\par

现在，若提出问题的人否认植物是动物，我们将表明他作出了自相矛盾的论断。因为，他以营养和生长来定义动物，却又断言植物不是动物，这无异于说：受到营养和生长的事物既是动物又不是动物。\par

那么，让他说出他想要了解什么。是他子宮内的东西是否生长并受养育？还是它是否具有任何感觉或由冲动引起的运动？因为，按照\Person{柏拉图}，植物是有灵的，且是动物；但按照\Person{亚里士多德}，它不是动物，因为它缺少感觉，却是有灵的。因此，根据他，动物是一个有灵的具有感觉的存在。但按照\Person{斯多葛学派}，植物既不是有灵的，也不是动物；因为动物是有灵的存在。那么，如果动物是有灵的，而生命是有感觉的本性，那么显然，有灵的就是有感觉的。因此，若提出问题的人再次被反问，他是否仍因胎儿受到营养和生长而称其为动物，他便得到了答案。\par

但假如他说，他所问的问题是：胎儿是否已经有感觉，或能否因任何冲动而自行运动，那么对这一问题的调查就清楚了，名称中的谬误不再存在。但如果他不回答追问，也不愿说出他的意图，或他提出问题时使用\Quote{动物}一词是基于何种考虑，反而要求我们自行定义，那么就应将他视为好辩之人。\par

既然有两种方法，一是问答法，二是阐述法，如果他拒绝前者，就让他听我们阐述与问题相关的一切。然后当我们完成后，他可以对每一点依次处理。但如果他试图通过提问打断探究，那他显然是不想听。\par

但若他选择回答，就应首先问他：他将\OriginalTerm{动物}{animal}这个名称用于什么事物。当他回答后，再问他：他如何理解胎儿，是指子宮内的东西，还是已经成形并活着的东西？再则，胎儿是指所播下的种子，还是只有在肢体和形状形成后才应使用\OriginalTerm{胚胎}{embryo}这个名称？在他回答之后，就应按照适当顺序对当前问题进行推理、得出结论并予以说明。\par

但若他希望我们说话而他不回答，就让他听。既然你不愿说出你以何种意义提出你所陈述的问题（否则我就不会这样讨论含义，而会直接考察事物本身），要知道，你的所作所为就像你提出了这样一个问题：一只狗是否是动物？因为我本可正当地问：你说的是哪一只狗？因为我要谈论陆地上的狗、海中的狗、天上的星座，还有第欧根尼，以及其他所有狗，依次排列。因为我无法猜出你是在询问所有狗还是某一只。你随后要做的是现在学习，并清楚地说出你的问题是什么。如果你在名称上闪烁其词，那么人人清楚，\Emph{胎儿}这个名称既不是动物也不是植物，而是一个名称、一个声音、一个物体、一个存在，以及任何东西和一切东西，唯独不是动物。如果你提出的正是这个问题，那么你已经得到了回答。\par

但\Emph{胎儿}这个名称所指代的东西也不是动物。因为那个东西是非物质的，可以被称为一个事物、一个概念，以及任何东西，唯独不是动物。动物的本性是不同的。因为关于该问题本身——我指的是胚胎的本性如何——已被清楚地说明。而关于\OriginalTerm{动物}{animal}这个名称所表达的含义的问题则是另一回事。\par

那么我说，如果你肯定动物是具有感觉能力和由欲望驱动自身运动的东西，那么动物就不仅仅是因欲望而运动并拥有感觉的东西。因为它还能睡眠，或者当感觉对象不存在时不运用感觉能力。但欲望或感觉的自然能力是动物的标志。因为通过这些事物指示了某种这样的性质。第一，胎儿是否不具有感觉能力或由欲望引起的运动；这是所提出的考虑点。另一点是：胎儿是否可能有时运用感觉能力或通过欲望运动。在这个意义上，无人质疑这一点，因为这是显而易见的。\par

但问题在于，胚胎已是一个动物，还是仍是一株植物。然后为了清晰起见，将\Quote{动物}一词归入定义。但在发现它凭借感觉和因欲求而运动而与不是动物的东西相区别之后，我们再次将其与附属物分离；断言：一方面，那种尚未拥有感觉和运动能力但将来会拥有的东西，只是潜在地如此；另一方面，已经现实地如此的是另一回事；而对于后者，发挥其能力是一回事，能够发挥但处于静止或睡眠状态是另一回事。这就是问题所在。\par

因为胚胎不可因它被滋养的事实而被称作动物；这是那些偏离问题本质、将心思转向其他情况之人的主张。但在所有据说被发现的结论中，都普遍应用了证明，证明就是\OriginalTerm{话语}{\GreekText{λ}{o}\GreekText{γος}}，从其他事物确立一事物。但用以确立所问之点的根据，必须为学习者所承认和知晓。而这一切的基础，就是感官和理智所明见的事物。\par

因此，首要的证明由所有这些构成。但那种由此已证明之点推出另一结论的证明，其可靠性并不亚于前者。它不能被称为首要的，因为其结论并非从作为前提的首要原理推导出来。\par

因此，三种不同类问题的第一种已被展示——我指的是，在已知本质的情况下，其能力或属性中的某一种未知。第二种命题是，我们都知晓其能力与属性，但不知道其本质；例如，灵魂的主要官能位于身体的哪一部分。\par

% structure: p0045 source=h2 target=subsection reason=relative-heading-rank; base=h2

\subsection{第五章 将证明应用于怀疑论的判断悬置}

现在，同样适用于证明的处理也适用于以下问题。\par

例如，有人说一个动物不能有多个起源因。一个动物不可能有多个同质的起源因；但可以有多个异质的，这并不荒谬。\par

假设皮浪式的判断悬置，如他们所说，[观念]即没有什么是确定的：显然，它始于自身，首先否定自身。它要么承认某些事是真的，即你不必对一切事物悬置判断；要么坚持说没有什么是真的。而很明显，首先它不会是真的。因为它要么肯定什么是真的，要么不肯定什么是真的。但如果它肯定什么是真的，它就不情愿地承认了某些事是真的。如果它不肯定什么是真的，它就让那些它本想消除的东西仍然为真。因为，就那种破坏性的怀疑论被证明为假而言，被破坏的立场就被证明为真；就像那个说所有梦都是假的梦。因为通过驳斥自身，它证实了其他立场。\par

最后，如果它是真的，它就会从自身开始，并且首先不是对别的事物的怀疑，而是对自身的怀疑。那么，如果[这样一个人]领会到他是一个人，或者他是怀疑的，显然他就不是怀疑的。他又将如何回答询问呢？因为显然他在这一点上不是怀疑论者。不，他甚至肯定他确实在怀疑。\par

而且，如果我们必须被说服对一切事物悬置判断，我们首先就应对我们的判断悬置本身悬置判断，无论我们是否相信它。\par

而如果这个立场——即我们不知道什么是真的——是真的，那么它就不允许任何东西为真。但如果他说，甚至这一点也是可质疑的，即我们是否知道什么是真的；通过这个陈述本身，他承认真理是可认知的，就在那看似建立对真理的怀疑的行动中。\par

但如果一个哲学学派是一种对教义的倾向，或按某些人的说法，是对一系列彼此一致且与现象一致、导向正确生活的教义的倾向；而教义是一个逻辑概念，概念是心灵的一种状态和同意：那么不仅怀疑论者，而且每个制定教义的人都习惯于在某些事情上悬置判断，或是由于心智力量的缺乏，或是由于事物不够清晰，或是由于理由势力相当。\par

% structure: p0053 source=h2 target=subsection reason=relative-heading-rank; base=h2

\subsection{第六章 定义、属与种}

问题的导言和来源是关于这些点并且在它们之中。\par

但在定义、证明和划分之前，必须说明问题是以何种方式陈述的；并且要处理歧义术语；同义词要根据其含义准确陈述。\par

然后应探究，这个命题是属于那些相对于其他事物而被考虑的要点，还是被独立看待。进一步，它是否存在、它是什么、它有什么属性；或者同样地，它是否存在、它是什么、它为什么存在。而对这些要点的考察，有赖于对个体与普遍、先行与差异及其划分的知识。\par

现在，归纳旨在概括和定义；而划分则涉及种、事物之所是以及个体。对\Quote{如何}的沉思引出对独特之物的假定；而怀疑带来具体的差异与证明，并在其他方面增强推理及其结果；整体的结果便是科学知识与真理。\par

再次，由划分得出的总括成为定义。因为定义在划分之前和之后都被采用：之前，当它被承认或陈述时；之后，当它被证明时。而通过感觉，普遍从个体被总括出来。因为归纳的起点是感觉；终点是普遍。\par

因此，归纳不显示\Emph{事物之所是}，而显示\Emph{它是否存在}。划分显示它是什么；定义与划分类似，教导本质和事物之所是，但不显示它是否存在；而证明解释三个要点：它是否存在、它是什么以及它为什么存在。\par

也有包含原因的界定。既然当我们见到原因时，我们才能知道；而原因有四——质料、动力因、形式、目的；那么界定也将是四重的。\par

因此，我们必先取属，其中包含最接近上位的点；然后取下一个种差。种差的连续切割与划分，便完成了\Quote{是什么}。没有必要表达每一事物的所有种差，只表达那些构成种的即可。\par

几何分析与综合类似于逻辑划分与界定；通过划分，我们回到更简单、更初级的要素。因此，我们把所考虑对象的属划分为它所包含的种；例如，就人而言，我们把动物这个属划分为它所显现的种，即有死的和不死的。如此，通过不断把看似复合的属划分为更简单的种，我们到达研究的主体，即不能再进一步划分之处。\par

因为，在把\Quote{动物}划分为有死的和不死的后，再分为陆生的和水生的；陆生动物又分为飞行的和行走的；如此划分最接近所求对象的种，它也包含所寻求的，我们便通过划分达到最简单的种，它不包含别的，只包含所寻求者。\par

因为，我们又把行走的分为有理性的和无理性的；然后从通过划分把握的种中选取那些最接近人的，并把这些组合成一个公式，就成为人的定义：人是有理性的、有死的、陆生的、行走的动物。\par

因此，划分提供了质料之类，为界定寻求名称的单纯性；而匠人和制造者的界定，通过组合与构造，呈现事物的知识本身，而非那些我们拥有一般概念之物的知识。对这些概念，我们说解释性表述属于它们。因为划分也适用于这些概念。\par

有一种划分把被划分者分为种，如属；另一种分为部分，如整体；另一种分为依附属性。\par

因此，整体分为部分，多半是根据大小来设想的；分为依附属性则永远不能完全阐明，如果本质必然内在存在于每个存在者之中的话。\par

因此，这两种划分都应被拒绝，只有属划分为种才被认可，藉此属中的同一性被表征，而存在于种差中的多样性也被表征。\par

种总是在部分中被设想。然而，反之，如果某物是另一物的部分，它并不也是种。因为手是人的部分，但并非种。属存在于种中，因为（属）既在人中也在牛中。但整体不在部分中，因为人不在他的脚中。因此，种比部分更重要；凡是述说属的，也都全部述说种。\par

因此，最好把属划分为两个种，如若不能，则三个。这样划分的种更为属别地，由同一与差异来表征。然后被划分后，它们由属别指明的点来表征。\par

因为每一类种要么是实体，如当我们说有些实体是有形的，有些是无形的；要么是数量，要么是关系，要么是地点，要么是时间，要么是动作，要么是承受。\par

所以，一个人将给出他所拥有知识的任何事物的定义；因为若不能以言语把握并界定它，便绝不可能熟悉它。由于对定义的无知，结果产生许多争论和欺骗。因为，如果知道事物的人心中有关于它的知识，并能用言语解释他所设想的；而思想的解释就是定义；那么知道事物的人必然也能给出定义。\par

在定义中，种差被采用，它在定义中占据标记的位置。因此，把笑的能力加到人的定义中，便构成了整体——有理性的、有死的、陆生的、行走的、会笑的动物。因为以种差的方式加到定义中的东西是事物特性的标记，但并不显示事物自身的本性。他们说，种差是对特有属性的指定；既然拥有种差的事物区别于所有其余，那么只属于它且可与之换位述说的事物，在定义中必须由最初的属作为主要和基础来采用。\par

因此，在较大的定义中，所发现的种的数目在十个范畴之内；而在最小的定义中，采用最接近种的主要点，标记事物的本质与本性。但最小的定义由三项组成：属和两个本质上必要的种。这样做是为了简洁。\par

因此，我们说：人是会笑的动物。我们应取主要发生于被界定者身上的东西，或其特有的德性，或其特有的功能，等等。\par

因此，尽管定义解释事物的本质，却不能精确地把握其本性。藉由主要种，定义对本质作出阐述，并几乎在性质中拥有本质。\par

% structure: p0077 source=h2 target=subsection reason=relative-heading-rank; base=h2

\subsection{第七章　论怀疑或同意的原因}

产生怀疑的原因主要有两点：一是人类心智的变动与不稳定，其本性易于产生分歧，或人与人之间，或人自身；二是事物本身的不一致性，这正如所料，易于产生怀疑。\par

因为，既不能相信所有观点（因它们相互冲突），也不能不信所有观点（因为说一切皆不可信的那个论断本身也被包括在那些不可信之中），也不能因均势而相信一些、不信另一些，我们便被引向怀疑。\par

但在怀疑论的主要成因中，包括心灵的不稳定，它导致分歧。而分歧是怀疑的近因。因此，生活中充满了法庭和议会；并且，最终，对所谓善与恶的选择，这些都是心灵处于怀疑、因软弱而在矛盾之事上踌躇不定的标记。此外，图书馆里充满了书籍，以及那些在教义上相异、自信自己知晓事物中真理者的汇编和论著。\par

% structure: p0081 source=h2 target=subsection reason=relative-heading-rank; base=h2

\subsection{第八章 分类事物与名称的方法}

在语言中有三样东西：名称，它们主要是观念的符号，从而也是主体的符号。第二，观念，它们是主体的肖象和印象。因此，在一切人中，观念是相同的；因为主体在所有人中产生相同的印象。但名称则不然，由于语言的差异。第三，主体事物，观念借此印在我们身上。\par

名称通过文法被还原为二十四个一般元素；因为元素必须确定。因为对个别事物没有科学知识，因它们是无限的。但科学的特性在于基于一般和确定的原则。因此，个别事物也被分解为普遍物。哲学研究涉及观念和实在的主体。但由于这些中的个别事物是无限的，已经发现了一些元素，每个研究对象都被归入其中；如果它被证明进入任何一个或多个元素，我们就证明它存在；但如果它避开所有元素，则证明它不存在。\par

在陈述中，有些是不连贯的陈述；例如，\Quote{人}和\Quote{跑}，以及任何不构成一个完整句子的东西，它或真或假。而在连贯的陈述中，有些指示\OriginalTerm{本质}{essence}，有些指示\OriginalTerm{性质}{quality}，有些指示\OriginalTerm{数量}{quantity}，有些指示\OriginalTerm{关系}{relation}，有些指示\OriginalTerm{地点}{where}，有些指示\OriginalTerm{时间}{when}，有些指示\OriginalTerm{位置}{position}，有些指示\OriginalTerm{状态}{possession}，有些指示\OriginalTerm{动作}{action}，有些指示\OriginalTerm{承受}{suffering}，我们称这些为质料事物的元素，紧随第一原理之后。因为这些能被理性所沉思。\par

但非物质的事物能单独通过理智，以直接的应用而被把握。\par

在归于十个范畴的事物中，有些是自谓的（如九个范畴），有些是相对于某物而言的。\par

此外，在这十个范畴所包含的事物中，有些是\OriginalTerm{单义}{Univocal}的，例如牛和人，就各自是动物而言。因为单义词是那些共有一个共同名称——动物——以及同一原理，即定义，即有灵本质的词。而\OriginalTerm{异名}{Heteronyms}是那些以不同名称指称同一主体的词，如上和下的道路；因为道路无论是向上还是向下都是同一条。异名的另一种，如马和黑色，是那些彼此名称和定义不同，且不具有同一主体的词。这些应被称为不同，而非异名。而\OriginalTerm{多名}{Polyonyms}是那些具有相同定义但不同名称的词，如军刀、剑、弯刀。\OriginalTerm{派生词}{Paronyms}是那些从不同事物得名的词，如\Quote{有男子气概的}来自\Quote{男子气概}。\par

\OriginalTerm{歧义词}{Equivocal}具有相同名称但不同定义，如人——既指动物也指画像。在歧义词中，有些是\Quote{偶然}得名的，如埃阿斯，洛克里斯人和萨拉米斯人；有些是出于意图；其中，有些出于\Quote{相似}，如活着的人和画中的人；有些出于\Quote{类比}，如伊达山的脚和我们脚，因为它们都在下方；有些出于动作，如船的脚（即龙骨），船借此航行，和我们的脚，我们借此移动。歧义词从同一事物出发并指向同一事物；如书和手术刀被称为外科的，既从使用它们的外科医生，也从外科事物本身而言。\par

% structure: p0089 source=h2 target=subsection reason=relative-heading-rank; base=h2

\subsection{第九章 论不同种类的原因}

在原因中，有些是\OriginalTerm{启动因}{Procatarctic}，有些是\OriginalTerm{合因}{Synectic}，有些是\OriginalTerm{协助因}{Co\-operating}，有些是\OriginalTerm{必要条件}{sine quâ non}原因。\par

那些率先为任何事物的起源提供契机的原因，是\OriginalTerm{启动因}{Procatarctic}；如同纵欲者眼中的美貌是爱的原因；因为当他们看到美貌时，它独自产生爱慕的倾向，但并不必然。\par

\OriginalTerm{合因}{Synectic}（也是单义地自身完美的）是这样一种原因：当它独立地自身便能产生效果时。\par

现在，所有原因都可以依次在学习者的例子中显示。父亲是学习的\OriginalTerm{启动因}{Procatarctic}，教师是\OriginalTerm{合因}{Synectic}，学习者的本性是\OriginalTerm{协助因}{co\-operating cause}，而时间具有\OriginalTerm{必要条件}{Cause sine quâ non}的关系。\par

真正被称为原因的是能够主动产生某种效果的东西；因为我们说钢能够切割，不仅当它正在切割时，而且当它尚未切割时也是如此。因此，\OriginalTerm{产生能力}{\GreekText{τὸ} \GreekText{παρεκτικόν}}表示两者：既指现在正在作用的东西，也指尚未作用但拥有作用能力的东西。\par

因此，有些人说原因是物体的属性；另一些人说是非实体实体的属性；另一些人说，严格来说物体是原因，而非实体的东西只是借用地、准原因地是原因。另一些人则颠倒过来，说非实体实体严格来说是原因，而物体是不严格地是原因；例如，切割是一种动作，是非实体的，并且是切割动作（非实体的）的原因，而对于物体，则是被切割的原因——例如剑和被它切割的东西。\par

事物的原因以三重方式被述谓。一，什么是原因，如雕刻家；二，它成为什么的原因，即雕像；三，它对什么而言是原因，例如质料：因为他是对青铜而言成为雕像的原因。产生和被切割，作为对其所属事物而言的原因，是动作，因而是非实体的。\par

根据这一原则，原因属于谓词类（\OriginalTerm{谓词}{\GreekText{κατηγορημάτων}}），或者，如其他人所言，属于言说类（\OriginalTerm{言说}{\GreekText{λεκτῶν}}）（因为克莱安特斯和阿切德穆斯称谓词为言说）；或者更确切地说，有些原因归于谓词类，如被切割之物，其情形是被切割；有些则归于公理类——例如，建造船只，其情形是船在建造中。亚里士多德将诸如房屋、船舶、燃烧、切割这类事物的名称称为称谓。但情形被认为是无形体的。因此，以下诡辩得以解决：你所说的话经过你的口。这是真的。你提到房屋。因此房屋经过你的口。这是假的。因为我们说的不是作为形体的房屋，而是情形，房屋在其中，那是无形体的。\par

我们说建造者建造房屋，是针对将要产生的东西而言。同样，我们说斗篷被编织，因为制造者是运作的指示。制造者并非属于一物、原因属于另一物，而是属于同一物，无论斗篷还是房屋皆然。因为一个人是某物产生的原因，也同样是其制造者。因此，原因、制造者和\OriginalTerm{通过它}{\GreekText{δἰ} \GreekText{ὅ}}是同一的。因此，如果某物是\Quote{原因}和\Quote{造成者}，它当然也是\Quote{通过它}。但如果某物是\Quote{通过它}，绝不意味着它也是\Quote{原因}。例如，许多事物汇聚于同一结果，通过它们达到目的，但并非全都是原因。因为美狄亚若非愤怒，不会杀子；若非嫉妒，不会愤怒；若非爱恋，不会嫉妒；若非伊阿宋航行到科尔基斯，不会爱恋；若非阿尔戈号建造，此事不会发生；若非从佩利翁砍伐木材，此事不会发生。尽管所有这些都有\Quote{通过它}的情形，但它们并不都是杀子的\Quote{原因}，只有美狄亚是原因。因此，不阻碍者不行动。因此，不阻碍者不是原因，阻碍者才是。因为原因是在行动和做某事中被理解的。\par

此外，不阻碍者与所发生之事分离；但原因与事件相关。因此，不阻碍者不能成为原因。那么，之所以成就，是因为那能阻碍者不在场。因果关系则在四种方式中谓述：动力因，如雕塑家；质料因，如铜；形式因，如形态；目的因，如体育场长的荣誉。\par

\OriginalTerm{必要条件因}{sine quâ non}的关系由铜相对于雕像的制作所持有；同样它也是一个[真正的]原因。因为没有它结果就不可能产生的一切，必然是一个原因，但并非绝对原因。因为必要条件因并非\OriginalTerm{联结因}{Synectic}，而是协作因。每一行动者都结合受动者的适宜性产生结果。因为原因安排一切。但每一事物按其自然构造受影响；适宜性作为原因，占据着必要条件因的位置。因此，没有适宜性，原因便无效；它不是原因，而是协作因素。因为一切因果关系都在行动中被理解。大地不能使自己成为自己，所以它不能是自己的原因。如果说火不是燃烧的原因，而是木柴，或剑不是切割的原因，而是肉，或对手的力量不是运动员被击败的原因，而是其自身的软弱，那将是可笑的。\par

联结因不需要时间。因为烧灼器一接触皮肉便产生疼痛。至于\OriginalTerm{引发因}{Procatarctic}，有些需要时间直到结果产生，有些则不需要，如骨折的情况。\par

难道这些不是被称为不依赖时间，不是作为缺失，而是作为缩减，如同突然之事，而非非时间中发生之事？\par

每一个原因，被心智领悟为原因，都专注于某物，并关联于某物而被理解；即某个结果，如剑用于切割；以及某个对象，如具有适宜性，如火之于木材。因为它不会烧钢铁。原因属于与某物相关的事物。因为它是在与另一物的关系中被理解的。因此我们专注于两者，以便将原因理解为原因。\par

同样的关系也存在于创造者、制造者和父亲中。一物不能是自己的原因。也没有人是自己的父亲。因为那样第一个就会成为第二个。原因行动并影响。原因所产生的东西被行动并受影响。但同一事物就其本身而言不能既行动又受影响，也不能既是儿子又是父亲。否则，原因在存在上先于其结果，如剑先于切割。同一事物不能在同一个瞬间既在质料上在先，作为原因，同时又是在后和后续的，作为结果。\par

现在，\Quote{存在}与\Quote{生成}不同，正如原因与结果、父亲与儿子不同。因为同一事物不能在同一个瞬间既\Quote{存在}又\Quote{生成}；因此它不能是自己的原因。事物不是彼此的原因，而是为彼此的原因。因为先前的脾脏病变不是发热的原因，而是发热发生的原因；先前的发热不是脾脏的原因，而是病变加剧的原因。\par

同样，德性相互成为原因，因为由于它们的相互对应，它们不能被分开。拱门中的石头是保持这种范畴的原因，但它们不是彼此的原因。教师和学生在谓词方面互为进步的原因。\par

相互的及交互的原因被述谓，有些是同一事物的原因，如商人与零售商是利润的原因；有时一个是某一事物的原因，另一个是另一事物的原因，如剑与肉：剑是肉被割的原因，肉是剑割的原因。［说得好：］\BibleQuote{以眼还眼，以命还命。}因为凡使人受致命伤者，对他而言是死亡或死亡发生的原因。但当他反过来被人致命伤时，他成为对方的原因，不是就对方作为原因而言，而是在另一方面的原因。因为他成为对方死亡的原因，并非死亡还击了致命的一击，而是受伤者本人。因此，他是一个事物的原因，而又另有原因。行恶之人成为他人（即被冤枉之人）的原因。但命令施加惩罚的法律，并非伤害的原因，而是对一人为报应，对另一人为惩戒。所以，作为原因的事物，并非互为原因。\par

仍然有人问：是否许多事物联合成为一事物的许多原因？因为一起拖船的人是船被拖下来的原因；但与其他人的合作中，除非联合原因也是原因。\par

另有人说，如果有许多原因，每个本身成为一事物的原因。例如，诸德性（它们是许多的）是幸福（它是单一的）的原因；同样，温暖与痛苦有许多原因。那么，许多德性在能力上不是一吗？温暖与痛苦的根源不也是如此吗？诸德性的众多，因其类上为一，难道不成为单一结果（幸福）的原因吗？\par

但事实上，\OriginalTerm{肇因}{Procatarctic cause}在属与种上都不止一个；例如，寒冷、虚弱、疲劳、消化不良、醉酒，就属而言是任何疾病的原因；就种而言是热病的原因。而\OriginalTerm{总因}{Synectic cause}则仅就属而言是如此，而非就种而言。\par

因为对于悦人的气味（它在属上是单一的），有许多种的原因，如乳香、玫瑰、番红花、苏合香、没药、膏油。因为玫瑰的甜香与没药并非同一种。\par

同一事物成为相反结果的原因：有时由于原因的巨大及其能力，有时由于受作用者的感受性。根据力的性质，同一根琴弦，因其张弛，发出尖细或低沉的声音。蜂蜜对健康者是甜的，对发烧者是苦的，取决于受作用者的感受状态。同一种酒使一些人狂怒，使另一些人欢愉。同一太阳融化蜡而硬化黏土。\par

再者，原因中有的是明显的，有的是靠推理把握的，有的是隐藏的，有的是通过类比推断的。\par

在隐藏的原因中，有些是暂时隐藏的，一时隐蔽，另一时又清楚可见；有些则本性隐藏，永远不能显现。在本性隐藏的原因中，有些可以被领会；这些有些人不会称为隐藏，因为它们通过标记以类比而被领会，例如感官通道的对称性，由理性所默观。有些则不能被领会，它们绝不可能以任何方式落入领会之中，按其定义即是隐藏的。\par

现在，有些是\OriginalTerm{肇因}{Procatarctic cause}，有些是\OriginalTerm{总因}{Synectic cause}，有些是\OriginalTerm{联合原因}{Joint-cause}，有些是\OriginalTerm{合作原因}{Co-operating cause}。有些符合本性，有些超乎本性。有些是疾病的原因，有些是偶然的原因，有些是感觉的原因，有些是这些的强度的原因，有些是时间和季节的原因。\par

\OriginalTerm{肇因}{Procatarctic cause}被移除后，效果仍然存在。但\OriginalTerm{总因}{Synectic cause}是：当其存在时，效果存在；当其被移除时，效果也移除。\par

\OriginalTerm{总因}{Synectic cause}也被同义词称为\Quote{本身圆满}。因为它本身足以产生效果。\par

如果原因显出一种本身充足的运作，那么\OriginalTerm{合作原因}{Co-operating cause}表明与另一方一起的辅助和帮助。因此，如果它不产生任何效果，它甚至不会被称作合作原因；如果它确实产生某种效果，那它完全就是这效果的原因，即它所产生的东西的原因。所以，合作原因就是：当其存在时，效果产生了——明显的明显地，隐藏地隐藏地。\par

\OriginalTerm{联合原因}{Joint-cause}也属于原因的类，如同伴是士兵，如同龄人是青年。\par

\OriginalTerm{合作原因}{Co-operating cause}进一步辅助\OriginalTerm{总因}{Synectic cause}，以加强其所产生的结果的方式。但\OriginalTerm{联合原因}{Joint-cause}不属于同一概念。因为一个东西可以是联合原因，即使它不是总因。因为联合原因被理解为与另一个原因结合，而另一个原因本身不能产生效果，它是与原因一起的原因。合作原因与联合原因的区别在于：联合原因产生效果时，所作用的对象本身不行动；而合作原因本身不产生任何效果，但它加入本身行动的原因，与之合作，以便最强烈地产生效果。尤其是，从\OriginalTerm{肇因}{Procatarctic cause}变成合作原因者，在加强原因的力量方面尤为有效。\par

\SourceRecord{Translated by William Wilson.；Ante-Nicene Fathers；Vol. 2.；Edited by Alexander Roberts, James Donaldson, and A. Cleveland Coxe.；Buffalo, NY: Christian Literature Publishing Co.,；1885.；<http://www.newadvent.org/fathers/02108.htm>.}
